% !TeX program = pdflatex
% Integracja audytu II: 35 jednostek, rozdziały 1–33 + moduły 15A i 15B.
% Wersje samodzielne zachowano w katalogu wersje_samodzielne/.
\documentclass[11pt,a4paper]{article}
\usepackage[T1]{fontenc}
\usepackage[utf8]{inputenc}
\usepackage[polish]{babel}
\usepackage{lmodern}
\usepackage{amsmath,amssymb,amsthm,mathtools}
\usepackage{microtype}
\usepackage{geometry}
\usepackage{enumitem}
\usepackage{tocloft}
\usepackage{xcolor}
\usepackage{listings}
\usepackage{tikz}
\usetikzlibrary{arrows.meta}
\usepackage[unicode]{hyperref}
\geometry{margin=2.5cm}
\hypersetup{colorlinks=true,linkcolor=blue!45!black,urlcolor=blue!45!black,pdftitle={Analiza funkcjonalna - wersja scalona}}
\setlength{\parindent}{0pt}
\setlength{\parskip}{4pt}
\setlength{\emergencystretch}{2em}
\setlength{\cftsecnumwidth}{3em}
\setlength{\cftsubsecnumwidth}{3.5em}
\setlist[enumerate]{topsep=4pt,itemsep=2pt}
\setlist[itemize]{topsep=4pt,itemsep=2pt}
\lstset{basicstyle=\ttfamily\footnotesize,frame=single,breaklines=true,columns=fullflexible,showstringspaces=false}
\numberwithin{equation}{section}
\theoremstyle{plain}
\newtheorem{twierdzenie}{Twierdzenie}[section]
\newtheorem{stwierdzenie}[twierdzenie]{Stwierdzenie}
\newtheorem{lemat}[twierdzenie]{Lemat}
\newtheorem{wniosek}[twierdzenie]{Wniosek}
\theoremstyle{definition}
\newtheorem{definicja}[twierdzenie]{Definicja}
\newtheorem{przyklad}[twierdzenie]{Przykład}
\theoremstyle{remark}
\newtheorem{uwaga}[twierdzenie]{Uwaga}
\newcommand{\R}{\mathbb R}
\newcommand{\C}{\mathbb C}
\newcommand{\N}{\mathbb N}
\newcommand{\K}{\mathbb K}
\newcommand{\cl}[1]{\overline{#1}}
\newcommand{\TVS}{\mathrm{TVS}}
\DeclareMathOperator{\lin}{lin}
\DeclareMathOperator{\supp}{supp}
\DeclareMathOperator{\id}{id}
\title{Analiza funkcjonalna\\\large Przestrzenie liniowo-topologiczne, dystrybucje i geometria przestrzeni Banacha}
\author{}
\date{Wersja scalona po drugiej rundzie audytu -- 19 listopada 2026}
\begin{document}
\maketitle
\tableofcontents
\bigskip
% Wersja źródłowa: audyt II, jednostka 01
\section{Przestrzenie liniowo-topologiczne: definicje i pierwsze przykłady}\label{sec:tvs}

\subsection{Cel i konwencje}

W analizie funkcjonalnej bada się przestrzenie liniowe, w których można
wykonywać przejścia graniczne, oraz odwzorowania liniowe zgodne z tymi
przejściami. Norma jest tylko jednym ze sposobów wprowadzenia topologii.
Na przykład topologia zbieżności punktowej funkcji czy naturalna topologia
funkcji gładkich nie muszą pochodzić od pojedynczej normy. Dlatego zaczniemy
od aksjomatów, które wiążą strukturę liniową z topologiczną.

W całym skrypcie $\K$ oznacza $\R$ albo $\C$ z ich zwykłą topologią.
Przestrzenie liniowe rozpatrujemy nad $\K$. Przez otoczenie punktu $x$
rozumiemy zbiór zawierający pewien zbiór otwarty zawierający $x$;
w szczególności otoczenie nie musi być otwarte. Przez $\mathcal N_X(x)$
oznaczamy rodzinę wszystkich otoczeń $x$ w przestrzeni topologicznej $X$.
Słowo „ciągły” zawsze odnosi się do rozważanych topologii, nie zaś
do niejawnie wybranej normy.

\begin{definicja}[Przestrzeń liniowo-topologiczna]\label{def:tvs}
Przestrzenią liniowo-topologiczną nad $\K$ nazywamy przestrzeń liniową
$X$ z topologią $\tau$, dla której odwzorowania
\[
 +\colon X\times X\longrightarrow X,\qquad (x,y)\longmapsto x+y,
 \qquad
 m\colon\K\times X\longrightarrow X,\qquad (\lambda,x)\longmapsto\lambda x
\]
są ciągłe, gdy w dziedzinach przyjmiemy topologie produktowe.
Nie wymagamy na razie warunku Hausdorffa.
\end{definicja}

Warunek ciągłości mnożenia przez skalar jest \emph{łączny} względem
$\lambda$ i $x$. Sama ciągłość każdej funkcji $x\mapsto\lambda x$,
dla ustalonego $\lambda$, jest warunkiem słabszym.

\begin{przyklad}[Topologia trywialna]
Niech $X$ będzie dowolną niezerową przestrzenią liniową.
Topologia $\{\varnothing,X\}$ czyni $X$ przestrzenią
liniowo-topologiczną: każde odwzorowanie o wartościach w $X$ jest
ciągłe. Nie jest ona jednak Hausdorffa. Przykład wyjaśnia,
dlaczego rozdzielamy aksjomaty przestrzeni liniowo-topologicznej
od aksjomatu rozdzielania punktów.
\end{przyklad}

\subsection{Operacje liniowe jako homeomorfizmy}

\begin{stwierdzenie}\label{prop:translation}
Niech $X$ będzie przestrzenią liniowo-topologiczną nad $\K$.
Dla każdego $a\in X$ przesunięcie $T_a(x)=x+a$ jest
homeomorfizmem $X$ na $X$. Dla każdego $\lambda\in\K\setminus\{0\}$
odwzorowanie $M_\lambda(x)=\lambda x$ jest homeomorfizmem $X$ na $X$.
\end{stwierdzenie}
\begin{proof}
Ciągłość $T_a$ wynika ze złożenia ciągłego odwzorowania
$x\mapsto(x,a)$ z dodawaniem. Jego odwrotnością jest ciągłe $T_{-a}$.
Analogicznie $M_\lambda$ jest ciągłe jako funkcja jednej zmiennej
otrzymana z mnożenia przez skalar, a jego odwrotnością jest
$M_{1/\lambda}$.
\end{proof}

\begin{wniosek}\label{cor:zero-neighborhood}
Dla dowolnej przestrzeni liniowo-topologicznej $X$, punktu $x\in X$
i podzbioru $U\subseteq X$ zachodzi równoważność:
$U\in\mathcal N_X(x)$ wtedy i tylko wtedy, gdy
$U-x\in\mathcal N_X(0)$. W szczególności topologia przestrzeni
liniowo-topologicznej jest całkowicie określona przez otoczenia zera.
\end{wniosek}
\begin{proof}
Korzystamy z tego, że $T_{-x}$ jest homeomorfizmem i
$T_{-x}(x)=0$.
\end{proof}

\begin{uwaga}
Jednorodność topologiczna nie oznacza, że każda przestrzeń
liniowo-topologiczna jest metryzowalna. Problem metryzowalności
będziemy rozstrzygać za pomocą baz otoczeń zera.
\end{uwaga}

\subsection{Otoczenia zera, zbiory zrównoważone i pochłaniające}

\begin{definicja}\label{def:balanced-absorbing}
Niech $X$ będzie przestrzenią liniową nad $\K$, a $A\subseteq X$.
Mówimy, że:
\begin{enumerate}[label=\textup{(\roman*)}]
 \item $A$ jest \emph{zrównoważony}, jeżeli $\lambda A\subseteq A$
       dla każdego $\lambda\in\K$ z $|\lambda|\leq1$;
 \item $A$ jest \emph{pochłaniający}, jeżeli dla każdego $x\in X$
       istnieje $t>0$ takie, że $x\in tA$;
 \item $A$ jest \emph{wypukły}, jeżeli $sx+(1-s)y\in A$
       dla wszystkich $x,y\in A$ i $s\in[0,1]$.
\end{enumerate}
W drugim punkcie $tA=\{ta:a\in A\}$, gdzie $t$ jest rzeczywistą
liczbą dodatnią, również gdy $\K=\C$.
\end{definicja}

\begin{uwaga}
Wszystkie trzy warunki są algebraiczne. Zrównoważenie nie
implikuje wypukłości. W analizie lokalnie wypukłej będziemy
konstruować bazy otoczeń zera złożone ze zbiorów jednocześnie
wypukłych i zrównoważonych.
\end{uwaga}

\begin{twierdzenie}[Podstawowe własności otoczeń zera]\label{thm:zero-neighborhoods}
Niech $X$ będzie przestrzenią liniowo-topologiczną nad $\K$.
Dla każdego $U\in\mathcal N_X(0)$ zachodzą następujące własności:
\begin{enumerate}[label=\textup{(\roman*)}]
 \item istnieje $V\in\mathcal N_X(0)$ takie, że $V+V\subseteq U$;
 \item istnieje otwarte, symetryczne otoczenie $V$ zera
       ($V=-V$) zawarte w $U$;
 \item istnieje otwarte, zrównoważone otoczenie $W$ zera
       zawarte w $U$;
 \item $U$ jest zbiorem pochłaniającym.
\end{enumerate}
\end{twierdzenie}
\begin{proof}
Wybierzmy otwarte otoczenie $U_0$ zera z $U_0\subseteq U$.
Ciągłość dodawania w punkcie $(0,0)$ daje otoczenia otwarte
$V_1,V_2$ zera takie, że $V_1+V_2\subseteq U_0$.
Dla $V=V_1\cap V_2$ otrzymujemy (i). Zbiór
$U_0\cap(-U_0)$ jest otwartym, symetrycznym otoczeniem zera;
to dowodzi (ii).

Dla (iii) oznaczmy przez $D=\{\lambda\in\K:|\lambda|\leq1\}$
domknięty dysk jednostkowy (dla $\K=\R$ przedział $[-1,1]$).
Zbiór $D$ jest zwarty. Dla każdego $\lambda\in D$, z ciągłości
mnożenia przez skalar w $(\lambda,0)$, istnieją otoczenia otwarte
$O_\lambda\subseteq\K$ punktu $\lambda$ oraz $V_\lambda\subseteq X$
zera takie, że $O_\lambda V_\lambda\subseteq U_0$.
Wybieramy skończone podpokrycie
$D\subseteq O_{\lambda_1}\cup\dots\cup O_{\lambda_m}$ i kładziemy
$V=\bigcap_{j=1}^m V_{\lambda_j}$. Wtedy $DV\subseteq U_0$.
Zbiór $A=DV$ jest zrównoważony, ponieważ $\mu D\subseteq D$
dla $|\mu|\leq1$. Ponieważ $1\in D$, zachodzi $V\subseteq A$,
a zatem $A$ jest otoczeniem zera. Wreszcie $W=\operatorname{int}A$
jest otwarte i zawiera zero. Jeśli $0<|\mu|\leq1$, to
$\mu W$ jest otwartym podzbiorem $\mu A\subseteq A$, wobec czego
$\mu W\subseteq\operatorname{int}A=W$. Dla $\mu=0$ mamy
$\mu W=\{0\}\subseteq W$. Tak więc $W$ jest zrównoważone.

Dla (iv) ustalmy $x\in X$. Funkcja
$\K\ni\lambda\mapsto\lambda x\in X$ jest ciągła i przyjmuje
wartość zero dla $\lambda=0$. Istnieje zatem rzeczywiste
$\varepsilon>0$ takie, że $\varepsilon x\in U$.
Wtedy $x\in\varepsilon^{-1}U$, czyli $U$ jest pochłaniający.
\end{proof}

\begin{przyklad}[Zrównoważenie bez wypukłości]
Zbiór $A=\{(x,y)\in\R^2:|x|\leq1,\ |y|\leq1\}\cup
\{(x,y)\in\R^2:|x|\leq2,\ |y|\leq\frac12\}$
jest zrównoważony, lecz nie jest wypukły:
$(1,1),(2,\tfrac12)\in A$, podczas gdy
$(\tfrac32,\tfrac34)\notin A$.
\end{przyklad}

\subsection{Warunek Hausdorffa i domknięcie zera}

\begin{twierdzenie}\label{thm:hausdorff-zero}
Niech $X$ będzie przestrzenią liniowo-topologiczną nad $\K$.
Następujące warunki są równoważne:
\begin{enumerate}[label=\textup{(\roman*)}]
 \item $X$ jest przestrzenią Hausdorffa;
 \item zbiór $\{0\}$ jest domknięty w $X$;
 \item $\bigcap_{U\in\mathcal N_X(0)}U=\{0\}$.
\end{enumerate}
Ponadto $N=\cl{\{0\}}$ jest domkniętą podprzestrzenią liniową
$X$ oraz
\[
 N=\bigcap_{U\in\mathcal N_X(0)}U.
\]
\end{twierdzenie}
\begin{proof}
Implikacja (i)$\Rightarrow$(ii) zachodzi w każdej przestrzeni
Hausdorffa. Aby wykazać wzór na $N$, ustalmy $x\in X$.
Jeżeli $x\in N$, to każde otoczenie $x$ zawiera zero.
Dla dowolnego symetrycznego otoczenia $V$ zera oznacza to,
że $0\in x+V$, czyli $-x\in V$, a z symetrii także $x\in V$.
Ponieważ symetryczne otoczenia tworzą bazę, $x$ należy do
wszystkich otoczeń zera.

Odwrotnie, jeśli $x$ należy do wszystkich otoczeń zera,
to dla każdego otoczenia $x+V$ punktu $x$, gdzie $V$ jest
symetrycznym otoczeniem zera, mamy $-x\in V$, skąd
$0\in x+V$. Oznacza to $x\in N$. Udowodniliśmy wzór.

Jeżeli (ii), to $N=\{0\}$ i powyższy wzór daje (iii).
Jeżeli (iii), to dla dowolnych $x\ne y$ istnieje otoczenie
$U$ zera, do którego nie należy $y-x$. Dobierzmy symetryczne
otoczenie $V$ zera takie, że $V+V\subseteq U$. Otoczenia
$x+V$ i $y+V$ są rozłączne: w przeciwnym razie
$y-x\in V-V=V+V\subseteq U$. To daje (i).

Pozostaje wykazać, że $N$ jest podprzestrzenią.
Jeśli $x,y\in N$ i $U$ jest otoczeniem zera, wybierzmy
otoczenie $V$ z $V+V\subseteq U$. Wtedy $x,y\in V$,
a więc $x+y\in U$. Ponieważ $U$ było dowolne, $x+y\in N$.
Dla $\lambda\ne0$ i $x\in N$, zbiór $\lambda^{-1}U$
jest otoczeniem zera, więc $x\in\lambda^{-1}U$,
czyli $\lambda x\in U$; dla $\lambda=0$ jest to oczywiste.
Zatem $N$ jest podprzestrzenią. Jest domknięta z definicji.
\end{proof}

\begin{uwaga}
Od tego miejsca \emph{przestrzeń liniowo-topologiczna} będzie,
o ile nie zaznaczymy inaczej, oznaczać przestrzeń liniowo-topologiczną
Hausdorffa. Warto jednak pamiętać, że wymaganie Hausdorffa jest
założeniem, a nie konsekwencją ciągłości działań liniowych.
\end{uwaga}

\subsection{Ciągłe odwzorowania liniowe}

\begin{twierdzenie}[Kryterium ciągłości w zerze]\label{thm:linear-continuity-zero}
Niech $X,Y$ będą przestrzeniami liniowo-topologicznymi nad tym samym
ciałem $\K$ i niech $T\colon X\to Y$ będzie odwzorowaniem liniowym.
Następujące warunki są równoważne:
\begin{enumerate}[label=\textup{(\roman*)}]
 \item $T$ jest ciągłe na $X$;
 \item $T$ jest ciągłe w punkcie $0$;
 \item dla każdego $V\in\mathcal N_Y(0)$ istnieje
       $U\in\mathcal N_X(0)$ takie, że $T(U)\subseteq V$.
\end{enumerate}
\end{twierdzenie}
\begin{proof}
Implikacja (i)$\Rightarrow$(ii) jest oczywista.
Warunek (iii) jest dokładnym zapisem definicji ciągłości w zerze,
ponieważ $T(0)=0$.
Jeśli (iii) oraz $x\in X$, to dla dowolnego otoczenia
$T(x)+V$ punktu $T(x)$ istnieje otoczenie $U$ zera takie,
że $T(U)\subseteq V$. Z liniowości
$T(x+U)=T(x)+T(U)\subseteq T(x)+V$.
Oznacza to ciągłość $T$ w punkcie $x$, a więc (i).
\end{proof}

\begin{wniosek}\label{cor:kernel-closed}
Niech $X$ będzie przestrzenią liniowo-topologiczną, $Y$
przestrzenią liniowo-topologiczną Hausdorffa, a $T\colon X\to Y$
ciągłym odwzorowaniem liniowym. Wtedy $\ker T$ jest domkniętą
podprzestrzenią liniową $X$.
\end{wniosek}
\begin{proof}
Jest podprzestrzenią z liniowości $T$. Ponieważ $Y$ jest
Hausdorffa, zbiór $\{0\}$ jest domknięty na mocy
twierdzenia~\ref{thm:hausdorff-zero}. Mamy
$\ker T=T^{-1}(\{0\})$, zatem jądro jest domknięte.
\end{proof}

\begin{przyklad}[Ciągła bijekcja liniowa bez ciągłej odwrotności]
Na $X=C([0,1],\R)$ rozważmy dwie normy:
$\|f\|_\infty=\max_{x\in[0,1]}|f(x)|$ oraz
$\|f\|_1=\int_{[0,1]}|f(x)|\,\mathrm d\lambda(x)$,
gdzie $\lambda$ jest miarą Lebesgue'a.
Identyczność $I\colon(X,\|\cdot\|_\infty)\to(X,\|\cdot\|_1)$
jest liniową ciągłą bijekcją, bo $\|f\|_1\leq\|f\|_\infty$.
Odwrotność nie jest ciągła. Istotnie, funkcje
$f_n(x)=\max\{1-nx,0\}$ spełniają $\|f_n\|_\infty=1$
i $\|f_n\|_1=1/(2n)$.
Nie przeczy to twierdzeniu o odwzorowaniu otwartym:
przestrzeń $(X,\|\cdot\|_1)$ nie jest zupełna. Rzeczywiście,
można dobrać ciąg funkcji ciągłych $g_n$ przyjmujących wartości
w $[0,1]$, równych zeru na $[0,\tfrac12-\tfrac1n]$ oraz jeden
na $[\tfrac12+\tfrac1n,1]$, liniowych pomiędzy tymi przedziałami
(dla $n\geq3$). Wtedy
$\|g_n-\mathbf1_{[1/2,1]}\|_1\leq2/n\to0$.
W szczególności $(g_n)$ jest ciągiem Cauchy'ego w normie
$\|\cdot\|_1$. Gdyby zbiegał do $g\in C([0,1])$,
jednoznaczność granicy w $L^1$ dawałaby
$g=\mathbf1_{[1/2,1]}$ prawie wszędzie. Z ciągłości $g$
wynikałoby wtedy $g=0$ na $[0,1/2)$ oraz $g=1$ na
$(1/2,1]$, co jest sprzeczne z ciągłością w $1/2$.
Zatem $(X,\|\cdot\|_1)$ nie jest zupełna.
\end{przyklad}

\subsection{Zbiory ograniczone bez użycia normy}

\begin{definicja}[Ograniczoność w sensie von Neumanna]\label{def:tvs-bounded}
Niech $X$ będzie przestrzenią liniowo-topologiczną nad $\K$.
Zbiór $B\subseteq X$ nazywamy \emph{ograniczonym}, jeżeli
dla każdego otoczenia $U$ zera istnieje $t>0$ takie,
że $B\subseteq tU$. W definicji nie wymaga się, aby $B$
był domknięty ani aby był otoczeniem zera.
\end{definicja}

\begin{lemat}\label{lem:bounded-test}
Niech $X$ będzie przestrzenią liniowo-topologiczną, a $B\subseteq X$.
Zbiór $B$ jest ograniczony wtedy i tylko wtedy, gdy dla każdego
zrównoważonego otoczenia $U$ zera istnieje $\varepsilon>0$ takie,
że $\lambda B\subseteq U$ dla wszystkich $\lambda\in\K$
z $|\lambda|\leq\varepsilon$.
\end{lemat}
\begin{proof}
Załóżmy, że $B$ jest ograniczony, i weźmy zrównoważone
otoczenie $U$ zera. Istnieje $t>0$ z $B\subseteq tU$.
Dla $\varepsilon=1/t$ i $|\lambda|\leq\varepsilon$
mamy $\lambda B\subseteq(\lambda t)U\subseteq U$.

Odwrotnie, niech $V$ będzie dowolnym otoczeniem zera.
Na mocy twierdzenia~\ref{thm:zero-neighborhoods} istnieje
zrównoważone otoczenie $U\subseteq V$. Z założenia
$\varepsilon B\subseteq U$ dla pewnego $\varepsilon>0$,
a zatem $B\subseteq\varepsilon^{-1}V$.
\end{proof}

\begin{twierdzenie}\label{thm:compact-bounded}
Niech $X$ będzie przestrzenią liniowo-topologiczną.
Wówczas:
\begin{enumerate}[label=\textup{(\roman*)}]
 \item każdy skończony podzbiór $X$ jest ograniczony;
 \item każdy zwarty podzbiór $X$ jest ograniczony;
 \item jeżeli $T\colon X\to Y$ jest ciągłym odwzorowaniem liniowym
       między przestrzeniami liniowo-topologicznymi, to obraz
       każdego ograniczonego zbioru $B\subseteq X$ jest ograniczony w $Y$.
\end{enumerate}
\end{twierdzenie}
\begin{proof}
Wystarczy udowodnić (ii) i (iii), ponieważ każdy zbiór
skończony jest zwarty. Ustalmy zwarty zbiór $K\subseteq X$
oraz otoczenie $U$ zera. Wybierzmy otwarte, zrównoważone
otoczenie $V\subseteq U$. Zbiór $V$ jest pochłaniający,
więc $X=\bigcup_{n\in\N,\,n\geq1}nV$; rodzina ta jest
wstępująca, ponieważ $V$ jest zrównoważony. Każdy $nV$
jest otwarty. Zwartość $K$ daje skończone podpokrycie,
a w konsekwencji pewne $n$ z $K\subseteq nV\subseteq nU$.
To dowodzi (ii).

Dla (iii) niech $W\in\mathcal N_Y(0)$.
Z ciągłości $T$ zbiór $T^{-1}(W)$ jest otoczeniem zera w $X$.
Istnieje $t>0$ z $B\subseteq tT^{-1}(W)$.
Z liniowości $T(B)\subseteq tW$, co kończy dowód.
\end{proof}

\begin{przyklad}[Ograniczoność normowa]
Niech $(X,\|\cdot\|)$ będzie przestrzenią unormowaną.
Zbiór $B\subseteq X$ jest ograniczony w sensie
definicji~\ref{def:tvs-bounded} wtedy i tylko wtedy, gdy
$\sup_{x\in B}\|x\|<\infty$ (przyjmujemy supremum zbioru
pustego równe zero). Istotnie, wystarczy sprawdzać
otoczenia postaci $\{x:\|x\|<r\}$, $r>0$.
\end{przyklad}

\subsection{Przykłady przestrzeni nienormowalnych}

\begin{przyklad}[Przestrzeń wszystkich ciągów]\label{ex:product-tvs}
Niech $X=\K^{\N}$ będzie przestrzenią wszystkich ciągów
$x=(x_1,x_2,\ldots)$, wyposażoną w topologię produktową.
Bazę otoczeń zera tworzą zbiory
\[
 U(F,\varepsilon)=\big\{x\in\K^{\N}:|x_j|<\varepsilon_j
          \text{ dla }j\in F\big\},
\]
gdzie $F\subseteq\N$ jest skończony, a $\varepsilon_j>0$ dla
$j\in F$. Dodawanie i mnożenie przez skalar są ciągłe,
ponieważ ich złożenia ze wszystkimi rzutami współrzędnymi
są ciągłe. Zatem $X$ jest przestrzenią liniowo-topologiczną Hausdorffa.

Nie jest ona normowalna: żadne otoczenie zera nie jest ograniczone.
Istotnie, niech $U(F,\varepsilon)$ będzie bazowym otoczeniem
zera i wybierzmy $j\notin F$. Weźmy otoczenie
$W=\{x:|x_j|<1\}$. Dla każdego $t>0$ ciąg
$x=(t+1)e_j$ należy do $U(F,\varepsilon)$, ale nie należy
do $tW$. Stąd $U(F,\varepsilon)$ nie jest ograniczone,
a więc żadne większe otoczenie zera nie jest ograniczone.
Gdyby topologia pochodziła od normy, otwarta kula jednostkowa
byłaby ograniczonym otoczeniem zera, co prowadzi do sprzeczności.

Warto zauważyć, że $e_n\to0$ w topologii produktowej:
dla każdego ustalonego $j$ współrzędna $(e_n)_j$
jest równa zero dla wszystkich dostatecznie dużych $n$.
\end{przyklad}

\begin{przyklad}[Topologia zbieżności punktowej funkcji]
Niech $S$ będzie zbiorem, a $X\subseteq\K^S$ podprzestrzenią
liniową, np. przestrzenią wszystkich funkcji ciągłych na
$S=[0,1]$. Topologia zbieżności punktowej na $X$
jest topologią podprzestrzeni odziedziczoną z produktu
$\K^S$. Ciąg (lub ogólniej sieć) funkcji zbiega w tej
topologii wtedy i tylko wtedy, gdy zbiega w każdym punkcie $S$.
W szczególności topologia ta nie wymaga normy supremum.
\end{przyklad}

\begin{przyklad}[Nieciągły funkcjonał na przestrzeni unormowanej]
Rozważmy $X=\R[t]$, przestrzeń wielomianów rzeczywistych
na $[0,1]$, z normą supremum. Funkcjonał liniowy
$L(p)=p'(0)$ nie jest ciągły. Dla $p_n(t)=(1-t)^n$
mamy $\|p_n\|_\infty=1$, ale $|L(p_n)|=n$.
Gdyby $L$ był ciągły w zerze, z liniowości otrzymalibyśmy
$|L(p)|\leq C\|p\|_\infty$ dla pewnego $C>0$ i wszystkich
$p\in X$, co przeczy powyższemu rachunkowi.
\end{przyklad}

\subsection{Przykład zasadniczy: \texorpdfstring{$L^p$}{Lp} dla \texorpdfstring{$0<p<1$}{0 < p < 1}}

W tym przykładzie zakładamy znajomość elementarnych pojęć miary
Lebesgue'a. Ma on pokazać, dlaczego lokalna wypukłość będzie
później odgrywać kluczową rolę.

\begin{twierdzenie}[Zupełna przestrzeń liniowo-topologiczna
bez niezerowych funkcjonałów ciągłych]\label{thm:lp-trivial-dual}
Niech $0<p<1$, a $X=L^p([0,1],\lambda)$ oznacza przestrzeń
klas funkcji mierzalnych $f\colon[0,1]\to\K$ takich, że
$\int_{[0,1]}|f(x)|^p\,\mathrm d\lambda(x)<\infty$,
gdzie utożsamiamy funkcje równe prawie wszędzie.
Wówczas wzór
\[
 d(f,g)=\int_{[0,1]}|f(x)-g(x)|^p\,\mathrm d\lambda(x)
\]
definiuje zupełną metrykę na $X$ zgodną ze strukturą
przestrzeni liniowo-topologicznej. Co więcej, każdy
ciągły funkcjonał $\K$-liniowy $\Lambda\colon X\to\K$
jest równy zeru.
\end{twierdzenie}
\begin{proof}
Dla liczb rzeczywistych $a,b\geq0$ i $0<p<1$ zachodzi
$(a+b)^p\leq a^p+b^p$. W połączeniu z nierównością
$|f-h|\leq|f-g|+|g-h|$ daje to nierówność trójkąta dla $d$.
Pozostałe aksjomaty metryki wynikają z własności całki.
Dodawanie jest ciągłe na mocy
$d(f_1+g_1,f_2+g_2)\leq d(f_1,f_2)+d(g_1,g_2)$.
Dla $\alpha,\beta\in\K$ oraz $f,g\in X$ mamy
\[
 d(\alpha f,\beta g)
 \leq |\alpha|^p d(f,g)+|\alpha-\beta|^p d(g,0).
\]
Jest to oszacowanie wystarczające do łącznej ciągłości
mnożenia przez skalar.

Wykażemy zupełność. Niech $(f_n)$ będzie ciągiem Cauchy'ego
w metryce $d$; wybierzmy mierzalnych reprezentantów jego wyrazów.
Możemy wybrać podciąg $(f_{n_k})$ tak, aby
$\sum_{k=1}^\infty d(f_{n_{k+1}},f_{n_k})<\infty$.
Z twierdzenia o monotonicznej zbieżności wynika, że
\[
 \sum_{k=1}^\infty|f_{n_{k+1}}(x)-f_{n_k}(x)|^p<\infty
 \quad\text{dla prawie każdego }x\in[0,1].
\]
Ponieważ $p<1$, zbieżność szeregu dodatnich liczb
$\sum_k a_k^p$ implikuje zbieżność $\sum_k a_k$
(wszystkie dostatecznie późne $a_k$ należą do $[0,1]$,
a wtedy $a_k\leq a_k^p$). Zatem szereg różnic
$\sum_k(f_{n_{k+1}}-f_{n_k})$ zbiega bezwzględnie
prawie wszędzie. Oznaczmy jego granicę powiększoną o
$f_{n_1}$ przez $f$. Z nierówności $(\sum a_k)^p\leq\sum a_k^p$
i lematu Fatou dostajemy
\[
 d(f_{n_k},f)
 \leq \sum_{j=k}^\infty d(f_{n_{j+1}},f_{n_j})\longrightarrow0.
\]
W szczególności $f\in L^p$ (bo $|f|^p\leq|f-f_{n_1}|^p+|f_{n_1}|^p$),
a ciąg Cauchy'ego $(f_n)$ także zbiega do $f$.

Pozostało zbadać dual. Niech $\Lambda\colon X\to\K$
będzie ciągłym funkcjonałem liniowym. Ciągłość w zerze
daje $\delta>0$ takie, że
$d(f,0)<\delta^p$ implikuje $|\Lambda(f)|<1$.
Z jednorodności $d(af,0)=|a|^p d(f,0)$ wynika istnienie
$C>0$ takie, że
\[
 |\Lambda(f)|\leq C\Big(\int_{[0,1]}|f|^p\,\mathrm d\lambda\Big)^{1/p}
 \qquad(f\in X).
\]
Rzeczywiście, dla $f\ne0$ wystarczy zastosować poprzedni
warunek do $a f$ z $|a|=\delta/(2d(f,0)^{1/p})$.

Ustalmy dowolny zbiór mierzalny $A\subseteq[0,1]$.
Dla każdego $n\geq1$ można rozbić $A$ na rozłączne
zbiory mierzalne $A_1,\ldots,A_n$ o równej mierze
$\lambda(A_i)=\lambda(A)/n$; wynika to z bezatomowości
miary Lebesgue'a. Stąd
\[
 |\Lambda(\mathbf1_A)|
 \leq\sum_{i=1}^n|\Lambda(\mathbf1_{A_i})|
 \leq C\lambda(A)^{1/p}n^{1-1/p}\longrightarrow0,
\]
ponieważ $1-1/p<0$. Funkcjonał $\Lambda$ znika więc
na funkcjach charakterystycznych, a z liniowości również
na funkcjach prostych. Uzasadnijmy również używaną gęstość funkcji prostych.
Dla $f\in L^p([0,1],\lambda)$ obcięcie
$f_N=f\mathbf1_{\{|f|\leq N\}}$ spełnia
$d(f_N,f)\to0$ na mocy zbieżności zmajoryzowanej.
Funkcja $f_N$ jest ograniczona; dzieląc kwadrat
$[-N,N]^2$ (dla $\K=\C$) albo przedział $[-N,N]$
(dla $\K=\R$) na dostatecznie małe komórki i zastępując
wartości $f_N$ reprezentantami komórek, otrzymujemy
mierzalne funkcje proste $s_k$ z
$|s_k-f_N|\leq1/k$ prawie wszędzie. Wtedy
$d(s_k,f_N)\leq k^{-p}\lambda([0,1])\to0$.
Funkcje proste są więc gęste w $L^p$; z ciągłości
$\Lambda=0$ na całym $X$.
\end{proof}

\begin{uwaga}\label{rem:local-convexity-purpose}
Na przestrzeni $L^p([0,1])$ z $0<p<1$ nie można oddzielić
niezerowego wektora od zera za pomocą ciągłego funkcjonału
liniowego: dual topologiczny jest trywialny. W dalszych
rozdziałach udowodnimy, że w każdej przestrzeni lokalnie
wypukłej Hausdorffa ciągłe funkcjonały liniowe oddzielają
punkty. Zatem powyższa przestrzeń \emph{nie jest lokalnie wypukła},
choć jest metryzowalna i zupełna.
\end{uwaga}

\subsection{Co zostało ustalone, a co wymaga dodatkowych założeń?}

W dowolnej przestrzeni liniowo-topologicznej możemy mówić o
ciągłości operatorów liniowych, zbiorach ograniczonych, otoczeniach
zera, podprzestrzeniach i zbieżności. Nie potrzebowaliśmy do tego
ani normy, ani iloczynu skalarnego, ani zupełności.

Nie możemy natomiast bez dalszych założeń zakładać, że istnieją
niezerowe ciągłe funkcjonały liniowe. Twierdzenie
\ref{thm:lp-trivial-dual} pokazuje, że jest to rzeczywista
przeszkoda. W następnym rozdziale zbadamy produkty, podprzestrzenie,
ilorazy, ograniczoność i zbieżność sieci; następnie wprowadzimy
seminormy i przestrzenie lokalnie wypukłe.

\begin{uwaga}[Dalszy kierunek skryptu]
Pochodne Gâteaux i Frécheta nie wymagają teorii dystrybucji.
Wprowadzimy je w osobnym bloku dotyczącym przestrzeni Banacha
i Hilberta, a następnie wykorzystamy w twierdzeniu o funkcji
uwikłanej, zasadach wariacyjnych i rachunku na przestrzeniach
funkcyjnych. Dla przestrzeni zespolonych wyraźnie odróżnimy
różniczkowalność w sensie rzeczywistym od pochodnej
zespolono-liniowej.
\end{uwaga}


% Wersja źródłowa: audyt II, jednostka 02
\bigskip
\setcounter{section}{1}
\section{Podprzestrzenie, produkty, ilorazy i topologie uniwersalne}\label{sec:konstrukcje}

\subsection{Dlaczego potrzebujemy konstrukcji nowych przestrzeni?}

Z jednej przestrzeni liniowo-topologicznej (TVS) powstają naturalnie nowe:
zbiór funkcji spełniających liniowy warunek, rodzina funkcji zależnych
od wielu parametrów, przestrzeń funkcji rozpatrywanych z dokładnością
do pewnej relacji równoważności. Odpowiadają im odpowiednio
\emph{podprzestrzenie}, \emph{produkty} oraz \emph{ilorazy}.
Istotą konstrukcji jest nie tylko zdefiniowanie zbioru i działań,
lecz także dobranie topologii tak, aby właściwe odwzorowania były
ciągłe. W dalszych rozdziałach ta idea posłuży do budowy
przestrzeni funkcji testowych i dystrybucji.

Przyjmujemy konwencje \hyperref[sec:tvs]{rozdziału 1}: $\K\in\{\R,\C\}$,
a wszystkie przestrzenie liniowo-topologiczne są Hausdorffa,
chyba że wyraźnie zaznaczymy inaczej. Pojęcie \emph{topologii
indukowanej} i \emph{topologii produktowej} rozumiemy w zwykłym sensie
topologii ogólnej. W tej części przypomnimy ich precyzyjne własności.

\subsection{Podprzestrzenie liniowo-topologiczne}

\begin{twierdzenie}[Topologia podprzestrzeni]\label{thm:subspace-tvs}
Niech $X$ będzie przestrzenią liniowo-topologiczną nad $\K$,
a $M\subseteq X$ podprzestrzenią liniową. Z topologią indukowaną
$\tau_M=\{O\cap M:O\subseteq X\text{ jest otwarty}\}$
przestrzeń $M$ jest przestrzenią liniowo-topologiczną nad $\K$.
Jeżeli $X$ jest Hausdorffa, to również $M$ jest Hausdorffa.
Otoczenia zera w $M$ mają bazę postaci $U\cap M$, gdzie
$U$ przebiega otoczenia zera w $X$.
\end{twierdzenie}
\begin{proof}
Ograniczenie dodawania $+_X$ do $M\times M$ przyjmuje wartości
w $M$, a topologia na $M\times M$ jest topologią podprzestrzeni
$X\times X$. Jeżeli $O\cap M$ jest otwarty w $M$, to
\[
 (+_M)^{-1}(O\cap M)=(+_X)^{-1}(O)\cap(M\times M),
\]
więc dodawanie w $M$ jest ciągłe. Identyczny argument dla
$\K\times M\subseteq\K\times X$ daje ciągłość mnożenia
przez skalar. Dziedziczenie własności Hausdorffa jest własnością
każdej podprzestrzeni topologicznej. Ostatnie zdanie wynika
wprost z definicji topologii indukowanej i faktu, że $0\in M$.
\end{proof}

\begin{lemat}[Domknięcie podprzestrzeni]\label{lem:closure-subspace}
Niech $M$ będzie podprzestrzenią liniową przestrzeni
liniowo-topologicznej $X$, a $\mathcal N_X(0)$ rodziną
otoczeń zera w $X$. Wówczas $\overline M$ jest podprzestrzenią
liniową oraz
\[
 \overline M=\bigcap_{U\in\mathcal N_X(0)}(M+U),
 \qquad M+U=\{m+u:m\in M,\ u\in U\}.
\]
\end{lemat}
\begin{proof}
Jeżeli $x,y\in\overline M$, to dla dowolnego otoczenia $U$
zera wybieramy symetryczne otoczenie $V$ z $V+V\subseteq U$.
Istnieją $m_x,m_y\in M$ takie, że $x-m_x\in V$ oraz
$y-m_y\in V$. Stąd $(x+y)-(m_x+m_y)\in U$,
więc $x+y\in\overline M$. Dla $\lambda\ne0$ ciągłość
mnożenia przez skalar (a nawet fakt, że jest homeomorfizmem)
daje $\lambda\overline M=\overline{\lambda M}=\overline M$;
przypadek $\lambda=0$ jest oczywisty.

Jeżeli $x\in\overline M$, to dla każdego $U$ możemy wybrać
symetryczne $V\subseteq U$; przecięcie $(x+V)\cap M$
jest niepuste, więc $x=m-v\in M+U$ dla odpowiednich
$m\in M$, $v\in V$. Odwrotnie, niech $x\in\bigcap_U(M+U)$.
Dla otoczenia $x+V$ punktu $x$ wybieramy symetryczne $V$;
z założenia $x=m+v$ z $m\in M$, $v\in V$,
a więc $m=x-v\in x+V$. Każde otoczenie $x$ spotyka $M$,
czyli $x\in\overline M$.
\end{proof}

\begin{przyklad}[Dwie topologie na wielomianach]
Niech $P=\K[t]$, a $X=C([0,1],\K)$ z normą supremum.
Topologia indukowana na $P$ jest unormowana, lecz $P$
nie jest domknięte w $X$: wielomiany są gęste w $X$
(z twierdzenia Weierstrassa o aproksymacji). W szczególności
zupełność $X$ nie przechodzi na dowolną podprzestrzeń.
W przyszłości udowodnimy, że podprzestrzeń przestrzeni Banacha
jest zupełna w normie indukowanej wtedy i tylko wtedy, gdy
jest domknięta.
\end{przyklad}

\begin{stwierdzenie}[Uniwersalna własność podprzestrzeni]
Niech $M\subseteq X$ będzie podprzestrzenią liniową przestrzeni
liniowo-topologicznej $X$, $i:M\hookrightarrow X$ naturalnym
włożeniem, $Z$ dowolną przestrzenią topologiczną oraz
$f:Z\to M$ odwzorowaniem. Wówczas $f$ jest ciągłe
wtedy i tylko wtedy, gdy $i\circ f:Z\to X$ jest ciągłe.
\end{stwierdzenie}
\begin{proof}
Konieczność jest natychmiastowa, gdyż $i$ jest ciągłe.
Jeśli $i\circ f$ jest ciągłe, to dla otwartego $O\cap M$
w $M$ mamy $f^{-1}(O\cap M)=(i\circ f)^{-1}(O)$,
więc $f$ jest ciągłe.
\end{proof}

\subsection{Produkty dowolnych rodzin przestrzeni}

\begin{definicja}[Produkt przestrzeni liniowo-topologicznych]
Niech $I$ będzie zbiorem indeksów, a $(X_i)_{i\in I}$ rodziną
przestrzeni liniowo-topologicznych nad $\K$. Na przestrzeni
liniowej $X=\prod_{i\in I}X_i$ działań współrzędnych
wprowadzamy \emph{topologię produktową}: jej bazę tworzą zbiory
$\prod_{i\in I}O_i$, gdzie każde $O_i\subseteq X_i$ jest
otwarte i $O_i=X_i$ dla wszystkich poza skończenie wieloma
indeksami. Rzut na $i$-tą współrzędną oznaczamy przez
$\pi_i:X\to X_i$.
\end{definicja}

\begin{twierdzenie}[Produkt i jego własność uniwersalna]\label{thm:product-tvs}
Niech $(X_i)_{i\in I}$ będzie rodziną przestrzeni
liniowo-topologicznych Hausdorffa nad $\K$.
Wówczas $X=\prod_{i\in I}X_i$ z topologią produktową
jest przestrzenią liniowo-topologiczną Hausdorffa, a każdy
$\pi_i$ jest ciągłym, liniowym i otwartym odwzorowaniem na $X_i$.
Dla dowolnej przestrzeni topologicznej $Z$ odwzorowanie
$f:Z\to X$ jest ciągłe wtedy i tylko wtedy, gdy
$\pi_i\circ f:Z\to X_i$ jest ciągłe dla każdego $i\in I$.
\end{twierdzenie}
\begin{proof}
Rzuty są ciągłe z definicji topologii produktowej. Dla zbioru
bazowego $\prod_i O_i$ zachodzi $\pi_j(\prod_i O_i)=O_j$,
ponieważ wszystkie pozostałe $O_i$ są niepuste, o ile
rozpatrujemy niepusty zbiór bazowy; rzut zbioru pustego
jest pusty. Zatem rzuty są otwarte.

Dodawanie $a:X\times X\to X$ jest ciągłe, bo każde
$\pi_i\circ a:(x,y)\mapsto x_i+y_i$ jest ciągłe
(jest to złożenie rzutowań i dodawania w $X_i$).
To samo rozumowanie stosuje się do mnożenia skalarnego
$m:\K\times X\to X$. Przestrzeń $X$ jest Hausdorffa,
ponieważ dwa różne elementy różnią się na pewnej współrzędnej,
a tę można rozdzielić otoczeniami w odpowiednim $X_i$.

Jeśli wszystkie $f_i=\pi_i\circ f$ są ciągłe,
to dla podzbioru bazowego $\prod_iO_i$ mamy
$f^{-1}(\prod_iO_i)=\bigcap_{i\in F}f_i^{-1}(O_i)$,
gdzie $F=\{i:O_i\ne X_i\}$ jest skończony;
jest to zbiór otwarty. Implikacja odwrotna wynika
ze złożenia odwzorowań ciągłych.
\end{proof}

\begin{uwaga}[Dlaczego tylko skończenie wiele warunków?]
W produkcie nieskończonym wymaganie, aby wszystkie
współrzędne jednocześnie należały do wybranych małych otoczeń,
może zniszczyć ciągłość działań. Na przykład w $\K^{\N}$
zbiory $\{x:|x_n|<1\text{ dla każdego }n\}$ nie są otwarte
w topologii produktowej. Ograniczanie skończenie wielu
współrzędnych jest kluczowe dla uniwersalnej własności produktu.
\end{uwaga}

\begin{twierdzenie}[Ograniczoność w produkcie]\label{thm:product-bounded}
Niech $(X_i)_{i\in I}$ będzie rodziną przestrzeni
liniowo-topologicznych i niech $B\subseteq\prod_{i\in I}X_i$.
Wówczas $B$ jest ograniczony w sensie von Neumanna wtedy
i tylko wtedy, gdy $\pi_i(B)$ jest ograniczony w $X_i$
dla każdego $i\in I$.
\end{twierdzenie}
\begin{proof}
W jedną stronę stosujemy rezultat z \hyperref[sec:tvs]{rozdziału 1}: ciągły operator liniowy
przenosi zbiory ograniczone na zbiory ograniczone. Stosujemy go do ciągłych rzutów $\pi_i$.
Odwrotnie, niech $U$ będzie otoczeniem zera w produkcie.
Możemy wybrać bazowe otoczenie $V=\prod_i V_i\subseteq U$,
gdzie $V_i=X_i$ poza skończonym zbiorem $F$, a dla $i\in F$
przyjmijmy $V_i$ zrównoważone. Dla każdego $i\in F$ istnieje
$t_i>0$ takie, że $\pi_i(B)\subseteq t_iV_i$.
Jeśli $F$ jest niepusty, ustalmy $t\geq\max(1,\max_{i\in F}t_i)$.
Zrównoważenie $V_i$ daje $t_iV_i\subseteq tV_i$, więc
$B\subseteq tV\subseteq tU$. Jeżeli $F$ jest pusty,
wystarczy przyjąć $t=1$.
\end{proof}

\begin{przyklad}[Dwa rodzaje zbieżności funkcji]
Niech $X=C([0,1],\K)$. Norma supremum wyznacza topologię
zbieżności jednostajnej; topologię zbieżności punktowej
otrzymujemy przez zanurzenie $X\subseteq\K^{[0,1]}$.
Identyczność z topologii normowej do punktowej jest ciągła.
Nie jest jednak homeomorfizmem: funkcje
$f_n(t)=\max\{0,1-n^2|t-1/n|\}$, $n\geq2$, są ciągłe,
zbiegają punktowo do zera, ale $\|f_n\|_\infty=1$.
Istotnie, dla każdego ustalonego $t\in[0,1]$ nośnik $f_n$
nie zawiera $t$ dla wszystkich dostatecznie dużych $n$.
\end{przyklad}

\begin{twierdzenie}[Funkcjonały na produkcie przestrzeni skalarnych]\label{thm:product-dual}
Niech $I$ będzie dowolnym zbiorem, a $X=\K^I$ wyposażmy
w topologię produktową. Każdy ciągły funkcjonał liniowy
$L:X\to\K$ ma postać
\[
 L(x)=\sum_{i\in F}a_i x_i
\]
przy pewnym skończonym zbiorze $F\subseteq I$ oraz
skalarach $a_i\in\K$. Odwrotnie, każdy taki funkcjonał
jest ciągły.
\end{twierdzenie}
\begin{proof}
Z ciągłości $L$ w zerze istnieje bazowe otoczenie zera
$U=\{x:|x_i|<\varepsilon_i\text{ dla }i\in F\}$
takie, że $|L(x)|<1$ dla $x\in U$.
Jeśli $x_i=0$ dla każdego $i\in F$, to $\lambda x\in U$
dla dowolnego $\lambda\in\K$. Stąd
$|\lambda L(x)|<1$ dla wszystkich $\lambda$, więc $L(x)=0$.
Dla dowolnego $x\in X$ odejmijmy skończoną sumę
$\sum_{i\in F}x_i e_i$, gdzie $e_i$ oznacza funkcję
równą $1$ na $i$ i $0$ poza nim. Pozostały wektor
ma zera na $F$, dlatego $L(x)=\sum_{i\in F}L(e_i)x_i$.
Odwrotność wynika z ciągłości rzutów i działań skalarnych.
\end{proof}

\begin{uwaga}
Twierdzenie~\ref{thm:product-dual} jest pierwszą wyraźną
manifestacją tego, że \emph{dual topologiczny} może być
znacznie mniejszy od dualu algebraicznego. Nie należy go
mylić z analogiczną własnością dowolnych produktów:
ciągły operator $\K^{\N}\to\K^{\N}$ może zależeć
od nieskończenie wielu współrzędnych (np. identyczność).
\end{uwaga}

\subsection{Przestrzeń ilorazowa i topologia ilorazowa}

\begin{definicja}[Iloraz przez podprzestrzeń]
Niech $X$ będzie przestrzenią liniowo-topologiczną nad $\K$,
a $M\subseteq X$ podprzestrzenią liniową. Na przestrzeni
wektorowej $X/M=\{x+M:x\in X\}$ wprowadzamy
\emph{topologię ilorazową}: zbiór $W\subseteq X/M$
jest otwarty wtedy i tylko wtedy, gdy $q^{-1}(W)$
jest otwarty w $X$, gdzie $q:X\to X/M$, $q(x)=x+M$.
\end{definicja}

\begin{twierdzenie}[Iloraz jako TVS]\label{thm:quotient-tvs}
Niech $X$ będzie dowolną przestrzenią liniowo-topologiczną,
niekoniecznie Hausdorffa, a $M\subseteq X$ podprzestrzenią.
Wówczas:
\begin{enumerate}[label=\textup{(\roman*)}]
\item odwzorowanie $q:X\to X/M$ jest liniowe, ciągłe,
      surjektywne i otwarte;
\item topologia ilorazowa czyni $X/M$ przestrzenią
      liniowo-topologiczną;
\item $X/M$ jest Hausdorffa wtedy i tylko wtedy,
      gdy $M$ jest domknięta w $X$.
\end{enumerate}
\end{twierdzenie}
\begin{proof}
Liniowość i surjektywność $q$ wynikają z konstrukcji
ilorazu, a ciągłość wynika z definicji topologii ilorazowej.
Jeżeli $O\subseteq X$ jest otwarty, to
\[
 q^{-1}(q(O))=O+M=\bigcup_{m\in M}(O+m)
\]
jest otwarty, ponieważ każde przesunięcie $O+m$
jest otwarte. Zatem $q(O)$ jest otwarty w $X/M$.

Wykażemy ciągłość dodawania w ilorazie. Odwzorowanie
$q\times q:X\times X\to(X/M)\times(X/M)$ jest ciągłą,
otwartą surjekcją: otwartość wynika z faktu, że obrazy
zbiorów bazowych $O_1\times O_2$ są produktami
otwartych $q(O_1)\times q(O_2)$. Każda otwarta surjekcja
ciągła jest odwzorowaniem ilorazowym. Mamy równość
$+_{X/M}\circ(q\times q)=q\circ+_X$, a prawa strona
jest ciągła. Z własności odwzorowania ilorazowego wynika
ciągłość $+_{X/M}$. Podobnie odwzorowanie
$\id_\K\times q:\K\times X\to\K\times(X/M)$ jest
otwartą surjekcją, więc z równości
$m_{X/M}\circ(\id_\K\times q)=q\circ m_X$
wynika ciągłość mnożenia skalarnego. Tym samym $X/M$
jest TVS, także gdy $M$ nie jest domknięta.

W dowolnej TVS singleton $\{0\}$ jest domknięty
wtedy i tylko wtedy, gdy przestrzeń jest Hausdorffa
(por. twierdzenie o domknięciu zera z \hyperref[sec:tvs]{rozdziału 1}).
Ponieważ $q^{-1}(\{0+M\})=M$, definicja topologii
ilorazowej pokazuje, że $\{0+M\}$ jest domknięty
w $X/M$ wtedy i tylko wtedy, gdy $M$ jest domknięta w $X$.
\end{proof}

\begin{uwaga}[Otoczenia w ilorazie]\label{rem:q-neighborhood}
Ponieważ $q$ jest otwarty, zbiory $q(U)$, gdzie
$U\in\mathcal N_X(0)$, tworzą bazę otoczeń zera
w $X/M$. Dokładniej: dla każdego otoczenia $W$
punktu $0+M$ zbiór $q^{-1}(W)$ jest otoczeniem zera,
a $q(q^{-1}(W))=W$; z drugiej strony każdy $q(U)$
jest otoczeniem zera, bo zawiera $q(O)$ dla pewnego
otwartego otoczenia $O\subseteq U$ zera.
\end{uwaga}

\begin{center}
\begin{tikzpicture}[x=0.95cm,y=0.82cm,>=Latex]
\draw[->] (-0.5,0)--(5.35,0) node[right] {$x$};
\draw[->] (0,-0.55)--(0,3.75) node[above] {$y$};
\foreach \yy in {0.8,1.7,2.65} {
  \draw[blue!55!black,thick] (0.2,\yy)--(4.7,\yy);
}
\filldraw[blue!55!black] (1.0,1.7) circle (1.7pt);
\filldraw[blue!55!black] (3.3,1.7) circle (1.7pt);
\node[anchor=west] at (4.75,1.7) {$x_0+M$};
\node[below] at (2.4,-0.10) {$M=\R\times\{0\}$};
\draw[->,very thick] (5.15,2.8) -- (6.4,2.8)
  node[midway,above] {$q$};
\draw[->] (7.2,-0.5)--(7.2,3.65);
\foreach \yy in {0.8,1.7,2.65} {
 \filldraw[blue!55!black] (7.2,\yy) circle (2.1pt);
}
\node[right] at (7.4,1.7) {$[x_0]$};
\node[below] at (7.2,-0.65) {$\R^2/M\simeq\R$};
\end{tikzpicture}

{\small Rysunek 2.1. Iloraz $\R^2/M$, gdzie $M$ jest osią poziomą.
Każda pozioma prosta jest jedną klasą równoważności
(prawa strona przedstawia zbiór klas).}
\end{center}

\begin{twierdzenie}[Własność uniwersalna ilorazu]\label{thm:quotient-universal}
Niech $X$ będzie TVS nad $\K$, $M\subseteq X$ podprzestrzenią,
a $q:X\to X/M$ rzutowaniem. Niech $Z$ będzie dowolną
przestrzenią topologiczną, a $g:X/M\to Z$ odwzorowaniem.
Wówczas $g$ jest ciągłe wtedy i tylko wtedy,
gdy $g\circ q$ jest ciągłe. Ponadto jeśli $Z$ jest TVS,
a $T:X\to Z$ jest liniowe i $M\subseteq\ker T$, to
istnieje dokładnie jedno liniowe $\widetilde T:X/M\to Z$
takie, że $T=\widetilde T\circ q$; jest ono ciągłe
dokładnie wtedy, gdy $T$ jest ciągłe.
\end{twierdzenie}
\begin{proof}
Pierwsza równoważność wynika z definicji topologii ilorazowej:
$g^{-1}(O)$ jest otwarty w $X/M$ dokładnie wtedy,
gdy $(g\circ q)^{-1}(O)$ jest otwarty w $X$, dla każdego
otwartego $O\subseteq Z$. Jeśli $M\subseteq\ker T$,
to wzór $\widetilde T(x+M)=T(x)$ jest dobrze określony:
z równości $x+M=y+M$ wynika $x-y\in M$, czyli $T(x)=T(y)$.
Wzór jest liniowy, a surjektywność $q$ zapewnia jednoznaczność.
Równoważność ciągłości stosujemy do $g=\widetilde T$.
\end{proof}

\begin{wniosek}[Pierwsze twierdzenie o izomorfizmie --- uwaga topologiczna]
Niech $T:X\to Y$ będzie ciągłym, liniowym odwzorowaniem
między przestrzeniami liniowo-topologicznymi Hausdorffa.
Wówczas odwzorowanie
$\widetilde T:X/\ker T\to T(X)$, $x+\ker T\mapsto T(x)$,
jest ciągłą bijekcją liniową, gdzie $T(X)$ ma topologię
indukowaną z $Y$. Jest izomorfizmem liniowo-topologicznym
wtedy i tylko wtedy, gdy $T:X\to T(X)$ jest odwzorowaniem
otwartym (przy topologii podprzestrzeni na obrazie).
\end{wniosek}
\begin{proof}
Jądro $T$ jest domknięte, ponieważ $T$ jest ciągłe,
a $Y$ jest Hausdorffa; dlatego iloraz jest Hausdorffa.
Istnienie i ciągłość $\widetilde T$ daje
\ref{thm:quotient-universal}; jest ona różnowartościowa,
bo $\ker T$ jest dokładnie jądrem $T$.
Jeśli $T$ jest otwarte na swój obraz, to
$\widetilde T(q(O))=T(O)$ jest otwarty dla każdego otwartego
$O\subseteq X$, a z otwartości $q$ i opisu zbiorów
otwartych w ilorazie wynika otwartość $\widetilde T$.
Odwrotnie, jeżeli $\widetilde T$ jest homeomorfizmem,
to $T=\widetilde T\circ q$ jest otwarte na obraz jako
złożenie dwóch odwzorowań otwartych.
\end{proof}

\begin{przyklad}[Dlaczego ciągła bijekcja nie wystarcza]
Niech $X=C([0,1],\R)$ z normą supremum, a $Y$ to ta sama
przestrzeń liniowa z normą $\|f\|_1=\int_{[0,1]}|f(t)|\,\mathrm d\lambda(t)$.
Identyczność $T:X\to Y$ jest ciągłą bijekcją liniową,
ale nie jest otwarta. Jej jądro jest zerowe, a mimo to
$X/\ker T\simeq X$ nie jest topologicznie izomorficzne z $Y$
poprzez odwzorowanie indukowane przez $T$. Dowód braku
ciągłości odwrotności podaliśmy w \hyperref[sec:tvs]{rozdziale 1}.
\end{przyklad}

\begin{przyklad}[Iloraz przez jądro wartościowania]
Niech $X=C([0,1],\K)$ z normą supremum, a
$E_0(f)=f(0)$. Jądro $M=\{f:f(0)=0\}$ jest domknięte.
Rzut $X/M\to\K$, $[f]\mapsto f(0)$, jest izomorfizmem
liniowo-topologicznym. Istotnie, dla $a\in\K$ stała
funkcja $s(a)(t)=a$ jest ciągłym liniowym prawym odwrotnym
do $E_0$, a odwrotne odwzorowanie na iloraz ma postać
$a\mapsto q(s(a))$. Jest więc ciągłe jako złożenie ciągłych.
\end{przyklad}

\subsection{Norma ilorazowa: ważny przypadek szczególny}

\begin{twierdzenie}[Norma w ilorazie]\label{thm:quotient-norm}
Niech $(X,\|\cdot\|)$ będzie przestrzenią unormowaną,
a $M\subseteq X$ domkniętą podprzestrzenią liniową. Wzór
\[
 \|x+M\|_{X/M}:=\inf_{m\in M}\|x-m\|=\operatorname{dist}(x,M)
\]
definiuje normę na $X/M$ generującą topologię ilorazową.
Dla każdego $r>0$ zachodzi
$q(B_X(0,r))=B_{X/M}(0,r)$.
Jeżeli dodatkowo $X$ jest przestrzenią Banacha,
to $X/M$ jest przestrzenią Banacha.
\end{twierdzenie}
\begin{proof}
Zastąpienie $x$ przez $x+m_0$, $m_0\in M$, nie zmienia
zbioru wartości $\|x+m_0-m\|$ dla $m\in M$, więc wzór
jest dobrze określony. Jednorodność wynika z liniowości $M$;
dla $\lambda=0$ obie strony są równe zeru. Dla dowolnych
$m_1,m_2\in M$ nierówność trójkąta w $X$ daje
$\|(x+y)-(m_1+m_2)\|\leq\|x-m_1\|+\|y-m_2\|$.
Biorąc infima, otrzymujemy nierówność trójkąta w ilorazie.
Jeśli $\operatorname{dist}(x,M)=0$, to $x\in\overline M=M$,
więc $x+M=0+M$. Zatem mamy normę.

Równość kul zachodzi, bo $\|x+M\|_{X/M}<r$ dokładnie
wtedy, gdy istnieje $m\in M$ z $\|x-m\|<r$.
Wobec otwartości $q$ wynika z tego, że obrazy kul
są otwarte. Dla dowolnego zbioru $W\subseteq X/M$
otwartego w topologii normowej jego przeciwobraz
$q^{-1}(W)$ jest otwarty, gdyż
$\|q(x)-q(y)\|_{X/M}\leq\|x-y\|$.
W drugą stronę, jeśli $W$ jest otwarty w topologii ilorazowej,
to dla każdego $[x]\in W$ istnieje $r>0$ takie, że
$B_X(x,r)\subseteq q^{-1}(W)$; wtedy
$B_{X/M}([x],r)=q(B_X(x,r))\subseteq W$.
Topologie są więc identyczne.

Niech teraz $X$ będzie zupełne, a $(u_n)$ ciągiem Cauchy'ego
w $X/M$. Wybierzmy podciąg $(u_{n_k})$ taki, że
$\|u_{n_{k+1}}-u_{n_k}\|_{X/M}<2^{-k-1}$.
Ustalmy przedstawiciela $x_1$ klasy $u_{n_1}$.
Indukcyjnie wybieramy przedstawiciela $x_{k+1}$ klasy
$u_{n_{k+1}}$ z $\|x_{k+1}-x_k\|<2^{-k}$.
Jest to możliwe, gdyż definicja normy ilorazowej
pozwala, po ustaleniu $x_k$, dobrać przedstawiciela
kolejnej klasy tak, aby odległość była dowolnie bliska
normie różnicy klas. Szereg odległości jest zbieżny,
więc $(x_k)$ jest ciągiem Cauchy'ego w $X$ i zbiega
do pewnego $x\in X$. Z ciągłości $q$ mamy
$u_{n_k}=q(x_k)\to q(x)$. Ponieważ cały $(u_n)$
jest ciągiem Cauchy'ego, również zbiega do $q(x)$.
\end{proof}

\begin{uwaga}
Gdy $M$ nie jest domknięta, ten sam wzór definiuje
jedynie seminormę: $\|x+M\|_{X/M}=0$ może zachodzić
dla niezerowej klasy. W szczególności taki wzór nie
może generować topologii Hausdorffa na $X/M$.
\end{uwaga}

\subsection{Topologie początkowe: najoszczędniejszy wybór topologii}

\begin{definicja}[Topologia początkowa]
Niech $S$ będzie zbiorem, a $f_i:S\to Y_i$, $i\in I$,
rodziną odwzorowań do przestrzeni topologicznych $Y_i$.
\emph{Topologią początkową} na $S$ względem $(f_i)$
nazywamy najmniejszą topologię, względem której wszystkie
$f_i$ są ciągłe. Jest ona generowana przez podbazę
$\{f_i^{-1}(O):i\in I,\ O\subseteq Y_i\text{ otwarty}\}$.
\end{definicja}

\begin{twierdzenie}[Początkowa topologia liniowa]\label{thm:initial-tvs}
Niech $X$ będzie przestrzenią liniową nad $\K$, $(Y_i)_{i\in I}$
rodziną przestrzeni liniowo-topologicznych nad $\K$,
a $T_i:X\to Y_i$ odwzorowaniami liniowymi.
Topologia początkowa względem $(T_i)$ czyni $X$
przestrzenią liniowo-topologiczną, niekoniecznie Hausdorffa.
Dla dowolnej przestrzeni topologicznej $Z$ odwzorowanie
$f:Z\to X$ jest ciągłe wtedy i tylko wtedy, gdy
$T_i\circ f$ jest ciągłe dla każdego $i\in I$.
Jeśli wszystkie $Y_i$ są Hausdorffa, to $X$ jest Hausdorffa
dokładnie wtedy, gdy $\bigcap_{i\in I}\ker T_i=\{0\}$.
\end{twierdzenie}
\begin{proof}
Ciągłość odwzorowań $T_i$ wynika z definicji topologii.
Ponieważ $T_i\circ+_X=+_{Y_i}\circ(T_i\times T_i)$,
a odwzorowania $T_i\times T_i$ są ciągłe względem
topologii początkowej na $X$, wszystkie złożenia
$T_i\circ+_X$ są ciągłe. Z uniwersalnej własności
topologii początkowej wynika ciągłość dodawania na $X$.
Tak samo $T_i\circ m_X=m_{Y_i}\circ(\id_\K\times T_i)$
daje ciągłość mnożenia skalarnego. Równoważność dotycząca
$f$ wynika wprost z preobrazów podbazowych zbiorów otwartych.

Jeśli $x\ne y$ i istnieje $i$ z $T_i(x)\ne T_i(y)$,
to dzięki własności Hausdorffa $Y_i$ istnieją rozłączne
otoczenia tych obrazów, których przeciwobrazy rozdzielają
$x$ i $y$. Jeśli natomiast $0\ne z\in\bigcap_i\ker T_i$,
to każde bazowe otoczenie zera zawiera $z$, a więc
$X$ nie jest Hausdorffa na mocy kryterium Hausdorffa
z \hyperref[sec:tvs]{rozdziału 1}.
\end{proof}

\begin{przyklad}[Produkty i podprzestrzenie jako przypadki topologii początkowej]
Topologia produktowa na $\prod_iX_i$ jest topologią
początkową względem rzutów $\pi_i$.
Topologia podprzestrzeni na $M\subseteq X$ jest
początkową względem włożenia $M\hookrightarrow X$.
W szczególności naturalna topologia na $C([0,1])$
zbieżności punktowej pochodzi od rodziny funkcjonałów
wartościowania $f\mapsto f(t)$, $t\in[0,1]$.
\end{przyklad}

\subsection{Topologie końcowe: przypadek delikatniejszy}

\begin{definicja}[Topologia końcowa zbiorów]
Niech $S$ będzie zbiorem, a $f_i:Y_i\to S$, $i\in I$,
rodziną odwzorowań z przestrzeni topologicznych.
\emph{Topologią końcową} na $S$ nazywamy najsilniejszą
topologię, przy której wszystkie $f_i$ są ciągłe:
\[
 \tau_{\mathrm{fin}}=\{O\subseteq S:
    f_i^{-1}(O)\text{ jest otwarty w }Y_i
    \text{ dla wszystkich }i\in I\}.
\]
Dla przestrzeni ilorazowej $X/M$ otrzymujemy właśnie
tę topologię przy jednym odwzorowaniu $q:X\to X/M$.
\end{definicja}

\begin{uwaga}[Niebezpieczeństwo: topologia końcowa nie musi być liniowa]
Jeżeli $S$ jest przestrzenią liniową, a $f_i$ są liniowe,
topologia końcowa w kategorii \emph{przestrzeni topologicznych}
nie musi czynić $S$ przestrzenią liniowo-topologiczną.
Przykład poniżej pokazuje, że nie można bez dowodu
przenosić własności ilorazu na dowolne rodziny map.
\end{uwaga}

\begin{przyklad}[Dwie osie i błędna topologia końcowa]\label{ex:final-cross}
Niech $S=\R^2$, $Y_1=Y_2=\R$ z topologiami zwykłymi,
a $f_1(t)=(t,0)$ i $f_2(t)=(0,t)$. W topologii końcowej
zbiór jednoelementowy $\{(1,1)\}$ jest otwarty,
gdyż oba jego przeciwobrazy są puste. Jednak $\{(0,0)\}$
nie jest otwarty: jego przeciwobraz przez $f_1$ jest
zbiorem $\{0\}\subseteq\R$. Przesunięcie o $(-1,-1)$
przenosi pierwszy zbiór na drugi, więc nie jest
homeomorfizmem. Tymczasem w każdej TVS wszystkie
przesunięcia są homeomorfizmami. A zatem ta topologia
końcowa nie jest topologią przestrzeni liniowo-topologicznej.
\end{przyklad}

\begin{definicja}[Końcowa topologia liniowo-topologiczna]
Niech $X$ będzie przestrzenią liniową nad $\K$, a
$T_i:X_i\to X$, $i\in I$, rodziną odwzorowań
liniowych z przestrzeni liniowo-topologicznych.
\emph{Końcową topologią liniowo-topologiczną}
nazywamy najsilniejszą topologię na $X$, która czyni $X$
przestrzenią liniowo-topologiczną i wszystkie $T_i$
odwzorowaniami ciągłymi. W tej definicji nie zakładamy
warunku Hausdorffa.
\end{definicja}

\begin{twierdzenie}[Istnienie końcowej topologii liniowo-topologicznej]
\label{thm:final-tvs}
Dla każdej przestrzeni liniowej $X$ nad $\K$ i rodziny
liniowych odwzorowań $T_i:X_i\to X$ z TVS istnieje
końcowa topologia liniowo-topologiczna. Ma ona następującą
własność: jeżeli $Z$ jest TVS, a $S:X\to Z$ odwzorowaniem
liniowym, to $S$ jest ciągłe dokładnie wtedy, gdy
$S\circ T_i$ jest ciągłe dla wszystkich $i\in I$.
\end{twierdzenie}
\begin{proof}
Rozpatrzmy rodzinę $\mathfrak T$ wszystkich topologii
liniowo-topologicznych na $X$, względem których wszystkie
$T_i$ są ciągłe. Nie jest ona pusta: zawiera topologię
trywialną. Niech $\tau$ będzie topologią generowaną
przez sumę wszystkich topologii z $\mathfrak T$.
Aby nie pomylić topologii na dziedzinach działań,
ustalmy otwarty w pewnej $\sigma\in\mathfrak T$ zbiór
$O\subseteq X$. Wówczas
$+_X^{-1}(O)$ jest otwarty w $(X,\sigma)\times(X,\sigma)$,
a więc również w $(X,\tau)\times(X,\tau)$.
Analogicznie $m_X^{-1}(O)$ jest otwarty w
$\K\times(X,\tau)$. Ponieważ takie zbiory $O$
tworzą podbazę $\tau$, oba działania są łącznie ciągłe.
Dla każdego $i$ przeciwobrazy zbiorów podbazowych przez
$T_i$ są otwarte w $X_i$; ich skończone przecięcia
i dowolne sumy również. Tak więc $T_i:X_i\to(X,\tau)$
są ciągłe. Każda dopuszczalna topologia jest zawarta
w $\tau$, co dowodzi maksymalności. W szczególności
w tej konstrukcji nie trzeba zakładać, że obrazy $T_i$
generują algebraicznie $X$.

Jeśli $S:X\to Z$ jest liniowe, a wszystkie $S\circ T_i$
są ciągłe, to topologia początkowa $\sigma_S$ na $X$
indukowana przez $S$ jest topologią liniowo-topologiczną
(twierdzenie~\ref{thm:initial-tvs}); wszystkie $T_i$
są ciągłe do $(X,\sigma_S)$, a zatem $\sigma_S\subseteq\tau$.
Z definicji $\sigma_S$ wynika, że $S:(X,\tau)\to Z$
jest ciągłe. W drugą stronę wystarczy składać odwzorowania
ciągłe.
\end{proof}

\begin{uwaga}[Wersja lokalnie wypukła]
Jeżeli wszystkie $X_i$ są lokalnie wypukłe, można
analogicznie rozważać \emph{najsilniejszą lokalnie wypukłą}
topologię na $X$, przy której $T_i$ są ciągłe.
Będzie to szczególnie ważne dla przestrzeni funkcji
testowych $\mathcal D(\Omega)$; tam użyjemy granicy
induktywnej w kategorii przestrzeni lokalnie wypukłych.
Nie należy bez dodatkowych argumentów utożsamiać
tej topologii z końcową topologią zwykłych przestrzeni
topologicznych.
\end{uwaga}

\subsection{Sumy proste: produkt skończony i nieskończony}

\begin{definicja}[Algebraiczna suma prosta]
Dla rodziny przestrzeni liniowych $(X_i)_{i\in I}$
przez $\bigoplus_{i\in I}X_i$ oznaczamy przestrzeń
rodzin $(x_i)\in\prod_iX_i$, w których tylko skończenie
wiele współrzędnych jest niezerowych. Naturalne włożenia
$j_i:X_i\to\bigoplus_iX_i$ umieszczają element w $i$-tej
współrzędnej, a na pozostałych wpisują zero.
\end{definicja}

\begin{twierdzenie}[Skończone sumy proste]
Niech $X_1,\dots,X_n$ będą przestrzeniami liniowo-topologicznymi.
Na $X_1\oplus\dots\oplus X_n$ utożsamionym algebraicznie
z $X_1\times\dots\times X_n$ topologia produktowa jest
końcową topologią liniowo-topologiczną względem włożeń $j_i$.
\end{twierdzenie}
\begin{proof}
Włożenia $j_i$ są ciągłe w topologii produktowej.
Jeśli $Z$ jest TVS i $S:X_1\oplus\dots\oplus X_n\to Z$
jest liniowe, a każde $S\circ j_i$ jest ciągłe,
to dla $(x_1,\dots,x_n)$ mamy
$S(x_1,\dots,x_n)=\sum_{i=1}^n(S\circ j_i)(x_i)$.
Skończona suma ciągłych odwzorowań jest ciągła.
Z uniwersalnej własności końcowej topologii
liniowo-topologicznej (twierdzenie~\ref{thm:final-tvs})
wynika teza.
\end{proof}

\begin{przyklad}[Przestrzeń ciągów o skończonym nośniku]\label{ex:direct-sum}
Niech $\varphi=\bigoplus_{n\geq1}\K e_n$ oznacza
przestrzeń ciągów o skończonym nośniku. Topologia
indukowana z produktu $\K^{\N}$ i \emph{topologia
lokalnie wypukłej sumy prostej} są różne.
Dla jasności: \emph{seminorma} to funkcja $p:X\to[0,\infty)$
spełniająca $p(\lambda x)=|\lambda|p(x)$ i
$p(x+y)\leq p(x)+p(y)$; może przyjmować wartość zero
na niezerowych wektorach. Rodzina seminorm definiuje topologię,
w której bazę otoczeń zera stanowią skończone przecięcia
zbiorów $\{x:p(x)<\varepsilon\}$, $\varepsilon>0$.
Jest to topologia \emph{lokalnie wypukła}: zero posiada bazę
wypukłych otoczeń. Teorię seminorm rozwiniemy w \hyperref[sec:lokalnie-wypukle]{rozdziale 5}.

Na $\varphi$ rozważmy rodzinę seminorm
\[
 p_a(x)=\sum_{n\geq1}a_n|x_n|,
 \qquad a=(a_n)_{n\geq1}\in[0,\infty)^{\N}.
\]
Każda suma jest skończona dla ustalonego $x\in\varphi$.
Wszystkie $p_a$ są seminormami, a ich rodzina
rozdziela punkty, więc określają lokalnie wypukłą
topologię Hausdorffa. Niech $Z$ będzie przestrzenią
liniowo-topologiczną z topologią zadaną przez rodzinę
seminorm. Dla liniowego $S:\varphi\to Z$ oraz
dowolnej seminormy $p$ z tej rodziny zachodzi
\[
 p(Sx)\leq\sum_{n\geq1}|x_n|\,p(S e_n)
         =p_a(x),\qquad a_n=p(S e_n).
\]
Wynika stąd ciągłość każdego liniowego $S:\varphi\to Z$,
gdyż każda podbazowa nierówność $p(Sx)<\varepsilon$
przyjmuje postać otwartego warunku na $\varphi$.
W \hyperref[sec:lokalnie-wypukle]{rozdziale 5} udowodnimy, że topologie wszystkich
przestrzeni lokalnie wypukłych można opisać seminormami;
wtedy powyższy argument da pełną własność uniwersalną
sumy prostej \emph{w kategorii przestrzeni lokalnie wypukłych}.
Nie dowodzi on analogicznej własności względem wszystkich
przestrzeni liniowo-topologicznych bez lokalnej wypukłości.

W topologii produktowej odziedziczonej z $\K^{\N}$
ciąg $e_n$ zbiega do zera. W topologii sumy prostej
nie zbiega: $p_{(1,1,\dots)}(e_n)=1$ dla każdego $n$.
To kolejny przykład, w którym ten sam zbiór liniowy
niesie dwie naturalne, ale istotnie różne topologie.
\end{przyklad}

\subsection{Iloraz Hausdorffa dowolnej TVS}

\begin{twierdzenie}[Hausdorffyzacja]\label{thm:hausdorffization}
Niech $X$ będzie przestrzenią liniowo-topologiczną
(niekoniecznie Hausdorffa), a $N=\overline{\{0\}}$.
Wówczas $N$ jest domkniętą podprzestrzenią, $X/N$
jest przestrzenią liniowo-topologiczną Hausdorffa
oraz dla każdego ciągłego liniowego $T:X\to Y$,
gdzie $Y$ jest TVS Hausdorffa, istnieje dokładnie
jedno ciągłe liniowe $\widetilde T:X/N\to Y$
spełniające $T=\widetilde T\circ q$.
\end{twierdzenie}
\begin{proof}
Domkniętość i liniowość $N$ wykazano w \hyperref[sec:tvs]{rozdziale 1}. Wobec
\ref{thm:quotient-tvs} iloraz $X/N$ jest Hausdorffa.
Ponieważ $T$ jest ciągłe i $Y$ jest Hausdorffa,
$\ker T=T^{-1}(\{0\})$ jest domknięte i zawiera zero,
a więc $N\subseteq\ker T$. Z własności uniwersalnej
ilorazu (twierdzenie~\ref{thm:quotient-universal})
istnieje jednoznaczne ciągłe $\widetilde T$.
\end{proof}

\subsection{Podsumowanie: kiedy topologia jest właściwa?}

W konstrukcjach tego rozdziału występują cztery pojęcia:
\begin{enumerate}[label=\textup{(\roman*)}]
\item topologia podprzestrzeni jest \emph{najsłabszą} topologią,
      przy której włożenie do większej przestrzeni jest ciągłe;
\item topologia produktowa jest \emph{najsłabszą} topologią,
      przy której wszystkie rzuty współrzędne są ciągłe;
\item topologia ilorazowa jest \emph{najsilniejszą} topologią,
      przy której rzutowanie ilorazowe jest ciągłe; w tym
      przypadku pozostaje topologią przestrzeni liniowo-topologicznej;
\item ogólna topologia końcowa dla rodziny odwzorowań liniowych
      nie musi być zgodna z działaniami liniowymi; trzeba
      rozróżniać wersję topologiczną, liniowo-topologiczną
      i lokalnie wypukłą.
\end{enumerate}

Warto zapamiętać dwie przeszkody: iloraz przez podprzestrzeń
niedomkniętą nie jest Hausdorffa, a ciągła liniowa bijekcja
nie musi być izomorfizmem topologicznym.
W następnym rozdziale zajmiemy się sieciami, strukturą jednostajną
przestrzeni liniowo-topologicznych, zupełnością oraz uzupełnieniami.


% Wersja źródłowa: audyt II, jednostka 03
\bigskip
\setcounter{section}{2}
\section{Sieci, struktura jednostajna i zupełność przestrzeni liniowo-topologicznych}\label{sec:zupelnosc}

\subsection{Dlaczego ciągi nie wystarczają?}

W przestrzeni metrycznej ciągi wystarczają do charakteryzowania domknięcia,
ciągłości i zupełności. W dowolnej przestrzeni topologicznej jest inaczej:
zbiór może zawierać granice wszystkich swoich zbieżnych ciągów, a mimo to
nie być domknięty. Dlatego zamiast samych ciągów będziemy posługiwać się
\emph{sieciami}. Później pojęcie sieci Cauchy'ego pozwoli określić
zupełność bez wybierania normy czy metryki.

Dla wygody przypomnijmy: $U\subseteq X$ jest otoczeniem
zera, jeżeli zawiera otwarty zbiór zawierający zero;
jest zrównoważony, gdy $\lambda U\subseteq U$
dla $|\lambda|\leq1$.
W tym rozdziale $\K\in\{\R,\C\}$, a TVS jest, jeśli nie zaznaczono
inaczej, przestrzenią liniowo-topologiczną Hausdorffa. Otoczenie punktu
nie musi być otwarte. Gdy piszemy $U\in\mathcal N_X(0)$, mamy na myśli
otoczenie zera w $X$.

\begin{definicja}[Zbiór skierowany i sieć]\label{def:net}
\emph{Zbiorem skierowanym} nazywamy niepusty zbiór $D$ z relacją
refleksyjną i przechodnią $\preceq$, w którym dla dowolnych $d,e\in D$
istnieje $f\in D$ takie, że $d\preceq f$ i $e\preceq f$.
\emph{Sieć} w zbiorze $X$ to odwzorowanie $D\to X$,
$d\mapsto x_d$, gdzie $D$ jest zbiorem skierowanym.
Jeżeli $X$ jest przestrzenią topologiczną, mówimy, że $x_d\to x\in X$,
gdy dla każdego otoczenia $U$ punktu $x$ istnieje $d_0\in D$ takie,
że $x_d\in U$ dla wszystkich $d\succeq d_0$.
\end{definicja}

Każdy ciąg jest siecią, w której $D=\N$ z naturalnym porządkiem.
Innym przykładem jest rodzina $\mathcal F(I)$ wszystkich skończonych
podzbiorów ustalonego zbioru $I$, skierowana przez inkluzję.
Jeśli $X$ jest przestrzenią topologiczną i $x\in X$, to rodzina
otwartych otoczeń $x$, skierowana przez \emph{odwrotną} inkluzję,
również jest zbiorem skierowanym: wspólnym późniejszym indeksem
jest przecięcie dwóch otoczeń.

\begin{twierdzenie}[Domknięcie i ciągłość przez sieci]\label{thm:nets-topology}
Niech $X,Y$ będą przestrzeniami topologicznymi, $A\subseteq X$,
$x\in X$ oraz $f:X\to Y$. Wtedy:
\begin{enumerate}[label=\textup{(\roman*)}]
\item $x\in\overline A$ dokładnie wtedy, gdy istnieje sieć
      $(a_d)_{d\in D}$ elementów $A$ zbieżna do $x$;
\item $f$ jest ciągłe w $x$ dokładnie wtedy, gdy dla każdej sieci
      $x_d\to x$ zachodzi $f(x_d)\to f(x)$.
\end{enumerate}
\end{twierdzenie}
\begin{proof}
Jeśli $a_d\in A$ i $a_d\to x$, każde otoczenie $x$ zawiera element $A$,
zatem $x\in\overline A$. Odwrotnie, niech $x\in\overline A$.
Skierujmy rodzinę otwartych otoczeń $x$ przez odwrotną inkluzję,
a dla każdego takiego otoczenia $U$ wybierzmy $a_U\in A\cap U$.
Dla danego otoczenia $V$ wszystkie indeksy $U\subseteq V$
spełniają $a_U\in V$. Stąd $a_U\to x$.

Jeśli $f$ jest ciągłe w $x$, zachowuje granice sieci zbieżnych do $x$.
Gdyby $f$ nie było ciągłe w $x$, istniałoby otwarte otoczenie $W$
punktu $f(x)$, dla którego w każdym otwartym otoczeniu $U$ punktu $x$
można wybrać $x_U\in U$ z $f(x_U)\notin W$. Wówczas $x_U\to x$,
ale $f(x_U)$ nie zbiega do $f(x)$, co daje sprzeczność.
\end{proof}

\begin{przyklad}[Domknięcie niewykrywalne przez ciągi]\label{ex:sigma-net}
Niech $I$ będzie zbiorem nieprzeliczalnym, a
$X=\K^I$ z topologią produktową. Oznaczmy przez
\[
 \Sigma(I)=\{x\in\K^I:\supp x=\{i\in I:x_i\ne0\}
                 \text{ jest przeliczalny}\}.
\]
Jest to podprzestrzeń liniowa, bo suma dwóch zbiorów przeliczalnych
jest przeliczalna. Jest gęsta w $X$: dla $x\in X$ sieć
$x_F=\sum_{i\in F}x_i e_i$, gdzie $F$ przebiega skończone podzbiory
$I$ uporządkowane przez inkluzję, zbiega współrzędnie do $x$.
Nie jest jednak równa $X$: wektor $\mathbf1=(1)_{i\in I}$
ma nieprzeliczalny nośnik.

Co istotne, $\Sigma(I)$ jest \emph{ciągowo domknięta} w $X$.
Jeśli $x_n\in\Sigma(I)$, to zbiór $J=\bigcup_n\supp x_n$
jest przeliczalny. Jeżeli $x_n\to x$ współrzędnie, to $x_i=0$
dla $i\notin J$, a więc $x\in\Sigma(I)$. Zatem żaden ciąg
z $\Sigma(I)$ nie może zbiegać do $\mathbf1$, choć pewna sieć może.
\end{przyklad}

\subsection{Filtry i ich zbieżność}

Filtry są alternatywnym językiem opisu zbieżności, szczególnie
wygodnym w konstrukcji uzupełnień.

\begin{definicja}[Filtr]\label{def:filter}
\emph{Filtrem} na niepustym zbiorze $X$ nazywamy niepustą rodzinę
$\mathcal F$ podzbiorów $X$ taką, że:
\begin{enumerate}[label=\textup{(\roman*)}]
\item $\varnothing\notin\mathcal F$;
\item $A,B\in\mathcal F$ pociąga $A\cap B\in\mathcal F$;
\item $A\in\mathcal F$ i $A\subseteq B\subseteq X$ pociąga
      $B\in\mathcal F$.
\end{enumerate}
Filtr $\mathcal F$ \emph{zbiega} do $x$ w przestrzeni topologicznej $X$,
gdy $\mathcal N_X(x)\subseteq\mathcal F$.
\end{definicja}

Z każdej sieci $(x_d)_{d\in D}$ powstaje filtr jej \emph{ogonów}:
jest to rodzina wszystkich zbiorów zawierających pewien ogon
$\{x_d:d\succeq d_0\}$. Sieć zbiega do $x$ dokładnie wtedy,
gdy filtr ogonów zbiega do $x$. Jeśli filtr $\mathcal F$ jest zadany,
można w drugą stronę utworzyć sieć: indeksami są pary $(A,a)$,
gdzie $A\in\mathcal F$ i $a\in A$, a
$(A,a)\preceq(B,b)$, gdy $B\subseteq A$.
Wartością sieci w $(A,a)$ jest $a$. Przecięcie dwóch elementów filtra
jest niepuste, więc podana relacja jest skierowana. Ta sieć zbiega
do $x$ dokładnie wtedy, gdy $\mathcal F$ zbiega do $x$.

\subsection{Kanoniczna struktura jednostajna TVS}

W przestrzeni metrycznej stwierdzenie, że $x$ i $y$ są blisko,
wyrażamy nierównością $d(x,y)<\varepsilon$. W TVS rolę tej nierówności
pełni warunek $x-y\in U$, gdzie $U$ jest otoczeniem zera.

\begin{definicja}[Otoczenie przekątnej i struktura jednostajna]
Niech $X$ będzie przestrzenią liniowo-topologiczną. Dla otoczenia
$U$ zera określamy
\[
 E_U=\{(x,y)\in X\times X:x-y\in U\}.
\]
Rodzinę zbiorów $W\subseteq X\times X$ zawierających przynajmniej
jeden $E_U$ nazywamy \emph{kanoniczną strukturą jednostajną} $X$.
Jej elementy to \emph{otoczenia przekątnej} (ang. \emph{entourages}).
\end{definicja}

\begin{twierdzenie}[Własności kanonicznej struktury jednostajnej]
\label{thm:uniform-tvs}
Niech $X$ będzie TVS. Rodzina $\{E_U:U\in\mathcal N_X(0)\}$
spełnia następujące warunki:
\begin{enumerate}[label=\textup{(\roman*)}]
\item każde $E_U$ zawiera przekątną $\Delta_X=\{(x,x):x\in X\}$;
\item dla otoczeń $U,V$ zachodzi $E_{U\cap V}=E_U\cap E_V$;
\item dla każdego $U$ istnieje $V$ takie, że $E_V^{-1}\subseteq E_U$;
\item dla każdego $U$ istnieje $V$ takie, że
      $E_V\circ E_V\subseteq E_U$, gdzie
      $E\circ F=\{(x,z):\exists y\ ((x,y)\in E,(y,z)\in F)\}$.
\end{enumerate}
Topologia wyznaczona przez tę strukturę jednostajną jest oryginalną
topologią $X$: przekrój $E_U[x]=\{y:(x,y)\in E_U\}$
tworzy wraz z $U$ bazę otoczeń $x$.
\end{twierdzenie}
\begin{proof}
Własności (i) i (ii) są bezpośrednie. Do (iii) wybieramy symetryczne
otoczenie $V\subseteq U$; wtedy $(x,y)\in E_V^{-1}$ oznacza
$y-x\in V$, zatem $x-y\in-V=V\subseteq U$.
Dla (iv) wybieramy $V$ z $V+V\subseteq U$. Jeżeli
$x-y\in V$ i $y-z\in V$, to $x-z\in U$.
Wreszcie $E_U[x]=x-U$, a symetryczne otoczenia zera tworzą bazę,
więc otrzymujemy dokładnie topologię $X$.
\end{proof}

\begin{uwaga}
Struktura jednostajna pozwala definiować \emph{jednostajną
ciągłość} również bez metryki: odwzorowanie $f:X\to Y$
między przestrzeniami z ustalonymi strukturami jednostajnymi
jest jednostajnie ciągłe, gdy dla każdego otoczenia przekątnej
$W$ w $Y$ istnieje otoczenie przekątnej $V$ w $X$ takie,
że $(x,y)\in V$ pociąga $(f(x),f(y))\in W$.
Przestrzenie liniowo-topologiczne mają kanoniczne struktury
jednostajne; dowolna przestrzeń topologiczna nie musi ich posiadać.
\end{uwaga}

\begin{twierdzenie}[Ciągłość liniowa implikuje jednostajną]
\label{thm:linear-uniform}
Jeśli $X,Y$ są TVS, a $T:X\to Y$ jest liniowe i ciągłe,
to $T$ jest jednostajnie ciągłe względem kanonicznych
struktur jednostajnych.
\end{twierdzenie}
\begin{proof}
Niech $W$ będzie otoczeniem zera w $Y$. Z ciągłości $T$ w zerze
wybieramy otoczenie $U$ zera w $X$ takie, że $T(U)\subseteq W$.
Jeśli $x-y\in U$, to $T(x)-T(y)=T(x-y)\in W$, co jest
żądaną jednostajną ciągłością.
\end{proof}

\subsection{Sieci Cauchy'ego i zupełność}

\begin{definicja}[Sieć Cauchy'ego, filtr Cauchy'ego i zupełność]
\label{def:tvs-complete}
Niech $X$ będzie TVS. Sieć $(x_d)_{d\in D}$ w $X$ nazywamy
\emph{siecią Cauchy'ego}, jeżeli dla każdego $U\in\mathcal N_X(0)$
istnieje $d_0\in D$ takie, że
$ x_d-x_e\in U $ dla wszystkich $d,e\succeq d_0$.
Filtr $\mathcal F$ na $X$ nazywamy \emph{filtrem Cauchy'ego},
jeżeli dla każdego $U\in\mathcal N_X(0)$ istnieje
$A\in\mathcal F$ takie, że $A-A\subseteq U$.
Przestrzeń $X$ jest \emph{zupełna}, jeżeli każda sieć Cauchy'ego
w $X$ ma granicę w $X$. Jest \emph{ciągowo zupełna},
jeżeli każdy jej ciąg Cauchy'ego ma granicę w $X$.
\end{definicja}

\begin{twierdzenie}[Równoważność opisu sieciowego i filtrowego]
\label{thm:filter-cauchy}
Niech $X$ będzie TVS. Sieć jest siecią Cauchy'ego dokładnie wtedy,
gdy jej filtr ogonów jest filtrem Cauchy'ego. Ponadto $X$ jest
zupełna dokładnie wtedy, gdy każdy filtr Cauchy'ego na $X$
jest zbieżny.
\end{twierdzenie}
\begin{proof}
Jeśli sieć jest Cauchy'ego, to dla otoczenia $U$ wybieramy
ogon $A$, w którym różnice dowolnych dwóch wyrazów należą do $U$;
wtedy $A-A\subseteq U$ i $A$ należy do filtra ogonów.
W drugą stronę, jeśli filtr ogonów jest Cauchy'ego, wybieramy
$A$ z $A-A\subseteq U$ oraz ogon zawarty w $A$. Różnice
wyrazów tego ogona należą do $U$.

Załóżmy, że $X$ jest zupełna, i niech $\mathcal F$ będzie filtrem
Cauchy'ego. Rozważmy związaną z nim sieć indeksowaną parami
$(A,a)$, opisaną po definicji~\ref{def:filter}. Dla otoczenia
$U$ wybieramy $A\in\mathcal F$ z $A-A\subseteq U$.
Wszystkie wyrazy sieci o indeksach późniejszych niż $(A,a)$
należą do $A$, więc sieć jest Cauchy'ego. Niech $x$ będzie jej
granicą. Dla dowolnego otoczenia $V$ punktu $x$ istnieje
$(A,a)$ takie, że wszystkie późniejsze wyrazy należą do $V$.
W szczególności wszystkie punkty $b\in A$ (otrzymane
z indeksów $(A,b)$) należą do $V$. Stąd $A\subseteq V$,
a więc $V\in\mathcal F$, co oznacza zbieżność filtra.
Odwrotna implikacja wynika z zastosowania założenia o filtrach
do filtra ogonów dowolnej sieci Cauchy'ego.
\end{proof}

\begin{stwierdzenie}[Granice i jednoznaczność]
\label{prop:complete-basic}
Niech $X$ będzie TVS. Każda sieć zbieżna w $X$ jest Cauchy'ego.
Każda sieć Cauchy'ego ma co najwyżej jedną granicę.
Jeżeli sieć Cauchy'ego ma podsieć zbieżną do $x\in X$, to sama
zbiega do $x$.
\end{stwierdzenie}
\begin{proof}
Jeśli $x_d\to x$ i $U$ jest otoczeniem zera, wybieramy
symetryczne otoczenie $V$ z $V+V\subseteq U$.
Dla późnych indeksów $x_d-x\in V$, a zatem
$x_d-x_e\in V-V\subseteq U$. Jednoznaczność granicy
wynika z własności Hausdorffa.

Załóżmy, że sieć jest Cauchy'ego, a jej podsieć zbiega do $x$.
Dla otoczenia $U$ wybieramy symetryczne $V$ z $V+V\subseteq U$.
Istnieje $d_0$ takie, że $x_d-x_e\in V$ dla $d,e\succeq d_0$.
Zbieżna podsieć ma wyraz $x_{e}$ o indeksie $e\succeq d_0$
i z $x_e-x\in V$. Wówczas dla każdego $d\succeq d_0$
mamy $x_d-x\in V+V\subseteq U$.
\end{proof}

\begin{uwaga}[Podsieć]
Podsieć sieci $(x_d)_{d\in D}$ rozumiemy jako sieć
$(x_{\phi(e)})_{e\in E}$, gdzie $E$ jest skierowany,
a $\phi:E\to D$ ma własność: dla każdego $d_0\in D$
istnieje $e_0\in E$ takie, że $\phi(e)\succeq d_0$
dla wszystkich $e\succeq e_0$. Ta własność, zwana
\emph{współkońcowością}, zapewnia przenoszenie warunków
prawdziwych od pewnego indeksu na podsieć.
\end{uwaga}

\begin{twierdzenie}[Zupełność ciągowa w przestrzeniach metryzowalnych]
\label{thm:metric-complete}
Niech $X$ będzie metryzowalną przestrzenią liniowo-topologiczną
Hausdorffa. Wtedy $X$ jest zupełna dokładnie wtedy,
gdy jest ciągowo zupełna.
\end{twierdzenie}
\begin{proof}
Jedna implikacja jest oczywista. W drugą stronę załóżmy,
że wszystkie ciągi Cauchy'ego są zbieżne. Niech $(x_d)$
będzie dowolną siecią Cauchy'ego. Wybierzmy malejącą
przeliczalną bazę symetrycznych otoczeń zera $(U_n)$
taką, że $U_{n+1}+U_{n+1}\subseteq U_n$.
Dla każdego $n$ istnieje $d_n$ takie, że
$x_d-x_e\in U_n$ dla $d,e\succeq d_n$.
Rekurencyjnie wybierzmy indeksy $e_n$ spełniające
$e_n\succeq d_1,\ldots,d_n,e_{n-1}$.
Wówczas $(x_{e_n})$ jest ciągiem Cauchy'ego: dla
$m,n\geq N$ oba indeksy są późniejsze od $d_N$,
więc $x_{e_n}-x_{e_m}\in U_N$.
Niech $x_{e_n}\to x$. Dla zadanego otoczenia $W$ zera
wybierzmy $N$ tak duże, że $U_N+U_N\subseteq W$.
Dla każdego $d\succeq d_N$ i $n\geq N$ mamy
$x_d-x_{e_n}\in U_N$. Wybierając $n$ dostatecznie duże,
aby $x_{e_n}-x\in U_N$, otrzymujemy $x_d-x\in W$.
Zatem oryginalna sieć zbiega do $x$.
\end{proof}

\subsection{Operacje na przestrzeniach zupełnych}

\begin{twierdzenie}[Podprzestrzenie i produkty zupełne]
\label{thm:closed-complete}
Niech $X$ będzie TVS Hausdorffa.
\begin{enumerate}[label=\textup{(\roman*)}]
\item Jeśli $X$ jest zupełna, to każda domknięta
      podprzestrzeń liniowa $M\subseteq X$ jest zupełna.
\item Jeśli $M\subseteq X$ jest podprzestrzenią liniową zupełną
      w topologii indukowanej, to $M$ jest domknięta w $X$.
\item Dla rodziny TVS Hausdorffa $(X_i)_{i\in I}$ produkt
      $\prod_{i\in I}X_i$ jest zupełny wtedy i tylko wtedy,
      gdy wszystkie $X_i$ są zupełne.
\end{enumerate}
\end{twierdzenie}
\begin{proof}
Sieć Cauchy'ego w $M$ jest Cauchy'ego w $X$.
Jeżeli $X$ jest zupełna, ma ona granicę w $X$; domkniętość
$M$ zapewnia, że granica należy do $M$. To daje (i).
Jeśli $x\in\overline M$, twierdzenie~\ref{thm:nets-topology}
daje sieć w $M$ zbieżną do $x$ w $X$. Jest ona Cauchy'ego
w $M$, więc z założenia ma granicę $m\in M$.
Jednoznaczność granicy w $X$ daje $x=m$, co dowodzi (ii).

Niech teraz wszystkie $X_i$ będą zupełne, a $(x_d)$
siecią Cauchy'ego w produkcie. Dla każdego $i$ sieć
$(\pi_i x_d)$ jest Cauchy'ego w $X_i$, gdyż rzuty
są ciągłe i liniowe, zatem jednostajnie ciągłe.
Istnieją więc $x_i\in X_i$ takie, że $\pi_i x_d\to x_i$.
Z uniwersalnej własności produktu $x_d\to(x_i)_{i\in I}$.
W drugą stronę, dla każdego $j\in I$ kopia $X_j$
złożona z wektorów o zerowych współrzędnych poza $j$
jest domkniętą podprzestrzenią produktu (pozostałe czynniki
są Hausdorffa). Jeśli produkt jest zupełny, część (i)
pokazuje zupełność $X_j$.
\end{proof}

\begin{przyklad}[Sekwencyjnie zupełna, ale niezupełna TVS]
\label{ex:seq-not-complete}
Niech $I$ będzie nieprzeliczalny, $X=\K^I$ z topologią
produktową, a $M=\Sigma(I)$ z przykładu~\ref{ex:sigma-net}.
Produkt $X$ jest zupełny, ponieważ $\K$ jest zupełne.
Każdy ciąg Cauchy'ego $(x_n)$ w $M$ jest Cauchy'ego w $X$,
więc zbiega tam do pewnego $x\in X$. Nośnik $x$ zawiera
się w przeliczalnej sumie $\bigcup_n\supp x_n$, toteż $x\in M$.
Zatem $M$ jest ciągowo zupełna. Nie jest zupełna,
ponieważ jest gęstą właściwą podprzestrzenią zupełnej
przestrzeni $X$ i nie jest w niej domknięta; do
niezupełności wystarcza też część (ii) twierdzenia
\ref{thm:closed-complete}.

Konkretnym świadkiem jest sieć $x_F=\sum_{i\in F}e_i$,
$F\in\mathcal F(I)$: jest Cauchy'ego w $M$, gdyż
zbiega do $\mathbf1$ w $X$, lecz nie zbiega w $M$.
\end{przyklad}

\subsection{Uzupełnienie TVS: konstrukcja bez normy}

Chcemy skonstruować przestrzeń zupełną $\widehat X$
i ciągłe liniowe zanurzenie $j:X\to\widehat X$ o gęstym obrazie,
bez zakładania lokalnej wypukłości czy metryzowalności.
Poniższy dowód jest konstrukcją opartą na sieciach Cauchy'ego.

\begin{twierdzenie}[Istnienie i własność uniwersalna uzupełnienia]
\label{thm:tvs-completion}
Niech $X$ będzie przestrzenią liniowo-topologiczną Hausdorffa
nad $\K$. Istnieją zupełna przestrzeń liniowo-topologiczna
Hausdorffa $\widehat X$ oraz liniowe zanurzenie topologiczne
$j:X\to\widehat X$ o gęstym obrazie. Dla każdej zupełnej
przestrzeni liniowo-topologicznej Hausdorffa $Y$ oraz każdego
ciągłego liniowego $T:X\to Y$ istnieje dokładnie jedno
ciągłe liniowe $\widehat T:\widehat X\to Y$ takie,
że $\widehat T\circ j=T$. Para $(\widehat X,j)$ jest
jednoznaczna z dokładnością do jedynego izomorfizmu
liniowo-topologicznego zachowującego osadzenie $X$.
\end{twierdzenie}
\begin{proof}
\emph{Krok 1: klasy sieci Cauchy'ego.}
Aby otrzymać rzeczywisty zbiór klas, a nie klasę właściwą,
ustalmy rodzinę $\mathcal U$ wszystkich otwartych zrównoważonych
otoczeń zera w $X$ oraz zbiór skierowany
$D_*=\mathcal F(\mathcal U)$ jej skończonych podzbiorów,
uporządkowany przez inkluzję. Każda sieć Cauchy'ego $(x_d)_{d\in D}$
ma równoważną sieć indeksowaną przez $D_*$: dla każdego
$U\in\mathcal U$ wybierzmy $d_U$ takie, że
$x_d-x_e\in U$ dla $d,e\succeq d_U$; dla skończonego
$F\subseteq\mathcal U$ wybierzmy $d_F$ późniejszy od wszystkich
$d_U$ z $U\in F$. Pokażemy, że sieć $(x_{d_F})_{F\in D_*}$
jest Cauchy'ego i równoważna pierwotnej. Dla otoczenia
$W$ zera wybierzmy $U\in\mathcal U$ z $U\subseteq W$.
Gdy $F,G\supseteq\{U\}$, mamy $d_F,d_G\succeq d_U$,
a więc $x_{d_F}-x_{d_G}\in U\subseteq W$.
Podobnie dla $F\supseteq\{U\}$ i $d\succeq d_U$
zachodzi $x_{d_F}-x_d\in W$; jest to dokładnie
wymagana równoważność dwóch sieci. Nie wymagamy,
aby wybór $F\mapsto d_F$ był monotoniczny: podane
oszacowanie zapewnia wszystkie potrzebne własności.
W konstrukcji $\widehat X$ wystarczy zatem uwzględniać
sieci o ustalonym zbiorze indeksów $D_*$. Tworzą one
zbiór, co usuwa problem klas właściwych.

Dwie sieci Cauchy'ego $(x_d)_{d\in D}$ i $(y_e)_{e\in E}$
nazwiemy równoważnymi, gdy dla każdego otoczenia $U$ zera
istnieją $d_0,e_0$ takie, że
$x_d-y_e\in U$ dla $d\succeq d_0$ i $e\succeq e_0$.
Relacja ta jest zwrotna, ponieważ sieci są Cauchy'ego,
symetryczna dzięki symetrycznym otoczeniom i przechodnia:
dla danego $U$ wybieramy symetryczne $V$ z $V+V\subseteq U$
i łączymy dwa oszacowania przez wspólny późny wyraz
sieci pośredniej. Relację ograniczamy formalnie do zbioru sieci indeksowanych
przez $D_*$, a przez $[x_d]$ oznaczamy klasę
równoważności; $\widehat X$ jest zbiorem tych klas.
Dzięki powyższemu przeindeksowaniu zapis $[x_d]$
pozostaje jednoznaczny również dla sieci o innym zbiorze
indeksów. Algebraicznie sumujemy klasy, przechodząc do sieci
$(x_d+y_e)_{(d,e)\in D\times E}$, a następnie w razie potrzeby
do równoważnego przedstawiciela na $D_*$;
mnożenie przez ustalony skalar określamy przez
$\lambda[x_d]=[\lambda x_d]$. Liniowość działań w $X$
i ich ciągłość w zerze zapewniają, że są to działania
dobrze określone na klasach sieci Cauchy'ego.
Definiujemy $j(x)$ jako klasę stałej sieci o wartości $x$.

\emph{Krok 2: topologia.}
Wybierzmy otwarte zrównoważone otoczenie $U$ zera w $X$.
Niech $\widehat U$ oznacza zbiór klas $[x_d]$, dla których
istnieje otwarte otoczenie $V$ zera oraz indeks $d_0$ takie, że
\[
       x_d+V\subseteq U\qquad(d\succeq d_0).
\]
Definicja nie zależy od przedstawiciela: jeśli $(y_e)$
jest równoważna $(x_d)$, wybieramy otoczenie $W$ z
$W+W\subseteq V$. Dla dostatecznie późnych $d,e$
mamy $y_e-x_d\in W$, a więc
$y_e+W\subseteq x_d+(W+W)\subseteq x_d+V\subseteq U$.

Rodzina $\{\widehat U\}$ jest bazą otoczeń zera w
$\widehat X$. Istotnie, dla $V+V\subseteq U$ zachodzi
$\widehat V+\widehat V\subseteq\widehat U$,
a dla $U_1,U_2$ można wybrać otwarte zrównoważone
$V\subseteq U_1\cap U_2$. Wyjaśnijmy dokładnie, dlaczego przesunięcia tych zbiorów
określają topologię grupy addytywnej. Jeśli
$z=[x_d]\in\widehat U$ i $x_d+W\subseteq U$ dla późnych $d$,
wybierzmy otwarte zrównoważone $V\subseteq W$.
Dla $w=[y_e]\in\widehat V$ mamy
$y_e+W_e\subseteq V$ przy dostatecznie późnych $e$
i pewnym otwartym otoczeniu $W_0$ zera (niezależnym od $e$).
Wówczas dla późnych $d,e$ zachodzi
$x_d+y_e+W_0\subseteq x_d+V\subseteq U$,
a zatem $z+\widehat V\subseteq\widehat U$.
Tak więc każda klasa należąca do $\widehat U$ ma
otoczenie będące przesunięciem pewnego $\widehat V$
i mieszczące się w $\widehat U$; w połączeniu z
$\widehat V+\widehat V\subseteq\widehat U$ oraz
stabilnością na skończone przecięcia uzasadnia to
bazę otoczeń i ciągłość dodawania.
Zbiory $\widehat U$ są zrównoważone, gdyż dla
$0<|\lambda|\leq1$ mnożenie przez $\lambda$
przenosi $x_d+V\subseteq U$ w
$\lambda x_d+\lambda V\subseteq\lambda U\subseteq U$,
a przypadek $\lambda=0$ jest oczywisty.
Dlatego dodawanie i zmiana znaku są ciągłe. Ciągłość
mnożenia przez skalar uzasadnimy po wykazaniu gęstości $j(X)$.

Odwzorowanie $j$ jest liniowe, a z definicji
$j^{-1}(\widehat U)=U$: jeśli $x\in U$, to
$x+V\subseteq U$ dla pewnego otwartego otoczenia $V$ zera,
a implikacja odwrotna jest natychmiastowa.
Stąd $j$ jest zanurzeniem topologicznym. Co więcej,
$\widehat X$ jest Hausdorffa. Jeśli
$z=[x_d]\in\bigcap_U\widehat U$, to dla każdego otoczenia
$U$ można wybrać otwarte zrównoważone $V\subseteq U$
i $z\in\widehat V$, skąd ostatecznie $x_d\in V\subseteq U$.
Zatem $x_d\to0$ i $z=0$.

\emph{Krok 3: gęstość oraz granice sieci z $X$.}
Jeżeli $z=[x_d]$ i $U$ jest otoczeniem zera, wybieramy
otwarte zrównoważone $V,W$ z $V+W\subseteq U$.
Z warunku Cauchy'ego dla $(x_d)$ istnieje $d_0$ takie, że
$x_d-x_e\in V$ dla $d,e\succeq d_0$.
Dla każdego $e\succeq d_0$ sieć $(x_d-x_e)_d$ jest
ostatecznie zawarta w $V$, a wraz z $W$ mieści się w $U$.
Stąd $z-j(x_e)\in\widehat U$. Dowodzi to zarówno
gęstości $j(X)$, jak i zbieżności $j(x_d)\to z$.

Teraz domykamy dowód ciągłości mnożenia przez skalar.
Dla ustalonego $z_0\in\widehat X$ i zrównoważonego
$\widehat U$ wybierzmy zrównoważone $\widehat V$
takie, że $\widehat V+\widehat V\subseteq\widehat U$.
Z gęstości istnieje $x\in X$ takie, że
$z_0-j(x)\in\widehat V$. Dla $|\mu|\leq1$ mamy
$\mu(z_0-j(x))\in\widehat V$, natomiast
$\mu j(x)=j(\mu x)\to0$ przy $\mu\to0$.
Zatem $\mu z_0\to0$. W dowolnym punkcie
$(\lambda_0,z_0)$ korzystamy z tożsamości
$\lambda z-\lambda_0z_0=\lambda(z-z_0)+
(\lambda-\lambda_0)z_0$. Dla $|\lambda|\leq M$
i otwartego zrównoważonego $\widehat W$ z
$M\widehat W\subseteq\widehat V$ pierwszy składnik
należy do $\widehat V$, gdy $z-z_0\in\widehat W$,
a drugi przy $\lambda\to\lambda_0$ również należy
do $\widehat V$. To daje ciągłość łączną.

\emph{Krok 4: zupełność $\widehat X$.}
Niech $(z_\alpha)_{\alpha\in A}$ będzie siecią Cauchy'ego
w $\widehat X$. Niech $\mathcal U$ będzie rodziną otwartych,
zrównoważonych otoczeń zera w $X$, skierowaną przez odwrotną
inkluzję. Z gęstości dla każdej pary $(\alpha,U)$ wybieramy
$x_{\alpha,U}\in X$ tak, aby
$z_\alpha-j(x_{\alpha,U})\in\widehat U$.
Indeksami nowej sieci są pary $(\alpha,U)$, uporządkowane
wzrostem $\alpha$ i malejącym $U$.
Ta sieć w $X$ jest Cauchy'ego. Dla danego otwartego
zrównoważonego $V$ wybieramy otwarte zrównoważone $W$
z $W+W+W\subseteq V$. Dla dostatecznie późnych
$\alpha,\beta$ mamy $z_\alpha-z_\beta\in\widehat W$,
a dla $U_1,U_2\subseteq W$ zachodzi
\[
 j(x_{\alpha,U_1}-x_{\beta,U_2})
 \in\widehat W+\widehat W+\widehat W
 \subseteq\widehat V.
\]
Ponieważ $j^{-1}(\widehat V)=V$, różnica wyrazów należy do $V$.
Oznaczmy klasę tej sieci przez $z\in\widehat X$.
Na mocy kroku 3 zachodzi $j(x_{\alpha,U})\to z$.
Pokażemy, że $z_\alpha\to z$. Ustalmy otoczenie
$V$ zera w $X$ i wybierzmy $W$ z $W+W+W\subseteq V$.
Dla $\alpha,\beta$ dostatecznie późnych
$z_\alpha-z_\beta\in\widehat W$.
Wybieramy parę $(\beta,U)$ z późnym $\beta$ i $U\subseteq W$
na tyle późną, aby ponadto
$j(x_{\beta,U})-z\in\widehat W$. Ponieważ
$z_\beta-j(x_{\beta,U})\in\widehat U\subseteq\widehat W$,
dla wszystkich odpowiednio późnych $\alpha$ zachodzi
\[
 z_\alpha-z=(z_\alpha-z_\beta)
 +(z_\beta-j(x_{\beta,U}))+(j(x_{\beta,U})-z)
 \in\widehat V.
\]
Przestrzeń $\widehat X$ jest więc zupełna.

\emph{Krok 5: przedłużanie i jednoznaczność.}
Niech $T:X\to Y$ będzie jak w tezie. Jednostajna ciągłość
$T$ zapewnia, że dla każdej sieci Cauchy'ego $(x_d)$
sieć $(Tx_d)$ jest Cauchy'ego w $Y$, a więc zbiega.
Definiujemy
$\widehat T([x_d])=\lim_d Tx_d$.
Równoważne sieci mają tę samą granicę, gdyż ich różnice
na produkcie indeksów zbiegają do zera; liniowość jest
natychmiastowa. Aby wykazać ciągłość, niech $W$ będzie
otwartym otoczeniem zera w $Y$. Wybierzmy otwarte
symetryczne $V$ z $V+V\subseteq W$. Wówczas
$\overline V\subseteq V+V\subseteq W$: dla $y\in\overline V$
zbiór $y+V$ spotyka $V$, więc $y\in V-V$.
Z ciągłości $T$ istnieje otwarte zrównoważone otoczenie
$U$ zera w $X$ takie, że $T(U)\subseteq V$.
Dla $[x_d]\in\widehat U$ wyrazy $x_d$ o późnych indeksach
należą do $U$, a więc
$\widehat T([x_d])=\lim_d Tx_d\in\overline V\subseteq W$.
Otrzymujemy ciągłość $\widehat T$ w zerze, a z liniowości
wszędzie. Ponieważ $j(X)$ jest gęsta, przedłużenie jest
jedyne: dwa ciągłe odwzorowania do przestrzeni Hausdorffa
zgodne na zbiorze gęstym są równe. Jeśli
$(\widehat X',j')$ jest drugą parą o podanych własnościach,
przedłużamy $j'$ do $\widehat X\to\widehat X'$ oraz $j$
do $\widehat X'\to\widehat X$. Z jednoznaczności ich
złożenia są identycznościami, a więc otrzymujemy żądany
izomorfizm.
\end{proof}

\begin{center}
\begin{tikzpicture}[>=Latex, every node/.style={font=\small}]
  \node (X) at (0,1.3) {$X$};
  \node (H) at (5.0,1.3) {$\widehat X$};
  \node (Y) at (5.0,-0.45) {$Y$};
  \draw[->,thick] (X)--node[above] {$j$ (gęsty obraz)} (H);
  \draw[->,thick] (X)--node[below left] {$T$} (Y);
  \draw[->,thick] (H)--node[right] {$\widehat T$} (Y);
\end{tikzpicture}

{\small Rysunek 3.1. Każdy ciągły operator liniowy $T$ z $X$
do zupełnej TVS Hausdorffa $Y$ przedłuża się jednoznacznie
na uzupełnienie $\widehat X$.}
\end{center}

\begin{uwaga}[Dlaczego w konstrukcji pojawia się dodatkowe $V$?]
Nie wystarczyłoby zdefiniować $\widehat U$ przez wymaganie,
aby reprezentująca sieć \emph{ostatecznie należała} do $U$.
W $\R$ sieci $1-1/n$ i $1+1/n$ reprezentują tę samą klasę,
ale tylko pierwsza jest ostatecznie zawarta w $(-1,1)$.
Warunek $x_d+V\subseteq U$ wymaga odpowiedniego zapasu
do brzegu $U$ i nie zależy od reprezentanta.
\end{uwaga}

\subsection{Jak rozpoznawać uzupełnienie w praktyce?}

\begin{stwierdzenie}[Kryterium uzupełnienia]\label{prop:completion-criterion}
Niech $X$ będzie TVS Hausdorffa, $Y$ zupełną TVS Hausdorffa,
a $T:X\to Y$ liniowym zanurzeniem topologicznym o gęstym obrazie.
Wtedy $Y$, razem z $T$, jest uzupełnieniem $X$ w sensie
twierdzenia~\ref{thm:tvs-completion}.
\end{stwierdzenie}
\begin{proof}
Niech $(\widehat X,j)$ będzie uzupełnieniem z poprzedniego
twierdzenia. Z jego własności uniwersalnej otrzymujemy ciągłe
liniowe $S:\widehat X\to Y$ takie, że $S\circ j=T$.
Zbudujemy odwrotność. Odwzorowanie
$R_0:T(X)\to\widehat X$, $R_0(Tx)=j(x)$, jest liniowe,
ciągłe i jednostajnie ciągłe, ponieważ $T$ jest zanurzeniem
topologicznym. Dla $y\in Y$ wybierzmy sieć
$(y_d)$ punktów $T(X)$ zbieżną do $y$ (istnieje dzięki
gęstości). Jest to sieć Cauchy'ego, zatem $(R_0y_d)$
jest Cauchy'ego w zupełnej $\widehat X$. Definiujemy
$R(y)=\lim_d R_0(y_d)$.
Jeżeli dwie sieci z $T(X)$ zbiegają do tego samego $y$,
to ich różnice na produkcie indeksów zbiegają do zera.
Z jednostajnej ciągłości $R_0$ wynika równość granic;
zatem definicja nie zależy od wyboru sieci. Z tego samego
powodu $R$ jest liniowe.

Sprawdzimy ciągłość $R$ w zerze. Niech $U$ będzie otwartym
otoczeniem zera w $\widehat X$. Wybierzmy otwarte
symetryczne otoczenie $V$ z $V+V\subseteq U$;
wówczas $\overline V\subseteq U$, co wynika z argumentu
użytego w kroku 5 poprzedniego dowodu. Z ciągłości $R_0$
istnieje otwarte otoczenie $W$ zera w $Y$ takie, że
$R_0(W\cap T(X))\subseteq V$.
Dla $y\in W$ możemy dobrać sieć punktów z $W\cap T(X)$
zbieżną do $y$ (otwartość $W$ i gęstość $T(X)$).
Wtedy $R(y)\in\overline V\subseteq U$, zatem $R$
jest ciągłe. Na gęstym zbiorze $T(X)$ zachodzi
$S\circ R=\mathrm{id}_Y$, a na gęstym $j(X)$ zachodzi
$R\circ S=\mathrm{id}_{\widehat X}$. Z ciągłości obie
równości rozciągają się na całe dziedziny. Wobec tego
$S$ jest izomorfizmem liniowo-topologicznym.
\end{proof}

\begin{przyklad}[Trzy klasyczne uzupełnienia]
\begin{enumerate}[label=\textup{(\roman*)}]
\item Przestrzeń $c_{00}$ z normą $\ell^p$, $1\leq p<\infty$,
      ma uzupełnienie $\ell^p$: obcięcia ciągów do pierwszych
      $n$ współrzędnych są zbieżne w normie $\ell^p$.
\item Przestrzeń $C([0,1],\K)$ z normą
      $\|f\|_1=\int_{[0,1]}|f(t)|\,\mathrm d\lambda(t)$
      ma uzupełnienie $L^1([0,1],\lambda)$: funkcje ciągłe
      są gęste w $L^1$, a $L^1$ jest zupełna.
\item Przestrzeń $c_{00}$ z topologią odziedziczoną
      z produktu $\K^{\N}$ ma uzupełnienie $\K^{\N}$.
      W tym przypadku nie używa się jednej normy.
\end{enumerate}
\end{przyklad}

\subsection{Co wykorzystamy w dalszej części skryptu?}

Zupełność nie jest własnością wyłącznie przestrzeni Banacha.
Przykładowo $\K^{\N}$ z topologią produktową jest zupełna,
ale nie ma ograniczonego otoczenia zera i nie jest normowalna.
Następnie przejdziemy do własności rodzin operatorów, warunków
metryzowalności i lokalnej wypukłości. To przygotuje narzędzia
do badania przestrzeni $C^\infty(\Omega)$ oraz $\mathcal D(\Omega)$,
a docelowo ich dualnych przestrzeni dystrybucji.


% Wersja źródłowa: audyt II, jednostka 04
\setcounter{section}{3}
\section{Zbiory ograniczone, rodziny operatorów i metryzowalność TVS}\label{sec:rodziny-metryzowalnosc}

\subsection{Motywacja i przyjęte założenia}

W przestrzeni unormowanej liczba $\sup_{x\in B}\|x\|$ mierzy
rozmiar zbioru $B$. W ogólnej przestrzeni liniowo-topologicznej
nie mamy wyróżnionej normy. Możemy jednak określić, czy wszystkie
punkty zbioru $B$ mieszczą się w odpowiednim wielokrotnym
powiększeniu dowolnego otoczenia zera. W \hyperref[sec:tvs]{rozdziale 1} przyjęliśmy
następującą definicję: $B\subseteq X$ jest \emph{ograniczony}, gdy
dla każdego otoczenia $U$ zera istnieje $t>0$ takie, że
$B\subseteq tU$. Jest to ograniczoność w sensie von Neumanna.

W całym rozdziale $\K\in\{\R,\C\}$, a $X,Y$ oznaczają
przestrzenie liniowo-topologiczne nad $\K$; jeśli nie wskazano
inaczej, są one Hausdorffa. Otoczenia zera nie muszą być otwarte.
Przez $\mathcal L(X,Y)$ rozumiemy \emph{zbiór ciągłych odwzorowań
liniowych} $X\to Y$. Dla rodziny odwzorowań $\mathcal T$
zapis $\mathcal T(B)$ oznacza zbiór
$\{T(x):T\in\mathcal T,\ x\in B\}$.

\subsection{Sekwencyjne kryterium ograniczoności}

\begin{twierdzenie}[Kryterium ograniczoności przez skalarne ciągi]
\label{thm:bounded-sequential-test}
Niech $X$ będzie przestrzenią liniowo-topologiczną, a
$B\subseteq X$. Następujące warunki są równoważne:
\begin{enumerate}[label=\textup{(\roman*)}]
\item $B$ jest ograniczony w sensie von Neumanna;
\item dla każdego ciągu $(x_n)$ w $B$ i każdego ciągu
      $(\lambda_n)$ w $\K$ zbieżnego do zera zachodzi
      $\lambda_nx_n\to0$ w $X$.
\end{enumerate}
\end{twierdzenie}
\begin{proof}
Załóżmy (i). Niech $U$ będzie otoczeniem zera; wybierzmy
zrównoważone otoczenie $V\subseteq U$.
Istnieje $t>0$ takie, że $B\subseteq tV$.
Gdy $|\lambda_n|t\le1$, dostajemy
$\lambda_nx_n\in\lambda_ntV\subseteq V\subseteq U$.
Zatem $\lambda_nx_n\to0$.

Jeżeli (i) nie zachodzi, istnieje zrównoważone otoczenie
$U$ zera takie, że $B\not\subseteq tU$ dla żadnego $t>0$.
Dla każdego $n\ge1$ wybierzmy $x_n\in B\setminus nU$
i połóżmy $\lambda_n=1/n$. Wtedy $\lambda_n\to0$,
ale $\lambda_nx_n\notin U$ dla żadnego $n$, co przeczy (ii).
\end{proof}

\begin{stwierdzenie}[Podstawowe działania na zbiorach ograniczonych]
\label{prop:bounded-operations}
Niech $X$ będzie TVS, a $B,C\subseteq X$ zbiorami ograniczonymi.
Wówczas $B\cup C$, $B+C$, $\lambda B$ dla $\lambda\in\K$
oraz domknięcie $\overline B$ są ograniczone. Każdy podzbiór
zbioru ograniczonego jest ograniczony.
\end{stwierdzenie}
\begin{proof}
Wybierzmy zrównoważone otoczenie $V$ zera z
$V+V\subseteq U$, gdzie $U$ jest zadanym otoczeniem zera.
Istnieją liczby $s,t>0$ z $B\subseteq sV$ i $C\subseteq tV$.
Jeżeli $r=\max(s,t)$, to $B\cup C\subseteq rV$, ponieważ
$V$ jest zrównoważony. Ponadto
$B+C\subseteq sV+tV\subseteq (s+t)(V+V)\subseteq(s+t)U$;
wystarcza też użyć $r$ i sumy $V+V$. Ograniczoność
$\lambda B$ i podzbiorów wynika bezpośrednio z definicji.

Dla domknięcia wybierzmy otwarte zrównoważone otoczenie
$V$ z $V+V\subseteq U$. Zachodzi $\overline V\subseteq V+V$:
gdy $x\in\overline V$, otwarty zbiór $x+V$ przecina $V$,
zatem $x\in V-V=V+V$. Skoro $B\subseteq tV$ dla pewnego
$t>0$, mamy
$\overline B\subseteq t\overline V\subseteq tU$.
\end{proof}

\begin{uwaga}[Ważna pułapka: otoczka wypukła]
Dla $B\subseteq X$ przez $\operatorname{conv}(B)$ rozumiemy
zbiór wszystkich skończonych kombinacji wypukłych punktów $B$
(przyjmując $\operatorname{conv}(\varnothing)=\varnothing$).
W \emph{dowolnej} TVS otoczka wypukła zbioru ograniczonego
nie musi być ograniczona. Własność ta stanie się prawdziwa
po założeniu lokalnej wypukłości: dla dowolnego otoczenia
$U$ zera wybieramy wypukłe, zrównoważone otoczenie
$V\subseteq U$. Jeżeli $B\subseteq tV$, to
$\operatorname{conv}(B)\subseteq tV\subseteq tU$, więc
otoczka wypukła jest ograniczona.
\end{uwaga}

\begin{przyklad}[Nieograniczona otoczka wypukła kuli jednostkowej]
\label{ex:quasinorm-convex-hull}
Niech $0<p<1$ oraz $X=L^p([0,1])$ z metryką
$d(f,g)=\int_{[0,1]}|f(t)-g(t)|^p\,\mathrm d\lambda(t)$.
Zbiór $B=\{f:\int_{[0,1]}|f(t)|^p\,\mathrm d\lambda(t)\le1\}$ jest
ograniczony: dla dowolnego $\varepsilon>0$ jego odpowiednio
mała homotetia mieści się w
$\{f:\int_{[0,1]}|f(t)|^p\,\mathrm d\lambda(t)<\varepsilon\}$.
Dla $n\ge1$ podzielmy $[0,1]$ na przedziały $I_{n,j}$,
$j=1,\ldots,n$, o długości $1/n$ i połóżmy
$f_{n,j}=n^{1/p}\mathbf1_{I_{n,j}}$. Wtedy
$\int_{[0,1]}|f_{n,j}(t)|^p\,\mathrm d\lambda(t)=1$,
natomiast ich średnia arytmetyczna
$g_n=n^{-1}\sum_{j=1}^n f_{n,j}=n^{1/p-1}\mathbf1_{[0,1]}$
należy do $\operatorname{conv}(B)$ i spełnia
\[
 \int_{[0,1]}|g_n(t)|^p\,\mathrm d\lambda(t)=n^{1-p}\longrightarrow\infty.
\]
Otoczka wypukła $B$ nie jest więc ograniczona. Dowodzi to
również, że $X$ nie ma bazy otoczeń zera złożonej ze
zbiorów wypukłych: w przeciwnym razie otoczka wypukła
każdego zbioru ograniczonego byłaby ograniczona.
\end{przyklad}

\subsection{Lokalna ograniczoność i jej granice}

\begin{definicja}[Przestrzeń lokalnie ograniczona]
Przestrzeń liniowo-topologiczną $X$ nazywamy
\emph{lokalnie ograniczoną}, jeżeli istnieje ograniczone
otoczenie zera w $X$.
\end{definicja}

\begin{stwierdzenie}[Kryterium lokalnej ograniczoności]
\label{prop:locally-bounded}
Dla TVS $X$ następujące warunki są równoważne:
\begin{enumerate}[label=\textup{(\roman*)}]
\item $X$ ma ograniczone otoczenie zera;
\item istnieje otoczenie zera $U$ takie, że dla każdego
otoczenia zera $V$ istnieje $t>0$ z $U\subseteq tV$;
\item istnieje baza otoczeń zera postaci $\{tU:t>0\}$
dla pewnego zrównoważonego otoczenia $U$.
\end{enumerate}
\end{stwierdzenie}
\begin{proof}
Równoważność (i) i (ii) jest treścią definicji.
Załóżmy (i) i wybierzmy zrównoważone otoczenie
$U$ zawarte w ograniczonym otoczeniu. Jest ono również
ograniczone. Każdy zbiór $tU$, $t>0$, jest otoczeniem
zera, a ograniczoność $U$ daje: dla dowolnego otoczenia
$V$ istnieje $s>0$ takie, że $U\subseteq sV$, czyli
$s^{-1}U\subseteq V$. Dowodzi to (iii).
Implikacja (iii)$\Rightarrow$(i) wynika z warunku, że
dla danego otoczenia $V$ pewna homotetia $tU$ jest w $V$.
\end{proof}

\begin{przyklad}[Dwa odmienne przypadki]
Przestrzeń $\K^{\N}$ z topologią produktową nie jest
lokalnie ograniczona. Każde jej otoczenie zera zawiera
zbiór, w którym tylko skończenie wiele współrzędnych jest
ograniczonych; poza tym zbiorem współrzędnych można wybrać
współrzędną $j$ całkowicie swobodną. Takie otoczenie nie
mieści się w żadnym wielokrotnym powiększeniu zbioru
$\{x:|x_j|<1\}$.

Z kolei $L^p([0,1])$, $0<p<1$, jest lokalnie ograniczona:
metryczna kula zera określona przez
$\int_{[0,1]}|f(t)|^p\,\mathrm d\lambda(t)<1$ jest ograniczona.
Nie jest jednak normowalna, ponieważ nie jest lokalnie
wypukła (przykład~\ref{ex:quasinorm-convex-hull}).
Tak więc \emph{lokalna ograniczoność nie oznacza normowalności}.
\end{przyklad}

\subsection{Rodziny operatorów: ograniczoność punktowa i równociągłość}

\begin{definicja}[Ograniczoność punktowa i równociągłość]
\label{def:operator-family}
Niech $X,Y$ będą TVS i niech
$\mathcal T\subseteq\mathcal L(X,Y)$.
Rodzinę $\mathcal T$ nazywamy \emph{punktowo ograniczoną}, jeśli
dla każdego $x\in X$ zbiór $\{Tx:T\in\mathcal T\}$
jest ograniczony w $Y$. Rodzinę nazywamy
\emph{równociągłą} (w zerze), jeśli dla każdego otoczenia
$V$ zera w $Y$ istnieje otoczenie $U$ zera w $X$ takie,
że $T(U)\subseteq V$ dla wszystkich $T\in\mathcal T$.
\end{definicja}

\begin{twierdzenie}[Równociągłość i zbiory ograniczone]
\label{thm:equi-bounded}
Niech $X,Y$ będą TVS i $\mathcal T\subseteq\mathcal L(X,Y)$.
Jeśli $\mathcal T$ jest równociągła, to dla każdego
ograniczonego $B\subseteq X$ zbiór $\mathcal T(B)$
jest ograniczony w $Y$. W szczególności każda rodzina
równociągła jest punktowo ograniczona.
\end{twierdzenie}
\begin{proof}
Ustalmy ograniczony $B\subseteq X$ i otoczenie $V$ zera
w $Y$. Z równociągłości istnieje otoczenie $U$ zera
w $X$ z $T(U)\subseteq V$ dla wszystkich $T\in\mathcal T$.
Ponieważ $B\subseteq tU$ dla pewnego $t>0$, mamy
$T(B)\subseteq tT(U)\subseteq tV$ dla każdego $T$.
Zatem $\mathcal T(B)\subseteq tV$.
\end{proof}

\begin{przyklad}[Punktowa ograniczoność bez równociągłości]
\label{ex:c00-uniform-boundedness-failure}
Niech $X=c_{00}$ będzie przestrzenią ciągów o skończonym
nośniku z normą $\|x\|_\infty=\sup_{k\ge1}|x_k|$.
Dla $n\ge1$ określmy ciągły funkcjonał
$T_n:X\to\K$, $T_nx=nx_n$. Każdy $x\in c_{00}$
ma skończony nośnik, zatem $T_nx=0$ dla wszystkich
dostatecznie dużych $n$. Rodzina $\{T_n:n\ge1\}$
jest więc punktowo ograniczona (a nawet $T_nx\to0$
dla każdego $x$). Nie jest równociągła: zbiór
$B=\{e_n:n\ge1\}$ (gdzie $e_n$ jest $n$-tym wektorem
standardowej bazy współrzędnościowej) jest ograniczony w $X$, lecz
$T_n(e_n)=n$ i $\mathcal T(B)$ jest nieograniczony w $\K$.

Ten przykład pokazuje, że punktowa ograniczoność rodzin
operatorów \emph{nie wystarcza} do równociągłości w dowolnej
przestrzeni normowanej. Zupełność dziedziny będzie kluczowa
w późniejszym twierdzeniu Banacha--Steinhausa.
\end{przyklad}

\begin{stwierdzenie}[Rodziny na przestrzeniach unormowanych]
\label{prop:equi-normed}
Niech $X,Y$ będą przestrzeniami unormowanymi, a
$\mathcal T\subseteq\mathcal L(X,Y)$. Wtedy $\mathcal T$
jest równociągła dokładnie wtedy, gdy
$\sup_{T\in\mathcal T}\|T\|<\infty$,
gdzie $\|T\|=\sup_{\|x\|\le1}\|Tx\|$.
\end{stwierdzenie}
\begin{proof}
Jeśli rodzina jest równociągła, to dla otoczenia
$V=\{y:\|y\|<1\}$ istnieje $r>0$ takie, że
$\|x\|<r$ pociąga $\|Tx\|<1$ dla wszystkich $T$.
Dla $\|x\|\le1$ zastosujmy to do $(r/2)x$:
$\|Tx\|<2/r$. Stąd supremum norm operatorowych jest
skończone.
Odwrotnie, jeśli $M=\sup_T\|T\|<\infty$, to dla
$\varepsilon>0$ wystarczy wybrać
$U=\{x:\|x\|<\varepsilon/(M+1)\}$.
\end{proof}

\subsection{Topologie zbieżności operatorów}

\begin{definicja}[Zbieżność punktowa i jednostajna na ograniczonych]
\label{def:operator-topologies}
Niech $X,Y$ będą TVS Hausdorffa. Na przestrzeni
$\mathcal L(X,Y)$ rozważamy dwie topologie:
\begin{enumerate}[label=\textup{(\roman*)}]
\item topologię $\tau_s$ \emph{zbieżności punktowej},
wyznaczoną odwzorowaniami $T\mapsto Tx$, $x\in X$;
\item topologię $\tau_b$ \emph{zbieżności jednostajnej na
zbiorach ograniczonych}, której bazę otoczeń zera
stanowią skończone przecięcia zbiorów
\[
 \mathcal N(B,V)=\{T\in\mathcal L(X,Y):T(B)\subseteq V\},
\]
gdzie $B\subseteq X$ jest ograniczony, a $V$ jest
otoczeniem zera w $Y$.
\end{enumerate}
\end{definicja}

\begin{twierdzenie}[Struktura liniowo-topologiczna przestrzeni operatorów]
\label{thm:op-topologies}
Niech $X,Y$ będą TVS Hausdorffa. Topologie $\tau_s$ i
$\tau_b$ czynią $\mathcal L(X,Y)$ przestrzenią
liniowo-topologiczną Hausdorffa, a $\tau_b$ jest silniejsza
od $\tau_s$. Dla sieci $(T_d)$ i $T\in\mathcal L(X,Y)$
zachodzi
\begin{align*}
 T_d\xrightarrow{\tau_s}T
 &\iff T_dx\to Tx\ \text{dla każdego }x\in X,\\
 T_d\xrightarrow{\tau_b}T
 &\iff \text{dla każdego ograniczonego }B\subseteq X
 \text{ i }V\in\mathcal N_Y(0)\\
 &\hspace{18mm}\text{ostatecznie }(T_d-T)(B)\subseteq V.
\end{align*}
\end{twierdzenie}
\begin{proof}
Topologia $\tau_s$ jest topologią podprzestrzeni
otrzymaną przez zanurzenie $T\mapsto(Tx)_{x\in X}$
w produkt $Y^X$, więc jest Hausdorffa i zgodna z
działaniami liniowymi.

Dla $\tau_b$ zbiory $\mathcal N(B,V)$ zawierają zero.
Wystarczy rozważać otwarte, zrównoważone otoczenia $V$,
ponieważ tworzą one bazę otoczeń zera w $Y$.
Dla takiego $V$ można dobrać otwarte, zrównoważone
$W$ z $W+W\subseteq V$;
wtedy
$\mathcal N(B,W)+\mathcal N(B,W)\subseteq\mathcal N(B,V)$.
Przecięcia wyraża się przez skończone układy warunków.
Zbiory te są zrównoważone, jeżeli $V$ jest zrównoważony.
Dla każdego $S\in\mathcal N(B,W)$ zachodzi
$S+\mathcal N(B,W)\subseteq\mathcal N(B,V)$.
Warunki te zapewniają istnienie topologii grupowej
z podaną bazą \emph{otoczeń} zera. Nie twierdzimy, że każdy
zbiór $\mathcal N(B,V)$ jest otwarty; może nim nie być.

Sprawdźmy łączną ciągłość mnożenia przez skalar.
Ustalmy $\lambda_0\in\K$, $T_0\in\mathcal L(X,Y)$,
ograniczony $B\subseteq X$ i otoczenie $V$ zera w $Y$.
Wybierzmy zrównoważone otoczenie $W$ z $W+W\subseteq V$.
Ponieważ $T_0(B)$ jest ograniczony, dla $\lambda$
dostatecznie bliskich $\lambda_0$ mamy
$(\lambda-\lambda_0)T_0(B)\subseteq W$.
W otoczeniu $\lambda_0$ liczby $|\lambda|$ są ograniczone
przez pewne $M>0$. Dobierzmy zrównoważone otoczenie
$U$ takie, że $\lambda U\subseteq W$ dla $|\lambda|\le M$.
Jeżeli $T-T_0\in\mathcal N(B,U)$, to
\[
 (\lambda T-\lambda_0T_0)(B)
 \subseteq\lambda(T-T_0)(B)+(\lambda-\lambda_0)T_0(B)
 \subseteq W+W\subseteq V.
\]
To dowodzi ciągłości łącznej. Ponieważ każdy singleton
$\{x\}\subseteq X$ jest ograniczony, otoczenia
$\mathcal N(\{x\},V)$ należą do bazy $\tau_b$,
a zatem $\tau_b\supseteq\tau_s$. Warunek Hausdorffa
wynika z faktu, że jeżeli $T\ne0$, to dla pewnego $x$
mamy $Tx\ne0$ i można oddzielić $Tx$ od zera w $Y$.
Opisy zbieżności wynikają wprost z definicji.
\end{proof}

\begin{stwierdzenie}[Ograniczoność rodzin w topologiach operatorów]
\label{prop:op-family-bounded}
Niech $X,Y$ będą TVS Hausdorffa i
$\mathcal T\subseteq\mathcal L(X,Y)$. Wówczas:
\begin{enumerate}[label=\textup{(\roman*)}]
\item $\mathcal T$ jest ograniczona w $(\mathcal L(X,Y),\tau_s)$
dokładnie wtedy, gdy jest punktowo ograniczona;
\item $\mathcal T$ jest ograniczona w $(\mathcal L(X,Y),\tau_b)$
dokładnie wtedy, gdy dla każdego ograniczonego $B\subseteq X$
zbiór $\mathcal T(B)$ jest ograniczony w $Y$.
\end{enumerate}
\end{stwierdzenie}
\begin{proof}
W (i) sprawdzamy ograniczoność we wszystkich ewaluacjach
$T\mapsto Tx$ i korzystamy z kryterium ograniczoności
w produkcie. W (ii) z definicji ograniczoności w
$\tau_b$ dla danego $B,V$ istnieje $t>0$ takie,
że $\mathcal T\subseteq t\mathcal N(B,V)$; jest to
równoważne zawieraniu $\mathcal T(B)\subseteq tV$.
Skończone przecięcia warunków nie zmieniają tego kryterium,
ponieważ bierzemy największą z odpowiednich stałych
skalowania po uprzednim zastąpieniu otoczeń zrównoważonymi.
\end{proof}

\begin{przyklad}[Dwie różne zbieżności operatorów]
Niech $X=c_{00}$ z normą $\ell^\infty$, a $Y=\K$.
Funkcjonały $T_nx=nx_n$ spełniają
$T_n\xrightarrow{\tau_s}0$, ale nie
$T_n\xrightarrow{\tau_b}0$, ponieważ
$\{e_n:n\ge1\}$ jest ograniczony i
$\sup_{x\in\{e_n:n\ge1\}}|T_nx|=n$.
\end{przyklad}

\begin{twierdzenie}[Przypadek przestrzeni unormowanych]
\label{thm:op-b-norm}
Jeżeli $X,Y$ są przestrzeniami unormowanymi, to topologia
$\tau_b$ na $\mathcal L(X,Y)$ jest topologią normy
operatorowej. Jeżeli $X$ ma dodatkowo wymiar skończony,
również $\tau_s=\tau_b$.
\end{twierdzenie}
\begin{proof}
Kula jednostkowa $B_X=\{x:\|x\|\le1\}$ jest ograniczona.
Warunek $(T_d-T)(B_X)\subseteq\{y:\|y\|<\varepsilon\}$
pociąga $\|T_d-T\|\le\varepsilon$. Wybierając $\varepsilon/2$ zamiast $\varepsilon$ widzimy, że $\tau_b$ jest co najmniej tak silna jak topologia normy. Odwrotnie, jeśli
$B\subseteq X$ jest ograniczony, to
$B\subseteq RB_X$ dla pewnego $R>0$, więc
$\|S\|<\varepsilon/R$ implikuje $S(B)\subseteq
\{y:\|y\|<\varepsilon\}$. W ten sposób otrzymujemy
tę samą topologię.

Jeśli $\dim X=0$, to $\mathcal L(X,Y)=\{0\}$
i obie topologie są jedyną topologią tej przestrzeni.
Załóżmy więc $\dim X=m\geq1$ i wybierzmy bazę
$e_1,\ldots,e_m$. Dla każdego $x=\sum_i a_ie_i$
zachodzą oszacowania
$|a_i|\le C\|x\|$ dla pewnej wspólnej stałej $C$:
aby to uzasadnić, rozważmy zbiór
$S=\{(a_1,\ldots,a_m):\max_i|a_i|=1\}$.
Jest on zwarty w $\K^m$, a ciągła funkcja
$(a_i)\mapsto\|\sum_i a_ie_i\|$ ma na nim dodatnie
minimum $c$. Skalowanie daje
$\|\sum_i a_ie_i\|\ge c\max_i|a_i|$, co
wystarcza przy $C=1/c$.
Jeżeli $\|S e_i\|<\delta$ dla każdego $i$, to
$\|Sx\|\le mC\delta\|x\|$, a zatem ograniczenie
operatora kontroluje skończenie wiele ewaluacji.
Wynika stąd $\tau_s=\tau_b$.
\end{proof}

\subsection{Twierdzenie o metryzowalności przestrzeni liniowo-topologicznych}

\begin{definicja}[Pierwszy aksjomat przeliczalności]
Przestrzeń topologiczną $X$ nazywamy
\emph{pierwszoprzeliczalną}, jeżeli każdy punkt ma
przeliczalną bazę otoczeń. W TVS wystarcza, aby
przeliczalna baza otoczeń istniała w zerze.
\end{definicja}

\begin{twierdzenie}[Kryterium metryzowalności TVS]
\label{thm:tvs-metrizable}
Niech $X$ będzie przestrzenią liniowo-topologiczną
Hausdorffa nad $\K$. Następujące warunki są równoważne:
\begin{enumerate}[label=\textup{(\roman*)}]
\item $X$ jest metryzowalna;
\item $X$ ma przeliczalną bazę otoczeń zera;
\item topologia $X$ jest wyznaczona przez pewną
      metrykę niezmienniczą względem przesunięć,
      tj. $d(x+a,y+a)=d(x,y)$ dla wszystkich $x,y,a\in X$.
\end{enumerate}
Co więcej, metryka z (iii) wyznacza kanoniczną strukturę
jednostajną $X$, a więc zupełność tej metryki jest
równoważna zupełności TVS.
\end{twierdzenie}
\begin{proof}
Implikacje $(\mathrm{iii})\Rightarrow(\mathrm{i})\Rightarrow(\mathrm{ii})$ są oczywiste:
w przestrzeni metrycznej kule o promieniach $1/n$
tworzą przeliczalną bazę otoczeń zera.

Załóżmy (ii). Wybierzmy malejącą bazę otwartych otoczeń
zera $(W_n)_{n\ge1}$ i skonstruujmy rekurencyjnie
otwarte, symetryczne i zrównoważone otoczenia $U_n$,
$n\ge0$, takie, że
$U_0=X$, $U_n\subseteq W_n$ dla $n\ge1$ oraz
$U_{n+1}+U_{n+1}+U_{n+1}\subseteq U_n$.
Jest to możliwe dzięki ciągłości dodawania i istnieniu
zrównoważonych otoczeń zera. Mamy $\bigcap_nU_n=\{0\}$:
przeciwne założenie przeczyłoby własności Hausdorffa.

Dla $x,y\in X$ połóżmy
\[
 d(x,y)=\inf\Big\{\sum_{j=1}^{m}2^{-n_j}:\
 m\ge1,\ x=x_0,\ldots,x_m=y,\
 x_{j-1}-x_j\in U_{n_j},\ n_j\ge0\Big\}.
\]
Infimum dotyczy wszystkich skończonych łańcuchów.
Ponieważ $U_0=X$, zbiór, z którego bierzemy infimum,
jest niepusty. Symetria zbiorów $U_n$ daje
$d(x,y)=d(y,x)$, a konkatenacja łańcuchów nierówność
trójkąta. Wzór zależy tylko od różnic $x-y$, więc
$d(x+a,y+a)=d(x,y)$.

Sprawdzimy, że $d$ jest dodatnio określona i generuje
topologię $X$. Potrzebujemy następującego rachunku:
\emph{jeżeli łańcuch ma łączną wagę mniejszą niż $2^{-n}$,
$n\ge1$, to różnica jego końców należy do $U_{n-1}$.}
Dowód prowadzimy indukcyjnie względem liczby jego
odcinków. Dla jednego odcinka warunek
$2^{-n_1}<2^{-n}$ daje $n_1>n$, czyli różnica należy do
$U_{n+1}\subseteq U_{n-1}$. Dla łańcucha dłuższego
wybierzmy odcinek przecinający połowę jego łącznej
wagi. Suma wag odcinków poprzedzających ten odcinek
oraz suma wag odcinków następujących po nim są
obie mniejsze niż $2^{-(n+1)}$. Z założenia
indukcyjnego ich całkowite różnice należą do $U_n$.
Waga wyróżnionego odcinka jest potęgą $2^{-n_j}$
i jest mniejsza niż $2^{-n}$, skąd $n_j\geq n+1$;
jego różnica należy więc do $U_{n_j}\subseteq U_n$.
Różnica końców całego łańcucha leży zatem w
$U_n+U_n+U_n\subseteq U_{n-1}$.

Dla każdego $n\ge1$ wynika stąd
\[
 U_{n+1}\subseteq\{x:d(x,0)<2^{-n}\}\subseteq U_{n-1}.
\]
Rzeczywiście pierwsza inkluzja odpowiada łańcuchowi
jednoodcinkowemu, a druga -- powyższemu rachunkowi.
Jeżeli $d(x,0)=0$, to $x$ należy do każdego $U_n$,
więc $x=0$. Metryka jest zatem dodatnio określona.
Dwie inkluzje dowodzą zarazem, że jej kule stanowią
tę samą bazę otoczeń zera co $U_n$; niezmienniczość
względem przesunięć daje zgodność topologii w każdym punkcie.

Otoczenia przekątnej dla metryki $d$ są postaci
$\{(x,y):d(x,y)<r\}$. Powyższe inkluzje pokazują,
że tworzą one dokładnie tę samą strukturę jednostajną,
co otoczenia $E_U=\{(x,y):x-y\in U\}$.
Sieci Cauchy'ego obu struktur są więc identyczne.
To dowodzi ostatniej części tezy.
\end{proof}

\begin{przyklad}[Metryzowalna, zupełna, ale nienormowalna]
\label{ex:metrizable-not-normable}
Na $X=\K^{\N}$ z topologią produktową określmy
\[
 d(x,y)=\sum_{n=1}^{\infty}2^{-n}\min\{1,|x_n-y_n|\}.
\]
Szereg jest zbieżny, $d$ jest metryką niezmienniczą
względem przesunięć i opisuje zbieżność współrzędną.
Przestrzeń $X$ jest zupełna jako produkt zupełnych
przestrzeni $\K$ (co wynika z zupełności produktów),
lecz nie jest normowalna, ponieważ nie ma ograniczonego
otoczenia zera. To pokazuje, że metryzowalność i
zupełność nie zastępują normowalności.
\end{przyklad}

\subsection{Jakie założenia będą potrzebne w następnych rozdziałach?}

Różnicę między ograniczonością punktową i równociągłością
usuwa w ważnych sytuacjach zasada jednostajnej ograniczoności;
jednym z jej naturalnych kontekstów są przestrzenie beczkowe.
Twierdzenie to wymaga dodatkowej struktury i zostanie
udowodnione później, a nie założone milcząco.

W kolejnym rozdziale wprowadzimy lokalną wypukłość,
seminormy i funkcjonał Minkowskiego. Dowiemy się między
innymi, dlaczego wypukłość otoczeń pozwala zamieniać
ograniczoność na oszacowania przez seminormy oraz
kiedy lokalna ograniczoność prowadzi już do normowalności.


% Wersja źródłowa: audyt II, jednostka 05
\setcounter{section}{4}
\section{Przestrzenie lokalnie wypukłe i seminormy}\label{sec:lokalnie-wypukle}

\subsection{Motywacja: od topologii do nierówności}

W poprzednich rozdziałach badaliśmy przestrzenie liniowo-topologiczne
bez wymagania, aby ich otoczenia były wypukłe. Okazało się, że nawet
zupełna i metryzowalna TVS może zachowywać się odmiennie od przestrzeni
unormowanej: w $L^p([0,1])$, $0<p<1$, otoczka wypukła ograniczonej
kuli nie jest ograniczona, a każdy ciągły funkcjonał liniowy jest zerowy.
Wprowadzenie \emph{lokalnej wypukłości} pozwoli nam powiązać topologię
z oszacowaniami przez seminormy, a następnie z teorią dualności.

W całym rozdziale $\K\in\{\R,\C\}$. Jeżeli nie zaznaczono inaczej,
$X$ i $Y$ oznaczają przestrzenie liniowo-topologiczne nad $\K$;
nie zakładamy milcząco ani normowalności, ani metryzowalności.
Przez $\mathcal N_X(0)$ oznaczamy rodzinę otoczeń zera w $X$.
Zbiór $A\subseteq X$ nazywamy \emph{pochłaniającym}, jeśli
$\forall x\in X\ \exists t>0:\ x\in tA$.
W TVS każde otoczenie zera jest pochłaniające, ponieważ
$\lambda x\to0$ dla $\lambda\to0$.

\begin{definicja}[Wypukłość i zrównoważenie]
Niech $X$ będzie przestrzenią liniową nad $\K$.
Zbiór $A\subseteq X$ nazywamy \emph{wypukłym}, gdy
$(1-t)x+ty\in A$ dla wszystkich $x,y\in A$, $t\in[0,1]$.
Nazywamy go \emph{zrównoważonym}, gdy $\lambda A\subseteq A$
dla każdej $\lambda\in\K$ z $|\lambda|\le1$.
Zbiór jednocześnie wypukły i zrównoważony nazywamy
\emph{absolutnie wypukłym}.
\end{definicja}

\begin{stwierdzenie}[Wypukłość absolutna jako nierówność na sumę]
\label{prop:absolutna}
Niech $X$ będzie przestrzenią liniową nad $\K$ i
$A\subseteq X$ zbiorem niepustym. Jest on absolutnie wypukły
wtedy i tylko wtedy, gdy dla wszystkich $x,y\in A$ oraz
$\alpha,\beta\in\K$ spełniających $|\alpha|+|\beta|\le1$
zachodzi $\alpha x+\beta y\in A$.
\end{stwierdzenie}
\begin{proof}
Jeśli $A$ jest absolutnie wypukły, to
$\alpha x+\beta y=(|\alpha|+|\beta|)z$ przy
$z=\frac{|\alpha|}{|\alpha|+|\beta|}\frac{\alpha}{|\alpha|}x+
\frac{|\beta|}{|\alpha|+|\beta|}\frac{\beta}{|\beta|}y$
(wyrazy z zerowym współczynnikiem pomijamy).
Wypukłość i zrównoważenie dają $z\in A$ oraz
$\alpha x+\beta y\in A$. Gdy oba współczynniki są zerowe,
$0\in A$ wynika ze zrównoważenia. Odwrotnie, kładąc
$\beta=0$ otrzymujemy zrównoważenie, a kładąc
$\alpha=t$, $\beta=1-t$ dla $t\in[0,1]$, wypukłość.
\end{proof}

\begin{definicja}[Przestrzeń lokalnie wypukła]
\label{def:lcs}
Przestrzeń liniowo-topologiczną $X$ nazywamy
\emph{lokalnie wypukłą} (LCS), jeśli zero ma bazę otoczeń
złożoną ze zbiorów wypukłych. Przez \emph{lokalnie wypukłą
przestrzeń Hausdorffa} rozumiemy LCS, której topologia jest
Hausdorffa; rozróżnienie to jest istotne w twierdzeniach o normowalności.
\end{definicja}

\begin{stwierdzenie}[Absolutnie wypukła baza otoczeń]
\label{prop:absconv-base}
Niech $X$ będzie TVS. Przestrzeń $X$ jest lokalnie wypukła
wtedy i tylko wtedy, gdy zero ma bazę otoczeń absolutnie wypukłych.
Bazę można wybrać ze zbiorów otwartych.
\end{stwierdzenie}
\begin{proof}
Załóżmy lokalną wypukłość i ustalmy otoczenie $U$ zera.
Wybierzmy wypukłe otoczenie $V\subseteq U$, a następnie
otwarte zrównoważone otoczenie $O\subseteq V$.
Zbiór $\operatorname{conv}(O)$ jest zawarty w $V$,
bo $V$ jest wypukły; jest też zrównoważony,
ponieważ dla $|\lambda|\le1$ zachodzi
$\lambda\operatorname{conv}(O)=\operatorname{conv}(\lambda O)
\subseteq\operatorname{conv}(O)$.
Zbiór $\operatorname{conv}(O)$ jest absolutnie wypukły.
Ponadto jest otwarty: jeżeli $x\in\operatorname{conv}(O)$,
to $x+O\subseteq\operatorname{conv}(O)+O\subseteq
2\operatorname{conv}(O)$ nie wystarcza do żądanej tezy;
można natomiast zapisać $x=\sum_{j=1}^m t_jx_j$
z $t_j\ge0$, $\sum t_j=1$, i wybrać $j$ z $t_j>0$.
Ponieważ $O$ jest otwarty, dla dostatecznie małego
otoczenia $N$ zera mamy $x_j+t_j^{-1}N\subseteq O$,
stąd $x+N\subseteq\operatorname{conv}(O)$.
Druga implikacja wynika bezpośrednio z definicji.
\end{proof}

\begin{uwaga}[Dlaczego lokalna wypukłość jest istotna?]
W przestrzeni $L^p([0,1])$ dla $0<p<1$ nie ma bazy
wypukłych otoczeń zera; dowód podano w
przykładzie z \hyperref[sec:rodziny-metryzowalnosc]{rozdziału 4}.
Tak więc zupełność i metryzowalność nie wymuszają lokalnej wypukłości.
\end{uwaga}

\subsection{Seminormy i topologie zadane przez rodziny seminorm}

\begin{definicja}[Seminorma]
Niech $X$ będzie przestrzenią liniową nad $\K$.
\emph{Seminormą} na $X$ nazywamy funkcję
$p:X\to[0,\infty)$ spełniającą
$p(\lambda x)=|\lambda|p(x)$ dla $\lambda\in\K$, $x\in X$,
oraz $p(x+y)\le p(x)+p(y)$ dla $x,y\in X$.
Jeżeli dodatkowo $p(x)=0$ implikuje $x=0$, to $p$ jest normą.
\end{definicja}

\begin{przyklad}[Pierwsze seminormy]
Na $\K^{\N}$ funkcje $p_n(x)=|x_n|$ są seminormami, ale
nie normami. Na przestrzeni $C^\infty([0,1])$ seminormami są
$p_m(f)=\max_{0\le j\le m}\sup_{t\in[0,1]}|f^{(j)}(t)|$,
$m\in\N\cup\{0\}$. Na $C(\Omega)$ dla otwartego
$\Omega\subseteq\R^d$ seminormą jest
$p_K(f)=\sup_{x\in K}|f(x)|$ dla dowolnego zwartego $K\subseteq\Omega$.
\end{przyklad}

\begin{twierdzenie}[Topologia generowana przez seminormy]
\label{thm:seminorm-topology}
Niech $X$ będzie przestrzenią liniową nad $\K$, a
$\mathcal P$ niepustą rodziną seminorm na $X$. Istnieje
najmniejsza topologia $\tau_{\mathcal P}$ \emph{zgodna ze
strukturą przestrzeni liniowo-topologicznej}, w której
wszystkie seminormy $p\in\mathcal P$ są ciągłe.
Topologia ta czyni $X$ przestrzenią lokalnie wypukłą,
a bazą otoczeń zera są zbiory
\[
 U(p_1,\ldots,p_m;\varepsilon)
 =\{x\in X:\max_{1\le j\le m}p_j(x)<\varepsilon\},
 \quad m\ge1,\quad p_j\in\mathcal P,\quad\varepsilon>0.
\]
Jest ona Hausdorffa wtedy i tylko wtedy, gdy
$\bigcap_{p\in\mathcal P}\ker p=\{0\}$.
Ponadto sieć $(x_d)$ zbiega do $x$ w $\tau_{\mathcal P}$
dokładnie wtedy, gdy $p(x_d-x)\to0$ dla każdego $p\in\mathcal P$.
\end{twierdzenie}
\begin{proof}
Zbiory z tezy są absolutnie wypukłe, pochłaniające i
zamknięte na skończone przecięcia (po ewentualnym
zmniejszeniu $\varepsilon$). Dla zbioru z promieniem
$\varepsilon$ suma dwóch analogicznych zbiorów z promieniem
$\varepsilon/2$ mieści się w wyjściowym zbiorze.
Aby uzasadnić ciągłość mnożenia przez skalar w punkcie
$(\lambda_0,x_0)$, dla każdego z ustalonych $p_j$
używamy oszacowania
$p_j(\lambda x-\lambda_0x_0)\le
|\lambda|p_j(x-x_0)+|\lambda-\lambda_0|p_j(x_0)$.
Wystarczy ograniczyć $|\lambda|$ w małym otoczeniu
$\lambda_0$, a następnie wybrać odpowiednio małe
$|\lambda-\lambda_0|$ i $p_j(x-x_0)$ dla wszystkich
skończenie wielu indeksów $j$.
W ten sposób bazy przesuniętych zbiorów definiują
topologię TVS.
Każde $p\in\mathcal P$ jest ciągłe w zerze, a z nierówności
$|p(x)-p(y)|\le p(x-y)$ wynika jej ciągłość wszędzie.
Każda inna topologia TVS, w której seminormy są ciągłe,
zawiera wszystkie powyższe otoczenia, zatem jest silniejsza.

Jeżeli $x\ne0$ i $p(x)>0$ dla pewnego $p$, otoczenie
$\{y:p(y)<p(x)/2\}$ oddziela $0$ od $x$.
Jeśli natomiast $p(x)=0$ dla wszystkich $p$, to $x$
należy do każdego otoczenia zera. W konsekwencji
$\tau_{\mathcal P}$ jest Hausdorffa dokładnie przy
podanym warunku. Opis zbieżności wynika z bazy otoczeń.
\end{proof}

\subsection{Funkcjonał Minkowskiego}

\begin{definicja}[Funkcjonał Minkowskiego]
\label{def:gauge}
Niech $X$ będzie przestrzenią liniową, a $A\subseteq X$
niepustym zbiorem pochłaniającym. Funkcjonałem Minkowskiego
(\emph{funkcjonałem cechującym}) zbioru $A$ nazywamy
\[
 p_A(x)=\inf\{t>0:x\in tA\},\qquad x\in X.
\]
Pochłanianie zapewnia, że infimum jest skończone.
\end{definicja}

\begin{twierdzenie}[Od zbioru absolutnie wypukłego do seminormy]
\label{thm:minkowski}
Niech $X$ będzie TVS nad $\K$, a $A\subseteq X$
pochłaniającym zbiorem absolutnie wypukłym.
Wówczas $p_A$ jest seminormą i
$\{x:p_A(x)<1\}\subseteq A\subseteq\{x:p_A(x)\le1\}$.
Jeśli dodatkowo $A$ jest \emph{otwarty}, to
$A=\{x:p_A(x)<1\}$ i $p_A$ jest ciągła.
\end{twierdzenie}
\begin{proof}
Pochłanianie oraz $0\in A$ dają $p_A(x)<\infty$ i
$p_A(0)=0$. Dla $\lambda\ne0$ mamy
$x\in tA$ dokładnie wtedy, gdy $\lambda x\in|\lambda|tA$:
przy zespolonym $\lambda$ korzystamy ze zrównoważenia
$A$ i z równości $\frac{\lambda}{|\lambda|}A=A$.
Otrzymujemy $p_A(\lambda x)=|\lambda|p_A(x)$;
przypadek $\lambda=0$ jest osobny i natychmiastowy.

Niech $s>p_A(x)$ oraz $t>p_A(y)$. Z definicji infimum
i z własności $rA\subseteq sA$ dla $0<r\le s$
wynika $x\in sA$ i $y\in tA$. Wypukłość daje
$sA+tA\subseteq(s+t)A$, zatem
$p_A(x+y)\le s+t$. Przechodząc z $s$ i $t$ do
odpowiednich infimów otrzymujemy podaddytywność.
Jeżeli $p_A(x)<1$, to $x\in tA$ dla pewnego $t<1$,
więc $x\in A$ dzięki zrównoważeniu.
Jeżeli $x\in A$, to $p_A(x)\le1$.

Niech teraz $A$ będzie otwarty i niech $x\in A$.
Ciągłość $\lambda\mapsto\lambda x$ w $\lambda=1$
daje $\lambda>1$ dostatecznie bliskie $1$ z
$\lambda x\in A$; stąd $p_A(x)\le1/\lambda<1$.
Zatem $A=\{p_A<1\}$.
Ponieważ $\{p_A<\varepsilon\}=\varepsilon A$
dla każdego $\varepsilon>0$, seminorma jest ciągła
w zerze, a nierówność
$|p_A(x)-p_A(y)|\le p_A(x-y)$ zapewnia ciągłość wszędzie.
\end{proof}

\begin{center}
\begin{tikzpicture}[scale=.85,>=Latex]
 \draw[fill=blue!5,draw=blue!50!black,thick] (0,0) ellipse (2.3 and 1.35);
 \draw[->,gray] (-2.8,0)--(4.35,0) node[right] {$x_1$};
 \draw[->,gray] (0,-1.65)--(0,1.9) node[above] {$x_2$};
 \fill (0,0) circle (1.8pt) node[below left] {$0$};
 \draw[->,thick] (0,0)--(3.45,1.12) node[above right] {$x$};
 \fill[blue!70!black] (2.012,.653) circle (2.3pt);
 \node[blue!60!black,above] at (1.50,.84) {$x/p_A(x)$};
 \node[blue!60!black] at (-1.45,.7) {$A$};
\end{tikzpicture}

\small Rysunek: dla absolutnie wypukłego otwartego otoczenia $A$
wartość $p_A(x)$ jest współczynnikiem skalowania potrzebnym,
aby umieścić $x$ na brzegu odpowiedniej homotetii $A$.
\end{center}

\begin{przyklad}[Seminorma nie musi być normą]
Dla $X=\R^2$ i $A=\{(x,y):|x|<1\}$ zachodzi
$p_A(x,y)=|x|$. Zbiór $A$ jest otwarty, absolutnie
wypukły i pochłaniający, lecz nie jest ograniczony,
a $p_A(0,1)=0$. Takie zdegenerowanie nie przeczy
Twierdzeniu~\ref{thm:minkowski}.
\end{przyklad}

\begin{przyklad}[Pochłanianie nie gwarantuje ciągłości]
Niech $X=c_{00}$ z normą $\|x\|_\infty$ oraz
$A=\{x\in c_{00}:\sup_{n\ge1}n|x_n|\le1\}$.
Każdy $x\in c_{00}$ ma skończony nośnik, więc $A$
jest pochłaniający i absolutnie wypukły. Jego funkcjonał
to $p_A(x)=\sup_{n\ge1}n|x_n|$. Nie jest on ciągły
w topologii normy supremum: $x^{(n)}=e_n/n\to0$,
ale $p_A(x^{(n)})=1$. W szczególności $A$ nie jest
otoczeniem zera w wyjściowej topologii.
\end{przyklad}

\subsection{Charakteryzacja lokalnej wypukłości przez seminormy}

\begin{twierdzenie}[Charakteryzacja lokalnej wypukłości]
\label{thm:lcs-seminorms}
Dla dowolnej TVS $X$ nad $\K$ następujące warunki są
równoważne: (i) $X$ jest lokalnie wypukła;
(ii) topologia $X$ jest generowana przez rodzinę
ciągłych seminorm. W przypadku Hausdorffa rodzinę
można wybrać rozdzielającą punkty.
\end{twierdzenie}
\begin{proof}
Implikację (ii)$\Rightarrow$(i) daje
Twierdzenie~\ref{thm:seminorm-topology}.
Załóżmy (i) i weźmy rodzinę $\mathcal U$ wszystkich
otwartych absolutnie wypukłych otoczeń zera.
Dla każdego $U\in\mathcal U$ określmy funkcjonał
Minkowskiego $p_U$ z następnej podsekcji. Dowód
Twierdzenia~\ref{thm:minkowski} pokazuje, że $p_U$
jest ciągłą seminormą i $U=\{p_U<1\}$. Zatem
rodzina $(p_U)_{U\in\mathcal U}$ generuje wyjściową
topologię. Jeśli $X$ jest Hausdorffa, dla $x\ne0$
możemy wybrać $U\in\mathcal U$ niezawierające $x$,
a wówczas $p_U(x)\ge1$. Żaden niezerowy wektor
nie należy do wszystkich jąder.
\end{proof}


\subsection{Ciągłość odwzorowań liniowych i ograniczoność}

\begin{twierdzenie}[Kryterium ciągłości przez seminormy]
\label{thm:continuous-sem}
Niech $X,Y$ będą przestrzeniami lokalnie wypukłymi,
$\mathcal P$ i $\mathcal Q$ rodzinami seminorm generującymi
odpowiednio ich topologie, a $T:X\to Y$ odwzorowaniem liniowym.
Następujące warunki są równoważne:
\begin{enumerate}[label=\textup{(\roman*)}]
\item $T$ jest ciągłe;
\item dla każdej $q\in\mathcal Q$ istnieją $m\ge1$,
$p_1,\ldots,p_m\in\mathcal P$ i $C>0$ takie, że
$q(Tx)\le C\max_{1\le j\le m}p_j(x)$ dla wszystkich $x\in X$.
\end{enumerate}
\end{twierdzenie}
\begin{proof}
Załóżmy ciągłość $T$. Dla $q\in\mathcal Q$
znajdźmy $p_1,\ldots,p_m$ i $\delta>0$ takie, że
$\max_jp_j(x)<\delta$ implikuje $q(Tx)<1$.
Oznaczmy $r(x)=\max_jp_j(x)$. Jeżeli $r(x)=0$,
to $r(tx)=0$ dla każdego $t>0$, a więc
$tq(Tx)<1$ dla każdego $t>0$; stąd $q(Tx)=0$.
Jeżeli $r(x)>0$, to dla
$t=\delta/(2r(x))$ mamy $r(tx)=\delta/2$,
a zatem $q(Tx)<2r(x)/\delta$. Uzyskujemy (ii)
z $C=2/\delta$. Odwrotnie, każda z nierówności
z (ii) zapewnia ciągłość $q\circ T$ w zerze;
opis bazy otoczeń z Twierdzenia~\ref{thm:seminorm-topology}
w $Y$ daje ciągłość $T$.
\end{proof}

\begin{twierdzenie}[Ograniczoność przez seminormy]
\label{thm:bounded-sem}
Niech $X$ będzie przestrzenią lokalnie wypukłą,
a $\mathcal P$ rodziną seminorm generującą topologię $X$.
Zbiór $B\subseteq X$ jest ograniczony w sensie von Neumanna
wtedy i tylko wtedy, gdy
$\sup_{x\in B}p(x)<\infty$ dla każdego $p\in\mathcal P$
(przyjmujemy $\sup\varnothing=0$).
\end{twierdzenie}
\begin{proof}
Jeśli $B$ jest ograniczony, to dla danego $p$ i
$V=\{x:p(x)<1\}$ istnieje $t>0$ z $B\subseteq tV$,
więc $\sup_Bp\le t$.
Odwrotnie, niech $U$ będzie otoczeniem zera.
Wybierzmy $p_1,\ldots,p_m\in\mathcal P$ i
$\varepsilon>0$ z
$\{x:\max_jp_j(x)<\varepsilon\}\subseteq U$.
Niech $M=\max_j\sup_{x\in B}p_j(x)<\infty$.
Dla $t>(M/\varepsilon)$ mamy
$B\subseteq t\{x:\max_jp_j(x)<\varepsilon\}\subseteq tU$.
\end{proof}

\begin{wniosek}[Otoczki ograniczone]
\label{cor:absconv-bounded}
Jeżeli $X$ jest przestrzenią lokalnie wypukłą,
a $B\subseteq X$ jest ograniczony, to jego otoczka wypukła
i jego otoczka absolutnie wypukła są ograniczone.
\end{wniosek}
\begin{proof}
Dla każdej ciągłej seminormy $p$ i dla kombinacji
$z=\sum_{j=1}^m\lambda_jx_j$ z $x_j\in B$ oraz
$\sum_j|\lambda_j|\le1$ zachodzi
$p(z)\le\sup_{x\in B}p(x)$. Takie kombinacje
wypełniają otoczkę absolutnie wypukłą $B$.
Stosujemy Twierdzenie~\ref{thm:bounded-sem}.
\end{proof}

\subsection{Twierdzenie Kolmogorowa o normowalności}

\begin{twierdzenie}[Kryterium normowalności]
\label{thm:kolmogorov-normability}
Niech $X$ będzie lokalnie wypukłą przestrzenią
liniowo-topologiczną Hausdorffa nad $\K$.
Następujące warunki są równoważne:
\begin{enumerate}[label=\textup{(\roman*)}]
\item topologia $X$ pochodzi od pewnej normy na $X$;
\item $X$ ma ograniczone otoczenie zera;
\item $X$ ma ograniczone, otwarte i absolutnie wypukłe otoczenie zera.
\end{enumerate}
\end{twierdzenie}
\begin{proof}
Jeżeli $X$ jest unormowana, jej otwarta kula jednostkowa
jest otoczeniem ograniczonym. To daje (i)$\Rightarrow$(ii).
Dzięki lokalnej wypukłości wewnątrz ograniczonego
otoczenia możemy wybrać otwarte absolutnie wypukłe
otoczenie $U$, które jest ograniczone jako podzbiór
zbioru ograniczonego. Mamy (ii)$\Rightarrow$(iii).

Załóżmy (iii). Funkcjonał Minkowskiego $p_U$ jest
ciągłą seminormą i $U=\{p_U<1\}$.
Pokażemy, że jest normą. Jeśli $p_U(x)=0$,
to dla każdego $s>0$ zachodzi $x\in sU$.
Niech $V$ będzie dowolnym otoczeniem zera w $X$.
Ograniczoność $U$ daje $t>0$ z $U\subseteq tV$,
zatem $x\in t^{-1}U\subseteq V$.
Ponieważ $X$ jest Hausdorffa, przecięcie wszystkich
otoczeń zera jest równe $\{0\}$, czyli $x=0$.

Topologia normy $p_U$ nie jest silniejsza od wyjściowej,
ponieważ $p_U$ jest ciągła. Odwrotnie, dla dowolnego
otoczenia $V$ ograniczoność $U$ daje
$t^{-1}U\subseteq V$, a zbiór ten jest kulą
$\{p_U<1/t\}$. Zatem obie topologie są równe.
\end{proof}

\begin{uwaga}[Gdzie konieczna jest własność Hausdorffa?]
Jeżeli lokalna wypukłość jest spełniona, ale przestrzeń
nie jest Hausdorffa, funkcjonał Minkowskiego ograniczonego
otoczenia może być tylko seminormą z niezerowym jądrem.
Topologia generowana przez seminormę nieoddzielającą punktów
nie może pochodzić od normy.
\end{uwaga}

\subsection{Metryzowalność lokalnie wypukła i przestrzenie Frécheta}

\begin{twierdzenie}[Przeliczalna rodzina seminorm]
\label{thm:countable-seminorms}
Niech $X$ będzie lokalnie wypukłą TVS Hausdorffa.
Jest ona metryzowalna wtedy i tylko wtedy, gdy jej topologię
generuje przeliczalna rodzina seminorm $(p_n)_{n\ge1}$.
W takim przypadku zgodną, niezmienniczą względem przesunięć
metryką jest
\[
 d(x,y)=\sum_{n=1}^{\infty}2^{-n}\min\{1,p_n(x-y)\}.
\]
\end{twierdzenie}
\begin{proof}
Jeżeli istnieje przeliczalna rodzina $(p_n)$,
to każda funkcja $\min\{1,p_n(x-y)\}$ spełnia
nierówność trójkąta, a szereg jest jednostajnie zbieżny
(bo jego wyrazy nie przekraczają $2^{-n}$).
Ponieważ seminormy rozdzielają punkty, $d$ jest metryką.
Jeśli $x_k\to x$ w topologii seminorm, to dla ustalonego
$N$ pierwszych $N$ składników sumy dążą do zera,
a ogon jest jednostajnie mniejszy od $2^{-N}$;
stąd $d(x_k,x)\to0$. W drugą stronę z
$d(x_k,x)\ge2^{-n}\min\{1,p_n(x_k-x)\}$
wynika zbieżność każdej seminormy. Ten sam argument
zamiast ciągów stosuje się do otoczeń i dowodzi
równości topologii.

Jeżeli $X$ jest metryzowalna, posiada przeliczalną bazę
otoczeń zera (kryterium metryzowalności TVS z \hyperref[sec:rodziny-metryzowalnosc]{rozdziału 4}).
Do każdego jej elementu wybierzmy, korzystając z
Twierdzenia~\ref{thm:lcs-seminorms}, skończoną rodzinę
ciągłych seminorm i wspólny promień tak, aby odpowiadający
zbiór był zawarty w danym otoczeniu. Suma tych skończonych
rodzin jest przeliczalna i generuje topologię $X$.
\end{proof}

\begin{definicja}[Przestrzeń Frécheta]
\label{def:frechet-space}
\emph{Przestrzenią Frécheta} nazywamy przestrzeń
liniowo-topologiczną, która jest lokalnie wypukła,
metryzowalna, zupełna i Hausdorffa. Nie wymagamy, by
metryka generująca topologię była zadana normą.
\end{definicja}

\begin{przyklad}[Produkt, który nie jest normowalny]
\label{ex:product-lcs}
Niech $X=\K^{\N}$ z topologią produktową.
Generują ją seminormy $p_n(x)=|x_n|$.
Topologia jest Hausdorffa, metryzowalna, a przestrzeń
jest zupełna jako produkt zupełnych przestrzeni $\K$.
Jest zatem przestrzenią Frécheta. Nie jest natomiast
normowalna: żadne otoczenie zera nie jest ograniczone,
co pokazaliśmy w \hyperref[sec:rodziny-metryzowalnosc]{rozdziale 4}.
\end{przyklad}

\begin{przyklad}[Gładkie funkcje i nieskończenie wiele pochodnych]
\label{ex:smooth-frechet-preview}
Niech $X=C^\infty([0,1])$ z topologią generowaną przez
$p_m(f)=\max_{0\le j\le m}\sup_{t\in[0,1]}|f^{(j)}(t)|$,
$m=0,1,\ldots$. Seminormy te rozdzielają punkty,
a topologia jest metryzowalna i lokalnie wypukła.
Zupełność wykażemy w następnym stwierdzeniu. Już teraz
można dowieść \emph{nienormowalności}. Gdyby $U$ było ograniczonym
otoczeniem zera, zawierałoby
$V=\{f:p_m(f)<\delta\}$ dla pewnych $m$ i $\delta>0$;
zbiór $V$ byłby więc ograniczony. Tymczasem dla
$N\ge2\pi$ funkcje
$f_N(t)=\frac{\delta}{2N^m}\sin(Nt)$
spełniają $p_m(f_N)\le\delta/2$, natomiast
$p_{m+1}(f_N)\ge\delta N/2\to\infty$.
To przeczy Twierdzeniu~\ref{thm:bounded-sem}.
Przestrzeń $C^\infty([0,1])$ jest więc lokalnie wypukła
oraz metryzowalna, ale nie normowalna.
\end{przyklad}


\subsection{Sieci Cauchy'ego i zupełność w języku seminorm}

\begin{twierdzenie}[Kryterium Cauchy'ego w przestrzeniach lokalnie wypukłych]
\label{thm:lcs-cauchy}
Niech $X$ będzie przestrzenią lokalnie wypukłą, a
$\mathcal P$ rodziną seminorm generującą topologię $X$.
Sieć $(x_d)_{d\in D}$ w $X$ jest siecią Cauchy'ego
w kanonicznej strukturze jednostajnej TVS dokładnie wtedy,
gdy dla każdej $p\in\mathcal P$ i każdego $\varepsilon>0$
istnieje $d_0\in D$ takie, że $p(x_d-x_e)<\varepsilon$
dla wszystkich $d,e\succeq d_0$.
W szczególności $X$ jest zupełna dokładnie wtedy, gdy każda
taka sieć jest zbieżna w $X$.
\end{twierdzenie}
\begin{proof}
Jeżeli sieć jest Cauchy'ego, zastosujmy definicję z
\hyperref[sec:zupelnosc]{rozdziału 3} do otoczenia $\{x:p(x)<\varepsilon\}$.
Otrzymujemy żądany warunek dla każdej seminormy.
Na odwrót, niech $U$ będzie dowolnym otoczeniem zera.
Znajdźmy seminormy $p_1,\ldots,p_m$ i $\varepsilon>0$
takie, że $\{x:\max_jp_j(x)<\varepsilon\}\subseteq U$.
Dla każdego $j$ mamy indeks $d_j$, po którym
$p_j(x_d-x_e)<\varepsilon$; skierowanie $D$ daje wspólny
indeks $d_0\succeq d_j$, $j=1,\ldots,m$.
Wtedy $x_d-x_e\in U$ dla $d,e\succeq d_0$.
Ostatnie zdanie jest dokładnie definicją zupełności TVS.
\end{proof}

\begin{twierdzenie}[Zupełność przestrzeni funkcji gładkich]
\label{thm:smooth-complete}
Niech $X=C^\infty([0,1])$ nad $\K$ z topologią generowaną
przez seminormy $p_m(f)=\max_{0\le j\le m}
\sup_{t\in[0,1]}|f^{(j)}(t)|$, $m=0,1,\ldots$.
Wtedy $X$ jest przestrzenią Frécheta, ale nie jest normowalna.
\end{twierdzenie}
\begin{proof}
Lokalna wypukłość, warunek Hausdorffa, metryzowalność
i brak normowalności zostały wykazane powyżej.
Dla zupełności, na mocy twierdzenia o metryzowalnych
TVS z \hyperref[sec:zupelnosc]{rozdziału 3}, wystarczy wykazać zupełność sekwencyjną.

Niech $(f_n)$ będzie ciągiem Cauchy'ego względem wszystkich
$p_m$. Dla każdego $j\ge0$ ciąg funkcji ciągłych
$(f_n^{(j)})$ jest jednostajnie Cauchy'ego na $[0,1]$.
Zupełność $C([0,1])$ w normie supremum daje
funkcję ciągłą $g_j$ taką, że $f_n^{(j)}\to g_j$
jednostajnie. Korzystając z podstawowego twierdzenia
rachunku całkowego, dla $t\in[0,1]$ mamy
\[
 f_n^{(j)}(t)-f_n^{(j)}(0)
 =\int_{[0,t]}f_n^{(j+1)}(s)\,\mathrm d\lambda(s).
\]
Jednostajna zbieżność całkowanych funkcji i wartości
w zerze pozwala przejść do granicy:
$g_j(t)-g_j(0)=\int_{[0,t]}g_{j+1}(s)\,\mathrm d\lambda(s)$.
Ponieważ $g_{j+1}$ jest ciągła, wynika stąd
$g_j'=g_{j+1}$. Indukcyjnie $g_0$ jest gładka i
$g_0^{(j)}=g_j$ dla wszystkich $j\ge0$.
Ponadto $p_m(f_n-g_0)\to0$ dla każdego $m$.
Oznacza to, że ciąg Cauchy'ego ma granicę w $X$,
a więc $X$ jest zupełna.
\end{proof}

\begin{uwaga}[Rzut na skończenie wiele pochodnych]
Z punktu widzenia konstrukcji przestrzeni funkcyjnych
$C^\infty([0,1])$ można uważać za zgodną rodzinę
funkcji należących do przestrzeni Banacha $C^m([0,1])$,
$m\ge0$, z normami $p_m$. Każdy operator
$f\mapsto f^{(j)}(t_0)$, $t_0\in[0,1]$, jest ciągłym
funkcjonałem liniowym na tej przestrzeni, gdyż
$|f^{(j)}(t_0)|\le p_j(f)$. Tę strukturę będziemy
wykorzystywać przy definiowaniu dystrybucji.
\end{uwaga}

\subsection{Podprzestrzenie, produkty i ilorazy}

\begin{stwierdzenie}[Trwałość lokalnej wypukłości]
\label{prop:lcs-constructions}
Niech $(X_i)_{i\in I}$ będzie rodziną przestrzeni lokalnie
wypukłych nad $\K$. Wówczas:
\begin{enumerate}[label=\textup{(\roman*)}]
\item każda podprzestrzeń liniowa $M\subseteq X_i$ z topologią
podprzestrzeni jest lokalnie wypukła;
\item produkt $\prod_{i\in I}X_i$ z topologią produktową
jest lokalnie wypukły;
\item dla dowolnej podprzestrzeni liniowej $M\subseteq X_i$
iloraz $X_i/M$ z topologią ilorazową jest lokalnie wypukły.
Iloraz jest Hausdorffa wtedy i tylko wtedy, gdy $M$ jest domknięta.
\end{enumerate}
\end{stwierdzenie}
\begin{proof}
W (i) przecinamy wypukłe otoczenia zera w $X_i$
z $M$. W (ii) podstawowe otoczenia zera są produktami,
w których jedynie skończenie wiele współrzędnych
jest ograniczonych do wypukłych otoczeń; takie produkty
są wypukłe. W (iii) naturalne rzutowanie
$q:X_i\to X_i/M$ jest otwarte, a obrazy wypukłych
otoczeń zera są wypukłymi otoczeniami zera w ilorazie.
Kryterium Hausdorffa udowodniono w
\hyperref[sec:konstrukcje]{rozdziale 2}.
\end{proof}

\begin{uwaga}[Dalszy kierunek]
Związek topologii z seminormami jest punktem wyjścia
teorii dualności. W następnym rozdziale sformułujemy i
udowodnimy twierdzenie Hahna--Banacha, które pozwoli
przejść od seminorm i rozdzielania punktów do istnienia
nietrywialnych ciągłych funkcjonałów liniowych.
Później wprowadzimy topologie słabe, przestrzenie Frécheta,
przestrzenie funkcji testowych $\mathcal D(\Omega)$ oraz ich
ciągły dual $\mathcal D'(\Omega)$.
\end{uwaga}


% Wersja źródłowa: audyt II, jednostka 06
\setcounter{section}{5}
\section{Twierdzenie Hahna--Banacha i ciągły dual}\label{sec:hahn-banach}

\subsection{Dlaczego potrzebujemy twierdzenia o przedłużaniu?}

W skończonym wymiarze każdy funkcjonał liniowy jest ciągły, a jego
przedłużenie z podprzestrzeni można skonstruować, uzupełniając bazę.
W ogólnej przestrzeni liniowo-topologicznej postępowanie takie daje
jedynie funkcjonał \emph{algebraiczny}: jego ciągłość nie jest
zagwarantowana. Twierdzenie Hahna--Banacha pozwala zachować
oszacowanie wyrażone przez funkcję podliniową, a w przestrzeniach
lokalnie wypukłych --- także ciągłość funkcjonału.

W całym rozdziale $\K\in\{\R,\C\}$. Przestrzeń liniowa nad $\K$
nie musi być wyposażona w topologię, dopóki nie zaznaczymy tego
jawnie. Nie przyjmujemy bez dodatkowego założenia, że przestrzeń
liniowo-topologiczna jest Hausdorffa.

\begin{definicja}[Dual algebraiczny i topologiczny]\label{def:dual-topological}
Niech $X$ będzie przestrzenią liniową nad $\K$. Przez $X^*$
oznaczamy przestrzeń wszystkich $\K$-liniowych funkcjonałów
$f:X\to\K$ (\emph{dual algebraiczny}). Jeśli $X$ jest TVS,
przez $X'$ oznaczamy przestrzeń funkcjonałów liniowych ciągłych
(\emph{dual topologiczny}). Wtedy $X'\subseteq X^*$.
\end{definicja}

\begin{stwierdzenie}[Elementarne kryterium ciągłości]
\label{prop:dual-continuity}
Niech $X$ będzie TVS nad $\K$ i $f:X\to\K$ funkcjonałem
liniowym. Wtedy $f$ jest ciągły wtedy i tylko wtedy, gdy
istnieje otoczenie $U$ zera w $X$ takie, że $|f(x)|<1$
dla każdego $x\in U$.
\end{stwierdzenie}
\begin{proof}
Ciągłość w zerze zapewnia istnienie otoczenia
$U=f^{-1}(\{z\in\K:|z|<1\})$ z wymaganą własnością.
W drugą stronę, dla $\varepsilon>0$ zbiór
$\varepsilon U$ jest otoczeniem zera i dla $x\in\varepsilon U$
mamy $|f(x)|<\varepsilon$. To dowodzi ciągłości w zerze.
Liniowość daje następnie ciągłość w każdym punkcie.
\end{proof}

\begin{przyklad}[Algebraiczny dual może być znacznie większy]
\label{ex:discontinuous-functional}
Niech $X$ będzie nieskończenie wymiarową przestrzenią
unormowaną. Wybierzmy liniowo niezależne wektory $(e_n)_{n\ge1}$
z $\|e_n\|=1$. Na ich rozpiętości określmy
$f(\sum_{j=1}^m a_je_j)=\sum_{j=1}^m j a_j$.
Po uzupełnieniu tej rodziny do bazy Hamela przestrzeni $X$
przedłużamy $f$ liniowo na całe $X$. Wtedy $f\in X^*$,
ale $f\notin X'$, ponieważ $|f(e_n)|=n\to\infty$,
podczas gdy $\|e_n\|=1$. Konstrukcja bazy Hamela
wykorzystuje aksjomat wyboru.
\end{przyklad}

\subsection{Wersja rzeczywista: funkcje podliniowe}

\begin{definicja}[Funkcja podliniowa]\label{def:sublinear}
Niech $E$ będzie przestrzenią liniową nad $\R$.
Funkcję $p:E\to\R$ nazywamy \emph{podliniową}, jeśli
$p(x+y)\le p(x)+p(y)$ dla wszystkich $x,y\in E$ oraz
$p(tx)=tp(x)$ dla wszystkich $x\in E$ i $t\ge0$.
Nie wymagamy symetrii $p(-x)=p(x)$ ani nieujemności $p(x)$.
Z definicji wynika $p(0)=0$.
\end{definicja}

\begin{lemat}[Przedłużenie o jeden wymiar]
\label{lem:hb-one-step}
Niech $E$ będzie rzeczywistą przestrzenią liniową,
$p:E\to\R$ funkcją podliniową, $M\subseteq E$ podprzestrzenią,
a $f:M\to\R$ funkcjonałem liniowym spełniającym
$f(m)\le p(m)$ dla każdego $m\in M$. Dla każdego
$x_0\in E\setminus M$ istnieje funkcjonał liniowy
$g:M+\R x_0\to\R$, który przedłuża $f$ i spełnia
$g(x)\le p(x)$ dla każdego $x\in M+\R x_0$.
\end{lemat}
\begin{proof}
Każdy $x\in M+\R x_0$ ma jednoznaczny zapis $x=m+tx_0$.
Szukamy liczby $a\in\R$ takiej, by wzór
$g(m+tx_0)=f(m)+ta$ spełniał $g\le p$.
Dla $t>0$ warunek ten, po podzieleniu przez $t$,
jest równoważny $a\le p(z+x_0)-f(z)$ dla każdego $z\in M$.
Dla $t<0$ analogicznie otrzymujemy
$a\ge f(y)-p(y-x_0)$ dla każdego $y\in M$.
Zdefiniujmy więc
\[
 L=\sup_{y\in M}\bigl(f(y)-p(y-x_0)\bigr),\qquad
 R=\inf_{z\in M}\bigl(p(z+x_0)-f(z)\bigr).
\]
Dla dowolnych $y,z\in M$, z liniowości $f$ oraz
podaddytywności $p$ mamy
\[
 f(y)+f(z)=f(y+z)\le p(y+z)
 \le p(y-x_0)+p(z+x_0).
\]
Stąd $L\le R$. Ponieważ $0\in M$, mamy też
$L\ge-p(-x_0)$ i $R\le p(x_0)$, zatem oba krańce
są skończone. Wybierzmy dowolne $a\in[L,R]$.
Wówczas wyprowadzone nierówności dają $g(m+tx_0)\le
p(m+tx_0)$, odpowiednio dla $t>0$ i $t<0$.
Dla $t=0$ korzystamy z założenia $f\le p$.
\end{proof}

\begin{twierdzenie}[Hahn--Banach, wersja rzeczywista]
\label{thm:hb-real}
Niech $E$ będzie przestrzenią liniową nad $\R$,
$p:E\to\R$ funkcją podliniową, $M\subseteq E$ podprzestrzenią,
a $f:M\to\R$ funkcjonałem liniowym takim, że
$f(m)\le p(m)$ dla każdego $m\in M$.
Wtedy istnieje liniowy $F:E\to\R$ spełniający
$F|_M=f$ oraz $F(x)\le p(x)$ dla każdego $x\in E$.
\end{twierdzenie}
\begin{proof}
Rozważmy zbiór par $(N,g)$, gdzie $N$ jest podprzestrzenią
$E$ zawierającą $M$, a $g:N\to\R$ jest liniowy, przedłuża
$f$ i spełnia $g\le p$ na $N$. Jest on niepusty, ponieważ
zawiera $(M,f)$. Porządkujemy pary przez przedłużanie:
$(N_1,g_1)\preceq(N_2,g_2)$, gdy $N_1\subseteq N_2$
i $g_2|_{N_1}=g_1$.

Jeżeli $\mathcal C$ jest łańcuchem, to $N_\mathcal C=
\bigcup_{(N,g)\in\mathcal C}N$ jest podprzestrzenią:
dwa dowolne jej elementy należą do jednego z dwóch
porównywalnych członów łańcucha. Funkcjonały z łańcucha
zgadzają się na przecięciach, więc określają wspólny
funkcjonał $g_\mathcal C$ na $N_\mathcal C$.
Jest on liniowy, przedłuża $f$ i nadal spełnia
$g_\mathcal C\le p$. Zatem każdy łańcuch ma ograniczenie
górne. Lemat Kuratowskiego--Zorna daje element maksymalny
$(N_0,g_0)$. Gdyby $N_0\ne E$, lemat
\ref{lem:hb-one-step} przedłużyłby $g_0$ do
$N_0+\R x_0$ dla $x_0\notin N_0$, co przeczyłoby
maksymalności. Wobec tego $N_0=E$ i $F=g_0$.
\end{proof}

\begin{uwaga}[O roli aksjomatu wyboru]
Dowód odwołuje się do lematu Kuratowskiego--Zorna,
który w zwykłej teorii mnogości ZFC wynika z aksjomatu
wyboru. Przy konstrukcji przedłużeń w skończonym wymiarze
wystarcza skończona iteracja Lematu~\ref{lem:hb-one-step}.
\end{uwaga}

\subsection{Wersja zespolona i dominacja przez seminormę}

W przypadku zespolonym nie porównujemy liczb zespolonych
nierównością $F(x)\le p(x)$. Zamiast tego wymagamy
$|F(x)|\le p(x)$, gdzie $p$ jest seminormą.

\begin{lemat}[Odtworzenie funkcjonału zespolonego]
\label{lem:complex-real-part}
Niech $E$ będzie przestrzenią zespoloną traktowaną także
jako przestrzeń rzeczywista. Każdemu funkcjonałowi
rzeczywiście liniowemu $u:E\to\R$ odpowiada dokładnie
jeden zespolono-liniowy funkcjonał $G:E\to\C$ spełniający
$\operatorname{Re}G=u$. Jest on dany wzorem
$G(x)=u(x)-i\,u(ix)$.
\end{lemat}
\begin{proof}
Z liniowości rzeczywistej $u$ wynika
$G(ix)=u(ix)-i u(-x)=iG(x)$ oraz
$G(ax+by)=aG(x)+bG(y)$ dla $a,b\in\R$.
Każdy skalar zespolony ma postać $a+ib$, więc $G$ jest
zespolono-liniowy. Ponadto $\operatorname{Re}G(x)=u(x)$.
Jeśli $H$ jest innym funkcjonałem zespolono-liniowym
o tej samej części rzeczywistej, to
$\operatorname{Im}H(x)=-\operatorname{Re}H(ix)=-u(ix)$;
stąd $H(x)=G(x)$.
\end{proof}

\begin{twierdzenie}[Hahn--Banach, dominacja przez seminormę]
\label{thm:hb-complex-seminorm}
Niech $E$ będzie przestrzenią liniową nad
$\K\in\{\R,\C\}$, $p:E\to[0,\infty)$ seminormą,
$M\subseteq E$ podprzestrzenią liniową, a
$f:M\to\K$ funkcjonałem $\K$-liniowym spełniającym
$|f(m)|\le p(m)$ dla każdego $m\in M$.
Wtedy istnieje $\K$-liniowy $F:E\to\K$ taki, że
$F|_M=f$ i $|F(x)|\le p(x)$ dla każdego $x\in E$.
\end{twierdzenie}
\begin{proof}
Dla $\K=\R$ stosujemy Twierdzenie~\ref{thm:hb-real}
do $p$ i $f$; dominacja $F(x)\le p(x)$ zastosowana
również do $-x$ daje $|F(x)|\le p(x)$.

Niech $\K=\C$. Na przestrzeni rzeczywistej
$E_{\R}$ funkcja $u=\operatorname{Re}f$ jest liniowa
na $M_{\R}$ i spełnia $u(m)\le p(m)$.
Twierdzenie~\ref{thm:hb-real} zapewnia przedłużenie
rzeczywiście liniowe $v:E_{\R}\to\R$ z $v\le p$.
Z Lematu~\ref{lem:complex-real-part} otrzymujemy
$F(x)=v(x)-iv(ix)$. Na $M$ funkcjonały $F$ i $f$
mają tę samą część rzeczywistą, zatem są równe.

Pozostaje sprawdzić \emph{moduł} oszacowania. Dla
$x\in E$ z $F(x)\ne0$ wybierzmy $\alpha\in\C$,
$|\alpha|=1$, w taki sposób, aby
$\alpha F(x)=|F(x)|$. Wówczas
\[
 |F(x)|=\operatorname{Re}F(\alpha x)
 =v(\alpha x)\le p(\alpha x)=p(x).
\]
Gdy $F(x)=0$, nierówność jest oczywista.
\end{proof}

\begin{uwaga}[Co zostaje zachowane przy przedłużaniu?]
Przedłużenie zachowuje dominację przez \emph{tę samą}
seminormę, bez dodatkowego współczynnika. Samo twierdzenie
jest algebraiczne: $E$ nie musi być przestrzenią topologiczną.
Dopiero ciągłość $p$ daje automatycznie ciągłość $F$.
\end{uwaga}

\subsection{Hahn--Banach w przestrzeniach lokalnie wypukłych}

\begin{stwierdzenie}[Ciągłość funkcjonału a seminorma]
\label{prop:dual-seminorm-criterion}
Niech $X$ będzie lokalnie wypukłą TVS nad $\K$,
a $f:X\to\K$ funkcjonałem liniowym. Wtedy
$f\in X'$ dokładnie wtedy, gdy istnieje ciągła
seminorma $p$ na $X$ taka, że
$|f(x)|\le p(x)$ dla każdego $x\in X$.
\end{stwierdzenie}
\begin{proof}
Jeśli $f$ jest ciągły, to $p(x)=|f(x)|$ jest
ciągłą seminormą i $|f|\le p$. Odwrotnie, gdy
$|f|\le p$ i $p$ jest ciągła, to
$f^{-1}(\{z:|z|<\varepsilon\})$ zawiera otwarte
otoczenie $\{x:p(x)<\varepsilon\}$ dla każdego
$\varepsilon>0$. Wynika stąd ciągłość $f$.
\end{proof}

\begin{twierdzenie}[Topologiczny Hahn--Banach]\label{thm:hb-topological}
Niech $X$ będzie lokalnie wypukłą przestrzenią
liniowo-topologiczną nad $\K\in\{\R,\C\}$,
$M\subseteq X$ podprzestrzenią liniową z topologią
indukowaną z $X$, a $f:M\to\K$ ciągłym funkcjonałem
liniowym. Istnieje $F\in X'$ takie, że $F|_M=f$.
Nie trzeba zakładać, że $M$ jest domknięta, a $X$
jest Hausdorffa lub zupełna.
\end{twierdzenie}
\begin{proof}
Ciągłość $f$ w zerze oznacza, że istnieje otoczenie
$W$ zera w $X$ takie, iż $|f(m)|<1$ dla wszystkich
$m\in M\cap W$. Korzystając z lokalnej wypukłości,
wybieramy otwarte absolutnie wypukłe otoczenie
$U\subseteq W$. Jego funkcjonał Minkowskiego $p_U$
jest ciągłą seminormą na $X$ na mocy
twierdzenia o funkcjonale Minkowskiego z \hyperref[sec:lokalnie-wypukle]{rozdziału 5}.

Jeżeli $m\in M$ i $t>p_U(m)$, to $m\in tU$,
zatem $m/t\in M\cap U$ i $|f(m)|<t$.
Ponieważ jest to prawdą dla każdego $t>p_U(m)$,
uzyskujemy $|f(m)|\le p_U(m)$.
Stosując Twierdzenie~\ref{thm:hb-complex-seminorm}
do $p_U$ i podprzestrzeni $M$, otrzymujemy funkcjonał
$F:X\to\K$ z $F|_M=f$ i $|F|\le p_U$.
Ciągłość $p_U$ implikuje ciągłość $F$.
\end{proof}

\begin{wniosek}[Przedłużanie z podprzestrzeni gęstej]
\label{cor:dense-extension}
Jeśli w założeniach Twierdzenia~\ref{thm:hb-topological}
podprzestrzeń $M$ jest gęsta w $X$, to ciągłe
przedłużenie $f$ jest jednoznaczne.
\end{wniosek}
\begin{proof}
Istnienie wynika z twierdzenia. Jeśli $F_1,F_2\in X'$
przedłużają $f$, to ciągły funkcjonał $F_1-F_2$
zeruje się na gęstym zbiorze $M$. Jego zbiór zerowy
jest domknięty (bo $\K$ jest Hausdorffa), więc
$F_1=F_2$ na całym $X$.
\end{proof}

\begin{przyklad}[Sumowanie w $\ell^1$ i jednoznaczność]
Niech $X=\ell^1(\N)$ z normą
$\|x\|_1=\sum_{n\ge1}|x_n|$, a $M=c_{00}\subseteq X$
niech oznacza przestrzeń ciągów o skończonym nośniku.
Przez $e_n$ oznaczamy ciąg mający $1$ na miejscu $n$
i zera na pozostałych miejscach. Funkcjonał $f(x)=\sum_{n\ge1}x_n$ jest dobrze określony
na $M$ i ma normę $1$, bo $|f(x)|\le\|x\|_1$
i $f(e_1)=\|e_1\|_1=1$. Podprzestrzeń $M$ jest gęsta
w $X$: obcięcia $x^{(N)}=(x_1,\ldots,x_N,0,\ldots)$
spełniają $\|x^{(N)}-x\|_1\to0$. Jedynym ciągłym
przedłużeniem $f$ jest więc
$F(x)=\sum_{n=1}^\infty x_n$ na $\ell^1$.
Dla porównania, ten sam wzór na $c_{00}$ z normą
supremum nie jest ciągły: dla
$x^{(N)}=\sum_{n=1}^N e_n$ zachodzi
$\|x^{(N)}\|_\infty=1$ oraz $f(x^{(N)})=N$.
\end{przyklad}

\begin{przyklad}[Funkcjonały na przestrzeni funkcji gładkich]
Niech $X=C^\infty([0,1])$ z topologią seminorm
$p_m$ z \hyperref[sec:lokalnie-wypukle]{rozdziału 5}. Dla $t_0\in[0,1]$ i
$j\in\N\cup\{0\}$ funkcjonał
$f\mapsto f^{(j)}(t_0)$ jest ciągły, gdyż
$|f^{(j)}(t_0)|\le p_j(f)$. Każdy taki funkcjonał
będzie później przykładem dystrybucji na odpowiedniej
przestrzeni funkcji testowych.
\end{przyklad}

\subsection{Przestrzenie unormowane: zachowanie normy}

\begin{twierdzenie}[Hahn--Banach dla przestrzeni unormowanych]
\label{thm:hb-normed}
Niech $X$ będzie przestrzenią unormowaną nad $\K$,
$M\subseteq X$ podprzestrzenią liniową, a
$f:M\to\K$ ograniczonym funkcjonałem liniowym
(względem normy odziedziczonej z $X$).
Istnieje funkcjonał ograniczony $F:X\to\K$ taki, że
$F|_M=f$ i $\|F\|=\|f\|$, gdzie
$\|f\|=\sup\{|f(m)|:\|m\|\le1,\ m\in M\}$.
\end{twierdzenie}
\begin{proof}
W Twierdzeniu~\ref{thm:hb-complex-seminorm}
weźmy $p(x)=\|f\|\,\|x\|$. Otrzymujemy $F$
z oszacowaniem $|F(x)|\le\|f\|\,\|x\|$,
czyli $\|F\|\le\|f\|$. Ponieważ $F|_M=f$,
mamy także $\|F\|\ge\|f\|$. Zatem normy są równe.
Nie wymagamy zupełności $X$ ani domkniętości $M$.
\end{proof}

\begin{wniosek}[Funkcjonał normujący wektor]\label{cor:norming}
Niech $X$ będzie przestrzenią unormowaną nad $\K$,
a $x_0\in X\setminus\{0\}$. Istnieje $f\in X'$
taki, że $\|f\|=1$ i $f(x_0)=\|x_0\|$.
W szczególności, dla każdego $x\in X$ zachodzi
\[
 \|x\|=\sup\bigl\{|f(x)|:f\in X',\ \|f\|\le1\bigr\}.
\]
\end{wniosek}
\begin{proof}
Na $M=\K x_0$ określmy
$f_0(\lambda x_0)=\lambda\|x_0\|$.
Wówczas $\|f_0\|=1$ i z
Twierdzenia~\ref{thm:hb-normed} otrzymujemy przedłużenie
$f$ o normie $1$ i wartości $f(x_0)=\|x_0\|$.
Dla dowolnego $x$ nierówność $|g(x)|\le\|x\|$
dla $\|g\|\le1$ wynika z definicji normy operatorowej,
a dla $x\ne0$ równość zostaje osiągnięta przez
skonstruowany funkcjonał. Dla $x=0$ obie strony są zerowe.
\end{proof}

\begin{przyklad}[Przedłużenie nie musi być jednoznaczne]
Niech $X=\R^2$ z normą $\|(x,y)\|_\infty
=\max(|x|,|y|)$, $M=\{(t,t):t\in\R\}$, a
$f(t,t)=t$. Wtedy $\|f\|=1$. Dla każdego
$a\in[0,1]$ wzór
$F_a(x,y)=a x+(1-a)y$ określa przedłużenie $f$
spełniające $\|F_a\|=|a|+|1-a|=1$.
Twierdzenie Hahna--Banacha nie daje zatem
jednoznaczności bez dodatkowych założeń.
\end{przyklad}

\subsection{Separacja punktów, podprzestrzeni i rola lokalnej wypukłości}

\begin{twierdzenie}[Seminorma jako supremum funkcjonałów]
\label{thm:seminorm-as-sup}
Niech $X$ będzie lokalnie wypukłą TVS nad $\K$,
$p$ ciągłą seminormą na $X$, a $x_0\in X$.
Wtedy
\[
 p(x_0)=\sup\bigl\{|f(x_0)|:f\in X',\quad
 |f(x)|\le p(x)\ \text{dla wszystkich }x\in X\bigr\}.
\]
Ponadto istnieje $f\in X'$ spełniający $|f|\le p$
i $f(x_0)=p(x_0)$.
\end{twierdzenie}
\begin{proof}
Jeśli $p(x_0)=0$, wybieramy funkcjonał zerowy.
Gdy $p(x_0)>0$, na $M=\K x_0$ określamy
$f_0(\lambda x_0)=\lambda p(x_0)$.
Z jednorodności seminormy wynika
$|f_0(\lambda x_0)|=p(\lambda x_0)$.
Twierdzenie~\ref{thm:hb-complex-seminorm} daje
przedłużenie $f$ z $|f|\le p$, a ciągłość $p$
gwarantuje $f\in X'$. Żaden funkcjonał w podanym
supremum nie może mieć wartości większej niż $p(x_0)$,
a skonstruowany osiąga tę wartość.
\end{proof}

\begin{wniosek}[Dual rozdziela punkty]\label{cor:dual-separates}
Niech $X$ będzie lokalnie wypukłą TVS Hausdorffa
nad $\K$. Dla każdych różnych $x,y\in X$ istnieje
$f\in X'$ takie, że $f(x)\ne f(y)$.
\end{wniosek}
\begin{proof}
Dla $z=x-y\ne0$ istnieje ciągła seminorma $p$
z $p(z)>0$, ponieważ seminormy generujące topologię
rozdzielają punkty (twierdzenie o topologii seminorm z \hyperref[sec:lokalnie-wypukle]{rozdziału 5}).
Twierdzenie~\ref{thm:seminorm-as-sup} daje $f\in X'$
z $f(z)=p(z)>0$, co kończy dowód.
\end{proof}

\begin{twierdzenie}[Oddzielanie punktu od domkniętej podprzestrzeni]
\label{thm:point-closed-subspace}
Niech $X$ będzie lokalnie wypukłą TVS nad $\K$,
$M\subseteq X$ domkniętą podprzestrzenią liniową,
a $x_0\in X\setminus M$. Istnieje $f\in X'$
takie, że $f|_M=0$ i $f(x_0)=1$.
Nie jest potrzebne założenie Hausdorffa całej przestrzeni $X$.
\end{twierdzenie}
\begin{proof}
Ponieważ $M$ jest domknięta, istnieje otoczenie
$W$ zera takie, że $(x_0+W)\cap M=\varnothing$.
Wybierzmy otwarte absolutnie wypukłe otoczenie
$U\subseteq W$ i oznaczmy jego ciągłą seminormę
Minkowskiego przez $p=p_U$.
Dla każdego $m\in M$ mamy $x_0-m\notin U$,
a więc $p(x_0-m)\ge1$, bo $U=\{x:p(x)<1\}$.
Na podprzestrzeni $M+\K x_0$ określmy
$g(m+\lambda x_0)=\lambda$. Dla $\lambda\ne0$,
jednorodność $p$ oraz $-m/\lambda\in M$ dają
\[
 p(m+\lambda x_0)=|\lambda|\,
 p\bigl(x_0+m/\lambda\bigr)\ge|\lambda|
 =|g(m+\lambda x_0)|.
\]
Dla $\lambda=0$ nierówność też zachodzi.
Z Twierdzenia~\ref{thm:hb-complex-seminorm} otrzymujemy
przedłużenie $f$ z $|f|\le p$. Jest ono ciągłe,
zeruje się na $M$ i spełnia $f(x_0)=1$.
\end{proof}

\begin{wniosek}[Charakteryzacja domkniętych podprzestrzeni]
\label{cor:closed-as-annihilator}
Niech $X$ będzie lokalnie wypukłą TVS nad $\K$,
a $M\subseteq X$ podprzestrzenią liniową. Wtedy
\[
 \overline M=\bigcap\{\ker f:f\in X',\ f|_M=0\}.
\]
W szczególności $M$ jest domknięta dokładnie wtedy,
gdy stanowi przecięcie jąder pewnej rodziny ciągłych
funkcjonałów liniowych.
\end{wniosek}
\begin{proof}
Dla każdego $f\in X'$ zerującego się na $M$
jego domknięte jądro zawiera $\overline M$.
Jeśli $x_0\notin\overline M$, zastosujmy
Twierdzenie~\ref{thm:point-closed-subspace}
do domkniętej podprzestrzeni $\overline M$.
Otrzymujemy funkcjonał z $f(x_0)=1$, który
wyklucza $x_0$ z podanego przecięcia.
\end{proof}

\begin{przyklad}[Dlaczego nie możemy pominąć lokalnej wypukłości?]
Niech $X=L^p([0,1])$, $0<p<1$, z topologią metryki
$d(f,g)=\int_{[0,1]}|f(t)-g(t)|^p\,\mathrm d\lambda(t)$.
Z \hyperref[sec:tvs]{rozdziału 1} wiemy, że $X'=\{0\}$.
Tymczasem $M=\K\mathbf 1$ jest jednowymiarową
podprzestrzenią, a funkcjonał $f(c\mathbf1)=c$
jest na niej ciągły: $d(c\mathbf1,0)=|c|^p$.
Nie istnieje ciągłe przedłużenie $f$ na $X$.
To pokazuje, że topologiczna wersja Hahna--Banacha
wymaga dodatkowych założeń; lokalna wypukłość jest
wystarczającym, lecz nie zawsze koniecznym warunkiem.
\end{przyklad}

\begin{stwierdzenie}[Dual topologiczny przestrzeni produktowej]
\label{prop:product-dual}
Niech $I$ będzie dowolnym zbiorem, a $X=\K^I$
przestrzenią liniowo-topologiczną z topologią produktową.
Każdy funkcjonał $f\in X'$ ma postać
$f(x)=\sum_{i\in I_0}a_i x_i$, gdzie $I_0\subseteq I$
jest skończony, a $a_i\in\K$. Odwrotnie, każdy taki
funkcjonał jest ciągły.
\end{stwierdzenie}
\begin{proof}
Niech $f\in X'$. Z ciągłości w zerze istnieje
podstawowe otoczenie produktowe $U$ zera, na którym
$|f(x)|<1$. Możemy wybrać je tak, aby ograniczało
współrzędne tylko dla indeksów z pewnego skończonego
$I_0\subseteq I$. Jeśli $x_i=0$ dla $i\in I_0$,
to $tx\in U$ dla każdego $t\in\K$, zatem
$|t f(x)|<1$ dla wszystkich $t$; koniecznie $f(x)=0$.
Niech $e_i\in\K^I$ oznacza wektor o współrzędnej
$1$ na pozycji $i$ i zerach na pozostałych pozycjach.
Każde $x\in X$ jest sumą wektora
$\sum_{i\in I_0}x_i e_i$ i wektora o zerowych
współrzędnych w $I_0$. Stąd
$f(x)=\sum_{i\in I_0}x_i f(e_i)$.
W drugą stronę, rzuty współrzędnych $x\mapsto x_i$
są ciągłe z definicji topologii produktowej;
skończona kombinacja liniowa też jest ciągła.
\end{proof}

\begin{uwaga}[Jak silne jest to ograniczenie?]
Gdy $I=\N$, na przestrzeni wszystkich ciągów
$\K^\N$ z topologią produktową żaden ciągły funkcjonał
nie może zależeć istotnie od nieskończenie wielu wyrazów.
Inaczej jest na $\ell^1$: funkcjonał sumowania z
poprzedniego przykładu jest ciągły.
Różnica pochodzi z topologii, nie z samej liniowości.
\end{uwaga}

\subsection{Norma odległości i anihilatory}

\begin{definicja}[Anihilator]\label{def:annihilator}
Niech $X$ będzie TVS nad $\K$, a $M\subseteq X$
podprzestrzenią liniową. Jej \emph{anihilatorem}
w dualu topologicznym nazywamy
$M^\perp=\{f\in X':f(m)=0\ \text{dla wszystkich }m\in M\}$.
Zapis $M^\perp$ nie wymaga iloczynu skalarnego.
\end{definicja}

\begin{stwierdzenie}[Dual ilorazu]\label{prop:quotient-dual}
Niech $X$ będzie TVS nad $\K$, a $M\subseteq X$
podprzestrzenią liniową. Niech $q:X\to X/M$
będzie rzutowaniem z topologią ilorazową.
Odwzorowanie $q^*:(X/M)'\to X'$ określone przez
$q^*(\varphi)=\varphi\circ q$ jest liniową bijekcją
z $(X/M)'$ na $M^\perp$.
\end{stwierdzenie}
\begin{proof}
Jeśli $\varphi$ jest ciągłym funkcjonałem na ilorazie,
to $\varphi\circ q$ jest ciągły i zeruje się na $M$.
Odwzorowanie $q^*$ jest różnowartościowe, bo $q$
jest surjekcją. Jeżeli $f\in M^\perp$, to wzór
$\varphi(x+M)=f(x)$ jest dobrze określony, liniowy
i spełnia $\varphi\circ q=f$.
Z własności topologii ilorazowej (\hyperref[sec:konstrukcje]{rozdział 2})
wynika ciągłość $\varphi$, gdyż $f$ jest ciągły.
\end{proof}

\begin{twierdzenie}[Odległość od podprzestrzeni i dual]\label{thm:distance-dual}
Niech $X$ będzie przestrzenią unormowaną nad $\K$,
$M\subseteq X$ domkniętą podprzestrzenią liniową,
a $x_0\in X$. Wtedy
\[
 \operatorname{dist}(x_0,M):=\inf_{m\in M}\|x_0-m\|
 =\sup\{|f(x_0)|:f\in M^\perp,\ \|f\|\le1\}.
\]
Jeśli $x_0\notin M$, istnieje $f\in M^\perp$
o normie $1$, dla którego $f(x_0)=\operatorname{dist}(x_0,M)$.
\end{twierdzenie}
\begin{proof}
Niech $d=\operatorname{dist}(x_0,M)$.
Dla $f\in M^\perp$ z $\|f\|\le1$ i $m\in M$
mamy $|f(x_0)|=|f(x_0-m)|\le\|x_0-m\|$.
Po wzięciu infimum po $m$ otrzymujemy $|f(x_0)|\le d$.
Jeśli $x_0\in M$, to $d=0$ i funkcjonał zerowy
wystarcza do równości.

Załóżmy $x_0\notin M$. Ponieważ $M$ jest domknięta,
$d>0$. Na $M+\K x_0$ określmy
$g(m+\lambda x_0)=\lambda d$.
Dla $\lambda\ne0$ mamy
$\|m+\lambda x_0\|=|\lambda|\,
\|x_0+m/\lambda\|\ge|\lambda|d=|g(m+\lambda x_0)|$;
dla $\lambda=0$ jest to oczywiste.
Z Twierdzenia~\ref{thm:hb-normed} otrzymujemy
przedłużenie $f$ z $\|f\|\le1$, $f|_M=0$ i
$f(x_0)=d$. Równość $\|f\|=1$ wynika z faktu, że
$\|g\|=1$: dla $\lambda=1$ i $m\in M$ ilorazy
$|g(x_0-m)|/\|x_0-m\|=d/\|x_0-m\|$
mają supremum $1$ z definicji $d$.
\end{proof}

\begin{uwaga}[Następne rozdziały]
Rozdzielanie \emph{ogólnych zbiorów wypukłych}, twierdzenie
bipolarne i topologie słabe wymagają dodatkowych konstrukcji.
Omówimy je osobno, opierając się na udowodnionych
wersjach twierdzenia Hahna--Banacha. W części o przestrzeniach
dystrybucji ciągły dual $\mathcal D(\Omega)'$ otrzyma
kanoniczną interpretację jako przestrzeń dystrybucji
$\mathcal D'(\Omega)$.
\end{uwaga}


% Wersja źródłowa: audyt II, jednostka 07
\setcounter{section}{6}
\section{Oddzielanie zbiorów wypukłych, polary i twierdzenie bipolarne}\label{sec:separacja-bipolarne}

\subsection{Przejście do rzeczywistych funkcjonałów}

W \hyperref[sec:hahn-banach]{rozdziale 6} udowodniliśmy twierdzenie Hahna--Banacha dla funkcjonałów liniowych.
Teraz wykorzystamy je do opisu geometrii zbiorów wypukłych i topologii,
którą można odtworzyć za pomocą przestrzeni dualnej. W szczególności odróżnimy
\emph{nierówności ścisłe dla każdego punktu} od \emph{oddzielenia z jednolitą luką}.

Przez cały rozdział $\K\in\{\R,\C\}$, a $X$ jest przestrzenią liniowo-topologiczną
nad $\K$. Gdy jest to wyraźnie zaznaczone, zakładamy lokalną wypukłość
oraz własność Hausdorffa. W topologii zespolonej oddzielanie zbiorów
odbywa się za pomocą \emph{rzeczywistych części} funkcjonałów z $X'$.

\begin{definicja}[Rzeczywiste funkcjonały i rodzaje separacji]\label{def:separacja}
Niech $A,B\subseteq X$ będą niepustymi zbiorami, a
$\ell:X\to\R$ niezerowym, ciągłym, rzeczywiście liniowym funkcjonałem
(na przestrzeni $X$ rozpatrywanej nad $\R$). Mówimy, że
\emph{oddziela} $A$ i $B$, gdy
$\sup_{a\in A}\ell(a)\le\inf_{b\in B}\ell(b)$.
Oddziela je \emph{ściśle punktowo}, gdy $\ell(a)<\ell(b)$
dla wszystkich $a\in A$, $b\in B$, a \emph{silnie}, gdy istnieje
$\delta>0$ takie, że $\sup_A\ell+\delta\le\inf_B\ell$.
Hiperpłaszczyzną afiniczną funkcjonału $\ell$ na poziomie $c\in\R$
nazywamy $H_{\ell,c}=\{x\in X:\ell(x)=c\}$.
\end{definicja}

\begin{lemat}[Rzeczywista część funkcjonału zespolonego]\label{lem:realification}
Niech $X$ będzie zespoloną TVS. Każdy ciągły, rzeczywiście liniowy
$\ell:X\to\R$ ma dokładnie jedno przedstawienie
$\ell=\operatorname{Re}F$, gdzie $F\in X'$. Mianowicie
$F(x)=\ell(x)-i\ell(ix)$.
\end{lemat}
\begin{proof}
Liniowość zespolona i jednoznaczność zostały wykazane w
lemacie o odtworzeniu funkcjonału zespolonego z \hyperref[sec:hahn-banach]{rozdziału 6}. Ponieważ $x\mapsto ix$
jest ciągły, także funkcjonał $F$ jest ciągły.
\end{proof}

\subsection{Oddzielanie zbioru otwartego od punktu}

\begin{lemat}[Funkcjonał Minkowskiego dla zbioru wypukłego]\label{lem:open-gauge}
Niech $E$ będzie rzeczywistą TVS, a $V\subseteq E$ otwartym
wypukłym otoczeniem zera (niekoniecznie symetrycznym).
Wówczas $V$ jest pochłaniający, a jego funkcjonał
$p_V(x)=\inf\{t>0:x\in tV\}$ jest skończony, podliniowy i ciągły.
Ponadto $V=\{x:p_V(x)<1\}$.
\end{lemat}
\begin{proof}
Pochłanianie wynika z ciągłości $t\mapsto tx$ w zerze,
a podliniowość z wypukłości i pochłaniania, tak jak dla
funkcjonału Minkowskiego z \hyperref[sec:lokalnie-wypukle]{rozdziału 5}; argument nie korzysta
z symetrii $V$. Z wypukłości i $0\in V$ wynika
$tV\subseteq V$ dla $0<t\le1$. Jeżeli $x\in V$, to
z otwartości $V$ i ciągłości $s\mapsto sx$ w $s=1$
można wybrać $s>1$ takie, że $sx\in V$, czyli
$p_V(x)\le1/s<1$. Odwrotnie, jeżeli $p_V(x)<1$,
można znaleźć $t<1$ z $x\in tV\subseteq V$.
Stąd $V=\{p_V<1\}$.

Dla $\varepsilon>0$ zbiór $\{p_V<\varepsilon\}=\varepsilon V$
jest otwarty. Zatem $p_V$ jest górnie półciągły. Aby
wykazać ciągłość, dla każdego $x,y\in E$ korzystamy z
$\,p_V(x)\le p_V(y)+p_V(x-y)$ i
$p_V(y)\le p_V(x)+p_V(y-x)$.
Jeżeli $x-y\in\varepsilon V\cap(-\varepsilon V)$,
to obie różnice są mniejsze niż $\varepsilon$.
Takie zbiory tworzą otoczenia zera, co daje ciągłość.
\end{proof}

\begin{twierdzenie}[Oddzielanie punktu od otwartego zbioru wypukłego]
\label{thm:open-point-sep}
Niech $X$ będzie lokalnie wypukłą TVS nad $\K$,
$C\subseteq X$ niepustym otwartym zbiorem wypukłym
(w sensie rzeczywistym), a $x_0\in X\setminus C$.
Wówczas istnieje niezerowy ciągły, rzeczywiście liniowy
$\ell:X\to\R$ taki, że $\ell(c)<\ell(x_0)$
dla każdego $c\in C$. Jeżeli $\K=\C$, można przyjąć
$\ell=\operatorname{Re}F$ z $F\in X'$.
Nie twierdzimy ogólnie, że $\sup_C\ell<\ell(x_0)$.
\end{twierdzenie}
\begin{proof}
Traktujemy $X$ jako rzeczywistą TVS; lokalna wypukłość
jest zachowana. Zbiór $D=C-x_0$ jest niepusty, otwarty
i wypukły, a $0\notin D$. Wybierzmy $a\in D$ i połóżmy
$V=D-a$. Zbiór $V$ jest otwartym, wypukłym otoczeniem zera,
a $-a\notin V$, stąd $p_V(-a)\ge1$ na mocy
Lematu~\ref{lem:open-gauge}. Na prostej $\R(-a)$
określmy $g(t(-a))=t$. Dla $t\ge0$ mamy
$g(t(-a))=t\le p_V(t(-a))$. Dla $t<0$ nierówność
także zachodzi, bo $p_V\ge0$. Rzeczywiste twierdzenie
Hahna--Banacha daje liniowe $\ell:X\to\R$ spełniające
$\ell\le p_V$ i $\ell(-a)=1$. Z ciągłości $p_V$
i nierówności $-p_V(-x)\le\ell(x)\le p_V(x)$
wynika ciągłość $\ell$.

Dla dowolnego $d\in D$ mamy $d-a\in V$, więc
$\ell(d-a)\le p_V(d-a)<1$. Ponieważ $\ell(a)=-1$,
otrzymujemy $\ell(d)<0$, czyli
$\ell(c)<\ell(x_0)$ dla $c\in C$. Równość
$\ell(-a)=1$ dowodzi, że $\ell\ne0$. W przypadku
zespolonym stosujemy Lemat~\ref{lem:realification}.
\end{proof}

\begin{twierdzenie}[Oddzielanie dwóch zbiorów wypukłych]
\label{thm:open-two-sep}
Niech $X$ będzie lokalnie wypukłą TVS nad $\K$,
a $A,B\subseteq X$ niepustymi, rozłącznymi zbiorami wypukłymi.
Jeśli $A$ jest otwarty, istnieje niezerowy, ciągły, rzeczywiście
liniowy funkcjonał $\ell$ spełniający
$\ell(a)<\ell(b)$ dla wszystkich $a\in A$, $b\in B$.
W szczególności istnieje $c\in\R$ takie, że
$\sup_A\ell\le c\le\inf_B\ell$.
\end{twierdzenie}
\begin{proof}
Zbiór $D=A-B=\{a-b:a\in A,\ b\in B\}$ jest wypukły
i otwarty (jest sumą zbioru otwartego i dowolnego),
a $0\notin D$. Z Twierdzenia~\ref{thm:open-point-sep}
dostajemy $\ell(d)<0$ dla wszystkich $d\in D$,
co daje nierówność punktową. Dla wybranego $b_0\in B$
mamy $\sup_A\ell\le\ell(b_0)<\infty$; analogicznie,
dla $a_0\in A$ mamy $\inf_B\ell\ge\ell(a_0)>-\infty$.
Stąd obie liczby są skończone i możemy wybrać $c$
pomiędzy nimi.
\end{proof}

\begin{przyklad}[Dlaczego otwartość nie daje jednolitej luki?]
Dla $X=\R$, $A=(-1,0)$, $B=\{0\}$ funkcjonał
$\ell(x)=x$ spełnia $\ell(a)<\ell(b)$ dla wszystkich
$a\in A$, $b\in B$, lecz $\sup_A\ell=\inf_B\ell=0$.
Nie istnieje silna separacja tych zbiorów, ponieważ
$0\in\overline A\cap B$.
\end{przyklad}

\subsection{Oddzielanie silne i hiperpłaszczyzny podpierające}

\begin{lemat}[Wspólne otoczenie dla zbioru zwartego]
\label{lem:compact-neighborhood}
Niech $X$ będzie TVS, $K\subseteq X$ zbiorem zwartym,
a $F\subseteq X$ zbiorem domkniętym, przy czym $K\cap F=\varnothing$.
Wtedy istnieje otoczenie $U$ zera takie, że
$(K+U)\cap F=\varnothing$. Jeżeli $X$ jest lokalnie wypukła,
można przyjąć $U$ otwarte i wypukłe.
\end{lemat}
\begin{proof}
Dla każdego $k\in K$ istnieje otoczenie $V_k$ zera
z $(k+V_k)\cap F=\varnothing$. Wybieramy otoczenie
$W_k$ zera takie, że $W_k+W_k\subseteq V_k$,
a w przypadku lokalnie wypukłym możemy żądać
jego wypukłości i otwartości. Zwartość $K$ daje
pokrycie skończone $K\subseteq\bigcup_{j=1}^m(k_j+W_{k_j})$.
Wybierzmy otoczenie $U\subseteq\bigcap_j W_{k_j}$.
Jeżeli $x\in K$ należy do $k_j+W_{k_j}$, to
$x+U\subseteq k_j+V_{k_j}$, zatem $(K+U)\cap F=\varnothing$.
\end{proof}

\begin{twierdzenie}[Silna separacja zbioru zwartego i domkniętego]
\label{thm:strong-separation}
Niech $X$ będzie lokalnie wypukłą TVS nad $\K$,
$A,B\subseteq X$ niepustymi, rozłącznymi zbiorami wypukłymi.
Jeżeli $A$ jest zwarty, a $B$ domknięty, to istnieje
niezerowy ciągły, rzeczywiście liniowy funkcjonał $\ell$
oraz $\delta>0$ takie, że
\[
\sup_{a\in A}\ell(a)+\delta\le\inf_{b\in B}\ell(b).
\]
W przypadku zespolonym $\ell=\operatorname{Re}F$ dla pewnego $F\in X'$.
\end{twierdzenie}
\begin{proof}
Na mocy Lematu~\ref{lem:compact-neighborhood} istnieje
otwarte wypukłe otoczenie $U$ zera takie, że
$(A+U)\cap B=\varnothing$. Zbiór $A+U$ jest otwarty
i wypukły. Z Twierdzenia~\ref{thm:open-two-sep}
otrzymujemy niezerowy $\ell$ z
$\ell(a+u)<\ell(b)$ dla wszystkich
$a\in A$, $u\in U$, $b\in B$. Ponieważ $\ell\ne0$
i $U$ jest otoczeniem zera, istnieje $u_0\in U$
z $\ell(u_0)>0$. Kładąc $\delta=\ell(u_0)$, uzyskujemy
$\ell(a)+\delta<\ell(b)$ dla wszystkich $a,b$,
a zatem $\sup_A\ell+\delta\le\inf_B\ell$.
\end{proof}

\begin{wniosek}[Separacja punktu i domkniętego zbioru wypukłego]
\label{cor:point-closed-convex}
Niech $X$ będzie lokalnie wypukłą TVS, $C\subseteq X$
niepustym, domkniętym i wypukłym zbiorem oraz
$x_0\in X\setminus C$. Wówczas istnieją
ciągły rzeczywiście liniowy $\ell\ne0$ i $\delta>0$
takie, że $\ell(x_0)+\delta\le\inf_{c\in C}\ell(c)$.
\end{wniosek}
\begin{proof}
Stosujemy Twierdzenie~\ref{thm:strong-separation}
do $A=\{x_0\}$ i $B=C$.
\end{proof}

\begin{twierdzenie}[Hiperpłaszczyzna podpierająca]
\label{thm:supporting-hyperplane}
Niech $X$ będzie lokalnie wypukłą TVS nad $\K$,
a $C\subseteq X$ niepustym, domkniętym i wypukłym
zbiorem o niepustym wnętrzu. Dla każdego $x_0\in\partial C$
istnieje niezerowy, ciągły, rzeczywiście liniowy
$\ell:X\to\R$ taki, że $\ell(x)\le\ell(x_0)$
dla wszystkich $x\in C$.
\end{twierdzenie}
\begin{proof}
Wnętrze $G=\operatorname{int} C$ jest wypukłe.
Ponieważ $x_0\in\partial C$, $x_0\notin G$.
Twierdzenie~\ref{thm:open-point-sep} daje $\ell\ne0$
z $\ell(g)<\ell(x_0)$ dla $g\in G$.
Dla $x\in C$ i ustalonego $g_0\in G$ punkty
$x_t=(1-t)x+tg_0$ należą do $G$ dla $0<t\le1$:
małe otoczenie zera $V$ z $g_0+V\subseteq C$
spełnia $x_t+tV\subseteq C$ na mocy wypukłości.
Ponieważ $x_t\to x$ przy $t\downarrow0$, ciągłość
$\ell$ daje $\ell(x)\le\ell(x_0)$.
\end{proof}

\begin{uwaga}[Czego nie należy wnioskować]
Dla zbioru domkniętego wypukłego o pustym wnętrzu
funkcjonał podpierający w zadanym punkcie nie musi
istnieć bez dodatkowych hipotez. Podobnie z istnienia
hiperpłaszczyzny podpierającej nie wynika, że jest ona
jednoznaczna. W przestrzeni zespolonej pojęcie
hiperpłaszczyzny podpierającej odnosi się do równania
$\operatorname{Re}F(x)=c$, a nie do $F(x)=c$.
\end{uwaga}

\subsection{Zbiory absolutnie wypukłe i polary}

\begin{definicja}[Otoczka absolutnie wypukła]
\label{def:absolutely-convex-hull}
Niech $X$ będzie przestrzenią liniową nad $\K$, a $A\subseteq X$.
Definiujemy
\[
\operatorname{aco}(A)=\left\{\sum_{j=1}^n\lambda_j a_j:
 n\in\N,\ a_j\in A,\ \lambda_j\in\K,
 \sum_{j=1}^n|\lambda_j|\le1\right\}\cup\{0\}.
\]
Jest to najmniejszy zbiór wypukły i zrównoważony zawierający $A$.
Przyjmujemy $\operatorname{aco}(\varnothing)=\{0\}$.
\end{definicja}

\begin{definicja}[Polary w parze $(X,X')$]\label{def:polar}
Niech $X$ będzie TVS nad $\K$ i $X'$ przestrzenią dualną topologiczną.
Dla $A\subseteq X$ i $B\subseteq X'$ definiujemy
\[
 \begin{aligned}
 A^\circ&=\{f\in X':|f(a)|\le1\ \text{dla każdego }a\in A\},\\
 B^\circ&=\{x\in X:|f(x)|\le1\ \text{dla każdego }f\in B\}.
 \end{aligned}
\]
W szczególności $A^{\circ\circ}=(A^\circ)^\circ\subseteq X$.
Powyższa biegunowość jest \emph{absolutna}: w przypadku
zespolonym również używamy modułu, nie samej części rzeczywistej.
\end{definicja}

\begin{stwierdzenie}[Podstawowe własności polary]\label{prop:polar-properties}
Niech $X$ będzie TVS nad $\K$ i $A,A_1,A_2\subseteq X$.
Wówczas:
\begin{enumerate}
\item $A_1\subseteq A_2$ implikuje $A_2^\circ\subseteq A_1^\circ$;
\item $A^\circ=(\operatorname{aco}A)^\circ$ oraz
      $A\subseteq A^{\circ\circ}$;
\item $A^\circ$ jest wypukły, zrównoważony i domknięty
      w topologii punktowej zbieżności na $X'$, tj. $\sigma(X',X)$;
\item $A^{\circ\circ}$ jest wypukły, zrównoważony i domknięty
      w oryginalnej topologii $X$.
\end{enumerate}
\end{stwierdzenie}
\begin{proof}
Pierwsza własność wynika bezpośrednio z kwantyfikatora
w definicji. Dla $f\in A^\circ$ i
$x=\sum_j\lambda_j a_j$ z $\sum_j|\lambda_j|\le1$
mamy $|f(x)|\le\sum_j|\lambda_j||f(a_j)|\le1$,
co dowodzi równości polar. Każdy $a\in A$ spełnia
$|f(a)|\le1$ dla $f\in A^\circ$, zatem $a\in A^{\circ\circ}$.
Zbiory $\{f:|f(a)|\le1\}$ są domknięte, wypukłe i
zrównoważone w topologii punktowej; ich przecięcie
jest zbiorem $A^\circ$. Analogicznie przecięcie
$\{x:|f(x)|\le1\}$ po $f\in A^\circ$ daje własność (4).
\end{proof}

\begin{twierdzenie}[Twierdzenie bipolarne]\label{thm:bipolar}
Niech $X$ będzie lokalnie wypukłą TVS Hausdorffa nad
$\K\in\{\R,\C\}$, a $A\subseteq X$ dowolnym zbiorem.
Wówczas
\[
 A^{\circ\circ}=\overline{\operatorname{aco}(A)},
\]
gdzie domknięcie jest brane w oryginalnej topologii $X$.
\end{twierdzenie}
\begin{proof}
Z poprzedniego stwierdzenia wynika
$\overline{\operatorname{aco}(A)}\subseteq A^{\circ\circ}$.
Oznaczmy $C=\overline{\operatorname{aco}(A)}$.
Jest to zbiór niepusty, domknięty, wypukły i zrównoważony.
Jeśli $x_0\notin C$, Wniosek~\ref{cor:point-closed-convex}
zastosowany do $C$ daje ciągły rzeczywiście liniowy
$\ell$ z $\ell(x_0)>\sup_{c\in C}\ell(c)$
(po ewentualnej zmianie znaku).
Niech $F=\ell$ w przypadku rzeczywistym, a
$F(x)=\ell(x)-i\ell(ix)$ w zespolonym.
Wówczas $F\in X'$, $\operatorname{Re}F=\ell$.
Ze zrównoważenia $C$ wynika
\[
 m:=\sup_{c\in C}|F(c)|=\sup_{c\in C}\operatorname{Re}F(c)
 =\sup_{c\in C}\ell(c)<\ell(x_0)\le|F(x_0)|.
\]
W szczególności $m$ jest skończone i nieujemne.
Jeżeli $m>0$, to $G=F/m$ spełnia $|G(a)|\le1$
dla wszystkich $a\in A$, lecz $|G(x_0)|>1$.
Jeżeli $m=0$, wybieramy dodatnie $t$ z
$t|F(x_0)|>1$ i kładziemy $G=tF$.
W obu przypadkach $G\in A^\circ$, a $x_0\notin A^{\circ\circ}$.
Otrzymujemy więc odwrotną inkluzję.
\end{proof}

\begin{wniosek}[Domknięcie wypukłe nie zmienia się po osłabieniu topologii]
\label{cor:weak-convex-closure}
Niech $X$ będzie lokalnie wypukłą TVS Hausdorffa,
$C\subseteq X$ zbiorem wypukłym, a $\sigma(X,X')$
najsłabszą topologią, w której wszystkie $f\in X'$
są ciągłe. Wtedy domknięcie $C$ w $\sigma(X,X')$
jest równe domknięciu $C$ w oryginalnej topologii $X$.
\end{wniosek}
\begin{proof}
Jeśli $C=\varnothing$, teza jest oczywista. W dalszej części
zakładamy $C\ne\varnothing$. Topologia $\sigma(X,X')$ jest słabsza od wyjściowej,
zatem domknięcie słabe zawiera domknięcie wyjściowe.
Jeżeli $x_0$ nie należy do domknięcia wyjściowego
$C$, Wniosek~\ref{cor:point-closed-convex}
daje ciągły rzeczywiście liniowy $\ell$
i $\alpha\in\R$ takie, że
$\ell(x_0)<\alpha<\inf_{c\in C}\ell(c)$.
Jego reprezentacja przez część rzeczywistą funkcjonału
$F\in X'$ pokazuje, że zbiór
$\{x:\operatorname{Re}F(x)<\alpha\}$ jest słabo otwartym
otoczeniem $x_0$ rozłącznym z $C$.
\end{proof}

\subsection{Przykłady i pułapki}

\begin{przyklad}[Polary kul jednostkowych w $\R^2$]\label{ex:polar-balls}
W parze $X=X'=\R^2$ z parowaniem
$\langle a,x\rangle=a_1x_1+a_2x_2$ polara
$A=\{x:|x_1|+|x_2|\le1\}$ jest zbiorem
$A^\circ=\{a:\max(|a_1|,|a_2|)\le1\}$.
Jest to konsekwencja tożsamości
$\sup_{\|x\|_1\le1}|\langle a,x\rangle|=\|a\|_\infty$.
\begin{center}
\begin{tikzpicture}[scale=1.07,>=Stealth]
\begin{scope}[shift={(-2.65,0)}]
\draw[->,gray] (-1.55,0)--(1.55,0) node[right] {$x_1$};
\draw[->,gray] (0,-1.48)--(0,1.5) node[above] {$x_2$};
\fill[blue!12] (1,0)--(0,1)--(-1,0)--(0,-1)--cycle;
\draw[blue!65!black,very thick] (1,0)--(0,1)--(-1,0)--(0,-1)--cycle;
\node at (0,-1.85) {$A:\ \|x\|_1\le1$};
\end{scope}
\draw[->,thick] (-.72,0)--(.75,0) node[midway,above] {$\circ$};
\begin{scope}[shift={(2.65,0)}]
\draw[->,gray] (-1.55,0)--(1.55,0) node[right] {$a_1$};
\draw[->,gray] (0,-1.48)--(0,1.5) node[above] {$a_2$};
\fill[blue!12] (-1,-1) rectangle (1,1);
\draw[blue!65!black,very thick] (-1,-1) rectangle (1,1);
\node at (0,-1.85) {$A^\circ:\ \|a\|_\infty\le1$};
\end{scope}
\end{tikzpicture}
\end{center}
\end{przyklad}

\begin{przyklad}[Polara w produkcie $\K^\N$]\label{ex:polar-product}
Niech $X=\K^\N$ z topologią produktową, a
$A=\{e_n:n\in\N\}$, gdzie $e_n$ ma jedynkę na
$n$-tej współrzędnej. Z \hyperref[sec:hahn-banach]{rozdziału 6} wiemy, że $X'$
składa się z funkcjonałów $f_a(x)=\sum_n a_nx_n$
o skończonym nośniku. Zatem
$A^\circ=\{f_a:\ a\in c_{00},\ |a_n|\le1\ \text{dla każdego }n\}$,
a więc
\[
 A^{\circ\circ}=
 \left\{x\in\K^\N:\sum_{n=1}^\infty|x_n|\le1\right\}.
\]
Istotnie, supremum $|\sum_n a_nx_n|$ po
skończenie wspartych ciągach $(a_n)$ z $|a_n|\le1$
wynosi $\sup_{m\in\N}\sum_{n=1}^m|x_n|$. Z drugiej strony
$\operatorname{aco}(A)$ jest kulą $\ell^1$ w $c_{00}$,
a jej domknięcie w topologii produktowej jest właśnie
powyższą kulą w $\ell^1$, zgodnie z twierdzeniem bipolarnym.
\end{przyklad}

\begin{przyklad}[Dlaczego lokalna wypukłość ma znaczenie?]
Niech $X=L^p([0,1])$, $0<p<1$, z naturalną topologią
metryczną. Jak wiemy, $X'=\{0\}$. Dlatego dla
$A=\{0\}$ zachodzi $A^\circ=\{0\}$ oraz
$A^{\circ\circ}=X$, mimo że
$\overline{\operatorname{aco}(A)}=\{0\}$.
Twierdzenie bipolarne nie zachodzi więc w dowolnych TVS.
\end{przyklad}

\begin{uwaga}[Co dalej?]
W następnym rozdziale szczegółowo zbadamy pary
przestrzeni dualnych i topologie $\sigma(X,X')$ oraz
$\sigma(X',X)$. Dzięki temu pojęcia polary i
zamknięcia słabego otrzymają naturalną interpretację
jako narzędzia opisu przestrzeni funkcjonałów ciągłych,
a później także dystrybucji.
\end{uwaga}


% Wersja źródłowa: audyt II, jednostka 08
\setcounter{section}{7}
\section{Pary dualne i topologie słabe}\label{sec:topologie-slabe}

\subsection{Dlaczego potrzebujemy kolejnej topologii?}

Dotąd topologia przestrzeni liniowo-topologicznej $X$ była dana z góry,
a przestrzeń dualna topologiczna $X'$ składała się z tych
funkcjonałów liniowych, które są względem niej ciągłe.
Można postąpić odwrotnie: wybrać rodzinę funkcjonałów i skonstruować
\emph{najsłabszą} topologię, w której są one ciągłe.
Z tego pomysłu powstają topologie słabe, kluczowe w badaniu zbieżności
funkcji, operatorów i dystrybucji.

W rozdziale $\K\in\{\R,\C\}$. Wszystkie przestrzenie liniowe
są nad $\K$. Przez $E^*$ oznaczamy przestrzeń dualną algebraiczną,
a przez $E'$ przestrzeń dualną topologiczną, gdy $E$ jest TVS.

\begin{definicja}[Para dualna]\label{def:para-dualna}
Parą dualną nazywamy dwie przestrzenie liniowe $E,F$ wraz
z dwuliniowym parowaniem $\langle\cdot,\cdot\rangle:E\times F\to\K$
takim, że
\begin{enumerate}
\item dla każdego $x\in E\setminus\{0\}$ istnieje $f\in F$
      spełniający $\langle x,f\rangle\ne0$;
\item dla każdego $f\in F\setminus\{0\}$ istnieje $x\in E$
      spełniający $\langle x,f\rangle\ne0$.
\end{enumerate}
Mówimy, że parowanie jest \emph{niezdegenerowane} po obu stronach.
Nie zakładamy istnienia topologii na $E$ ani $F$.
\end{definicja}

\begin{przyklad}[Para kanoniczna]
Jeżeli $E$ jest lokalnie wypukłą TVS Hausdorffa, to
$(E,E')$, z parowaniem $\langle x,f\rangle=f(x)$,
jest parą dualną. Pierwsza niezdegenerowaność wynika
z konsekwencji twierdzenia Hahna--Banacha (\hyperref[sec:hahn-banach]{rozdział 6}),
a druga z faktu, że niezerowy funkcjonał ma niezerową wartość
na pewnym wektorze.
\end{przyklad}

\begin{przyklad}[Para ciągów]
Niech $E=c_0$, a $F=\ell^1$. Parowanie
$\langle x,a\rangle=\sum_{n=1}^{\infty}x_na_n$
jest poprawnie określone, gdyż
$\sum_n|x_na_n|\le\|x\|_\infty\|a\|_1$.
Jest niezdegenerowane: wektory współrzędnych $e_n$ należą
zarówno do $c_0$, jak i do $\ell^1$, i wykrywają odpowiednie
współrzędne.
\end{przyklad}

\subsection{Topologie słabe jako topologie początkowe}

\begin{definicja}[Topologia słaba pary dualnej]\label{def:slaba-para}
Niech $(E,F,\langle\cdot,\cdot\rangle)$ będzie parą dualną.
Topologią $\sigma(E,F)$ nazywamy najsłabszą topologię na $E$,
w której każdy funkcjonał $x\mapsto\langle x,f\rangle$,
$f\in F$, jest ciągły. Analogicznie definiujemy
$\sigma(F,E)$, wykorzystując funkcjonały
$f\mapsto\langle x,f\rangle$ dla $x\in E$.
\end{definicja}

\begin{twierdzenie}[Seminormy i baza otoczeń]\label{thm:slaba-baza}
Niech $(E,F)$ będzie parą dualną. Wtedy $\sigma(E,F)$
jest lokalnie wypukłą topologią Hausdorffa, generowaną przez
seminormy $p_f(x)=|\langle x,f\rangle|$, $f\in F$.
Bazę otoczeń zera stanowią zbiory
\[
 U(f_1,\ldots,f_m;\varepsilon)
 =\{x\in E:|\langle x,f_j\rangle|<\varepsilon,
 \ j=1,\ldots,m\},
\]
gdzie $m\ge1$, $f_1,\ldots,f_m\in F$ i $\varepsilon>0$.
W szczególności sieć $(x_i)_{i\in I}$ zbiega do $x\in E$
w $\sigma(E,F)$ wtedy i tylko wtedy, gdy
$\langle x_i,f\rangle\to\langle x,f\rangle$ dla każdego $f\in F$.
Analogiczne fakty zachodzą dla $\sigma(F,E)$.
\end{twierdzenie}
\begin{proof}
Z definicji topologii początkowej względem rodziny
odwzorowań do $\K$ otrzymujemy zbiory opisane skończenie
wieloma nierównościami. Każda funkcja $p_f$ jest seminormą;
przecięcia jej podpoziomic są absolutnie wypukłe, więc
otrzymana topologia jest lokalnie wypukła. Jeżeli $x\ne0$,
niezdegenerowaność parowania daje $f$ z $p_f(x)>0$.
Można zatem oddzielić $x$ od zera otoczeniem zera,
co dowodzi własności Hausdorffa. Kryterium zbieżności
jest własnością definiującą topologię początkową.
\end{proof}

\begin{stwierdzenie}[Słaba ograniczoność]\label{prop:slaba-ograniczonosc}
Niech $(E,F)$ będzie parą dualną, a $B\subseteq E$.
Zbiór $B$ jest ograniczony w TVS $(E,\sigma(E,F))$
wtedy i tylko wtedy, gdy dla każdego $f\in F$
zachodzi $\sup_{x\in B}|\langle x,f\rangle|<\infty$.
\end{stwierdzenie}
\begin{proof}
Jeżeli $B$ jest ograniczony, jego obraz przez każdy ciągły
funkcjonał jest ograniczonym zbiorem w $\K$.
Odwrotnie, niech $U=U(f_1,\ldots,f_m;\varepsilon)$
będzie podstawowym otoczeniem zera. Z założeń wynika, że
$M=\max_{1\le j\le m}\sup_{x\in B}|\langle x,f_j\rangle|$
jest skończone. Dla każdego $t>\max\{1,M/\varepsilon\}$
mamy $B\subseteq tU$, co jest definicją ograniczoności.
\end{proof}

\begin{twierdzenie}[Przestrzeń dualna topologii słabej]\label{thm:slaby-dual}
Niech $(E,F)$ będzie parą dualną. Odwzorowanie
$F\to(E,\sigma(E,F))'$, przypisujące $f$ funkcjonał
$x\mapsto\langle x,f\rangle$, jest izomorfizmem
algebraicznym. W szczególności topologia słaba nie wytwarza
nowych ciągłych funkcjonałów liniowych poza elementami $F$.
\end{twierdzenie}
\begin{proof}
Niezdegenerowaność zapewnia różnowartościowość, a
ciągłość funkcjonałów $f\in F$ wynika z definicji.
Niech teraz $g:E\to\K$ będzie liniowy i ciągły w
$\sigma(E,F)$. Istnieją $f_1,\ldots,f_m\in F$
i $\varepsilon>0$ takie, że
$|g(x)|<1$ dla $x\in U(f_1,\ldots,f_m;\varepsilon)$.
Połóżmy $T(x)=(\langle x,f_1\rangle,\ldots,
\langle x,f_m\rangle)\in\K^m$. Jeżeli $x\in\ker T$,
to $tx\in U$ dla każdego $t\in\K$; nierówność
$|t g(x)|<1$ dla wszystkich $t$ oznacza $g(x)=0$.
Zatem $g$ faktoryzuje się przez liniową mapę
$\widetilde g:T(E)\to\K$. Przedłużamy ją algebraicznie
z podprzestrzeni $T(E)$ na $\K^m$. Każdy funkcjonał
liniowy na $\K^m$ jest kombinacją współrzędnych,
więc $g(x)=\sum_{j=1}^m a_j\langle x,f_j\rangle
=\langle x,\sum_{j=1}^m a_j f_j\rangle$.
To kończy dowód surjektywności.
\end{proof}

\begin{wniosek}[Topologia wyjściowa i topologia słaba]\label{cor:slaba-wyjsciowa}
Jeżeli $E$ jest lokalnie wypukłą TVS Hausdorffa z topologią
$\tau$ i przestrzenią dualną topologiczną $E'$, to
$\sigma(E,E')\subseteq\tau$ (inkluzja topologii),
a przestrzenie dualne topologiczne $(E,\sigma(E,E'))'$ i
$(E,\tau)'$ są algebraicznie identyczne.
\end{wniosek}
\begin{proof}
Wszystkie elementy $E'$ są ciągłe względem $\tau$,
zatem minimalność topologii początkowej daje inkluzję.
Druga teza wynika z Twierdzenia~\ref{thm:slaby-dual}.
\end{proof}

\subsection{Słabe domknięcia i zjawiska nieskończenie wymiarowe}

\begin{twierdzenie}[Domknięcie podprzestrzeni przez anihilator]
\label{thm:anihilator-domkniecie}
Niech $E$ będzie lokalnie wypukłą TVS Hausdorffa,
$M\subseteq E$ podprzestrzenią liniową, a
$M^\perp=\{f\in E': f(m)=0\ \text{dla wszystkich }m\in M\}$.
Wtedy
\[
\overline M^{\,\tau}=\overline M^{\,\sigma(E,E')}
=\{x\in E:f(x)=0\ \text{dla wszystkich }f\in M^\perp\}.
\]
\end{twierdzenie}
\begin{proof}
Każdy funkcjonał z $M^\perp$ znika na obu domknięciach,
co daje jedną inkluzję. Jeżeli $x_0\notin\overline M^{\,\tau}$,
to w \hyperref[sec:hahn-banach]{rozdziale 6} oddzieliliśmy punkt od domkniętej
podprzestrzeni za pomocą $f\in E'$ spełniającego
$f|_M=0$ i $f(x_0)=1$. Tak więc $x_0$ nie leży w
zbiorze po prawej stronie. Ponieważ domknięcie słabe
zawiera domknięcie w $\tau$, otrzymujemy wszystkie równości.
\end{proof}

\begin{stwierdzenie}[Otoczenia słabe w nieskończonym wymiarze]
\label{prop:slabe-otoczenia}
Niech $(E,F)$ będzie parą dualną, przy czym
$\dim E=\infty$. Wtedy każde otoczenie zera w
$\sigma(E,F)$ zawiera niezerową podprzestrzeń liniową,
a więc nie jest ograniczone w tej topologii.
W konsekwencji $(E,\sigma(E,F))$ nie jest normowalna.
\end{stwierdzenie}
\begin{proof}
Dane otoczenie zawiera $U(f_1,\ldots,f_m;\varepsilon)$,
a ono zawiera $M=\bigcap_{j=1}^m\ker f_j$.
Podprzestrzeń $M$ ma kowymiar nie większy niż $m$,
więc jest niezerowa. Jeżeli $v\in M\setminus\{0\}$,
to $\K v\subseteq U$, lecz zbiór $\K v$ nie jest
słabo ograniczony: wybierając $f\in F$ z
$\langle v,f\rangle\ne0$, otrzymujemy nieograniczone
wartości $|\langle tv,f\rangle|$.
Przestrzeń normowalna ma ograniczone otoczenie zera,
co przeczy powyższemu faktowi.
\end{proof}

\begin{przyklad}[Słaba zbieżność bez zbieżności w normie]
Niech $E=\ell^2$ nad $\R$ albo $\C$ z topologią normową,
a $e_n$ oznacza $n$-ty wektor jednostkowy. Wtedy
$e_n\to0$ słabo, ponieważ każdy funkcjonał ciągły na
$\ell^2$ ma postać $x\mapsto\sum_k x_k a_k$
z pewnym $a\in\ell^2$ (w przypadku zespolonym
odpowiednio wybieramy współrzędne liniowe), a $a_n\to0$.
Jednocześnie $\|e_n\|_2=1$ dla każdego $n$.
\end{przyklad}

\begin{twierdzenie}[Słabe domknięcie sfery jednostkowej]
\label{thm:sfera-slaba}
Niech $E$ będzie nieskończenie wymiarową przestrzenią
unormowaną nad $\K$, $S=\{x\in E:\|x\|=1\}$,
$B=\{x\in E:\|x\|\le1\}$. Wtedy
$\overline S^{\,\sigma(E,E')}=B$.
\end{twierdzenie}
\begin{proof}
Kula $B$ jest słabo domknięta. Istotnie, jeżeli
$\|x_0\|>1$, twierdzenie Hahna--Banacha daje
$f\in E'$ z $\|f\|=1$ i $f(x_0)=\|x_0\|$,
a zbiór $\{x:\operatorname{Re}f(x)>1\}$ jest słabo
otwarty, zawiera $x_0$ i nie przecina $B$.
Stąd $\overline S^{\,\sigma(E,E')}\subseteq B$.

Niech $\|x\|<1$ i $U$ będzie podstawowym słabym
otoczeniem $x$, określonym przez skończenie wiele
funkcjonałów $f_1,\ldots,f_m$. Z nieskończonego wymiaru
wynika istnienie $v\ne0$ we wspólnym jądrze tych
funkcjonałów. Funkcja $t\mapsto\|x+tv\|$, $t\ge0$,
jest ciągła, w zerze mniejsza od $1$, a dla $t\to\infty$
dąży do nieskończoności. Istnieje więc $t_0>0$
z $\|x+t_0v\|=1$. Punkt $x+t_0v$ należy do
$S\cap U$, gdyż funkcjonały definiujące $U$ mają
na nim te same wartości co na $x$. To dowodzi
$B\subseteq\overline S^{\,\sigma(E,E')}$.
\end{proof}

\begin{uwaga}[Co stało się z intuicją geometryczną?]
W przestrzeni skończenie wymiarowej sfera jest domknięta.
W nieskończenie wymiarowej przestrzeni unormowanej topologia
słaba nie rozróżnia, za pomocą skończenie wielu funkcjonałów,
punktów sfery od punktów wewnątrz kuli. To jednak nie znaczy,
że kula jest słabo zwarta: zwartość omówimy osobno.
\end{uwaga}

\subsection{Topologia słaba-* na przestrzeni dualnej}

\begin{definicja}[Topologia słaba-*]\label{def:slaba-gwiazdka}
Niech $E$ będzie lokalnie wypukłą TVS Hausdorffa.
Na przestrzeni dualnej topologicznej $E'$ topologią
\emph{słabą-*} nazywamy $\sigma(E',E)$, tj. topologię
punktowej zbieżności funkcjonałów na $E$. Dla sieci
$(f_i)_{i\in I}\subseteq E'$ zapis
$f_i\overset{*}{\rightharpoonup}f$ oznacza, że
$f_i(x)\to f(x)$ dla każdego $x\in E$.
\end{definicja}

\begin{stwierdzenie}[Słaba i słaba-* topologia na przestrzeni dualnej]
\label{prop:slaba-weakstar}
Niech $E$ będzie przestrzenią Banacha. Na przestrzeni
Banacha $E'$ rozważmy topologię słabą
$\sigma(E',E'')$ i słabą-* $\sigma(E',E)$,
gdzie $E$ utożsamiamy z obrazem kanonicznego zanurzenia
$J:E\to E''$, $J(x)(f)=f(x)$.
Wtedy $\sigma(E',E)\subseteq\sigma(E',E'')$,
a równość zachodzi wtedy i tylko wtedy, gdy $E$ jest
refleksywna (tj. $J$ jest surjekcją).
\end{stwierdzenie}
\begin{proof}
Pierwsza inkluzja wynika z $J(E)\subseteq E''$.
Jeżeli $J(E)=E''$, rodziny definiujących funkcjonałów
są jednakowe, więc topologie są równe.
Odwrotnie, jeżeli topologie są równe, ich przestrzenie
dualne topologiczne są równe. Z Twierdzenia
\ref{thm:slaby-dual}, zastosowanego do par
$(E',J(E))$ i $(E',E'')$, wynika odpowiednio,
że te przestrzenie dualne są algebraicznie równe
$J(E)$ i $E''$. Zatem $J(E)=E''$.
\end{proof}

\begin{przyklad}[Zbieżność słaba-* nie musi być słaba]
Niech $E=c_0$, $E'=\ell^1$ z naturalnym parowaniem.
Ciąg wektorów jednostkowych $e_n\in\ell^1$ spełnia
$e_n\to0$ w $\sigma(\ell^1,c_0)$, ponieważ dla każdego
$x\in c_0$ mamy $\langle x,e_n\rangle=x_n\to0$.
Nie jest to jednak zbieżność słaba w $\ell^1$:
funkcjonał z $\ell^\infty=(\ell^1)'$ odpowiadający
ciągowi $(1,1,\ldots)$ ma na każdym $e_n$ wartość $1$.
\end{przyklad}

\begin{przyklad}[Topologia produktowa jako topologia słaba]
Niech $E=\K^{\N}$ z topologią produktową i $F=c_{00}$
z parowaniem $\langle x,a\rangle=\sum_n x_na_n$
(suma jest skończona). Ponieważ współrzędne
$x\mapsto x_n$ odpowiadają wektorom $e_n\in F$,
topologia produktowa zawiera $\sigma(E,F)$ i odwrotnie:
każdy funkcjonał z $F$ jest skończoną kombinacją
odwzorowań współrzędnych. Zatem topologia produktowa
jest dokładnie $\sigma(E,F)$.
\end{przyklad}

\subsection{Zupełność topologii słabej i przestrzeń dualna algebraiczna}

\begin{twierdzenie}[Uzupełnienie przestrzeni ze słabą topologią]
\label{thm:slabe-uzupelnienie}
Niech $(E,F)$ będzie parą dualną. Niech $F^*$ oznacza
przestrzeń dualną algebraiczną przestrzeni liniowej $F$,
wyposażoną w topologię punktowej zbieżności na $F$.
Wówczas odwzorowanie $J:E\to F^*$, określone wzorem
$J(x)(f)=\langle x,f\rangle$, jest liniowym zanurzeniem
topologicznym przestrzeni $(E,\sigma(E,F))$ w $F^*$.
Obraz $J(E)$ jest gęsty, a $F^*$ jest zupełną TVS
Hausdorffa. Tym samym $F^*$ jest uzupełnieniem
$(E,\sigma(E,F))$ w sensie \hyperref[sec:zupelnosc]{rozdziału 3}.
\end{twierdzenie}
\begin{proof}
Traktujemy $F^*$ jako podzbiór produktu $\K^F$.
Jest to zbiór tych funkcji $\Phi:F\to\K$, które spełniają
$\Phi(f+g)=\Phi(f)+\Phi(g)$ i
$\Phi(af)=a\Phi(f)$ dla $f,g\in F$, $a\in\K$.
Każdy taki warunek definiuje zbiór domknięty w topologii
produktowej, zatem $F^*$ jest domkniętą podprzestrzenią
$\K^F$. Produkt zupełnych przestrzeni liniowo-topologicznych
jest zupełny, a podprzestrzeń domknięta przestrzeni
zupełnej jest zupełna. Topologia jest Hausdorffa.

Niezdegenerowaność parowania daje różnowartościowość $J$,
a definicje topologii początkowych oznaczają, że jest to
zanurzenie topologiczne. Aby wykazać gęstość, weźmy
$\Phi\in F^*$ i podstawowe otoczenie punktu $\Phi$,
wyznaczone przez $f_1,\ldots,f_m\in F$ oraz $\varepsilon>0$.
Niech $G=\operatorname{lin}\{f_1,\ldots,f_m\}\subseteq F$.
Odwzorowanie $T:E\to G^*$, $T(x)(f)=\langle x,f\rangle$,
jest surjekcją. Istotnie, gdyby jego obraz był właściwą
podprzestrzenią skończenie wymiarowej przestrzeni $G^*$,
istniałby niezerowy element $g\in G$ z
$\langle x,g\rangle=0$ dla każdego $x\in E$,
co przeczy niezdegenerowaności parowania.
Wybieramy zatem $x\in E$ takie, że
$T(x)=\Phi|_G$. Wtedy $J(x)(f_j)=\Phi(f_j)$
dla wszystkich $j$, a więc $J(x)$ należy do rozważanego
otoczenia. To dowodzi gęstości.
\end{proof}

\begin{wniosek}[Zaskakująca niezupełność topologii słabej]
\label{cor:slabe-niezupełne}
Jeśli $E$ jest nieskończenie wymiarową przestrzenią
unormowaną, to przestrzeń $(E,\sigma(E,E'))$ nie jest
zupełna, nawet gdy $E$ jest przestrzenią Hilberta.
\end{wniosek}
\begin{proof}
Przestrzeń $E'$ jest nieskończenie wymiarową przestrzenią
Banacha. Istnieje na niej nieciągły funkcjonał liniowy
$\Phi\in(E')^*$ (przykład z \hyperref[sec:hahn-banach]{rozdziału 6}).
Każdy $J(x)$, $x\in E$, jest jednak ciągły względem normy
operatorowej na $E'$, ponieważ
$|J(x)(f)|=|f(x)|\le\|x\|\|f\|$.
Stąd $\Phi\notin J(E)$. Twierdzenie
\ref{thm:slabe-uzupelnienie} mówi, że $J(E)$ jest
gęstą, właściwą podprzestrzenią zupełnej TVS Hausdorffa
$F^*=(E')^*$. Zatem $J(E)$ nie jest zupełne.
\end{proof}

\begin{uwaga}
Powyższy wniosek nie przeczy zupełności $E$ w normie:
zupełność zależy od topologii i jej struktury jednostajnej.
Nie należy też mylić przestrzeni dualnej algebraicznej
$(E')^*$ z bidualem topologicznym $E''=(E')'$.
\end{uwaga}

\subsection{Topologia silna i zapowiedź topologii Mackeya}

\begin{definicja}[Topologia silna na przestrzeni dualnej]
\label{def:silna-dualna}
Niech $E$ będzie lokalnie wypukłą TVS Hausdorffa.
Dla każdego niepustego ograniczonego zbioru $B\subseteq E$ i $f\in E'$
kładziemy $p_B(f)=\sup_{x\in B}|f(x)|$.
Każdy $p_B$ jest skończony, ponieważ ciągły funkcjonał
przekształca zbiory ograniczone na ograniczone zbiory skalarów.
Topologię na $E'$ generowaną przez rodzinę takich seminorm
nazywamy topologią silną i oznaczamy $\beta(E',E)$.
\end{definicja}
\begin{stwierdzenie}[Przypadek przestrzeni unormowanych]
\label{prop:silna-unormowana}
Jeśli $E$ jest przestrzenią unormowaną, to topologia
$\beta(E',E)$ jest topologią normy operatorowej
$\|f\|=\sup_{\|x\|\le1}|f(x)|$ na $E'$.
\end{stwierdzenie}
\begin{proof}
Dla ograniczonego $B$ mamy
$\sup_{x\in B}\|x\|=M<\infty$, więc
$p_B(f)\le M\|f\|$. Z drugiej strony kula jednostkowa
$B_E$ jest ograniczona, a $p_{B_E}(f)=\|f\|$.
Obie rodziny seminorm definiują zatem tę samą topologię.
\end{proof}

\begin{definicja}[Topologia Mackeya]\label{def:mackey}
Niech $(E,F)$ będzie parą dualną. Dla każdego absolutnie
wypukłego, niepustego i zwartego w $\sigma(F,E)$ zbioru $K\subseteq F$
określamy $p_K(x)=\sup_{f\in K}|\langle x,f\rangle|$.
Topologię lokalnie wypukłą generowaną przez wszystkie takie
seminormy nazywamy \emph{topologią Mackeya} $\tau(E,F)$.
\end{definicja}

\begin{stwierdzenie}[Poprawność definicji Mackeya]\label{prop:mackey-basic}
Dla dowolnej pary dualnej $(E,F)$ funkcje $p_K$ są
skończonymi seminormami, a topologia $\tau(E,F)$
jest lokalnie wypukła, Hausdorffa i silniejsza od
$\sigma(E,F)$.
\end{stwierdzenie}
\begin{proof}
Dla ustalonego $x\in E$ funkcja
$f\mapsto|\langle x,f\rangle|$ jest ciągła względem
$\sigma(F,E)$, a zatem ograniczona na każdym zwartym $K$.
Nierówność trójkąta i jednorodność dają własności seminormy.
Dla każdego $f\in F$ zbiór
$K_f=\{af:|a|\le1\}$ jest absolutnie wypukły i zwarty
jako obraz zwartego domkniętego koła w $\K$.
Mamy $p_{K_f}(x)=|\langle x,f\rangle|$.
Topologia Mackeya zawiera więc topologię słabą,
a ponieważ ta ostatnia jest Hausdorffa, taką własność
ma również $\tau(E,F)$.
\end{proof}

\begin{uwaga}[Twierdzenie Mackeya--Arensa: zapowiedź]
Twierdzenie Mackeya--Arensa orzeka, że $\tau(E,F)$ jest
najsilniejszą topologią lokalnie wypukłą na $E$, dla której
przestrzeń dualna topologiczna jest algebraicznie równa $F$
(przy podanym parowaniu). Wymaga to jednak twierdzeń
o zwartości i zbiorach polarnych, które rozwiniemy w następnym
rozdziale. Nie traktujemy tej charakterystyki jako
udowodnionej w bieżącym miejscu.
\end{uwaga}

\begin{uwaga}[W stronę dystrybucji]
Dla przestrzeni funkcji testowych $\mathcal D(\Omega)$
topologia słaba-* $\sigma(\mathcal D'(\Omega),\mathcal D(\Omega))$
opisuje zbieżność dystrybucji na każdej ustalonej funkcji
testowej. Topologia silna przestrzeni dualnej opisuje
natomiast zbieżność jednostajną na wszystkich zbiorach
ograniczonych funkcji testowych. Ich porównanie będzie
istotne w dalszych rozdziałach.
\end{uwaga}


% Wersja źródłowa: audyt II, jednostka 09
\setcounter{section}{8}
\section{Banach--Alaoglu, Krein--Milman i topologia Mackeya}\label{sec:alaoglu-mackey}

\subsection{Zwartość słaba-* i równociągłość}

W przestrzeni nieskończenie wymiarowej kula jednostkowa nie jest zwarta
w topologii normowej. Zmiana topologii na słabą-* pozwala odzyskać zwartość,
ale zasadniczą rolę odgrywają wtedy \emph{zbiory polarne otoczeń zera}.
Najpierw dowiedziemy twierdzenia w klasie dowolnych przestrzeni
liniowo-topologicznych, bez zakładania lokalnej wypukłości i zupełności.

W całym rozdziale $\K\in\{\R,\C\}$. Dla TVS $E$ symbol $E'$ oznacza
przestrzeń dualną topologiczną, a $E^*$ przestrzeń dualną algebraiczną.
Zbiorem polarnym $A\subseteq E$ nazywamy
$A^\circ=\{f\in E':|f(x)|\le1\ \text{dla każdego }x\in A\}$.
Na $E'$ rozpatrujemy topologię $\sigma(E',E)$ punktowej zbieżności na $E$.

\begin{definicja}[Równociągła rodzina funkcjonałów]\label{def:equicont-family}
Niech $E$ będzie TVS i $H\subseteq E'$. Rodzinę $H$ nazywamy
\emph{równociągłą}, jeżeli dla każdego $\varepsilon>0$ istnieje otoczenie
$U$ zera w $E$ takie, że $|f(u)|<\varepsilon$ dla wszystkich
$f\in H$ i $u\in U$.
\end{definicja}

\begin{stwierdzenie}[Równociągłość i zbiory polarne]\label{prop:equicont-polar}
Niech $E$ będzie TVS, a $H\subseteq E'$. Rodzina $H$ jest równociągła
wtedy i tylko wtedy, gdy istnieje otoczenie $U$ zera w $E$ takie,
że $H\subseteq U^\circ$. Każda rodzina równociągła jest punktowo
ograniczona, tzn. $\sup_{f\in H}|f(x)|<\infty$ dla każdego $x\in E$.
\end{stwierdzenie}
\begin{proof}
Jeżeli $H$ jest równociągła, bierzemy w definicji $\varepsilon=1$;
odpowiednie $U$ spełnia $H\subseteq U^\circ$. Odwrotnie,
jeżeli $H\subseteq U^\circ$, to dla ustalonego $\varepsilon>0$
otoczenie $(\varepsilon/2)U$ daje
$|f(v)|\le\varepsilon/2<\varepsilon$ dla $v\in(\varepsilon/2)U$
i $f\in H$. Dla $x\in E$ pochłanianie przez $U$ zapewnia
$x\in tU$ dla pewnego $t>0$, zatem $|f(x)|\le t$ dla $f\in H$.
\end{proof}

\begin{uwaga}[Wykorzystywane twierdzenie topologiczne]
W dowodzie twierdzenia Banacha--Alaoglu zastosujemy \emph{twierdzenie
Tychonowa}: iloczyn dowolnej rodziny zwartych przestrzeni topologicznych,
wyposażony w topologię produktową, jest zwarty. Jest to klasyczne
wynikanie z aksjomatu wyboru; nie będziemy tu ponownie dowodzić
tego ogólnotopologicznego rezultatu. Stosujemy go jedynie do iloczynów
zwartych przedziałów lub domkniętych kół w $\K$.
\end{uwaga}

\begin{twierdzenie}[Banach--Alaoglu--Bourbaki]\label{thm:banach-alaoglu-general}
Niech $E$ będzie przestrzenią liniowo-topologiczną nad $\K$, a
$U\subseteq E$ otoczeniem zera. Wówczas zbiór polarny $U^\circ$
jest zwarty w topologii $\sigma(E',E)$.
Nie zakładamy, że $E$ jest Hausdorffa, lokalnie wypukła, metryzowalna
lub zupełna.
\end{twierdzenie}
\begin{proof}
Umieszczamy $E^*$ w $\K^E$, utożsamiając funkcjonał z rodziną
jego wartości $(f(x))_{x\in E}$. Topologia punktowej zbieżności
jest topologią indukowaną przez topologię produktową.
Ponieważ $U$ jest pochłaniające, dla każdego $x\in E$ można wybrać
$t_x>0$ takie, że $x\in t_xU$. Każdy $f\in U^\circ$ spełnia wtedy
$|f(x)|\le t_x$. Niech $D_x=\{z\in\K:|z|\le t_x\}$ i
$P=\prod_{x\in E}D_x$. Zbiory $D_x$ są zwarte, więc $P$ jest zwarty
na mocy twierdzenia Tychonowa.

Zbiór $U^\circ$ jest domknięty w $P$. Rzeczywiście, warunki
$f(x+y)=f(x)+f(y)$, $f(ax)=af(x)$ oraz $|f(u)|\le1$ dla $u\in U$
opisują przecięcia zbiorów domkniętych w topologii produktowej.
Każda funkcja $f:E\to\K$ spełniająca te warunki jest liniowa,
a przy tym ciągła: dla $\varepsilon>0$ na otoczeniu
$(\varepsilon/2)U$ zachodzi $|f(v)|\le\varepsilon/2$.
Zatem otrzymany zbiór jest dokładnie $U^\circ\subseteq E'$.
Jako domknięty podzbiór zwartego $P$ jest on zwarty. Topologia
indukowana z $P$ na $U^\circ$ to $\sigma(E',E)$.
\end{proof}

\begin{wniosek}[Kula przestrzeni dualnej]\label{cor:alaoglu-norm}
Niech $E$ będzie przestrzenią unormowaną, niekoniecznie zupełną.
Dla każdego $R\ge0$ kula
$B_R'=\{f\in E':\|f\|\le R\}$ jest zwarta w $\sigma(E',E)$.
W szczególności normowa kula jednostkowa przestrzeni dualnej
każdej przestrzeni Banacha jest słabo-* zwarta.
\end{wniosek}
\begin{proof}
Dla $R>0$ niech $U=\{x\in E:\|x\|<1/R\}$.
Wtedy $U^\circ=B_R'$: warunek $|f(u)|\le1$ dla $\|u\|<1/R$
jest równoważny $\|f\|\le R$. Zastosowanie
Twierdzenia~\ref{thm:banach-alaoglu-general} kończy dowód.
Dla $R=0$ zbiór $B_0'=\{0\}$ jest zwarty.
\end{proof}

\begin{wniosek}[Domknięte podzbiory kuli]\label{cor:alaoglu-closed}
Niech $E$ będzie przestrzenią unormowaną, $R\ge0$, a
$A\subseteq B_R'$ będzie domknięty w topologii
$\sigma(E',E)$ indukowanej na $B_R'$. Wtedy $A$ jest
słabo-* zwarty.
\end{wniosek}
\begin{proof}
Domknięty podzbiór zwartej przestrzeni jest zwarty.
\end{proof}

\begin{uwaga}[Czego twierdzenie nie mówi?]
Twierdzenie Banacha--Alaoglu nie orzeka, że cała przestrzeń $E'$
jest słabo-* zwarta ani że kula jest zwarta w normie.
Zwartość nie daje też automatycznie możliwości zastąpienia
sieci przez ciągi: do tego potrzebna jest dodatkowa metryzowalność.
\end{uwaga}

\subsection{Metryzowalność i kryterium ośrodkowości}

\begin{twierdzenie}[Metryzowalność na rodzinie równociągłej]
\label{thm:weakstar-metrizable-equicont}
Niech $E$ będzie ośrodkową TVS Hausdorffa, a $H\subseteq E'$
rodziną równociągłą. Wówczas topologia indukowana na $H$
przez $\sigma(E',E)$ jest metryzowalna.
\end{twierdzenie}
\begin{proof}
Niech $D=\{x_n:n\in\N\}$ będzie zbiorem gęstym w $E$.
Określamy na $H$ funkcję
\[
 d(f,g)=\sum_{n=1}^\infty 2^{-n}\min\{1,|f(x_n)-g(x_n)|\}.
\]
Jest to metryka: nierówność trójkąta wynika z
$\min\{1,a+b\}\le\min\{1,a\}+\min\{1,b\}$,
a jeśli $d(f,g)=0$, to $f$ i $g$ pokrywają się na gęstym
$D$, zatem przez ciągłość na całym $E$.
Topologia metryki jest topologią punktowej zbieżności na $D$,
więc jest nie silniejsza od $\sigma(E',E)|_H$.

Pokażemy odwrotną inkluzję. Ustalmy $x\in E$, $f_0\in H$
i $\varepsilon>0$. Ponieważ $H$ jest równociągła,
istnieje zrównoważone otoczenie $V$ zera takie, że
$|h(v)|<\varepsilon/3$ dla wszystkich $h\in H$, $v\in V$.
Wybieramy $x_n\in D\cap(x+V)$. Dla $f\in H$ mamy
\[
 |(f-f_0)(x)|\le |f(x-x_n)|+|(f-f_0)(x_n)|
                      +|f_0(x_n-x)|
 <2\varepsilon/3+|(f-f_0)(x_n)|.
\]
Warunek $|(f-f_0)(x_n)|<\varepsilon/3$ jest otwarty
w metryce $d$. Zatem wartościowania w każdym $x\in E$
są ciągłe na $(H,d)$, co daje odwrotną inkluzję topologii.
\end{proof}

\begin{twierdzenie}[Kryterium dla kuli przestrzeni dualnej]
\label{thm:banach-ball-metrizable-iff-separable}
Niech $E$ będzie przestrzenią unormowaną nad $\K$,
a $B'=\{f\in E':\|f\|\le1\}$. Wówczas następujące
warunki są równoważne:
\begin{enumerate}
\item $E$ jest ośrodkowa w topologii normowej;
\item $B'$ z topologią $\sigma(E',E)$ jest metryzowalna.
\end{enumerate}
Jeśli warunki zachodzą, to $B'$ jest zwarta i sekwencyjnie zwarta.
\end{twierdzenie}
\begin{proof}
Kula $B'$ jest równociągła, bo $|f(x)|\le\|x\|$ dla $f\in B'$.
Jeżeli $E$ jest ośrodkowa, stosujemy
Twierdzenie~\ref{thm:weakstar-metrizable-equicont};
zwartość wynika z Twierdzenia~\ref{thm:banach-alaoglu-general}.
Każda zwarta przestrzeń metryczna jest sekwencyjnie zwarta.

Odwrotnie, załóżmy, że $B'$ jest metryzowalna. Istnieje
przeliczalna baza otoczeń zera w $B'$, a więc można wybrać
ciąg podstawowych otoczeń słabych-* zera
$V_n=\{f\in B':|f(x)|<\varepsilon_n\ \text{dla }x\in F_n\}$,
gdzie $F_n\subseteq E$ jest skończony, tak aby $\{V_n\}$
była bazą otoczeń zera w $B'$. Niech
$M=\operatorname{lin}(\bigcup_{n=1}^\infty F_n)$.
Jeśli $\overline M^{\|\cdot\|}\ne E$, twierdzenie Hahna--Banacha
daje niezerowy $f\in B'$ zerujący się na $M$.
Wtedy $f\in V_n$ dla każdego $n$. Ponieważ $B'$ jest Hausdorffa,
przecięcie wszystkich otoczeń zera w $B'$ jest równe $\{0\}$;
sprzeczność. Zatem $M$ jest gęsta w $E$, a jako przestrzeń
rozpięta nad $\K$ przez zbiór przeliczalny ma przeliczalny
gęsty podzbiór (kombinacje wymierne, a nad $\C$ -- z
wymiernymi częściami rzeczywistymi i urojonymi współczynników).
Stąd $E$ jest ośrodkowa.
\end{proof}

\begin{przyklad}[Kula przestrzeni dualnej $\ell^1$]
Dla $E=c_0$ mamy $E'=\ell^1$. Ponieważ $c_0$ jest ośrodkowa,
kula jednostkowa w $\ell^1$ jest zwarta i metryzowalna
w $\sigma(\ell^1,c_0)$. Nie jest natomiast zwarta w normie:
wektory $e_n$ należą do kuli, lecz $\|e_n-e_m\|_1=2$
dla $n\ne m$. Przy tym $e_n\to0$ słabo-*.
\end{przyklad}

\subsection{Dlaczego sieci są niezbędne?}

\begin{przyklad}[Zwartość bez metryzowalności]\label{ex:uncountable-product}
Niech $I$ będzie zbiorem nieprzeliczalnym,
$E=\ell^1(I)$, a $E'=\ell^\infty(I)$ z kanonicznym parowaniem
$\langle x,f\rangle=\sum_{i\in I}x_i f_i$.
Kula jednostkowa $B'$ przestrzeni $\ell^\infty(I)$
wyposażona w $\sigma(\ell^\infty(I),\ell^1(I))$
jest homeomorficzna z produktem
$\prod_{i\in I}\{z\in\K:|z|\le1\}$.
Istotnie, zbieżność słaba-* implikuje współrzędną przez
podstawienie $x=e_i$. Odwrotnie, jeśli sieć $f_\alpha\in B'$
zbiega współrzędnie do $f$, to dla $x\in\ell^1(I)$
wybieramy skończony $F\subseteq I$ z
$\sum_{i\notin F}|x_i|<\varepsilon/4$ i kontrolujemy
skończoną sumę $\sum_{i\in F}x_i(f_{\alpha,i}-f_i)$;
reszta jest mniejsza od $2\sum_{i\notin F}|x_i|<\varepsilon/2$.
Produkt jest zwarty, lecz nie jest metryzowalny, ponieważ
$\ell^1(I)$ nie jest ośrodkowa.

Wewnątrz $B'$ rozpatrzmy zbiór
$A=\{\mathbf 1_F:F\subseteq I\ \text{jest skończony}\}$,
gdzie $\mathbf 1_F$ jest funkcją charakterystyczną $F$.
Sieć $(\mathbf 1_F)$, indeksowana skończonymi podzbiorami $I$
i uporządkowana inkluzją, zbiega słabo-* do funkcji
$\mathbf 1_I$. Żaden \emph{ciąg} elementów $A$ nie zbiega
jednak słabo-* do $\mathbf 1_I$: suma jego przeliczalnie wielu
skończonych nośników jest przeliczalna, a więc istnieje
$i\in I$, na którym wszystkie wyrazy ciągu mają wartość zero.
\end{przyklad}

\begin{przyklad}[Słabo-* zwarta rodzina nierównociągła]
\label{ex:compact-not-equicont}
Niech $E=c_{00}$ będzie przestrzenią ciągów o skończonym nośniku,
z normą $\|x\|_\infty=\sup_n|x_n|$. Przestrzeń $E$ nie jest
zupełna. Dla $n\in\N$ określamy $f_n\in E'$ wzorem
$f_n(x)=n x_n$. Ponieważ każdy $x\in c_{00}$ ma skończony
nośnik, $f_n(x)=0$ dla wszystkich dostatecznie dużych $n$.
Wobec tego $f_n\to0$ w $\sigma(E',E)$, a zbiór
$K=\{0\}\cup\{f_n:n\in\N\}$ jest słabo-* zwarty:
każdy otwarty zbiór zawierający zero zawiera prawie wszystkie $f_n$.
Jednak $\|f_n\|=n$, zatem $K$ nie jest równociągły.
To pokazuje, że słaba-* zwartość nie oznacza równociągłości
bez dodatkowych założeń. W rozdziale o twierdzeniu
Banacha--Steinhausa omówimy warunki, pod którymi implikacja
staje się prawdziwa.
\end{przyklad}

\subsection{Punkty ekstremalne i twierdzenie Kreina--Milmana}
\label{subsec:krein-milman}

Twierdzenie Banacha--Alaoglu dostarcza zwartych zbiorów wypukłych
w topologiach słabych-*. Kolejne twierdzenie opisuje ich geometrię:
cały zbiór można odzyskać z jego punktów ekstremalnych. W tej
części, jeśli nie zaznaczono inaczej, $E$ oznacza przestrzeń
lokalnie wypukłą Hausdorffa nad $\K\in\{\R,\C\}$. W przypadku
zespolonym wypukłość rozumiemy w sensie rzeczywistym.

\begin{definicja}[Punkt ekstremalny i ściana]
\label{def:extreme-face9}
Niech $K\subseteq E$ będzie niepustym zbiorem wypukłym.
Punkt $e\in K$ jest \emph{ekstremalny}, jeżeli z
$e=tx+(1-t)y$, $x,y\in K$, $0<t<1$ wynika $x=y=e$.
Zbiór wszystkich takich punktów oznaczamy $\operatorname{Ext}(K)$.
Niepusty podzbiór $F\subseteq K$ jest \emph{ścianą} zbioru $K$,
jeżeli jest wypukły i dla $x,y\in K$ oraz $0<t<1$ z
$tx+(1-t)y\in F$ wynika $x,y\in F$.
Przez $\operatorname{conv}(A)$ oznaczamy otoczkę wypukłą $A$,
a przez $\overline A$ jej domknięcie w $E$.
\end{definicja}

\begin{przyklad}[Wielokąt i kula w przestrzeni Hilberta]
Wierzchołki nieosobliwego wielokąta w $\R^2$ są dokładnie jego
punktami ekstremalnymi. Jeżeli $H$ jest rzeczywistą przestrzenią
Hilberta, to punktami ekstremalnymi domkniętej kuli jednostkowej
$B_H$ są dokładnie punkty sfery $\{x:\|x\|=1\}$.
Istotnie, dla $\|x\|<1$ i dostatecznie małego $h\ne0$
mamy $x=\tfrac12(x+h)+\tfrac12(x-h)$, a oba punkty należą do kuli.
Natomiast jeśli $\|x\|=1$ i $x=(u+v)/2$ przy $u,v\in B_H$,
to tożsamość równoległoboku daje
$\|u-v\|^2=2\|u\|^2+2\|v\|^2-4\|x\|^2\le0$, więc $u=v=x$.
\end{przyklad}

\begin{center}
\begin{tikzpicture}[scale=0.90, every node/.style={font=\small}]
\filldraw[fill=blue!7,draw=blue!55!black,very thick]
(0.1,0.6)--(1.5,0.15)--(3.9,0.65)--(4.35,2.05)--(2.75,2.85)--(0.4,2.3)--cycle;
\foreach \x/\y in {0.1/0.6,1.5/0.15,3.9/0.65,4.35/2.05,2.75/2.85,0.4/2.3}
 {\fill[blue!65!black] (\x,\y) circle (2.7pt);}
\draw[orange!85!black,densely dashed,thick] (-0.1,2.85)--(4.75,2.85);
\fill[orange!80!black] (2.75,2.85) circle (3.5pt);
\node[anchor=west,align=left] at (4.83,2.9) {$\ell=\max_K\ell$};
\node at (2.2,1.5) {$K$};
\node[anchor=north] at (2.2,-0.05) {Wierzchołki są punktami ekstremalnymi.};
\end{tikzpicture}
\end{center}
\noindent
\emph{Ilustracja.} Linia przerywana jest poziomicą funkcjonału
liniowego podpierającą $K$; odpowiednia ściana składa się tutaj
z jednego punktu ekstremalnego.

\begin{lemat}[Ściana poziomicy maksymalnej]
\label{lem:max-face9}
Niech $K\subseteq E$ będzie niepusty, zwarty i wypukły,
a $\ell:E\to\R$ niech będzie ciągłym rzeczywistym
funkcjonałem liniowym. Wówczas zbiór
$F_\ell=\{x\in K:\ell(x)=\max_{y\in K}\ell(y)\}$
jest niepustą zwartą ścianą $K$.
Jeżeli $F$ jest ścianą $K$, a $G$ ścianą $F$,
to $G$ jest również ścianą $K$.
\end{lemat}
\begin{proof}
Ciągłość $\ell$ i zwartość $K$ zapewniają istnienie maksimum;
$F_\ell$ jest domkniętym, a więc zwartym podzbiorem $K$.
Jest wypukły z liniowości $\ell$. Jeśli
$tx+(1-t)y\in F_\ell$, gdzie $x,y\in K$ i $0<t<1$, to
$\max_K\ell=t\ell(x)+(1-t)\ell(y)$, podczas gdy
$\ell(x),\ell(y)\le\max_K\ell$. Stąd obie wartości
są maksymalne i $x,y\in F_\ell$. Druga część wynika
z dwukrotnego zastosowania definicji ściany.
\end{proof}

\begin{lemat}[Istnienie punktu ekstremalnego]
\label{lem:exist-extreme9}
Każdy niepusty, zwarty i wypukły podzbiór przestrzeni lokalnie
wypukłej Hausdorffa ma co najmniej jeden punkt ekstremalny.
Co więcej, każda jego niepusta, domknięta ściana zawiera
punkt ekstremalny całego zbioru.
\end{lemat}
\begin{proof}
Niech $K\subseteq E$ spełnia założenia. Rozważmy rodzinę
$\mathcal F$ wszystkich niepustych \emph{domkniętych} ścian $K$,
uporządkowaną przez inkluzję odwrotną.
Każdy łańcuch w $\mathcal F$ ma górne ograniczenie:
jego elementy są zagnieżdżonymi niepustymi zbiorami
domkniętymi w zwartym $K$, więc mają niepuste przecięcie.
Przecięcie jest domkniętą ścianą, ponieważ warunek definiujący
ścianę jest zachowywany przy przecięciach.
Lemat Kuratowskiego--Zorna zapewnia istnienie minimalnej
niepustej domkniętej ściany $F$.

Gdyby $F$ zawierała różne punkty $x,y$, lokalna wypukłość
i warunek Hausdorffa zapewniałyby ciągły rzeczywisty
funkcjonał liniowy $\ell$ rozróżniający $x$ i $y$.
Dla $\K=\C$ możemy wziąć część rzeczywistą ciągłego
zespolonego funkcjonału. Z Lematu~\ref{lem:max-face9}
wynika, że zbiór punktów, w których $\ell$ osiąga
maksimum na $F$, jest niepustą, domkniętą ścianą $K$,
prawdziwie zawartą w $F$. To przeczy minimalności.
Zatem $F=\{e\}$, a warunek bycia ścianą oznacza,
że $e\in\operatorname{Ext}(K)$.
Jeśli zaczniemy od dowolnej niepustej domkniętej ściany
$F_0$ zamiast od $K$, ten sam argument da jej punkt
ekstremalny, ekstremalny również w $K$.
\end{proof}

\begin{twierdzenie}[Krein--Milman]
\label{thm:krein-milman}
Niech $E$ będzie rzeczywistą lub zespoloną przestrzenią
lokalnie wypukłą Hausdorffa, a $K\subseteq E$ niepustym,
zwartym i wypukłym zbiorem. Wtedy
\[
 K=\overline{\operatorname{conv}(\operatorname{Ext}(K))}.
\]
Domknięcie rozumiemy w wyjściowej topologii $E$.
\end{twierdzenie}
\begin{proof}
Z Lematu~\ref{lem:exist-extreme9} zbiór
$\operatorname{Ext}(K)$ jest niepusty. Ponieważ $K$ jest
wypukły i domknięty, zbiór
$C=\overline{\operatorname{conv}(\operatorname{Ext}(K))}$
jest zawarty w $K$ i jest zwarty jako domknięty podzbiór $K$.

Załóżmy, że istnieje $x_0\in K\setminus C$.
Twierdzenie o ścisłym oddzielaniu punktu od domkniętego
zbioru wypukłego w przestrzeni lokalnie wypukłej
(twierdzenie Hahna--Banacha o oddzielaniu zbiorów wypukłych) daje ciągły
rzeczywisty funkcjonał liniowy $\ell:E\to\R$ spełniający
$\ell(x_0)>\sup_{c\in C}\ell(c)$. W przypadku zespolonym
$\ell$ można otrzymać jako część rzeczywistą funkcjonału
zespolonego.

Niech $M=\max_{x\in K}\ell(x)$.
Z Lematu~\ref{lem:max-face9} zbiór
$F_\ell=\{x\in K:\ell(x)=M\}$ jest niepustą zwartą
ścianą $K$. Z Lematu~\ref{lem:exist-extreme9}
ma on punkt $e\in\operatorname{Ext}(K)$.
Wtedy $e\in C$, a równocześnie
$\ell(e)=M\ge\ell(x_0)>\sup_C\ell$, co jest sprzecznością.
Stąd $C=K$.
\end{proof}

\begin{uwaga}[Dlaczego hipotezy są istotne?]
Zwartość nie może być pominięta: cała przestrzeń $E=\R$
ma niepusty, domknięty i wypukły zbiór $K=\R$,
ale nie ma punktów ekstremalnych. Lokalna wypukłość
gwarantuje istnienie dostatecznie wielu ciągłych
funkcjonałów do oddzielania punktów. W ogólnych TVS
bez lokalnej wypukłości powyższy dowód nie działa.
Twierdzenie nie mówi także, że każdy punkt $K$ jest
\emph{skończoną} kombinacją wypukłą punktów ekstremalnych;
w ogólności konieczne jest domknięcie.
\end{uwaga}

\begin{wniosek}[Zasada maksimum Bauera]
\label{cor:bauer9}
Niech $K$ będzie niepustym zwartym wypukłym podzbiorem
przestrzeni lokalnie wypukłej Hausdorffa $E$ nad $\K$.
Jeżeli $f:K\to\R$ jest wypukła i półciągła z góry,
to istnieje $e\in\operatorname{Ext}(K)$ takie,
że $f(e)=\max_{x\in K}f(x)$.
\end{wniosek}
\begin{proof}
Niech $M=\max_K f$; maksimum istnieje dzięki zwartości $K$
i półciągłości $f$ z góry. Zbiór $A=\{x\in K:f(x)=M\}$
jest niepusty i domknięty. Nie musi być wypukły, ale ma
\emph{własność ekstremalności}: jeśli $tx+(1-t)y\in A$
przy $x,y\in K$ i $0<t<1$, to $x,y\in A$. Rzeczywiście,
$M=f(tx+(1-t)y)\le tf(x)+(1-t)f(y)\le M$,
więc obie nierówności są równościami.

Wybierzmy, korzystając z lematu Kuratowskiego--Zorna,
minimalny niepusty, domknięty podzbiór $B\subseteq A$
mający własność ekstremalności względem $K$.
Argument z przecięciem łańcucha jest taki sam jak
w Lemacie~\ref{lem:exist-extreme9}; nie używa wypukłości $B$.
Gdyby $B$ zawierał różne punkty, istniałby ciągły rzeczywisty
funkcjonał $\ell$ niebędący na nim stałym. Niepusty zbiór
$B_\ell=\{x\in B:\ell(x)=\max_{y\in B}\ell(y)\}$
byłby zwartym, właściwym podzbiorem $B$.
Jeśli $z=tx+(1-t)y\in B_\ell$ dla $x,y\in K$, to
$x,y\in B$, gdyż $B$ jest ekstremalny, a z liniowości
$\ell$ wynika $x,y\in B_\ell$. Jest to sprzeczność
z minimalnością. Zatem $B=\{e\}$; własność ekstremalności
tej jednoelementowej rodziny mówi dokładnie,
że $e\in\operatorname{Ext}(K)$, a $e\in A$ kończy dowód.
\end{proof}

\begin{przyklad}[Produkt przedziałów]
\label{ex:krein-cube9}
Niech $I$ będzie dowolnym niepustym zbiorem,
a $K=[-1,1]^I$ z topologią produktową, uważany za
podzbiór $\R^I$. Jest to zwarty zbiór wypukły na mocy
twierdzenia Tychonowa. Jego punkty ekstremalne to dokładnie
$x\in K$ takie, że $x_i\in\{-1,1\}$ dla każdego $i\in I$.
Istotnie, jeżeli $|x_i|<1$, można zmienić tylko współrzędną
$i$ o małe wartości $\pm\varepsilon$, aby przedstawić
$x$ jako środek dwóch różnych punktów $K$.
Jeśli wszystkie współrzędne są równe $\pm1$, to każdy
rozkład $x=ty+(1-t)z$ w $K$ musi być równy $y=z=x$
w każdej współrzędnej.
Dla nieprzeliczalnego $I$ nie można tutaj zastępować
zwartości produktowej zwartym opisem sekwencyjnym.
\end{przyklad}

\begin{przyklad}[Słabo-* zwarta kula $\ell^1$]
\label{ex:krein-l1-weakstar9}
Na przestrzeni $\ell^1(\N;\R)=c_0(\N;\R)'$ rozważmy
słabą-* topologię $\sigma(\ell^1,c_0)$.
Kula $B=\{x\in\ell^1:\|x\|_1\le1\}$ jest zwarta
w tej topologii na mocy twierdzenia Banacha--Alaoglu.
Jej punktami ekstremalnymi są dokładnie $\pm e_n$.
Dla $\|x\|_1<1$ można rozdzielić $x$ jako średnią
$x\pm\varepsilon e_1$. Jeśli $\|x\|_1=1$ i $x$ ma
co najmniej dwie niezerowe współrzędne, odpowiednio małe
przeniesienie masy między nimi pozwala zapisać
$x=(y+z)/2$ dla różnych $y,z\in B$.
Z kolei $\pm e_n$ są ekstremalne, gdyż równość
$\pm1=(y_n+z_n)/2$ dla $y,z\in B$ wymusza
$y_n=z_n=\pm1$ i zerowanie się innych współrzędnych.
Wobec twierdzenia Kreina--Milmana kula $B$ jest słabo-*
domkniętą otoczką wypukłą tych punktów.
\end{przyklad}


\subsection{Twierdzenie Mackeya--Arensa}

Przypomnijmy, że dla pary dualnej $(E,F)$ topologia
$\tau(E,F)$ jest generowana przez seminormy
$p_K(x)=\sup_{f\in K}|\langle x,f\rangle|$, gdzie $K$
przebiega niepuste, absolutnie wypukłe i zwarte zbiory
w topologii $\sigma(F,E)$.

\begin{definicja}[Topologia zgodna z parą dualną]\label{def:compatible-duality}
Niech $(E,F,\langle\cdot,\cdot\rangle)$ będzie parą dualną.
Lokalnie wypukłą topologię Hausdorffa $\rho$ na $E$ nazywamy
\emph{zgodną z parą dualną $(E,F)$}, jeżeli odwzorowanie
$F\to(E,\rho)'$, $f\mapsto(x\mapsto\langle x,f\rangle)$,
jest izomorfizmem algebraicznym.
\end{definicja}

\begin{lemat}[Zwartość i dominacja funkcjonału]\label{lem:compact-domination}
Niech $(E,F)$ będzie parą dualną, a $K\subseteq F$ niepustym,
absolutnie wypukłym i zwartym zbiorem w $\sigma(F,E)$.
Jeżeli $g\in E^*$ spełnia
$|g(x)|\le p_K(x)$ dla każdego $x\in E$,
to istnieje $f\in K$ takie, że
$g(x)=\langle x,f\rangle$ dla wszystkich $x\in E$.
\end{lemat}
\begin{proof}
Dla każdego skończonego zbioru $S=\{x_1,\ldots,x_m\}\subseteq E$
rozważamy odwzorowanie
$T_S:K\to\K^m$, $T_S(f)=(\langle x_1,f\rangle,\ldots,
\langle x_m,f\rangle)$. Jego obraz jest zwarty i absolutnie wypukły.
Pokażemy, że $y=(g(x_1),\ldots,g(x_m))\in T_S(K)$.
Gdyby tak nie było, twierdzenie o silnym oddzielaniu punktu
od zwartego wypukłego zbioru w skończonym wymiarze dałoby
$a_1,\ldots,a_m\in\K$ takie, że
\[
 \operatorname{Re}\sum_{j=1}^m a_jg(x_j)
 >\sup_{f\in K}\operatorname{Re}
       \sum_{j=1}^m a_j\langle x_j,f\rangle.
\]
Połóżmy $x=\sum_{j=1}^m a_jx_j$.
Absolutna wypukłość $K$ oznacza, że zbiór wartości
$\{\langle x,f\rangle:f\in K\}$ jest zrównoważony,
a więc jego supremum części rzeczywistych jest równe
supremum modułów. Otrzymalibyśmy
$\operatorname{Re}g(x)>p_K(x)$, sprzecznie z założeniem.

Zatem dla każdego skończonego $S$ zbiór
$K_S=\{f\in K:\langle x,f\rangle=g(x)\ \text{dla }x\in S\}$
jest niepusty i domknięty w $K$. Rodzina $(K_S)$ ma
własność skończonych przecięć, ponieważ
$K_{S_1}\cap\cdots\cap K_{S_r}=K_{S_1\cup\cdots\cup S_r}$.
Zwartość $K$ daje punkt we wspólnym przecięciu,
co kończy dowód.
\end{proof}

\begin{twierdzenie}[Mackey--Arens]\label{thm:mackey-arens}
Niech $(E,F,\langle\cdot,\cdot\rangle)$ będzie parą dualną
przestrzeni liniowych nad $\K$. Wówczas topologia
$\tau(E,F)$ jest \emph{najsilniejszą} lokalnie wypukłą
topologią Hausdorffa na $E$ zgodną z tą parą dualną.
Dokładniej:
\begin{enumerate}
\item przestrzeń dualna topologiczna $(E,\tau(E,F))'$ jest
      algebraicznie równa $F$ przez dane parowanie;
\item jeżeli $\rho$ jest dowolną inną lokalnie wypukłą topologią
      Hausdorffa na $E$ zgodną z $(E,F)$, to $\rho\subseteq\tau(E,F)$.
\end{enumerate}
\end{twierdzenie}
\begin{proof}
Z definicji $\tau(E,F)$ każdy element $F$ wyznacza
funkcjonał ciągły na $E$; istotnie, zbiór
$\{af:|a|\le1\}$ jest absolutnie wypukły i zwarty
w $\sigma(F,E)$.

Niech $g\in(E,\tau(E,F))'$. Ciągłość w zerze daje
niepuste absolutnie wypukłe zbiory zwarte
$K_1,\ldots,K_m\subseteq F$ i stałą $C>0$ takie, że
$|g(x)|\le C\max_j p_{K_j}(x)$ dla wszystkich $x\in E$.
Aby uzyskać to oszacowanie, bierzemy bazowe otoczenie
$\{x:\max_j p_{K_j}(x)<\varepsilon\}$ zawarte w
$\{x:|g(x)|<1\}$ i skalujemy $x$; można przyjąć
$C=2/\varepsilon$. Niech $K=K_1+\cdots+K_m$.
Jest to niepusty, absolutnie wypukły zbiór zwarty
w $\sigma(F,E)$ jako obraz zwartego produktu przez dodawanie.
Ponieważ $0\in K_j$, mamy $p_K\ge\max_j p_{K_j}$.
Zatem $|g(x)|\le C p_K(x)=p_{CK}(x)$.
Lemat~\ref{lem:compact-domination}, zastosowany do zwartego
absolutnie wypukłego $CK$, zapewnia $g\in F$.
To dowodzi punktu (1).

Niech teraz $\rho$ będzie zgodną lokalnie wypukłą topologią.
Wystarczy sprawdzić, że każda $\rho$-ciągła seminorma
$p:E\to[0,\infty)$ jest $\tau(E,F)$-ciągła.
Określmy $U=\{x:p(x)\le1\}$ i
\[
 K_p=\{f\in F:|\langle x,f\rangle|\le p(x)
                 \ \text{dla wszystkich }x\in E\}.
\]
Zbiór $U$ jest otoczeniem zera dla $\rho$, a
$K_p=U^\circ$ przy danym utożsamieniu $F=(E,\rho)'$.
Istotnie, jeśli $|\langle x,f\rangle|\le1$ dla $p(x)\le1$,
to dla $p(x)>0$ stosujemy $x/p(x)$, a przy $p(x)=0$
skalujemy $nx$ i przechodzimy z $n\to\infty$.
Na mocy Twierdzenia~\ref{thm:banach-alaoglu-general}
zbiór $K_p$ jest zwarty w $\sigma(F,E)$;
jest również niepusty i absolutnie wypukły.

Pokażemy, że $p=p_{K_p}$. Nierówność $p_{K_p}\le p$
wynika z definicji. Dla ustalonego $x\ne0$ funkcjonał
na $\K x$, $\lambda x\mapsto\lambda p(x)$,
jest dominowany przez seminormę $p$ co do modułu.
Twierdzenie Hahna--Banacha (w wersji dla seminormy)
przedłuża go do funkcjonału $f\in E^*$ z
$|f(y)|\le p(y)$ dla każdego $y\in E$.
Funkcjonał ten jest $\rho$-ciągły, należy więc do $F$,
a w konsekwencji $f\in K_p$ i $|f(x)|=p(x)$.
Zachodzi zatem $p=p_{K_p}$, czyli $p$ jest
$\tau(E,F)$-ciągła. Ponieważ topologia $\rho$ jest
generowana przez swoje ciągłe seminormy, otrzymujemy
$\rho\subseteq\tau(E,F)$.
\end{proof}

\begin{wniosek}[Najsłabsza i najsilniejsza topologia zgodna]\label{cor:mackey-extremes}
Niech $(E,F)$ będzie parą dualną. Topologia słaba
$\sigma(E,F)$ jest najsłabszą, zaś topologia Mackeya
$\tau(E,F)$ najsilniejszą lokalnie wypukłą topologią
Hausdorffa zgodną z tą parą. W szczególności dla każdej
takiej topologii $\rho$ zachodzi
$\sigma(E,F)\subseteq\rho\subseteq\tau(E,F)$.
\end{wniosek}
\begin{proof}
Jeżeli $\rho$ jest zgodna, wszystkie funkcjonały z $F$
są $\rho$-ciągłe, więc minimalność topologii początkowej
daje $\sigma(E,F)\subseteq\rho$. Druga inkluzja wynika
z Twierdzenia~\ref{thm:mackey-arens}.
\end{proof}

\begin{przyklad}[Przestrzeń słaba i topologia Mackeya nie muszą być równe]
Niech $E$ będzie nieskończenie wymiarową przestrzenią Banacha,
a $F=E'$. Topologia $\sigma(E,E')$ nie jest normowalna
(co udowodniono w \hyperref[sec:topologie-slabe]{rozdziale 8}), a topologia normowa
jest zgodna z parą $(E,E')$. Z wniosku
\ref{cor:mackey-extremes} wynika więc, że
$\sigma(E,E')\subsetneq\tau(E,E')$.
Nie utożsamiamy tu jednak automatycznie $\tau(E,E')$
z topologią normową: odpowiednie twierdzenie dla przestrzeni
Banacha powiążemy w dalszej części skryptu z zasadą
jednostajnej ograniczoności.
\end{przyklad}

\begin{uwaga}[Dalsza perspektywa: przestrzenie dystrybucji]
W teorii dystrybucji zbiory polarne otoczeń zera w przestrzeni
funkcji testowych są słabo-* zwarte, natomiast topologia silna
na przestrzeni dystrybucji oznacza zbieżność jednostajną
na zbiorach ograniczonych funkcji testowych. Różnica między
zbiorem ograniczonym, zwartym i równociągłym będzie podstawowa
w opisie topologii $\mathcal D'(\Omega)$.
\end{uwaga}

\subsection{Zastosowanie zwartości: nieskończenie wiele warunków liniowych}

\begin{twierdzenie}[Kryterium przedłużenia z wieloma warunkami]
\label{thm:helly-interpolation}
Niech $E$ będzie przestrzenią unormowaną nad $\K$,
$R\ge0$, $I$ dowolnym zbiorem, $(x_i)_{i\in I}$ rodziną
punktów $E$, a $(b_i)_{i\in I}$ rodziną skalarów.
Istnieje funkcjonał $f\in E'$ spełniający
$\|f\|\le R$ i $f(x_i)=b_i$ dla każdego $i\in I$
wtedy i tylko wtedy, gdy dla dowolnego $m\in\N$,
$i_1,\ldots,i_m\in I$ oraz $c_1,\ldots,c_m\in\K$
zachodzi
\[
 \left|\sum_{j=1}^{m}c_jb_{i_j}\right|
 \le R\left\|\sum_{j=1}^{m}c_jx_{i_j}\right\|.
\]
\end{twierdzenie}
\begin{proof}
Konieczność wynika z liniowości funkcjonału i oszacowania
$|f(x)|\le\|f\|\,\|x\|$.

Udowodnijmy dostateczność w sposób wykorzystujący zwartość.
Niech $B_R'=\{f\in E':\|f\|\le R\}$; jest to słabo-* zwarty
zbiór na mocy Wniosku~\ref{cor:alaoglu-norm}. Dla $i\in I$
połóżmy $C_i=\{f\in B_R':f(x_i)=b_i\}$.
Każdy $C_i$ jest słabo-* domknięty w $B_R'$.
Dla ustalonego skończonego $J\subseteq I$ definiujemy
na $M_J=\operatorname{lin}\{x_j:j\in J\}$ funkcjonał
$g_J(\sum_{j\in J}c_jx_j)=\sum_{j\in J}c_jb_j$.
Warunek z tezy zapewnia, że definicja nie zależy od
przedstawienia argumentu, a ponadto $|g_J(x)|\le R\|x\|$.
Twierdzenie Hahna--Banacha daje przedłużenie
$f_J\in E'$ z $\|f_J\|\le R$; stąd
$f_J\in\bigcap_{j\in J} C_j$.
Rodzina $(C_i)_{i\in I}$ ma więc własność skończonych przecięć.
Zwartość $B_R'$ daje $f\in\bigcap_{i\in I}C_i$,
który spełnia wszystkie żądane warunki.
\end{proof}

\begin{przyklad}[Przedłużanie wartości na wektorach współrzędnych]
Niech $E=\ell^1(\N)$, $x_n=e_n$ i $R>0$.
Dla dowolnej rodziny liczb $b_n$ ograniczonej przez
$|b_n|\le R$ otrzymujemy funkcjonał $f\in E'$
spełniający $f(e_n)=b_n$ i $\|f\|\le R$:
\[
 \left|\sum_{j=1}^m c_jb_{i_j}\right|
 \le R\sum_{n=1}^\infty
       \left|\sum_{j:i_j=n}c_j\right|
 =R\left\|\sum_{j=1}^m c_je_{i_j}\right\|_1.
\]
Jest to zgodne z identyfikacją $(\ell^1)'=\ell^\infty$.
\end{przyklad}


% Wersja źródłowa: audyt II, jednostka 10
\setcounter{section}{9}
\section{Przestrzenie beczkowe i bornologiczne. Twierdzenie Banacha--Steinhausa}\label{sec:beczkowe-bornologiczne}

\subsection{Od ograniczoności punktowej do równociągłości}

W rozdziale o twierdzeniu Banacha--Alaoglu pokazaliśmy, że rodzina
$H\subseteq E'$ jest równociągła dokładnie wtedy, gdy mieści się
w zbiorze polarnym pewnego otoczenia zera. Równociągłość implikuje
ograniczoność punktową, lecz przykład z funkcjonałami $f_n(x)=nx_n$ na $c_{00}$
na $c_{00}$ pokazuje, że implikacji nie można odwrócić w każdej
przestrzeni liniowo-topologicznej. Zbadamy teraz, jakie dodatkowe
własności przestrzeni dziedziny gwarantują tę implikację.

W całym rozdziale $\K\in\{\R,\C\}$. Przestrzenie $E$ i $F$
są lokalnie wypukłymi przestrzeniami liniowo-topologicznymi Hausdorffa,
chyba że wyraźnie zaznaczono inaczej. Symbol $E'$ oznacza przestrzeń
dualną topologiczną, a $\mathcal L(E,F)$ zbiór ciągłych operatorów
liniowych z $E$ do $F$. Przez zbiór ograniczony rozumiemy zbiór
ograniczony w sensie von Neumanna. Zbiór $C$ jest \emph{absolutnie
wypukły}, jeżeli jest wypukły i zrównoważony.

\begin{definicja}[Beczka i przestrzeń beczkowa]\label{def:beczka}
Niech $E$ będzie lokalnie wypukłą TVS. Zbiór $B\subseteq E$
nazywamy \emph{beczką}, jeżeli jest domknięty, absolutnie wypukły
i pochłaniający, to znaczy dla każdego $x\in E$ istnieje $t>0$
takie, że $x\in tB$. Przestrzeń $E$ jest \emph{beczkowa}, jeśli
każda beczka w $E$ jest otoczeniem zera.
\end{definicja}

Warunek pochłaniania dotyczy pojedynczych punktów, a nie całych
zbiorów ograniczonych. To rozróżnienie będzie istotne przy
wprowadzaniu przestrzeni bornologicznych.

\begin{lemat}[Zbiór polarny i ograniczoność punktowa]\label{lem:polar-barrel}
Niech $H\subseteq E'$ i niech
$H^\circ=\{x\in E:|f(x)|\le1\ \text{dla każdego }f\in H\}$.
Wówczas $H^\circ$ jest domknięty i absolutnie wypukły.
Ponadto $H^\circ$ jest pochłaniający wtedy i tylko wtedy, gdy
$H$ jest punktowo ograniczony, czyli
$\sup_{f\in H}|f(x)|<\infty$ dla każdego $x\in E$.
\end{lemat}
\begin{proof}
Każdy zbiór $\{x:|f(x)|\le1\}$ jest domknięty i absolutnie
wypukły, a przecięcie dowolnej rodziny takich zbiorów
zachowuje obie własności. Jeśli $H$ jest punktowo ograniczony,
to dla ustalonego $x$ liczba $M_x=\sup_{f\in H}|f(x)|$
jest skończona; przyjmujemy $M_x=0$, gdy $H=\varnothing$.
Dla $t>\max\{1,M_x\}$ mamy $x/t\in H^\circ$, czyli
$x\in tH^\circ$. Odwrotnie, jeśli $x\in tH^\circ$, to
$|f(x)|\le t$ dla wszystkich $f\in H$. W obu kierunkach
dopuszczamy również $x=0$ i rodzinę pustą.
\end{proof}

\begin{twierdzenie}[Charakteryzacja beczkowości]\label{thm:barrelled-char}
Niech $E$ będzie lokalnie wypukłą TVS Hausdorffa.
Następujące warunki są równoważne:
\begin{enumerate}[label=\textup{(\roman*)}]
\item $E$ jest beczkowa;
\item każda punktowo ograniczona rodzina $H\subseteq E'$
      jest równociągła;
\item każda rodzina $H\subseteq E'$ ograniczona w
      $\sigma(E',E)$ jest równociągła.
\end{enumerate}
\end{twierdzenie}
\begin{proof}
Warunki (ii) i (iii) są równoważne, gdyż topologia
$\sigma(E',E)$ jest topologią punktowej zbieżności,
a ograniczoność zbioru w produkcie jest równoważna
ograniczoności każdej jego projekcji na $\K$.

Załóżmy (i) i niech $H$ będzie punktowo ograniczona.
Z lematu~\ref{lem:polar-barrel} wynika, że $H^\circ$
jest beczką, zatem jest otoczeniem zera. Dla
$u\in H^\circ$ zachodzi $|f(u)|\le1$ dla wszystkich
$f\in H$, czyli $H$ jest równociągła.

Załóżmy (ii) i niech $B\subseteq E$ będzie beczką.
Zbiór polarny $B^\circ\subseteq E'$ jest punktowo ograniczony:
dla każdego $x$ istnieje $t>0$ z $x\in tB$, zatem
$|f(x)|\le t$ dla wszystkich $f\in B^\circ$.
Z (ii) otrzymujemy otoczenie $U$ takie, że
$B^\circ\subseteq U^\circ$. Wskutek odwracania inkluzji
przez operację przejścia do zbioru polarnego i twierdzenia
bipolarnego dla lokalnie wypukłej przestrzeni Hausdorffa
zachodzi
$U\subseteq U^{\circ\circ}\subseteq B^{\circ\circ}=B$.
Zatem $B$ jest otoczeniem zera.
\end{proof}

\begin{uwaga}[Dlaczego konieczna jest lokalna wypukłość?]
W dowodzie odwrotnej implikacji użyliśmy równości
$B^{\circ\circ}=B$ dla domkniętego zbioru absolutnie
wypukłego. Opiera się ona na twierdzeniu Hahna--Banacha
i może zawieść w ogólnej TVS. Charakteryzację
\ref{thm:barrelled-char} formułujemy więc w klasie
przestrzeni lokalnie wypukłych Hausdorffa.
\end{uwaga}

\subsection{Własność Baire'a i przestrzenie Frécheta}

\begin{definicja}[Przestrzeń Baire'a]
Przestrzeń topologiczną $X$ nazywamy \emph{przestrzenią Baire'a},
jeżeli każde przeliczalne przecięcie otwartych zbiorów gęstych
w $X$ jest gęste. Równoważnie: przeliczalna suma zbiorów
domkniętych o pustym wnętrzu ma puste wnętrze.
\end{definicja}

\begin{twierdzenie}[Baire: zupełność metryczna wystarcza]
\label{thm:baire-complete}
Każda niepusta zupełna przestrzeń metryczna $(X,d)$
jest przestrzenią Baire'a.
\end{twierdzenie}
\begin{proof}
Niech $G_1,G_2,\ldots$ będą otwartymi podzbiorami
gęstymi w $X$. Pokażemy, że ich przecięcie przecina
każdy niepusty zbiór otwarty $O\subseteq X$.
Wybieramy indukcyjnie punkty $x_n\in X$ oraz
liczby $0<r_n<2^{-n}$ tak, by
\[
 \overline{B}(x_1,r_1)\subseteq O\cap G_1,\qquad
 \overline{B}(x_{n+1},r_{n+1})\subseteq
 B(x_n,r_n)\cap G_{n+1}.
\]
W pierwszym kroku $O\cap G_1$ jest niepusty i otwarty,
więc zawiera domkniętą kulę o dodatnim promieniu.
Analogicznie $B(x_n,r_n)\cap G_{n+1}$ jest niepusty
i otwarty, ponieważ $G_{n+1}$ jest gęsty. Możemy
zatem kontynuować konstrukcję.

Dla $m>n$ mamy $x_m\in B(x_n,r_n)$, więc
$d(x_m,x_n)<r_n<2^{-n}$. Ciąg $(x_n)$ jest ciągiem
Cauchy'ego. Z zupełności otrzymujemy punkt $x\in X$
z $x_n\to x$. Dla ustalonego $n$ wszystkie punkty
$x_m$ o $m>n$ leżą w $\overline B(x_n,r_n)$,
a domkniętość tej kuli daje
$x\in\overline B(x_n,r_n)\subseteq G_n$.
Ponadto $x\in\overline B(x_1,r_1)\subseteq O$.
Wobec dowolności $O$ przecięcie $\bigcap_nG_n$
jest gęste w $X$.
\end{proof}

\begin{twierdzenie}[Przestrzenie Baire'a są beczkowe]\label{thm:baire-barrelled}
Niech $E$ będzie lokalnie wypukłą TVS Hausdorffa, której
przestrzeń topologiczna jest przestrzenią Baire'a.
Wówczas $E$ jest beczkowa.
\end{twierdzenie}
\begin{proof}
Niech $B$ będzie beczką w $E$. Ponieważ jest pochłaniająca
i zrównoważona, $E=\bigcup_{n=1}^\infty nB$. Każdy zbiór
$nB$ jest domknięty, bo mnożenie przez niezerowy skalar
jest homeomorfizmem. Z własności Baire'a otrzymujemy
$\operatorname{int}(n_0B)\ne\varnothing$ dla pewnego
$n_0\ge1$, a zatem $\operatorname{int}B\ne\varnothing$.

Wybierzmy $y\in\operatorname{int}B$. Zrównoważenie daje
$-y\in\operatorname{int}B$. Istnieje zrównoważone
otwarte otoczenie $V$ zera takie, że
$y+V\subseteq B$ i $-y+V\subseteq B$. Dla każdego
$v\in V$ wypukłość $B$ daje
$v=\frac12(y+v)+\frac12(-y+v)\in B$.
Zatem $V\subseteq B$ i $B$ jest otoczeniem zera.
\end{proof}

\begin{wniosek}[Przestrzenie Frécheta i Banacha]\label{cor:frechet-barrelled}
Każda przestrzeń Frécheta i każda przestrzeń Banacha
jest beczkowa. Również każda lokalnie wypukła TVS
Hausdorffa, której przestrzeń topologiczna jest Baire'a,
jest beczkowa.
\end{wniosek}
\begin{proof}
Z Twierdzenia~\ref{thm:baire-complete} wynika, że
przestrzeń Frécheta jest Baire'a.
Przestrzeń Banacha jest szczególnym przypadkiem przestrzeni
Frécheta. Pozostała część jest
Twierdzeniem~\ref{thm:baire-barrelled}.
\end{proof}

\begin{uwaga}[Dwa twierdzenia o nazwie Baire'a]
Należy odróżnić własność Baire'a przestrzeni topologicznej
od twierdzenia Baire'a, które pozwala ją uzyskać z
zupełności metrycznej. Dowód twierdzenia Baire'a podaliśmy powyżej;
w następnym rozdziale zastosujemy go do twierdzeń
o odwzorowaniu otwartym i wykresie domkniętym.
\end{uwaga}

\subsection{Twierdzenie Banacha--Steinhausa w przestrzeniach lokalnie wypukłych}

\begin{definicja}[Ograniczoność punktowa i równociągłość operatorów]
Niech $\mathcal T\subseteq\mathcal L(E,F)$.
Rodzina $\mathcal T$ jest \emph{punktowo ograniczona}, gdy
$\{Tx:T\in\mathcal T\}$ jest ograniczony w $F$ dla każdego
$x\in E$. Jest \emph{równociągła}, gdy dla każdego otoczenia
$V$ zera w $F$ istnieje otoczenie $U$ zera w $E$
takie, że $T(U)\subseteq V$ dla każdego $T\in\mathcal T$.
\end{definicja}

\begin{twierdzenie}[Banach--Steinhaus dla przestrzeni beczkowych]
\label{thm:banach-steinhaus-lcs}
Niech $E$ będzie beczkową lokalnie wypukłą TVS Hausdorffa,
$F$ lokalnie wypukłą TVS Hausdorffa, a
$\mathcal T\subseteq\mathcal L(E,F)$ rodziną operatorów.
Jeżeli $\mathcal T$ jest punktowo ograniczona, to jest
równociągła.
\end{twierdzenie}
\begin{proof}
Ustalmy otoczenie $V$ zera w $F$. Z lokalnej wypukłości
możemy wybrać ciągłe seminormy $q_1,\ldots,q_m$ na $F$
i $\varepsilon>0$ takie, aby
$W=\{y:\max_jq_j(y)\le\varepsilon\}\subseteq V$.
Zbiór $W$ jest domknięty, absolutnie wypukły i stanowi
otoczenie zera w $F$. Określmy
\[
 B=\bigcap_{T\in\mathcal T}T^{-1}(W).
\]
Jest to zbiór domknięty i absolutnie wypukły w $E$.
Dla ustalonego $x\in E$ zbiór $\{Tx:T\in\mathcal T\}$
jest ograniczony w $F$. Istnieje więc $t_x>0$ takie,
że $Tx\in t_x W$ dla każdego $T\in\mathcal T$,
a zatem $x\in t_xB$. Zbiór $B$ jest beczką.
Beczkowość $E$ daje otoczenie zera $U\subseteq B$,
więc $T(U)\subseteq W\subseteq V$ dla wszystkich
$T\in\mathcal T$.
\end{proof}

\begin{wniosek}[Klasyczna zasada jednostajnej ograniczoności]
\label{cor:uniform-boundedness}
Niech $E$ będzie przestrzenią Banacha, a $F$ przestrzenią
unormowaną nad $\K$. Jeżeli rodzina
$\mathcal T\subseteq\mathcal L(E,F)$ spełnia
$\sup_{T\in\mathcal T}\|Tx\|_F<\infty$ dla każdego
$x\in E$, to $\sup_{T\in\mathcal T}\|T\|<\infty$.
Nie zakładamy zupełności $F$.
\end{wniosek}
\begin{proof}
Przestrzeń $E$ jest beczkowa. Stosujemy
Twierdzenie~\ref{thm:banach-steinhaus-lcs} do kuli
$V=\{y:\|y\|_F<1\}$. Istnieje $\delta>0$ takie, że
$\|Tx\|_F<1$ dla wszystkich $T\in\mathcal T$
i $\|x\|_E<\delta$. Przez skalowanie otrzymujemy
$\|T\|\le 2/\delta$ dla każdego $T\in\mathcal T$.
\end{proof}

\begin{proof}[Drugi dowód wniosku: bezpośrednio z twierdzenia Baire'a]
Dla $n\in\N$ niech
$A_n=\{x\in E:\sup_{T\in\mathcal T}\|Tx\|_F\le n\}$.
Są to zbiory domknięte, ponieważ każdy $T$ jest ciągły.
Ograniczoność punktowa oznacza
$E=\bigcup_{n=1}^\infty A_n$. Na mocy twierdzenia Baire'a
pewien $A_N$ ma niepuste wnętrze. Istnieją więc
$x_0\in E$ i $r>0$ takie, że
$x_0+\{h:\|h\|_E<r\}\subseteq A_N$.
Dla $\|h\|_E<r$ i każdego $T\in\mathcal T$ mamy
$\|Th\|_F\le\|T(x_0+h)\|_F+\|Tx_0\|_F\le2N$.
Skalowanie daje $\|T\|\le4N/r$, niezależnie od $T$.
W dowodzie nie wystąpiła zupełność przestrzeni wartości $F$.
\end{proof}

\begin{przyklad}[Dlaczego zupełność dziedziny jest istotna?]
\label{ex:c00-not-barrelled}
Niech $E=c_{00}$ z normą $\|x\|_\infty=\sup_n|x_n|$.
Dla $n\ge1$ określmy $f_n(x)=nx_n$. Każdy $f_n\in E'$,
ponieważ $|f_n(x)|\le n\|x\|_\infty$.
Dla ustalonego $x\in c_{00}$ ciąg $(f_n(x))$ jest
ostatecznie zerowy, więc rodzina $H=\{f_n:n\ge1\}$ jest
punktowo ograniczona. Jednak $\|f_n\|=n\to\infty$,
zatem rodzina nie jest równociągła.

Wprost widać też beczkę niebędącą otoczeniem zera:
\[
 B=\{x\in c_{00}: n|x_n|\le1\ \text{dla każdego }n\ge1\}.
\]
Jest ona domknięta, absolutnie wypukła i pochłaniająca,
gdyż każdy $x$ ma skończony nośnik. Nie jest otoczeniem
zera: dla $x^{(n)}=n^{-1/2}e_n$ zachodzi
$\|x^{(n)}\|_\infty\to0$, podczas gdy
$n|x_n^{(n)}|=\sqrt n>1$ dla $n>1$. Tym samym
$E$ nie jest beczkowa.
\end{przyklad}

\begin{przyklad}[Topologia słaba na przestrzeni Banacha]
\label{ex:weak-not-barrelled}
Niech $E$ będzie nieskończenie wymiarową przestrzenią Banacha
i wyposażmy ją w topologię $\sigma(E,E')$. Domknięta kula
normowa $B=\{x:\|x\|\le1\}$ jest słabo domknięta
(z twierdzenia Hahna--Banacha), absolutnie wypukła i pochłaniająca.
Nie jest jednak słabym otoczeniem zera, gdyż każde słabe
otoczenie zawiera niezerową prostą, która nie mieści się
w $B$. Zatem $(E,\sigma(E,E'))$ nie jest beczkowa,
mimo że $E$ jest zupełna w swojej oryginalnej normie.
O beczkowości decyduje \emph{wybrana topologia}, a nie tylko
struktura algebraiczna przestrzeni.
\end{przyklad}

\subsection{Trzy konsekwencje zasady jednostajnej ograniczoności}

\begin{twierdzenie}[Punktowa zbieżność operatorów]\label{thm:pointwise-limit-cont}
Niech $E$ będzie beczkową lokalnie wypukłą TVS Hausdorffa,
a $F$ lokalnie wypukłą TVS Hausdorffa. Załóżmy, że
$(T_i)_{i\in I}$ jest siecią w $\mathcal L(E,F)$,
rodzina $\{T_i:i\in I\}$ jest punktowo ograniczona,
a dla każdego $x\in E$ istnieje granica
$T(x)=\lim_iT_i(x)$ w $F$. Wówczas $T:E\to F$
jest ciągłym operatorem liniowym, a rodzina
$\{T_i:i\in I\}$ jest równociągła.
\end{twierdzenie}
\begin{proof}
Liniowość $T$ wynika z ciągłości dodawania i mnożenia
przez skalar w $F$. Z ograniczoności punktowej i
Twierdzenia~\ref{thm:banach-steinhaus-lcs} otrzymujemy
równociągłość rodziny. Dla danego otoczenia $V$ zera
w $F$ wybieramy domknięte otoczenie $W\subseteq V$.
Istnieje otoczenie $U$ zera w $E$ takie, że
$T_i(U)\subseteq W$ dla wszystkich $i$.
Przechodząc do granicy przy ustalonym $x\in U$, otrzymujemy
$T(x)\in W\subseteq V$. Wobec dowolności $V$ operator
$T$ jest ciągły w zerze, a z liniowości także na całym $E$.
\end{proof}

\begin{wniosek}[Granica punktowa ciągu operatorów]
\label{cor:sequence-operators}
Niech $E$ i $F$ spełniają założenia
Twierdzenia~\ref{thm:pointwise-limit-cont}, a
$T_n\in\mathcal L(E,F)$ dla $n\ge1$. Jeśli
$T_n(x)\to T(x)$ dla każdego $x\in E$, to $T$ jest
ciągłym operatorem liniowym, a rodzina $\{T_n:n\ge1\}$
jest równociągła.
\end{wniosek}
\begin{proof}
Dla każdego $x\in E$ zbiór wartości zbieżnego ciągu
$(T_n(x))$ jest ograniczony w $F$ (jego ogon zbiega do zera
po odjęciu granicy, a skończony początek jest ograniczony).
Stosujemy Twierdzenie~\ref{thm:pointwise-limit-cont}.
\end{proof}

\begin{uwaga}[Dlaczego nie wystarcza sama zbieżność sieci?]
Dla sieci zbieżnej zbiór wszystkich jej wartości nie musi
być ograniczony. Niech $I$ oznacza zbiór skończonych
podzbiorów $\N$, uporządkowany inkluzją, i niech
$a_F=|F|$, gdy $1\notin F$, oraz $a_F=0$, gdy $1\in F$.
Wtedy $a_F\to0$ jako sieć w $\R$, ponieważ na całym
ogonie zaczynającym się od $\{1\}$ mamy $a_F=0$.
Jednocześnie zbiór wszystkich $a_F$ jest nieograniczony.
Przyjmując $E=F=\R$ i $T_F(x)=a_Fx$, otrzymujemy sieć
ciągłych operatorów liniowych zbieżną punktowo do zera,
która nie jest równociągła. Dlatego w twierdzeniu dla
ogólnych sieci jawnie zakładamy ograniczoność punktową.
\end{uwaga}

\begin{twierdzenie}[Jednostajność zbieżności na zbiorach zwartych]
\label{thm:equicont-compact-conv}
Niech $E$ i $F$ będą TVS Hausdorffa,
$\mathcal T\subseteq\mathcal L(E,F)$ rodziną równociągłą,
a $(T_i)_{i\in I}$ siecią z $\mathcal T$ zbieżną punktowo
do operatora $T\in\mathcal L(E,F)$.
Wówczas $T_i\to T$ jednostajnie na każdym zwartym
zbiorze $K\subseteq E$: dla każdego otoczenia $V$ zera
w $F$ istnieje $i_0\in I$ takie, że
$(T_i-T)(K)\subseteq V$ dla $i\ge i_0$.
\end{twierdzenie}
\begin{proof}
Ustalmy $K$ i otoczenie $V$. Wybierzmy zrównoważone
otoczenie $W$ z $W+W+W\subseteq V$. Równociągłość
rodziny $\mathcal T$ i ciągłość $T$ dają wspólne
otoczenie zera $U$ w $E$, które możemy wybrać
zrównoważone, tak że $T_i(U)\subseteq W$ dla wszystkich
$i$ oraz $T(U)\subseteq W$.
Wybieramy $U$ dodatkowo otwarte. Zwartość $K$
pozwala pokryć go skończenie wieloma zbiorami
$x_1+U,\ldots,x_m+U$ o środkach $x_j\in K$.
Punktowa zbieżność zapewnia indeks $i_0$ z
$(T_i-T)x_j\in W$ dla wszystkich $j$ i $i\ge i_0$.
Jeśli $x\in K\cap(x_j+U)$, to
$(T_i-T)x=(T_i-T)x_j+T_i(x-x_j)-T(x-x_j)\in W+W+W
\subseteq V$, co kończy dowód.
\end{proof}

\begin{stwierdzenie}[Ograniczoność w topologiach operatorowych]
\label{prop:sot-bot-bounded}
Niech $E$ będzie beczkową lokalnie wypukłą TVS Hausdorffa,
a $F$ lokalnie wypukłą TVS Hausdorffa. Na $\mathcal L(E,F)$
rozważmy topologię $\tau_s$ zbieżności punktowej i
$\tau_b$ zbieżności jednostajnej na zbiorach ograniczonych.
Wtedy rodzina $\mathcal T\subseteq\mathcal L(E,F)$
jest ograniczona w $\tau_s$ wtedy i tylko wtedy,
gdy jest ograniczona w $\tau_b$.
\end{stwierdzenie}
\begin{proof}
Ponieważ singletony są ograniczone, $\tau_b$ jest
silniejsza od $\tau_s$, więc ograniczoność w $\tau_b$
implikuje ograniczoność w $\tau_s$.

Załóżmy odwrotnie, że $\mathcal T$ jest ograniczona
w $\tau_s$. Zbieżność punktowa oznacza, że dla każdego
$x\in E$ zbiór $\{Tx:T\in\mathcal T\}$ jest ograniczony
w $F$. Z twierdzenia Banacha--Steinhausa rodzina jest
równociągła. Dla ograniczonego $B\subseteq E$ i otoczenia
$V$ zera w $F$ istnieje otoczenie $U$ zera w $E$
takie, że $T(U)\subseteq V$ dla wszystkich $T$.
Ponieważ $B$ jest ograniczony, istnieje $t>0$
z $B\subseteq tU$. Stąd $T(B)\subseteq tV$ dla wszystkich
$T\in\mathcal T$. Oznacza to ograniczoność w $\tau_b$.
\end{proof}

\subsection{Przestrzenie bornologiczne}

\begin{definicja}[Zbiór pochłaniający zbiory ograniczone]
\label{def:bornivorous}
Niech $E$ będzie TVS. Zbiór $C\subseteq E$ nazywamy
\emph{pochłaniającym zbiory ograniczone}, jeśli dla każdego
ograniczonego $B\subseteq E$ istnieje $t>0$ takie, że
$B\subseteq tC$. W szczególności $C$ jest pochłaniający,
ponieważ wszystkie singletony są ograniczone.
\end{definicja}

\begin{definicja}[Przestrzeń bornologiczna]\label{def:bornological}
Lokalnie wypukłą TVS Hausdorffa $E$ nazywamy
\emph{bornologiczną}, jeżeli każdy absolutnie wypukły
zbiór pochłaniający zbiory ograniczone w $E$ jest
otoczeniem zera. Nie wymagamy jego domkniętości.
\end{definicja}

\begin{lemat}[Funkcjonał Minkowskiego i pochłanianie ograniczonych]
\label{lem:minkowski-bornivorous}
Niech $C\subseteq E$ będzie absolutnie wypukły i
pochłaniający. Jego funkcjonał Minkowskiego
$p_C(x)=\inf\{t>0:x\in tC\}$ jest skończoną seminormą.
Zbiór $C$ pochłania wszystkie zbiory ograniczone wtedy
i tylko wtedy, gdy $p_C$ jest ograniczona na każdym
zbiorze ograniczonym w $E$.
\end{lemat}
\begin{proof}
Własności seminormy wynikają z absolutnej wypukłości i
pochłaniania. Jeżeli $B\subseteq tC$, to $p_C(x)\le t$
dla $x\in B$. W drugą stronę, jeśli
$M=\sup_{x\in B}p_C(x)<\infty$, wybieramy $t>M$.
Dla każdego $x\in B$ definicja kresu dolnego zapewnia
$x\in s_x C$ dla pewnego $s_x<t$. Ponieważ $C$ jest
zrównoważony, $s_xC\subseteq tC$, stąd $B\subseteq tC$.
\end{proof}

\begin{twierdzenie}[Kryteria bornologiczności]\label{thm:bornological-equiv}
Niech $E$ będzie lokalnie wypukłą TVS Hausdorffa.
Następujące warunki są równoważne:
\begin{enumerate}[label=\textup{(\roman*)}]
\item $E$ jest bornologiczna;
\item każda seminorma $p:E\to[0,\infty)$, która jest
ograniczona na każdym ograniczonym zbiorze $B\subseteq E$,
jest ciągła;
\item dla każdej lokalnie wypukłej TVS Hausdorffa $F$
każdy liniowy operator $T:E\to F$, który odwzorowuje
zbiory ograniczone w $E$ na zbiory ograniczone w $F$,
jest ciągły.
\end{enumerate}
\end{twierdzenie}
\begin{proof}
(i)$\Rightarrow$(ii). Dla seminormy $p$ jak w (ii)
zbiór $C=\{x:p(x)\le1\}$ jest absolutnie wypukły.
Dla każdego ograniczonego $B$ mamy $M_B=\sup_Bp<\infty$,
więc $B\subseteq tC$ przy dowolnym $t>\max\{1,M_B\}$.
Bornologiczność oznacza, że $C$ jest otoczeniem zera,
co daje ciągłość $p$.

(ii)$\Rightarrow$(iii). Niech $T$ będzie operatorem
odwzorowującym zbiory ograniczone na ograniczone.
Każda ciągła seminorma $q$ na $F$ jest ograniczona
na ograniczonych podzbiorach $F$. Zatem $q\circ T$
jest seminormą ograniczoną na wszystkich ograniczonych
podzbiorach $E$. Z (ii) jest ciągła. Ponieważ topologię
$F$ wyznaczają ciągłe seminormy, operator $T$ jest ciągły.

(iii)$\Rightarrow$(i). Niech $C\subseteq E$ będzie
absolutnie wypukły i pochłaniający wszystkie zbiory
ograniczone. Jego funkcjonał Minkowskiego $p_C$
jest seminormą ograniczoną na każdym takim zbiorze.
Rozważmy przestrzeń $F$, której zbiorem punktów jest $E$,
a topologia jest generowana przez seminormy ciągłe
pierwotnej topologii oraz dodatkowo przez $p_C$.
Jest to lokalnie wypukła TVS Hausdorffa. Identyczność
$T:(E,\text{topologia pierwotna})\to F$ jest liniowa
i odwzorowuje zbiory ograniczone na ograniczone:
wszystkie seminormy wyznaczające topologię $F$
pozostają ograniczone na każdym ograniczonym $B\subseteq E$.
Z (iii) operator $T$ jest ciągły, czyli $p_C$ jest
ciągły w pierwotnej topologii. Z definicji funkcjonału
Minkowskiego wynika $\{x:p_C(x)<1\}\subseteq C$,
zatem $C$ jest otoczeniem zera. Dowiedliśmy (i).
\end{proof}

\begin{twierdzenie}[Przestrzenie unormowane są bornologiczne]
\label{thm:normed-bornological}
Każda przestrzeń unormowana nad $\K$, także niezupełna,
jest bornologiczna.
\end{twierdzenie}
\begin{proof}
Niech $E$ będzie przestrzenią unormowaną, a $C\subseteq E$
zbiorem absolutnie wypukłym, który pochłania wszystkie
zbiory ograniczone. W szczególności pochłania normową
kulę jednostkową $B=\{x:\|x\|\le1\}$, to znaczy
$B\subseteq tC$ dla pewnego $t>0$. Wtedy
$\{x:\|x\|<1/t\}\subseteq C$, więc $C$ jest
otoczeniem zera.
\end{proof}

\begin{przyklad}[Bornologiczna, ale niebeczkowa]
Przestrzeń $c_{00}$ z normą supremum jest bornologiczna
na mocy Twierdzenia~\ref{thm:normed-bornological}, lecz
nie jest beczkowa, jak pokazaliśmy w
Przykładzie~\ref{ex:c00-not-barrelled}.
Zatem bornologiczność nie wystarcza do zastosowania
twierdzenia Banacha--Steinhausa.
\end{przyklad}

\begin{uwaga}[Dwie różne zasady ciągłości]
Beczkowość pozwala przejść od \emph{punktowej ograniczoności
rodziny ciągłych operatorów} do jej równociągłości.
Bornologiczność pozwala zaś przejść od \emph{ograniczoności
pojedynczego operatora liniowego} (przenoszenia zbiorów
ograniczonych na ograniczone) do jego ciągłości.
Te dwa pojęcia mają różne cele i nie są zamienne.
\end{uwaga}

\subsection{Powrót do topologii Mackeya i przestrzeni dystrybucji}

\begin{twierdzenie}[Przestrzenie beczkowe są przestrzeniami Mackeya]
\label{thm:barrelled-mackey}
Niech $E$ będzie beczkową lokalnie wypukłą TVS Hausdorffa
z topologią $\rho$, a $E'$ jej przestrzenią dualną topologiczną.
Wówczas $\rho=\tau(E,E')$, to znaczy pierwotna topologia
jest topologią Mackeya odpowiadającą parze dualnej $(E,E')$.
\end{twierdzenie}
\begin{proof}
Topologia $\rho$ jest zgodna z $(E,E')$, więc z
twierdzenia Mackeya--Arensa
(twierdzenie Mackeya--Arensa) mamy
$\rho\subseteq\tau(E,E')$.
Niech $K\subseteq E'$ będzie absolutnie wypukłym zbiorem
zwartym w $\sigma(E',E)$. Dla każdego $x\in E$
zbiór $\{f(x):f\in K\}$ jest zwarty w $\K$, a zatem
ograniczony. Twierdzenie~\ref{thm:barrelled-char} daje
równociągłość $K$. Istnieje więc otoczenie $U$ zera,
na którym $p_K(x)=\sup_{f\in K}|f(x)|\le1$.
To dowodzi ciągłości $p_K$ względem $\rho$.
Wszystkie seminormy generujące $\tau(E,E')$ są więc
$\rho$-ciągłe, skąd $\tau(E,E')\subseteq\rho$.
\end{proof}

\begin{wniosek}[Topologia Mackeya przestrzeni Banacha]\label{cor:banach-mackey}
Niech $E$ będzie przestrzenią Banacha, a $E'$ jej
przestrzenią dualną topologiczną. Wówczas
$\tau(E,E')$ jest topologią normową na $E$.
\end{wniosek}
\begin{proof}
Przestrzeń Banacha jest beczkowa, zatem stosujemy
Twierdzenie~\ref{thm:barrelled-mackey}.
\end{proof}

\begin{przyklad}[Aproksymacja dystrybucji Diraca]\label{ex:delta-approx-barrel}
Niech $E=C^\infty([0,1])$ z seminormami
$p_m(\varphi)=\max_{0\le j\le m}\sup_{x\in[0,1]}|\varphi^{(j)}(x)|$,
$m\in\N_0$. Jest to przestrzeń Frécheta, a więc beczkowa.
Dla $n\ge1$ określmy
$T_n(\varphi)=n\int_{[0,1/n]}\varphi(x)\,\mathrm d\lambda(x)$.
Każdy $T_n$ jest ciągłym funkcjonałem liniowym i
$|T_n(\varphi)|\le p_0(\varphi)$, zatem rodzina
$\{T_n:n\ge1\}$ jest równociągła. Ponieważ
$|T_n(\varphi)-\varphi(0)|\le
\sup_{x\in[0,1/n]}|\varphi(x)-\varphi(0)|\to0$,
mamy $T_n(\varphi)\to\varphi(0)$. W późniejszej teorii
dystrybucji funkcjonał $\delta_0(\varphi)=\varphi(0)$
będzie przykładem dystrybucji Diraca; granica zachodzi
punktowo na $E$, a dzięki
Twierdzeniu~\ref{thm:equicont-compact-conv} również
jednostajnie na każdym zwartym zbiorze funkcji testowych
w topologii $E$.
\end{przyklad}

\begin{uwaga}[Granica zakresu obecnego rozdziału]
Nie należy utożsamiać zwartości i ograniczoności w
$C^\infty([0,1])$ ani przenosić automatycznie jej własności
na przestrzeń $\mathcal D(\Omega)$ z topologią LF.
Precyzyjne omówienie zbiorów ograniczonych w przestrzeniach
funkcji testowych nastąpi po konstrukcji granic induktywnych.
\end{uwaga}


% Wersja źródłowa: audyt II, jednostka 11
\setcounter{section}{10}
\section{Przestrzenie Frécheta. Twierdzenia o odwzorowaniu otwartym i wykresie domkniętym}\label{sec:frechet-open-graph}

\subsection{Dlaczego zupełność i metryzowalność występują razem?}

W \hyperref[sec:beczkowe-bornologiczne]{rozdziale 10} udowodniliśmy twierdzenie
Baire'a dla zupełnych przestrzeni metrycznych i wynikającą z niego
beczkowość przestrzeni Frécheta. W tym rozdziale wyprowadzimy dwa dalsze
fundamentalne rezultaty: ciągły liniowy operator surjektywny między
przestrzeniami Frécheta jest otwarty, a operator liniowy określony na
całej przestrzeni Frécheta jest ciągły, jeśli jego wykres jest domknięty
w iloczynie dwóch przestrzeni Frécheta.

W całym rozdziale $\K\in\{\R,\C\}$. Wszystkie przestrzenie
liniowo-topologiczne są Hausdorffa, chyba że podano inaczej.
Nie utożsamiamy zupełności topologicznej z zupełnością sekwencyjną
bez wcześniejszego sprawdzenia metryzowalności.

\begin{definicja}[Przestrzeń Frécheta]\label{def:frechet11}
Przestrzeń liniowo-topologiczną $E$ nad $\K$ nazywamy
\emph{przestrzenią Frécheta}, jeżeli $E$ jest lokalnie wypukła,
Hausdorffa, metryzowalna i zupełna względem swej kanonicznej
struktury jednostajnej. Oznaczenie $\mathcal L(E,F)$ oznacza
przestrzeń ciągłych operatorów liniowych $E\to F$.
\end{definicja}

\begin{stwierdzenie}[Metryka wyznaczona przez seminormy]\label{prop:frechet-metric11}
Niech $E$ będzie lokalnie wypukłą TVS Hausdorffa, której topologię
wyznacza przeliczalna rodzina ciągłych seminorm $(p_n)_{n\ge1}$.
Wtedy wzór
\[
 d(x,y)=\sum_{n=1}^{\infty}2^{-n}\min\{1,p_n(x-y)\},\qquad x,y\in E,
\]
określa metrykę niezmienniczą względem przesunięć, generującą
początkową topologię $E$. Metryka $d$ jest zupełna wtedy i tylko
wtedy, gdy $E$ jest zupełna jako TVS.
\end{stwierdzenie}
\begin{proof}
Każdy składnik jest nieujemny i niezmienniczy względem przesunięć.
Z nierówności $\min\{1,a+b\}\le \min\{1,a\}+\min\{1,b\}$
wynika nierówność trójkąta. Ponieważ seminormy rozdzielają punkty,
$d(x,y)=0$ implikuje $x=y$.

Jeśli $x_i\to x$ we wszystkich seminormach, to dla dowolnego
$\varepsilon>0$ wybieramy $N$ z $\sum_{n>N}2^{-n}<\varepsilon/2$.
Następnie pierwsze $N$ składników sumy definiującej $d(x_i,x)$
są dostatecznie małe, więc $d(x_i,x)<\varepsilon$.
Odwrotnie, z $d(x_i,x)\to0$ wynika dla ustalonego $n$,
że $\min\{1,p_n(x_i-x)\}\to0$, a więc $p_n(x_i-x)\to0$.
To dowodzi zgodności topologii, również dla sieci.

Z tego samego oszacowania wynika równoważność warunku Cauchy'ego
w metryce $d$ i warunku Cauchy'ego dla każdej seminormy $p_n$.
Metryka $d$ wyznacza więc kanoniczną strukturę jednostajną TVS,
a jej zupełność jest równoważna zupełności przestrzeni $E$.
\end{proof}

\begin{przyklad}[Podstawowe przestrzenie Frécheta]\label{ex:frechet-list11}
Każda przestrzeń Banacha jest przestrzenią Frécheta. Dalsze
przykłady to $\K^{\N}$ z topologią produktową oraz
$C^\infty([0,1])$ z seminormami
$p_n(f)=\max_{0\le j\le n}\sup_{t\in[0,1]}|f^{(j)}(t)|$,
gdzie pochodne na $(0,1)$ mają ciągłe przedłużenia na $[0,1]$.
W przestrzeni $\K^{\N}$ można wziąć seminormy
$p_n(x)=\max_{1\le j\le n}|x_j|$. Każda współrzędna ciągu
Cauchy'ego jest wtedy ciągiem Cauchy'ego w $\K$, a granice
współrzędnych dają granicę w topologii produktowej.
Przestrzeń $C^\infty([0,1])$ jest zupełna, ponieważ zbieżność
jednostajna każdej pochodnej na $[0,1]$ i twierdzenie podstawowe
rachunku całkowego pozwalają wykazać, że granica pochodnych
rzędu $j+1$ jest pochodną granicy rzędu $j$.
\end{przyklad}

\subsection{Podprzestrzenie, iloczyny i przestrzenie ilorazowe}

\begin{twierdzenie}[Stabilność klasy przestrzeni Frécheta]\label{thm:frechet-stability11}
Niech $E$ będzie przestrzenią Frécheta, a $M\subseteq E$
domkniętą podprzestrzenią liniową. Wtedy $M$ oraz $E/M$
ze swoimi naturalnymi topologiami są przestrzeniami Frécheta.
Ponadto przeliczalny produkt przestrzeni Frécheta z topologią
produktową jest przestrzenią Frécheta.
\end{twierdzenie}
\begin{proof}
Niech $d$ będzie zgodną, zupełną i niezmienniczą względem
przesunięć metryką na $E$, jak w Stwierdzeniu~\ref{prop:frechet-metric11}.
Domknięta podprzestrzeń $M$ jest zupełna względem metryki
indukowanej z $d$, a lokalna wypukłość i metryzowalność są
dziedziczne dla podprzestrzeni.

Niech $q:E\to E/M$ będzie rzutowaniem ilorazowym.
Z domkniętości $M$ wynika, że
\[
 \bar d\bigl(q(x),q(y)\bigr)=\inf_{m\in M}d(x-y+m,0)
\]
jest metryką: nie zależy od przedstawicieli klas, spełnia
nierówność trójkąta, a $\bar d(q(x),q(y))=0$ oznacza
$x-y\in\overline M=M$. Ponadto
$q(B_d(x,r))=B_{\bar d}(q(x),r)$, więc $\bar d$ wyznacza
topologię ilorazową (rzutowanie $q$ jest otwarte).

Wykażemy zupełność. Niech $(u_n)$ będzie ciągiem Cauchy'ego
w $E/M$. Wybierzmy rosnący podciąg indeksów $(n_k)$ tak,
aby $\bar d(u_{n_{k+1}},u_{n_k})<2^{-k-1}$.
Wybierzmy przedstawiciela $z_1\in E$ klasy $u_{n_1}$.
Indukcyjnie wybieramy przedstawiciela $z_{k+1}$ klasy
$u_{n_{k+1}}$ z $d(z_{k+1},z_k)<2^{-k}$;
istnienie wynika z definicji $\bar d$. Ciąg $(z_k)$
jest Cauchy'ego w $E$, a więc zbiega do pewnego $z\in E$.
W konsekwencji $u_{n_k}=q(z_k)\to q(z)$.
Ponieważ całe $(u_n)$ jest ciągiem Cauchy'ego,
również $u_n\to q(z)$. Zatem $E/M$ jest zupełna.
Jej lokalna wypukłość wynika z obrazów otwartych,
wypukłych otoczeń zera przez $q$.

Niech teraz $(E_j)_{j\ge1}$ będzie ciągiem przestrzeni Frécheta.
Wybierzmy w każdym $E_j$ zgodną zupełną metrykę
$d_j$, niezmienniczą względem przesunięć. Metryka
\[
 D(x,y)=\sum_{j=1}^\infty 2^{-j}\min\{1,d_j(x_j,y_j)\}
\]
wyznacza topologię produktową na $\prod_{j\ge1}E_j$.
Ciąg $D$-Cauchy'ego jest Cauchy'ego na każdej współrzędnej;
po przejściu do granic współrzędnych i oszacowaniu ogona szeregu
otrzymujemy zbieżność względem $D$. Lokalna wypukłość
wynika z lokalnej wypukłości czynników.
\end{proof}

\begin{uwaga}[Znaczenie domkniętości podprzestrzeni]
Jeśli $M$ nie jest domknięta, iloraz $E/M$ z topologią
ilorazową nie jest Hausdorffa. Dlatego nie jest przestrzenią
Frécheta. Warunek domkniętości w Twierdzeniu~\ref{thm:frechet-stability11}
nie może zostać pominięty.
\end{uwaga}

\subsection{Twierdzenie o odwzorowaniu otwartym}

\begin{lemat}[Krok Baire'a]\label{lem:baire-open11}
Niech $E$ będzie TVS, $F$ lokalnie wypukłą TVS Hausdorffa
będącą przestrzenią Baire'a, a $T:E\to F$ ciągłym,
surjektywnym operatorem liniowym. Jeśli $U\subseteq E$
jest otwartym, absolutnie wypukłym otoczeniem zera, to
$\overline{T(U)}$ jest otoczeniem zera w $F$.
\end{lemat}
\begin{proof}
Ponieważ $U$ jest pochłaniający, a $T$ jest surjektywny,
$F=\bigcup_{n=1}^\infty nT(U)$. Zbiór $nT(U)$ jest
wypukły i zrównoważony, a jego domknięcie również.
Z własności Baire'a wynika, że pewien zbiór
$\overline{nT(U)}=n\overline{T(U)}$ ma niepuste wnętrze.
Stąd także $\overline{T(U)}$ ma niepuste wnętrze.
Jeżeli wypukły, zrównoważony zbiór $C$ ma niepuste wnętrze,
to zawiera otoczenie zera: dla $a\in\operatorname{int}C$
mamy także $-a\in\operatorname{int}C$, a wypukłość
daje otoczenie zera zawarte w $C$, jako zbiór
$\tfrac12((a+V)+(-a+V))=V$ dla dostatecznie małego,
otwartego, zrównoważonego otoczenia $V$ zera.
Stosujemy to do $C=\overline{T(U)}$.
\end{proof}

\begin{twierdzenie}[Twierdzenie o odwzorowaniu otwartym]\label{thm:open-mapping11}
Niech $E$ i $F$ będą przestrzeniami Frécheta nad $\K$.
Jeżeli $T:E\to F$ jest ciągłym i surjektywnym
operatorem liniowym, to $T$ jest odwzorowaniem otwartym:
dla każdego otwartego $O\subseteq E$ zbiór $T(O)$
jest otwarty w $F$.
\end{twierdzenie}
\begin{proof}
Wystarczy pokazać, że obraz dowolnego otoczenia zera
w $E$ zawiera otoczenie zera w $F$. Ustalmy takie
otoczenie $U$ i zmniejszmy je, jeśli trzeba, aby było
otwarte i absolutnie wypukłe. Wybierzmy otwarte,
absolutnie wypukłe otoczenie $U_0$ zera takie, że
$\overline{U_0}\subseteq U$. Istnienie wynika z lokalnej
wypukłości i ciągłości seminorm wyznaczających topologię.

Indukcyjnie konstruujemy otwarte, absolutnie wypukłe
otoczenia $(U_n)_{n\ge1}$ zera w $E$, takie że
$U_n+U_n\subseteq U_{n-1}$ oraz $U_n\subseteq A_n$,
gdzie $(A_n)$ jest ustaloną malejącą przeliczalną bazą
otoczeń zera w $E$. Możemy to zrobić, zmniejszając
kolejne otoczenia. Wówczas $(U_n)$ jest bazą otoczeń zera.

Z Lematu~\ref{lem:baire-open11}, zastosowanego osobno
do każdego $U_n$, wynika, że
$\overline{T(U_n)}$ zawiera otoczenie zera w $F$.
Wybieramy zatem otwarte, absolutnie wypukłe otoczenia
$V_n$ zera w $F$ takie, że
$V_n\subseteq\overline{T(U_n)}\cap B_n$, gdzie $(B_n)$
jest malejącą przeliczalną bazą otoczeń zera w $F$.

Niech $y\in V_1$. Wobec $y\in\overline{T(U_1)}$
i otwartości $V_2$ istnieje $x_1\in U_1$ takie, że
$y-Tx_1\in V_2$. Ponieważ $V_2\subseteq\overline{T(U_2)}$,
możemy wybrać $x_2\in U_2$ tak, aby
$y-T(x_1+x_2)\in V_3$. Kontynuując konstrukcję,
otrzymujemy $x_n\in U_n$ spełniające
\[
 r_n:=y-T\Bigl(\sum_{k=1}^n x_k\Bigr)\in V_{n+1}
 \quad (n\ge1).
\]
Ponieważ $V_{n+1}\subseteq B_{n+1}$, mamy $r_n\to0$ w $F$.

Pokażemy, że szereg $\sum_{n=1}^\infty x_n$ zbiega w $E$.
Z inkluzji $U_n+U_n\subseteq U_{n-1}$ wynika także
$U_n\subseteq U_{n-1}$. Indukcja wsteczna dowodzi, że
$\sum_{k=m}^N x_k\in U_{m-1}$ dla dowolnych $N\ge m\ge1$.
Ponieważ $(U_n)$ jest bazą otoczeń zera, ciąg sum częściowych
jest Cauchy'ego. Zupełność $E$ daje
$x=\lim_{N\to\infty}\sum_{k=1}^N x_k$.
Dla $m=1$ sumy częściowe należą do $U_0$, więc
$x\in\overline{U_0}\subseteq U$. Z ciągłości $T$ i
$r_n\to0$ otrzymujemy $Tx=y$. Udowodniliśmy
$V_1\subseteq T(U)$.

Dla dowolnego zbioru otwartego $O\subseteq E$ i punktu
$y=Tx\in T(O)$ wybieramy otoczenie zera $U_x$ z
$x+U_x\subseteq O$. Wówczas $T(O)$ zawiera otoczenie
$y$ postaci $y+V$, gdzie $V\subseteq T(U_x)$.
Zatem $T(O)$ jest otwarty.
\end{proof}

\begin{uwaga}[Co jest rzeczywiście potrzebne?]\label{rem:open-assumptions11}
Dowód wykorzystuje zupełność i metryzowalność dziedziny,
własność Baire'a oraz metryzowalność przestrzeni wartości,
a także lokalną wypukłość obu przestrzeni. W konsekwencji
ten sam argument działa dla $E$ będącej przestrzenią Frécheta
oraz $F$ będącej metryzowalną lokalnie wypukłą TVS Hausdorffa
o własności Baire'a, nawet jeśli zupełności $F$ nie założono.
Nie twierdzimy natomiast, że dowolny ciągły operator
liniowy między dwiema przestrzeniami Frécheta jest otwarty:
konieczna jest surjektywność.
\end{uwaga}

\begin{wniosek}[Twierdzenie o ciągłości odwrotności]\label{cor:inverse-mapping11}
Niech $E,F$ będą przestrzeniami Frécheta, a
$T\in\mathcal L(E,F)$ liniową bijekcją.
Wtedy $T^{-1}:F\to E$ jest ciągłym operatorem liniowym.
Zatem $T$ jest izomorfizmem liniowo-topologicznym.
\end{wniosek}
\begin{proof}
Z Twierdzenia~\ref{thm:open-mapping11} operator $T$ jest otwarty.
Dla każdego otwartego $U\subseteq E$ zbiór
$(T^{-1})^{-1}(U)=T(U)$ jest zatem otwarty w $F$.
\end{proof}

\begin{wniosek}[Oszacowanie normy operatora odwrotnego]\label{cor:bounded-inverse11}
Niech $E,F$ będą przestrzeniami Banacha i niech
$T:E\to F$ będzie liniowym, ciągłym i bijektywnym
operatorem. Istnieje stała $C>0$ taka, że
$\|T^{-1}y\|_E\le C\|y\|_F$ dla każdego $y\in F$.
Równoważnie, $\|Tx\|_F\ge C^{-1}\|x\|_E$ dla $x\in E$.
\end{wniosek}
\begin{proof}
Z Wniosku~\ref{cor:inverse-mapping11} odwrotność jest ciągła.
Liniowy operator między przestrzeniami unormowanymi jest
ciągły wtedy i tylko wtedy, gdy jest ograniczony; przyjmujemy
$C=\max\{1,\|T^{-1}\|\}$.
\end{proof}

\begin{wniosek}[Izomorfizm z przestrzenią ilorazową]\label{cor:quotient-open11}
Niech $E,F$ będą przestrzeniami Frécheta oraz
$T\in\mathcal L(E,F)$ operatorem surjektywnym.
Wtedy $M=\ker T$ jest domkniętą podprzestrzenią $E$,
a operator
$\widetilde T:E/M\to F$, $\widetilde T(x+M)=Tx$,
jest izomorfizmem liniowo-topologicznym.
\end{wniosek}
\begin{proof}
Ciągłość $T$ i domkniętość $\{0\}\subset F$ dają
domkniętość $M$. Własność uniwersalna ilorazu zapewnia
ciągłość $\widetilde T$. Jest ona liniową bijekcją.
Ponieważ $E/M$ i $F$ są przestrzeniami Frécheta,
Wniosek~\ref{cor:inverse-mapping11} daje ciągłość odwrotności.
\end{proof}

\subsection{Twierdzenie o wykresie domkniętym}

\begin{definicja}[Wykres operatora]\label{def:graph11}
Jeżeli $E,F$ są TVS, a $T:D(T)\subseteq E\to F$
odwzorowaniem, to jego \emph{wykresem} nazywamy
$G(T)=\{(x,Tx):x\in D(T)\}\subseteq E\times F$.
Dla operatora liniowego $G(T)$ jest podprzestrzenią
liniową produktu.
\end{definicja}

\begin{stwierdzenie}[Kryterium sekwencyjne domknięcia wykresu]
\label{prop:graph-sequential11}
Niech $E,F$ będą metryzowalnymi TVS Hausdorffa, a
$T:E\to F$ operatorem liniowym określonym na całym $E$.
Wykres $G(T)$ jest domknięty w $E\times F$ wtedy i tylko
wtedy, gdy dla każdego ciągu $(x_n)$ w $E$ i każdego
$x\in E$, $y\in F$ spełniających $x_n\to x$ i
$Tx_n\to y$ zachodzi $Tx=y$.
\end{stwierdzenie}
\begin{proof}
Produkt $E\times F$ jest metryzowalny. W przestrzeni
metrycznej zbiór jest domknięty dokładnie wtedy, gdy
zawiera granicę każdego zbieżnego ciągu swoich punktów.
Zastosowanie do ciągu $(x_n,Tx_n)$ daje tezę.
\end{proof}

\begin{twierdzenie}[Twierdzenie o wykresie domkniętym]\label{thm:closed-graph11}
Niech $E,F$ będą przestrzeniami Frécheta, a
$T:E\to F$ operatorem liniowym określonym na całym $E$.
Wówczas $T$ jest ciągły wtedy i tylko wtedy,
gdy jego wykres $G(T)$ jest domknięty w $E\times F$.
\end{twierdzenie}
\begin{proof}
Jeżeli $T$ jest ciągły, to odwzorowanie
$(x,y)\mapsto y-Tx$ jest ciągłe z $E\times F$ do $F$.
Jego przeciwobraz zbioru domkniętego $\{0\}$ to $G(T)$,
a więc wykres jest domknięty.

Odwrotnie, załóżmy, że $G(T)$ jest domknięty.
Iloczyn $E\times F$ jest przestrzenią Frécheta
na mocy Twierdzenia~\ref{thm:frechet-stability11},
a jego domknięta podprzestrzeń $G(T)$ również.
Rzutowanie
$P:G(T)\to E$, $P(x,Tx)=x$, jest ciągłą liniową
bijekcją między przestrzeniami Frécheta. Z Wniosku
\ref{cor:inverse-mapping11} wynika ciągłość $P^{-1}$.
Stąd $T=\pi_F\circ P^{-1}$ jest ciągły, gdzie
$\pi_F:E\times F\to F$ oznacza drugie rzutowanie.
\end{proof}

\begin{center}
\begin{tikzpicture}[>=Stealth, every node/.style={font=\small}]
  \node (g) at (0,1.65) {$G(T)\subseteq E\times F$};
  \node (e) at (-2.75,0) {$E$};
  \node (f) at (2.75,0) {$F$};
  \draw[->] (g) -- node[above left] {$P$} (e);
  \draw[->] (g) -- node[above right] {$\pi_F|_{G(T)}$} (f);
  \draw[->,dashed] (e) -- node[below] {$T=\pi_F\circ P^{-1}$} (f);
\end{tikzpicture}

\smallskip
\small Schemat dowodu twierdzenia o wykresie domkniętym:
$P$ jest ciągłą bijekcją, a jego odwrotność jest ciągła
na mocy twierdzenia o odwzorowaniu otwartym.
\end{center}

\begin{uwaga}[Pełna dziedzina operatora jest niezbędna]
Twierdzenie~\ref{thm:closed-graph11} wymaga, by $T$ był
określony na całym $E$. Operatory różniczkowe jako
operatory nieograniczone w przestrzeniach Banacha mają
zwykle właściwe dziedziny $D(T)\subsetneq E$ i mogą
mieć domknięte wykresy bez ciągłości w normie $E$.
Przykład podamy poniżej.
\end{uwaga}

\subsection{Kontrprzykłady i zastosowania}

\begin{przyklad}[Brak zupełności przeciwdziedziny: bijekcja nieotwarta]
\label{ex:incomplete-codomain11}
Na zbiorze $\ell^1$ rozważmy dwie normy
$\|x\|_1=\sum_{n\ge1}|x_n|$ oraz
$\|x\|_2=(\sum_{n\ge1}|x_n|^2)^{1/2}$.
Operator identycznościowy
$J:(\ell^1,\|\cdot\|_1)\to(\ell^1,\|\cdot\|_2)$
jest ciągłą liniową bijekcją, bo $\|x\|_2\le\|x\|_1$.
Jego odwrotność nie jest ciągła. Dla wektora $u^{(n)}$
o pierwszych $n$ współrzędnych równych $1/n$
i pozostałych równych zeru mamy
$\|u^{(n)}\|_1=1$ i $\|u^{(n)}\|_2=n^{-1/2}\to0$.
Zatem $J$ nie jest otwarty. Przeciwdziedzina
$(\ell^1,\|\cdot\|_2)$ nie jest zupełna: jej
uzupełnieniem jest $\ell^2$, ponieważ $c_{00}$ jest
gęste w $\ell^2$.
\end{przyklad}

\begin{przyklad}[Wykres domknięty nie zapewnia ciągłości bez zupełności]
\label{ex:incomplete-domain-graph11}
Niech $E=(\ell^1,\|\cdot\|_2)$, $F=(\ell^1,\|\cdot\|_1)$,
a $T:E\to F$ będzie odwzorowaniem identycznościowym.
Nie jest ono ciągłe, jak pokazuje ciąg $u^{(n)}$ z
Przykładu~\ref{ex:incomplete-codomain11}. Ma jednak
wykres domknięty w $E\times F$. Istotnie, jeśli
$x_n\to x$ w normie $\ell^2$ oraz
$Tx_n=x_n\to y$ w normie $\ell^1$, to $x_n\to y$
także w normie $\ell^2$. Jednoznaczność granicy
w $\ell^2$ daje $x=y$, czyli $(x,y)\in G(T)$.
Kryterium sekwencyjne domknięcia wykresu kończy dowód.
Tutaj $E$ jest niezupełna, mimo że $F$ jest przestrzenią Banacha.
\end{przyklad}

\begin{przyklad}[Domknięty nieograniczony operator różniczkowania]
\label{ex:closed-derivative11}
Niech $E=F=C([0,1])$ z normą supremum i
$D(T)=C^1([0,1])$. Określmy $Tf=f'$. Jest to liniowy
operator nieograniczony względem normy $E$:
$f_n(t)=n^{-1}\sin(nt)$ spełnia $\|f_n\|_\infty\to0$,
podczas gdy $\|Tf_n\|_\infty=1$.
Wykres $G(T)$ jest jednak domknięty w $E\times F$.
Jeżeli $f_n\to f$ i $f_n'\to g$ jednostajnie,
to dla $t\in[0,1]$ z twierdzenia podstawowego rachunku
całkowego mamy
\[
 f_n(t)-f_n(0)=\int_{[0,t]}f_n'(s)\,\mathrm d\lambda(s).
\]
Przechodząc do granicy, otrzymujemy
$f(t)-f(0)=\int_{[0,t]}g(s)\,\mathrm d\lambda(s)$.
Ponieważ $g$ jest ciągła, funkcja $f$ jest klasy $C^1$
i $f'=g$. Nie przeczy to Twierdzeniu~\ref{thm:closed-graph11},
ponieważ $D(T)\ne E$.
\end{przyklad}

\begin{przyklad}[Operator pochodnej w przestrzeni Frécheta]
\label{ex:frechet-derivative11}
Niech $E=F=C^\infty([0,1])$ z seminormami $p_n$
z Przykładu~\ref{ex:frechet-list11} i niech $Df=f'$.
Operator $D:E\to F$ jest ciągły, ponieważ
$p_n(Df)\le p_{n+1}(f)$ dla każdego $n$.
Jest również surjektywny: dla każdego $g\in F$ funkcja
$Rg(t)=\int_{[0,t]}g(s)\,\mathrm d\lambda(s)$ należy do $E$
i spełnia $D(Rg)=g$. Jądro $D$ składa się z funkcji stałych.
Zatem Twierdzenie~\ref{cor:quotient-open11} daje
izomorfizm liniowo-topologiczny
$C^\infty([0,1])/\K\cong C^\infty([0,1])$,
gdzie $\K$ oznacza podprzestrzeń funkcji stałych.
W istocie $R$ jest ciągły, ponieważ
$p_n(Rg)\le 2p_n(g)$ dla wszystkich $n\ge0$.
\end{przyklad}

\begin{przyklad}[Gęsty, niedomknięty obraz operatora]
\label{ex:dense-range11}
Naturalne włożenie $\iota:\ell^1\to\ell^2$ jest
ciągłe i różnowartościowe. Jego obraz jest gęsty, gdyż
zawiera $c_{00}$, ale nie jest cały: ciąg
$(n^{-3/4})_{n\ge1}$ należy do $\ell^2\setminus\ell^1$.
Odwrotność $\iota^{-1}:\iota(\ell^1)\to\ell^1$
nie jest ciągła przy topologii indukowanej z $\ell^2$.
Twierdzenie o odwzorowaniu otwartym nie ma zastosowania
do operatora $\ell^1\to\iota(\ell^1)$, ponieważ
przeciwdziedzina z topologią indukowaną nie jest zupełna.
\end{przyklad}

\begin{wniosek}[Równoważność zupełnych norm porównywalnych]
\label{cor:two-banach-norms11}
Niech $V$ będzie przestrzenią wektorową nad $\K$,
a $\|\cdot\|_a$ oraz $\|\cdot\|_b$ normami na $V$
takimi, że obie przestrzenie $(V,\|\cdot\|_a)$ i
$(V,\|\cdot\|_b)$ są Banacha. Jeśli istnieje $C>0$
z $\|x\|_b\le C\|x\|_a$ dla wszystkich $x\in V$,
to istnieje $D>0$ takie, że
$\|x\|_a\le D\|x\|_b$ dla wszystkich $x\in V$.
\end{wniosek}
\begin{proof}
Operator identycznościowy między podanymi przestrzeniami
Banacha jest ciągłą bijekcją. Jego odwrotność jest ciągła
na mocy Wniosku~\ref{cor:inverse-mapping11}, co jest
równoważne żądanemu oszacowaniu.
\end{proof}

\begin{uwaga}[Pomost do rachunku różniczkowego]
W dalszych rozdziałach będziemy badać pochodne Gâteaux
i Frécheta odwzorowań między przestrzeniami Banacha.
W twierdzeniu o lokalnej funkcji odwrotnej zakłada się
odwracalność pochodnej Frécheta, która jest ograniczonym
operatorem liniowym. Wniosek~\ref{cor:bounded-inverse11}
pokazuje, dlaczego algebraiczna odwracalność takiej
pochodnej między przestrzeniami Banacha zapewnia
ciągłość jej odwrotności. Nie należy mylić tego z
odwracalnością samego odwzorowania nieliniowego.
\end{uwaga}


% Wersja źródłowa: audyt II, jednostka 12
\setcounter{section}{11}
\section{Klasyczne przestrzenie Banacha i ich przestrzenie dualne}
\label{sec:klasyczne-banach12}

\subsection{Przejście od ogólnych przestrzeni liniowo-topologicznych do norm}

Dotychczas prowadziliśmy rozważania dla przestrzeni liniowo-topologicznych,
przestrzeni lokalnie wypukłych i przestrzeni Frécheta. Teraz zawężamy
założenia, aby otrzymać konkretne wzory na normy, ciągłe funkcjonały
i przestrzenie dualne. Będziemy pracować nad $\K\in\{\R,\C\}$.
Niech $E'$ oznacza przestrzeń ciągłych funkcjonałów $\K$-liniowych na $E$.
Jeżeli $E$ jest unormowana, oznaczamy przez $\|\cdot\|_{E'}$ normę
operatorową na $E'$. Dla $1<p<\infty$ liczbę $q=p/(p-1)$ nazywamy
\emph{wykładnikiem sprzężonym} do $p$; dla $p=1$ przyjmujemy $q=\infty$.

\begin{definicja}[Przestrzeń unormowana i przestrzeń Banacha]
\label{def:banach12}
Przestrzenią unormowaną nazywamy parę $(E,\|\cdot\|)$, gdzie $E$
jest przestrzenią wektorową nad $\K$, a norma spełnia
$\|x\|=0\iff x=0$, $\|\alpha x\|=|\alpha|\|x\|$ oraz
$\|x+y\|\le\|x\|+\|y\|$. Przestrzeń unormowaną nazywamy
\emph{przestrzenią Banacha}, jeżeli jest zupełna względem metryki
$d(x,y)=\|x-y\|$.
\end{definicja}

\begin{stwierdzenie}[Norma operatora i przestrzeń dualna]
\label{prop:operator-norm12}
Niech $E$ będzie przestrzenią unormowaną i $F$ przestrzenią Banacha.
Wtedy przestrzeń $\mathcal L(E,F)$ wszystkich ciągłych operatorów
liniowych $E\to F$, z normą
$\|T\|=\sup_{\|x\|_E\le1}\|Tx\|_F$, jest przestrzenią Banacha.
W szczególności $E'=\mathcal L(E,\K)$ jest zawsze przestrzenią
Banacha, nawet gdy $E$ nie jest zupełna.
\end{stwierdzenie}
\begin{proof}
Ciągłość liniowego $T$ w zerze jest równoważna istnieniu $C\ge0$
takiego, że $\|Tx\|_F\le C\|x\|_E$: dla implikacji nietrywialnej
wystarczy przeskalować niezerowe $x$ do kuli, na której $T$ jest
ograniczony. Norma operatorowa jest zatem skończona i spełnia
aksjomaty normy.

Jeśli $(T_n)$ jest ciągiem Cauchy'ego w normie operatorowej,
to dla każdego $x\in E$ ciąg $(T_nx)$ jest Cauchy'ego w $F$, ponieważ
$\|T_nx-T_mx\|\le\|T_n-T_m\|\|x\|$.
Niech $Tx=\lim_nT_nx$. Granice zachowują liniowość. Ustalmy
$\varepsilon>0$ i wybierzmy $N$ tak, że
$\|T_n-T_m\|<\varepsilon$ dla $m,n\ge N$.
Po przejściu do granicy $m\to\infty$ otrzymujemy
$\|(T_n-T)x\|\le\varepsilon\|x\|$ dla $n\ge N$.
Stąd $T$ jest ciągły i $\|T_n-T\|\to0$.
\end{proof}

\subsection{Przestrzenie ciągowe}

\begin{definicja}[Przestrzenie $\ell^p$, $c_0$ i $\ell^\infty$]
\label{def:sequence-spaces12}
Dla $1\le p<\infty$ określamy
$\ell^p=\{x=(x_n)_{n\ge1}\in\K^{\N}:\sum_{n\ge1}|x_n|^p<\infty\}$
i $\|x\|_p=(\sum_{n\ge1}|x_n|^p)^{1/p}$.
Ponadto $\ell^\infty$ jest przestrzenią ciągów ograniczonych
z normą $\|x\|_\infty=\sup_{n\ge1}|x_n|$, a $c_0$ oznacza
podprzestrzeń ciągów zbieżnych do zera, z tą samą normą.
Przez $c_{00}$ oznaczamy przestrzeń ciągów o skończonym nośniku.
\end{definicja}

\begin{lemat}[Nierówność H\"oldera dla ciągów]
\label{lem:holder-seq12}
Dla $1<p<\infty$, $q=p/(p-1)$, $x\in\ell^p$ i $y\in\ell^q$
zachodzi
$\sum_{n\ge1}|x_ny_n|\le\|x\|_p\|y\|_q$.
Dla $p=1$, $q=\infty$, ta sama nierówność jest prawdziwa.
\end{lemat}
\begin{proof}
Dla $a,b\ge0$ nierówność Younga
$ab\le a^p/p+b^q/q$ wynika z wypukłości funkcji
$t\mapsto e^t$ (równoważnie, z minimalizacji
$t^p/p-tb$ względem $t\ge0$). Jeżeli obie normy są niezerowe,
stosujemy ją do $a=|x_n|/\|x\|_p$ i
$b=|y_n|/\|y\|_q$, sumujemy po $n\le N$,
a następnie przechodzimy do granicy $N\to\infty$.
Przypadki z normą zerową są oczywiste. Dla $p=1$
oszacowujemy bezpośrednio każdy wyraz przez
$|x_n|\|y\|_\infty$.
\end{proof}

\begin{twierdzenie}[Zupełność klasycznych przestrzeni ciągowych]
\label{thm:seq-banach12}
Dla każdego $1\le p\le\infty$ przestrzeń $\ell^p$ jest Banacha.
Przestrzeń $c_0$ jest domkniętą podprzestrzenią $\ell^\infty$,
a zatem także przestrzenią Banacha. Ponadto $c_{00}$ jest gęste
w $\ell^p$ dla $1\le p<\infty$ i w $c_0$, lecz nie jest gęste
w $\ell^\infty$.
\end{twierdzenie}
\begin{proof}
Dla $p=\infty$ ciąg Cauchy'ego $(x^{(k)})$ jest Cauchy'ego
jednostajnie względem współrzędnych. Istnieją granice
$x_n=\lim_kx_n^{(k)}$, a przejście do granicy w oszacowaniu
$\|x^{(k)}-x^{(m)}\|_\infty\le\varepsilon$ daje
$\|x^{(k)}-x\|_\infty\le\varepsilon$. Ciąg $x$ jest ograniczony.

Nierówność Minkowskiego dla skończonych sum ciągów
wynika z nierówności H\"oldera tak samo jak dla całek:
$\|u+v\|_p^p\le(\|u\|_p+\|v\|_p)\|u+v\|_p^{p-1}$
dla $p>1$, zaś dla $p=1$ jest natychmiastowa.

Dla $1\le p<\infty$ wybierzmy z ciągu Cauchy'ego
podciąg $(x^{(k_j)})$ spełniający
$\|x^{(k_{j+1})}-x^{(k_j)}\|_p\le2^{-j}$.
Niech $h_j=x^{(k_{j+1})}-x^{(k_j)}$.
Nierówność Minkowskiego dla sum skończonych daje
$\|\sum_{j=1}^N|h_j|\|_p\le\sum_{j=1}^N\|h_j\|_p\le1$.
Przez monotoniczność sum współrzędnych i ich granic
otrzymujemy $s=\sum_{j\ge1}|h_j|\in\ell^p$.
Szereg $\sum_jh_j$ jest więc bezwzględnie zbieżny
na każdej współrzędnej i określa $h\in\ell^p$;
ponadto norma jego ogona nie przekracza $\sum_{j\ge N}2^{-j}$.
Wobec tego $x^{(k_j)}\to x^{(k_1)}+h$ w $\ell^p$.
Cały ciąg Cauchy'ego ma tę samą granicę.

Jeśli $x^{(k)}\in c_0$ i $x^{(k)}\to x$ jednostajnie,
to dla $\varepsilon>0$ wybieramy $k$ z
$\|x-x^{(k)}\|_\infty<\varepsilon/2$, a następnie $N$
z $|x_n^{(k)}|<\varepsilon/2$ dla $n\ge N$.
Wtedy $x\in c_0$. Obcięcie ciągu do pierwszych $N$
współrzędnych daje zbieżność w $\ell^p$ dla $p<\infty$
i w $c_0$. W $\ell^\infty$ odległość stałego ciągu
$(1,1,\ldots)$ od $c_{00}$ wynosi co najmniej $1$.
\end{proof}

\begin{twierdzenie}[Przestrzenie dualne $\ell^p$ i $c_0$]
\label{thm:dual-sequences12}
Niech $1\le p<\infty$ i niech $q$ będzie wykładnikiem sprzężonym.
Każdy $F\in(\ell^p)'$ ma dokładnie jedną reprezentację
\[
 F(x)=\sum_{n=1}^\infty a_nx_n,\qquad x\in\ell^p,
 \quad a\in\ell^q,
\]
a odwzorowanie $a\mapsto F$ jest izometrycznym izomorfizmem
$\ell^q\cong(\ell^p)'$ (dla $p=1$ przyjmujemy $q=\infty$).
Ponadto każdy $F\in(c_0)'$ ma dokładnie jedną reprezentację
tego samego typu z $a\in\ell^1$ i $\|F\|=\|a\|_1$.
\end{twierdzenie}
\begin{proof}
Niech $e_n$ będzie $n$-tym wektorem standardowym i niech
$a_n=F(e_n)$. Dla $1<p<\infty$ i $N\ge1$ oznaczmy
$S_N=\sum_{n=1}^N|a_n|^q$. Jeśli $S_N>0$, wybieramy
$x_n=\overline{a_n}|a_n|^{q-2}S_N^{-1/p}$
dla $n\le N$ (a dla $a_n=0$ przyjmujemy $x_n=0$)
i $x_n=0$ dla $n>N$. Wtedy
$\|x\|_p=1$ i $F(x)=S_N^{1/q}$, ponieważ
$(q-1)p=q$. Zatem $S_N^{1/q}\le\|F\|$ dla każdego $N$.
Stąd $a\in\ell^q$ i $\|a\|_q\le\|F\|$.
Dla $p=1$ mamy bezpośrednio
$|a_n|\le\|F\|\|e_n\|_1=\|F\|$.

Ponieważ $c_{00}$ jest gęste w $\ell^p$ dla $p<\infty$,
na mocy ciągłości i Lematu~\ref{lem:holder-seq12}
$F(x)=\sum_na_nx_n$ dla każdego $x\in\ell^p$.
Nierówność H\"oldera daje $\|F\|\le\|a\|_q$,
a powyższe oszacowania dają równość.

Dla $F\in(c_0)'$ zdefiniujmy ponownie $a_n=F(e_n)$.
Dla każdego $N$ wektor $x\in c_{00}$ o współrzędnych
$x_n=\overline{a_n}/|a_n|$ przy $a_n\ne0$ i
$x_n=0$ przy $a_n=0$, $n\le N$, ma normę co najwyżej $1$.
Zatem $\sum_{n=1}^N|a_n|\le\|F\|$.
Wobec tego $a\in\ell^1$, a gęstość $c_{00}$ w $c_0$
oraz nierówność $|\sum_na_nx_n|\le\|a\|_1\|x\|_\infty$
kończą dowód wraz z identycznością norm.
Jednoznaczność w obu przypadkach wynika z podstawienia $x=e_n$.
\end{proof}

\begin{przyklad}[Dlaczego $(\ell^\infty)'$ jest większa niż $\ell^1$?]
\label{ex:linfty-dual-seq12}
Na podprzestrzeni $c=c_0+\K\mathbf1\subseteq\ell^\infty$
ciągów zbieżnych określmy $L(x)=\lim_{n\to\infty}x_n$.
Jest to ciągły funkcjonał liniowy o normie $1$.
Z twierdzenia Hahna--Banacha przedłużamy go do
$\widetilde L\in(\ell^\infty)'$ z tą samą normą.
Mamy $\widetilde L(e_n)=0$ dla każdego $n$, ale
$\widetilde L(\mathbf1)=1$. Gdyby $\widetilde L$
był reprezentowany przez pewne $a\in\ell^1$ wzorem
$\widetilde L(x)=\sum_na_nx_n$, to z podstawienia $e_n$
wynikałoby $a_n=0$ dla wszystkich $n$, co przeczy
$\widetilde L(\mathbf1)=1$.
\end{przyklad}

\subsection{Przestrzenie \texorpdfstring{$L^p$}{Lp} na przestrzeni z miarą skończoną}

W tej podsekcji $(\Omega,\mathcal A,\mu)$ jest przestrzenią
z miarą dodatnią i skończoną, tj. $\mu(\Omega)<\infty$.
Rozważamy funkcje mierzalne o wartościach w $\K$, utożsamiane,
gdy są równe $\mu$-prawie wszędzie. Zapis
$\int_\Omega f(x)\,\mathrm d\mu(x)$ zawsze oznacza całkę
względem wskazanej miary. Jeżeli $\mu(\Omega)=0$, wszystkie
przestrzenie $L^p(\mu)$ są zerowe i twierdzenia pozostają prawdziwe.

\begin{definicja}[Przestrzenie $L^p(\mu)$]
\label{def:lp12}
Dla $1\le p<\infty$ niech $L^p(\mu)$ oznacza przestrzeń
klas funkcji $f$ spełniających
$\int_\Omega|f(x)|^p\,\mathrm d\mu(x)<\infty$, z normą
$\|f\|_p=(\int_\Omega|f(x)|^p\,\mathrm d\mu(x))^{1/p}$.
Przestrzeń $L^\infty(\mu)$ składa się z klas funkcji istotnie
ograniczonych z normą
$\|f\|_\infty=\operatorname*{ess\,sup}_{x\in\Omega}|f(x)|$.
\end{definicja}

\begin{twierdzenie}[H\"older, Minkowski i zupełność $L^p$]
\label{thm:lp-banach12}
Dla $1\le p\le\infty$ przestrzeń $L^p(\mu)$ jest Banacha.
Dla $1<p<\infty$ i $q=p/(p-1)$ zachodzi nierówność H\"oldera
$\int_\Omega|f(x)g(x)|\,\mathrm d\mu(x)\le\|f\|_p\|g\|_q$,
a dla $p=1$ analogiczny wzór z $q=\infty$.
Dla każdego $1\le p\le\infty$ zachodzi nierówność Minkowskiego
$\|f+g\|_p\le\|f\|_p+\|g\|_p$.
\end{twierdzenie}
\begin{proof}
Dowód H\"oldera powtarza argument z nierównością Younga
z Lematu~\ref{lem:holder-seq12}, po zastąpieniu sumy całką.
Dla $1<p<\infty$ nierówność Minkowskiego otrzymujemy z tożsamości
$|f+g|^p\le(|f|+|g|)|f+g|^{p-1}$, nierówności H\"oldera
i podzielenia przez $\|f+g\|_p^{p-1}$, o ile norma ta jest
niezerowa. Dla $p=1$ i $p=\infty$ jest to bezpośrednia
nierówność trójkąta.

Niech $(f_n)$ będzie ciągiem Cauchy'ego w $L^p$ dla
$1\le p<\infty$. Wybierzmy podciąg $(f_{n_k})$ taki, że
$\|f_{n_{k+1}}-f_{n_k}\|_p\le2^{-k}$.
Dla reprezentantów $h_k=f_{n_{k+1}}-f_{n_k}$,
Minkowski i twierdzenie o zbieżności monotonicznej dają
\[
 s(x)=\sum_{k\ge1}|h_k(x)|<\infty\quad\text{dla prawie każdego }x,
 \qquad \|s\|_p\le\sum_{k\ge1}2^{-k}.
\]
Szereg $\sum_kh_k$ zbiega więc prawie wszędzie do mierzalnej
funkcji $h\in L^p$, a normy ogonów są oszacowane przez
$\sum_{k\ge N}2^{-k}$. Stąd $f_{n_k}\to f_{n_1}+h$ w $L^p$.
Cauchy'ego całego ciągu zapewnia tę samą granicę.

Dla $p=\infty$ z warunku Cauchy'ego wybieramy podciąg
$(f_{n_k})$ z $\|f_{n_{k+1}}-f_{n_k}\|_\infty\le2^{-k}$.
Po usunięciu przeliczalnej sumy zbiorów zerowych
reprezentantów mamy punktową jednostajną zbieżność
szeregu różnic do istotnie ograniczonej funkcji.
Reszta dowodu przebiega jak dla $p<\infty$.
\end{proof}

\begin{lemat}[Gęstość funkcji prostych]
\label{lem:simple-dense12}
Dla $1\le p<\infty$ funkcje proste są gęste w $L^p(\mu)$.
Każda funkcja prosta jest kombinacją liniową wskaźników
zbiorów mierzalnych.
\end{lemat}
\begin{proof}
Dla rzeczywistej funkcji nieujemnej wybieramy rosnące funkcje
proste $s_n$ z $0\le s_n\uparrow f$ prawie wszędzie
(np. przez obcinanie wartości do $n$ i zaokrąglanie w dół
do wielokrotności $2^{-n}$). Wtedy
$|f-s_n|^p\le|f|^p$ i $|f-s_n|^p\to0$ prawie wszędzie.
Twierdzenie o zbieżności zmajoryzowanej daje
$\|f-s_n\|_p\to0$. Dla funkcji rzeczywistych ogólnych
i zespolonych rozdzielamy części dodatnie i ujemne
oraz części rzeczywiste i urojone.
\end{proof}

\begin{twierdzenie}[Przestrzeń dualna $L^p$ dla $1\le p<\infty$]
\label{thm:lp-dual12}
Niech $(\Omega,\mathcal A,\mu)$ będzie przestrzenią z miarą
skończoną, niech $1\le p<\infty$ i $q$ będzie wykładnikiem
sprzężonym. Dla każdego $F\in(L^p(\mu))'$ istnieje dokładnie
jedna klasa $g\in L^q(\mu)$ taka, że
\[
 F(f)=\int_\Omega f(x)g(x)\,\mathrm d\mu(x)
 \quad\text{dla wszystkich }f\in L^p(\mu).
\]
Ponadto $\|F\|=\|g\|_q$, więc
$(L^p(\mu))'\cong L^q(\mu)$ izometrycznie. W przypadku
zespolonym jest to identyfikacja liniowa przy przyjętej konwencji
wzoru bez sprzężenia $g$.
\end{twierdzenie}
\begin{proof}
Nierówność H\"oldera pokazuje, że dowolne $g\in L^q$
definiuje ciągły funkcjonał o normie co najwyżej $\|g\|_q$.
Dla danego $F\in(L^p)'$ definiujemy miarę zespoloną
(lub znakowaną w przypadku rzeczywistym)
$\nu(A)=F(\mathbf1_A)$, $A\in\mathcal A$.
Wskaźniki należą do $L^p$, ponieważ $\mu$ jest skończona.
Jeżeli $A_j$ są parami rozłączne, to
$\mathbf1_{\cup_{j=1}^NA_j}\to\mathbf1_{\cup_{j\ge1}A_j}$
w $L^p$ na mocy zbieżności monotonicznej; ciągłość $F$
daje przeliczalną addytywność $\nu$.
Dla dowolnego skończonego rozkładu $\Omega$ na zbiory
mierzalne $A_j$ wybierzmy liczby $\theta_j\in\K$ o module
nie większym niż $1$ tak, aby
$\theta_j\nu(A_j)=|\nu(A_j)|$. Wówczas
\[
 \sum_j|\nu(A_j)|
 =F\Bigl(\sum_j\theta_j\mathbf1_{A_j}\Bigr)
 \le\|F\|\,\mu(\Omega)^{1/p}.
\]
Zatem $\nu$ ma skończoną wariację i jest absolutnie ciągła
względem $\mu$. Z twierdzenia Radona--Nikodyma dla miar
skończonych otrzymujemy $g\in L^1(\mu)$ takie, że
$\nu(A)=\int_Ag(x)\,\mathrm d\mu(x)$ dla $A\in\mathcal A$.
Przez liniowość i Lemat~\ref{lem:simple-dense12} trzeba jeszcze
wykazać, że $g\in L^q$.

Jeżeli $p=1$, to dla dowolnego $A\in\mathcal A$ stosujemy
nierówność normową do funkcji
$f=\mathbf1_A\overline g/|g|$ na $\{g\ne0\}$,
a poza tym zbiorem przyjmujemy $f=0$. Z rozszerzenia wzoru
z funkcji prostych na funkcje ograniczone (przez jednostajną
aproksymację funkcjami prostymi) otrzymujemy
$\int_A|g|\,\mathrm d\mu\le\|F\|\mu(A)$.
Wybór $A=\{|g|>\|F\|+\varepsilon\}$ pokazuje,
że $\|g\|_\infty\le\|F\|$.

Jeżeli $1<p<\infty$, niech
$g_N=g\mathbf1_{\{|g|\le N\}}$ i niech
$f_N=\overline{g_N}|g_N|^{q-2}$, z wartością $0$ w
punktach, gdzie $g_N=0$. Ponieważ $\mu$ jest skończona,
$f_N$ jest ograniczona, a więc należy do $L^p$.
Wzór reprezentacyjny obowiązuje dla niej, gdyż
jest ona granicą jednostajną funkcji prostych; wobec tego
\[
 \int_\Omega|g_N(x)|^q\,\mathrm d\mu(x)
 =F(f_N)\le\|F\|\left(\int_\Omega|g_N(x)|^q\,\mathrm d\mu(x)\right)^{1/p}.
\]
Jeśli całka po lewej stronie jest dodatnia, dzielimy przez
jej potęgę $1/p$ i otrzymujemy $\|g_N\|_q\le\|F\|$.
Zbieżność monotoniczna daje $g\in L^q$ i $\|g\|_q\le\|F\|$.
Dla $p=1$ wykazaliśmy analogiczną nierówność.
Teraz równość $F(f)=\int fg\,\mathrm d\mu$ wynika z gęstości
funkcji prostych i ciągłości obu stron na $L^p$.
Nierówność H\"oldera daje $\|F\|\le\|g\|_q$, czyli
równość norm. Jednoznaczność $g$ wynika z podstawiania
wskaźników zbiorów mierzalnych.
\end{proof}

\begin{uwaga}[Dlaczego ograniczenie $p<\infty$ jest istotne?]
\label{rem:linfty-dual12}
Przy $p=\infty$ przestrzeń dualna $(L^\infty(\mu))'$ jest
na ogół istotnie większa od $L^1(\mu)$. W przypadku miary
Lebesgue'a na $[0,1]$ wybierzmy parami rozłączne zbiory mierzalne
$A_n$ o dodatniej mierze, a $A=\bigcup_nA_n$. Operator
$P:L^\infty([0,1])\to\ell^\infty$ dany wzorem
$P(f)_n=\lambda(A_n)^{-1}\int_{A_n}f(x)\,\mathrm d\lambda(x)$
jest ciągły, liniowy i $\|P\|\le1$.
Jeżeli $\widetilde L$ jest przedłużeniem funkcjonału granicy
z Przykładu~\ref{ex:linfty-dual-seq12}, to $\Phi=\widetilde L\circ P$
spełnia $\Phi(\mathbf1_{A_n})=0$ dla wszystkich $n$, lecz
$\Phi(\mathbf1_A)=1$.
Gdyby $\Phi(f)=\int_{[0,1]}f(x)g(x)\,\mathrm d\lambda(x)$
dla pewnego $g\in L^1$, przeliczalna addytywność całki dałaby
$\Phi(\mathbf1_A)=\sum_n\Phi(\mathbf1_{A_n})=0$, sprzeczność.
\end{uwaga}

\subsection{Przestrzeń \texorpdfstring{$C(K)$}{C(K)} i twierdzenie Riesza--Markowa}

\begin{definicja}[Przestrzeń funkcji ciągłych i miary regularne]
\label{def:ck-measures12}
Niech $K$ będzie niepustą zwartą przestrzenią metryczną.
Przez $C(K;\K)$ oznaczamy przestrzeń ciągłych funkcji
$K\to\K$ z normą $\|f\|_\infty=\max_{x\in K}|f(x)|$.
Skończoną miarę borelowską $\nu$ na $K$ nazywamy
\emph{regularną}, jeśli dla każdego borelowskiego $B\subseteq K$
i $\varepsilon>0$ istnieje zbiór domknięty $F$ i zbiór otwarty
$U$ takie, że $F\subseteq B\subseteq U$ oraz
$|\nu|(U\setminus F)<\varepsilon$, gdzie $|\nu|$ oznacza
miarę wariacji całkowitej. Zbiór takich miar znakowanych
lub zespolonych oznaczamy przez $\mathcal M(K;\K)$.
\end{definicja}

\begin{stwierdzenie}[Zupełność $C(K)$]
\label{prop:ck-banach12}
Dla każdej niepustej zwartej przestrzeni metrycznej $K$
przestrzeń $C(K;\K)$ jest przestrzenią Banacha.
\end{stwierdzenie}
\begin{proof}
Ciąg Cauchy'ego w normie supremum jest jednostajnie Cauchy'ego;
jego granica punktowa jest jednostajna, ponieważ $\K$ jest zupełne.
Jednostajna granica funkcji ciągłych jest ciągła, więc należy do $C(K)$.
\end{proof}

\begin{lemat}[Reprezentacja dodatniego funkcjonału]
\label{lem:riesz-positive12}
Niech $K$ będzie niepustą zwartą przestrzenią metryczną,
a $T:C(K;\R)\to\R$ dodatnim funkcjonałem liniowym,
tj. $T(f)\ge0$ dla każdej $f\ge0$.
Wówczas istnieje dokładnie jedna skończona regularna dodatnia miara
borelowska $\mu_T$ na $K$, dla której
$T(f)=\int_Kf(x)\,\mathrm d\mu_T(x)$ dla wszystkich $f\in C(K;\R)$.
\end{lemat}
\begin{proof}
Dodatniość daje $|T(f)|\le T(\mathbf1)\|f\|_\infty$,
więc $T$ jest ciągły. Dla otwartego $U\subseteq K$ ustalmy
\[
 m(U)=\sup\{T(f):f\in C(K;\R),\ 0\le f\le1,
 \ \operatorname{supp}f\subseteq U\}.
\]
Mamy $m(K)=T(\mathbf1)$, $m(\varnothing)=0$, a $m$
jest monotoniczna. Dla każdej rodziny otwartych
$U_j$ zachodzi $m(\bigcup_jU_j)\le\sum_jm(U_j)$:
jeśli $f$ występuje w definicji lewej strony,
to zwarty nośnik $f$ ma skończone podpokrycie przez $U_j$.
Wybieramy ciągłe funkcje $\psi_j\ge0$, o nośnikach zawartych
w odpowiednich $U_j$, których suma jest równa $1$
na zwartym nośniku $f$. Można je skonstruować z funkcji
odległości od dopełnień zbiorów tworzących skończone
podpokrycie i znormalizować ich sumę na otoczeniu nośnika. Piszemy
$f=\sum_jf\psi_j$, przy czym wszystkie niezerowe składniki
mają nośniki zawarte w odpowiednich $U_j$ i należą do $[0,1]$.
Jeżeli $U,V$ są rozłącznymi zbiorami otwartymi, ten sam
argument, zastosowany do podziału na $U$ i $V$, wraz z
monotonicznością daje $m(U\cup V)=m(U)+m(V)$.

Niech dla dowolnego $A\subseteq K$
$\mu^*(A)=\inf\{m(U):A\subseteq U,\ U\text{ otwarty}\}$.
Własności $m$ dowodzą, że $\mu^*$ jest miarą zewnętrzną.
Ponadto jest miarą zewnętrzną metryczną: jeśli
$\operatorname{dist}(A,B)>0$, rozdzielamy $A,B$ rozłącznymi
otwartymi otoczeniami wewnątrz dowolnego otwartego
zbioru pokrywającego $A\cup B$ i korzystamy z addytywności
$m$ na rozłącznych zbiorach otwartych. Otrzymujemy
$\mu^*(A\cup B)=\mu^*(A)+\mu^*(B)$.
Z twierdzenia Carath\'eodory'ego o metrycznych miarach
zewnętrznych wynika, że wszystkie zbiory borelowskie są
$\mu^*$-mierzalne. Ograniczenie $\mu_T=\mu^*|_{\mathcal B(K)}$
jest miarą skończoną i z definicji zewnętrznie regularną.
Jest także wewnętrznie regularna: dla borelowskiego $B$
aproksymujemy $K\setminus B$ z zewnątrz zbiorem otwartym
$V$, wtedy $K\setminus V$ jest domknięty i zawiera się w $B$,
a $\mu_T(B\setminus(K\setminus V))$ jest dowolnie mała.

Pokażemy reprezentację. Ustalmy $f\in C(K;\R)$, $f\ge0$.
Dla $\delta>0$ i $N\delta\ge\|f\|_\infty$ definiujemy
$h_j(x)=\min\{1,\max\{0,(f(x)-j\delta)/\delta\}\}$,
$j=0,\dots,N-1$. Wtedy punktowo
$f=\delta\sum_{j=0}^{N-1}h_j$. Z definicji $m$
i z przybliżeń funkcji $h_j$ funkcjami o nośniku
zawartym w $\{f>j\delta\}$ wynika
\[
 \mu_T(\{f>(j+1)\delta\})\le T(h_j)
 \le\mu_T(\{f>j\delta\}).
\]
Istotnie, dolne oszacowanie dostajemy z $h_j=1$
na $\{f\ge(j+1)\delta\}$, a górne, przybliżając $h_j$
jednostajnie funkcjami zerującymi się na
$\{f\le j\delta+\varepsilon\}$ i korzystając z
ciągłości $T$. Po zsumowaniu i pomnożeniu przez $\delta$
otrzymujemy dolną i górną sumę dla całki
$\int_0^{\|f\|_\infty}\mu_T(\{f>t\})\,\mathrm dt$;
różnią się one o co najwyżej $\delta\mu_T(K)$.
Warstwowy wzór całkowania daje
$T(f)=\int_Kf\,\mathrm d\mu_T$. Liniowość rozszerza
tę równość na wszystkie funkcje rzeczywiste.

Jeśli dwie regularne miary reprezentują $T$, to ich całki
przeciwko wszystkim funkcjom ciągłym są równe.
Dla zbioru otwartego $U$ regularność i konstrukcja ciągłych
funkcji $0\le f\le1$ o nośniku w $U$ dają
$\mu(U)=\sup\{\int_Kf\,\mathrm d\mu:0\le f\le1,
\operatorname{supp}f\subseteq U\}$.
Miary są zatem równe na zbiorach otwartych,
a zewnętrzna regularność daje ich równość na zbiorach borelowskich.
\end{proof}

\begin{twierdzenie}[Twierdzenie Riesza--Markowa na zwartej przestrzeni metrycznej]
\label{thm:riesz-markov12}
Niech $K$ będzie niepustą zwartą przestrzenią metryczną.
Dla każdego $F\in C(K;\K)'$ istnieje dokładnie jedna
regularna skończona miara borelowska $\nu\in\mathcal M(K;\K)$
taka, że
\[
 F(f)=\int_K f(x)\,\mathrm d\nu(x)
 \quad\text{dla wszystkich }f\in C(K;\K).
\]
Odwzorowanie $\nu\mapsto F$ jest izometrycznym izomorfizmem
$\mathcal M(K;\K)\cong C(K;\K)'$, przy normie
$\|\nu\|=|\nu|(K)$.
\end{twierdzenie}
\begin{proof}
Najpierw $\K=\R$. Dla $f\ge0$ definiujemy
$F^+(f)=\sup\{F(g):g\in C(K;\R),\ 0\le g\le f\}$.
Wartość supremum należy do $[0,\|F\|\|f\|_\infty]$.
Funkcja $F^+$ jest addytywna na stożku dodatnim.
Jedna nierówność wynika z wyboru $g_1\le f_1$ i
$g_2\le f_2$. Dla przeciwnej, jeśli
$0\le h\le f_1+f_2$, rozkładamy
$h=\min\{h,f_1\}+(h-\min\{h,f_1\})$;
oba składniki są ciągłe, nieujemne i ograniczone odpowiednio
przez $f_1$ i $f_2$. Po przejściu do supremum mamy
$F^+(f_1+f_2)=F^+(f_1)+F^+(f_2)$.
Rozszerzamy zatem $F^+$ na rzeczywisty funkcjonał liniowy
przez $F^+(f)=F^+(f_+)-F^+(f_-)$, gdzie
$f_+=\max\{f,0\}$, $f_-=\max\{-f,0\}$.
Jest on dodatni i ciągły; również $F^-=F^+-F$ jest dodatni,
ponieważ $F^+(f)\ge F(f)$ dla $f\ge0$.
Lemat~\ref{lem:riesz-positive12} daje dodatnie regularne miary
$\nu^+$ i $\nu^-$ reprezentujące odpowiednio $F^+$ i $F^-$.
Zatem $F(f)=\int_Kf\,\mathrm d(\nu^+-\nu^-)$.

Dla $\K=\C$ bierzemy części rzeczywistą i urojoną
ograniczenia $F$ do rzeczywistych funkcji ciągłych,
reprezentujemy je dwiema miarami rzeczywistymi i składamy
$\nu=\nu_1+i\nu_2$. Zespolona liniowość $F$ zapewnia,
że ten sam wzór działa na wszystkich $f\in C(K;\C)$.
Jednoznaczność wynika z jednoznaczności reprezentacji
części rzeczywistych i urojonych.

Dla dowolnej miary o skończonej wariacji mamy
$|\int f\,\mathrm d\nu|\le\|f\|_\infty|\nu|(K)$,
czyli $\|F\|\le|\nu|(K)$. W drugą stronę niech
$u=d\nu/d|\nu|$ będzie pochodną Radona--Nikodyma;
$|u|=1$ $|\nu|$-prawie wszędzie. Regularność miary
$|\nu|$ i własność aproksymacji wskaźników zbiorów
borelowskich funkcjami ciągłymi (przez zbiory domknięte
i otwarte oraz funkcje Urysohna) dają gęstość $C(K;\K)$
w $L^1(|\nu|)$. Aproksymując $\overline u$ funkcjami
ciągłymi o module nie większym niż $1$, otrzymujemy
$\|F\|\ge\lim|\int f_nu\,\mathrm d|\nu||=|\nu|(K)$.
To dowodzi izometrii.
\end{proof}

\begin{przyklad}[Miary Diraca i funkcjonały ewaluacji]
\label{ex:dirac12}
Dla $x_0\in K$ miara Diraca $\delta_{x_0}$ jest zdefiniowana
przez $\delta_{x_0}(B)=1$, gdy $x_0\in B$, i $0$ w przeciwnym
przypadku. Funkcjonał $F_{x_0}(f)=f(x_0)$ spełnia
$F_{x_0}(f)=\int_Kf\,\mathrm d\delta_{x_0}$ i ma normę $1$.
Jeśli $x_0\ne x_1$, funkcjonały ewaluacji są różne,
ponieważ metryka $K$ pozwala skonstruować funkcję ciągłą
rozdzielającą $x_0$ i $x_1$.
\end{przyklad}

\begin{uwaga}[Przygotowanie do przestrzeni dystrybucji]
\label{rem:bridge-distribution12}
Miara $\nu$ działa liniowo na funkcjach testowych przez całkowanie,
a miara Diraca działa przez ewaluację. Te same mechanizmy wystąpią
w przestrzeni dystrybucji $\mathcal D'(\Omega)$, jednak tam
przestrzeń funkcji testowych ma topologię granicy induktywnej,
a nie normę supremum. W szczególności nie wolno bez dodatkowych
argumentów utożsamiać wszystkich dystrybucji z miarami.
\end{uwaga}


% Wersja źródłowa: audyt II, jednostka 13
\setcounter{section}{12}
\section{Przestrzenie Hilberta: geometria, reprezentacja Riesza i formy półtoraliniowe}
\label{sec:hilbert13}

\subsection{Iloczyn skalarny i zupełność}

\begin{definicja}[Iloczyn skalarny i przestrzeń Hilberta]
\label{def:iloczyn-skalarny13}
Niech $H$ będzie przestrzenią wektorową nad $\K\in\{\R,\C\}$.
\emph{Iloczynem skalarnym} na $H$ nazywamy odwzorowanie
$\langle\cdot,\cdot\rangle:H\times H\to\K$ spełniające, dla
$x,y,z\in H$ i $\alpha,\beta\in\K$:
$\langle\alpha x+\beta y,z\rangle=\alpha\langle x,z\rangle+
\beta\langle y,z\rangle$,
$\langle x,y\rangle=\overline{\langle y,x\rangle}$ oraz
$\langle x,x\rangle>0$ dla $x\ne0$.
\emph{W całym rozdziale iloczyn skalarny jest liniowy w pierwszym
argumencie} i antyliniowy w drugim. Przestrzenią Hilberta nazywamy
przestrzeń z iloczynem skalarnym zupełną w normie
$\|x\|=\sqrt{\langle x,x\rangle}$.
\end{definicja}

\begin{lemat}[Nierówność Cauchy'ego--Schwarza]
\label{lem:cs13}
Dla dowolnych $x,y$ w przestrzeni z iloczynem skalarnym zachodzi
$|\langle x,y\rangle|\le\|x\|\|y\|$.
Równość dla $y\ne0$ zachodzi dokładnie wtedy, gdy $x$ jest
wielokrotnością $y$.
\end{lemat}
\begin{proof}
Dla $y\ne0$ przyjmijmy $\alpha=\langle x,y\rangle/\|y\|^2$.
Wówczas $\langle x-\alpha y,y\rangle=0$ oraz
$0\le\|x-\alpha y\|^2=\|x\|^2-|\langle x,y\rangle|^2/\|y\|^2$.
Równość oznacza $x=\alpha y$. Dla $y=0$ teza jest oczywista.
Nierówność trójkąta dla normy wynika teraz z rozwinięcia
$\|x+y\|^2$ i nierówności Cauchy'ego--Schwarza.
\end{proof}

\begin{stwierdzenie}[Tożsamość równoległoboku i polaryzacja]
\label{prop:parallelogram13}
W każdej przestrzeni z iloczynem skalarnym zachodzi
$\|x+y\|^2+\|x-y\|^2=2\|x\|^2+2\|y\|^2$.
Dla $\K=\C$ iloczyn skalarny jest odtworzony z normy wzorem
\[
\langle x,y\rangle=\frac14\bigl(\|x+y\|^2-\|x-y\|^2
+i\|x+iy\|^2-i\|x-iy\|^2\bigr).
\]
Dla $\K=\R$ wystarczą pierwsze dwa wyrazy.
\end{stwierdzenie}
\begin{proof}
Rozwijamy kwadraty norm. Przy naszej konwencji
$\|x+iy\|^2-\|x-iy\|^2=4\operatorname{Im}\langle x,y\rangle$;
różnica dla $x\pm y$ to $4\operatorname{Re}\langle x,y\rangle$.
\end{proof}

\begin{przyklad}[Przestrzenie $\ell^2$ i $L^2$]
\label{ex:hilbert-l2-13}
Na $\ell^2$ przyjmujemy $\langle x,y\rangle=\sum_{n\ge1}x_n\overline{y_n}$.
Na $L^2(\Omega,\mu)$ przyjmujemy
$\langle f,g\rangle=\int_\Omega f(x)\overline{g(x)}\,\mathrm d\mu(x)$.
Nierówność Cauchy'ego--Schwarza zapewnia zbieżność i ciągłość tych
wyrażeń, zaś zupełność wynika z
twierdzeń o zupełności przestrzeni $\ell^2$ i $L^2$
(drugi przykład przy założeniach o mierze z \hyperref[sec:klasyczne-banach12]{rozdziału~12}).
Przestrzeń $C([0,1])$ z normą supremum nie jest przestrzenią Hilberta w tej normie:
dla $f(t)=1$ i $g(t)=t$ mamy
$\|f+g\|_\infty^2+\|f-g\|_\infty^2=5$,
podczas gdy $2\|f\|_\infty^2+2\|g\|_\infty^2=4$.
\end{przyklad}

\subsection{Rzut na zbiór wypukły i ortogonalność}

\begin{twierdzenie}[Twierdzenie o rzucie metrycznym]
\label{thm:projection-convex13}
Niech $H$ będzie rzeczywistą lub zespoloną przestrzenią Hilberta,
a $C\subseteq H$ niepustym, domkniętym i wypukłym zbiorem.
Dla każdego $x\in H$ istnieje dokładnie jeden punkt $p=P_Cx\in C$
spełniający $\|x-p\|=\inf_{z\in C}\|x-z\|$.
Punkt $p\in C$ jest tym minimizatorem wtedy i tylko wtedy, gdy
\[
\operatorname{Re}\langle x-p,z-p\rangle\le0\qquad(z\in C).
\]
Ponadto odwzorowanie $P_C:H\to C$ jest $1$-lipschitzowskie.
\end{twierdzenie}
\begin{proof}
Niech $d=\inf_{z\in C}\|x-z\|$ i wybierzmy $z_n\in C$ z
$\|x-z_n\|^2\to d^2$. Ponieważ $(z_n+z_m)/2\in C$,
tożsamość równoległoboku zastosowana do $x-z_n$ i $x-z_m$ daje
\[
\|z_n-z_m\|^2
=2\|x-z_n\|^2+2\|x-z_m\|^2
-4\left\|x-\frac{z_n+z_m}{2}\right\|^2
\le2\|x-z_n\|^2+2\|x-z_m\|^2-4d^2\longrightarrow0.
\]
Zatem $(z_n)$ jest ciągiem Cauchy'ego. Zupełność $H$ i domkniętość
$C$ dają granicę $p\in C$, a ciągłość normy daje $\|x-p\|=d$.
Jeśli $p,q$ są minimizatorami, ten sam wzór daje $p=q$.

Niech $z\in C$ i $t\in[0,1]$. Minimalność $p$ oznacza
$\|x-p\|^2\le\|x-(p+t(z-p))\|^2$,
a więc
$0\le-2t\operatorname{Re}\langle x-p,z-p\rangle+t^2\|z-p\|^2$.
Dzieląc przez $t>0$ i przechodząc z $t$ do zera otrzymujemy
warunek wariacyjny. Odwrotnie, warunek ten i rozwinięcie
$\|x-z\|^2=\|x-p\|^2+\|z-p\|^2-
2\operatorname{Re}\langle x-p,z-p\rangle$ implikują minimalność.

Dla $p=P_Cx$, $q=P_Cy$ zastosujmy warunek kolejno z $z=q$
i $z=p$. Otrzymujemy
$\operatorname{Re}\langle x-p,q-p\rangle\le0$ oraz
$\operatorname{Re}\langle y-q,p-q\rangle\le0$. Po dodaniu
$\|p-q\|^2\le\operatorname{Re}\langle x-y,p-q\rangle
\le\|x-y\|\|p-q\|$, skąd $\|p-q\|\le\|x-y\|$.
\end{proof}

\begin{definicja}[Dopełnienie ortogonalne]
\label{def:orth-complement13}
Dla $M\subseteq H$ definiujemy
$M^\perp=\{x\in H:\langle x,m\rangle=0\text{ dla wszystkich }m\in M\}$.
Jest to domknięta podprzestrzeń liniowa $H$; istotnie, to
przecięcie jąder ciągłych funkcjonałów $x\mapsto\langle x,m\rangle$.
\end{definicja}

\begin{twierdzenie}[Rozkład ortogonalny]
\label{thm:orth-decomposition13}
Niech $H$ będzie przestrzenią Hilberta, a $M\subseteq H$ domkniętą
podprzestrzenią liniową. Wówczas dla każdego $x\in H$ istnieje
jedyny rozkład $x=m+n$, $m\in M$, $n\in M^\perp$.
Odwzorowanie $P_Mx=m$ jest liniowym operatorem o normie co najwyżej
$1$, spełnia $P_M^2=P_M=P_M^*$ oraz
$\|x\|^2=\|P_Mx\|^2+\|x-P_Mx\|^2$.
Ponadto $M^{\perp\perp}=M$.
Jeśli $M\ne\{0\}$, to $\|P_M\|=1$.
\end{twierdzenie}
\begin{proof}
Z Twierdzenia~\ref{thm:projection-convex13} istnieje $m=P_Mx$.
Warunek wariacyjny zastosowany do $m+tv$ dla każdego $v\in M$
i $t\in\R$, a dla przestrzeni zespolonej także do $m+itv$,
daje $\langle x-m,v\rangle=0$. Stąd $n=x-m\in M^\perp$.
Jeśli $x=m_1+n_1=m_2+n_2$, to $m_1-m_2=n_2-n_1$ należy
do $M\cap M^\perp=\{0\}$. Jednoznaczność rozkładu zapewnia
liniowość $P_M$; twierdzenie Pitagorasa daje jego ograniczoność.

Dla $x=m+n$ i $y=u+v$ mamy
$\langle P_Mx,y\rangle=\langle m,u\rangle
=\langle x,P_My\rangle$; to oznacza $P_M=P_M^*$,
przy czym istnienie operatora sprzężonego wykażemy formalnie
w Twierdzeniu~\ref{thm:adjoint13} poniżej.
Z rozkładu $H=M\oplus M^\perp$ wynika, że
$M^{\perp\perp}=M$. Dla $m\in M\setminus\{0\}$ zachodzi
$P_Mm=m$, zatem norma operatorowa nie może być mniejsza niż $1$.
\end{proof}

\begin{przyklad}[Rzut na funkcje stałe]
\label{ex:projection-mean13}
W przestrzeni $H=L^2([0,1],\lambda)$ niech
$M=\{\text{klasy funkcji stałych}\}$. Wtedy
$P_Mf=(\int_{[0,1]}f(t)\,\mathrm d\lambda(t))\chi_{[0,1]}$.
Rzeczywiście, $\langle f-P_Mf,\chi_{[0,1]}\rangle=0$.
Ogólniej, rzut na linię $\K u$, $u\ne0$, ma wzór
$P_{\K u}x=\langle x,u\rangle u/\|u\|^2$.
\end{przyklad}

\subsection{Twierdzenie Riesza o reprezentacji funkcjonałów}

\begin{twierdzenie}[Twierdzenie Riesza dla przestrzeni Hilberta]
\label{thm:riesz-hilbert13}
Niech $H$ będzie przestrzenią Hilberta nad $\K\in\{\R,\C\}$.
Dla każdego ciągłego $\K$-liniowego funkcjonału $F:H\to\K$
istnieje dokładnie jeden $y\in H$ taki, że
$F(x)=\langle x,y\rangle$ dla wszystkich $x\in H$.
Ponadto $\|F\|_{H'}=\|y\|$.
Odwzorowanie $J:H\to H'$, $J(y)(x)=\langle x,y\rangle$,
jest izometrycznym izomorfizmem \emph{antyliniowym} dla $\K=\C$
i liniowym dla $\K=\R$.
\end{twierdzenie}
\begin{proof}
Jeżeli $F=0$, wybieramy $y=0$. Załóżmy $F\ne0$.
Podprzestrzeń $M=\ker F$ jest domknięta, a $M^\perp\ne\{0\}$:
gdyby $M^\perp=\{0\}$, to z Twierdzenia~\ref{thm:orth-decomposition13}
wynikałoby $M=H$.
Niech $z\in M^\perp$, $z\ne0$. Wtedy $F(z)\ne0$, gdyż
$z\in M\cap M^\perp$ implikowałoby $z=0$.
Dla dowolnego $x\in H$ mamy
$x-F(x)z/F(z)\in M$, więc
\[
0=\left\langle x-\frac{F(x)}{F(z)}z,z\right\rangle
=\langle x,z\rangle-\frac{F(x)}{F(z)}\|z\|^2.
\]
Zatem $F(x)=F(z)\langle x,z\rangle/\|z\|^2$.
Przy naszej konwencji należy wziąć
$y=\overline{F(z)}z/\|z\|^2$, ponieważ
$\langle x,cy\rangle=\overline c\langle x,y\rangle$.
Nierówność Cauchy'ego--Schwarza daje $\|F\|\le\|y\|$,
a podstawienie $x=y/\|y\|$ dla $y\ne0$ daje równość.
Jeśli $\langle x,y_1-y_2\rangle=0$ dla wszystkich $x$, to
$x=y_1-y_2$ zapewnia $y_1=y_2$. Antyliniowość $J$
wynika bezpośrednio z $J(\alpha y+\beta z)
=\overline\alpha J(y)+\overline\beta J(z)$, ponieważ
iloczyn skalarny jest antyliniowy w drugim argumencie.
\end{proof}

\begin{uwaga}[Dwa różne twierdzenia Riesza]
\label{rem:two-riesz13}
Twierdzenie~\ref{thm:riesz-hilbert13} utożsamia przestrzeń dualną
przestrzeni Hilberta z nią samą (antyliniowo w przypadku zespolonym).
Natomiast twierdzenie Riesza--Markowa z teorii przestrzeni $C(K)$ identyfikuje
przestrzeń dualną $C(K)$ z przestrzenią miar. Nie są to
odmienne wersje tej samej reprezentacji w sensie konstrukcyjnym.
\end{uwaga}

\subsection{Rodziny ortonormalne i bazy przestrzeni Hilberta}

\begin{definicja}[Rodzina ortonormalna i baza przestrzeni Hilberta]
\label{def:on-basis13}
Niech $H$ będzie przestrzenią Hilberta, a $I$ zbiorem indeksów.
Rodzina $(e_i)_{i\in I}\subseteq H$ jest \emph{ortonormalna}, jeśli
$\langle e_i,e_j\rangle=\delta_{ij}$.
Nazywamy ją \emph{bazą ortonormalną przestrzeni Hilberta}, jeśli
$\overline{\lin\{e_i:i\in I\}}=H$.
Nie należy utożsamiać bazy ortonormalnej przestrzeni Hilberta z algebraiczną bazą Hamela.
\end{definicja}

\begin{lemat}[Bessel i ortogonalne sumy skończone]
\label{lem:bessel13}
Niech $H$ będzie przestrzenią Hilberta, $(e_i)_{i\in I}$ rodziną
ortonormalną, a $x\in H$. Dla każdego skończonego $A\subseteq I$
zachodzi
\[
\left\|x-\sum_{i\in A}\langle x,e_i\rangle e_i\right\|^2
=\|x\|^2-\sum_{i\in A}|\langle x,e_i\rangle|^2.
\]
W szczególności $\sum_{i\in A}|\langle x,e_i\rangle|^2\le\|x\|^2$
dla każdego takiego $A$ (nierówność Bessela).
\end{lemat}
\begin{proof}
Wektor $x-\sum_{i\in A}\langle x,e_i\rangle e_i$ jest
prostopadły do każdego $e_j$, $j\in A$. Z twierdzenia Pitagorasa
dla jego rozkładu z sumą ortogonalną wynika wzór.
\end{proof}


\begin{twierdzenie}[Przeliczalność nośnika współczynników Fouriera]
\label{thm:countable-fourier-support13}
Niech $H$ będzie przestrzenią Hilberta nad $\K\in\{\R,\C\}$,
$(e_i)_{i\in I}$ dowolną rodziną ortonormalną (zbiór indeksów $I$
może być nieprzeliczalny), a $x\in H$. Wówczas zbiór
$ I_x:=\{i\in I:\langle x,e_i\rangle\ne0\}$
jest co najwyżej przeliczalny. Dokładniej, dla każdego $n\in\N$,
$n\ge1$, zbiór
$ I_{x,n}:=\{i\in I:|\langle x,e_i\rangle|\ge1/n\}$
jest skończony i ma co najwyżej $n^2\|x\|^2$ elementów.
\end{twierdzenie}
\begin{proof}
Niech $F\subseteq I_{x,n}$ będzie zbiorem skończonym.
Z nierówności Bessela (Lemat~\ref{lem:bessel13}) otrzymujemy
$|F|/n^2\le\sum_{i\in F}|\langle x,e_i\rangle|^2\le\|x\|^2$.
Gdyby $I_{x,n}$ miał więcej niż $n^2\|x\|^2$ elementów,
można byłoby wybrać skończony podzbiór o liczności
przekraczającej tę liczbę, co przeczyłoby otrzymanemu oszacowaniu.
Każdy $I_{x,n}$ jest zatem skończony. Dla każdego
$i\in I_x$ liczba $|\langle x,e_i\rangle|$ jest dodatnia, więc
przekracza $1/n$ dla pewnego $n\ge1$. Stąd
$I_x=\bigcup_{n\ge1}I_{x,n}$ jest sumą przeliczalnej rodziny
zbiorów skończonych.
\end{proof}

\begin{uwaga}[Dlaczego warunek jest silniejszy od samego Bessela?]
\label{rem:fourier-countability13}
Nie zakładaliśmy ośrodkowości $H$ ani przeliczalności $I$.
Twierdzenie dotyczy \emph{każdego pojedynczego wektora}:
przy nieprzeliczalnej bazie ortonormalnej różne wektory mogą
mieć różne przeliczalne nośniki współczynników. Nie wynika
stąd, że całą bazę można ponumerować jednym ciągiem.
\end{uwaga}

\begin{twierdzenie}[Rozwinięcie Fouriera, tożsamość Parsevala i zbieżność bezwarunkowa]
\label{thm:parseval13}
\label{cor:fourier-enumeration13}
Niech $H$ będzie przestrzenią Hilberta nad $\K\in\{\R,\C\}$,
$(e_i)_{i\in I}$ rodziną ortonormalną, a
$M=\overline{\lin\{e_i:i\in I\}}$. Dla każdego $x\in H$
sieć sum $s_A(x)=\sum_{i\in A}\langle x,e_i\rangle e_i$,
indeksowana skończonymi podzbiorami $A\subseteq I$ uporządkowanymi
przez inkluzję, zbiega w normie do $P_Mx$. Ponadto
\[
 \|P_Mx\|^2=\sum_{i\in I}|\langle x,e_i\rangle|^2,
 \qquad
 \|x\|^2=\|x-P_Mx\|^2+
               \sum_{i\in I}|\langle x,e_i\rangle|^2,
\]
gdzie suma po dowolnym zbiorze indeksów oznacza supremum
sum skończonych. Rodzina $(e_i)_{i\in I}$ jest bazą ortonormalną
przestrzeni Hilberta wtedy i tylko wtedy, gdy dla każdego $x\in H$
zachodzi tożsamość Parsevala $\|x\|^2=\sum_{i\in I}
|\langle x,e_i\rangle|^2$; równoważnie, gdy
$M^\perp=\{0\}$.

Jeśli $(e_i)_{i\in I}$ jest bazą ortonormalną, to zbiór
$I_x=\{i:\langle x,e_i\rangle\ne0\}$ jest co najwyżej
przeliczalny. Gdy $I_x$ jest nieskończony, dla każdej bijekcji
$\nu:\N\to I_x$ zachodzi
\[
 x=\sum_{n=1}^{\infty}\langle x,e_{\nu(n)}\rangle e_{\nu(n)}
 \quad\text{w normie $H$},\qquad
 \|x\|^2=\sum_{n=1}^{\infty}|\langle x,e_{\nu(n)}\rangle|^2.
\]
Suma nie zależy od numeracji $I_x$ (zbieżność bezwarunkowa).
Jeśli $I_x$ jest skończony, obowiązuje analogiczna suma skończona.
Dla każdego skończonego $A\subseteq I$ mamy dokładny wzór na błąd
\[
 \left\|x-\sum_{i\in A}\langle x,e_i\rangle e_i\right\|^2
 =\sum_{i\in I\setminus A}|\langle x,e_i\rangle|^2.
\]
\end{twierdzenie}
\begin{proof}
Dla skończonych $A\subseteq B\subseteq I$ ortonormalność daje
$\|s_B-s_A\|^2=\sum_{i\in B\setminus A}|\langle x,e_i\rangle|^2$.
Niech $S=\sup_A\sum_{i\in A}|\langle x,e_i\rangle|^2
\le\|x\|^2$ (nierówność Bessela). Dla $\varepsilon>0$ wybieramy
$A_0$ o sumie większej niż $S-\varepsilon^2$.
Dla $A,B\supseteq A_0$ porównanie z $A\cup B$ daje
$\|s_A-s_B\|<2\varepsilon$, zatem $(s_A)$ jest siecią
Cauchy'ego. Z zupełności $H$ wynika jej zbieżność do pewnego
$s\in M$. Dla każdego $j\in I$ od chwili, gdy $j\in A$,
$\langle x-s_A,e_j\rangle=0$. Po przejściu do granicy
$x-s\in M^\perp$, więc $s=P_Mx$. Przejście do granicy
w $\|s_A\|^2=\sum_{i\in A}|\langle x,e_i\rangle|^2$
oraz twierdzenie Pitagorasa dają obie równości.
Charakterystyka bazy wynika z równoważności $M^\perp=\{0\}$
i $M=H$.

Przeliczalność $I_x$ została już niezależnie udowodniona
w Twierdzeniu~\ref{thm:countable-fourier-support13}.
Jeśli $I_x$ jest nieskończony, to zbiory
$A_N=\{\nu(1),\ldots,\nu(N)\}$ są współkońcowe w rodzinie
skończonych podzbiorów $I_x$: każdy skończony
$A\subseteq I_x$ zawiera się w pewnym $A_N$. Współczynniki
spoza $I_x$ są zerowe, więc sieć sum po podzbiorach $I$ i
sieć sum po podzbiorach $I_x$ mają tę samą granicę.
Stąd $s_{A_N}\to x$ (dla bazy $M=H$), niezależnie od $\nu$.
Suma kwadratów wynika z Parsevala. Wreszcie wektor
$x-s_A(x)$ jest prostopadły do $s_A(x)$, więc
$\|x-s_A(x)\|^2=\|x\|^2-\sum_{i\in A}|\langle x,e_i\rangle|^2$;
ponowne użycie Parsevala daje wzór na błąd.
\end{proof}

\begin{przyklad}[Różnica między rozpięciem algebraicznym i domkniętym]
\label{ex:hilbert-series-not-hamel13}
Niech $(e_n)_{n\ge1}$ będzie ciągiem ortonormalnym w $H$.
Szereg $x=\sum_{n\ge1}2^{-n}e_n$ zbiega w $H$, ponieważ
$\sum_{n\ge1}4^{-n}<\infty$. Dla każdego $n$ mamy
$\langle x,e_n\rangle=2^{-n}\ne0$, więc
$x\notin\lin\{e_n:n\ge1\}$, choć należy do domknięcia
tej podprzestrzeni. Baza ortonormalna przestrzeni Hilberta pozwala zatem na rozwinięcia
\emph{normowo zbieżne}, a baza Hamela dopuszcza wyłącznie
\emph{skończone} kombinacje liniowe.
\end{przyklad}

\begin{twierdzenie}[Istnienie baz ortonormalnych]
\label{thm:hilbert-basis13}
Każda przestrzeń Hilberta ma bazę ortonormalną (dopuszczając pustą
bazę dla $H=\{0\}$). Przestrzeń Hilberta jest ośrodkowa wtedy i tylko
wtedy, gdy posiada bazę ortonormalną co najwyżej przeliczalną.
\end{twierdzenie}
\begin{proof}
W zbiorze rodzin ortonormalnych, uporządkowanym przez rozszerzenie,
każdy łańcuch ma ograniczenie górne dane sumą rodzin.
Lemat Zorna zapewnia rodzinę maksymalną $(e_i)_{i\in I}$.
Jeśli $x\perp e_i$ dla każdego $i$ i $x\ne0$, można ją powiększyć
o $x/\|x\|$, sprzeczność. Zatem jest bazą.

Jeśli $H$ ma przeliczalną bazę, skończone kombinacje liniowe z
współczynnikami w $\mathbb Q$ (dla $\R$) lub
$\mathbb Q+i\mathbb Q$ (dla $\C$) tworzą zbiór przeliczalny i gęsty.
Odwrotnie, jeśli $H$ jest ośrodkowa, nie może zawierać
nieprzeliczalnej rodziny ortonormalnej: kule
$B(e_i,1/3)$ są rozłączne, a każdy taki otwarty zbiór zawiera
punkt ustalonego przeliczalnego zbioru gęstego.
\end{proof}


\begin{twierdzenie}[Nieskończenie wymiarowe przestrzenie Hilberta a wymiar Hamela]
\label{thm:hilbert-hamel-uncountable13}
Niech $H$ będzie nieskończenie wymiarową przestrzenią Hilberta
nad $\K\in\{\R,\C\}$. Wtedy każda baza Hamela przestrzeni $H$
jest nieprzeliczalna. Co więcej, $H$ zawiera liniowo niezależną
rodzinę mocy continuum, zatem
$\dim_{\mathrm{Hamel}}H\ge 2^{\aleph_0}$.
W szczególności twierdzenie dotyczy również ośrodkowych,
nieskończenie wymiarowych przestrzeni Hilberta, które mają
przeliczalne bazy ortonormalne.
\end{twierdzenie}
\begin{proof}
Wybierzmy nieskończony ciąg ortonormalny $(e_n)_{n\ge1}$
w $H$; istnieje on, ponieważ $H$ jest nieskończenie wymiarowa
(ewentualnie wybieramy go z bazy ortonormalnej otrzymanej
w Twierdzeniu~\ref{thm:hilbert-basis13}). Dla $t\in(0,1)$
określmy
\[
 v_t:=\sum_{n=1}^{\infty}t^{n-1}e_n.
\]
Szereg ten zbiega w normie dzięki twierdzeniu Pitagorasa i
$\sum_{n\ge1}t^{2(n-1)}=(1-t^2)^{-1}<\infty$.
Pokażemy, że rodzina $(v_t)_{t\in(0,1)}$ jest liniowo niezależna.
Niech $t_1,\ldots,t_m\in(0,1)$ będą parami różne, a
$\sum_{j=1}^m a_jv_{t_j}=0$. Biorąc iloczyn skalarny tej
równości z $e_n$ dla $n=1,\ldots,m$, dostajemy
$\sum_{j=1}^m a_jt_j^{n-1}=0$. Współczynniki $a_j$
spełniają jednorodny układ o macierzy Vandermonde'a
$(t_j^{n-1})_{1\le n,j\le m}$. Jej wyznacznik
$\prod_{1\le j<k\le m}(t_k-t_j)$ jest niezerowy, więc
$a_1=\cdots=a_m=0$. Rodzina $(v_t)_{t\in(0,1)}$ ma moc
$2^{\aleph_0}$ i jest liniowo niezależna. Każda baza Hamela
musi mieć moc co najmniej tak dużą.
\end{proof}

\begin{uwaga}[Drugi argument: twierdzenie Baire'a]
\label{rem:hilbert-hamel-baire13}
Sama nieprzeliczalność bazy Hamela wynika nawet dla każdej
nieskończenie wymiarowej przestrzeni \emph{Banacha} $X$.
Gdyby $X$ miała przeliczalną bazę Hamela $(h_n)_{n\ge1}$,
to $X=\bigcup_{n\ge1}F_n$, gdzie
$F_n=\lin\{h_1,\ldots,h_n\}$. Każda skończenie wymiarowa
podprzestrzeń $F_n$ jest domknięta (w przestrzeni unormowanej)
i ma puste wnętrze: podprzestrzeń liniowa o niepustym wnętrzu
jest całą przestrzenią, a $X$ nie jest skończenie wymiarowa.
Twierdzenie Baire'a wyklucza pokrycie zupełnej przestrzeni
metrycznej przeliczalną sumą takich zbiorów.
Argument ten dowodzi nieprzeliczalności, natomiast poprzednia
konstrukcja daje również dolne oszacowanie mocy bazy Hamela.
\end{uwaga}

\begin{uwaga}[Dwie wielkości, których nie należy mylić]
\label{rem:hamel-vs-hilbert13}
\emph{Wymiar hilbertowski} to moc bazy ortonormalnej,
a \emph{wymiar Hamela} to moc bazy algebraicznej.
Dla $H=\ell^2(\N)$ pierwszy jest przeliczalny, drugi
co najmniej continuum. Jeśli baza ortonormalna jest
nieprzeliczalna, każdy $x\in H$ nadal ma najwyżej
przeliczalnie wiele niezerowych współczynników.
\end{uwaga}

\begin{stwierdzenie}[Izometryczne współrzędne w $\ell^2(I)$]
\label{prop:hilbert-l2I13}
Niech $(e_i)_{i\in I}$ będzie bazą ortonormalną przestrzeni
Hilberta $H$. Wówczas $U:H\to\ell^2(I)$,
$Ux=(\langle x,e_i\rangle)_{i\in I}$, jest izometrycznym
izomorfizmem liniowym, gdzie
$\ell^2(I)=\{a:I\to\K:\sup_{A\subseteq I\text{ skończony}}
\sum_{i\in A}|a_i|^2<\infty\}$.
\end{stwierdzenie}
\begin{proof}
Liniowość i izometria wynikają z Parsevala. Dla dowolnego
$a\in\ell^2(I)$ sieć $\sum_{i\in A}a_ie_i$ jest Cauchy'ego:
tak jak w dowodzie Twierdzenia~\ref{thm:parseval13},
odrzucenie dostatecznie dużego skończonego zbioru współrzędnych
daje dowolnie mały ogon sumy kwadratów. Niech $x$ będzie granicą;
wtedy $\langle x,e_j\rangle=a_j$ dla każdego $j\in I$.
To dowodzi surjektywności.
\end{proof}

\begin{przyklad}[Gram--Schmidt i szeregi Fouriera]
\label{ex:gram-schmidt13}
Dla liniowo niezależnych $v_1,\ldots,v_n\in H$ konstrukcja
$u_k=v_k-\sum_{j<k}\langle v_k,e_j\rangle e_j$,
$e_k=u_k/\|u_k\|$, tworzy rodzinę ortonormalną o tym samym
rozpięciu liniowym. Dla nieskończonego ciągu liniowo niezależnego
ortogonalizujemy kolejne wyrazy i otrzymujemy bazę ortonormalną domkniętego rozpięcia tych wektorów. Dla $H=L^2([0,2\pi],\lambda)$
układ $e_n(t)=(2\pi)^{-1/2}e^{int}$, $n\in\mathbb Z$, jest
ortonormalny; jego zupełność wymaga dodatkowego twierdzenia
o gęstości wielomianów trygonometrycznych, które zostanie
omówione przy transformacji Fouriera. Nie używamy tej zupełności
bez dowodu.
\end{przyklad}

\subsection{Operatory sprzężone i reprezentacja form}

\begin{twierdzenie}[Operator sprzężony w przestrzeni Hilberta]
\label{thm:adjoint13}
Niech $H,G$ będą przestrzeniami Hilberta nad tym samym ciałem $\K$,
a $T\in\mathcal L(H,G)$. Istnieje dokładnie jeden operator
$T^*\in\mathcal L(G,H)$ taki, że
$\langle Tx,y\rangle_G=\langle x,T^*y\rangle_H$
dla wszystkich $x\in H$, $y\in G$.
Ponadto $\|T^*\|=\|T\|$, $(T^*)^*=T$ i
$(ST)^*=T^*S^*$ dla każdego $S\in\mathcal L(G,W)$,
gdzie $W$ jest kolejną przestrzenią Hilberta.
\end{twierdzenie}
\begin{proof}
Dla ustalonego $y\in G$ funkcja $x\mapsto\langle Tx,y\rangle_G$
jest ciągłym liniowym funkcjonałem na $H$. Z
Twierdzenia~\ref{thm:riesz-hilbert13} istnieje jedyny $z_y\in H$
z $\langle Tx,y\rangle_G=\langle x,z_y\rangle_H$.
Definiujemy $T^*y=z_y$. Z antyliniowości w drugim argumencie
po obu stronach, z uwzględnieniem sprzężenia skalarów, wynika
$T^*(\alpha y_1+\beta y_2)=\alpha T^*y_1+\beta T^*y_2$.
Równość norm funkcjonałów reprezentujących daje
$\|T^*y\|\le\|T\|\|y\|$, więc $T^*$ jest ograniczony.
Tożsamość dla $(T^*)^*$ wynika z hermitowskości iloczynu
skalarnego i jednoznaczności operatora sprzężonego.
Stąd $\|T\|=\|(T^*)^*\|\le\|T^*\|\le\|T\|$.
Wreszcie
$\langle STx,w\rangle=\langle Tx,S^*w\rangle
=\langle x,T^*S^*w\rangle$, co daje wzór dla iloczynu.
\end{proof}

\begin{definicja}[Formy półtoraliniowe]
\label{def:sesquilinear13}
Dla przestrzeni Hilberta $H$ nad $\K\in\{\R,\C\}$
\emph{formą półtoraliniową} nazywamy odwzorowanie
$a:H\times H\to\K$ liniowe w pierwszym i antyliniowe
w drugim argumencie. W przypadku rzeczywistym są to formy
dwuliniowe. Dla $\K=\C$ warunek antyliniowości oznacza
$a(x,\alpha y+\beta z)=\overline\alpha a(x,y)+\overline\beta a(x,z)$. Forma jest \emph{ograniczona}, jeśli
$|a(x,y)|\le M\|x\|\|y\|$ dla pewnego $M\ge0$ i wszystkich
$x,y\in H$. Najmniejsze takie $M$ oznaczamy $\|a\|$.
Formę nazywamy \emph{koercywną}, jeśli istnieje $\alpha>0$
spełniające $\operatorname{Re}a(x,x)\ge\alpha\|x\|^2$
dla wszystkich $x\in H$.
\end{definicja}

\begin{twierdzenie}[Reprezentacja ograniczonych form]
\label{thm:form-repr13}
Niech $H$ będzie przestrzenią Hilberta nad $\K$,
a $a:H\times H\to\K$ formą półtoraliniową ograniczoną.
Wówczas istnieje dokładnie jeden $A\in\mathcal L(H)$ taki, że
$a(x,y)=\langle Ax,y\rangle$ dla wszystkich $x,y\in H$.
Ponadto $\|A\|=\|a\|$.
Forma jest hermitowska ($a(x,y)=\overline{a(y,x)}$) wtedy
i tylko wtedy, gdy $A=A^*$.
\end{twierdzenie}
\begin{proof}
Dla ustalonego $x\in H$ funkcja
$y\mapsto\overline{a(x,y)}$ jest ciągłym liniowym funkcjonałem,
zatem istnieje dokładnie jeden $Ax\in H$ spełniający
$\overline{a(x,y)}=\langle y,Ax\rangle$ dla każdego $y$.
Po sprzężeniu otrzymujemy $a(x,y)=\langle Ax,y\rangle$.
Jednoznaczność daje liniowość $x\mapsto Ax$.
Równość $\|Ax\|=\sup_{\|y\|\le1}|a(x,y)|\le\|a\|\|x\|$
dowodzi ograniczoności i $\|A\|\le\|a\|$.
Z drugiej strony
$|a(x,y)|=|\langle Ax,y\rangle|\le\|A\|\|x\|\|y\|$,
stąd równość norm. Równoważność hermitowskości i $A=A^*$
wynika bezpośrednio z definicji operatora sprzężonego.
\end{proof}

\begin{twierdzenie}[Twierdzenie Laxa--Milgrama]
\label{thm:lax-milgram13}
Niech $H$ będzie przestrzenią Hilberta nad $\K\in\{\R,\C\}$,
a $a:H\times H\to\K$ formą półtoraliniową ograniczoną
i koercywną ze stałą $\alpha>0$. Dla każdego ciągłego
\emph{antyliniowego} funkcjonału $\Lambda:H\to\K$ istnieje
dokładnie jeden $u\in H$ taki, że
$a(u,v)=\Lambda(v)$ dla wszystkich $v\in H$.
Ponadto $\|u\|\le\|\Lambda\|/\alpha$, gdzie
$\|\Lambda\|=\sup_{\|v\|\le1}|\Lambda(v)|$.
\end{twierdzenie}
\begin{proof}
Niech $A\in\mathcal L(H)$ reprezentuje $a$ jak w
Twierdzeniu~\ref{thm:form-repr13}. Z koercywności i nierówności
Cauchy'ego--Schwarza mamy
$\alpha\|x\|^2\le\operatorname{Re}\langle Ax,x\rangle
\le\|Ax\|\|x\|$, zatem $\|Ax\|\ge\alpha\|x\|$.
Stąd $A$ jest injektywny i ma domknięty obraz:
$Ax_n\to y$ implikuje, że $x_n$ jest ciągiem Cauchy'ego,
a więc $x_n\to x$ i $Ax=y$.

Jeżeli $y\perp\operatorname{ran}A$, to $A^*y=0$.
Wtedy
$\alpha\|y\|^2\le\operatorname{Re}\langle Ay,y\rangle
=\operatorname{Re}\langle y,A^*y\rangle=0$,
więc $y=0$. Zatem obraz $A$ jest gęsty i jednocześnie
domknięty, czyli $\operatorname{ran}A=H$.
Ponieważ $\Lambda$ jest antyliniowy, funkcja
$v\mapsto\overline{\Lambda(v)}$ jest ciągłym liniowym
funkcjonałem. Twierdzenie~\ref{thm:riesz-hilbert13} daje
jedyny $f\in H$ spełniający $\Lambda(v)=\langle f,v\rangle$.
Bierzemy $u=A^{-1}f$; wtedy $a(u,v)=\langle Au,v\rangle
=\Lambda(v)$. Jednoznaczność wynika z injektywności $A$,
a oszacowanie z nierówności $\|Au\|\ge\alpha\|u\|$
i $\|f\|=\|\Lambda\|$.
\end{proof}

\begin{przyklad}[Jednowymiarowy problem brzegowy]
\label{ex:weak-poisson13}
Niech $H=\{u:[0,1]\to\C:\ u\text{ jest absolutnie ciągła},
u(0)=u(1)=0,\ u'\in L^2([0,1],\lambda)\}$,
przy czym pochodną rozumiemy prawie wszędzie.
Definiujemy $\langle u,v\rangle_H=
\int_{[0,1]}u'(t)\overline{v'(t)}\,\mathrm d\lambda(t)$.
Warunki brzegowe zapewniają dodatnią określoność. Przestrzeń $H$
jest zupełna: odwzorowanie $u\mapsto u'$ jest izometryczną bijekcją
$H\to L^2_0:=\{g\in L^2([0,1],\lambda):
\int_{[0,1]}g(t)\,\mathrm d\lambda(t)=0\}$,
którego obraz jest domknięty. Odwrotność ma wzór
$g\mapsto(x\mapsto\int_{[0,x]}g(t)\,\mathrm d\lambda(t))$.
Dla $u\in H$ nierówność Cauchy'ego--Schwarza daje
$|u(x)|\le\sqrt{x}\,\|u'\|_2$, a więc
$\|u\|_2\le\|u'\|_2/\sqrt2$.

Dla $f\in L^2([0,1],\lambda)$ forma
$a(u,v)=\int_{[0,1]}u'(t)\overline{v'(t)}\,\mathrm d\lambda(t)$
jest ograniczona i koercywna ze stałą $1$, zaś
$\Lambda(v)=\int_{[0,1]}f(t)\overline{v(t)}\,\mathrm d\lambda(t)$
jest ciągła i antyliniowa na $H$.
Twierdzenie~\ref{thm:lax-milgram13} daje jedyne $u\in H$
spełniające $a(u,v)=\Lambda(v)$ dla wszystkich $v\in H$.
Jest to słaba postać zagadnienia $-u''=f$, $u(0)=u(1)=0$;
konkretnie $u'(t)=c-\int_{[0,t]}f(s)\,\mathrm d\lambda(s)$,
gdzie $c=\int_{[0,1]}(1-s)f(s)\,\mathrm d\lambda(s)$.
Weryfikacja przez całkowanie przez części jest bezpośrednia
najpierw dla gładkich funkcji $v$ o zerowych wartościach na brzegu,
a następnie przez gęstość w $H$. Istotnie, dla $g=v'\in L^2_0$
wybieramy $g_n\in C^\infty([0,1])$ z $g_n\to g$ w $L^2$
i kładziemy $\widetilde g_n=g_n-\int_{[0,1]}g_n(t)\,\mathrm d\lambda(t)$.
Wówczas $\widetilde g_n\to g$ w $L^2_0$, zaś funkcje
pierwotne $\widetilde g_n$ $v_n(x)=\int_{[0,x]}\widetilde g_n(t)\,\mathrm d\lambda(t)$
należą do $C^\infty([0,1])$, znikają na brzegu i zbiegają do $v$
w normie $H$.
\end{przyklad}

\begin{uwaga}[Droga do pochodnych Gâteaux i Frécheta]
\label{rem:gateaux-bridge13}
Jeśli $H$ jest rzeczywistą przestrzenią Hilberta, to funkcjonał
$\Phi(x)=\frac12\|x\|^2$ ma pochodną Frécheta
$D\Phi(x)h=\langle h,x\rangle$. Formalnie wynika to z
$\Phi(x+h)-\Phi(x)-\langle h,x\rangle=\frac12\|h\|^2$.
Definicję pochodnej Frécheta i jej relację do pochodnej Gâteaux
wprowadzimy w następnym rozdziale. W przestrzeni zespolonej
$\Phi$ jest funkcją rzeczywistą i rozważamy jej pochodną
\emph{rzeczywistą}, równą $h\mapsto\operatorname{Re}\langle h,x\rangle$.
\end{uwaga}


% Wersja źródłowa: audyt II, jednostka 14
\setcounter{section}{13}
\section{Pochodne Gâteaux i Frécheta w przestrzeniach Banacha i Hilberta}
\label{sec:gateaux-frechet14}

\subsection{Przygotowanie: kierunki, przyrosty i liniowe przybliżenia}

W analizie funkcji jednej zmiennej pochodna opisuje najlepsze liniowe przybliżenie funkcji. Dla odwzorowania między przestrzeniami unormowanymi możliwe są dwie naturalne interpretacje: przybliżenie \emph{osobno w każdym ustalonym kierunku} oraz przybliżenie \emph{jednostajne względem wszystkich małych przyrostów}. Pierwsza prowadzi do pochodnej Gâteaux, druga do pochodnej Frécheta. Różnica jest istotna już w $\R^2$.

Niech $E,F$ będą rzeczywistymi przestrzeniami Banacha, $U\subseteq E$ zbiorem otwartym i $f:U\to F$. Przez $\mathcal L(E,F)$ oznaczamy przestrzeń ograniczonych operatorów liniowych, wyposażoną w normę operatorową. Otwartość $U$ zapewnia, że dla $x\in U$ i $h\in E$ punkty $x+th$ należą do $U$ dla wszystkich dostatecznie małych $|t|$.

\begin{definicja}[Pochodna kierunkowa]
\label{def:directional14}
Dla $x\in U$ i $h\in E$ mówimy, że $f$ ma \emph{dwustronną pochodną kierunkową} w punkcie $x$ w kierunku $h$, jeśli w $F$ istnieje granica
$\partial_h f(x)=\lim_{t\to0,\,t\in\R}(f(x+th)-f(x))/t$.
Nie wymagamy, aby zależność $h\mapsto\partial_h f(x)$ była liniowa ani ciągła. Istnienie wszystkich pochodnych kierunkowych samo w sobie nie oznacza różniczkowalności Gâteaux.
\end{definicja}

\begin{definicja}[Różniczkowalność Gâteaux]
\label{def:gateaux14}
Odwzorowanie $f$ jest \emph{różniczkowalne w sensie Gâteaux} w $x\in U$, jeśli istnieje $A\in\mathcal L(E,F)$ takie, że dla każdego $h\in E$ zachodzi $\partial_h f(x)=Ah$. Operator $A$, jednoznacznie wyznaczony przez te granice, oznaczamy $D_Gf(x)$. Przyjęliśmy definicję wymagającą ciągłości pochodnej jako operatora liniowego; w niektórych źródłach spotyka się słabszą konwencję, dopuszczającą nieograniczone odwzorowanie liniowe.
\end{definicja}

\begin{definicja}[Różniczkowalność Frécheta]
\label{def:frechet14}
Odwzorowanie $f$ jest \emph{różniczkowalne w sensie Frécheta} w $x\in U$, jeśli istnieje $A\in\mathcal L(E,F)$ oraz funkcja $r$ określona dla małych $h\in E$ taka, że
\[
f(x+h)=f(x)+Ah+r(h),\qquad
\lim_{\|h\|_E\to0,\,h\ne0}\frac{\|r(h)\|_F}{\|h\|_E}=0.
\]
Operator $A$ oznaczamy $Df(x)$ albo $D_Ff(x)$; piszemy również $r(h)=o(\|h\|_E)$. Symbol $Df(x)h$ oznacza wartość operatora $Df(x)$ na wektorze $h$, nie jego iloczyn skalarny z $h$.
\end{definicja}

\begin{uwaga}[Założenia o zupełności]
\label{rem:completeness14}
Do samych definicji wystarczą przestrzenie unormowane i otwarta dziedzina. Zupełność potrzebna będzie między innymi do twierdzenia o całkowaniu funkcji o wartościach w przestrzeni Banacha oraz w późniejszych twierdzeniach o funkcji odwrotnej. Stąd w niniejszym rozdziale pracujemy przede wszystkim w przestrzeniach Banacha.
\end{uwaga}

\subsection{Implikacje między pochodnymi}

\begin{twierdzenie}[Fréchet implikuje Gâteaux]
\label{thm:frechet-implies-gateaux14}
Niech $E,F$ będą rzeczywistymi przestrzeniami unormowanymi, $U\subseteq E$ zbiorem otwartym, a $f:U\to F$ odwzorowaniem różniczkowalnym w sensie Frécheta w punkcie $x\in U$. Wówczas $f$ jest ciągłe w $x$, jest różniczkowalne w sensie Gâteaux w $x$ oraz $D_Gf(x)=Df(x)$. Pochodna Frécheta, jeśli istnieje, jest jednoznaczna.
\end{twierdzenie}
\begin{proof}
Niech $A=Df(x)$ i $f(x+h)-f(x)=Ah+r(h)$, gdzie $\|r(h)\|=o(\|h\|)$. Ponieważ $A$ jest ograniczony, $\|Ah+r(h)\|\le\|A\|\|h\|+o(\|h\|)\to0$, co daje ciągłość. Dla ustalonego $h\ne0$ i małego rzeczywistego $t$ mamy
$\|(f(x+th)-f(x))/t-Ah\|=\|r(th)\|/|t|=\|h\|\,\|r(th)\|/\|th\|\to0$.
Dla $h=0$ pochodna kierunkowa wynosi zero. Jeżeli $A$ i $B$ spełniają definicję Frécheta, powyższy argument daje $Ah=Bh$ dla każdego $h$, a więc $A=B$.
\end{proof}

\begin{stwierdzenie}[Kryterium jednostajnego przybliżenia kierunkowego]
\label{prop:uniform-directions14}
Niech $E,F$ będą rzeczywistymi przestrzeniami unormowanymi, $U\subseteq E$ otwarty, $x\in U$ i $A\in\mathcal L(E,F)$. Odwzorowanie $f:U\to F$ ma pochodną Frécheta $Df(x)=A$ wtedy i tylko wtedy, gdy
\[
\lim_{t\to0,\,t\in\R}\ \sup_{\|h\|_E=1}
\left\|\frac{f(x+th)-f(x)}{t}-Ah\right\|_F=0.
\]
W wyrażeniu supremum rozpatrujemy wyłącznie $t\ne0$ o module na tyle małym, że $x+tB_E\subseteq U$; zakładamy też $E\ne\{0\}$.
\end{stwierdzenie}
\begin{proof}
Dla reszty $r(k)=f(x+k)-f(x)-Ak$ mamy
$\sup_{\|h\|=1}\|r(th)\|/|t|=\sup_{\|k\|=|t|}\|r(k)\|/\|k\|$.
Warunek Frécheta zapewnia, że te suprema dążą do zera. Odwrotnie, dla $k\ne0$ połóżmy $t=\|k\|$ i $h=k/\|k\|$; uzyskana jednostajna granica daje $\|r(k)\|/\|k\|\to0$.
\end{proof}

\begin{uwaga}
\label{rem:quantifiers14}
W definicji Gâteaux kwantyfikatory mają kolejność: \emph{dla każdego ustalonego $h$ granica przy $t\to0$}. W definicji Frécheta wymagamy wspólnej kontroli błędu dla wszystkich kierunków. Zamiana kolejności kwantyfikatorów nie jest uprawniona, nawet w skończonym wymiarze.
\end{uwaga}

\subsection{Kontrprzykłady: czego pochodna Gâteaux nie gwarantuje}

\begin{przyklad}[Wszystkie pochodne kierunkowe i brak ciągłości]
\label{ex:directional-discontinuous14}
Dla $(x,y)\ne(0,0)$ połóżmy
$g(x,y)=x^3y/(x^6+y^2)$, zaś $g(0,0)=0$. Dla każdego $(a,b)\in\R^2$ granica $\partial_{(a,b)}g(0,0)$ istnieje i jest równa $0$: jeśli $b\ne0$, to $g(ta,tb)/t=a^3b\,t/(a^6t^4+b^2)\to0$; jeśli $b=0$, iloraz jest zerowy. Zatem $D_Gg(0,0)=0\in\mathcal L(\R^2,\R)$. Jednak na krzywej $y=x^3$, $x\ne0$, mamy $g(x,x^3)=1/2$. Funkcja nie jest ciągła w $(0,0)$, więc nie jest tam różniczkowalna w sensie Frécheta.
\end{przyklad}

\begin{przyklad}[Ciągłość i pochodna Gâteaux nie implikują pochodnej Frécheta]
\label{ex:continuous-gateaux-not-frechet14}
Niech $f(0,0)=0$, zaś dla $(x,y)\ne(0,0)$ niech
\[
f(x,y)=\frac{x^3y}{x^4+y^2}.
\]
Z nierówności $2|x^2y|\le x^4+y^2$ wynika $|f(x,y)|\le|x|/2$; zatem $f$ jest ciągła w $(0,0)$. Dla ustalonego $(a,b)$, gdy $b\ne0$, otrzymujemy
$f(ta,tb)/t=a^3b\,t/(a^4t^2+b^2)\to0$, a dla $b=0$ iloraz jest zerowy. Stąd $D_Gf(0,0)=0$. Natomiast $f(x,x^2)=x/2$ dla $x\ne0$ oraz
$|f(x,x^2)|/\sqrt{x^2+x^4}=1/(2\sqrt{1+x^2})\to1/2$ przy $x\to0^+$. Reszta względem zerowego operatora nie jest $o(\|(x,y)\|_2)$: pochodna Frécheta nie istnieje.
\end{przyklad}

\begin{przyklad}[Sama pochodna kierunkowa nie musi być liniowa]
\label{ex:nonlinear-directional14}
Dla $f(x,y)=|x|+|y|$ w zerze istnieje tylko jednostronna pochodna kierunkowa, równa $|a|+|b|$ w kierunku $(a,b)$; dwustronna pochodna kierunkowa przy $h\ne0$ nie istnieje. Dla przykładu z \emph{dwustronnymi} pochodnymi kierunkowymi, lecz nieliniową zależnością od kierunku, rozważmy $q(0,0)=0$ oraz $q(x,y)=x^3/(x^2+y^2)$ dla $(x,y)\ne0$. Wówczas $\partial_{(a,b)}q(0,0)=a^3/(a^2+b^2)$ dla $(a,b)\ne(0,0)$, a w kierunku zerowym pochodna wynosi $0$. Odwzorowanie to jest jednorodne względem mnożenia kierunku przez rzeczywiste skalary, lecz nie jest addytywne: $\partial_{(1,0)}q+\partial_{(0,1)}q=1$, podczas gdy $\partial_{(1,1)}q=1/2$. Dlatego $q$ nie ma pochodnej Gâteaux według Definicji~\ref{def:gateaux14}.
\end{przyklad}

\begin{stwierdzenie}[Brak reguły łańcuchowej dla samej pochodnej Gâteaux]
\label{prop:chain-fails14}
Istnieją ciągłe odwzorowanie $f:\R^2\to\R$, różniczkowalne w sensie Gâteaux w zerze, oraz gładka krzywa $\gamma:\R\to\R^2$ z $\gamma(0)=0$ takie, że $\frac{\mathrm d}{\mathrm dt}(f\circ\gamma)(0)\ne D_Gf(0)\gamma'(0)$.
\end{stwierdzenie}
\begin{proof}
Niech $f$ będzie funkcją z Przykładu~\ref{ex:continuous-gateaux-not-frechet14}, a $\gamma(t)=(t,t^2)$. Wówczas $f(\gamma(t))=t/2$ dla $t\ne0$, a także dla $t=0$ z ciągłości. Zatem $(f\circ\gamma)'(0)=1/2$, podczas gdy $D_Gf(0)=0$. Niezgodność wynika z niestabilności kierunkowego przybliżenia przy zmiennym kierunku przyrostu.
\end{proof}

\begin{center}
\begin{tikzpicture}[x=2.15cm,y=1.65cm,>=Stealth]
  \draw[->,thin] (-0.17,0)--(1.7,0) node[below] {$x$};
  \draw[->,thin] (0,-0.42)--(0,1.48) node[left] {$y$};
  \draw[gray,densely dashed] (0,0)--(1.32,0.62);
  \draw[gray,densely dashed] (0,0)--(0.79,1.22);
  \draw[gray,densely dashed] (0,0)--(1.42,-0.29);
  \draw[very thick,blue!55!black,domain=0:1.15,samples=70,smooth]
       plot (\x,{\x*\x});
  \fill (0,0) circle (1.05pt);
  \node[anchor=west,blue!55!black] at (1.17,1.22) {\footnotesize $y=x^2$};
  \node[anchor=west] at (0.83,-0.37) {\scriptsize kierunki stałe};
\end{tikzpicture}

{\small Proste przez zero odpowiadają ustalonym kierunkom w definicji
Gâteaux. Krzywa $y=x^2$ prowadzi do zera po zmieniającym się
kierunku; na niej ujawnia się brak różniczkowalności Frécheta.}
\end{center}

\subsection{Pochodna Gâteaux normy w \texorpdfstring{$\ell^1$}{ell1} bez pochodnej Frécheta}

\begin{twierdzenie}[Kontrprzykład w nieskończonym wymiarze]
\label{thm:l1-gateaux14}
Niech $E=\ell^1(\N;\R)$ i $N:E\to\R$, $N(x)=\|x\|_1$. Jeśli $x=(x_n)\in E$ spełnia $x_n\ne0$ dla wszystkich $n$, to $N$ jest różniczkowalne w sensie Gâteaux w $x$ i
$D_GN(x)h=\sum_{n=1}^\infty \operatorname{sgn}(x_n)h_n$.
Funkcja $N$ nie jest jednak różniczkowalna w sensie Frécheta w żadnym takim punkcie $x$.
\end{twierdzenie}
\begin{proof}
Dla ustalonego $h\in\ell^1$ i $t\ne0$ oznaczmy $a_n(t)=(|x_n+th_n|-|x_n|)/t$. Ponieważ $x_n\ne0$, mamy $a_n(t)\to\operatorname{sgn}(x_n)h_n$ dla każdego $n$, przy czym $|a_n(t)|\le|h_n|$. Zbieżność zmajoryzowana dla szeregu (albo argument z obcięciem sumy i oszacowaniem ogona przez $\sum_{n>N}|h_n|$) daje
$\lim_{t\to0}(N(x+th)-N(x))/t=\sum_n \operatorname{sgn}(x_n)h_n$.
Jest to funkcjonał liniowy i ciągły o normie $1$.

Jeżeli istniałaby pochodna Frécheta, musiałaby być równa $D_GN(x)$. Dla $h^{(n)}=-2x_ne_n$ mamy $\|h^{(n)}\|_1=2|x_n|\to0$, ponieważ $x\in\ell^1$, a zarazem $N(x+h^{(n)})=N(x)$ i $D_GN(x)h^{(n)}=-2|x_n|$. Wobec tego
\[
\frac{|N(x+h^{(n)})-N(x)-D_GN(x)h^{(n)}|}{\|h^{(n)}\|_1}=1
\]
dla każdego $n$. Jest to sprzeczne z warunkiem Frécheta.
\end{proof}

\begin{uwaga}
Jeśli pewna współrzędna $x_n$ znika, norma $\ell^1$ nie ma w tym punkcie nawet dwustronnej pochodnej kierunkowej w kierunku $e_n$. Powyższy kontrprzykład pokazuje więc ścisłą różnicę między Gâteaux a Fréchetem również w punktach, w których wszystkie pochodne kierunkowe tworzą ciągły funkcjonał liniowy.
\end{uwaga}

\subsection{Reguły różniczkowania w przestrzeniach Banacha}

\begin{stwierdzenie}[Suma i operator liniowy]
\label{prop:sum-linear14}
Niech $E,F,G$ będą rzeczywistymi przestrzeniami Banacha, $U\subseteq E$ otwarty i $x\in U$. Jeżeli $f,g:U\to F$ są różniczkowalne w sensie Frécheta w $x$, a $\alpha,\beta\in\R$, to $D(\alpha f+\beta g)(x)=\alpha Df(x)+\beta Dg(x)$. Jeśli $T\in\mathcal L(F,G)$, to $D(T\circ f)(x)=T\circ Df(x)$. Ponadto dla $L\in\mathcal L(E,F)$ i $b\in F$ funkcja $f(y)=Ly+b$ ma $Df(y)=L$ w każdym $y\in E$.
\end{stwierdzenie}
\begin{proof}
Wszystkie równości wynikają przez podstawienie rozwinięć $f(x+h)=f(x)+Df(x)h+r_f(h)$ oraz $g(x+h)=g(x)+Dg(x)h+r_g(h)$ i użycie $\|Tr_f(h)\|\le\|T\|\|r_f(h)\|$. W przypadku funkcji afinicznej reszta wynosi zero.
\end{proof}

\begin{twierdzenie}[Reguła łańcuchowa Frécheta]
\label{thm:chain14}
Niech $E,F,G$ będą rzeczywistymi przestrzeniami Banacha, $U\subseteq E$, $V\subseteq F$ otwarte, $f:U\to F$ oraz $g:V\to G$ i $f(U)\subseteq V$. Jeśli $f$ jest różniczkowalne w sensie Frécheta w $x\in U$, a $g$ w $f(x)$, to $g\circ f$ jest różniczkowalne w sensie Frécheta w $x$, a
$D(g\circ f)(x)=Dg(f(x))\circ Df(x)$.
\end{twierdzenie}
\begin{proof}
Oznaczmy $A=Df(x)$, $B=Dg(f(x))$. Wówczas $f(x+h)=f(x)+Ah+r(h)$, $\|r(h)\|=o(\|h\|)$, i dla $k\to0$ mamy $g(f(x)+k)=g(f(x))+Bk+s(k)$, gdzie $\|s(k)\|=o(\|k\|)$. Ponieważ $f$ jest ciągłe w $x$, dla małego $h$ można przyjąć $k=Ah+r(h)$, a $\|k\|\le(\|A\|+1)\|h\|$. Otrzymujemy
$g(f(x+h))-g(f(x))-BAh=Br(h)+s(Ah+r(h))$.
Pierwszy składnik jest $o(\|h\|)$ przez ograniczoność $B$. Drugi także: jeśli $k\ne0$, to
$\|s(k)\|/\|h\|=(\|s(k)\|/\|k\|)(\|k\|/\|h\|)\to0$;
jeśli $k=0$, to $s(k)=0$. Dowód nie wymaga ciągłości pochodnych w otoczeniu rozważanych punktów.
\end{proof}

\begin{twierdzenie}[Reguła dla odwzorowania dwuliniowego]
\label{thm:bilinear14}
Niech $E_1,E_2,G$ będą rzeczywistymi przestrzeniami Banacha i $B:E_1\times E_2\to G$ ciągłym odwzorowaniem dwuliniowym, tj. $\|B(u,v)\|\le C\|u\|\|v\|$. Wówczas $B$ jest różniczkowalne w sensie Frécheta na przestrzeni produktowej $E_1\times E_2$ i
$DB(u,v)(h,k)=B(h,v)+B(u,k)$.
Jeżeli $f:U\to E_1$, $g:U\to E_2$ są różniczkowalne w $x\in U\subseteq E$ otwartym, to
$D(B(f,g))(x)h=B(Df(x)h,g(x))+B(f(x),Dg(x)h)$.
\end{twierdzenie}
\begin{proof}
Dokładna tożsamość
$B(u+h,v+k)-B(u,v)-B(h,v)-B(u,k)=B(h,k)$
daje $\|B(h,k)\|\le C\|h\|\|k\|\le C\max(\|h\|,\|k\|)^2$. Reszta jest więc $o(\max(\|h\|,\|k\|))$, a operator liniowy wskazany w tezie jest ograniczony. Ostatni wzór wynika z reguły łańcuchowej zastosowanej do odwzorowania $z\mapsto(f(z),g(z))$ i do $B$.
\end{proof}

\begin{przyklad}[Operatory i funkcje kwadratowe]
\label{ex:quadratic14}
Niech $E$ będzie rzeczywistą przestrzenią Banacha, $F$ rzeczywistą przestrzenią Banacha i $B:E\times E\to F$ ciągłym odwzorowaniem dwuliniowym. Dla $Q(x)=B(x,x)$ zachodzi
$DQ(x)h=B(h,x)+B(x,h)$.
W szczególności, jeśli $H$ jest rzeczywistą przestrzenią Hilberta i $\Phi(x)=\frac12\|x\|^2$, to
$D\Phi(x)h=\langle h,x\rangle$, a reszta ma postać $\frac12\|h\|^2$.
\end{przyklad}

\subsection{Różniczkowalność klasy \texorpdfstring{$C^1$}{C1} i twierdzenie o przyrostach}

\begin{definicja}[Klasa $C^1$]
\label{def:c1-14}
Niech $E,F$ będą rzeczywistymi przestrzeniami Banacha, $U\subseteq E$ otwarty i $f:U\to F$. Piszemy $f\in C^1(U,F)$, jeśli $f$ jest różniczkowalne w sensie Frécheta w każdym punkcie $U$ oraz odwzorowanie $Df:U\to\mathcal L(E,F)$ jest ciągłe w normie operatorowej.
\end{definicja}

\begin{lemat}[Podstawowe twierdzenie rachunku dla krzywych Banacha]
\label{lem:ftc-banach14}
Niech $F$ będzie rzeczywistą przestrzenią Banacha, $a<b$ i $\gamma:[a,b]\to F$ funkcją klasy $C^1$ (pochodna w końcach jednostronna i ciągle przedłużona). Wówczas $\gamma'$ jest całkowalna w sensie Riemanna z wartościami w $F$ oraz
$\gamma(b)-\gamma(a)=\int_a^b\gamma'(t)\,\mathrm dt$.
\end{lemat}
\begin{proof}
Ciągłość $\gamma'$ na zwartym przedziale daje jej jednostajną ciągłość, a przez kryterium Cauchy'ego sum Riemanna i zupełność $F$ zapewnia istnienie całki. Dla podziału $a=t_0<\cdots<t_m=b$ wybierzmy dowolne $\xi_j\in[t_{j-1},t_j]$. Z definicji pochodnej i jednostajnej ciągłości $\gamma'$ wynika, że przy dostatecznie małej średnicy podziału
$\|\gamma(t_j)-\gamma(t_{j-1})-(t_j-t_{j-1})\gamma'(\xi_j)\|\le\varepsilon(t_j-t_{j-1})$.
Ściślej, dla dowolnego rzeczywistego funkcjonału $\ell\in F'$ o normie co najwyżej $1$ stosujemy skalarne twierdzenie o wartości średniej do $t\mapsto\ell(\gamma(t)-t\gamma'(\xi_j))$ na $[t_{j-1},t_j]$. Otrzymujemy górne oszacowanie modułu przyrostu przez $(t_j-t_{j-1})\sup_{t\in[t_{j-1},t_j]}\|\gamma'(t)-\gamma'(\xi_j)\|$. Z twierdzenia Hahna--Banacha, po supremum po takich $\ell$, dostajemy to samo oszacowanie normy wektora. Jednostajna ciągłość $\gamma'$ zapewnia żądaną nierówność dla dostatecznie drobnych podziałów. Sumowanie oraz przejście z średnicą podziału do zera kończy dowód.
\end{proof}

\begin{twierdzenie}[Wzór całkowy i oszacowanie przyrostu]
\label{thm:mean-value14}
Niech $E,F$ będą rzeczywistymi przestrzeniami Banacha, $U\subseteq E$ zbiorem otwartym, $f\in C^1(U,F)$, a $x,y\in U$ takimi punktami, że $[x,y]=\{x+t(y-x):t\in[0,1]\}\subseteq U$. Wówczas
\[
f(y)-f(x)=\int_0^1 Df(x+t(y-x))(y-x)\,\mathrm dt,
\]
gdzie całka ma wartości w $F$. W szczególności, jeśli $\sup_{z\in[x,y]}\|Df(z)\|\le M$, to $\|f(y)-f(x)\|\le M\|y-x\|$.
\end{twierdzenie}
\begin{proof}
Definiujemy $\gamma(t)=f(x+t(y-x))$. Reguła łańcuchowa daje $\gamma'(t)=Df(x+t(y-x))(y-x)$; ciągłość $Df$ zapewnia ciągłość $\gamma'$. Lemat~\ref{lem:ftc-banach14} daje wzór. Oszacowanie normy całki przez całkę normy wynika z oszacowania sum Riemanna i przejścia do granicy.
\end{proof}

\begin{twierdzenie}[Ciągła pochodna Gâteaux daje pochodną Frécheta]
\label{thm:continuous-gateaux14}
Niech $E,F$ będą rzeczywistymi przestrzeniami Banacha, $U\subseteq E$ otwarty, $f:U\to F$. Załóżmy, że $f$ jest różniczkowalne w sensie Gâteaux w każdym punkcie $U$ i odwzorowanie $D_Gf:U\to\mathcal L(E,F)$ jest ciągłe w normie operatorowej. Wtedy $f\in C^1(U,F)$ oraz $Df(x)=D_Gf(x)$ dla wszystkich $x\in U$.
\end{twierdzenie}
\begin{proof}
Ustalmy $x\in U$ i weźmy małe $h$ takie, że $[x,x+h]\subseteq U$. Funkcja $\gamma(t)=f(x+th)$ jest różniczkowalna dla $t\in(0,1)$ i $\gamma'(t)=D_Gf(x+th)h$, co wynika z definicji kierunkowej. Jej pochodna jest ciągła dzięki założeniu, a $\gamma$ jest ciągła również w końcach przedziału, ponieważ istnienie pochodnej kierunkowej w tych punktach zapewnia ciągłość wzdłuż tej prostej. Z Lematu~\ref{lem:ftc-banach14} mamy
$f(x+h)-f(x)=\int_0^1 D_Gf(x+th)h\,\mathrm dt$.
Po odjęciu $D_Gf(x)h$ i oszacowaniu normy dostajemy
\[
\frac{\|f(x+h)-f(x)-D_Gf(x)h\|}{\|h\|}
\le\sup_{0\le t\le1}\|D_Gf(x+th)-D_Gf(x)\|\longrightarrow0.
\]
To dowodzi różniczkowalności Frécheta; ciągłość pochodnej otrzymujemy z założenia.
\end{proof}

\begin{przyklad}[Pochodna istniejąca, ale nieciągła]
\label{ex:derivative-noncontinuous14}
Na $\R$ niech $f(0)=0$ oraz $f(x)=x^2\sin(1/x^2)$ dla $x\ne0$. Wtedy $f'(0)=0$, natomiast dla $x\ne0$ mamy $f'(x)=2x\sin(1/x^2)-2x^{-1}\cos(1/x^2)$. Zatem $f$ jest różniczkowalna Frécheta w każdym punkcie, lecz nie należy do $C^1(\R,\R)$. Nie należy utożsamiać istnienia pochodnej z jej ciągłością.
\end{przyklad}

\subsection{Norma i gradient w rzeczywistych i zespolonych przestrzeniach Hilberta}

\begin{twierdzenie}[Różniczkowanie kwadratu normy i normy]
\label{thm:hilbert-norm14}
Niech $H$ będzie przestrzenią Hilberta nad $\K\in\{\R,\C\}$, traktowaną jako \emph{rzeczywista} przestrzeń Banacha, i niech $\Phi(x)=\frac12\|x\|^2$. W każdym $x\in H$ istnieje rzeczywista pochodna Frécheta $D_{\R}\Phi(x)h=\operatorname{Re}\langle h,x\rangle$. Dla $x\ne0$ norma $N(x)=\|x\|$ ma rzeczywistą pochodną Frécheta
$D_{\R}N(x)h=\operatorname{Re}\langle h,x\rangle/\|x\|$.
Jeżeli $H\ne\{0\}$, norma nie jest różniczkowalna Gâteaux w $0$ jako funkcja na rzeczywistej przestrzeni Banacha.
\end{twierdzenie}
\begin{proof}
Z rozwinięcia iloczynu skalarnego wynika dokładna tożsamość
$\Phi(x+h)-\Phi(x)-\operatorname{Re}\langle h,x\rangle=\frac12\|h\|^2$, a funkcjonał $h\mapsto\operatorname{Re}\langle h,x\rangle$ jest rzeczywisto-liniowy i ograniczony. To dowodzi pierwszej tezy. Dla $x\ne0$ stosujemy rzeczywistą regułę łańcuchową do $N(x)=\sqrt{2\Phi(x)}$, ponieważ pierwiastek jest gładki na $(0,\infty)$. W zerze iloraz kierunkowy wynosi $\|th\|/t=(|t|/t)\|h\|$; dla $h\ne0$ nie istnieje granica dwustronna.
\end{proof}

\begin{definicja}[Gradient na rzeczywistej przestrzeni Hilberta]
\label{def:gradient14}
Niech $H$ będzie rzeczywistą przestrzenią Hilberta, $U\subseteq H$ otwarty i $f:U\to\R$ różniczkowalna w sensie Frécheta w $x\in U$. \emph{Gradientem} $f$ w $x$ nazywamy jedyny wektor $\nabla f(x)\in H$ spełniający $Df(x)h=\langle h,\nabla f(x)\rangle$ dla wszystkich $h\in H$. Istnienie i jednoznaczność daje twierdzenie Riesza. Dla zespolonej przestrzeni Hilberta i funkcjonału rzeczywistego używamy struktury rzeczywistej i iloczynu skalarnego $(h,k)_{\R}=\operatorname{Re}\langle h,k\rangle$.
\end{definicja}

\begin{przyklad}[Gradient energii kwadratowej]
\label{ex:quadratic-energy14}
Niech $H$ będzie rzeczywistą lub zespoloną przestrzenią Hilberta, $A\in\mathcal L(H)$ operatorem samosprzężonym, $b\in H$, a $J(x)=\frac12\langle Ax,x\rangle-\operatorname{Re}\langle x,b\rangle$. Z samosprzężoności wynika, że $J$ ma wartości rzeczywiste. Dla $h\in H$ otrzymujemy
$D_{\R}J(x)h=\operatorname{Re}\langle h,Ax-b\rangle$, więc $\nabla J(x)=Ax-b$ w rzeczywistej strukturze Hilberta. Jeżeli $A$ jest dodatnio określony w sensie $\langle Ax,x\rangle\ge\alpha\|x\|^2$ dla pewnego $\alpha>0$, to $J$ ma dokładnie jeden minimizator: rozwiązanie $Ax=b$; jego istnienie zapewnia przypadek samosprzężony twierdzenia Laxa--Milgrama.
\end{przyklad}

\begin{przyklad}[Funkcjonał całkowy i pochodna na $L^2$]
\label{ex:integral-functional14}
Niech $(\Omega,\mathcal A,\mu)$ będzie przestrzenią z miarą skończoną i $H=L^2(\Omega,\mu;\R)$. Dla $g\in H$ określmy $J(u)=\frac12\int_\Omega |u(t)-g(t)|^2\,\mathrm d\mu(t)$. Wtedy
$DJ(u)h=\int_\Omega (u(t)-g(t))h(t)\,\mathrm d\mu(t)$,
a $\nabla J(u)=u-g$. Wynik można otrzymać przez rozwinięcie kwadratu oraz nierówność Cauchy'ego--Schwarza; reszta ma dokładnie postać $\frac12\|h\|_2^2$.
\end{przyklad}

\subsection{Rzeczywista a zespolona różniczkowalność}

\begin{definicja}[Pochodna zespolona Frécheta i Gâteaux]
\label{def:complex-differentiability14}
Niech $E,F$ będą zespolonymi przestrzeniami Banacha, $U\subseteq E$ zbiorem otwartym i $f:U\to F$. \emph{Zespoloną pochodną Frécheta} $D_{\C}f(x)$ nazywamy operator $A\in\mathcal L_{\C}(E,F)$ spełniający $f(x+h)=f(x)+Ah+o(\|h\|)$. \emph{Zespolona pochodna Gâteaux} w $x$ oznacza istnienie operatora $A\in\mathcal L_{\C}(E,F)$ takiego, że dla każdego $h\in E$
$\lim_{z\to0,\,z\in\C}(f(x+zh)-f(x))/z=Ah$.
W obu definicjach liniowość wymagana jest nad $\C$, a dla pochodnej Gâteaux parametr $z$ przebiega płaszczyznę zespoloną, nie tylko oś rzeczywistą.
\end{definicja}

\begin{stwierdzenie}[Pochodna zespolona a rzeczywista]
\label{prop:complex-real14}
Niech $E,F$ będą zespolonymi przestrzeniami Banacha, $U\subseteq E$ otwarty i $f:U\to F$. Jeśli $f$ jest zespolenie różniczkowalne w sensie Frécheta w $x\in U$, to jest też rzeczywiście różniczkowalne w sensie Frécheta w $x$, a $D_{\R}f(x)$ jest operatorem rzeczywistym odpowiadającym $D_{\C}f(x)$. Odwrotnie, jeśli $D_{\R}f(x)$ istnieje, to $D_{\C}f(x)$ istnieje wtedy i tylko wtedy, gdy $D_{\R}f(x)(ih)=iD_{\R}f(x)h$ dla wszystkich $h\in E$.
\end{stwierdzenie}
\begin{proof}
Pochodna zespolona spełnia oszacowanie reszty wymagane dla pochodnej rzeczywistej, ponieważ normy nie zmieniają się po zapomnieniu struktury zespolonej. Jeśli istnieje $D_{\R}f(x)=A$, to identyczna reszta $o(\|h\|)$ daje zespoloną pochodną dokładnie wtedy, gdy rzeczywisto-liniowy operator $A$ jest zespolenie liniowy. Dla operatora rzeczywisto-liniowego $A$ warunek $A(ih)=iA(h)$ jest równoważny liniowości zespolonej.
\end{proof}

\begin{przyklad}[Sprzężenie zespolone i kwadrat modułu]
\label{ex:complex-nonholomorphic14}
Funkcja $f:\C\to\C$, $f(z)=\overline z$, jest rzeczywiście różniczkowalna Frécheta w każdym punkcie, z $D_{\R}f(z)h=\overline h$, ale nie jest zespolenie różniczkowalna w żadnym punkcie. Podobnie $z\mapsto|z|^2$ ma rzeczywistą pochodną $h\mapsto2\operatorname{Re}(h\overline z)$, nie ma zaś zespolonej pochodnej Frécheta w punktach $z\ne0$. W punkcie $z=0$ zespolona pochodna Frécheta funkcji $z\mapsto|z|^2$ istnieje i wynosi zero, gdyż $|h|^2=o(|h|)$.
\end{przyklad}

\begin{uwaga}[Poziomy ogólności]
Związki między kierunkową a pełną różniczkowalnością, normą, topologią operatorową i różniczkowalnością zespoloną wykorzystują normowalność dziedziny. Istnieją pojęcia różniczkowania na ogólnych przestrzeniach lokalnie wypukłych (m.in. różniczkowalność Michala--Bastiani i rachunek wygodny), ale wymagają odrębnego aparatu. Nie utożsamiamy ich bez dodatkowych definicji z rachunkiem Frécheta w przestrzeniach Banacha.
\end{uwaga}


% Wersja źródłowa: audyt II, jednostka 15
\setcounter{section}{14}
\section{Pochodne wyższych rzędów i lokalne twierdzenia analizy nieliniowej}
\label{sec:wyzsze-pochodne15}

\subsection{Odwzorowania wieloliniowe i definicja wyższych pochodnych}

W podstawowym rachunku różniczkowym różniczkę $Df(x)$ traktowaliśmy jako ograniczony operator liniowy. Przy dalszym różniczkowaniu otrzymujemy odwzorowania wieloliniowe. Ta perspektywa pozwala sformułować twierdzenie Taylora i twierdzenia o funkcji odwrotnej oraz uwikłanej bez redukowania przestrzeni Banacha do współrzędnych.

\begin{definicja}[Ciągłe odwzorowania wieloliniowe]
\label{def:multilinear15}
Niech $E_1,\ldots,E_m,F$ będą rzeczywistymi przestrzeniami unormowanymi, $m\ge1$. Przez $\mathcal L^m(E_1,\ldots,E_m;F)$ oznaczamy przestrzeń odwzorowań $A:E_1\times\cdots\times E_m\to F$, liniowych w każdym argumencie i ciągłych, wyposażoną w normę
$\|A\|=\sup\{\|A(h_1,\ldots,h_m)\|:\|h_j\|\le1,\ 1\le j\le m\}$. Jeśli wszystkie dziedziny są równe $E$, piszemy $\mathcal L^m(E;F)$, a dla $m=1$ używamy dotychczasowego oznaczenia $\mathcal L(E,F)$. Przez $\mathcal L^m_s(E;F)$ oznaczamy podprzestrzeń odwzorowań symetrycznych, tj. niezmienniczych przy permutacjach argumentów.
\end{definicja}

\begin{stwierdzenie}[Kryterium ograniczoności i zupełność]
\label{prop:multi-banach15}
Niech $E_1,\ldots,E_m$ będą przestrzeniami unormowanymi, a $F$ przestrzenią Banacha. Odwzorowanie $m$-liniowe $A:E_1\times\cdots\times E_m\to F$ jest ciągłe wtedy i tylko wtedy, gdy istnieje $C\ge0$ takie, że
$\|A(h_1,\ldots,h_m)\|\le C\prod_{j=1}^m\|h_j\|$ dla wszystkich $h_j\in E_j$. Przestrzeń $\mathcal L^m(E_1,\ldots,E_m;F)$ jest przestrzenią Banacha.
\end{stwierdzenie}
\begin{proof}
Jeśli $A$ jest ciągłe w zerze, istnieje $\delta>0$ takie, że $\|A(v_1,\ldots,v_m)\|\le1$ dla $\|v_j\|<\delta$. Dla niezerowych $h_j$ podstawiamy $v_j=\delta h_j/(2\|h_j\|)$, co po wykorzystaniu wieloliniowości daje oszacowanie z $C=(2/\delta)^m$. Gdy któryś argument jest zerowy, nierówność jest oczywista. W drugą stronę nierówność zapewnia ciągłość w zerze, a wieloliniowość i rozwinięcie $A(h_1+k_1,\ldots,h_m+k_m)$ w skończoną sumę zapewniają ciągłość w dowolnym punkcie.

Niech $(A_n)$ będzie ciągiem Cauchy'ego w normie operatorowej. Dla ustalonych $h_1,\ldots,h_m$ ciąg $A_n(h_1,\ldots,h_m)$ jest Cauchy'ego w $F$, więc definiuje $A(h_1,\ldots,h_m)$. Granica zachowuje wieloliniowość. Ponieważ normy $\|A_n\|$ są ograniczone przez wspólną stałą, $A$ spełnia oszacowanie wieloliniowe. Przechodząc do granicy w nierówności $\|(A_n-A_p)(h_1,\ldots,h_m)\|\le\varepsilon\prod\|h_j\|$ dla $n,p\ge N$, otrzymujemy $\|A_n-A\|\le\varepsilon$. To dowodzi zupełności.
\end{proof}

\begin{lemat}[Utożsamienie operatorów wieloliniowych]
\label{lem:currying15}
Niech $E,F$ będą rzeczywistymi przestrzeniami unormowanymi, a $m\ge1$. Naturalne odwzorowanie
$\mathcal L(E,\mathcal L^m(E;F))\to\mathcal L^{m+1}(E;F)$, zadane przez $T\mapsto\widetilde T$ oraz $\widetilde T(h_0,h_1,\ldots,h_m)=(Th_0)(h_1,\ldots,h_m)$, jest izometrycznym izomorfizmem liniowym.
\end{lemat}
\begin{proof}
Liniowość, wieloliniowość i ciągłość wynikają z oszacowania $\|\widetilde T(h_0,\ldots,h_m)\|\le\|T\|\prod_{j=0}^m\|h_j\|$. Ponadto, z definicji norm operatorowych,
\[
\begin{aligned}
\|\widetilde T\|&=\sup_{\|h_0\|\le1}\sup_{\substack{\|h_j\|\le1\\1\le j\le m}}\|(Th_0)(h_1,\ldots,h_m)\|\\
&=\|T\|.
\end{aligned}
\]
Dla $A\in\mathcal L^{m+1}(E;F)$ definiujemy $(T_Ah_0)(h_1,\ldots,h_m)=A(h_0,h_1,\ldots,h_m)$. Wówczas $T_A\in\mathcal L(E,\mathcal L^m(E;F))$ oraz $\widetilde T_A=A$, co daje surjektywność.
\end{proof}

\begin{definicja}[Klasa $C^m$ i pochodne wyższych rzędów]
\label{def:ck15}
Niech $E,F$ będą rzeczywistymi przestrzeniami Banacha, $U\subseteq E$ zbiorem otwartym, $m\ge1$. Mówimy, że $f:U\to F$ jest klasy $C^m$, gdy $f\in C^1$ oraz rekurencyjnie $D^{j-1}f:U\to\mathcal L^{j-1}(E;F)$ jest klasy $C^1$ dla $2\le j\le m$. Pochodną $D^jf(x)$ utożsamiamy przez Lemat~\ref{lem:currying15} z elementem $\mathcal L^j(E;F)$. Przyjmujemy $D^0f=f$ oraz $C^\infty(U,F)=\bigcap_{m\ge1}C^m(U,F)$. Ciągłość $D^jf$ rozumiemy w normie $\mathcal L^j(E;F)$.
\end{definicja}

\begin{przyklad}[Odwzorowanie wieloliniowe i wielomiany]
\label{ex:multi-poly15}
Niech $B\in\mathcal L^m(E;F)$ i $P(x)=B(x,\ldots,x)$. Wówczas $P\in C^\infty(E,F)$ oraz
$DP(x)h=\sum_{j=1}^m B(x,\ldots,x,\underset{j}{h},x,\ldots,x)$.
Istotnie, różnica $P(x+h)-P(x)$ składa się z wyrazów, w których $h$ występuje co najmniej raz; wyrazy z dokładnie jednym $h$ tworzą wskazany operator, a suma pozostałych jest $O(\|h\|^2)$ przy ustalonym $x$. Powtarzanie argumentu daje wyższe pochodne, które zerują się od rzędu $m+1$. Dla symetrycznego $B$ zachodzi $DP(x)h=mB(x,\ldots,x,h)$, a $D^mP(x)=m!B$.
\end{przyklad}

\subsection{Symetria drugiej pochodnej i reguły rachunku}

\begin{twierdzenie}[Symetria drugiej pochodnej]
\label{thm:schwarz15}
Niech $E,F$ będą rzeczywistymi przestrzeniami Banacha, $U\subseteq E$ otwarty i $f\in C^2(U,F)$. Dla każdego $x\in U$ i $h,k\in E$ zachodzi
$D^2f(x)[h,k]=D^2f(x)[k,h]$. W szczególności $D^2f(x)\in\mathcal L_s^2(E;F)$.
\end{twierdzenie}
\begin{proof}
Dla małych rzeczywistych $s,t$ połóżmy
$\Delta(s,t)=f(x+sh+tk)-f(x+sh)-f(x+tk)+f(x)$. Z dwukrotnego zastosowania wzoru całkowego na odcinku dla funkcji klasy $C^1$ otrzymujemy
\[
\Delta(s,t)=st\int_0^1\int_0^1 D^2f(x+\alpha sh+\beta tk)[h,k] \,\mathrm d\beta\,\mathrm d\alpha.
\]
Można bowiem najpierw scałkować po kierunku $h$, a następnie różnicę wyrażeń $Df$ po kierunku $k$. Odwrócenie kolejności tych dwóch kroków daje ten sam przyrost $\Delta(s,t)$ oraz wyrażenie z $D^2f(\cdot)[k,h]$. Dla $s,t\ne0$ dzielimy przez $st$ i przechodzimy z $(s,t)$ do $(0,0)$. Ciągłość $D^2f$ w normie operatorowej zapewnia jednostajną zbieżność funkcji podcałkowych na kwadracie $[0,1]^2$ odpowiednio do $D^2f(x)[h,k]$ i $D^2f(x)[k,h]$. Obie granice są równe, co kończy dowód.
\end{proof}


\begin{wniosek}[Symetria wszystkich pochodnych wyższych rzędów]
\label{cor:symmetric-higher15}
Niech $E,F$ będą rzeczywistymi przestrzeniami Banacha, $U\subseteq E$ otwarty, $m\ge2$ i $f\in C^m(U,F)$. Wówczas $D^mf(x)\in\mathcal L_s^m(E;F)$ dla każdego $x\in U$.
\end{wniosek}
\begin{proof}
Stosujemy indukcję po $m$. Dla $m=2$ jest to Twierdzenie~\ref{thm:schwarz15}. Jeśli $m\ge3$, to odwzorowanie $D^{m-2}f:U\to\mathcal L^{m-2}(E;F)$ jest klasy $C^2$, a więc zastosowanie tego twierdzenia do $D^{m-2}f$ daje symetrię $D^mf(x)$ względem pierwszych dwóch argumentów. Z założenia indukcyjnego $D^{m-1}f(z)$ jest symetryczne dla każdego $z\in U$. Różniczkując tożsamości wyrażające tę symetrię w kierunku pierwszego argumentu, stwierdzamy, że $D^mf(x)$ jest symetryczne także względem ostatnich $m-1$ argumentów. Łącznie otrzymujemy niezmienniczość przy wszystkich zamianach sąsiednich argumentów, a te generują wszystkie permutacje.
\end{proof}

\begin{stwierdzenie}[Druga pochodna złożenia]
\label{prop:second-chain15}
Niech $E,F,G$ będą rzeczywistymi przestrzeniami Banacha, $U\subseteq E$ i $V\subseteq F$ otwarte, $f\in C^2(U,F)$, $g\in C^2(V,G)$ oraz $f(U)\subseteq V$. Wówczas $g\circ f\in C^2(U,G)$ i dla $x\in U$ oraz $h,k\in E$ zachodzi
\[
\begin{aligned}
D^2(g\circ f)(x)[h,k]
={}&D^2g(f(x))[Df(x)h,Df(x)k]\\
 &+Dg(f(x))\bigl(D^2f(x)[h,k]\bigr).
\end{aligned}
\]
\end{stwierdzenie}
\begin{proof}
Z reguły łańcuchowej $D(g\circ f)(x)=Dg(f(x))\circ Df(x)$. Kompozycja operatorów jest ciągłym odwzorowaniem dwuliniowym $\mathcal L(F,G)\times\mathcal L(E,F)\to\mathcal L(E,G)$. Stosując regułę różniczkowania ciągłego odwzorowania dwuliniowego do odwzorowań $x\mapsto Dg(f(x))$ i $x\mapsto Df(x)$, otrzymujemy po zadziałaniu na $k$ dwa składniki. Pierwszy, dzięki regule łańcuchowej zastosowanej do $Dg$, ma postać $D^2g(f(x))[Df(x)h,Df(x)k]$, a drugi jest równy $Dg(f(x))(D^2f(x)[h,k])$. Ciągłość wszystkich występujących pochodnych dowodzi przynależności złożenia do $C^2$.
\end{proof}

\begin{przyklad}[Hesjan funkcjonału na Hilbercie]
\label{ex:hessian15}
Niech $H$ będzie rzeczywistą przestrzenią Hilberta, $A\in\mathcal L(H)$ i $b\in H$. Dla $J(x)=\frac12\langle Ax,x\rangle-\langle x,b\rangle$ mamy
$DJ(x)h=\langle h,(A+A^*)x/2-b\rangle$ oraz
$D^2J(x)[h,k]=\langle h,(A+A^*)k/2\rangle$.
Przy konwencji liniowości iloczynu skalarnego w pierwszym argumencie wzór jest symetryczny dla rzeczywistego $H$, ponieważ $(A+A^*)/2$ jest samosprzężony. Jeśli $A=A^*$, to $D^2J(x)[h,k]=\langle h,Ak\rangle$.
\end{przyklad}

\subsection{Twierdzenie Taylora w przestrzeniach Banacha}

\begin{twierdzenie}[Wzór Taylora z resztą całkową]
\label{thm:taylor15}
Niech $E,F$ będą rzeczywistymi przestrzeniami Banacha, $U\subseteq E$ otwarty, $m\ge1$, $f\in C^m(U,F)$, a $x\in U$ i $h\in E$ takie, że odcinek $\{x+th:0\le t\le1\}$ zawiera się w $U$. Wtedy
\[
 f(x+h)=\sum_{j=0}^{m-1}\frac1{j!}D^jf(x)[h^j]
 +\frac1{(m-1)!}\int_0^1(1-t)^{m-1}D^mf(x+th)[h^m] \,\mathrm dt,
\]
gdzie $D^0f(x)[h^0]=f(x)$, a $D^jf(x)[h^j]$ oznacza $D^jf(x)[h,\ldots,h]$. Ponadto
$f(x+h)=\sum_{j=0}^{m}\frac1{j!}D^jf(x)[h^j]+r_m(h)$, gdzie
\[
\|r_m(h)\|\le \frac{\|h\|^m}{m!}\sup_{0\le t\le1}\|D^mf(x+th)-D^mf(x)\|.
\]
W szczególności $\|r_m(h)\|=o(\|h\|^m)$ dla $h\to0$.
\end{twierdzenie}
\begin{proof}
Ustalmy $x,h$ i zdefiniujmy $\gamma(t)=f(x+th)$. Z reguły łańcuchowej i indukcji mamy $\gamma^{(j)}(t)=D^jf(x+th)[h^j]$ dla $1\le j\le m$. Pierwsza pochodna jest ciągła, zatem stosujemy podstawowe twierdzenie rachunku całkowego dla funkcji o wartościach w $F$. Iterowanie wzoru $\gamma(t)-\gamma(0)=\int_0^t\gamma'(s)\,\mathrm ds$ prowadzi do jednokrotnej całki z jądrem $(1-t)^{m-1}/(m-1)!$: krok indukcyjny wynika z zamiany kolejności całkowania na trójkącie $0\le s\le t\le1$, dopuszczalnej dla ciągłej funkcji o wartościach w Banachu (można ją uzasadnić najpierw dla sum Riemanna). Otrzymujemy pierwszy wzór.

Odejmujemy od wyrażenia pod całką stałą $D^mf(x)[h^m]$, której całka z wagą $(1-t)^{m-1}/(m-1)!$ jest równa $D^mf(x)[h^m]/m!$. Następnie stosujemy normę operatorową, oszacowanie normy całki przez całkę normy oraz $\int_0^1(1-t)^{m-1}\,\mathrm dt=1/m$. Ciągłość $D^mf$ w $x$ zapewnia, że supremum w oszacowaniu reszty dąży do zera przy $h\to0$.
\end{proof}

\begin{wniosek}[Wzór drugiego rzędu]
\label{cor:taylor2-15}
Przy założeniach Twierdzenia~\ref{thm:taylor15} dla $m=2$ zachodzi
\[
 f(x+h)=f(x)+Df(x)h+\tfrac12D^2f(x)[h,h]
 +\int_0^1(1-t)\bigl(D^2f(x+th)-D^2f(x)\bigr)[h,h] \,\mathrm dt.
\]
Jeżeli na odcinku $[x,x+h]$ zachodzi $\|D^2f(z)\|\le M$, to reszta po wyrazach rzędu pierwszego ma normę co najwyżej $M\|h\|^2/2$.
\end{wniosek}

\begin{stwierdzenie}[Warunki konieczne i wystarczające minimum]
\label{prop:extremum15}
Niech $E$ będzie rzeczywistą przestrzenią Banacha, $U\subseteq E$ otwarty, $x_0\in U$, a $f\in C^2(U,\R)$. Jeżeli $x_0$ jest lokalnym minimum $f$, to $Df(x_0)=0$ i $D^2f(x_0)[h,h]\ge0$ dla każdego $h\in E$. Odwrotnie, jeżeli $Df(x_0)=0$ oraz istnieje $c>0$ takie, że $D^2f(x_0)[h,h]\ge c\|h\|^2$ dla wszystkich $h\in E$, to $x_0$ jest ścisłym lokalnym minimum $f$; dokładniej, dla dostatecznie małych $h$ zachodzi $f(x_0+h)\ge f(x_0)+(c/4)\|h\|^2$.
\end{stwierdzenie}
\begin{proof}
Dla dowolnego $h$ rozważamy rzeczywistą funkcję $\varphi(t)=f(x_0+th)$. Ma ona minimum w $0$, więc $\varphi'(0)=Df(x_0)h=0$ i $\varphi''(0)=D^2f(x_0)[h,h]\ge0$. W drugim kierunku ciągłość $D^2f$ daje promień $r>0$ taki, że $\|D^2f(x_0+u)-D^2f(x_0)\|\le c/2$ dla $\|u\|<r$. Stąd $D^2f(x_0+u)[h,h]\ge(c/2)\|h\|^2$. Z całkowego wzoru Taylora drugiego rzędu, przy $Df(x_0)=0$, otrzymujemy $f(x_0+h)-f(x_0)\ge(c/2)\|h\|^2\int_0^1(1-t)\,\mathrm dt=c\|h\|^2/4$ dla $\|h\|<r$.
\end{proof}

\begin{przyklad}[Dodatniość bez jednostajności nie wystarcza]
\label{ex:positive-hessian15}
W rzeczywistej przestrzeni Hilberta $H=\ell^2$ niech
$Q(x)=\frac12\sum_{n=1}^\infty x_n^2/n$, zaś $f(x)=Q(x)-\|x\|_2^3$. Funkcja $f$ jest klasy $C^2$ w otoczeniu zera. Mamy $Df(0)=0$ i $D^2f(0)[h,h]=\sum_{n\ge1}h_n^2/n>0$ dla $h\ne0$, ale $0$ nie jest minimum lokalnym. Istotnie, $f(te_n)=t^2/(2n)-t^3<0=f(0)$, gdy $t>1/(2n)$, a takie dodatnie $t$ można wybierać dowolnie małe przez zwiększanie $n$. Funkcja $G(x)=\|x\|_2^3$ jest klasy $C^2$ również w zerze. Dla $x\ne0$ mamy
$DG(x)h=3\|x\|\langle x,h\rangle$ oraz
$D^2G(x)[h,k]=3\|x\|\langle h,k\rangle+3\langle x,h\rangle\langle x,k\rangle/\|x\|$.
Stąd $\|D^2G(x)\|\le6\|x\|$. Ponadto $DG(0)=0$ i $\|DG(x)\|\le3\|x\|^2$, co dowodzi $D^2G(0)=0$ i ciągłości drugiej pochodnej w zerze.
\end{przyklad}

\subsection{Szereg Neumanna i odwracanie operatorów}

\begin{lemat}[Szereg Neumanna]
\label{lem:neumann15}
Niech $E$ będzie rzeczywistą przestrzenią Banacha i $T\in\mathcal L(E)$ spełnia $\|T\|<1$. Wtedy $I-T$ jest odwracalny oraz
$(I-T)^{-1}=\sum_{n=0}^\infty T^n$ w normie operatorowej, przy czym $\|(I-T)^{-1}\|\le(1-\|T\|)^{-1}$.
\end{lemat}
\begin{proof}
Szereg jest absolutnie zbieżny w przestrzeni Banacha $\mathcal L(E)$. Dla sum częściowych $S_N=\sum_{n=0}^N T^n$ zachodzi $(I-T)S_N=S_N(I-T)=I-T^{N+1}$; przechodząc do granicy, otrzymujemy obustronną odwrotność. Oszacowanie normy wynika z szeregu geometrycznego.
\end{proof}

\begin{stwierdzenie}[Otwartość grupy operatorów odwracalnych]
\label{prop:inverse-operator15}
Niech $E$ będzie rzeczywistą przestrzenią Banacha i $\mathrm{GL}(E)$ zbiorem odwracalnych operatorów w $\mathcal L(E)$. Wówczas $\mathrm{GL}(E)$ jest otwarty, a $\iota:\mathrm{GL}(E)\to\mathcal L(E)$, $\iota(A)=A^{-1}$, jest klasy $C^\infty$. Jej pochodna ma postać
$D\iota(A)H=-A^{-1}HA^{-1}$.
\end{stwierdzenie}
\begin{proof}
Dla $A\in\mathrm{GL}(E)$ i $\|A^{-1}H\|<1$ mamy $A+H=A(I+A^{-1}H)$, więc Lemat~\ref{lem:neumann15} daje odwracalność oraz lokalną jednostajną ograniczoność odwrotności. Z tożsamości
$(A+H)^{-1}-A^{-1}=-(A+H)^{-1}HA^{-1}$ wynika
\[
(A+H)^{-1}-A^{-1}+A^{-1}HA^{-1}
=A^{-1}H(A+H)^{-1}HA^{-1}.
\]
Jej prawą stronę ograniczamy przez $C_A\|H\|^2$ dla małych $H$, skąd wynika pochodna Frécheta. Jej ciągłość otrzymujemy z ciągłości odwracania i ciągłości kompozycji operatorów. Powtarzając różniczkowanie wzoru $D\iota(A)H=-A^{-1}HA^{-1}$ z regułą wieloliniową, otrzymujemy rekurencyjnie pochodne wszystkich rzędów i ich ciągłość.
\end{proof}

\subsection{Twierdzenie o funkcji odwrotnej}

\begin{twierdzenie}[Twierdzenie o funkcji odwrotnej dla przestrzeni Banacha]
\label{thm:inverse15}
Niech $E,F$ będą rzeczywistymi przestrzeniami Banacha, $U\subseteq E$ otwarty, $f\in C^1(U,F)$ i $x_0\in U$. Załóżmy, że $A=Df(x_0):E\to F$ jest izomorfizmem liniowym i topologicznym, tj. ma ograniczoną odwrotność $A^{-1}:F\to E$. Wówczas istnieją otoczenia otwarte $U_0\subseteq U$ punktu $x_0$ i $V_0\subseteq F$ punktu $y_0=f(x_0)$ takie, że $f|_{U_0}:U_0\to V_0$ jest bijekcją, a odwrotność $g:V_0\to U_0$ należy do $C^1(V_0,E)$ i spełnia
$Dg(y)=(Df(g(y)))^{-1}$ dla każdego $y\in V_0$.
\end{twierdzenie}
\begin{proof}
\emph{Krok 1: jednostajne przybliżenie liniowe.} Z ciągłości $Df$ w $x_0$ wybieramy $r>0$ tak małe, aby domknięta kula $\overline B(x_0,r)$ zawierała się w $U$ i
$\|I-A^{-1}Df(x)\|\le q<1$ dla $x\in\overline B(x_0,r)$; można przyjąć $q=1/2$. Dla $y\in F$ określamy $T_y(x)=x-A^{-1}(f(x)-y)$. Dla $x,z$ w tej kuli wzór całkowy na odcinku daje
\[
\|T_y(x)-T_y(z)\|\le q\|x-z\|.
\]
Ponadto $T_y(x_0)-x_0=A^{-1}(y-y_0)$. Dla $\|A^{-1}(y-y_0)\|\le(1-q)r$ operator $T_y$ przeprowadza $\overline B(x_0,r)$ w siebie, bo
$\|T_y(x)-x_0\|\le q\|x-x_0\|+\|A^{-1}(y-y_0)\|\le r$.

\emph{Krok 2: punkt stały.} Przypomnijmy zasadę kontrakcji: jeżeli $T$ jest odwzorowaniem $q$-lipschitzowskim pełnej przestrzeni metrycznej w siebie, $q<1$, to ciąg $x_{n+1}=T(x_n)$ jest Cauchy'ego, ponieważ $d(x_{n+1},x_n)\le q^n d(x_1,x_0)$; jego granica jest jedynym punktem stałym. Stosujemy tę zasadę do $T_y$ na domkniętej kuli w $E$. Otrzymujemy jedyny $x$ z $T_y(x)=x$, co jest równoważne $f(x)=y$. Jeżeli $y$ spełnia nierówność ostrą $\|A^{-1}(y-y_0)\|<(1-q)r$, punkt stały leży we wnętrzu kuli. Dodatkowo każdy $x\in B(x_0,r)$ ma na tym obszarze co najwyżej jeden punkt o tej samej wartości $f(x)$: jeżeli $f(x)=f(z)$, to $x,z$ są punktami stałymi tej samej kontrakcji $T_{f(x)}$.

\emph{Krok 3: otwarte otoczenia i ciągłość odwrotności.} Ustalmy $0<r_0<r$ i połóżmy $U_0=B(x_0,r_0)$ oraz $V_0=f(U_0)$. Z oszacowania
$\|(x-z)-A^{-1}(f(x)-f(z))\|\le q\|x-z\|$ wynika
\[
\|x-z\|\le\frac{\|A^{-1}\|}{1-q}\,\|f(x)-f(z)\| \qquad(x,z\in B(x_0,r)).
\]
Zatem $f|_{U_0}$ jest różnowartościowe, a odwrotność $g:V_0\to U_0$ jest lokalnie Lipschitzowska. Zbiór $V_0$ jest otwarty: dla dowolnego $x\in U_0$ wybieramy kulę $\overline B(x,\delta)\subset U_0$; ta sama kontrakcja $T_y$ zachowuje $\overline B(x,\delta)$ przy $y$ dostatecznie bliskim $f(x)$, ponieważ oszacowanie stałej $q$ obowiązuje w całej kuli $B(x_0,r)$. W konsekwencji otoczenie $f(x)$ zawiera się w $V_0$.

\emph{Krok 4: różniczkowalność odwrotności.} Zmniejszając $r$, jeśli trzeba, mamy $\|I-A^{-1}Df(x)\|<1$ dla $x\in U_0$, więc każdy $Df(x)$ jest odwracalny z Lematu~\ref{lem:neumann15}. Ustalmy $y=f(x)\in V_0$ i rozważmy $y+k=f(x+h)$, gdzie $h=g(y+k)-g(y)\to0$. Z różniczkowalności $f$ w $x$:
$k=Df(x)h+\rho(h)$ z $\|\rho(h)\|=o(\|h\|)$. Po pomnożeniu przez $(Df(x))^{-1}$ dostajemy
$h-(Df(x))^{-1}k=-(Df(x))^{-1}\rho(h)=o(\|h\|)=o(\|k\|)$,
gdyż otrzymana w kroku 3 nierówność zapewnia $\|h\|\le C\|k\|$. Stąd $Dg(y)=(Df(x))^{-1}$. Ciągłość $Dg$ wynika z ciągłości $Df$, $g$ oraz odwracania w $\mathrm{GL}(E,F)$ (można zastosować ten sam argument Neumanna do $(Df(x))^{-1}Df(x')$). Tak więc $g\in C^1$.
\end{proof}

\begin{wniosek}[Wersja $C^m$]
\label{cor:inverse-ck15}
Jeżeli przy założeniach Twierdzenia~\ref{thm:inverse15} odwzorowanie $f$ jest klasy $C^m$ dla $m\ge2$ albo klasy $C^\infty$, to lokalna odwrotność $g$ ma tę samą regularność.
\end{wniosek}
\begin{proof}
Wzór $Dg=(Df\circ g)^{-1}$ oraz Stwierdzenie~\ref{prop:inverse-operator15} pozwalają prowadzić indukcję: jeśli $g\in C^j$ i $f\in C^{j+1}$, to $Df\circ g\in C^j$, zatem także $Dg\in C^j$, a więc $g\in C^{j+1}$. Zaczynamy od $g\in C^1$.
\end{proof}

\begin{przyklad}[Nieliniowa mapa na przestrzeni Hilberta]
\label{ex:inverse-hilbert15}
Niech $H$ będzie rzeczywistą przestrzenią Hilberta, $v\in H$ ustalony i $f:H\to H$ dane przez $f(x)=x+\|x\|^2v$. Wtedy $Df(x)h=h+2\langle x,h\rangle v$, a $Df(0)=I$. Twierdzenie~\ref{thm:inverse15} zapewnia lokalną odwrotność w otoczeniu zera. Jeżeli $v\ne0$, nie wnioskujemy stąd, że $f$ ma odwrotność na całym $H$.
\end{przyklad}

\subsection{Twierdzenie o funkcji uwikłanej}

\begin{twierdzenie}[Twierdzenie o funkcji uwikłanej dla przestrzeni Banacha]
\label{thm:implicit15}
Niech $X,Y,Z$ będą rzeczywistymi przestrzeniami Banacha, $\mathcal U\subseteq X\times Y$ zbiorem otwartym, $F\in C^1(\mathcal U,Z)$, $(x_0,y_0)\in\mathcal U$ i $F(x_0,y_0)=0$. Załóżmy, że pochodna częściowa $D_yF(x_0,y_0):Y\to Z$ jest izomorfizmem liniowym i topologicznym. Wtedy istnieją otwarte otoczenia $U_0\subseteq X$ punktu $x_0$, $V_0\subseteq Y$ punktu $y_0$ oraz jedyna funkcja $\varphi:U_0\to V_0$ klasy $C^1$ taka, że dla $(x,y)\in U_0\times V_0$ zachodzi
$F(x,y)=0\ \Longleftrightarrow\ y=\varphi(x)$. Ponadto
\[
D\varphi(x)=-\bigl(D_yF(x,\varphi(x))\bigr)^{-1}D_xF(x,\varphi(x))\qquad(x\in U_0).
\]
Jeżeli $F\in C^m$ dla $m\ge2$ lub $C^\infty$, to również $\varphi$ ma odpowiednią regularność.
\end{twierdzenie}
\begin{proof}
Definiujemy $\Psi:\mathcal U\to X\times Z$ wzorem $\Psi(x,y)=(x,F(x,y))$. Jest to odwzorowanie klasy $C^1$ i
$D\Psi(x_0,y_0)(h,k)=(h,Bh+Ak)$, gdzie $A=D_yF(x_0,y_0)$ i $B=D_xF(x_0,y_0)$. Operator ten jest odwracalny: odwrotny przeprowadza $(h,z)$ w $(h,A^{-1}(z-Bh))$ i jest ograniczony. Twierdzenie~\ref{thm:inverse15} daje lokalny dyfeomorfizm $\Psi$ między otoczeniami $(x_0,y_0)$ oraz $(x_0,0)$. Po zmniejszeniu otoczeń możemy wybrać prostokąt $U_0\times W_0$ zawarty w obrazie $\Psi$ oraz $U_0\times V_0$ zawarty w dziedzinie lokalnego dyfeomorfizmu. Zapiszmy odwrotność $\Psi^{-1}(x,z)=(x,G(x,z))$; pierwsza współrzędna ma tę postać, ponieważ $\Psi$ zachowuje $x$. Połóżmy $\varphi(x)=G(x,0)$. Wówczas $F(x,\varphi(x))=0$.

Dla zapewnienia równoważności na całym $U_0\times V_0$ zwężamy $V_0$ i $U_0$ tak, by $\Psi$ była różnowartościowa na tym prostokącie i $\varphi(U_0)\subset V_0$. Jeśli $F(x,y)=0$, to $\Psi(x,y)=(x,0)=\Psi(x,\varphi(x))$, zatem $y=\varphi(x)$. Różniczkując tożsamość $F(x,\varphi(x))=0$ i używając reguły łańcuchowej, otrzymujemy
$D_xF(x,\varphi(x))+D_yF(x,\varphi(x))D\varphi(x)=0$.
Po ewentualnym zmniejszeniu otoczeń $D_yF$ jest odwracalne na grafie $\varphi$, bo zbiór operatorów odwracalnych jest otwarty. Stąd wzór na $D\varphi$. Wersja wyższej regularności wynika z Wniosku~\ref{cor:inverse-ck15}.
\end{proof}

\begin{przyklad}[Rodzina równań operatorowych]
\label{ex:implicit-operator15}
Niech $E$ będzie rzeczywistą przestrzenią Banacha, $K\in\mathcal L(E)$ i $g\in E$. Na $\R\times E$ definiujemy $F(\lambda,u)=u-\lambda Ku-g$. W punkcie $(0,g)$ mamy $F(0,g)=0$ oraz $D_uF(0,g)=I$, zatem istnieje lokalnie jednoznaczne rozwiązanie $u(\lambda)$ klasy $C^\infty$ równania $(I-\lambda K)u(\lambda)=g$. Z wzoru na pochodną funkcji uwikłanej wynika
$u'(\lambda)=(I-\lambda K)^{-1}Ku(\lambda)$, w szczególności $u'(0)=Kg$. Różniczkując raz jeszcze, otrzymujemy $u''(0)=2K^2g$. Dla $|\lambda|\,\|K\|<1$ rozwiązanie można również zapisać w postaci zbieżnego szeregu Neumanna $u(\lambda)=\sum_{n\ge0}\lambda^nK^ng$.
\end{przyklad}

\begin{uwaga}[Dlaczego przestrzenie Banacha i dlaczego pochodna Frécheta?]
\label{rem:limits-inverse15}
Dowód twierdzenia o funkcji odwrotnej wykorzystuje zupełność domkniętej kuli (zasada kontrakcji) oraz jednostajną kontrolę operatora $Df(x)$ w normie. Sama odwracalność pochodnej Gâteaux w punkcie, bez stosownych założeń o ciągłości i pełnej różniczkowalności, nie daje takiego wniosku. W ogólnych przestrzeniach lokalnie wypukłych istnieją odrębne twierdzenia o funkcji odwrotnej, lecz nie wolno przenosić powyższego dowodu na dowolną TVS przez formalne zastąpienie normy seminormami.
\end{uwaga}


% Wersja źródłowa: audyt II, jednostka 15A
\setcounter{section}{15}
\noindent\textbf{Wymagane wyniki.} Korzystamy z twierdzenia Hahna--Banacha, twierdzenia Banacha--Steinhausa, podstawowego twierdzenia rachunku dla krzywych o wartościach w przestrzeni Banacha i faktu, że ciągłe funkcjonały liniowe rozdzielają punkty lokalnie wypukłej przestrzeni Hausdorffa. Pełną teorię całek na takich przestrzeniach rozwijamy poniżej.
\medskip
% \hyperref[sec:integral-lcs15a]{Moduł 15A}. Celowo nie zmieniamy numeracji \hyperref[sec:wariacyjny16]{rozdziałów 16}--33.
\begingroup
\renewcommand{\thesection}{15A}
\renewcommand{\theHsection}{15A}
\section{Całkowanie funkcji o wartościach w przestrzeniach lokalnie wypukłych}
\label{sec:integral-lcs15a}

W \hyperref[sec:gateaux-frechet14]{rozdziale~14} udowodniliśmy podstawowe twierdzenie
Newtona--Leibniza dla krzywych w przestrzeni Banacha, potrzebne w rachunku
różniczkowym. Teraz budujemy szerszą teorię: obiektem całkowanym może być
funkcja przyjmująca wartości w przestrzeni lokalnie wypukłej, której topologia
nie musi pochodzić od jednej normy. Nie należy utożsamiać tego zagadnienia
z całkowaniem funkcji mierzalnych: całki Bochnera, Pettisa, Dunforda i Gelfanda
będą przedmiotem \hyperref[sec:integral-vector15b]{modułu~15B}.

W całym module $\K\in\{\R,\C\}$, a $E$ jest przestrzenią liniowo-topologiczną
Hausdorffa nad $\K$, lokalnie wypukłą, z topologią wyznaczoną przez rodzinę
ciągłych seminorm $\mathcal P$. \emph{Zupełność ciągowa} oznacza, że każdy
ciąg Cauchy'ego w $E$ ma granicę w $E$; nie mylimy jej z ogólną zupełnością
względem sieci Cauchy'ego. Jeśli nie zaznaczono inaczej, $a<b$ są liczbami
rzeczywistymi. Funkcja na odcinku jest rozumiana z topologią podprzestrzeni.

\subsection{Sumy Riemanna i twierdzenie o istnieniu}

\begin{definicja}[Całka Riemanna o wartościach w przestrzeni liniowo-topologicznej]
\label{def:riemann-lcs15a}
Niech $f:[a,b]\to E$. \emph{Podziałem z punktami pośrednimi} nazywamy
parę $P=(a=t_0<\cdots<t_N=b;\,\xi_j\in[t_{j-1},t_j])$.
Piszemy $|P|=\max_j(t_j-t_{j-1})$ i
$S(f,P)=\sum_{j=1}^N(t_j-t_{j-1})f(\xi_j)\in E$.
Funkcja $f$ jest całkowalna w sensie Riemanna w $E$, jeżeli istnieje $I\in E$
takie, że dla każdej seminormy $p\in\mathcal P$ i $\varepsilon>0$ istnieje
$\delta>0$ taka, iż $p(S(f,P)-I)<\varepsilon$ dla \emph{każdego} podziału
z punktami pośrednimi spełniającego $|P|<\delta$.
Wówczas oznaczamy $I=\int_a^b f(t)\,\mathrm dt$.
Dla $a=b$ przyjmujemy całkę równą zeru, a dla $b<a$ całkę
przeciwną do $\int_b^a f(t)\,\mathrm dt$.
\end{definicja}

\begin{uwaga}[Znaczenie topologii]
Definicja wymaga zbieżności \emph{w topologii $E$}, a nie tylko po
zastosowaniu każdego funkcjonału $\ell\in E'$. Tę drugą własność nazywamy
słabą całkowalnością skalarną i sama nie gwarantuje istnienia wektora całki.
Powyższa granica jest granicą rodziny wszystkich sum Riemanna przy
$|P|\to0$; nie sprowadzamy jej z definicji do jednego wybranego ciągu podziałów.
\end{uwaga}

\begin{twierdzenie}[Istnienie całki przy zupełności ciągowej]
\label{thm:riemann-existence15a}
Niech $E$ będzie przestrzenią lokalnie wypukłą Hausdorffa,
zupełną ciągowo. Wtedy każda funkcja ciągła $f:[a,b]\to E$
jest całkowalna w sensie Riemanna. W szczególności wynik
zachodzi w każdej zupełnej przestrzeni lokalnie wypukłej Hausdorffa.
\end{twierdzenie}
\begin{proof}
Dla $p\in\mathcal P$ definiujemy moduł ciągłości
$\omega_p(r)=\sup\{p(f(t)-f(s)):s,t\in[a,b],\ |s-t|\le r\}$.
Ciągłość $f$ na zwartym odcinku oznacza jednostajną ciągłość
w topologii $E$, zatem $\omega_p(r)\to0$ dla $r\downarrow0$.

Niech $P,Q$ będą dwoma podziałami z punktami pośrednimi, o średnicach
co najwyżej $r$. Rozważmy wspólne rozdrobnienie ich podziałów
przedziału i wybierzmy na nim dowolne punkty pośrednie, otrzymując $R$.
Każdy przedział $R$ mieści się w jednym przedziale $P$, a odpowiadające
mu punkty pośrednie są odległe o najwyżej $r$; analogicznie dla $Q$.
Po rozpisaniu sum względem wspólnego rozdrobnienia dostajemy
\[
 p(S(f,P)-S(f,R))\le(b-a)\omega_p(r),\qquad
 p(S(f,Q)-S(f,R))\le(b-a)\omega_p(r).
\]
Zatem $p(S(f,P)-S(f,Q))\le2(b-a)\omega_p(r)$.

Wybierzmy ciąg podziałów dyadycznych $P_n$, na przykład ze środkami
przedziałów jako punktami pośrednimi, o średnicach dążących do zera.
Powyższa nierówność pokazuje, że $(S(f,P_n))$ jest ciągiem Cauchy'ego
w $E$. Z zupełności ciągowej istnieje $I\in E$, do którego zbiega.
Dla dowolnego $P$ o średnicy mniejszej od $r$ stosujemy ostatnie
oszacowanie do $P$ i $P_n$ przy dostatecznie dużym $n$, a następnie
przechodzimy do granicy $n\to\infty$:
$p(S(f,P)-I)\le2(b-a)\omega_p(r)$.
Ponieważ dotyczy to każdej seminormy, zachodzi warunek z definicji.
Jednoznaczność $I$ wynika z hausdorffowskości $E$.
\end{proof}

\begin{stwierdzenie}[Liniowość i oszacowania seminormowe]
\label{prop:riemann-seminorm15a}
Niech $E$ będzie przestrzenią lokalnie wypukłą Hausdorffa,
a $f,g:[a,b]\to E$ funkcjami ciągłymi i całkowalnymi w $E$.
Dla $\alpha,\beta\in\K$ i $c\in[a,b]$ zachodzą
\[
 \int_a^b(\alpha f(t)+\beta g(t))\,\mathrm dt
 =\alpha\int_a^b f(t)\,\mathrm dt+\beta\int_a^b g(t)\,\mathrm dt,
 \qquad \int_a^b f(t)\,\mathrm dt=\int_a^c f(t)\,\mathrm dt
 +\int_c^b f(t)\,\mathrm dt.
\]
Dla każdej ciągłej seminormy $p$ na $E$ mamy
\[
 p\left(\int_a^b f(t)\,\mathrm dt\right)
 \le\int_a^b p(f(t))\,\mathrm dt
 \le (b-a)\sup_{t\in[a,b]}p(f(t)).
\]
\end{stwierdzenie}
\begin{proof}
Liniowość wynika z liniowości każdej sumy Riemanna i ciągłości
działań w $E$. Dla addytywności względem przedziału można brać
podziały zawierające punkt $c$; sumy rozpadają się wtedy na dwie
części. Zbieżność całki dla wszystkich drobnych podziałów
pozwala przejść do granicy niezależnie od ich doboru.
Nierówność trójkąta i dodatnia jednorodność seminormy dają
$p(S(f,P))\le\sum_j(t_j-t_{j-1})p(f(\xi_j))$.
Po przejściu do granicy lewa strona zbiega do
$p(\int_a^b f(t)\,\mathrm dt)$, a prawa do skalarnej
całki funkcji ciągłej $p\circ f$. Ostatnia nierówność jest skalarna.
\end{proof}

\begin{twierdzenie}[Przemienność z odwzorowaniami liniowymi ciągłymi]
\label{thm:riemann-functorial15a}
Niech $E,F$ będą przestrzeniami lokalnie wypukłymi Hausdorffa,
$f:[a,b]\to E$ funkcją ciągłą i całkowalną w $E$, a
$T:E\to F$ odwzorowaniem liniowym ciągłym. Wtedy $T\circ f$
jest całkowalna w $F$ i
\[
 T\left(\int_a^b f(t)\,\mathrm dt\right)
 =\int_a^b T(f(t))\,\mathrm dt.
\]
W szczególności równość zachodzi dla każdego $\ell\in E'$, a
całkę w $E$ wyznaczają jednoznacznie jej wartości skalarne
$\ell(\int_a^b f(t)\,\mathrm dt)$ dla $\ell\in E'$.
Nie zakładamy zupełności $F$.
\end{twierdzenie}
\begin{proof}
Dla każdego $P$ mamy $S(T\circ f,P)=T(S(f,P))$.
Ciągłość $T$ przenosi granicę sum na $T(I)$; ten wektor jest
całką w $F$, nawet jeśli $F$ nie jest zupełna.
Funkcjonały z $E'$ rozdzielają punkty $E$ na mocy
twierdzenia Hahna--Banacha i hausdorffowskości $E$,
więc równości wszystkich obrazów skalarnych wyznaczają $I$.
\end{proof}

\begin{wniosek}[Środek ciężkości krzywej]
\label{cor:barycenter15a}
Niech $E$ będzie zupełną ciągowo przestrzenią lokalnie wypukłą
Hausdorffa i niech $f:[a,b]\to E$ będzie ciągła. Wówczas
\[
 (b-a)^{-1}\int_a^b f(t)\,\mathrm dt
 \in\overline{\operatorname{conv}}\bigl(f([a,b])\bigr),
\]
gdzie domknięcie i otoczkę wypukłą rozumiemy w topologii $E$.
\end{wniosek}
\begin{proof}
Dla każdego podziału $P$ wektor $(b-a)^{-1}S(f,P)$ jest
kombinacją wypukłą wartości $f(\xi_j)$, ponieważ jego
współczynniki $(t_j-t_{j-1})/(b-a)$ są nieujemne i sumują się do $1$.
Zamkniętość wskazanego zbioru i zbieżność sum dowodzą tezy.
\end{proof}

\subsection{Rola zupełności i całka niewłaściwa}

\begin{przyklad}[Funkcja ciągła bez całki w przestrzeni niezupełnej]
\label{ex:noncomplete-integral15a}
Rozważmy $E=c_{00}$, przestrzeń skończenie podpartych ciągów
rzeczywistych, wyposażoną w normę odziedziczoną z $\ell^2$.
Jest ona niezupełna, a jej uzupełnieniem jest $\ell^2$.
Wybierzmy nieujemną, niezerową funkcję ciągłą $\eta:\R\to\R$
o zwartym nośniku zawartym w $(0,1)$.
Dla $t\in[0,1]$ niech
\[
 f(t)=\sum_{n=1}^{\infty}2^{-n}\eta(2^{n+1}t-1)e_n.
\]
Przedziały, na których kolejne składniki są niezerowe,
leżą w rozłącznych zbiorach $(2^{-(n+1)},2^{-n})$.
Dlatego dla każdego $t$ suma ma co najwyżej jeden niezerowy
wyraz i $f(t)\in c_{00}$. Funkcja jest ciągła na $(0,1]$,
a w zerze jest ciągła dzięki oszacowaniu
$\|f(t)\|_2\le2^{-n}\|\eta\|_\infty$ na $n$-tym przedziale.
Całka istnieje w $\ell^2$, lecz jej $n$-ta współrzędna to
\[
 \int_0^1 2^{-n}\eta(2^{n+1}t-1)\,\mathrm dt
 =2^{-2n-1}\int_0^1\eta(s)\,\mathrm ds>0.
\]
Ma ona więc nieskończenie wiele niezerowych współrzędnych i nie
należy do $c_{00}$. Gdyby $f$ była całkowalna w $E$, przemienność
z ciągłymi funkcjonałami współrzędnych wymusiłaby tę samą
nieskończenie podpartą wartość; sprzeczność.
\end{przyklad}

\begin{stwierdzenie}[Całka w uzupełnieniu]
\label{prop:completion-integral15a}
Niech $E$ będzie przestrzenią lokalnie wypukłą Hausdorffa, a
$j:E\to\widehat E$ jej kanonicznym zanurzeniem do uzupełnienia.
Dla ciągłej $f:[a,b]\to E$ funkcja $j\circ f$ jest
całkowalna w $\widehat E$. Całka $f$ w $E$ istnieje wtedy
i tylko wtedy, gdy $\int_a^b j(f(t))\,\mathrm dt\in j(E)$;
wówczas $j(\int_a^b f(t)\,\mathrm dt)
=\int_a^b j(f(t))\,\mathrm dt$.
\end{stwierdzenie}
\begin{proof}
Uzupełnienie $\widehat E$ jest zupełne, zatem pierwsza teza
wynika z twierdzenia o istnieniu. Jeżeli całka w $E$ istnieje,
przemienność z $j$ daje konieczność. W drugą stronę załóżmy,
że całka w $\widehat E$ jest postaci $j(v)$ dla $v\in E$.
Rodzina sum Riemanna należy do $j(E)$ i zbiega w $\widehat E$
do $j(v)$. Topologia $E$ jest topologią indukowaną przez
zanurzenie $j$, więc sumy te zbiegają do $v$ w $E$.
\end{proof}

\begin{twierdzenie}[Kryterium całki niewłaściwej]
\label{thm:improper-lcs15a}
Niech $E$ będzie przestrzenią lokalnie wypukłą Hausdorffa,
zupełną ciągowo, a $f:[a,\infty)\to E$ funkcją ciągłą.
Załóżmy, że dla każdej ciągłej seminormy $p$ zachodzi
$\int_a^\infty p(f(t))\,\mathrm dt<\infty$ w sensie
niewłaściwej całki skalarnej. Wówczas istnieje
$\lim_{R\to\infty}\int_a^R f(t)\,\mathrm dt\in E$, oznaczana przez
$\int_a^\infty f(t)\,\mathrm dt$, oraz
$p(\int_a^\infty f(t)\,\mathrm dt)
\le\int_a^\infty p(f(t))\,\mathrm dt$.
\end{twierdzenie}
\begin{proof}
Dla $R'>R\ge a$ z oszacowania seminormowego i addytywności
całki wynika
$p(\int_a^{R'}f(t)\,\mathrm dt-\int_a^R f(t)\,\mathrm dt)
\le\int_R^{R'}p(f(t))\,\mathrm dt$.
Ogon ostatniej całki dąży do zera, dlatego ciąg
$\bigl(\int_a^{a+n}f(t)\,\mathrm dt\bigr)_{n\ge1}$
jest ciągiem Cauchy'ego w $E$ i ma granicę $I$.
Dla dowolnego $R\ge a+n$ to samo oszacowanie względem
$a+n$ zapewnia zbieżność do $I$ także przy rzeczywistym
$R\to\infty$. Oszacowanie końcowe wynika z przejścia do granicy.
\end{proof}

\subsection{Twierdzenie Newtona--Leibniza i rachunek z parametrem}

\begin{twierdzenie}[Podstawowe twierdzenie rachunku całkowego]
\label{thm:ftc-lcs15a}
Niech $E$ będzie zupełną ciągowo przestrzenią lokalnie wypukłą
Hausdorffa, a $f:[a,b]\to E$ funkcją ciągłą. Wtedy funkcja
$F:[a,b]\to E$, $F(x)=\int_a^x f(t)\,\mathrm dt$, jest klasy
$C^1$ na $(a,b)$, ma jednostronne pochodne na końcach i
$F'(x)=f(x)$ (z odpowiednimi pochodnymi jednostronnymi).
Jeżeli $G:[a,b]\to E$ jest klasy $C^1$, to
\[
 G(b)-G(a)=\int_a^b G'(t)\,\mathrm dt.
\]
W szczególności jest to uogólnienie wyniku z \hyperref[sec:gateaux-frechet14]{rozdziału~14},
a nie jego ponowny dowód w przypadku przestrzeni Banacha.
\end{twierdzenie}
\begin{proof}
Dla $x,x+h\in[a,b]$ i $h\ne0$ z addytywności całki mamy
$h^{-1}(F(x+h)-F(x))-f(x)
=h^{-1}\int_x^{x+h}(f(t)-f(x))\,\mathrm dt$.
Oszacowanie seminormowe (również dla $h<0$ po zmianie
orientacji całki) daje
\[
 p\left(h^{-1}(F(x+h)-F(x))-f(x)\right)
 \le\sup_{|t-x|\le|h|}p(f(t)-f(x))\longrightarrow0.
\]
Stąd pochodna istnieje w topologii $E$ i jest ciągła.

Dla drugiej tezy nie musimy odtwarzać dowodu z przestrzeni Banacha.
Dla każdego $\ell\in E'$ złożenie $\ell\circ G$ jest skalarne
klasy $C^1$ z pochodną $\ell(G'(t))$. Skalarne twierdzenie
Newtona--Leibniza i przemienność całki z $\ell$ dają
$\ell(G(b)-G(a))=\ell(\int_a^b G'(t)\,\mathrm dt)$.
Ponieważ $E'$ rozdziela punkty, zachodzi żądana równość w $E$.
\end{proof}

\begin{twierdzenie}[Różniczkowanie całki zależnej od parametru]
\label{thm:parametric-integral15a}
Niech $E$ będzie zupełną ciągowo przestrzenią lokalnie wypukłą
Hausdorffa, $J\subseteq\R$ otwartym przedziałem i
$f:[a,b]\times J\to E$ funkcją ciągłą. Załóżmy, że dla
każdego $(t,s)\in[a,b]\times J$ istnieje pochodna
$\partial_s f(t,s)$ w topologii $E$ i że funkcja
$(t,s)\mapsto\partial_s f(t,s)$ jest ciągła.
Wtedy $F(s)=\int_a^b f(t,s)\,\mathrm dt$ jest funkcją klasy $C^1$
o wartościach w $E$ oraz
$F'(s)=\int_a^b\partial_s f(t,s)\,\mathrm dt$ dla $s\in J$.
\end{twierdzenie}
\begin{proof}
Ciągłość $F$ wynika z jednostajnej, względem $t\in[a,b]$,
ciągłości $f(t,s)$ w sąsiedztwie ustalonego parametru.
Dla ustalonego $s$ i małego $h$ skalarne twierdzenie
Newtona--Leibniza, zastosowane do $\ell\circ f(t,\cdot)$ dla
$\ell\in E'$, wraz z twierdzeniem Hahna--Banacha, daje dla
każdej ciągłej seminormy $p$ oszacowanie
\[
 p\left(\frac{f(t,s+h)-f(t,s)}h-\partial_s f(t,s)\right)
 \le\sup_{u\text{ między }s\text{ i }s+h}
 p(\partial_s f(t,u)-\partial_s f(t,s)).
\]
Po scałkowaniu względem $t$ i zastosowaniu oszacowania
seminormowego prawa strona jest nie większa od
$(b-a)$ razy supremum po $t\in[a,b]$ i $u$ między $s$ i $s+h$.
Ciągłość $\partial_s f$ na zwartym prostokącie wokół
$[a,b]\times\{s\}$ sprawia, że supremum dąży do zera.
Daje to pochodną $F'$ o wskazanej postaci; jej ciągłość
wynika z tego samego oszacowania seminormowego.
\end{proof}

\begin{twierdzenie}[Twierdzenie Fubiniego dla funkcji ciągłych]
\label{thm:fubini-continuous15a}
Niech $E$ będzie zupełną ciągowo przestrzenią lokalnie wypukłą
Hausdorffa i $f:[a,b]\times[c,d]\to E$ funkcją ciągłą,
gdzie $a<b$ i $c<d$. Wówczas obie funkcje
$s\mapsto\int_c^d f(s,t)\,\mathrm dt$ oraz
$t\mapsto\int_a^b f(s,t)\,\mathrm ds$ są ciągłe i
\[
 \int_a^b\left(\int_c^d f(s,t)\,\mathrm dt\right)\mathrm ds
 =\int_c^d\left(\int_a^b f(s,t)\,\mathrm ds\right)\mathrm dt.
\]
Nie zakładamy mierzalności Bochnerowskiej ani twierdzenia Fubiniego
dla ogólnych funkcji mierzalnych.
\end{twierdzenie}
\begin{proof}
Dla każdej seminormy $p$ jednostajna ciągłość $f$
na zwartym prostokącie daje ciągłość obydwu funkcji
całek iterowanych: na przykład
$p(\int_c^d(f(s,t)-f(s_0,t))\,\mathrm dt)
\le(d-c)\sup_{t\in[c,d]}p(f(s,t)-f(s_0,t))$.

Dla podziałów $P$ odcinka $[a,b]$ i $Q$ odcinka $[c,d]$
suma Riemanna funkcji
$s\mapsto\int_c^d f(s,t)\,\mathrm dt$ różni się od sumy
podwójnej $\sum_{j,k}\Delta s_j\Delta t_k f(\xi_j,\eta_k)$
o co najwyżej $(b-a)(d-c)\omega_p(|Q|)$ w seminormie $p$,
gdzie $\omega_p$ jest modułem ciągłości względem drugiej
zmiennej, jednostajnym względem pierwszej. Analogiczna
nierówność dotyczy drugiej kolejności całkowania.
Przepuszczając średnice obu podziałów do zera, otrzymujemy
identyczne granice wszystkich podwójnych sum. Z
hausdorffowskości $E$ wynika równość iterowanych całek.
\end{proof}

\subsection{Przykłady funkcyjne i całkowanie rodzin operatorów}

\begin{przyklad}[Całkowanie w $C^\infty(\Omega)$]
\label{ex:c-infty-integral15a}
Niech $\Omega\subset\R^d$ będzie niepustym zbiorem otwartym.
Przestrzeń $C^\infty(\Omega)$ wyposażamy w seminormy
$p_{K,m}(u)=\max_{|\alpha|\le m}\sup_{x\in K}|\partial^\alpha u(x)|$,
gdzie $K\Subset\Omega$ jest zwarty, a $m\in\N_0$.
Jest to zupełna przestrzeń lokalnie wypukła, co szczegółowo
udowodnimy w \hyperref[sec:gladkie-schwartz17]{rozdziale~17}. Jeżeli $t\mapsto u_t$ jest
ciągłe z $[a,b]$ do tej przestrzeni, to
$u=\int_a^b u_t\,\mathrm dt\in C^\infty(\Omega)$ i
\[
 \partial^\alpha u(x)
 =\int_a^b\partial^\alpha u_t(x)\,\mathrm dt
 \qquad(x\in\Omega,\ \alpha\in\N_0^d).
\]
Wynika to z ciągłości operatora $\partial^\alpha$ na
$C^\infty(\Omega)$ i ciągłości ewaluacji w punkcie $x$,
bez nieuzasadnionego różniczkowania pod całką skalarną.
\end{przyklad}

\begin{przyklad}[Całkowanie w przestrzeni Schwartza]
\label{ex:schwartz-integral15a}
Niech $\mathcal S(\R^d)$ oznacza przestrzeń funkcji gładkich,
dla których $q_{m,k}(u)=\max_{|\alpha|\le k}\sup_{x\in\R^d}
(1+|x|)^m|\partial^\alpha u(x)|$ jest skończone
przy wszystkich $m,k\in\N_0$; tę topologię również
omówimy w \hyperref[sec:gladkie-schwartz17]{rozdziale~17}. Przestrzeń ta jest zupełna.
Dla ciągłego $t\mapsto u_t\in\mathcal S(\R^d)$ całka
$u=\int_a^b u_t\,\mathrm dt$ należy do $\mathcal S(\R^d)$
i spełnia $q_{m,k}(u)\le\int_a^b q_{m,k}(u_t)\,\mathrm dt$.
Różniczkowanie jest na tej przestrzeni ciągłym
odwzorowaniem liniowym: dla każdego wielowskaźnika $\alpha$
seminorma $q_{m,k}(\partial^\alpha u)$ jest ograniczona
przez $q_{m,k+|\alpha|}(u)$. Stąd
$\partial^\alpha(\int_a^b u_t\,\mathrm dt)
=\int_a^b\partial^\alpha u_t\,\mathrm dt$
(jako równość w przestrzeni Schwartza).
\end{przyklad}

\begin{twierdzenie}[Silna całka rodziny operatorów]
\label{thm:strong-operator-integral15a}
Niech $X,Y$ będą przestrzeniami Banacha i
$T:[a,b]\to\mathcal L(X,Y)$ rodziną operatorów taką, że
$t\mapsto T(t)x$ jest ciągła w $Y$ dla każdego $x\in X$.
Wtedy wzór
$Ax=\int_a^b T(t)x\,\mathrm dt$ określa operator
$A\in\mathcal L(X,Y)$, a
\[
 \|A\|\le(b-a)\sup_{t\in[a,b]}\|T(t)\|<\infty.
\]
Nie wynika stąd ciągłość $t\mapsto T(t)$ w normie operatorowej.
\end{twierdzenie}
\begin{proof}
Dla każdego $x$ zbiór $\{T(t)x:t\in[a,b]\}$ jest ograniczony,
gdyż jest obrazem zwartym przez funkcję ciągłą. Twierdzenie
Banacha--Steinhausa zapewnia
$M=\sup_{t\in[a,b]}\|T(t)\|<\infty$.
Istnienie całki wektorowej w przestrzeni Banacha $Y$
wynika z twierdzenia o istnieniu, a liniowość z
liniowości całki. Dla $x\in X$ oszacowanie seminormowe
daje $\|Ax\|\le\int_a^b\|T(t)x\|\,\mathrm dt
\le(b-a)M\|x\|$, co dowodzi ciągłości $A$.
\end{proof}

\begin{przyklad}[Ciągłość silna nie oznacza ciągłości normowej]
\label{ex:strong-not-norm-cont15a}
W przestrzeni Hilberta $H=L^2(\R,\lambda_1)$ niech
$(U_t f)(x)=e^{itx}f(x)$ dla $f\in H$ i $t\in\R$.
Dla każdego $f\in H$ mamy $\|U_tf-f\|_2\to0$ przy
$t\to0$ na mocy skalarnej zbieżności zmajoryzowanej zastosowanej
do $|e^{itx}-1|^2|f(x)|^2$.
Jednak $\|U_t-I_H\|=\operatorname*{ess\,sup}_{x\in\R}
|e^{itx}-1|=2$ dla każdego $t\ne0$.
Można więc całkować wektory $U_tf$ względem $t$ w $H$,
ale nie wolno bez dodatkowego dowodu zakładać ciągłości
$t\mapsto U_t$ w normie $\mathcal L(H)$.
\end{przyklad}

\begin{uwaga}[Granice zakresu]
Definicja sum Riemanna wymaga jedynie topologii i działań liniowych,
ale powyższe twierdzenia o istnieniu i oszacowaniach wykorzystują
lokalną wypukłość, seminormy i zupełność ciągową. W przestrzeni
niezupełnej ciągłość nie gwarantuje istnienia całki; w ogólnej
przestrzeni liniowo-topologicznej nielokalnie wypukłej nie można
bez uzasadnienia stosować oszacowań przez rodzinę seminorm.
Dla funkcji mierzalnych konieczne są dodatkowe pojęcia mierzalności;
teorię tę rozwiniemy osobno w \hyperref[sec:integral-vector15b]{module~15B}.
\end{uwaga}

\addtocounter{section}{-1}
\endgroup


% Wersja źródłowa: audyt II, jednostka 15B
\setcounter{section}{15}
\noindent\textbf{Wymagane wyniki.} Wykorzystujemy podstawy skalarnej całki Lebesgue'a, twierdzenia o zbieżności zmajoryzowanej i Fubiniego, twierdzenia Hahna--Banacha i o wykresie domkniętym oraz definicje przestrzeni Banacha i jej przestrzeni dualnej.\medskip
% \hyperref[sec:integral-vector15b]{Moduł 15B}; umieszczony po 15A, bez przesuwania \hyperref[sec:wariacyjny16]{rozdziałów 16}--33.
\begingroup
\renewcommand{\thesection}{15B}
\renewcommand{\theHsection}{15B}
\section{Całki Bochnera, Pettisa, Dunforda i Gelfanda oraz miary wektorowe}
\label{sec:integral-vector15b}

W \hyperref[sec:integral-lcs15a]{module 15A} całkowaliśmy \emph{funkcje ciągłe na odcinku} w topologii
przestrzeni lokalnie wypukłej. Teraz dopuszczamy funkcje mierzalne;
różne pojęcia całki wynikają z tego, czy granica istnieje w normie,
czy najpierw całkujemy wszystkie obrazy skalarne, i w jakiej przestrzeni
musi leżeć wektor reprezentujący te całki. Nie powtarzamy konstrukcji
skalarnej całki Lebesgue'a.

W całym module $(S,\Sigma,\mu)$ jest \emph{zupełną przestrzenią z miarą
skończoną}, $\K\in\{\R,\C\}$, zaś $X$ jest przestrzenią Banacha nad
$\K$, chyba że jawnie przyjmiemy inne założenia. Przez $X'$ i $X''$
oznaczamy odpowiednio przestrzeń dualną i bidual; $J_X:X\to X''$
jest zanurzeniem kanonicznym. Równości funkcji mierzalnych rozumiemy
z dokładnością do zbioru miary zero. Dla $A\in\Sigma$ zapis
$\int_A f(s)\,\mathrm d\mu(s)$ oznacza $\int_S\chi_A(s)f(s)\,\mathrm d\mu(s)$.

\subsection{Mierzalność silna i całka Bochnera}

\begin{definicja}[Mierzalność silna i słaba]
\label{def:measurability15b}
Funkcję $f:S\to X$ nazywamy \emph{silnie mierzalną}, jeżeli istnieją
funkcje proste $s_n=\sum_{j=1}^{m_n}\chi_{A_{n,j}}x_{n,j}$,
$A_{n,j}\in\Sigma$, $x_{n,j}\in X$, takie że $s_n(s)\to f(s)$
w normie dla prawie każdego $s\in S$. Nazywamy ją
\emph{słabo mierzalną}, jeżeli $x'\circ f:S\to\K$ jest
mierzalna dla każdego $x'\in X'$. Funkcję $f$ nazywamy
\emph{istotnie ośrodkowo-wartościową}, jeśli poza pewnym zbiorem
miary zero jej wartości leżą w ośrodkowej podprzestrzeni domkniętej
przestrzeni $X$.
\end{definicja}

\begin{twierdzenie}[Twierdzenie Pettisa o mierzalności]
\label{thm:pettis-meas15b}
Funkcja $f:S\to X$ jest silnie mierzalna wtedy i tylko wtedy, gdy
jest słabo mierzalna oraz istotnie ośrodkowo-wartościowa.
\end{twierdzenie}
\begin{proof}
Jeżeli $s_n\to f$ prawie wszędzie, to dla każdego $x'\in X'$
mamy $x'(s_n)\to x'(f)$, co daje mierzalność słabą.
Wartości wszystkich $s_n$ tworzą zbiór przeliczalny;
jego domknięta otoczka liniowa jest ośrodkową podprzestrzenią,
w której leży $f(s)$ poza wspólnym zbiorem miary zero.

Odwrotnie, po usunięciu zbioru miary zero można przyjąć, że $f$
przyjmuje wartości w ośrodkowej przestrzeni Banacha $Y\subset X$.
Wybierzmy gęsty ciąg $(y_j)$ w $Y$. Istnieje przeliczalna
rodzina $(\ell_k)$ elementów $X'$ o normie co najwyżej $1$
taka, że dla każdego $y\in Y$ zachodzi
$\|y\|=\sup_k|\ell_k(y)|$: otrzymujemy ją, dobierając
funkcjonały normujące gęsty zbiór sfery jednostkowej $Y$
i przedłużając je za pomocą twierdzenia Hahna--Banacha.
W konsekwencji $s\mapsto\|f(s)-y_j\|$ jest mierzalna,
ponieważ jest supremum przeliczalnej rodziny mierzalnych
funkcji skalarnych.

Dla $n\ge1$ definiujemy rozłączne, mierzalne zbiory $B_{n,j}$,
przypisując $s$ pierwszemu indeksowi $j$, dla którego
$\|f(s)-y_j\|<1/n$. Zbiory te pokrywają $S$ z dokładnością
do zbioru miary zero, a funkcja przyjmująca wartość $y_j$
na $B_{n,j}$ przybliża $f$ z błędem mniejszym od $1/n$.
Ponieważ $\mu(S)<\infty$, można wybrać $m_n$ tak, aby
$\mu(\bigcup_{j>m_n}B_{n,j})<2^{-n}$, i przyjąć wartość
$0$ na tym ogonie. Powstaje funkcja prosta $s_n$;
z lematu Borela--Cantellego otrzymujemy $s_n(s)\to f(s)$
prawie wszędzie. Jest to mierzalność silna.
\end{proof}

\begin{definicja}[Całka Bochnera]
\label{def:bochner15b}
Funkcję silnie mierzalną $f:S\to X$ nazywamy
\emph{całkowalną w sensie Bochnera}, jeśli istnieje ciąg
funkcji prostych $(s_n)$ taki, że
$\int_S\|f(s)-s_n(s)\|\,\mathrm d\mu(s)\to0$.
Wtedy definiujemy
$\int_S f(s)\,\mathrm d\mu(s)
=\lim_{n\to\infty}\int_S s_n(s)\,\mathrm d\mu(s)$,
gdzie całka funkcji prostej $\sum_j\chi_{A_j}x_j$ jest
równa $\sum_j\mu(A_j)x_j$ dla rozłącznych zbiorów $A_j$.
\end{definicja}

\begin{twierdzenie}[Kryterium całkowalności i oszacowanie]
\label{thm:bochner-criterion15b}
Funkcja $f:S\to X$ jest całkowalna w sensie Bochnera wtedy
i tylko wtedy, gdy jest silnie mierzalna i
$\int_S\|f(s)\|\,\mathrm d\mu(s)<\infty$.
Całka nie zależy od ciągu przybliżeń, jest liniowa i spełnia
\[
 \left\|\int_A f(s)\,\mathrm d\mu(s)\right\|
 \le\int_A\|f(s)\|\,\mathrm d\mu(s)\qquad(A\in\Sigma).
\]
Dla każdej przestrzeni Banacha $Y$ i operatora
$T\in\mathcal L(X,Y)$ zachodzi
$T(\int_A f(s)\,\mathrm d\mu(s))=\int_A T(f(s))\,\mathrm d\mu(s)$.
\end{twierdzenie}
\begin{proof}
Dla funkcji prostych $h$ mamy
$\|\int_A h(s)\,\mathrm d\mu(s)\|\le\int_A\|h(s)\|\,\mathrm d\mu(s)$.
Zatem jeżeli $s_n\to f$ w $L^1(S;X)$, całki $s_n$ są
ciągiem Cauchy'ego w zupełnej przestrzeni $X$, a ich granica
nie zależy od aproksymacji. To dowodzi też liniowości i
nierówności po przejściu do granicy. Ponadto
$|\|f\|-\|s_n\||\le\|f-s_n\|$, stąd całkowalność normy.

Jeśli z kolei $f$ jest silnie mierzalna i ma całkowalną normę,
wybierzmy funkcje proste $u_n	o f$ prawie wszędzie.
Funkcja $s_n$ jest równa $u_n$ tam, gdzie
$\|u_n(s)\|\le2\|f(s)\|$, a zeru w pozostałych punktach.
Oba zbiory są mierzalne, zatem $s_n$ jest funkcją prostą,
$\|s_n\|\le2\|f\|$, a $s_n	o f$ prawie wszędzie:
dla $f(s)
e0$ nierówność definiująca $s_n$ jest
spełniona od pewnego miejsca, a dla $f(s)=0$ mamy
$s_n(s)=0$. Stąd $\|s_n-f\|\le3\|f\|$ i skalarne
twierdzenie o zbieżności zmajoryzowanej zapewnia
$\int_S\|f(s)-s_n(s)\|\,\mathrm d\mu(s)\to0$.
Dla operatora $T$ każda funkcja $Ts_n$ jest prosta,
$\int\|Tf-Ts_n\|\le\|T\|\int\|f-s_n\|$, a równość
$T\int s_n=\int Ts_n$ przechodzi do granicy.
\end{proof}

\begin{twierdzenie}[Zbieżność zmajoryzowana Bochnera]
\label{thm:bochner-dct15b}
Niech $f_n:S\to X$ będą silnie mierzalne,
$f_n(s)\to f(s)$ w normie dla prawie każdego $s\in S$.
Jeżeli istnieje $g\in L^1(S,\mu)$, $g\ge0$, taka że
$\|f_n(s)\|\le g(s)$ prawie wszędzie dla wszystkich $n$,
to $f$ i $f_n$ są całkowalne w sensie Bochnera oraz
$\int_S\|f_n(s)-f(s)\|\,\mathrm d\mu(s)\to0$, a więc
$\int_S f_n(s)\,\mathrm d\mu(s)\to\int_S f(s)\,\mathrm d\mu(s)$ w normie.
\end{twierdzenie}
\begin{proof}
Granica prawie wszędzie funkcji silnie mierzalnych jest
silnie mierzalna: wystarczy dobrać prostą aproksymację
każdego $f_n$ tak, aby przekątna aproksymowała $f$ prawie
wszędzie. Ponadto $\|f\|\le g$ i
$\|f_n-f\|\le2g$ prawie wszędzie. Skalarne twierdzenie
o zbieżności zmajoryzowanej i kryterium całkowalności
Bochnera kończą dowód.
\end{proof}

\begin{twierdzenie}[Twierdzenie Fubiniego dla całki Bochnera]
\label{thm:bochner-fubini15b}
Niech $(S,\Sigma,\mu)$ i $(T,\mathcal T,\nu)$ będą
zupełnymi przestrzeniami z miarami skończonymi, a
$f:S\times T\to X$ niech będzie silnie mierzalna względem
uzupełnienia miary produktowej i spełnia
$\int_{S\times T}\|f(s,t)\|\,\mathrm d(\mu\otimes\nu)(s,t)<\infty$.
Wtedy prawie wszystkie przekroje są całkowalne w sensie
Bochnera; całki iterowane istnieją i
\[
 \int_{S\times T}f\,\mathrm d(\mu\otimes\nu)
 =\int_S\left(\int_Tf(s,t)\,\mathrm d\nu(t)\right)\mathrm d\mu(s)
 =\int_T\left(\int_Sf(s,t)\,\mathrm d\mu(s)\right)\mathrm d\nu(t).
\]
\end{twierdzenie}
\begin{proof}
Najpierw sprawdzamy równości dla funkcji prostych postaci
$\chi_{A\times B}(s,t)x$. Z liniowości zachodzą dla
skończonych sum takich funkcji. Ich rozpiętość jest gęsta
w $L^1(S\times T;X)$: prostokąty generują produktową
$\sigma$-algebrę, a standardowy argument klasy monotonicznej
przybliża wskaźniki zbiorów mierzalnych w $L^1$;
uzupełnienie miary nie zmienia klas funkcji prawie wszędzie.
Wybierając aproksymacje $f_n$ zbieżne do $f$ w $L^1$,
z oszacowania Bochnera i skalarnego twierdzenia Fubiniego
dla $\|f_n-f_m\|$ wynika, że całki po przekrojach są
ciągami Cauchy'ego w odpowiednich przestrzeniach $L^1$.
Wybór podciągu z sumowalnymi błędami zapewnia zbieżność
dla prawie wszystkich przekrojów, a przejście do granicy
w obu całkach iterowanych daje tezę.
\end{proof}

\begin{stwierdzenie}[Zgodność z całką Riemanna]
\label{prop:riemann-bochner15b}
Jeśli $f:[a,b]\to X$ jest ciągła, to jej całka Riemanna
w przestrzeni Banacha z \hyperref[sec:integral-lcs15a]{modułu 15A} jest równa całce Bochnera
względem miary Lebesgue'a na $[a,b]$. W drugą stronę
całkowalność Bochnera nie implikuje całkowalności Riemanna.
\end{stwierdzenie}
\begin{proof}
Ciągłość na zwartym odcinku daje jednostajną ciągłość.
Funkcje proste, stałe na odcinkach drobnego podziału,
przybliżają $f$ jednostajnie; ich całki Bochnera są
odpowiednimi sumami Riemanna, a zbieżność zachodzi w normie.
Kontrprzykładem jest $f=\chi_{\mathbb Q\cap[0,1]}x_0$
dla niezerowego $x_0\in X$: funkcja jest równa zeru prawie
wszędzie, więc ma zerową całkę Bochnera, lecz nie jest
całkowalna w sensie Riemanna.
\end{proof}

\subsection{Całki wyznaczane przez wartości skalarne}

\begin{definicja}[Całka Dunforda i całka Pettisa]
\label{def:dunford-pettis15b}
Niech $f:S\to X$ będzie słabo mierzalna i niech
$x'\circ f\in L^1(S,\mu)$ dla każdego $x'\in X'$.
\emph{Całką Dunforda} funkcji $f$ na $A\in\Sigma$ nazywamy
element $D_A\in X''$ określony równością
$D_A(x')=\int_Ax'(f(s))\,\mathrm d\mu(s)$.
Funkcja $f$ jest \emph{całkowalna w sensie Pettisa}, gdy
dla każdego $A\in\Sigma$ istnieje $x_A\in X$ spełniający
$x'(x_A)=\int_Ax'(f(s))\,\mathrm d\mu(s)$ dla każdego
$x'\in X'$. Wtedy $x_A=\int_A^{\mathrm P} f(s)\,\mathrm d\mu(s)$.
Nie wystarcza wymagać reprezentacji tylko dla $A=S$.
\end{definicja}

\begin{twierdzenie}[Istnienie całki Dunforda]
\label{thm:dunford-closedgraph15b}
Przy założeniach powyższej definicji wszystkie $D_A$
istnieją jako elementy $X''$, nawet gdy $f$ nie jest
całkowalna w sensie Bochnera. Odwzorowanie
$Q_f:X'\to L^1(S,\mu)$, $Q_f(x')=x'\circ f$, jest
liniowe i ograniczone, a
$\|D_A\|\le\|Q_f\|$ dla każdego $A\in\Sigma$.
\end{twierdzenie}
\begin{proof}
Jeżeli $x_n'\to x'$ w normie $X'$ i
$Q_f(x_n')\to h$ w $L^1$, to pewien podciąg
$Q_f(x_{n_k}')$ zbiega do $h$ prawie wszędzie.
Równocześnie dla każdego $s\in S$ mamy
$x_{n_k}'(f(s))\to x'(f(s))$, więc $h=Q_f(x')$ prawie
wszędzie. Wykres operatora $Q_f$ jest zatem domknięty.
Twierdzenie o wykresie domkniętym dla przestrzeni Banacha
daje $\|Q_f\|<\infty$. Stąd
$|D_A(x')|\le\int_A|x'(f(s))|\,\mathrm d\mu(s)
\le\|Q_f\|\|x'\|$, co definiuje element bidualu.
\end{proof}

\begin{twierdzenie}[Hierarchia trzech całek]
\label{thm:hierarchy-integrals15b}
Jeżeli $f:S\to X$ jest całkowalna w sensie Bochnera,
to jest całkowalna w sensie Pettisa i Dunforda, a dla
każdego $A\in\Sigma$ zachodzi
$D_A=J_X(\int_A f(s)\,\mathrm d\mu(s))$.
Jeżeli $f$ jest całkowalna w sensie Pettisa, to dla każdego
$A$ jej całka Dunforda należy do $J_X(X)$.
Odwrotne implikacje są w ogólności fałszywe.
\end{twierdzenie}
\begin{proof}
Dla funkcji prostej równość
$x'(\int_A f(s)\,\mathrm d\mu(s))=\int_Ax'(f(s))\,\mathrm d\mu(s)$
jest bezpośrednia. Przybliżenie w $L^1(S;X)$
i $|x'(v)|\le\|x'\|\|v\|$ pozwalają przejść do granicy
w przypadku Bochnera. Druga teza jest dokładnie definicją
całki Pettisa i zanurzenia $J_X$. Zapowiedziane
kontrprzykłady podajemy niżej.
\end{proof}

\begin{przyklad}[Całka Dunforda poza przestrzenią Banacha]
\label{ex:dunford-not-pettis15b}
Niech $S=\N=\{1,2,\ldots\}$ i
$\mu(\{n\})=2^{-n}$. W przestrzeni Banacha $X=c_0$
przyjmijmy $f(n)=2^ne_n$, gdzie $(e_n)$ jest standardową
bazą jednostkową $c_0$. Ponieważ $X'=\ell^1$, dla
$a=(a_n)\in\ell^1$ mamy
$\int_S|a(f(n))|\,\mathrm d\mu(n)
=\sum_{n\ge1}|a_n|<\infty$.
Całka Dunforda ma zatem współrzędne
$D_S=(1,1,\ldots)\in\ell^\infty=X''$ i nie leży w
$J_X(c_0)$. Funkcja nie jest całkowalna w sensie Pettisa
ani Bochnera; co więcej,
$\int_S\|f(n)\|\,\mathrm d\mu(n)=\sum_{n\ge1}1=\infty$.
\end{przyklad}

\begin{przyklad}[Całka Pettisa bez całki Bochnera]
\label{ex:pettis-not-bochner15b}
Niech $S=[0,1]$ z miarą Lebesgue'a, zaś
$X=\ell^2([0,1])$. Dla $t\in[0,1]$ definiujemy $f(t)=e_t$.
Każdy $a\in\ell^2([0,1])$ ma co najwyżej przeliczalnie
wiele niezerowych współrzędnych, dlatego funkcja
$t\mapsto\langle f(t),a\rangle$ jest równa zeru prawie
wszędzie. Całki skalarne po każdym $A\in\Sigma$ znikają,
więc $f$ ma całkę Pettisa równą $0$ na każdym $A$.
Nie jest jednak silnie mierzalna: zbiór jej wartości poza
jakimkolwiek zbiorem miary zero jest nieprzeliczalną rodziną
punktów odległych o $\sqrt 2$, a więc nie może leżeć
w ośrodkowej podprzestrzeni $X$. Całka Bochnera nie istnieje,
choć $\|f(t)\|=1$ dla każdego $t$.
\end{przyklad}

\begin{uwaga}[Refleksywność i całka Dunforda]
Jeżeli $X$ jest refleksywna, to $J_X(X)=X''$ i każdy
$D_A$ jest reprezentowany przez wektor z $X$. Zatem każda
słabo mierzalna, skalarnie całkowalna funkcja do refleksywnej
przestrzeni Banacha jest całkowalna w sensie Pettisa.
Nie oznacza to całkowalności Bochnera; przykład z
przestrzeni Hilberta $\ell^2([0,1])$ pozostaje ważny.
\end{uwaga}

\begin{twierdzenie}[Zbieżność zmajoryzowana dla całki Pettisa]
\label{thm:pettis-convergence15b}
Niech $X$ będzie refleksywną przestrzenią Banacha, a
$f_n:S\to X$ słabo mierzalnymi funkcjami takimi, że
$f_n(s)\to f(s)$ w normie dla prawie każdego $s$.
Załóżmy, że $\|f_n(s)\|\le g(s)$ dla wszystkich $n$
i prawie wszystkich $s$, gdzie $g\in L^1(S,\mu)$,
$g\ge0$, oraz że funkcje
$s\mapsto\|f_n(s)-f(s)\|$ są mierzalne.
Wtedy $f_n$ i $f$ są całkowalne w sensie Pettisa i dla
każdego $A\in\Sigma$ zachodzi
$\int_A^{\mathrm P}f_n(s)\,\mathrm d\mu(s)
\to\int_A^{\mathrm P}f(s)\,\mathrm d\mu(s)$ w normie $X$.
\end{twierdzenie}
\begin{proof}
Funkcja $f$ jest słabo mierzalna, bo dla każdego $x'\in X'$
jej skalaryzacja jest granicą mierzalnych $x'\circ f_n$.
Dominacja przez $g$ daje całkowalność skalarną $f_n$ i $f$.
Z refleksywności $X$ oraz twierdzenia o istnieniu całki
Dunforda wynika, że wszystkie te funkcje mają całkę Pettisa.
Dla każdego $A\in\Sigma$ i $x'\in X'$ o normie co najwyżej
$1$ zachodzi
$|x'(\int_A^{\mathrm P}(f_n(s)-f(s))\,\mathrm d\mu(s))|
\le\int_A\|f_n(s)-f(s)\|\,\mathrm d\mu(s)$.
Biorąc supremum po takich $x'$ otrzymujemy oszacowanie
normy różnicy całek. Skalarne twierdzenie o zbieżności
zmajoryzowanej stosuje się do mierzalnych funkcji
$\|f_n-f\|\le2g$; jego prawa strona dąży do zera.
\end{proof}

\subsection{Całka Gelfanda i miary wektorowe}

\begin{definicja}[Całka Gelfanda]
\label{def:gelfand15b}
Niech $Y$ będzie przestrzenią Banacha, a
$f:S\to Y'$ funkcją \emph{słabo-* mierzalną}, tzn.
$s\mapsto f(s)(y)$ jest mierzalna dla każdego $y\in Y$.
Zakładamy, że $f(\cdot)(y)\in L^1(S,\mu)$ dla każdego
$y\in Y$. \emph{Całką Gelfanda} na $A\in\Sigma$ nazywamy
jedyny $g_A\in Y'$ spełniający
$g_A(y)=\int_Af(s)(y)\,\mathrm d\mu(s)$ dla wszystkich
$y\in Y$. Stosujemy oznaczenie
$g_A=\int_A^{\mathrm G}f(s)\,\mathrm d\mu(s)$.
\end{definicja}

\begin{twierdzenie}[Istnienie całki Gelfanda i jej addytywność]
\label{thm:gelfand-exists15b}
Przy założeniach definicji całka Gelfanda istnieje dla
każdego $A\in\Sigma$. Odwzorowanie
$A\mapsto g_A\in Y'$ jest przeliczalnie addytywne
w topologii słabej-* $\sigma(Y',Y)$; bez dodatkowych
założeń nie musi być przeliczalnie addytywne w normie $Y'$.
\end{twierdzenie}
\begin{proof}
Operator $Q:Y\to L^1(S,\mu)$, $(Qy)(s)=f(s)(y)$,
jest domknięty. Istotnie, z $y_n\to y$ i $Qy_n\to h$
w $L^1$ wybieramy podciąg zbieżny prawie wszędzie;
punktowa ciągłość każdego $f(s)\in Y'$ daje $h=Qy$.
Z twierdzenia o wykresie domkniętym $Q$ jest ograniczony.
Zatem $|\int_Af(s)(y)\,\mathrm d\mu(s)|\le\|Q\|\|y\|$,
a wzór definiuje $g_A\in Y'$.
Dla parami rozłącznych $A_n$ i $y\in Y$ skalarna
addytywność całki Lebesgue'a daje
$g_{\bigcup_nA_n}(y)=\sum_ng_{A_n}(y)$.
Jest to dokładnie zbieżność w topologii $\sigma(Y',Y)$.
\end{proof}

\begin{przyklad}[Słaba-* addytywność bez addytywności normowej]
\label{ex:gelfand-weakstar15b}
Niech $Y=\ell^1$, $Y'=\ell^\infty$, a na $S=\N$ przyjmijmy
$\mu(\{n\})=2^{-n}$ oraz $f(n)=2^ne_n\in\ell^\infty$.
Dla $y=(y_n)\in\ell^1$ zachodzi
$\int_S|f(n)(y)|\,\mathrm d\mu(n)=\sum_n|y_n|<\infty$.
W konsekwencji $g_A=\chi_A\in\ell^\infty$, gdzie
$\chi_A=(\chi_A(n))_{n\ge1}$. Dla $A_m=\{m,m+1,\ldots\}$
mamy $\|g_{A_m}\|_\infty=1$ przy każdym $m$, mimo że
$g_{A_m}\to0$ w topologii $\sigma(\ell^\infty,\ell^1)$.
\end{przyklad}

\begin{definicja}[Miara wektorowa i całkowita wariacja]
\label{def:vector-measure15b}
Odwzorowanie $m:\Sigma\to X$ nazywamy \emph{miarą wektorową},
jeśli $m(\emptyset)=0$ i dla każdego ciągu parami rozłącznych
$A_n\in\Sigma$ zachodzi
$m(\bigcup_{n\ge1}A_n)=\lim_{N\to\infty}\sum_{n=1}^Nm(A_n)$
w normie $X$. Jej \emph{całkowitą wariacją} jest miara
dodatnia o wartościach w $[0,\infty]$ określona przez
$|m|(A)=\sup\{\sum_{j=1}^k\|m(A_j)\|:\ A_j\text{ są
parami rozłącznymi zbiorami mierzalnymi, }
\bigcup_{j=1}^kA_j=A\}$.
\end{definicja}

\begin{twierdzenie}[Miara generowana przez całkę Bochnera]
\label{thm:bochner-measure15b}
Jeśli $f:S\to X$ jest całkowalna w sensie Bochnera, to
$m_f(A)=\int_Af(s)\,\mathrm d\mu(s)$ jest miarą wektorową,
a jej całkowita wariacja spełnia
\[
 |m_f|(A)=\int_A\|f(s)\|\,\mathrm d\mu(s)
 \qquad(A\in\Sigma).
\]
W szczególności $|m_f|$ jest skończoną miarą bezwzględnie
ciągłą względem $\mu$.
\end{twierdzenie}
\begin{proof}
Dla parami rozłącznych $A_n$ i $A=\bigcup_nA_n$
z oszacowania całki otrzymujemy
$\|m_f(A)-\sum_{n=1}^Nm_f(A_n)\|
\le\int_{\bigcup_{n>N}A_n}\|f(s)\|\,\mathrm d\mu(s)\to0$.
Dla każdej skończonej partycji $A$ suma norm całek
jest ograniczona przez $\int_A\|f\|$, czyli
$|m_f|(A)\le\int_A\|f\|$.
Gdy $f=\sum_j\chi_{B_j}v_j$ jest prosta o rozłącznych
zbiorach poziomicowych, partycja $A$ przez $A\cap B_j$
daje równość. Dla ogólnej $f$ wybierzmy funkcje proste
$f_n\to f$ w $L^1$. Z nierówności
$\bigl||m_f|(A)-|m_{f_n}|(A)\bigr|
\le\int_A\|f(s)-f_n(s)\|\,\mathrm d\mu(s)$
(wynikającej z definicji wariacji i nierówności trójkąta)
oraz zbieżności norm w $L^1$ otrzymujemy równość.
\end{proof}

\begin{definicja}[Własność Radona--Nikodyma]
\label{def:rnp15b}
Przestrzeń Banacha $X$ ma \emph{własność Radona--Nikodyma}
(RNP), jeżeli dla każdej skończonej przestrzeni z miarą
$(S,\Sigma,\mu)$ oraz każdej miary wektorowej
$m:\Sigma\to X$ o skończonej całkowitej wariacji,
bezwzględnie ciągłej względem $\mu$, istnieje funkcja
$f\in L^1(S,\mu;X)$ taka, że
$m(A)=\int_Af(s)\,\mathrm d\mu(s)$ dla wszystkich $A\in\Sigma$.
\end{definicja}

\begin{przyklad}[Brak RNP w przestrzeni funkcji ciągłych]
\label{ex:rnp-counterexample15b}
Niech $X=C([0,1])$ z normą supremum i niech $\mu$ będzie
miarą Lebesgue'a na $[0,1]$. Definiujemy
$m(A)(t)=\mu(A\cap[0,t])$ dla $t\in[0,1]$.
Funkcja $m(A)$ jest ciągła, a $m$ jest miarą wektorową,
ponieważ dla parami rozłącznych $A_n$ zachodzi
$\|m(\bigcup_{n>N}A_n)\|_\infty
=\mu(\bigcup_{n>N}A_n)\to0$.
Co więcej, $\|m(A)\|_\infty=\mu(A)$, więc
$|m|(A)=\mu(A)$; wariacja jest skończona i
bezwzględnie ciągła względem $\mu$.

Gdyby $m(A)=\int_Af(s)\,\mathrm d\mu(s)$ dla pewnej funkcji
$f\in L^1([0,1];C([0,1]))$, to dla każdego wymiernego
$q\in[0,1]$ i wszystkich zbiorów mierzalnych $A$
mielibyśmy
$\int_A f(s)(q)\,\mathrm d\mu(s)
=\int_A\chi_{[0,q]}(s)\,\mathrm d\mu(s)$.
Stąd $f(s)(q)=\chi_{[0,q]}(s)$ dla prawie wszystkich $s$.
Po usunięciu jednego wspólnego zbioru miary zero równość
zachodzi dla wszystkich wymiernych $q$ jednocześnie.
Dla takiego $s\in(0,1)$ funkcja
$q\mapsto\chi_{[0,q]}(s)$ ma skok w $q=s$, a zatem
nie może być ograniczeniem ciągłej funkcji $q\mapsto f(s)(q)$.
Otrzymaliśmy sprzeczność: przestrzeń $C([0,1])$ nie ma RNP.
\end{przyklad}

\begin{uwaga}[Rozszerzenie na przestrzenie lokalnie wypukłe]
\label{rem:lcs-pettis15b}
Dla przestrzeni lokalnie wypukłej Hausdorffa $E$ sensowna
pozostaje \emph{całka Pettisa w $E$}: dla funkcji $f:S\to E$
wymagamy całkowalności wszystkich skalaryzacji $\ell\circ f$,
$\ell\in E'$, i dla każdego $A\in\Sigma$ istnienia $v_A\in E$
takiego, że $\ell(v_A)=\int_A\ell(f(s))\,\mathrm d\mu(s)$.
Rozdzielanie punktów przez $E'$ zapewnia jednoznaczność
$v_A$. Jeżeli $T:E\to F$ jest ciągłym odwzorowaniem liniowym
między przestrzeniami lokalnie wypukłymi Hausdorffa,
to $Tf$ jest całkowalna w sensie Pettisa w $F$ i
$T(v_A)=\int_A^{\mathrm P} T(f(s))\,\mathrm d\mu(s)$:
dla każdego $\varphi\in F'$ wystarczy zastosować
$\varphi\circ T\in E'$. Nie zakładamy jednak, że sama
zupełność $E$ wymusza istnienie całek Pettisa wszystkich
funkcji słabo mierzalnych: jest to własność mocniejsza.
\end{uwaga}

\begin{uwaga}[Zakres teorii]
W osobnym skrypcie teorii miary rozwijamy skalarne podstawy
całkowania, natomiast tu wyjaśniliśmy różnice pomiędzy
całkowaniem w normie, całkowaniem skalarnym z wartościami
w bidualu i całkowaniem słabym-* do przestrzeni dualnej.
Równania różniczkowe wykorzystujące te konstrukcje należą
do osobnego skryptu PDE. Ujęcie całek dla ogólnych przestrzeni
liniowo-topologicznych nie musi być oparte na seminormach;
bez lokalnej wypukłości argumenty Hahna--Banacha nie są
automatycznie dostępne.
\end{uwaga}

\addtocounter{section}{-1}
\endgroup


% Wersja źródłowa: audyt II, jednostka 16
\setcounter{section}{15}
\section{Funkcjonały wypukłe i podstawy rachunku wariacyjnego}
\label{sec:wariacyjny16}

\subsection{Od pochodnej Frécheta do problemu minimum}

W poprzednich rozdziałach badaliśmy lokalną strukturę funkcji
różniczkowalnych na przestrzeniach Banacha. Rachunek wariacyjny pyta
o istnienie i charakter minimizatorów funkcjonału $J:E\to\R$,
często na nieskończenie wymiarowej przestrzeni funkcji.
W tym rozdziale $H$ oznacza rzeczywistą przestrzeń Hilberta,
a $E$ rzeczywistą przestrzeń Banacha. Pisząc $\langle\cdot,\cdot\rangle$
w $H$, przyjmujemy liniowość w pierwszym argumencie.

\begin{definicja}[Wypukłość funkcjonału]
\label{def:convex-fun16}
Niech $C$ będzie wypukłym podzbiorem przestrzeni liniowej $E$.
Funkcjonał $J:C\to\R$ nazywamy \emph{wypukłym}, jeżeli
$J((1-t)x+ty)\le(1-t)J(x)+tJ(y)$ dla wszystkich $x,y\in C$
i $t\in[0,1]$. Jest \emph{ściśle wypukły}, jeśli nierówność
jest ostra dla $x\ne y$ i $0<t<1$.
\end{definicja}

\begin{twierdzenie}[Nierówność styczna]
\label{thm:convex-derivative16}
Niech $U\subseteq E$ będzie otwarty i wypukły, a
$J:U\to\R$ różniczkowalny Frécheta w każdym punkcie.
Następujące warunki są równoważne:
\begin{enumerate}
\item $J$ jest wypukły na $U$;
\item $J(y)\ge J(x)+DJ(x)(y-x)$ dla wszystkich $x,y\in U$.
\end{enumerate}
Jeżeli $DJ(x_0)=0$, gdzie $x_0\in U$, a $J$ jest wypukły,
to $x_0$ jest globalnym minimizatorem $J$ na $U$.
Odwrotnie, jeśli $x_0$ jest lokalnym minimizatorem
różniczkowalnego $J$ i $x_0\in U$, to $DJ(x_0)=0$.
\end{twierdzenie}
\begin{proof}
Dla ustalonych $x,y\in U$ niech
$\phi(t)=J(x+t(y-x))$, $0\le t\le1$.
Jeżeli $J$ jest wypukły, to dla $0<t\le1$
z definicji wypukłości
$[\phi(t)-\phi(0)]/t\le\phi(1)-\phi(0)$.
Przechodząc z $t$ do zera, otrzymujemy
$DJ(x)(y-x)\le J(y)-J(x)$.

Odwrotnie, ustalmy $z=(1-t)x+ty$.
Nierówność z punktu 2 zastosowana w $z$ do $x$ i do $y$
daje po przemnożeniu przez $1-t$ i $t$ oraz dodaniu
$(1-t)J(x)+tJ(y)\ge J(z)$; składniki liniowe się znoszą.
Jeżeli $DJ(x_0)=0$, nierówność styczna zapewnia
$J(y)\ge J(x_0)$. Jeśli zaś $x_0$ jest lokalnym minimum,
jednowymiarowa funkcja $t\mapsto J(x_0+th)$ ma minimum
przy $t=0$ dla każdego $h\in E$, więc $DJ(x_0)h=0$.
\end{proof}

\begin{definicja}[Silna wypukłość]
\label{def:strong-convex16}
Funkcjonał $J:H\to\R$ jest $\mu$-silnie wypukły
($\mu>0$), jeśli $x\mapsto J(x)-\frac\mu2\|x\|^2$
jest wypukły.
\end{definicja}

\begin{stwierdzenie}[Koercywność z silnej wypukłości]
\label{prop:strong-coercive16}
Niech $J:H\to\R$ będzie $\mu$-silnie wypukły i
różniczkowalny Frécheta na $H$. Wtedy dla $x,y\in H$
zachodzi
\[
J(y)\ge J(x)+DJ(x)(y-x)+\frac\mu2\|y-x\|^2.
\]
W szczególności $J$ jest koercywny, tj.
$J(y)\to+\infty$, gdy $\|y\|\to\infty$,
i ma co najwyżej jeden minimizator.
\end{stwierdzenie}
\begin{proof}
Stosujemy Twierdzenie~\ref{thm:convex-derivative16}
do funkcjonału $G(z)=J(z)-\mu\|z\|^2/2$.
Ponieważ $DG(x)h=DJ(x)h-\mu\langle x,h\rangle$,
po przekształceniu dostajemy żądaną nierówność.
Dla $x=0$ mamy
$J(y)\ge J(0)-\|DJ(0)\|\|y\|+\mu\|y\|^2/2$,
co zapewnia koercywność. Jeżeli $x\ne y$ są minimizatorami,
to silna wypukłość przy $t=1/2$ daje
$J((x+y)/2)\le [J(x)+J(y)]/2-\mu\|x-y\|^2/8$,
sprzecznie z minimalnością.
\end{proof}

\begin{przyklad}[Niezbędne jednostajne oszacowanie hesjanu]
\label{ex:variational-hessian16}
Jeżeli $J\in C^2(H,\R)$ i
$D^2J(x)[h,h]\ge\mu\|h\|^2$ dla wszystkich
$x,h\in H$ i pewnego $\mu>0$, to $J$ jest
$\mu$-silnie wypukły. Rzeczywiście, druga pochodna
funkcji $t\mapsto J(x+t(y-x))-
\frac\mu2\|x+t(y-x)\|^2$ jest nieujemna,
więc funkcja jest wypukła na każdym odcinku.
Sama dodatnia określoność $D^2J(0)$ nie wystarcza:
w analizie pochodnych wyższych rzędów można skonstruować
na $\ell^2$ funkcjonał z dodatnio określonym hesjanem
w zerze, dla którego zero nie jest minimum lokalnym.
\end{przyklad}

\subsection{Zbieżność słaba i metoda bezpośrednia}

\begin{definicja}[Półciągłość z dołu i koercywność]
\label{def:weak-lsc16}
Funkcjonał $J:H\to(-\infty,+\infty]$ jest
\emph{słabo półciągły z dołu}, jeśli dla każdej sieci
$x_\alpha\to x$ w topologii słabej $\sigma(H,H')$
zachodzi $J(x)\le\liminf_\alpha J(x_\alpha)$,
gdzie $\liminf$ sieci rozumiemy w zwykłym porządku liczb
rozszerzonych. Równoważnie, każdy zbiór podpoziomowy
$\{x:J(x)\le a\}$, $a\in\R$, jest słabo domknięty.
Funkcjonał jest \emph{koercywny}, jeśli $J(x)\to+\infty$
przy $\|x\|\to\infty$.
\end{definicja}

\begin{lemat}[Podpoziomowe zbiory wypukłe]
\label{lem:convex-weaklsc16}
Niech $J:H\to(-\infty,+\infty]$ będzie wypukły
i półciągły z dołu w topologii normowej.
Wtedy $J$ jest słabo półciągły z dołu.
\end{lemat}
\begin{proof}
Dla każdego $a\in\R$ zbiór $C_a=\{x:J(x)\le a\}$
jest normowo domknięty przez półciągłość z dołu
i wypukły przez wypukłość $J$. Z twierdzenia
o oddzielaniu zbiorów wypukłych (twierdzenie Hahna--Banacha o oddzielaniu zbiorów wypukłych)
każdy normowo domknięty zbiór wypukły jest domknięty
także w topologii słabej. Zatem wszystkie $C_a$
są słabo domknięte, co daje tezę.
\end{proof}

\begin{lemat}[Słaba zbieżność podciągu w przestrzeni Hilberta]
\label{lem:hilbert-weakseq16}
Niech $H$ będzie rzeczywistą przestrzenią Hilberta,
a $(x_n)$ ciągiem ograniczonym w $H$.
Istnieje podciąg $(x_{n_k})$ i $x\in H$ takie,
że $x_{n_k}\to x$ słabo w $H$.
Nie zakładamy ośrodkowości całej przestrzeni $H$.
\end{lemat}
\begin{proof}
Niech $M=\sup_n\|x_n\|<\infty$ i
$H_0=\overline{\operatorname{lin}\{x_n:n\in\N\}}$.
Jest to ośrodkowa przestrzeń Hilberta, więc ma
przeliczalną bazę ortonormalną $(e_j)$
(w przypadku skończonego wymiaru kończymy po skończenie wielu elementach).
Dla każdego $j$ ciąg $\langle x_n,e_j\rangle$ jest ograniczony.
Metodą diagonalną wybieramy podciąg $(x_{n_k})$,
na którym wszystkie te współrzędne mają granice.

Na przestrzeni skończonych kombinacji wektorów $e_j$
definiujemy $L(v)=\lim_k\langle x_{n_k},v\rangle$.
Istnienie granicy wynika z jej istnienia na każdym $e_j$,
a $|L(v)|\le M\|v\|$. Funkcjonał $L$ przedłuża się
jednoznacznie na $H_0$ jako ciągły funkcjonał liniowy.
Z twierdzenia Riesza istnieje $x\in H_0$ takie,
że $L(v)=\langle x,v\rangle$ dla $v\in H_0$.
Dla dowolnego $v\in H_0$ aproksymujemy $v$
skończonymi kombinacjami $e_j$ i używamy
$\|x_{n_k}\|\le M$, aby otrzymać
$\langle x_{n_k},v\rangle\to\langle x,v\rangle$.
Dla $v\in H_0^\perp$ obie strony są zerowe.
Rozkład $H=H_0\oplus H_0^\perp$ kończy dowód.
\end{proof}

\begin{twierdzenie}[Metoda bezpośrednia w przestrzeni Hilberta]
\label{thm:direct-method16}
Niech $H$ będzie rzeczywistą przestrzenią Hilberta,
$C\subseteq H$ niepustym, normowo domkniętym i wypukłym
zbiorem, a $J:C\to(-\infty,+\infty]$ funkcjonałem,
który jest właściwy (tzn. przyjmuje przynajmniej jedną
wartość skończoną), ograniczony z dołu, koercywny na $C$
i półciągły z dołu dla topologii słabej indukowanej na $C$.
Wtedy istnieje $u\in C$ takie, że
$J(u)=\inf_{v\in C}J(v)$.
Jeżeli dodatkowo $J$ jest ściśle wypukły i $C$ jest wypukły,
minimizator jest jedyny.
\end{twierdzenie}
\begin{proof}
Niech $m=\inf_C J$, który jest skończony na mocy
właściwości i ograniczoności z dołu. Wybierzmy
ciąg minimalizujący $(u_n)\subseteq C$ z $J(u_n)\to m$.
Koercywność oznacza, że $(u_n)$ jest ograniczony.
Na mocy Lematu~\ref{lem:hilbert-weakseq16}
pewien podciąg $u_{n_k}\to u$ słabo w $H$.
Ponieważ $C$ jest normowo domknięty i wypukły,
jest również słabo domknięty, więc $u\in C$.
Słaba półciągłość z dołu daje
$J(u)\le\liminf_k J(u_{n_k})=m$; z definicji infimum
zachodzi też $J(u)\ge m$. Jeśli $u\ne v$ są minimizatorami,
ścisła wypukłość daje $J((u+v)/2)<m$, sprzeczność.
\end{proof}

\begin{uwaga}[Czego nie daje sama koercywność?]
\label{rem:coercive16}
W nieskończonym wymiarze normowo domknięta kula
nie jest zwarta. Sama ciągłość i koercywność funkcjonału
nie wystarczają więc do uzyskania minimum metodą
Weierstrassa. W twierdzeniu wykorzystaliśmy
słabą zbieżność podciągów i słabą półciągłość z dołu.
W ogólnej przestrzeni Banacha naturalną rolę przestrzeni
Hilberta przejmują przestrzenie refleksywne, lecz dowód
w tej klasie wymaga dodatkowej teorii zwartości słabej.
\end{uwaga}

\subsection{Subgradienty i rzutowanie proksymalne}

\begin{definicja}[Subgradient i subróżniczka]
\label{def:subgrad16}
Niech $J:H\to(-\infty,+\infty]$ będzie właściwym
funkcjonałem wypukłym, a $u\in H$ punktem takim,
że $J(u)<+\infty$. Wektor $p\in H$ nazywamy
\emph{subgradientem} $J$ w $u$, jeśli
$J(v)\ge J(u)+\langle p,v-u\rangle$ dla każdego $v\in H$.
Zbiór wszystkich subgradientów oznaczamy $\partial J(u)$
i nazywamy \emph{subróżniczką} $J$ w $u$.
Dla $J(u)=+\infty$ przyjmujemy $\partial J(u)=\varnothing$.
\end{definicja}

\begin{stwierdzenie}[Warunek optymalności i monotoniczność]
\label{prop:subgrad16}
Jeżeli $J:H\to(-\infty,+\infty]$ jest właściwy i wypukły,
to $u$ minimalizuje $J$ wtedy i tylko wtedy, gdy
$0\in\partial J(u)$. Jeśli $p\in\partial J(u)$ oraz
$q\in\partial J(v)$, to
$\langle p-q,u-v\rangle\ge0$.
Jeżeli $J$ jest różniczkowalny Frécheta w $u$ i skończony
na otoczeniu $u$, to $\partial J(u)=\{\nabla J(u)\}$.
\end{stwierdzenie}
\begin{proof}
Pierwsza równoważność wynika wprost z definicji.
Z nierówności subgradientowych
$J(v)\ge J(u)+\langle p,v-u\rangle$ i
$J(u)\ge J(v)+\langle q,u-v\rangle$,
po dodaniu dostajemy monotoniczność.
Jeśli $J$ jest różniczkowalny, to dla $p\in\partial J(u)$
podstawiamy $v=u+th$ z dodatnim i ujemnym $t$,
po czym przechodzimy do zera. Otrzymujemy
$\langle p,h\rangle=DJ(u)h$ dla każdego $h\in H$.
Z kolei nierówność styczna dla wypukłego $J$
zapewnia, że $\nabla J(u)$ jest subgradientem.
\end{proof}

\begin{przyklad}[Subróżniczka normy w Hilbercie]
\label{ex:norm-subgrad16}
Dla $J(x)=\|x\|$ zachodzi
$\partial J(0)=\{p:\|p\|\le1\}$ i
$\partial J(x)=\{x/\|x\|\}$ dla $x\ne0$.
Pierwszy wzór wynika z nierówności Cauchy'ego--Schwarza
i podstawienia $v=tp/\|p\|$ dla $t>0$.
Drugi otrzymujemy, różniczkując normę w punkcie $x\ne0$.
W ogólnej przestrzeni Banacha analogiczny subgradient normy
jest funkcjonałem z przestrzeni dualnej, a jego istnienie
dla $x\ne0$ zapewnia twierdzenie Hahna--Banacha.
\end{przyklad}

\begin{twierdzenie}[Odwzorowanie proksymalne]
\label{thm:prox16}
Niech $J:H\to(-\infty,+\infty]$ będzie właściwy,
wypukły, normowo półciągły z dołu i ograniczony z dołu,
a $\lambda>0$. Dla każdego $y\in H$ istnieje dokładnie
jeden punkt $u=\operatorname{prox}_{\lambda J}(y)$
minimalizujący funkcjonał
$v\mapsto J(v)+\|v-y\|^2/(2\lambda)$.
Jest on scharakteryzowany równoważnością
$y-u\in\lambda\partial J(u)$.
Ponadto odwzorowanie
$\operatorname{prox}_{\lambda J}:H\to H$ jest
nieekspansywne:
$\|\operatorname{prox}_{\lambda J}(y)-
\operatorname{prox}_{\lambda J}(z)\|\le\|y-z\|$.
\end{twierdzenie}
\begin{proof}
Funkcjonał $\Phi_y(v)=J(v)+\|v-y\|^2/(2\lambda)$
jest właściwy, wypukły i normowo półciągły z dołu,
a więc słabo półciągły z dołu na mocy
Lematu~\ref{lem:convex-weaklsc16}.
Jest koercywny, gdyż $J$ jest ograniczony z dołu,
i jest ściśle wypukły dzięki składnikowi kwadratowemu.
Twierdzenie~\ref{thm:direct-method16} dla $C=H$
daje jednoznaczny minimizator $u$.

Jeśli $u$ minimalizuje $\Phi_y$, to dla $v$ z
$J(v)<+\infty$ i $0<t<1$, z wypukłości $J$
oraz minimalności mamy
\[
0\le\Phi_y(u+t(v-u))-\Phi_y(u)
\le t[J(v)-J(u)]+\frac t\lambda\langle u-y,v-u\rangle
+\frac{t^2}{2\lambda}\|v-u\|^2.
\]
Po podzieleniu przez $t$ i przejściu z $t$ do zera
otrzymujemy $J(v)\ge J(u)+\langle(y-u)/\lambda,v-u\rangle$,
a więc $(y-u)/\lambda\in\partial J(u)$.
Odwrotnie, ta nierówność oraz rozwinięcie
$\|v-y\|^2$ wokół $u$ dają
$\Phi_y(v)-\Phi_y(u)\ge\|v-u\|^2/(2\lambda)$,
co zapewnia minimalność.

Niech $u=\operatorname{prox}_{\lambda J}(y)$ oraz
$v=\operatorname{prox}_{\lambda J}(z)$.
Monotoniczność subróżniczki daje
$\langle (y-u)-(z-v),u-v\rangle\ge0$.
Stąd $\|u-v\|^2\le\langle y-z,u-v\rangle
\le\|y-z\|\,\|u-v\|$,
co po podzieleniu przez $\|u-v\|$ (gdy $u\ne v$)
daje tezę.
\end{proof}

\subsection{Przykład: słabe zagadnienie brzegowe}

Niech $\lambda$ oznacza miarę Lebesgue'a na $[0,1]$.
Wprowadzamy przestrzeń $H_0^1(0,1)$ przez opis elementarny:
\[
H_0^1(0,1)=\Bigl\{u:[0,1]\to\R:\ 
 u(x)=\int_0^x v(t)\,\mathrm dt\ (0\le x\le1),\ 
 v\in L^2(0,1),\ \int_0^1 v(t)\,\mathrm dt=0\Bigr\}.
\]
Funkcje $u$ są absolutnie ciągłe, ich pochodna $u'=v$
istnieje prawie wszędzie, a warunki brzegowe
$u(0)=u(1)=0$ są spełnione. Przyjmujemy
$\langle u,w\rangle_H=\int_0^1u'(t)w'(t)\,\mathrm dt$.
Odwzorowanie $u\mapsto u'$ jest liniową izometrią
$H_0^1(0,1)$ na domkniętą podprzestrzeń
$\{v\in L^2(0,1):\int_0^1v(t)\,\mathrm dt=0\}$,
więc $H_0^1(0,1)$ jest przestrzenią Hilberta.

\begin{lemat}[Nierówność Poincarégo na odcinku]
\label{lem:poincare16}
Dla każdego $u\in H_0^1(0,1)$ zachodzi
$\|u\|_{L^2}\le2^{-1/2}\|u'\|_{L^2}$.
\end{lemat}
\begin{proof}
Ponieważ $u(x)=\int_0^xu'(t)\,\mathrm dt$,
nierówność Cauchy'ego--Schwarza daje
$|u(x)|^2\le x\int_0^x|u'(t)|^2\,\mathrm dt
\le x\|u'\|_{L^2}^2$.
Całkując po $x\in[0,1]$, otrzymujemy
$\|u\|_{L^2}^2\le\frac12\|u'\|_{L^2}^2$.
\end{proof}

\begin{twierdzenie}[Zasada minimalizacji dla problemu Dirichleta]
\label{thm:dirichlet16}
Niech $f\in L^2(0,1;\R)$. Na $H=H_0^1(0,1)$
określmy funkcjonał
\[
J(u)=\frac12\int_0^1|u'(t)|^2\,\mathrm dt
-\int_0^1 f(t)u(t)\,\mathrm dt.
\]
Istnieje dokładnie jeden minimizator $u\in H$ funkcjonału $J$.
Jest on dokładnie jedynym elementem $H$ spełniającym
równanie słabe
\[
\int_0^1u'(t)v'(t)\,\mathrm dt
=\int_0^1f(t)v(t)\,\mathrm dt\qquad(v\in H).
\]
Dla $f(t)=1$ rozwiązaniem jest $u(t)=t(1-t)/2$.
\end{twierdzenie}
\begin{proof}
Funkcjonał $J$ jest ciągły i ściśle wypukły na $H$.
Z Lematu~\ref{lem:poincare16} i nierówności
Cauchy'ego--Schwarza wynika
$J(u)\ge\frac12\|u\|_H^2-
2^{-1/2}\|f\|_{L^2}\|u\|_H$, więc jest koercywny
i ograniczony z dołu. Ponieważ jest wypukły
i normowo ciągły, jest słabo półciągły z dołu.
Twierdzenie~\ref{thm:direct-method16} daje minimizator,
a ścisła wypukłość zapewnia jego jednoznaczność.

Z ciągłości formy kwadratowej i funkcjonału liniowego
wynika, że $J$ jest różniczkowalny Frécheta,
a jego pochodna ma postać
$DJ(u)v=\int_0^1u'(t)v'(t)\,\mathrm dt-
\int_0^1 f(t)v(t)\,\mathrm dt$.
Twierdzenie~\ref{thm:convex-derivative16} daje
równoważność $DJ(u)=0$ i globalnej minimalności.
Jeśli $f\equiv1$, funkcja $u(t)=t(1-t)/2$ spełnia
$u(0)=u(1)=0$ oraz $u''(t)=-1$; całkowanie przez części
dla absolutnie ciągłego $v\in H$ daje
$\int_0^1u'v'\,\mathrm dt=\int_0^1v\,\mathrm dt$.
Przez jednoznaczność jest to szukane rozwiązanie.
\end{proof}

\subsection{Wypukłe maksima i twierdzenie Kreina--Milmana}

Twierdzenie Kreina--Milmana i zasada maksimum Bauera orzekają,
że niepusty zwarty wypukły podzbiór przestrzeni lokalnie wypukłej
Hausdorffa jest domkniętą otoczką wypukłą swoich punktów
ekstremalnych, a każda wypukła i półciągła z góry funkcja
rzeczywista na tym zbiorze osiąga maksimum w punkcie ekstremalnym. Ich rola w optymalizacji jest odmienna od metody bezpośredniej:
\emph{zwartość wypukłego zbioru} daje możliwość badania maksimum
funkcjonału wypukłego na punktach ekstremalnych,
natomiast \emph{koercywność i słaba półciągłość z dołu}
służą do znajdowania minimum.

\begin{stwierdzenie}[Maksimum funkcjonału afinicznego]
\label{prop:affine-extreme16}
Niech $E$ będzie lokalnie wypukłą TVS Hausdorffa nad $\K$,
$K\subseteq E$ niepustym, zwartym i wypukłym zbiorem,
a $A:K\to\R$ ciągłą funkcją afiniczną,
czyli $A((1-t)x+ty)=(1-t)A(x)+tA(y)$.
Wtedy istnieje punkt ekstremalny $e\in K$, w którym
$A$ osiąga maksimum, oraz (być może inny) punkt ekstremalny,
w którym osiąga minimum.
\end{stwierdzenie}
\begin{proof}
Funkcje $A$ i $-A$ są wypukłe i ciągłe,
więc zastosowanie zasady maksimum Bauera do obu daje tezę.
\end{proof}

\begin{przyklad}[Miary prawdopodobieństwa i ich punkty ekstremalne]
\label{ex:prob-measures16}
Niech $K$ będzie niepustą zwartą przestrzenią metryczną,
a $\mathcal P(K)$ zbiorem borelowskich miar probabilistycznych na $K$.
Na mocy twierdzenia Riesza--Markowa identyfikujemy
$\mathcal P(K)$ z dodatnimi funkcjonałami na $C(K;\R)$,
przyjmującymi wartość $1$ na funkcji stałej $1$.
W topologii słabej-* $\sigma(C(K)',C(K))$ jest to zbiór
zwarty i wypukły: jako podzbiór kuli jednostkowej
jest domknięty przez warunki
$L(1)=1$ i $L(\varphi)\ge0$ dla każdej $\varphi\ge0$.

Punktami ekstremalnymi $\mathcal P(K)$ są dokładnie
miary Diraca $\delta_x$, $x\in K$.
Dla $\delta_x=t\mu+(1-t)\nu$, $0<t<1$, dodatniość
miar i równość ich mas wymuszają
$\mu(K\setminus\{x\})=\nu(K\setminus\{x\})=0$,
więc $\mu=\nu=\delta_x$.
Jeśli natomiast $\mu$ nie jest miarą Diraca,
istnieje zbiór borelowski $B$ o
$0<\mu(B)<1$: w przeciwnym razie nośnik miary
musiałby być jednopunktowy (dla zwartej metrycznej $K$
wynika to z rozdzielania różnych punktów nośnika
rozłącznymi otoczeniami o dodatniej mierze).
Wówczas
$\mu=\mu(B)\,\mu_B+(1-\mu(B))\,\mu_{K\setminus B}$,
gdzie $\mu_B$ i $\mu_{K\setminus B}$ są znormalizowanymi
ograniczeniami $\mu$, a obie miary są różne.
Twierdzenie Kreina--Milmana stwierdza więc, że
$\mathcal P(K)$ jest słabo-* domkniętą otoczką wypukłą
rodziny $\{\delta_x:x\in K\}$.
W szczególności dla $\varphi\in C(K;\R)$
maksimum funkcjonału $\mu\mapsto\int_K\varphi(x)\,\mathrm d\mu(x)$
można osiągnąć na pewnej mierze Diraca.
\end{przyklad}

\begin{uwaga}[Dalsza perspektywa]
Twierdzenie Kreina--Milmana jest rezultatem o \emph{zwartych
zbiorach wypukłych} w ogólnych przestrzeniach lokalnie wypukłych,
a nie jedynie o kulach jednostkowych przestrzeni Banacha.
Z kolei subróżniczki i operatory proksymalne prowadzą
do teorii operatorów monotonicznych i metod minimalizacji,
co wykracza poza podstawowy zakres obecnego rozdziału.
\end{uwaga}


% Wersja źródłowa: audyt II, jednostka 17
\setcounter{section}{16}
\section{Przestrzenie funkcji gładkich i przestrzeń Schwartza}
\label{sec:gladkie-schwartz17}

\subsection{Topologia przestrzeni \texorpdfstring{$C^\infty(\Omega)$}{Cinfty(Omega)}}

Niech $d\ge1$, niech $\Omega\subseteq\R^d$ będzie niepustym zbiorem
otwartym, a $\K\in\{\R,\C\}$. Dla wielowskaźnika
$\alpha=(\alpha_1,\dots,\alpha_d)\in\N_0^d$ oznaczamy
$|\alpha|=\sum_j\alpha_j$, $\alpha!=\prod_j\alpha_j!$,
$x^\alpha=\prod_jx_j^{\alpha_j}$ oraz
$\partial^\alpha=\partial_1^{\alpha_1}\cdots\partial_d^{\alpha_d}$.
Przez $C^\infty(\Omega;\K)$ rozumiemy przestrzeń funkcji
$\Omega\to\K$ posiadających ciągłe pochodne cząstkowe wszystkich rzędów.
Nie narzucamy ograniczoności funkcji ani ich pochodnych na całej $\Omega$.

\begin{definicja}[Seminormy zbieżności lokalnej wszystkich pochodnych]
\label{def:seminormy-cinf17}
Dla zwartego $K\subseteq\Omega$ i liczby $m\in\N_0$ określamy
\[
 p_{K,m}(f)=\max_{|\alpha|\le m}\sup_{x\in K}|\partial^\alpha f(x)|,
 \qquad f\in C^\infty(\Omega;\K).
\]
Topologią standardową na $C^\infty(\Omega;\K)$ nazywamy
najmniejszą topologię liniowo-topologiczną, w której wszystkie
$p_{K,m}$ są ciągłe. Zbieżność w tej topologii oznacza
jednostajną zbieżność każdej pochodnej na każdym zbiorze zwartym
zawartym w $\Omega$.
\end{definicja}

\begin{lemat}[Wyczerpanie zbiorami zwartymi]
\label{lem:exhaustion17}
Dla każdego niepustego otwartego $\Omega\subseteq\R^d$ istnieje
ciąg zbiorów zwartych $(K_n)_{n\ge1}$ taki, że
$K_n\subseteq\operatorname{int}K_{n+1}$, $\bigcup_nK_n=\Omega$
i każdy zwarty $K\subseteq\Omega$ jest zawarty w pewnym $K_n$.
\end{lemat}
\begin{proof}
Dla $\Omega\ne\R^d$ weźmy
$K_n=\{x\in\Omega:|x|\le n,\
\operatorname{dist}(x,\R^d\setminus\Omega)\ge1/n\}$;
dla $\Omega=\R^d$ przyjmijmy $K_n=\{x:|x|\le n\}$.
Zbiory te są zwarte: warunek odległości jest domknięty,
a dodatnia odległość od dopełnienia wymusza przynależność do $\Omega$.
Nierówności definiujące $K_n$ są ściślejsze niż analogiczne dla
$K_{n+1}$, stąd $K_n\subseteq\operatorname{int}K_{n+1}$.
Każdy $x\in\Omega$ ma dodatnią odległość od dopełnienia
i skończoną normę, więc należy do pewnego $K_n$.
Jeżeli zwarty $K\subseteq\Omega$ nie byłby zawarty w żadnym $K_n$,
można byłoby wybrać $x_n\in K\setminus K_n$ i podciąg zbieżny
do $x\in K$. Pewne $K_N$ zawiera $x$ w swoim wnętrzu;
dla dużych $n\ge N$ mielibyśmy $x_n\in K_N\subseteq K_n$,
sprzeczność. Gdy dopełnienie jest puste, argument jest identyczny.
\end{proof}

\begin{twierdzenie}[Przeliczalna rodzina seminorm]
\label{thm:countable-cinf17}
Niech $\Omega\subseteq\R^d$ będzie niepusty i otwarty,
a $(K_n)$ ciągiem z Lematu~\ref{lem:exhaustion17}.
Wówczas topologia z Definicji~\ref{def:seminormy-cinf17}
jest generowana przez rosnący ciąg seminorm
$q_n(f)=\max_{|\alpha|\le n}\sup_{x\in K_n}|\partial^\alpha f(x)|$,
$n\ge1$. Jest ona metryzowalna i Hausdorffa, a jej zbieżność
ciągów jest równoważna warunkowi
$q_n(f_j-f)\to0$ dla każdego $n$.
\end{twierdzenie}
\begin{proof}
Każde $q_n$ jest jedną z seminorm $p_{K,m}$.
Z drugiej strony dla ustalonych $K\Subset\Omega$ i $m$
wybierzmy $n\ge m$ z $K\subseteq K_n$.
Wtedy $p_{K,m}\le q_n$, więc obie rodziny wyznaczają tę samą
topologię. Ponieważ $K_n\subseteq K_{n+1}$, seminormy $q_n$
są niemalejące. Jeśli $f\ne0$, to $f(x)\ne0$ dla pewnego
$x\in\Omega$ i $x\in K_n$ dla pewnego $n$, zatem $q_n(f)>0$.
Topologia jest więc Hausdorffa. Jest metryzowalna na mocy
kryterium dla przeliczalnej rodziny seminorm; można użyć metryki
$\sum_{n\ge1}2^{-n}\min\{1,q_n(f-g)\}$.
Z definicji topologii wyznaczonej przez seminormy wynika
podane kryterium zbieżności.
\end{proof}

\begin{twierdzenie}[Zupełność $C^\infty(\Omega)$]
\label{thm:frechet-cinf17}
Dla każdego niepustego otwartego $\Omega\subseteq\R^d$ przestrzeń
$C^\infty(\Omega;\K)$ ze standardową topologią jest przestrzenią
Frécheta.
\end{twierdzenie}
\begin{proof}
Metryzowalność, lokalna wypukłość i warunek Hausdorffa wynikają
z Twierdzenia~\ref{thm:countable-cinf17}; pozostaje zupełność.
Niech $(f_j)$ będzie ciągiem Cauchy'ego dla wszystkich $q_n$.
Dla każdego wielowskaźnika $\alpha$ pochodne
$\partial^\alpha f_j$ są jednostajnie Cauchy'ego na każdym
$K_n$, gdy $n\ge|\alpha|$. Z zupełności przestrzeni
$C(K_n;\K)$ mają granice jednostajne $g_{\alpha,n}$.
Na przecięciach $K_n$ granice się zgadzają, więc określają
funkcje ciągłe $g_\alpha:\Omega\to\K$; w szczególności
$f_j\to g_0$ lokalnie jednostajnie.

Wykażemy, że $g_\alpha=\partial^\alpha g_0$.
Ustalmy $x\in\Omega$ i zamkniętą kulę $\overline B(x,r)\Subset\Omega$.
Dla $y$ na dostatecznie krótkim odcinku równoległym do osi $e_i$
wewnątrz tej kuli zachodzi
\[
 \partial^\alpha f_j(x+te_i)-\partial^\alpha f_j(x)
 =\int_0^t\partial^{\alpha+e_i}f_j(x+se_i)\,\mathrm ds.
\]
Wybierając $K_n$ zawierający tę kulę i przechodząc do granicy
jednostajnej, dostajemy tę samą równość dla $g_\alpha$
i $g_{\alpha+e_i}$. Stąd istnieje pochodna cząstkowa
$\partial_i g_\alpha=g_{\alpha+e_i}$.
Indukcja względem $|\alpha|$ daje $g_0\in C^\infty(\Omega)$
i wszystkie wymagane tożsamości. Jednostajna zbieżność
pochodnych na $K_n$ oznacza $q_n(f_j-g_0)\to0$ dla każdego $n$.
Zupełność sekwencyjna w przestrzeni metryzowalnej oznacza zupełność.
\end{proof}

\begin{stwierdzenie}[Operacje i ograniczoność]
\label{prop:operations-cinf17}
Niech $\Omega\subseteq\R^d$ będzie niepusty i otwarty.
W przestrzeni $C^\infty(\Omega;\K)$ różniczkowanie
$\partial^\alpha$ jest operatorem liniowym ciągłym,
a mnożenie $(f,g)\mapsto fg$ jest odwzorowaniem dwuliniowym
łącznie ciągłym. Zbiór $B\subseteq C^\infty(\Omega;\K)$
jest ograniczony w sensie TVS wtedy i tylko wtedy, gdy
$\sup_{f\in B}p_{K,m}(f)<\infty$ dla wszystkich zwartych
$K\Subset\Omega$ i wszystkich $m\in\N_0$.
\end{stwierdzenie}
\begin{proof}
Z nierówności $p_{K,m}(\partial^\alpha f)\le p_{K,m+|\alpha|}(f)$
wynika ciągłość różniczkowania. Wzór Leibniza daje
$p_{K,m}(fg)\le C_{d,m}p_{K,m}(f)p_{K,m}(g)$,
co dowodzi łącznej ciągłości mnożenia w zerze;
dwuliniowość i ciągłość przy dowolnej parze wynikają przez
rozwinięcie $(f+h)(g+k)-fg$. Kryterium ograniczoności
wynika z ogólnego kryterium seminormowego dla przestrzeni
lokalnie wypukłych: rodzina obrazów zbioru ograniczonego przez
każdą ciągłą seminormę jest ograniczona w $\R$ i odwrotnie.
\end{proof}

\subsection{Funkcje odcięcia i przestrzeń dualna topologiczna}

\begin{lemat}[Gładka funkcja odcięcia]
\label{lem:cutoffs17}
Jeżeli $K\subseteq\Omega\subseteq\R^d$, gdzie $K$ jest zwarty,
a $\Omega$ otwarty, to istnieje $\chi\in C_c^\infty(\Omega;\R)$
taka, że $0\le\chi\le1$ i $\chi=1$ w pewnym otoczeniu $K$.
Tu $C_c^\infty$ oznacza funkcje gładkie o zwartym nośniku.
\end{lemat}
\begin{proof}
Wybierzmy otwarte zbiory $V,W$ takie, że
$K\subseteq V$, $\overline V\subseteq W$, $\overline W\Subset\Omega$.
Można je uzyskać jako skończone sumy kul o dostatecznie małych
promieniach. Istnieje $\delta>0$ taka, że zbiór
$\{x:\operatorname{dist}(x,\overline V)\le\delta\}$
leży w $W$ oraz $\operatorname{dist}(K,\R^d\setminus V)>\delta$.
Niech $\rho\in C_c^\infty(\R^d)$ będzie nieujemna,
niech $\operatorname{supp}\rho\subseteq B(0,\delta)$ i
$\int_{\R^d}\rho(x)\,\mathrm d\lambda_d(x)=1$.
Taką funkcję otrzymujemy przez unormowanie
$x\mapsto\exp[-1/(1-|x|^2/\delta^2)]$ wewnątrz kuli,
przedłużonej zerem poza nią.
Definiujemy $\chi(x)=\int_{\R^d}\mathbf1_V(y)\rho(x-y)\,\mathrm d\lambda_d(y)$.
Różniczkowanie pod znakiem całki jest dozwolone, ponieważ
$\mathbf1_V$ jest ograniczona i $\rho$ ma zwarty nośnik;
stąd $\chi\in C^\infty$. Jest nieujemna i nie przekracza $1$.
Dla $x$ w dostatecznie małym otoczeniu $K$ wszystkie
$x-z$, $z\in\operatorname{supp}\rho$, należą do $V$,
więc $\chi(x)=1$. Nośnik $\chi$ leży w zwartych
podzbiorach $W\Subset\Omega$, co kończy dowód.
\end{proof}

\begin{wniosek}[Gęstość funkcji o zwartym nośniku]
\label{cor:density-cinf17}
Dla każdego niepustego otwartego $\Omega\subseteq\R^d$
przestrzeń $C_c^\infty(\Omega;\K)$ jest gęsta w
$C^\infty(\Omega;\K)$ ze standardową topologią.
\end{wniosek}
\begin{proof}
Dla wyczerpania $(K_n)$ wybierzmy z Lematu~\ref{lem:cutoffs17}
$\chi_n\in C_c^\infty(\Omega)$ równą $1$ na otoczeniu $K_n$.
Dla każdego $f\in C^\infty(\Omega)$ funkcje $\chi_n f$
mają zwarty nośnik. Dla ustalonego $K_m$ i dowolnego
rzędu pochodnych mamy $\chi_n f=f$ na otoczeniu $K_m$
dla wszystkich $n\ge m$. Zatem $\chi_n f\to f$
w topologii $C^\infty(\Omega)$.
\end{proof}

\begin{twierdzenie}[Kryterium ciągłości funkcjonału]
\label{thm:dual-cinf17}
Niech $\Omega\subseteq\R^d$ będzie niepusty i otwarty,
a $T:C^\infty(\Omega;\K)\to\K$ liniowy.
Następujące warunki są równoważne:
\begin{enumerate}
\item $T$ jest ciągły;
\item istnieją zwarty $K\Subset\Omega$, $m\in\N_0$ i $C>0$ takie, że
$|T(f)|\le C\max_{|\alpha|\le m}\sup_{x\in K}|\partial^\alpha f(x)|$
dla każdego $f\in C^\infty(\Omega;\K)$.
\end{enumerate}
W szczególności ciągły funkcjonał $T$ zależy tylko od
zachowania funkcji w pewnym zwartym podzbiorze $\Omega$:
jeśli $f$ znika w otoczeniu tego zbioru, to $T(f)=0$.
\end{twierdzenie}
\begin{proof}
Wybierzmy rosnącą rodzinę seminorm $q_n$ z
Twierdzenia~\ref{thm:countable-cinf17}.
Ciągłość $T$ w zerze oznacza, że pewne otoczenie
$\{f:q_n(f)<\varepsilon\}$ jest odwzorowywane w
$\{z:|z|<1\}$. Z liniowości, przez skalowanie $f$,
otrzymujemy $|T(f)|\le\varepsilon^{-1}q_n(f)$;
jeśli $q_n(f)=0$, skalowanie przez dowolny skalar daje $T(f)=0$.
To dowodzi warunku 2 z $K=K_n$ i $m=n$.
Odwrotna implikacja wynika z ciągłości seminormy
$p_{K,m}$. Ostatnie stwierdzenie jest bezpośrednim
wnioskiem z oszacowania.
\end{proof}

\begin{przyklad}[Ewaluacja i pochodne ewaluacji]
\label{ex:evaluation-cinf17}
Dla $a\in\Omega$ i $\alpha\in\N_0^d$ funkcjonały
$T_{a,\alpha}(f)=(-1)^{|\alpha|}\partial^\alpha f(a)$
są ciągłe na $C^\infty(\Omega;\K)$. Dla $\alpha=0$
przyjmują postać ewaluacji $f\mapsto f(a)$.
W teorii dystrybucji pojawią się jako pochodne
rozkładu Diraca. Obecnie są po prostu przykładami
elementów przestrzeni dualnej topologicznej.
\end{przyklad}

\subsection{Przestrzeń Schwartza \texorpdfstring{$\mathcal S(\R^d)$}{S(Rd)}}

\begin{definicja}[Przestrzeń Schwartza]
\label{def:schwartz17}
Dla $d\ge1$ definiujemy
\[
\mathcal S(\R^d;\K)=\left\{f\in C^\infty(\R^d;\K):
\sup_{x\in\R^d}|x^\alpha\partial^\beta f(x)|<\infty
\ \text{dla wszystkich }\alpha,\beta\in\N_0^d\right\}.
\]
Jej standardową topologię określa rodzina seminorm
$s_{\alpha,\beta}(f)=\sup_x|x^\alpha\partial^\beta f(x)|$.
Równoważną rodzinę tworzą
\[
 Q_{N,m}(f)=\max_{|\beta|\le m}\sup_{x\in\R^d}
 (1+|x|)^N|\partial^\beta f(x)|,
 \qquad N,m\in\N_0.
\]
Przyjmujemy też rosnący ciąg $Q_n=Q_{n,n}$, $n\ge1$.
\end{definicja}

\begin{lemat}[Równoważność seminorm Schwartza]
\label{lem:schwartz-seminorms17}
Seminormy $s_{\alpha,\beta}$ oraz $Q_{N,m}$ generują tę samą
topologię. Ponadto $Q_n$ są niemalejące, oddzielają punkty
i wyznaczają topologię przestrzeni Schwartza.
\end{lemat}
\begin{proof}
Ponieważ $|x^\alpha|\le |x|^{|\alpha|}\le(1+|x|)^{|\alpha|}$,
mamy $s_{\alpha,\beta}\le Q_{|\alpha|,|\beta|}$.
Z drugiej strony nierówność
$(1+|x|)^N\le C_{d,N}\bigl(1+\sum_{|\alpha|=N}|x^\alpha|\bigr)$
wynika z $|x|\le\sqrt d\max_j|x_j|$ i dwumianu Newtona.
Stąd $Q_{N,m}$ jest z góry ograniczone przez skończoną
sumę odpowiednich $s_{\alpha,\beta}$ (łącznie z $\alpha=0$).
Rodzina $Q_n$ jest współkońcowa względem $Q_{N,m}$,
gdy $n\ge\max\{N,m\}$. Oddzielanie punktów wynika już z $Q_1$.
\end{proof}

\begin{twierdzenie}[Zupełność przestrzeni Schwartza]
\label{thm:schwartz-frechet17}
Dla $\K=\R$ lub $\K=\C$ przestrzeń $\mathcal S(\R^d;\K)$
ze standardową topologią jest przestrzenią Frécheta.
Naturalne włożenie $\mathcal S(\R^d;\K)\hookrightarrow
C^\infty(\R^d;\K)$ jest ciągłe.
\end{twierdzenie}
\begin{proof}
Metryzowalność i lokalna wypukłość wynikają z
Lematu~\ref{lem:schwartz-seminorms17}.
Niech $(f_j)$ będzie ciągiem Cauchy'ego dla wszystkich $Q_n$.
Ponieważ na każdym zwartym zbiorze i dla każdego rzędu pochodnych
seminormy $C^\infty$ są kontrolowane przez odpowiednie $Q_n$,
ciąg jest Cauchy'ego w $C^\infty(\R^d)$.
Na mocy Twierdzenia~\ref{thm:frechet-cinf17} istnieje
$f\in C^\infty(\R^d)$, do której $f_j$ zbiega lokalnie
wraz ze wszystkimi pochodnymi.
Ustalmy $N,m$ i $\varepsilon>0$. Istnieje $j_0$ takie, że
$Q_{N,m}(f_j-f_k)\le\varepsilon$ dla $j,k\ge j_0$.
Przechodząc z $k$ do nieskończoności punktowo w każdej pochodnej,
dostajemy $Q_{N,m}(f_j-f)\le\varepsilon$ dla $j\ge j_0$.
Ponieważ $Q_{N,m}(f_j)<\infty$, wynika stąd
$Q_{N,m}(f)<\infty$. Wszystkie takie seminormy są skończone,
więc $f\in\mathcal S$, a $f_j\to f$ w topologii Schwartza.
Ciągłość włożenia wynika z nierówności
$p_{K,m}(f)\le Q_{0,m}(f)$ dla każdego $K\Subset\R^d$.
\end{proof}

\begin{przyklad}[Funkcje w przestrzeni Schwartza]
\label{ex:schwartz-functions17}
Każda funkcja z $C_c^\infty(\R^d)$ należy do $\mathcal S$.
Dla każdego wielomianu $P$ funkcja
$f(x)=P(x)e^{-|x|^2}$ także należy do $\mathcal S$:
każda jej pochodna ma postać wielomianu pomnożonego przez
$e^{-|x|^2}$. Natomiast $f(x)=(1+|x|^2)^{-1}$ nie należy
do $\mathcal S$, ponieważ dla dostatecznie dużego $N$
wartość $\sup_x(1+|x|)^N|f(x)|$ jest nieskończona.
\end{przyklad}

\begin{stwierdzenie}[Różniczkowanie, wielomiany i normy $L^p$]
\label{prop:schwartz-operators17}
Dla każdego $\alpha\in\N_0^d$ operatory
$f\mapsto\partial^\alpha f$ oraz $f\mapsto x^\alpha f$
są ciągłymi operatorami liniowymi na $\mathcal S(\R^d;\K)$.
Mnożenie $\mathcal S\times\mathcal S\to\mathcal S$ jest
łącznie ciągłe. Dla każdego $1\le p\le\infty$
kanoniczne włożenie $\mathcal S(\R^d;\K)\to L^p(\R^d;\K)$
jest ciągłe.
\end{stwierdzenie}
\begin{proof}
Różniczkowanie spełnia
$Q_{N,m}(\partial^\alpha f)\le Q_{N,m+|\alpha|}(f)$.
Dla mnożenia przez $x^\alpha$ stosujemy wzór Leibniza,
uzyskując oszacowanie przez skończoną sumę seminorm
$Q_{N+|\alpha|,m}(f)$. Analogiczny wzór daje
$Q_{N,m}(fg)\le C_{d,m}Q_{N,m}(f)Q_{0,m}(g)$,
a więc łączną ciągłość iloczynu.
Dla $p=\infty$ mamy $\|f\|_\infty\le Q_{0,0}(f)$.
Dla $1\le p<\infty$ wybieramy liczbę całkowitą $N>d/p$.
Wtedy
$\|f\|_{L^p}^p\le Q_{N,0}(f)^p
\int_{\R^d}(1+|x|)^{-Np}\,\mathrm d\lambda_d(x)$,
a całka jest skończona. To dowodzi ciągłości.
\end{proof}

\begin{twierdzenie}[Gęstość funkcji o zwartym nośniku w $\mathcal S$]
\label{thm:density-schwartz17}
Przestrzeń $C_c^\infty(\R^d;\K)$ jest gęsta w
$\mathcal S(\R^d;\K)$ dla standardowej topologii Schwartza.
\end{twierdzenie}
\begin{proof}
Weźmy $\chi\in C_c^\infty(\R^d;\R)$ taką, że
$0\le\chi\le1$, $\chi=1$ na $\overline B(0,1)$
i $\operatorname{supp}\chi\subseteq B(0,2)$.
Niech $\chi_R(x)=\chi(x/R)$ dla $R\ge1$.
Funkcja $\chi_R f$ ma zwarty nośnik. Dla wielowskaźnika $\beta$
ze wzoru Leibniza
\[
\partial^\beta((1-\chi_R)f)
=(1-\chi_R)\partial^\beta f+
\sum_{0<\gamma\le\beta}\binom{\beta}{\gamma}
 (-1)R^{-|\gamma|}(\partial^\gamma\chi)(x/R)
 \partial^{\beta-\gamma}f.
\]
Pierwszy składnik jest wsparty w $\{|x|\ge R\}$,
a pozostałe w $\{R\le|x|\le2R\}$.
Dla ustalonych $N,m$ wszystkie składniki po pomnożeniu przez
$(1+|x|)^N$ są więc ograniczone przez
$C_{N,m}(1+R)^{-1}Q_{N+1,m}(f)$
(z uwzględnieniem $R^{-|\gamma|}\le1$).
Stąd $Q_{N,m}(f-\chi_Rf)\to0$ przy $R\to\infty$.
\end{proof}

\begin{stwierdzenie}[Brak normowalności]
\label{prop:non-normable17}
Jeśli $d\ge1$ i $\Omega\subseteq\R^d$ jest niepusty i otwarty,
to ani $C^\infty(\Omega;\K)$, ani $\mathcal S(\R^d;\K)$
ze swoimi standardowymi topologiami nie są normowalne.
\end{stwierdzenie}
\begin{proof}
Wystarczy pokazać, że żadne otoczenie zera nie jest ograniczone,
a następnie zastosować kryterium normowalności Kolmogorowa.
Najpierw rozważmy funkcje zespolone. Każde otoczenie zera
w $C^\infty(\Omega)$ zawiera zbiór $U=\{f:q_N(f)<\varepsilon\}$.
Wybierzmy $\psi\in C_c^\infty(\Omega)$ różną od zera,
taką, że jej nośnik leży w pewnym $K_M$ z $M>N+1$.
Niech $f_j(x)=c j^{-N}\psi(x)e^{ijx_1}$.
Wzór Leibniza zapewnia $q_N(f_j)\le Cc$,
więc małym $c>0$ można zagwarantować $f_j\in U$.
Na otwartym podzbiorze, gdzie $\psi\ne0$, pochodna
$\partial_1^{N+1}f_j$ ma wartość bezwzględną rosnącą
co najmniej liniowo z $j$ (główny składnik ma rząd $j$,
a pozostałe są jednostajnie ograniczone). Stąd
$q_M(f_j)\to\infty$, więc $U$ nie jest ograniczony.
Dla przestrzeni Schwartza ten sam ciąg $f_j$, z ustaloną
$\psi\in C_c^\infty(\R^d)$, spełnia
$Q_N(f_j)\le Cc$ oraz $Q_{N+1}(f_j)\to\infty$.
Przypadek rzeczywisty otrzymujemy, stosując
$\cos(jx_1)$ zamiast $e^{ijx_1}$ i supremum na małej kuli,
gdzie $\psi$ jest stała i niezerowa; przy dużym $j$
cosinus lub sinus odpowiedniej pochodnej przyjmuje tam
wartość bezwzględną $1$.
\end{proof}

\subsection{Transformacja Fouriera w przestrzeni Schwartza}

W tej podsekcji pracujemy nad $\C$. Wzór z całką jest zwykłą
całką Lebesgue'a, bez potrzeby odwoływania się do dystrybucji.

\begin{definicja}[Transformacja Fouriera]
\label{def:fourier-schwartz17}
Dla $f\in\mathcal S(\R^d;\C)$ definiujemy
\[
\widehat f(\xi)=\mathcal Ff(\xi)=
(2\pi)^{-d/2}\int_{\R^d}e^{-ix\cdot\xi}f(x)\,
\mathrm d\lambda_d(x),\qquad\xi\in\R^d.
\]
Całka jest bezwzględnie zbieżna, gdyż $\mathcal S\subseteq L^1$.
\end{definicja}

\begin{twierdzenie}[Transformacja Fouriera zachowuje $\mathcal S$]
\label{thm:fourier-schwartz17}
Dla każdego $d\ge1$ transformacja Fouriera jest ciągłym
izomorfizmem liniowo-topologicznym
$\mathcal F:\mathcal S(\R^d;\C)\to\mathcal S(\R^d;\C)$.
Jej odwrotność wyraża wzór
\[
\mathcal F^{-1}g(x)=(2\pi)^{-d/2}
\int_{\R^d}e^{ix\cdot\xi}g(\xi)\,
\mathrm d\lambda_d(\xi).
\]
Dla każdego $f\in\mathcal S$ oraz $j=1,\dots,d$ zachodzą
$\mathcal F(\partial_jf)(\xi)=i\xi_j\widehat f(\xi)$
i $\mathcal F(x_jf)(\xi)=i\partial_{\xi_j}\widehat f(\xi)$.
\end{twierdzenie}
\begin{proof}
Różniczkowanie pod znakiem całki daje
$\partial_\xi^\beta\widehat f=(-i)^{|\beta|}
\mathcal F(x^\beta f)$. Całkowanie przez części, bez wyrazów
brzegowych dzięki szybkiemu zanikaniu funkcji, daje
$\mathcal F(\partial^\alpha f)=(i\xi)^\alpha\widehat f$.
Połączenie tych tożsamości prowadzi do
\[
\xi^\alpha\partial_\xi^\beta\widehat f(\xi)
=(-i)^{|\alpha|+|\beta|}
\mathcal F\bigl(\partial^\alpha(x^\beta f)\bigr)(\xi).
\]
Wartość bezwzględną prawej strony ograniczamy przez
$(2\pi)^{-d/2}\|\partial^\alpha(x^\beta f)\|_{L^1}$.
Wzór Leibniza i wybór wagi $(1+|x|)^{d+1}$ pokazują,
że każde takie wyrażenie jest ograniczone przez skończoną
kombinację seminorm Schwartza funkcji $f$.
Wszystkie seminormy $s_{\alpha,\beta}(\widehat f)$ są więc
skończone i zależą w sposób ciągły od $f$.

Pozostaje udowodnić wzór odwrotny. Dla $\varepsilon>0$ definiujemy
\[
I_\varepsilon(x)=(2\pi)^{-d/2}\int_{\R^d}
 e^{-\varepsilon|\xi|^2/2}e^{ix\cdot\xi}
 \widehat f(\xi)\,\mathrm d\lambda_d(\xi).
\]
Ponieważ $f\in L^1$ i czynnik Gaussa należy do $L^1$,
możemy zastosować twierdzenie Fubiniego. Jednowymiarową całkę Gaussa obliczamy następująco:
$I(u)=\int_\R e^{-\varepsilon t^2/2}e^{iut}\,\mathrm dt$
spełnia $I'(u)=-(u/\varepsilon)I(u)$ na mocy całkowania
przez części oraz $I(0)=\sqrt{2\pi/\varepsilon}$.
Zatem $I(u)=\sqrt{2\pi/\varepsilon}e^{-u^2/(2\varepsilon)}$.
Zastosowanie tego wzoru w każdej współrzędnej daje
\[
I_\varepsilon(x)=\int_{\R^d}f(y)
(2\pi\varepsilon)^{-d/2}e^{-|x-y|^2/(2\varepsilon)}
\,\mathrm d\lambda_d(y).
\]
Jądro po prawej ma całkę $1$ i jego masa poza każdą
kulą wokół zera dąży do zera przy $\varepsilon\downarrow0$.
Ponieważ $f$ jest ograniczona i ciągła, prawa strona
dąży do $f(x)$: dzielimy całkę na $|x-y|<\delta$,
gdzie korzystamy z ciągłości $f$, i dopełnienie,
gdzie korzystamy z ograniczoności $f$ i zanikania masy jądra.
Z drugiej strony $\widehat f\in L^1$, zatem
z twierdzenia o zbieżności zmajoryzowanej
$I_\varepsilon(x)$ dąży do całki ze wzoru odwrotnego.
Otrzymujemy $\mathcal F^{-1}\mathcal Ff=f$.
Ten sam argument, ze zmienionym znakiem wykładnika,
daje $\mathcal F\mathcal F^{-1}g=g$ dla $g\in\mathcal S$.
Ciągłość operatora odwrotnego wynika z identycznych
oszacowań seminorm po zmianie znaku wykładnika.
\end{proof}

\begin{przyklad}[Funkcja Gaussa i wizualizacja w Plotly]
\label{ex:fourier-gauss17}
Dla $a>0$ oraz $f_a(x)=e^{-ax^2/2}$, $x\in\R$,
otrzymujemy $\widehat f_a(\xi)=a^{-1/2}e^{-\xi^2/(2a)}$.
W szczególności zwiększenie skupienia funkcji $f_a$
odpowiada rozszerzeniu jej transformaty Fouriera.
Poniższy kod Python rysuje obie funkcje z użyciem biblioteki Plotly.
\begin{lstlisting}[language=Python,caption={Para funkcji Gaussa: funkcja i transformata Fouriera}]
import numpy as np
import plotly.graph_objects as go

a = 2.0
x = np.linspace(-5.0, 5.0, 1000)
f = np.exp(-a * x**2 / 2)
f_hat = a**(-0.5) * np.exp(-x**2 / (2 * a))

fig = go.Figure()
fig.add_scatter(x=x, y=f, mode="lines", name="f(x)")
fig.add_scatter(x=x, y=f_hat, mode="lines", name="Fourier(f)(x)")
fig.update_layout(title="Para Gaussa", xaxis_title="x / xi",
                  yaxis_title="wartosc")
fig.show()
\end{lstlisting}
\end{przyklad}

\begin{stwierdzenie}[Splot funkcji szybko malejących]
\label{prop:convolution-schwartz17}
Dla $f,g\in\mathcal S(\R^d;\C)$ definiujemy
$(f*g)(x)=\int_{\R^d}f(y)g(x-y)\,\mathrm d\lambda_d(y)$.
Wówczas $f*g\in\mathcal S$, odwzorowanie
$(f,g)\mapsto f*g$ jest dwuliniowe i łącznie ciągłe
na $\mathcal S\times\mathcal S$, a ponadto
$\mathcal F(f*g)=(2\pi)^{d/2}\widehat f\,\widehat g$.
\end{stwierdzenie}
\begin{proof}
Z $\mathcal S\subset L^1\cap L^\infty$ wynika
bezwzględna zbieżność splotu. Różniczkowanie pod całką
daje $\partial^\beta(f*g)=f*(\partial^\beta g)$.
Nierówność $1+|x|\le(1+|y|)(1+|x-y|)$ prowadzi do
\[
 Q_{N,m}(f*g)\le Q_{N,m}(g)
 \int_{\R^d}(1+|y|)^N|f(y)|\,\mathrm d\lambda_d(y)
 \le C_{d,N}Q_{N+d+1,0}(f)Q_{N,m}(g).
\]
To dowodzi jednocześnie szybkiego zanikania i ciągłości
splotu. Wzór Fouriera otrzymujemy z twierdzenia Fubiniego:
w całce podwójnej podstawiamy $z=x-y$ i rozdzielamy
czynniki $e^{-i(y+z)\cdot\xi}$; przy naszej normalizacji
pojawia się współczynnik $(2\pi)^{d/2}$.
\end{proof}

\begin{uwaga}[Przygotowanie teorii dystrybucji]
\label{rem:distribution-next17}
Ciągłe funkcjonały na $C^\infty(\Omega)$ mają oszacowania
na pochodne na jednym zwartym zbiorze (Twierdzenie~\ref{thm:dual-cinf17}).
Dla funkcjonału $T:\mathcal S(\R^d)\to\K$ analogiczny warunek
ciągłości ma postać $|T(f)|\le C Q_{N,m}(f)$
dla pewnych $C,N,m$ i wszystkich $f\in\mathcal S$.
Nie są to jednak te same przestrzenie dualne topologiczne.
W następnym rozdziale skonstruujemy $\mathcal D(\Omega)$
z topologią granicy induktywnej przestrzeni funkcji o nośniku
zawartym w ustalonym zbiorze zwartym. Dopiero wtedy zdefiniujemy
ogólną przestrzeń dystrybucji $\mathcal D'(\Omega)$.
\end{uwaga}


% Wersja źródłowa: audyt II, jednostka 18
\setcounter{section}{17}
\section{Funkcje testowe i ścisłe granice induktywne}
\label{sec:funkcje-testowe18}

\subsection{Funkcje o ustalonym zwartym nośniku}

W całym rozdziale zakładamy, że $d\ge1$, $\Omega\subseteq\R^d$
jest niepustym zbiorem otwartym, a $\K\in\{\R,\C\}$.
Dla $f\in C^\infty(\Omega;\K)$ definiujemy jego \emph{nośnik}
przez $\supp f=\overline{\{x\in\Omega:f(x)\ne0\}}^{\,\Omega}$,
tj. domknięcie względne w $\Omega$. Zapis $K\Subset\Omega$
oznacza, że $K$ jest zwartym podzbiorem $\Omega$.

\begin{definicja}[Przestrzeń $\mathcal D_K(\Omega)$]
\label{def:DK18}
Dla każdego zwartego $K\Subset\Omega$ kładziemy
\[
 \mathcal D_K(\Omega;\K)=\{\varphi\in C^\infty(\Omega;\K):
 \supp\varphi\subseteq K\}.
\]
Wyposażamy tę przestrzeń w topologię wyznaczoną przez seminormy
\[p_m(\varphi)=\max_{|\alpha|\le m}\sup_{x\in\Omega}
|\partial^\alpha\varphi(x)|,\qquad m\in\N_0.\]
Supremum jest skończone: wszystkie pochodne znikają poza $K$.
Jeżeli $\operatorname{int}K=\varnothing$, to $\mathcal D_K=\{0\}$.
\end{definicja}

\begin{twierdzenie}[Struktura etapów]
\label{thm:DK-frechet18}
Dla każdego zwartego $K\Subset\Omega$ przestrzeń
$\mathcal D_K(\Omega;\K)$ jest przestrzenią Fr\'echeta.
Jeżeli $K\subseteq L\Subset\Omega$ są zwarte, to
$\mathcal D_K$ jest domkniętą podprzestrzenią $\mathcal D_L$,
a włożenie $\mathcal D_K\hookrightarrow\mathcal D_L$
jest izomorfizmem na swój obraz z topologią względną.
\end{twierdzenie}
\begin{proof}
Z poprzedniego rozdziału wiemy, że
$C^\infty(\Omega;\K)$ jest przestrzenią Fr\'echeta.
Przestrzeń $\mathcal D_K$ jest w niej domknięta:
przy zbieżności $\varphi_j\to\varphi$ w $C^\infty(\Omega)$
i $\varphi_j\in\mathcal D_K$ mamy $\varphi(x)=0$ dla
$x\in\Omega\setminus K$, więc $\supp\varphi\subseteq K$.
Topologia względna na $\mathcal D_K$ pokrywa się z topologią
wyznaczoną przez $p_m$, ponieważ dla funkcji wspartych w $K$
wszystkie pochodne znikają poza $K$ i $p_m=p_{K,m}$
w oznaczeniach poprzedniego rozdziału.
Domknięta podprzestrzeń przestrzeni Fr\'echeta jest przestrzenią
Fr\'echeta. Ten sam argument, zastosowany do $K\subseteq L$,
daje domkniętość $\mathcal D_K$ w $\mathcal D_L$;
seminormy $p_m$ są na obu etapach określone tym samym wzorem.
\end{proof}

\begin{definicja}[Algebraiczna przestrzeń funkcji testowych]
\label{def:D-algebraic18}
Oznaczamy przez $\mathcal D(\Omega;\K)=C_c^\infty(\Omega;\K)$
przestrzeń wszystkich funkcji gładkich o zwartym nośniku.
Jeśli $(K_n)_{n\ge1}$ jest wyczerpaniem zwartymi zbiorami,
$K_n\subseteq\operatorname{int}K_{n+1}$ i każdy zwarty
$K\Subset\Omega$ jest zawarty w pewnym $K_n$, to
$\mathcal D(\Omega)=\bigcup_{n\ge1}\mathcal D_{K_n}(\Omega)$.
Takie wyczerpanie skonstruowano w
wcześniejszym lemacie o wyczerpaniu zwartym.
\end{definicja}

\subsection{Topologia granicy induktywnej}

\begin{definicja}[Lokalnie wypukła granica induktywna]
\label{def:LF18}
Niech $(E_n,\tau_n)$ będzie rosnącym ciągiem przestrzeni
lokalnie wypukłych nad $\K$, gdzie każde włożenie
$i_n:E_n\hookrightarrow E_{n+1}$ jest liniowe i ciągłe.
Na $E=\bigcup_nE_n$ rozważamy \emph{najsilniejszą lokalnie
wypukłą topologię}, dla której wszystkie odwzorowania
$E_n\hookrightarrow E$ są ciągłe. Oznaczamy ją
$\operatorname{ind}_n E_n$ i nazywamy topologią granicy induktywnej.
Jeżeli $E_n$ są przestrzeniami Fr\'echeta, otrzymujemy
przestrzeń typu $(LF)$; granica jest \emph{ścisła}, jeśli
każde $E_n$ ma topologię względną z $E_{n+1}$ i jest w nim
podprzestrzenią domkniętą.
\end{definicja}

\begin{twierdzenie}[Konstrukcja przez seminormy i własność uniwersalna]
\label{thm:LF-universal18}
Niech $E=\bigcup_nE_n$, gdzie każde $E_n$ jest przestrzenią
lokalnie wypukłą, a włożenia są ciągłe. Rozważmy rodzinę
$\mathcal P$ wszystkich seminorm $p:E\to[0,\infty)$ takich,
że $p|_{E_n}$ jest ciągła dla każdego $n$.
Topologia generowana przez $\mathcal P$ jest topologią
granicy induktywnej. Dla dowolnej przestrzeni lokalnie
wypukłej $F$ i odwzorowania liniowego $T:E\to F$ mamy
\[
 T\text{ jest ciągłe}\quad\Longleftrightarrow\quad
 T|_{E_n}:E_n\to F\text{ jest ciągłe dla każdego }n.
\]
\end{twierdzenie}
\begin{proof}
Dla $p\in\mathcal P$ włożenie $E_n\hookrightarrow E$ jest
ciągłe w topologii wyznaczonej przez $p$, więc również
w topologii wszystkich tych seminorm. Niech $\rho$ będzie
inną lokalnie wypukłą topologią na $E$, dla której włożenia
są ciągłe. Każda seminorma ciągła dla $\rho$ należy do
$\mathcal P$, gdyż jej ograniczenie do $E_n$ jest ciągłe.
Stąd $\rho$ jest nie silniejsza od naszej topologii.
Dla ostatniej równoważności wystarczy zastosować ten
argument do seminorm $q\circ T$, gdzie $q$ przebiega
wszystkie ciągłe seminormy na $F$.
\end{proof}

\begin{definicja}[Standardowa topologia $\mathcal D(\Omega)$]
\label{def:D-LF18}
Na $\mathcal D(\Omega;\K)$ przyjmujemy topologię
\[
 \mathcal D(\Omega;\K)
   =\operatorname{ind}_{n\to\infty}\mathcal D_{K_n}(\Omega;\K),
\]
gdzie $(K_n)$ jest wyczerpaniem z
Definicji~\ref{def:D-algebraic18}. Jest to \emph{standardowa
topologia funkcji testowych}. Nie należy jej utożsamiać
ani z topologią podprzestrzeni dziedziczoną po
$C^\infty(\Omega)$, ani automatycznie z końcową topologią
w kategorii wszystkich przestrzeni topologicznych.
\end{definicja}

\begin{stwierdzenie}[Niezależność od wyczerpania]
\label{prop:D-independence18}
Dwa wyczerpania $\Omega$ zbiorami zwartymi współkońcowe
wśród zwartych podzbiorów $\Omega$ prowadzą do tej samej
topologii na $\mathcal D(\Omega)$, nawet jeśli nie mają
identycznych wyrazów.
\end{stwierdzenie}
\begin{proof}
Niech $(K_n)$, $(L_j)$ będą takimi wyczerpaniami.
Seminorma $p$ na $\mathcal D(\Omega)$ jest ciągła na
każdym $\mathcal D_{K_n}$ wtedy i tylko wtedy, gdy jest
ciągła na każdym $\mathcal D_{L_j}$: dla ustalonego $K_n$
istnieje $L_j$ zawierający $K_n$, a ograniczenie seminormy
ciągłej na $\mathcal D_{L_j}$ do $\mathcal D_{K_n}$ jest
ciągłe. W drugą stronę argument jest symetryczny.
Korzystamy z Twierdzenia~\ref{thm:LF-universal18}.
\end{proof}

\begin{twierdzenie}[Topologia na każdym etapie i oddzielanie punktów]
\label{thm:D-stage18}
Dla dowolnego zwartego $K\Subset\Omega$ topologia względna
indukowana na $\mathcal D_K$ przez $\mathcal D(\Omega)$
jest topologią jego przestrzeni Fr\'echeta.
Ponadto $\mathcal D(\Omega)$ jest lokalnie wypukłą
przestrzenią liniowo-topologiczną Hausdorffa.
\end{twierdzenie}
\begin{proof}
Ciągłość włożenia $\mathcal D_K\to\mathcal D$ mówi, że
topologia względna nie jest silniejsza od topologii
$\mathcal D_K$. Dla każdego $m\in\N_0$ określmy na
całej $\mathcal D$ seminormę
$p_m(\varphi)=\max_{|\alpha|\le m}\sup_{x\in\Omega}
|\partial^\alpha\varphi(x)|$. Jej ograniczenie do
każdego $\mathcal D_{K_n}$ jest ciągłe, a zatem $p_m$
jest ciągła na $\mathcal D$ z
Twierdzenia~\ref{thm:LF-universal18}. Seminormy te
wyznaczają topologię względną $\mathcal D_K$, stąd
otrzymujemy także nierówność w przeciwnym kierunku.
Ponieważ $p_0(\varphi)>0$ dla $\varphi\ne0$, topologia
jest Hausdorffa; lokalna wypukłość wynika z konstrukcji.
\end{proof}

\begin{uwaga}[Przestrzeń dualna topologiczna]
\label{rem:D-dual18}
Własność uniwersalna jest szczególnie użyteczna dla
funkcjonałów: odwzorowanie liniowe $T:\mathcal D(\Omega)\to\K$
jest ciągłe wtedy i tylko wtedy, gdy każde ograniczenie
$T|_{\mathcal D_K}$ jest ciągłe. Ponieważ $\mathcal D_K$
ma topologię seminorm $p_m$, równoważnie: dla każdego
$K\Subset\Omega$ istnieją $C_K>0$ i $m_K\in\N_0$ takie, że
\[
 |T(\varphi)|\le C_K p_{m_K}(\varphi)
 \qquad(\varphi\in\mathcal D_K(\Omega)).
\]
Wielkości $C_K$ i $m_K$ mogą zależeć od $K$.
To będzie podstawowe kryterium w teorii dystrybucji.
\end{uwaga}


\begin{przyklad}[Dystrybucja o nieograniczonym rzędzie globalnym]
\label{ex:variable-order18}
Niech $\Omega=\R$. Dla $\varphi\in\mathcal D(\R)$ określamy
$T(\varphi)=\sum_{n=1}^\infty \varphi^{(n)}(n)$.
Suma jest skończona dla każdej $\varphi$, gdyż tylko skończenie
wiele liczb naturalnych może należeć do zwartego nośnika.
Na każdym $\mathcal D_K$ funkcjonał jest sumą skończonej
liczby ciągłych ewaluacji pochodnych, więc $T$ jest ciągły
na $\mathcal D(\R)$.

Nie istnieje jednak jeden rząd $m$ właściwy dla wszystkich
kompaktów, nawet z osobno dobranymi stałymi $C_K$.
Ustalmy $m$ i wybierzmy liczbę całkowitą $n>m$.
Niech $\chi\in C_c^\infty((-1/4,1/4))$ będzie równa $1$
w pobliżu zera, a $\psi(t)=t^n\chi(t)$. Dla $0<\varepsilon<1$
kładziemy $\varphi_\varepsilon(x)=\varepsilon^n
\psi((x-n)/\varepsilon)$. Wszystkie te funkcje mają
nośniki w jednym zwartym przedziale wokół $n$, przy czym
$T(\varphi_\varepsilon)=\psi^{(n)}(0)=n!$,
natomiast $p_m(\varphi_\varepsilon)\to0$.
To wyklucza oszacowanie z jednym globalnie ustalonym $m$.
\end{przyklad}

\subsection{Zbiory ograniczone, ciągi i pełna zupełność}

\begin{twierdzenie}[Dokładne kryterium ograniczoności]
\label{thm:D-bounded18}
Niech $B\subseteq\mathcal D(\Omega;\K)$. Następujące
warunki są równoważne:
\begin{enumerate}
\item $B$ jest ograniczony w topologii $\mathcal D(\Omega)$;
\item istnieje zwarty $K\Subset\Omega$, który zawiera
nośniki wszystkich $\varphi\in B$, a ponadto dla każdego
$m\in\N_0$ zachodzi $\sup_{\varphi\in B}p_m(\varphi)<\infty$.
\end{enumerate}
\end{twierdzenie}
\begin{proof}
Jeżeli $B\subseteq\mathcal D_K$ jest ograniczony w
$\mathcal D_K$, to jego obraz przez ciągłe włożenie
jest ograniczony w $\mathcal D$.

Na odwrót, niech $B$ będzie ograniczony w $\mathcal D$.
Ponieważ każda seminorma $p_m$ z dowodu
Twierdzenia~\ref{thm:D-stage18} jest ciągła, mamy
$\sup_{\varphi\in B}p_m(\varphi)<\infty$.
Pozostaje wykazać wspólność nośnika. Załóżmy przeciwnie,
że $B$ nie zawiera się w żadnym $\mathcal D_{K_n}$.
Dla każdego $n$ wybierzmy $\varphi_n\in B$ oraz
$x_n\in\Omega\setminus K_n$ z $\varphi_n(x_n)\ne0$.
Punkty $x_n$ opuszczają każdy zwarty podzbiór $\Omega$,
ponieważ wyczerpanie jest współkońcowe. Możemy więc
wybrać gładkie funkcje nieujemne $\theta_n$, równe $1$
w $x_n$, o nośnikach tworzących rodzinę lokalnie skończoną
w $\Omega$. Ustalmy $a_n>n/|\varphi_n(x_n)|$ i określmy
$w(x)=1+\sum_{n\ge1}a_n\theta_n(x)$. Suma jest lokalnie
skończona, zatem $w$ jest dodatnią funkcją gładką.
Seminorma $p_w(\varphi)=\sup_{x\in\Omega}w(x)|\varphi(x)|$
jest skończona dla każdej funkcji testowej. Dla ustalonego
$K$ mamy $p_w|_{\mathcal D_K}\le(\sup_Kw)p_0$, więc
z Twierdzenia~\ref{thm:LF-universal18} wynika jej ciągłość
na $\mathcal D$. Jednak $p_w(\varphi_n)>n$,
co przeczy ograniczoności $B$. Wobec tego istnieje
$K_n$ zawierający wszystkie nośniki, a oszacowania
$p_m$ kończą dowód.
\end{proof}

\begin{twierdzenie}[Zbieżność ciągów funkcji testowych]
\label{thm:D-sequences18}
Dla ciągu $(\varphi_j)\subseteq\mathcal D(\Omega;\K)$
i $\varphi\in\mathcal D(\Omega;\K)$ następujące warunki
są równoważne:
\begin{enumerate}
\item $\varphi_j\to\varphi$ w topologii $\mathcal D(\Omega)$;
\item istnieje zwarty $K\Subset\Omega$ zawierający
$\supp\varphi$ i wszystkie $\supp\varphi_j$ oraz
$p_m(\varphi_j-\varphi)\to0$ dla każdego $m\in\N_0$.
\end{enumerate}
\end{twierdzenie}
\begin{proof}
Ciąg zbieżny wraz z granicą jest zbiorem ograniczonym
w dowolnej przestrzeni liniowo-topologicznej: dla ustalonego
otoczenia zera najpierw wchłaniamy jego wyrazy od pewnego
indeksu, a następnie skończenie wiele pozostałych.
Twierdzenie~\ref{thm:D-bounded18} daje wspólny zwarty $K$.
Ponieważ topologia na $\mathcal D_K$ jest względna,
zbieżność w $\mathcal D$ oznacza zbieżność w $\mathcal D_K$,
a tę charakteryzują seminormy $p_m$.
Implikacja odwrotna wynika z ciągłości
$\mathcal D_K\hookrightarrow\mathcal D$.
\end{proof}

\begin{wniosek}[Zupełność sekwencyjna]
\label{cor:D-seqcomplete18}
Każdy ciąg Cauchy'ego w $\mathcal D(\Omega;\K)$ jest
zbieżny w tej przestrzeni.
\end{wniosek}
\begin{proof}
Każdy ciąg Cauchy'ego jest ograniczony w przestrzeni lokalnie
wypukłej: dla każdej ciągłej seminormy jego wartości tworzą
ciąg Cauchy'ego, a więc ograniczony.
Wybieramy wspólny $K$ z
Twierdzenia~\ref{thm:D-bounded18}. Ograniczenia topologii
na $\mathcal D_K$ i $\mathcal D$ są jednakowe, więc ciąg
jest Cauchy'ego w przestrzeni Fr\'echeta $\mathcal D_K$;
tam ma granicę, do której zbiega również w $\mathcal D$.
\end{proof}

\begin{twierdzenie}[Zupełność $\mathcal D(\Omega)$]
\label{thm:D-complete18}
Przestrzeń $\mathcal D(\Omega;\K)$ jest zupełna:
każdy filtr Cauchy'ego względem jej kanonicznej struktury
jednostajnej jest zbieżny. W szczególności jest zupełną
ścisłą przestrzenią typu $(LF)$.
\end{twierdzenie}
\begin{proof}
Przeprowadzimy dowód dla ogólnej rosnącej sumy
$E=\bigcup_nE_n$ przestrzeni zupełnych lokalnie wypukłych,
która ma topologię granicy induktywnej i indukuje na
każdym $E_n$ topologię wyjściową. W naszym przypadku
$E_n=\mathcal D_{K_n}$, a wszystkie te własności już
udowodniono.

Niech $\mathcal F$ będzie filtrem Cauchy'ego w $E$.
Przez $\mathcal N$ oznaczmy filtr otoczeń zera, a przez
$\mathcal G$ filtr generowany przez zbiory $F+V$,
$F\in\mathcal F$, $V\in\mathcal N$. Filtr $\mathcal G$
jest filtrem Cauchy'ego, słabszym od $\mathcal F$.
Pokażemy, że dla pewnego $n$ każdy element $\mathcal G$
przecina $E_n$.

Załóżmy przeciwnie. Dla każdego $n$ istnieją
$F_n\in\mathcal F$ i absolutnie wypukłe otoczenie
$V_n$ zera w $E$ takie, że
$(F_n+V_n)\cap E_n=\varnothing$. Zastępując $V_n$
odpowiednimi skończonymi przekrojami, możemy zakładać
$V_{n+1}\subseteq V_n$. Definiujemy
$W=\operatorname{aco}(\bigcup_n(V_n\cap E_n))$,
gdzie $\operatorname{aco}$ oznacza absolutnie wypukłą
otoczkę zbioru. Zbiór $W$ jest absolutnie wypukły,
pochłaniający w $E$, a dla każdego $n$ zbiór
$W\cap E_n$ zawiera otoczenie $V_n\cap E_n$ zera w $E_n$.
Z konstrukcji topologii granicy induktywnej wynika, że
$W$ jest otoczeniem zera: jego funkcjonał Minkowskiego
jest ciągły na każdym etapie i zatem na $E$.

Dla ustalonego $n$ składniki sumy z indeksami $k\le n$
należą do $E_n$, natomiast te z $k>n$ należą do $V_n$.
Dlatego $W\subseteq E_n+V_n$ i
$(F_n+W)\cap E_n=\varnothing$.
Ponieważ $\mathcal F$ jest filtrem Cauchy'ego, istnieje
$F_0\in\mathcal F$ takie, że $F_0-F_0\subseteq W$.
Wybierzmy $x\in F_0$; należy on do pewnego $E_n$.
Istnieje $y\in F_0\cap F_n$, bo oba zbiory należą do
filtru $\mathcal F$. Ale
$x=y+(x-y)\in(F_n+W)\cap E_n$, sprzeczność.

Zatem istnieje $n$ takie, że każdy element $\mathcal G$
przecina $E_n$. Przecięcia tych elementów z $E_n$
wyznaczają filtr Cauchy'ego na $E_n$, ponieważ topologia
$E_n$ jest indukowana przez $E$, a $\mathcal G$
jest Cauchy'ego. Zupełność $E_n$ daje jego granicę
$x\in E_n$. Punkt $x$ jest punktem skupienia filtru
$\mathcal G$ w $E$. Filtr Cauchy'ego mający punkt skupienia
jest zbieżny do tego punktu: dla dowolnego otoczenia
$U$ zera bierzemy $V$ z $V+V\subseteq U$ oraz
$G\in\mathcal G$ z $G-G\subseteq V$; skoro
$(x+V)\cap G\ne\varnothing$, to $G\subseteq x+U$.
Stąd $\mathcal G\to x$, a ponieważ $\mathcal F$ jest
silniejszy, również $\mathcal F\to x$.
\end{proof}

\subsection{Dlaczego topologia \texorpdfstring{$\mathcal D$}{D} jest istotna?}

\begin{przyklad}[Ucieczka nośników]
\label{ex:escaping18}
Wybierzmy ciąg niezerowych funkcji $\psi_n\in\mathcal D(\Omega)$
mających nośniki poza $K_n$. Każda taka funkcja istnieje,
ponieważ $\Omega\setminus K_n$ jest niepustym zbiorem otwartym.
Dla dowolnego zwartego $K\Subset\Omega$ i ustalonego $m$
mamy $p_{K,m}(\psi_n)=0$ dla dostatecznie dużych $n$.
Zatem $\psi_n\to0$ w $C^\infty(\Omega)$.
Nie zachodzi jednak $\psi_n\to0$ w $\mathcal D(\Omega)$,
ponieważ nośniki nie mieszczą się w żadnym wspólnym
zbiorze zwartym. Właśnie dlatego topologia funkcji testowych
jest ściśle silniejsza od topologii względnej z $C^\infty$.
\end{przyklad}

\begin{twierdzenie}[Niemetryzowalność]
\label{thm:D-nonmetric18}
Dla każdego niepustego otwartego $\Omega\subseteq\R^d$,
$d\ge1$, przestrzeń $\mathcal D(\Omega;\K)$ z jej
standardową topologią nie jest metryzowalna.
\end{twierdzenie}
\begin{proof}
Każda podprzestrzeń $E_n=\mathcal D_{K_n}$ jest domknięta
w $\mathcal D$: istotnie,
$E_n=\bigcap_{x\in\Omega\setminus K_n}\{\varphi:\varphi(x)=0\}$,
a ewaluacje $\varphi\mapsto\varphi(x)$ są ciągłe
na każdym etapie, więc i na $\mathcal D$.
Każde $E_n$ jest właściwe i jako podprzestrzeń liniowa
ma puste wnętrze w $\mathcal D$.
Gdyby $\mathcal D$ miała przeliczalną malejącą bazę
otoczeń zera $(U_n)$, moglibyśmy wybrać
$\varphi_n\in U_n\setminus E_n$. Wtedy
$\varphi_n\to0$, lecz nośniki $\varphi_n$ nie mieściłyby
się w żadnym wspólnym zbiorze zwartym, wbrew
Twierdzeniu~\ref{thm:D-sequences18}.
\end{proof}

\begin{twierdzenie}[Względna zwartość zbiorów ograniczonych]
\label{thm:D-Montel18}
Każdy zbiór ograniczony $B\subseteq\mathcal D(\Omega;\K)$
jest względnie zwarty w standardowej topologii.
\end{twierdzenie}
\begin{proof}
Z Twierdzenia~\ref{thm:D-bounded18} istnieje zwarty $K$
taki, że $B\subseteq\mathcal D_K$ i wszystkie seminormy
$p_m$ są jednostajnie ograniczone na $B$.
Przedłużmy funkcje z $\mathcal D_K$ zerem poza $\Omega$
do funkcji gładkich na $\R^d$. Dla każdej pochodnej
$\partial^\alpha$ pochodne rzędu $|\alpha|+1$ mają
jednostajną normę supremum, toteż rodzina
$\{\partial^\alpha\varphi:\varphi\in B\}$ jest
jednostajnie ograniczona i równociągła na $\R^d$
(z twierdzenia o wartości średniej). Na zwartym $K$
możemy zastosować twierdzenie Arzel\`a--Ascoliego.
Wybierając kolejno podciągi dla przeliczalnej rodziny
wielowskaźników i stosując argument diagonalny,
uzyskujemy podciąg dowolnego ciągu elementów $B$,
którego wszystkie pochodne są jednostajnie zbieżne na $K$.
Poza $K$ są zerowe. Z zupełności $\mathcal D_K$
wynika istnienie granicy w $\mathcal D_K$, a więc także
w $\mathcal D$. Skoro $\mathcal D_K$ jest metryzowalna,
sekwencyjna względna zwartość daje względną zwartość
w $\mathcal D_K$, a ciągłe włożenie do $\mathcal D$
przenosi zwartość domknięcia.
\end{proof}

\begin{stwierdzenie}[Naturalne operatory na funkcjach testowych]
\label{prop:D-operators18}
Niech $a\in C^\infty(\Omega;\K)$ i
$\alpha\in\N_0^d$. Odwzorowania
$\varphi\mapsto\partial^\alpha\varphi$ oraz
$\varphi\mapsto a\varphi$ są ciągłymi operatorami
liniowymi $\mathcal D(\Omega)\to\mathcal D(\Omega)$.
Dla $f\in L^1_{\mathrm{loc}}(\Omega;\K)$ odwzorowanie
$T_f(\varphi)=\int_\Omega f(x)\varphi(x)\,\mathrm d\lambda_d(x)$
jest ciągłym funkcjonałem liniowym na $\mathcal D(\Omega)$.
\end{stwierdzenie}
\begin{proof}
Każdy z dwóch pierwszych operatorów zachowuje
$\mathcal D_K$ i na nim jest ciągły: z reguły Leibniza
na $K$ dostajemy oszacowania za pomocą skończenie
wielu seminorm $p_m$. Ciągłość na $\mathcal D$ wynika
z Twierdzenia~\ref{thm:LF-universal18}.
Dla $\varphi\in\mathcal D_K$ mamy
$|T_f(\varphi)|\le
\bigl(\int_K|f(x)|\,\mathrm d\lambda_d(x)\bigr)p_0(\varphi)$,
a więc każde ograniczenie $T_f|_{\mathcal D_K}$
jest ciągłe. Ponownie stosujemy własność uniwersalną.
\end{proof}

\begin{center}
\begin{tikzpicture}[scale=0.82]
 \draw[thick,rounded corners=14pt] (-4,-1.7) rectangle (4,1.8);
 \draw[thick,rounded corners=12pt,blue!55!black] (-3.0,-1.25) rectangle (2.6,1.3);
 \draw[thick,rounded corners=10pt,blue!55!black] (-2.1,-0.85) rectangle (1.15,0.85);
 \draw[thick,rounded corners=7pt,blue!55!black] (-1.2,-0.43) rectangle (0.25,0.43);
 \node at (-0.48,0) {$K_1$};
 \node at (0.75,0.64) {$K_2$};
 \node at (-2.70,1.02) {$K_3$};
 \node at (-3.55,1.49) {$\Omega$};
 \fill[orange!75!black] (3.32,-1.02) circle (3pt);
 \node[orange!65!black,anchor=west,font=\footnotesize] at (3.41,-1.05) {$\psi_n$};
\end{tikzpicture}

\small\emph{Schemat:} funkcje testowe mogą mieć nośniki poza
kolejnymi zbiorami zwartymi, a wtedy nie tworzą ciągu zbieżnego
w topologii $\mathcal D(\Omega)$, nawet jeśli zbiegają do zera
w $C^\infty(\Omega)$.
\end{center}

\begin{uwaga}[Przejście do dystrybucji]
\label{rem:D-distributions-next18}
W kolejnym rozdziale zdefiniujemy przestrzeń dystrybucji
$\mathcal D'(\Omega)$ jako przestrzeń dualną topologiczną
$\mathcal D(\Omega)$. Podstawowe oszacowania funkcjonałów
uzyskaliśmy już w Uwadze~\ref{rem:D-dual18}; będziemy też
odróżniać topologię słabą-* od topologii zbieżności
jednostajnej na zbiorach ograniczonych funkcji testowych.
\end{uwaga}


% Wersja źródłowa: audyt II, jednostka 19
\setcounter{section}{18}
\section{Przestrzeń dystrybucji i jej topologie}
\label{sec:dystrybucje19}

\subsection{Definicja i pierwsze przykłady}

W całym rozdziale zakładamy, że $d\ge1$, $\Omega\subseteq\R^d$
jest niepustym zbiorem otwartym, a $\K\in\{\R,\C\}$.
Przestrzeń $\mathcal D(\Omega;\K)$ ma topologię skonstruowaną
w \hyperref[sec:funkcje-testowe18]{rozdziale~18}. W szczególności
jest lokalnie wypukła, Hausdorffa i zupełna, ale nie jest
metryzowalna. Dla $K\Subset\Omega$ zapisujemy
$\mathcal D_K=\{\varphi\in\mathcal D(\Omega):\supp\varphi\subseteq K\}$
oraz $p_m(\varphi)=\max_{|\alpha|\le m}\sup_{x\in\Omega}
|\partial^\alpha\varphi(x)|$. Symbolem $\mathcal D'(\Omega;\K)$
oznaczamy zbiór \emph{ciągłych $\K$-liniowych} funkcjonałów
na $\mathcal D(\Omega;\K)$.

\begin{definicja}[Dystrybucja]
\label{def:distribution19}
\emph{Dystrybucją} na $\Omega$ nazywamy element
$T\in\mathcal D'(\Omega;\K)$; wartość $T$ na funkcji
testowej $\varphi$ zapisujemy $\langle T,\varphi\rangle$.
Przyjmujemy konwencję $\K$-liniowości względem $\varphi$:
$\langle T,a\varphi+b\psi\rangle=a\langle T,\varphi\rangle+
b\langle T,\psi\rangle$. Nie wprowadzamy tutaj sprzężenia
zespolonego w drugim argumencie.
\end{definicja}

\begin{twierdzenie}[Lokalne oszacowanie charakteryzujące dystrybucje]
\label{thm:distribution-estimate19}
Niech $T:\mathcal D(\Omega;\K)\to\K$ będzie odwzorowaniem
liniowym. Następujące warunki są równoważne:
\begin{enumerate}
\item $T$ jest ciągły w standardowej topologii $\mathcal D(\Omega)$;
\item dla każdego $K\Subset\Omega$ istnieją $C_K\ge0$
i $m_K\in\N_0$ takie, że
\[|\langle T,\varphi\rangle|\le C_K p_{m_K}(\varphi)
\qquad(\varphi\in\mathcal D_K(\Omega)).\]
\end{enumerate}
\end{twierdzenie}
\begin{proof}
Z własności uniwersalnej granicy induktywnej
(twierdzenie o własności uniwersalnej z \hyperref[sec:funkcje-testowe18]{rozdziału~18}) ciągłość $T$
jest równoważna ciągłości każdego
$T|_{\mathcal D_K}$. Topologia $\mathcal D_K$ jest wyznaczona
przez rosnącą rodzinę seminorm $(p_m)_{m\in\N_0}$.
Jeśli $T|_{\mathcal D_K}$ jest ciągły, istnieje $m$ oraz
$\varepsilon>0$ takie, że $p_m(\varphi)<\varepsilon$
implikuje $|T(\varphi)|<1$. Dla $p_m(\varphi)>0$ stosujemy
tę implikację do odpowiedniego wielokrotnego skalara
$\varphi$ i uzyskujemy $|T(\varphi)|\le(2/\varepsilon)p_m(\varphi)$.
Jeśli $p_m(\varphi)=0$, skalowanie przez dowolną liczbę
pokazuje $T(\varphi)=0$. Odwrotnie, podane oszacowanie
zapewnia ciągłość na każdym etapie, a więc na całej
$\mathcal D(\Omega)$.
\end{proof}

\begin{definicja}[Rząd dystrybucji na zwartym zbiorze]
\label{def:order19}
Dla $T\in\mathcal D'(\Omega)$ i $K\Subset\Omega$ mówimy,
że $T$ ma \emph{rząd nie większy niż $m$ na $K$}, jeżeli
istnieje $C_K\ge0$ takie, że
$|\langle T,\varphi\rangle|\le C_K p_m(\varphi)$
dla $\varphi\in\mathcal D_K$. Mówimy, że $T$ ma
\emph{rząd skończony na $\Omega$}, jeżeli można dobrać
jedno $m$ właściwe dla wszystkich $K\Subset\Omega$
(przy czym $C_K$ może zależeć od $K$).
\end{definicja}

\begin{uwaga}[Zależność od zbioru zwartego]
\label{rem:order-local19}
Twierdzenie~\ref{thm:distribution-estimate19} gwarantuje
\emph{lokalnie} skończony rząd każdej dystrybucji,
nie zaś istnienie jednego rzędu na całej $\Omega$.
Przykład z \hyperref[sec:funkcje-testowe18]{rozdziału~18} konstruuje dystrybucję
$\langle T,\varphi\rangle=\sum_{n\ge1}\varphi^{(n)}(n)$
na $\R$, która nie ma skończonego rzędu globalnego.
\end{uwaga}

\begin{przyklad}[Dystrybucje regularne]
\label{ex:regular19}
Jeżeli $f\in L^1_{\mathrm{loc}}(\Omega;\K)$, to
\[\langle T_f,\varphi\rangle=\int_\Omega
f(x)\varphi(x)\,\mathrm d\lambda_d(x)\]
określa dystrybucję, gdyż dla $\varphi\in\mathcal D_K$
\[|\langle T_f,\varphi\rangle|\le
\left(\int_K|f(x)|\,\mathrm d\lambda_d(x)\right)p_0(\varphi).\]
Tak więc dystrybucje regularne mają rząd nie większy niż $0$.
Zwróćmy uwagę, że nie jest potrzebne założenie
$f\in L^1(\Omega)$; np. $f(x)=e^{|x|^2}$ na $\R^d$
definiuje dystrybucję.
\end{przyklad}

\begin{stwierdzenie}[Wierność utożsamienia funkcji z dystrybucją]
\label{prop:regular-injective19}
Jeśli $f,g\in L^1_{\mathrm{loc}}(\Omega;\K)$ i $T_f=T_g$,
to $f=g$ prawie wszędzie w $\Omega$ względem miary Lebesgue'a.
\end{stwierdzenie}
\begin{proof}
Po odjęciu wystarczy rozważyć $T_f=0$.
Wybierzmy nieujemną $\rho\in C_c^\infty(\R^d)$
o nośniku w $\overline{B(0,1)}$, unormowaną warunkiem
$\int_{\R^d}\rho(x)\,\mathrm d\lambda_d(x)=1$, i określmy
$\rho_\varepsilon(x)=\varepsilon^{-d}\rho(x/\varepsilon)$.
Dla $x\in\Omega$ i dostatecznie małego $\varepsilon>0$
funkcja $y\mapsto\rho_\varepsilon(x-y)$ należy do
$\mathcal D(\Omega)$, zatem
$(f*\rho_\varepsilon)(x)=0$.

Na każdym $K\Subset\Omega$ funkcje
$f*\rho_\varepsilon$ zbiegają do $f$ w $L^1(K)$.
Istotnie, wybierzmy otwarty $V$ o zwartym domknięciu,
$K\subseteq V\Subset\Omega$, i przedłużmy
$g=f\mathbf1_V$ przez zero na $\R^d$. Wtedy
$g\in L^1(\R^d)$, a dla dostatecznie małych
$\varepsilon$ sploty $f*\rho_\varepsilon$ i
$g*\rho_\varepsilon$ są równe na $K$. Ponadto
\[\|g*\rho_\varepsilon-g\|_{L^1(\R^d)}\le
\int_{\R^d}\rho(z)\,
\|g(\cdot-\varepsilon z)-g\|_{L^1(\R^d)}\,
\mathrm d\lambda_d(z)\longrightarrow0,\]
co wynika z ciągłości translacji w $L^1$ i twierdzenia
o zbieżności zmajoryzowanej (wyrażenia pod całką są
ograniczone przez $2\|g\|_{L^1}$).
Otrzymujemy $f=0$ prawie wszędzie na każdym $K$,
a więc w $\Omega$.
\end{proof}

\begin{przyklad}[Dystrybucje punktowe]
\label{ex:delta19}
Dla $a\in\Omega$ i wielowskaźnika $\alpha\in\N_0^d$
funkcjonały
\[\langle\delta_a,\varphi\rangle=\varphi(a),\qquad
\langle T_{a,\alpha},\varphi\rangle
=(-1)^{|\alpha|}\partial^\alpha\varphi(a)\]
są dystrybucjami. Mają rząd nie większy niż $|\alpha|$;
dla $\alpha\ne0$ dystrybucja $T_{a,\alpha}$ nie jest
postaci $T_f$ dla $f\in L^1_{\mathrm{loc}}$. Aby uzasadnić nieregularność dla $\alpha\ne0$,
wybierzmy $\psi\in C_c^\infty(B(0,1))$ z
$\partial^\alpha\psi(0)\ne0$ i połóżmy
$\varphi_\varepsilon(x)=\varepsilon^{|\alpha|}
\psi((x-a)/\varepsilon)$. Wtedy
$T_{a,\alpha}(\varphi_\varepsilon)$ nie zależy od
$\varepsilon$, natomiast dla każdego
$f\in L^1_{\mathrm{loc}}(\Omega)$ mamy
$|T_f(\varphi_\varepsilon)|\le
\varepsilon^{|\alpha|}\|\psi\|_\infty
\int_{B(a,\varepsilon)}|f(x)|\,\mathrm d\lambda_d(x)
\to0$. Podobny argument bez czynnika
$\varepsilon^{|\alpha|}$ dowodzi, że sama $\delta_a$
nie jest regularna. Znak wybrano zgodnie z definicją
pochodnej dystrybucyjnej w kolejnym rozdziale:
$T_{a,\alpha}=\partial^\alpha\delta_a$.
\end{przyklad}

\begin{stwierdzenie}[Dokładny rząd pochodnej delty]
\label{prop:delta-order19}
Jeśli $K\Subset\Omega$ zawiera otoczenie punktu $a$,
to $T_{a,\alpha}$ ma na $K$ dokładnie rząd $|\alpha|$.
\end{stwierdzenie}
\begin{proof}
Oszacowanie rzędu $|\alpha|$ wynika z definicji.
Niech $m<|\alpha|$ i wybierzmy
$\psi\in C_c^\infty(B(0,1))$ z
$\partial^\alpha\psi(0)\ne0$.
Dla małego $\varepsilon>0$ funkcja
$\varphi_\varepsilon(x)=\varepsilon^{|\alpha|}
\psi((x-a)/\varepsilon)$ ma nośnik w $K$. Ponadto
$|T_{a,\alpha}(\varphi_\varepsilon)|
=|\partial^\alpha\psi(0)|>0$, podczas gdy
$p_m(\varphi_\varepsilon)\le C_m\varepsilon^{|\alpha|-m}
\to0$. Nie może więc zachodzić oszacowanie rzędu $m$.
\end{proof}

\begin{przyklad}[Wartość główna Cauchy'ego]
\label{ex:principal-value19}
Na $\R$ definiujemy
\[\left\langle\operatorname{pv}\frac1x,\varphi\right\rangle
=\lim_{\varepsilon\downarrow0}
\int_{\{x\in\R:|x|>\varepsilon\}}
\frac{\varphi(x)}{x}\,\mathrm d\lambda_1(x)
=\int_0^\infty\frac{\varphi(t)-\varphi(-t)}{t}\,
\mathrm d\lambda_1(t).\]
Dla $\supp\varphi\subseteq[-R,R]$ oraz $R\ge1$
twierdzenie o wartości średniej daje
$|\varphi(t)-\varphi(-t)|\le2t\,p_1(\varphi)$ dla
$0<t\le1$, natomiast dla $1\le t\le R$ licznik
ma moduł nie większy niż $2p_0(\varphi)$. Dlatego
\[\left|\left\langle\operatorname{pv}\frac1x,
\varphi\right\rangle\right|\le
2p_1(\varphi)+2(\log R)p_0(\varphi).\]
Jeśli $R<1$, wystarcza pierwsza część oszacowania.
Jest to dystrybucja, chociaż $1/x\notin L^1_{\mathrm{loc}}(\R)$.
Symbol $\operatorname{pv}$ oznacza \emph{wartość główną},
a nie zwykłą, bezwzględnie zbieżną całkę.
\end{przyklad}

\subsection{Nośnik i lokalność}

\begin{definicja}[Zerowanie się dystrybucji na zbiorze otwartym]
\label{def:distribution-zero19}
Niech $T\in\mathcal D'(\Omega)$, a $U\subseteq\Omega$
będzie otwarty. Piszemy $T|_U=0$, jeśli
$\langle T,\varphi\rangle=0$ dla każdej
$\varphi\in C_c^\infty(U)$, przedłużonej zerem na $\Omega$.
Nośnikiem $T$ jest domknięty względnie zbiór
\[\supp T=\Omega\setminus\bigcup\{U\subseteq\Omega:
U\text{ jest otwarty i }T|_U=0\}.\]
\end{definicja}

\begin{lemat}[Lokalność dystrybucji]
\label{lem:distribution-local19}
Niech $T\in\mathcal D'(\Omega)$ i
$\varphi\in\mathcal D(\Omega)$ spełnia
$\supp\varphi\cap\supp T=\varnothing$. Wtedy
$\langle T,\varphi\rangle=0$.
\end{lemat}
\begin{proof}
Każdy punkt zwartego $\supp\varphi$ ma otoczenie otwarte
$U_x\subseteq\Omega$, na którym $T$ się zeruje.
Wybieramy skończone podpokrycie
$U_1,\ldots,U_N$, gładką funkcję odcięcia $\eta$ równą
$1$ na otoczeniu $\supp\varphi$ oraz gładki podział
jedności $(\theta_j)_{j=1}^N$ podporządkowany tym
otwartym zbiorom, taki że
$\sum_j\theta_j=1$ na $\supp\varphi$.
Wówczas $\varphi=\sum_j\theta_j\varphi$,
$\supp(\theta_j\varphi)\Subset U_j$, a więc
$\langle T,\theta_j\varphi\rangle=0$ dla każdego $j$.
Sumując, otrzymujemy tezę.
\end{proof}

\begin{stwierdzenie}[Nośniki przykładów]
\label{prop:support-examples19}
Dla $a\in\Omega$ i $\alpha\in\N_0^d$ mamy
$\supp\delta_a=\supp T_{a,\alpha}=\{a\}$.
Dla $f\in L^1_{\mathrm{loc}}(\Omega)$ nośnik $T_f$
jest dopełnieniem największego otwartego podzbioru
$U\subseteq\Omega$, na którym $f=0$ prawie wszędzie.
\end{stwierdzenie}
\begin{proof}
Gdy $U$ nie zawiera $a$, funkcjonały punktowe zerują się
na $\mathcal D(U)$. Każde otoczenie $a$ zawiera funkcję
testową z niezerową pochodną $\partial^\alpha$ w $a$.
Dla $T_f$ zastosujmy Stwierdzenie~\ref{prop:regular-injective19}
na dowolnym otwartym $U\subseteq\Omega$.
\end{proof}

\begin{twierdzenie}[Zwarty nośnik implikuje skończony rząd]
\label{thm:compact-support-order19}
Jeżeli $T\in\mathcal D'(\Omega;\K)$ i
$S=\supp T\Subset\Omega$, to istnieją
$m\in\N_0$, $C\ge0$ i $K\Subset\Omega$ takie, że
\[|\langle T,\varphi\rangle|\le C p_{m,K}(\varphi)
\quad(\varphi\in\mathcal D(\Omega)),\qquad
p_{m,K}(\varphi)=\max_{|\alpha|\le m}
\sup_{x\in K}|\partial^\alpha\varphi(x)|.\]
W szczególności $T$ ma skończony rząd na całej $\Omega$.
Ponadto $T$ przedłuża się w sposób jednoznaczny do
ciągłego funkcjonału liniowego na $C^\infty(\Omega;\K)$.
\end{twierdzenie}
\begin{proof}
Wybierzmy $\chi\in C_c^\infty(\Omega)$, równą $1$
na pewnym otoczeniu $S$. Wtedy
$\supp((1-\chi)\varphi)\cap S=\varnothing$, więc z
Lematu~\ref{lem:distribution-local19} wynika
$\langle T,\varphi\rangle=\langle T,\chi\varphi\rangle$.
Niech $K=\supp\chi$. Z Twierdzenia~\ref{thm:distribution-estimate19}
istnieją $m,C_0$ takie, że
$|\langle T,\psi\rangle|\le C_0p_m(\psi)$ dla
$\psi\in\mathcal D_K$. Reguła Leibniza daje
$p_m(\chi\varphi)\le C_1p_{m,K}(\varphi)$,
co dowodzi oszacowania globalnego.
Dla $u\in C^\infty(\Omega)$ definiujemy
$\widetilde T(u)=\langle T,\chi u\rangle$.
Podane oszacowanie zapewnia ciągłość w topologii
$C^\infty(\Omega)$, a wzór z poprzedniego akapitu
pokazuje zgodność z $T$ na funkcjach testowych.
Przestrzeń $\mathcal D(\Omega)$ jest gęsta w
$C^\infty(\Omega)$ w tej topologii (\hyperref[sec:gladkie-schwartz17]{rozdział~17}),
więc przedłużenie jest jednoznaczne.
\end{proof}

\subsection{Słaba-* i silna topologia na przestrzeni dystrybucji}

\begin{definicja}[Topologia słaba-*]
\label{def:weakstar19}
Na $\mathcal D'(\Omega)$ definiujemy
$\sigma(\mathcal D'(\Omega),\mathcal D(\Omega))$ jako
najsłabszą lokalnie wypukłą topologię, dla której
wszystkie funkcjonały $T\mapsto\langle T,\varphi\rangle$,
$\varphi\in\mathcal D(\Omega)$, są ciągłe.
Jest ona wyznaczona przez seminormy
$q_\varphi(T)=|\langle T,\varphi\rangle|$.
Sieć $(T_i)$ zbiega w tej topologii do $T$ dokładnie wtedy,
gdy $\langle T_i,\varphi\rangle\to\langle T,\varphi\rangle$
dla każdej $\varphi\in\mathcal D(\Omega)$.
\end{definicja}

\begin{definicja}[Topologia silna]
\label{def:strong19}
Na $\mathcal D'(\Omega)$ definiujemy
\emph{silną topologię przestrzeni dualnej}
$\beta(\mathcal D'(\Omega),\mathcal D(\Omega))$
jako topologię jednostajnej zbieżności na zbiorach
ograniczonych $B\subseteq\mathcal D(\Omega)$.
Wyznaczają ją seminormy
\[q_B(T)=\sup_{\varphi\in B}|\langle T,\varphi\rangle|.\]
Wartości $q_B$ są skończone, ponieważ ciągły funkcjonał
liniowy przenosi zbiory ograniczone na ograniczone.
Dla $B=\{\varphi\}$ otrzymujemy $q_B=q_\varphi$,
zatem topologia silna jest nie słabsza od słabej-*.
\end{definicja}

\begin{twierdzenie}[Beczkowość przestrzeni testowej]
\label{thm:D-barrelled19}
Przestrzeń $\mathcal D(\Omega;\K)$ jest beczkowa.
Każda rodzina $\mathcal A\subseteq\mathcal D'(\Omega;\K)$
ograniczona punktowo, tj. spełniająca
\[\sup_{T\in\mathcal A}|\langle T,\varphi\rangle|<\infty
\quad\text{dla każdego }\varphi\in\mathcal D(\Omega),\]
jest równociągła.
\end{twierdzenie}
\begin{proof}
Niech $U\subseteq\mathcal D(\Omega)$ będzie
beczką, tj. zbiorem domkniętym, absolutnie wypukłym
i pochłaniającym. Każdy $U\cap\mathcal D_{K_n}$ jest
beczką w przestrzeni Fr\'echeta $\mathcal D_{K_n}$.
Ta ostatnia jest beczkowa na mocy twierdzenia Baire'a
i twierdzenia Banacha--Steinhausa w wersji geometrycznej
(\hyperref[sec:beczkowe-bornologiczne]{rozdział~10}). Zatem
$U\cap\mathcal D_{K_n}$ jest otoczeniem zera
w $\mathcal D_{K_n}$.
Funkcjonał Minkowskiego $p_U$ ma więc ciągłe ograniczenie
na każdy etap $\mathcal D_{K_n}$. Własność uniwersalna
granicy induktywnej pokazuje, że $p_U$ jest ciągłą
seminormą na $\mathcal D(\Omega)$; zbiór
$\{\varphi:p_U(\varphi)<1\}\subseteq U$ jest otoczeniem
zera. Przestrzeń $\mathcal D(\Omega)$ jest więc beczkowa.
Druga część wynika z ogólnego twierdzenia
Banacha--Steinhausa dla przestrzeni beczkowych
(\hyperref[sec:beczkowe-bornologiczne]{rozdział~10}).
\end{proof}

\begin{wniosek}[Dwie postaci ograniczoności dystrybucji]
\label{cor:bounded-distributions19}
Dla rodziny $\mathcal A\subseteq\mathcal D'(\Omega)$
następujące warunki są równoważne:
\begin{enumerate}
\item $\mathcal A$ jest punktowo ograniczona na $\mathcal D(\Omega)$;
\item $\mathcal A$ jest równociągła na $\mathcal D(\Omega)$;
\item $\mathcal A$ jest ograniczona w silnej topologii
$\beta(\mathcal D',\mathcal D)$.
\end{enumerate}
\end{wniosek}
\begin{proof}
Pierwszy warunek implikuje drugi przez
Twierdzenie~\ref{thm:D-barrelled19}. Jeżeli
$\mathcal A$ jest równociągła, to dla każdego
$\varepsilon>0$ istnieje otoczenie $V$ zera takie,
że $|\langle T,\varphi\rangle|\le\varepsilon$ dla
$T\in\mathcal A$ i $\varphi\in V$.
Każdy zbiór ograniczony $B$ leży w $\lambda V$
dla pewnego $\lambda>0$, więc
$\sup_{T\in\mathcal A}q_B(T)\le\lambda\varepsilon$.
To daje ograniczoność silną. Wreszcie singletony
$\{\varphi\}$ są ograniczone, więc trzeci warunek
implikuje pierwszy.
\end{proof}

\begin{wniosek}[Jednostajne lokalne oszacowania rodziny dystrybucji]
\label{cor:uniform-order19}
Jeżeli rodzina $\mathcal A\subseteq\mathcal D'(\Omega)$ jest
punktowo ograniczona, to dla każdego $K\Subset\Omega$
istnieją $m_K\in\N_0$ i $C_K\ge0$ takie, że
$|\langle T,\varphi\rangle|\le C_Kp_{m_K}(\varphi)$
jednocześnie dla wszystkich $T\in\mathcal A$
i $\varphi\in\mathcal D_K$.
\end{wniosek}
\begin{proof}
Z Twierdzenia~\ref{thm:D-barrelled19} rodzina $\mathcal A$
jest równociągła na $\mathcal D(\Omega)$, a więc jej
ograniczenia do $\mathcal D_K$ są równociągłe.
Istnieje zatem otoczenie zera $U$ w $\mathcal D_K$,
na którym $|\langle T,\varphi\rangle|\le1$ dla wszystkich
$T\in\mathcal A$. Ponieważ seminormy $p_m$ są rosnące,
istnieją $m_K$ i $\varepsilon_K>0$ takie, że
$\{\varphi:p_{m_K}(\varphi)<\varepsilon_K\}\subseteq U$.
Skalowanie $\varphi$ daje tezę z $C_K=2/\varepsilon_K$.
\end{proof}

\begin{twierdzenie}[Zbieżność ciągów: słaba-* = silna]
\label{thm:weak-strong-sequences19}
Niech $(T_n)_{n\ge1}\subseteq\mathcal D'(\Omega;\K)$
i $T\in\mathcal D'(\Omega;\K)$. Następujące warunki
są równoważne:
\begin{enumerate}
\item $T_n\to T$ w $\sigma(\mathcal D',\mathcal D)$;
\item $T_n\to T$ w $\beta(\mathcal D',\mathcal D)$.
\end{enumerate}
\end{twierdzenie}
\begin{proof}
Implikacja (2)$\Rightarrow$(1) wynika z definicji.
Zakładamy (1). Rodzina $\{T_n-T:n\ge1\}$ jest
punktowo ograniczona, gdyż dla ustalonej funkcji
testowej ciąg liczb $\langle T_n-T,\varphi\rangle$
jest zbieżny. Z Wniosku~\ref{cor:bounded-distributions19}
rodzina ta jest równociągła. Ustalmy zbiór ograniczony
$B\subseteq\mathcal D(\Omega)$ i $\varepsilon>0$.
Jego domknięcie $K=\overline B^{\,\mathcal D}$ jest zwarte
na mocy twierdzenia o względnej zwartości zbiorów ograniczonych z \hyperref[sec:funkcje-testowe18]{rozdziału~18}.
Z równociągłości wybierzmy zrównoważone otoczenie
zera $V$ takie, że
$|\langle T_n-T,v\rangle|<\varepsilon/2$
za każdym razem, gdy $n\ge1$ i $v\in V$.
Zwartość $K$ pozwala wybrać
$\varphi_1,\ldots,\varphi_N\in K$ z
$K\subseteq\bigcup_{j=1}^N(\varphi_j+V)$.
Zbieżność punktowa zapewnia indeks $n_0$ taki, że
$|\langle T_n-T,\varphi_j\rangle|<\varepsilon/2$
dla wszystkich $n\ge n_0$ i $j\le N$.
Dla dowolnego $\varphi\in B$ piszemy
$\varphi=\varphi_j+v$ z $v\in V$, a więc
$|\langle T_n-T,\varphi\rangle|<\varepsilon$.
Otrzymujemy $q_B(T_n-T)\to0$, jak trzeba.
\end{proof}

\begin{uwaga}[Dlaczego nie wystarczy zbieżność sieci?]
\label{rem:net-distinction19}
Powyższy wynik dotyczy ciągów. Dla ogólnej sieci
zbieżność punktowa nie musi zapewniać punktowej
ograniczoności \emph{całej} rodziny jej wyrazów.
Nie wolno więc mechanicznie zastępować w tym dowodzie
ciągu przez sieć. Poniżej wykażemy, że obie topologie
są istotnie różne.
\end{uwaga}

\begin{twierdzenie}[Ścisła różnica topologii]
\label{thm:weak-strong-different19}
Jeżeli $\Omega\subseteq\R^d$ jest niepustym zbiorem
otwartym, to na $\mathcal D'(\Omega)$ topologia
$\beta(\mathcal D',\mathcal D)$ jest ściśle silniejsza
od $\sigma(\mathcal D',\mathcal D)$. Istnieje sieć
dystrybucji zbieżna słabo-* do zera, która nie
zbiega silnie do zera.
\end{twierdzenie}
\begin{proof}
Wybierzmy kulę $B(a,r)\Subset\Omega$ i funkcję
$\chi\in C_c^\infty(\Omega)$ dodatnią na tej kuli
i równą zeru poza jej domknięciem. Taka funkcja istnieje,
na przykład przez użycie standardowej funkcji wypukłej
w wykładniku $\exp(-1/(r^2-|x-a|^2))$ wewnątrz kuli.
Dla $j\ge0$ połóżmy
$\psi_j(x)=2^{-j}((x_1-a_1)/r)^j\chi(x)$.
Z reguły Leibniza, ponieważ $|(x_1-a_1)/r|\le1$
na nośniku $\chi$, dla każdego $m$ otrzymujemy
$\sup_jp_m(\psi_j)<\infty$: każdy składnik jest
oszacowany przez stałą razy $2^{-j}(1+j)^m$.
Wobec tego $B_0=\{\psi_j:j\ge0\}$ jest ograniczony
w $\mathcal D(\Omega)$. Funkcje $\psi_j$ są
liniowo niezależne, ponieważ wielomiany różnych stopni
są liniowo niezależne na otwartym zbiorze, gdzie
$\chi\ne0$.

Niech $I$ będzie zbiorem par $(F,n)$, gdzie $F$ jest
skończonym podzbiorem $\mathcal D(\Omega)$, a $n\in\N$,
z porządkiem $(F,n)\preceq(G,m)$ wtedy i tylko wtedy,
gdy $F\subseteq G$ i $n\le m$. Dla każdej pary wybieramy
$\psi\in B_0\setminus\operatorname{lin}(F)$.
Skończenie wymiarowa podprzestrzeń
$\operatorname{lin}(F)$ jest domknięta, zatem
z twierdzenia Hahna--Banacha o oddzielaniu punktu od
podprzestrzeni istnieje $T_{F,n}\in\mathcal D'(\Omega)$
spełniające $T_{F,n}|_{\operatorname{lin}(F)}=0$ i
$|\langle T_{F,n},\psi\rangle|\ge n$.

Dla ustalonej $\varphi\in\mathcal D(\Omega)$ i
$(F,n)\succeq(\{\varphi\},1)$ mamy
$\langle T_{F,n},\varphi\rangle=0$. Sieć
$(T_{F,n})_{(F,n)\in I}$ zbiega więc słabo-* do zera.
Jednak $q_{B_0}(T_{F,n})\ge n$, toteż nie zbiega
w topologii silnej. Dowodzi to również ścisłości
obu topologii.
\end{proof}

\subsection{Granice regularizacji i zastosowania}

\begin{przyklad}[Delty uciekające do nieskończoności]
\label{ex:delta-escape19}
Niech $\Omega=\R^d$, a $(a_n)\subseteq\R^d$
spełnia $|a_n|\to\infty$. Wtedy $\delta_{a_n}\to0$
\emph{silnie} w $\mathcal D'(\R^d)$.
Istotnie, każdy ograniczony $B\subseteq\mathcal D(\R^d)$
ma wspólny zwarty nośnik $K$, zatem dla dużych $n$
$\sup_{\varphi\in B}|\varphi(a_n)|=0$.
\end{przyklad}

\begin{przyklad}[Ciągłość delty względem położenia]
\label{ex:delta-continuity19}
Jeżeli $a_n\to a$ w $\Omega$, to
$\delta_{a_n}\to\delta_a$ silnie w $\mathcal D'(\Omega)$.
Dla ograniczonego $B\subseteq\mathcal D(\Omega)$
i $n$ dostatecznie dużych odcinek $[a,a_n]$ leży
w $\Omega$. Z twierdzenia o wartości średniej
\[q_B(\delta_{a_n}-\delta_a)\le
|a_n-a|\sup_{\varphi\in B}\sup_{x\in\Omega}|\nabla\varphi(x)|,
\]
a prawa strona dąży do zera.
\end{przyklad}

\begin{twierdzenie}[Aproksymacja jednostki i zbieżność do delty]
\label{thm:mollifier-delta19}
Niech $a\in\Omega$ i niech
$\rho\in C_c^\infty(\R^d)$ ma nośnik w
$\overline{B(0,1)}$, spełnia $\rho\ge0$ i
$\int_{\R^d}\rho(x)\,\mathrm d\lambda_d(x)=1$.
Dla dostatecznie dużych $n$ określmy
$f_n(x)=n^d\rho(n(x-a))$ dla $x\in\Omega$.
Wtedy $T_{f_n}\to\delta_a$ silnie w $\mathcal D'(\Omega)$.
Co więcej, dla każdego zbioru ograniczonego
$B\subseteq\mathcal D(\Omega)$ istnieje stała $C_\rho$
niezależna od $B$ i $n$ taka, że
\[q_B(T_{f_n}-\delta_a)\le
\frac{C_\rho}{n}\sup_{\varphi\in B}p_1(\varphi).\]
\end{twierdzenie}
\begin{proof}
Po podstawieniu $y=n(x-a)$ otrzymujemy
\[\langle T_{f_n},\varphi\rangle
=\int_{\R^d}\rho(y)\varphi(a+y/n)\,\mathrm d\lambda_d(y).\]
Dla $y\in\supp\rho$ oraz dostatecznie dużego $n$
odcinek od $a$ do $a+y/n$ leży w $\Omega$. Zatem
$|\varphi(a+y/n)-\varphi(a)|\le
|y|\,p_1(\varphi)/n$ (po ewentualnym zastąpieniu
$p_1$ równoważną seminormą kontrolującą normę gradientu;
w naszym zapisie można wziąć dodatkowy czynnik $\sqrt d$).
Stąd teza z
$C_\rho=\sqrt d\int_{\R^d}|y|\rho(y)\,
\mathrm d\lambda_d(y)$. Ograniczoność $B$ zapewnia
$\sup_{\varphi\in B}p_1(\varphi)<\infty$, więc prawa
strona dąży do zera.
\end{proof}



\begin{twierdzenie}[Zbieżność w $L^1_{\mathrm{loc}}$ a zbieżność dystrybucji]
\label{thm:L1-to-strong19}
Niech $f_n,f\in L^1_{\mathrm{loc}}(\Omega;\K)$ i załóżmy,
że $f_n\to f$ w $L^1(K)$ dla każdego $K\Subset\Omega$.
Wówczas $T_{f_n}\to T_f$ w silnej topologii
$\beta(\mathcal D'(\Omega),\mathcal D(\Omega))$.
\end{twierdzenie}
\begin{proof}
Niech $B\subseteq\mathcal D(\Omega)$ będzie ograniczony.
Z kryterium ograniczoności z \hyperref[sec:funkcje-testowe18]{rozdziału~18} istnieje
$K\Subset\Omega$ zawierający wszystkie nośniki funkcji
z $B$, przy czym $M_B=\sup_{\varphi\in B}p_0(\varphi)<\infty$.
Otrzymujemy
\[q_B(T_{f_n}-T_f)\le
M_B\int_K|f_n(x)-f(x)|\,\mathrm d\lambda_d(x)
\longrightarrow0.\]
Jest to dokładnie definicja zbieżności silnej.
\end{proof}

\begin{przyklad}[Silna zbieżność bez zbieżności lokalnej w $L^1$]
\label{ex:oscillations19}
Na $\Omega=\R$ weźmy $f_n(x)=\sin(nx)$.
Dla każdej $\varphi\in\mathcal D(\R)$ całkowanie przez części daje
\[\langle T_{f_n},\varphi\rangle=
\frac1n\int_\R\cos(nx)\varphi'(x)\,
\mathrm d\lambda_1(x).\]
Niech $B\subseteq\mathcal D(\R)$ będzie ograniczony.
Wspólny nośnik funkcji z $B$ mieści się w $[-R,R]$
dla pewnego $R>0$, a $\sup_{\varphi\in B}p_1(\varphi)<\infty$.
Dlatego
\[q_B(T_{f_n})\le\frac{2R}{n}
\sup_{\varphi\in B}p_1(\varphi)\longrightarrow0.\]
Zatem $T_{f_n}\to0$ \emph{silnie} w $\mathcal D'(\R)$,
ale $\int_{-1}^1|\sin(nx)|\,\mathrm d\lambda_1(x)$
nie dąży do zera (w istocie dąży do $4/\pi$).
Zbieżność dystrybucji jest więc znacznie słabsza niż
zbieżność reprezentujących je funkcji w $L^1_{\mathrm{loc}}$.
\end{przyklad}


% Wersja źródłowa: audyt II, jednostka 20
\setcounter{section}{19}
\section{Pochodne dystrybucyjne i transpozycja operatorów}
\label{sec:pochodne20}

\subsection{Definicja pochodnej i ciągłość w topologiach dualnych}

W tym rozdziale $d\ge1$, $\Omega\subseteq\R^d$ jest niepustym
zbiorem otwartym, a $\K\in\{\R,\C\}$. Dystrybucje są
$\K$-liniowymi funkcjonałami na $\mathcal D(\Omega;\K)$, zgodnie
z konwencją przyjętą w \hyperref[sec:dystrybucje19]{rozdziale~19}.
Dla wielowskaźnika $\alpha=(\alpha_1,\ldots,\alpha_d)\in\N_0^d$
piszemy $|\alpha|=\sum_j\alpha_j$,
$\partial^\alpha=\partial_1^{\alpha_1}\cdots\partial_d^{\alpha_d}$
oraz $\binom\alpha\beta=\prod_j\binom{\alpha_j}{\beta_j}$
dla $\beta\le\alpha$ współrzędnie.

\begin{lemat}[Różniczkowanie funkcji testowych]
\label{lem:test-derivative20}
Dla każdego $\alpha\in\N_0^d$ odwzorowanie
$D_\alpha:\mathcal D(\Omega;\K)\to\mathcal D(\Omega;\K)$,
$D_\alpha\varphi=\partial^\alpha\varphi$, jest ciągłe i liniowe.
Ponadto dla każdego zbioru ograniczonego $B\subseteq\mathcal D(\Omega)$
zbiór $D_\alpha(B)$ jest ograniczony.
\end{lemat}
\begin{proof}
Dla każdego $K\Subset\Omega$ odwzorowanie $D_\alpha$ zachowuje
$\mathcal D_K(\Omega)$, ponieważ różniczkowanie nie powiększa nośnika.
Przy seminormach z \hyperref[sec:funkcje-testowe18]{rozdziału~18} mamy
$p_m(D_\alpha\varphi)\le p_{m+|\alpha|}(\varphi)$.
Ograniczenie $D_\alpha|_{\mathcal D_K}$ jest więc ciągłe.
Własność uniwersalna standardowej topologii granicy induktywnej
zapewnia ciągłość na $\mathcal D(\Omega)$. Obraz zbioru ograniczonego
przez ciągłe odwzorowanie liniowe między przestrzeniami
liniowo-topologicznymi jest ograniczony.
\end{proof}

\begin{definicja}[Pochodna dystrybucyjna]
\label{def:dist-deriv20}
Niech $T\in\mathcal D'(\Omega;\K)$ i $\alpha\in\N_0^d$.
\emph{Pochodną dystrybucyjną} $\partial^\alpha T$ określamy wzorem
\[
 \langle\partial^\alpha T,\varphi\rangle
 =(-1)^{|\alpha|}\langle T,\partial^\alpha\varphi\rangle,
 \qquad\varphi\in\mathcal D(\Omega;\K).
\]
Dla $\alpha=e_j$ piszemy również $\partial_jT$.
\end{definicja}

\begin{twierdzenie}[Ciągłość i algebra pochodnych]
\label{thm:dist-deriv-cont20}
Niech $\alpha,\beta\in\N_0^d$. Wówczas:
\begin{enumerate}
\item $\partial^\alpha T$ jest dystrybucją dla każdego
$T\in\mathcal D'(\Omega)$;
\item $\partial^\alpha\partial^\beta T
=\partial^{\alpha+\beta}T$, w szczególności
$\partial_j\partial_kT=\partial_k\partial_jT$;
\item operator $\partial^\alpha:\mathcal D'(\Omega)\to
\mathcal D'(\Omega)$ jest ciągły zarówno w topologii słabej-*
$\sigma(\mathcal D',\mathcal D)$, jak i silnej
$\beta(\mathcal D',\mathcal D)$;
\item $\supp(\partial^\alpha T)\subseteq\supp T$.
Jeśli $T$ ma na $K\Subset\Omega$ rząd nie większy niż $m$,
to $\partial^\alpha T$ ma na $K$ rząd nie większy niż
$m+|\alpha|$.
\end{enumerate}
\end{twierdzenie}
\begin{proof}
Punkt pierwszy wynika z ciągłości $D_\alpha$ na
$\mathcal D(\Omega)$ i złożenia funkcjonałów.
Dla drugiego punktu stosujemy dwukrotnie definicję i równość
$\partial^\alpha\partial^\beta\varphi
=\partial^{\alpha+\beta}\varphi$ dla funkcji gładkich.
Dla ustalonej $\varphi$ zachodzi
$q_\varphi(\partial^\alpha T)=q_{\partial^\alpha\varphi}(T)$,
więc otrzymujemy ciągłość słabą-*.
Dla zbioru ograniczonego $B\subseteq\mathcal D(\Omega)$
mamy natomiast
\[
 q_B(\partial^\alpha T)=q_{D_\alpha(B)}(T).
\]
Z lematu~\ref{lem:test-derivative20} zbiór $D_\alpha(B)$ jest
ograniczony, co daje ciągłość silną. Gdy $T$ zeruje się na
otwartym $U\subseteq\Omega$, zeruje się tam także $\partial^\alpha T$:
dla $\varphi\in\mathcal D(U)$ pochodna $\partial^\alpha\varphi$
nadal należy do $\mathcal D(U)$. Stąd inkluzja nośników.
Wreszcie dla $\varphi\in\mathcal D_K$ mamy
$|\langle\partial^\alpha T,\varphi\rangle|
\le C_Kp_m(\partial^\alpha\varphi)
\le C_Kp_{m+|\alpha|}(\varphi)$.
\end{proof}

\begin{stwierdzenie}[Zgodność z pochodną klasyczną]
\label{prop:classical-weak20}
Niech $f\in C^m(\Omega;\K)$, $m\ge1$, a $\alpha\in\N_0^d$
niech spełnia $|\alpha|\le m$. Wtedy
$\partial^\alpha T_f=T_{\partial^\alpha f}$.
Dla $f,g\in L^1_{\mathrm{loc}}(\Omega)$ tożsamość
$\partial_jT_f=T_g$ jest równoważna temu, że $g$ jest
\emph{słabą pochodną} $f$ w kierunku $x_j$, tj.
\[
 \int_\Omega f(x)\,\partial_j\varphi(x)\,\mathrm d\lambda_d(x)
 =-\int_\Omega g(x)\varphi(x)\,\mathrm d\lambda_d(x)
 \quad(\varphi\in\mathcal D(\Omega)).
\]
\end{stwierdzenie}
\begin{proof}
W pierwszej części całkujemy przez części $|\alpha|$ razy.
Znikanie składników brzegowych wynika ze zwartego nośnika
funkcji testowej, a nie z warunków brzegowych na $f$.
W drugiej części równość jest dokładnie definicją pochodnej
dystrybucyjnej zastosowaną do regularnych dystrybucji
$T_f$ i $T_g$.
\end{proof}

\subsection{Skoki, funkcja Heaviside'a i osobliwości punktowe}

\begin{przyklad}[Heaviside i delta Diraca]
\label{ex:heaviside20}
Na $\R$ niech $H(x)=0$ dla $x<0$ oraz $H(x)=1$ dla $x>0$;
wartość w $0$ nie ma znaczenia dla wyznaczonej dystrybucji.
Dla dowolnej $\varphi\in\mathcal D(\R)$ otrzymujemy
\[
 \langle (T_H)',\varphi\rangle
 =-\int_0^\infty\varphi'(x)\,\mathrm d\lambda_1(x)
 =\varphi(0)=\langle\delta_0,\varphi\rangle.
\]
Zatem $(T_H)'=\delta_0$; w dalszym tekście piszemy krótko
$H'=\delta_0$. Również $(\operatorname{sgn}x)'=2\delta_0$,
$|x|'=\operatorname{sgn}x$ i $|x|''=2\delta_0$.
\end{przyklad}

\begin{center}
\begin{tikzpicture}[scale=0.82,>=Stealth]
\draw[->] (-3,0) -- (3.3,0) node[right] {$x$};
\draw[->] (0,-0.5) -- (0,1.9) node[above] {$H(x)$};
\draw[very thick,blue!60!black] (-2.9,0)--(-0.06,0);
\draw[very thick,blue!60!black] (0.06,1.3)--(3.0,1.3);
\draw[fill=white,draw=blue!60!black,thick] (0,0) circle (2.3pt);
\draw[fill=white,draw=blue!60!black,thick] (0,1.3) circle (2.3pt);
\node[left] at (-0.12,1.3) {$1$};
\node[below left] at (-0.08,-0.05) {$0$};
\draw[gray,dashed] (0,0.12)--(0,1.2);
\end{tikzpicture}

\small\emph{Skok funkcji $H$. Wartość w zerze jest dowolna;
pochodna $H'=\delta_0$ jest tożsamością dystrybucyjną.}
\end{center}

\begin{twierdzenie}[Pochodna funkcji ze skokiem]
\label{thm:jump20}
Niech $a\in\R$, a $f\in L^1_{\mathrm{loc}}(\R;\K)$ będzie
klasy $C^1$ na $(-\infty,a)$ i $(a,\infty)$. Załóżmy ponadto,
że funkcja ma skończone granice jednostronne
$f(a-),f(a+)$, a klasyczna pochodna określona poza $a$
należy do $L^1_{\mathrm{loc}}(\R;\K)$.
Niech $g(x)=f'(x)$ dla $x\ne a$ (wartość w $a$ jest dowolna).
Wówczas
\[
 (T_f)'=T_g+\bigl(f(a+)-f(a-)\bigr)\delta_a.
\]
\end{twierdzenie}
\begin{proof}
Dla $\varphi\in\mathcal D(\R)$ całkujemy przez części
oddzielnie na $(-\infty,a-\varepsilon)$ i
$(a+\varepsilon,\infty)$, a następnie przechodzimy do
granicy $\varepsilon\downarrow0$. Ponieważ $f$ ma skończone
granice jednostronne, jest ograniczona w pobliżu $a$,
więc wkład z $(a-\varepsilon,a+\varepsilon)$ do
$\int f\varphi'$ dąży do zera. Granice członów brzegowych
tworzą $f(a-)\varphi(a)-f(a+)\varphi(a)$.
Z założenia $g\in L^1_{\mathrm{loc}}$, zatem całki
$\int g\varphi$ po obu stronach są zbieżne i otrzymujemy
\[
 -\int_\R f(x)\varphi'(x)\,\mathrm d\lambda_1(x)
 =\int_\R g(x)\varphi(x)\,\mathrm d\lambda_1(x)
 +(f(a+)-f(a-))\varphi(a).
\]
Jest to żądana tożsamość dystrybucyjna.
\end{proof}

\begin{przyklad}[Funkcja częściowo gładka]
\label{ex:piecewise20}
Niech $f(x)=e^x$ dla $x<0$ i $f(x)=2e^{-x}$ dla $x>0$.
Wtedy $f(0+)-f(0-)=1$, a
\[
 (T_f)'=T_{e^x\mathbf1_{(-\infty,0)}
 -2e^{-x}\mathbf1_{(0,\infty)}}+\delta_0.
\]
Jest to równość dystrybucji na $\R$; jej prawa strona
zawiera składnik osobliwy, którego nie wykrywa klasyczne
różniczkowanie poza punktem skoku.
\end{przyklad}

\begin{przyklad}[Pochodne delty]
\label{ex:delta-derivative20}
Dla $a\in\Omega$ i $\alpha\in\N_0^d$ zachodzi
$\langle\partial^\alpha\delta_a,\varphi\rangle
=(-1)^{|\alpha|}\partial^\alpha\varphi(a)$.
Dla $\Omega=\R$ funkcjonał $\delta_0'$ jest więc różny
od zwykłej miary Diraca i ma rząd dokładnie $1$ na każdym
zwartym przedziale zawierającym otoczenie zera.
\end{przyklad}

\subsection{Mnożenie przez funkcję gładką i reguła Leibniza}

\begin{definicja}[Mnożenie dystrybucji]
\label{def:multiply20}
Dla $g\in C^\infty(\Omega;\K)$ oraz $T\in\mathcal D'(\Omega;\K)$
definiujemy $gT\in\mathcal D'(\Omega;\K)$ wzorem
$\langle gT,\varphi\rangle=\langle T,g\varphi\rangle$.
\end{definicja}

\begin{twierdzenie}[Własności mnożenia]
\label{thm:multiply20}
Niech $g,h\in C^\infty(\Omega;\K)$, $T\in\mathcal D'(\Omega;\K)$
i $\alpha\in\N_0^d$. Wtedy:
\begin{enumerate}
\item $gT$ jest dobrze określoną dystrybucją,
$g(hT)=(gh)T$ oraz $1T=T$;
\item $\supp(gT)\subseteq\supp g\cap\supp T$;
\item operator $M_g:T\mapsto gT$ jest ciągły w topologiach
słabej-* i silnej;
\item zachodzi \emph{wzór Leibniza}
\[
 \partial^\alpha(gT)=
 \sum_{\beta\le\alpha}\binom\alpha\beta
 (\partial^\beta g)\,\partial^{\alpha-\beta}T.
\]
\end{enumerate}
\end{twierdzenie}
\begin{proof}
Mnożenie przez $g$ przekształca $\mathcal D_K$ w siebie.
Na każdym takim etapie z reguły Leibniza wynika oszacowanie
\[
 p_m(g\varphi)\le C_{K,m,g}\,p_m(\varphi),
 \quad C_{K,m,g}=\max_{|\gamma|\le m}
 \sum_{\beta\le\gamma}\binom\gamma\beta
 \sup_{x\in K}|\partial^\beta g(x)|.
\]
Własność uniwersalna $\mathcal D(\Omega)$ daje ciągłość
$M_g:\varphi\mapsto g\varphi$. Wobec tego $gT=T\circ M_g$
jest dystrybucją. Tożsamości algebraiczne wynikają ze zwykłej
łączności mnożenia funkcji.
Jeśli $g$ zeruje się na otwartym $U$, to dla funkcji testowych
w $U$ mamy $g\varphi=0$. Jeśli $T|_U=0$, to również
$\langle T,g\varphi\rangle=0$. Stąd inkluzja nośników.
Ciągłość słaba-* wynika z tożsamości
$q_\varphi(gT)=q_{g\varphi}(T)$, a silna z
$q_B(gT)=q_{M_g(B)}(T)$ i ograniczoności $M_g(B)$.
Dla $\alpha=e_j$ obliczamy
\[
 \begin{aligned}
 \langle\partial_j(gT),\varphi\rangle
 &=-\langle T,g\partial_j\varphi\rangle\\
 &=\langle\partial_jT,g\varphi\rangle
   +\langle T,(\partial_jg)\varphi\rangle,
 \end{aligned}
\]
czyli $\partial_j(gT)=g\partial_jT+(\partial_jg)T$.
Indukcja względem $|\alpha|$ i tożsamość dwumianowa
dla współczynników wielowskaźnikowych kończą dowód.
\end{proof}

\begin{przyklad}[Mnożenie dystrybucji punktowych]
\label{ex:multiply-delta20}
Dla $g\in C^\infty(\R)$ mamy
\[
 g\delta_a=g(a)\delta_a,\qquad
 g\delta_a'=g(a)\delta_a'-g'(a)\delta_a.
\]
W szczególności $x\delta_0=0$ oraz $x\delta_0'=-\delta_0$.
Równości sprawdza się przez zastosowanie obu stron
do dowolnej funkcji testowej i zróżniczkowanie iloczynu
$g\varphi$ w punkcie $a$.
\end{przyklad}

\begin{uwaga}[Dlaczego nie mnożymy dowolnych dystrybucji?]
\label{rem:no-product20}
Powyższa definicja wymaga \emph{gładkiego} mnożnika.
Nie ustanawia iloczynu dwóch dowolnych dystrybucji.
Przykładowo wyrażenie $H\delta_0$ nie ma kanonicznego
znaczenia na podstawie samej struktury $\mathcal D'(\R)$:
po zastąpieniu $\delta_0$ przez
$\rho_\varepsilon(x)=\varepsilon^{-1}\rho(x/\varepsilon)$,
gdzie $\rho\in C_c^\infty(\R)$ oraz
$\int_\R\rho(x)\,\mathrm d\lambda_1(x)=1$,
otrzymujemy w sensie dystrybucji
\[
 H\rho_\varepsilon\longrightarrow c_\rho\delta_0,
 \qquad c_\rho=\int_0^\infty\rho(t)\,\mathrm d\lambda_1(t).
\]
Wartość $c_\rho$ zależy od rozmieszczenia masy
aproksymacji jedności po obu stronach zera.
Tym samym różne naturalne regularizacje mogą dawać różne wyniki.
\end{uwaga}

\subsection{Operatory różniczkowe i ich transpozycja}

\begin{definicja}[Formalny operator transponowany]
\label{def:formal-transpose20}
Niech $m\in\N_0$ i $a_\alpha\in C^\infty(\Omega;\K)$
dla $|\alpha|\le m$. Dla funkcji gładkich i dystrybucji
rozważamy operator $P=\sum_{|\alpha|\le m}a_\alpha\partial^\alpha$.
Jego \emph{formalny operator transponowany} na $\mathcal D(\Omega)$
definiujemy wzorem
\[
 P^{\mathrm t}\varphi=
 \sum_{|\alpha|\le m}(-1)^{|\alpha|}
 \partial^\alpha(a_\alpha\varphi).
\]
Nie jest to operator sprzężony w przestrzeni Hilberta:
przy naszej konwencji parowania nie występuje sprzężenie
współczynników zespolonych.
\end{definicja}

\begin{twierdzenie}[Transpozycja jako konstrukcja operatora na dystrybucjach]
\label{thm:transpose20}
Dla operatora $P$ z definicji~\ref{def:formal-transpose20}
oraz każdej dystrybucji $T\in\mathcal D'(\Omega;\K)$
zachodzą:
\begin{enumerate}
\item $P^{\mathrm t}:\mathcal D(\Omega)\to\mathcal D(\Omega)$
jest ciągły;
\item wzór $\langle PT,\varphi\rangle=
\langle T,P^{\mathrm t}\varphi\rangle$ wyznacza dystrybucję;
\item $P$ jest ciągły w topologiach słabej-* i silnej;
\item $\supp(PT)\subseteq\supp T$;
\item jeżeli $f\in C^m(\Omega)$, to $PT_f=T_{Pf}$.
\end{enumerate}
\end{twierdzenie}
\begin{proof}
Każdy składnik $P^{\mathrm t}$ jest złożeniem mnożenia
przez funkcję gładką z pochodną funkcji testowej, a obie
operacje są ciągłe. Zatem $P^{\mathrm t}$ jest ciągły;
zachowuje nośniki w sensie niepowiększania ich.
Przez złożenie z $T$ otrzymujemy ciągły funkcjonał liniowy.
Zależności
$q_\varphi(PT)=q_{P^{\mathrm t}\varphi}(T)$ oraz
$q_B(PT)=q_{P^{\mathrm t}(B)}(T)$
dają obie ciągłości, ponieważ $P^{\mathrm t}(B)$ jest ograniczony.
Jeśli $T$ zeruje się na otwartym $U$, to dla
$\varphi\in\mathcal D(U)$ również
$P^{\mathrm t}\varphi\in\mathcal D(U)$, stąd inkluzja nośników.
Dla $f\in C^m$ i $\varphi\in\mathcal D(\Omega)$
całkowanie przez części pokazuje
$\int_\Omega(Pf)(x)\varphi(x)\,\mathrm d\lambda_d(x)
=\int_\Omega f(x)(P^{\mathrm t}\varphi)(x)\,\mathrm d\lambda_d(x)$,
co daje ostatni punkt.
\end{proof}

\begin{przyklad}[Operator pierwszego rzędu]
\label{ex:first-order-P20}
W jednym wymiarze, gdy $P=a(x)\frac{\mathrm d}{\mathrm dx}+b(x)$,
$P^{\mathrm t}\varphi=-(a\varphi)'+b\varphi$.
W $\R^d$ dla pola $b=(b_1,\ldots,b_d)$ o składowych
gładkich, operator $P=b\cdot\nabla+c$ ma formalną
transpozycję $P^{\mathrm t}\varphi=-\operatorname{div}(b\varphi)
+c\varphi$. Przyjmujemy, że
$\operatorname{div}(b\varphi)=\sum_{j=1}^d\partial_j(b_j\varphi)$.
\end{przyklad}

\begin{wniosek}[Komutator pochodnej i mnożenia]
\label{cor:commutator20}
Dla $g\in C^\infty(\Omega)$ i $j\in\{1,\ldots,d\}$ mamy
$[\partial_j,M_g]=M_{\partial_jg}$ na $\mathcal D'(\Omega)$,
gdzie $[A,B]=AB-BA$ jest komutatorem operatorów.
\end{wniosek}
\begin{proof}
Jest to wzór Leibniza pierwszego rzędu:
$\partial_j(gT)-g\partial_jT=(\partial_jg)T$.
\end{proof}

\subsection{Równania różniczkowe w sensie dystrybucji}

\begin{twierdzenie}[Jądro pochodnej na przedziale]
\label{thm:zero-derivative20}
Niech $I\subseteq\R$ będzie niepustym przedziałem otwartym,
a $T\in\mathcal D'(I;\K)$. Wówczas $T'=0$
wtedy i tylko wtedy, gdy istnieje $c\in\K$ takie, że
$T=T_c$, czyli
$\langle T,\varphi\rangle=c\int_I\varphi(x)\,\mathrm d\lambda_1(x)$
dla każdej $\varphi\in\mathcal D(I)$.
\end{twierdzenie}
\begin{proof}
Wybierzmy $\eta\in\mathcal D(I)$ z
$\int_I\eta(x)\,\mathrm d\lambda_1(x)=1$.
Dla $\varphi\in\mathcal D(I)$ oznaczmy
$m_\varphi=\int_I\varphi(x)\,\mathrm d\lambda_1(x)$
i $h=\varphi-m_\varphi\eta$. Wtedy $\int_Ih=0$.
Określamy $\psi(x)=\int_{-\infty}^xh(t)\,\mathrm d\lambda_1(t)$,
przedłużając $h$ przez zero poza $I$.
Ponieważ $h$ ma zwarty nośnik zawarty w $I$ i całkę zero,
$\psi$ jest gładka i ma zwarty nośnik w $I$, przy czym
$\psi'=h$. Jeśli $T'=0$, to
$0=\langle T',\psi\rangle=-\langle T,\psi'\rangle$,
więc $\langle T,\varphi\rangle=m_\varphi\langle T,\eta\rangle$.
Wybieramy $c=\langle T,\eta\rangle$.
Odwrotnie, dla stałej dystrybucji $T_c$ mamy
$\langle T_c',\varphi\rangle=-c\int_I\varphi'(x)\,
\mathrm d\lambda_1(x)=0$.
\end{proof}

\begin{twierdzenie}[Istnienie pierwotnej dystrybucji]
\label{thm:primitive20}
Dla każdej dystrybucji $S\in\mathcal D'(\R;\K)$
istnieje $T\in\mathcal D'(\R;\K)$ takie, że $T'=S$.
Wszystkie takie pierwotne różnią się o stałą dystrybucję.
\end{twierdzenie}
\begin{proof}
Ustalmy $\eta\in\mathcal D(\R)$ o całce $1$ i
zdefiniujmy $m_\varphi=\int_\R\varphi(x)\,\mathrm d\lambda_1(x)$
oraz
\[
 R\varphi(x)=\int_{-\infty}^x
 \bigl(\varphi(t)-m_\varphi\eta(t)\bigr)\,\mathrm d\lambda_1(t).
\]
Jak w dowodzie twierdzenia~\ref{thm:zero-derivative20},
$R\varphi\in\mathcal D(\R)$. Odwzorowanie
$R:\mathcal D(\R)\to\mathcal D(\R)$ jest liniowe i ciągłe.
Istotnie, dla każdego zwartego $K\Subset\R$ jego ograniczenie
na $\mathcal D_K$ ma wartości w
$\mathcal D_L$, gdzie $L$ jest zwartym przedziałem
zawierającym $K\cup\supp\eta$.
Dla $\varphi\in\mathcal D_K$ mamy
$|m_\varphi|\le\lambda_1(K)p_0(\varphi)$, a pochodne
$R\varphi$ rzędu $j\ge1$ są pochodnymi rzędu $j-1$
funkcji $\varphi-m_\varphi\eta$.
Sama funkcja $R\varphi$ jest ograniczona przez długość $L$
razy supremum tej różnicy. Daje to dla każdego $m$ oszacowanie
$p_m(R\varphi)\le C_{K,m,\eta}p_m(\varphi)$
(po ewentualnym zastąpieniu stałej), zatem ciągłość na etapach
oraz na granicy induktywnej.
Definiujemy
$\langle T,\varphi\rangle=-\langle S,R\varphi\rangle$.
Wtedy $T\in\mathcal D'(\R)$.
Dla $\varphi\in\mathcal D(\R)$ zachodzi $m_{\varphi'}=0$
oraz $R(\varphi')=\varphi$, bo $\varphi$ znika daleko na lewo.
Stąd
$\langle T',\varphi\rangle=-\langle T,\varphi'\rangle
=\langle S,R\varphi'\rangle=\langle S,\varphi\rangle$.
Jednoznaczność z dokładnością do stałej wynika z
twierdzenia~\ref{thm:zero-derivative20}.
\end{proof}

\begin{wniosek}[Równanie pierwszego rzędu]
\label{cor:ode20}
Niech $a\in C^\infty(\R;\K)$ i
$S\in\mathcal D'(\R;\K)$. Wszystkie dystrybucyjne
rozwiązania równania $T'+aT=S$ mają postać
$T=e^{-A}V$, gdzie $A(x)=\int_0^xa(t)\,\mathrm d\lambda_1(t)$,
a $V$ jest dowolną pierwotną dystrybucją $e^AS$.
Dwa rozwiązania różnią się o $c e^{-A}$, $c\in\K$.
\end{wniosek}
\begin{proof}
Funkcja $e^A$ jest gładka i nigdzie nie znika.
Z reguły Leibniza $(e^AT)'=e^A(T'+aT)=e^AS$.
Z twierdzenia~\ref{thm:primitive20} istnieje pierwotna
$V$ dystrybucji $e^AS$ i każda inna różni się od niej
o stałą. Mnożenie przez $e^{-A}$ daje wszystkie rozwiązania.
\end{proof}

\begin{przyklad}[Rozwiązania dystrybucyjne i skoki]
\label{ex:forced-ode20}
Dla $\lambda\in\K$ rozwiązujemy równanie
$T'+\lambda T=\delta_0$ na $\R$.
Dystrybucja $E_\lambda=T_{e^{-\lambda x}H(x)}$ spełnia
$E_\lambda'+\lambda E_\lambda=\delta_0$, ponieważ
z twierdzenia o skoku
$(e^{-\lambda x}H)'=-\lambda e^{-\lambda x}H+\delta_0$.
Wszystkie rozwiązania mają więc postać
$T=E_\lambda+cT_{e^{-\lambda x}}$, $c\in\K$.
Dla $\lambda=0$ otrzymujemy równanie $T'=\delta_0$
i pierwotną $H$. Ponadto $x_+=xH(x)$ spełnia
$\frac{\mathrm d^2}{\mathrm dx^2}T_{x_+}=\delta_0$.
\end{przyklad}

\begin{twierdzenie}[Znikanie gradientu na zbiorze spójnym]
\label{thm:zero-gradient20}
Niech $\Omega\subseteq\R^d$ będzie niepustym spójnym
zbiorem otwartym i $T\in\mathcal D'(\Omega;\K)$.
Jeżeli $\partial_jT=0$ dla każdego $j=1,\ldots,d$,
to istnieje $c\in\K$ takie, że
$\langle T,\varphi\rangle=c\int_\Omega\varphi(x)\,
\mathrm d\lambda_d(x)$ dla każdej $\varphi\in\mathcal D(\Omega)$.
\end{twierdzenie}
\begin{proof}
Najpierw rozważmy prostokąt $Q=I_1\times\cdots\times I_d$
o domknięciu zwartym w $\Omega$. Dla każdego $j$ wybieramy
$\eta_j\in\mathcal D(I_j)$ o całce $1$ i określamy operator
\[
 (P_j\varphi)(x)=\eta_j(x_j)
 \int_{I_j}\varphi(x_1,\ldots,x_{j-1},t,x_{j+1},\ldots,x_d)
 \mathrm d\lambda_1(t).
\]
Operatory $P_j$ komutują. Funkcja $(I-P_j)\psi$ ma całkę
zero po zmiennej $x_j$ dla każdego ustalenia pozostałych
współrzędnych i ma zwarty nośnik w $Q$.
Przez całkowanie względem $x_j$ od lewego końca przedziału
otrzymujemy $\theta_j\in\mathcal D(Q)$ z
$\partial_j\theta_j=(I-P_j)\psi$. Ponieważ
$\partial_jT=0$, dla takich $\psi$ zachodzi
$\langle T,(I-P_j)\psi\rangle=0$.
Z tożsamości teleskopowej
\[
 I-P_1\cdots P_d
 =(I-P_1)+P_1(I-P_2)+\cdots+P_1\cdots P_{d-1}(I-P_d)
\]
wynika
$\langle T,\varphi\rangle=\langle T,P_1\cdots P_d\varphi\rangle$.
Ponieważ
$P_1\cdots P_d\varphi=(\int_Q\varphi\,\mathrm d\lambda_d)
\eta_1\cdots\eta_d$,
na $Q$ mamy stałą $c_Q=\langle T,\eta_1\cdots\eta_d\rangle$.
Dla dwóch nakładających się prostokątów $Q,Q'$ stałe
$c_Q,c_{Q'}$ są równe: wystarczy przetestować obie
reprezentacje na funkcji o niezerowej całce i nośniku
w $Q\cap Q'$. Ponieważ spójny zbiór otwarty w $\R^d$
jest łukowo spójny, łańcuch nakładających się prostokątów
łączy dowolne dwa lokalne otoczenia. Współczynniki
$c_Q$ są więc równe jednej stałej $c$.
Dla dowolnej $\varphi\in\mathcal D(\Omega)$ wybieramy
skończone pokrycie jej nośnika takimi prostokątami
i gładki podział jedności; zastosowanie lokalnego wzoru
do poszczególnych składników kończy dowód.
\end{proof}

\begin{uwaga}[Granice twierdzeń]
\label{rem:limits20}
Wielowymiarowe równanie $\partial_1T=0$ na obszarze
$\Omega=I\times V$ oznacza niezależność dystrybucji od
pierwszej współrzędnej (w odpowiednim sensie produktu
tensorowego); nie oznacza, że $T$ jest stała, jeśli
nie znikają pozostałe pochodne. Ponadto równania
$\partial_jT=S_j$ dla wszystkich $j$ wymagają warunków
zgodności $\partial_kS_j=\partial_jS_k$ i globalnych
warunków topologicznych dziedziny. Pełną teorię
omówimy przy słabych rozwiązaniach PDE.
\end{uwaga}

\begin{uwaga}[Związek z operatorami na TVS]
\label{rem:tvs-derivative20}
Konstrukcja $T\mapsto PT$ nie wymaga normy na
$\mathcal D(\Omega)$ ani na $\mathcal D'(\Omega)$.
Opiera się na ciągłości operatora $P^{\mathrm t}$ na
lokalnie wypukłej przestrzeni funkcji testowych.
W ten sposób pochodna dystrybucyjna jest szczególnym
przypadkiem transpozycji ciągłego operatora liniowego
między przestrzeniami liniowo-topologicznymi.
\end{uwaga}


% Wersja źródłowa: audyt II, jednostka 21
\setcounter{section}{20}
\section{Splot dystrybucji i regularyzacja}
\label{sec:splot21}

W całym rozdziale $d\ge1$, $\K\in\{\R,\C\}$, a dystrybucje
są ciągłymi $\K$-liniowymi funkcjonałami na
$\mathcal D(\R^d;\K)$. Rozpatrujemy splot na całej
$\R^d$; dla ogólnego zbioru otwartego $\Omega$ nie ma
jednej operacji splotu dwóch dystrybucji z
$\mathcal D'(\Omega)$ bez dodatkowych ustaleń.
Oznaczymy przez $\mathcal E'(\R^d)$ przestrzeń dystrybucji
o zwartym nośniku. Nie należy mylić jej ze zbiorem funkcji
testowych $\mathcal D(\R^d)$.

\subsection{Dystrybucje o zwartym nośniku i funkcje gładkie}

\begin{lemat}[Przedłużenie dystrybucji o zwartym nośniku]
\label{lem:compact-action21}
Niech $T\in\mathcal D'(\R^d;\K)$ i
$K=\supp T\Subset\R^d$. Wówczas dla każdej funkcji
$f\in C^\infty(\R^d;\K)$ istnieje dobrze określona wartość
$\langle T,f\rangle:=\langle T,\chi f\rangle$, gdzie
$\chi\in\mathcal D(\R^d)$ jest równa $1$ na otoczeniu $K$.
Wartość ta nie zależy od wyboru $\chi$. Istnieje także
zbiór zwarty $L$ i liczba $m\in\N_0$ oraz $C>0$ takie, że
\[
 |\langle T,f\rangle|\le
 C\max_{|\alpha|\le m}\sup_{x\in L}|\partial^\alpha f(x)|
 \quad(f\in C^\infty(\R^d)).
\]
W szczególności $T$ wyznacza ciągły funkcjonał na
przestrzeni Frécheta $C^\infty(\R^d)$.
\end{lemat}
\begin{proof}
Funkcja $\chi$ istnieje na mocy lematu o gładkim odcięciu.
Jeżeli $\chi_1,\chi_2$ są równe $1$ w otoczeniach $K$,
to $(\chi_1-\chi_2)f$ zeruje się na otoczeniu $K$;
z definicji nośnika dystrybucji $T$ zeruje się na tej funkcji.
Dla ustalonej $\chi$ wybieramy $L=\supp\chi$.
Ciągłość $T$ na $\mathcal D_L$ daje nierówność
$|\langle T,\psi\rangle|\le C_0p_m(\psi)$ dla funkcji
$\psi\in\mathcal D_L$. Podstawiamy $\psi=\chi f$
i stosujemy wzór Leibniza. Pochodne $\chi$ są ograniczone,
więc uzyskujemy wskazane oszacowanie, a zatem również
ciągłość na $C^\infty(\R^d)$.
\end{proof}

\begin{uwaga}[Przestrzeń dualna przestrzeni $C^\infty$]
\label{rem:Eprime21}
W drugą stronę, każdy ciągły funkcjonał na
$C^\infty(\R^d)$ definiuje dystrybucję o zwartym nośniku:
ciągłość w topologii seminorm zależnych od zbiorów zwartych
zapewnia oszacowanie wartości funkcjonału przez skończenie
wiele seminorm na jednym zwartym zbiorze. Funkcjonał
zeruje się więc na funkcjach testowych z nośnikiem
rozłącznym z tym zbiorem. Otrzymujemy naturalną identyfikację
$\mathcal E'(\R^d)=C^\infty(\R^d)'$, przy czym poniżej
nie będziemy utożsamiać bez wyjaśnienia różnych możliwych
topologii na tej przestrzeni dualnej.
\end{uwaga}

\subsection{Iterowane działania dystrybucji i definicja splotu}

\begin{lemat}[Przestawianie dwóch działań dystrybucji]
\label{lem:iterated-distributions21}
Niech $T\in\mathcal D'(\R^k)$,
$S\in\mathcal D'(\R^l)$ i
$\Phi\in\mathcal D(\R^{k+l})$. Funkcje
\[
 x\longmapsto\langle S_y,\Phi(x,y)\rangle,
 \qquad
 y\longmapsto\langle T_x,\Phi(x,y)\rangle
\]
należą odpowiednio do $\mathcal D(\R^k)$ i
$\mathcal D(\R^l)$. Ponadto
\[
 \langle T_x,\langle S_y,\Phi(x,y)\rangle\rangle
 =\langle S_y,\langle T_x,\Phi(x,y)\rangle\rangle.
\]
\end{lemat}
\begin{proof}
Różniczkowanie względem zewnętrznej zmiennej można
przenieść pod znak działania dystrybucji, ponieważ
odpowiednie rodziny funkcji testowych mają lokalnie wspólny
zwarty nośnik i zbiegają we wszystkich seminormach.
Nośniki obu otrzymanych funkcji mieszczą się w rzutach
zwartego zbioru $\supp\Phi$.

Dla $\Phi(x,y)=u(x)v(y)$ równość działań jest oczywista.
Takie sumy skończone są gęste w $\mathcal D(\R^{k+l})$:
wybieramy prostokąt z wnętrzem zawierającym $\supp\Phi$,
przedłużamy funkcję przez zero do gładkiej funkcji okresowej
na odpowiednim torusie, a następnie aproksymujemy ją
skończonymi sumami Fouriera. Sumy te zbiegają wraz ze
wszystkimi pochodnymi. Mnożąc je przez iloczyn dwóch
funkcji odcięcia równych $1$ na nośniku $\Phi$,
otrzymujemy skończone sumy $u_j(x)v_j(y)$ o wspólnym
zwartym nośniku, zbiegające do $\Phi$ w $\mathcal D$.
Obie strony żądanej równości są ciągłymi funkcjonałami
liniowymi na $\mathcal D(\R^{k+l})$ (wynika to z
oszacowań na odpowiednich etapach $\mathcal D_K$).
Zatem równość przechodzi na granicę.
\end{proof}

\begin{lemat}[Operator testowy indukowany dystrybucją]
\label{lem:convolution-transpose21}
Niech $T\in\mathcal E'(\R^d)$, $K=\supp T$, i niech
$\varphi\in\mathcal D(\R^d)$. Funkcja
\[
 (A_T\varphi)(y)=\langle T_x,\varphi(x+y)\rangle
\]
należy do $\mathcal D(\R^d)$ i spełnia
$\supp(A_T\varphi)\subseteq\supp\varphi-K$.
Operator $A_T:\mathcal D(\R^d)\to\mathcal D(\R^d)$
jest liniowy i ciągły, a obrazy zbiorów ograniczonych
są ograniczone.
\end{lemat}
\begin{proof}
W działaniu $T_x$ korzystamy z przedłużenia z
Lematu~\ref{lem:compact-action21}. Dla $y$ z ustalonego
zbioru zwartego funkcje $x\mapsto\varphi(x+y)$ mają
wspólny zwarty nośnik; można więc różniczkować po $y$:
$\partial_y^\beta(A_T\varphi)(y)=
\langle T_x,\partial^\beta\varphi(x+y)\rangle$.
Jeżeli $y\notin\supp\varphi-K$, funkcja
$\varphi(\cdot+y)$ zeruje się na otoczeniu $K$,
stąd wartość działania $T$ jest zerowa. Zatem obraz
jest funkcją testową.

Dla ustalonego $L\Subset\R^d$ i
$\varphi\in\mathcal D_L$ oszacowanie z
Lematu~\ref{lem:compact-action21}, zastosowane do
przesunięć funkcji $\varphi$, daje
\[
 p_m(A_T\varphi)\le C_{T,L,m}\,p_{m+r}(\varphi),
 \qquad \supp(A_T\varphi)\subseteq L-K,
\]
gdzie $r$ jest odpowiednio dużym, ustalonym rzędem
oszacowania dystrybucji $T$ w otoczeniu $K$.
To dowodzi ciągłości na każdym etapie $\mathcal D_L$,
a z własności uniwersalnej granicy induktywnej
wynika ciągłość na $\mathcal D$. Ciągła mapa liniowa
przenosi zbiory ograniczone na ograniczone.
\end{proof}

\begin{definicja}[Splot dystrybucji]
\label{def:convolution21}
Niech $T\in\mathcal E'(\R^d)$ i
$S\in\mathcal D'(\R^d)$. Ich \emph{splotem} nazywamy
dystrybucję $T*S\in\mathcal D'(\R^d)$ określoną wzorem
\[
 \langle T*S,\varphi\rangle
 =\langle S_y,\langle T_x,\varphi(x+y)\rangle\rangle
 =\langle S,A_T\varphi\rangle,
 \qquad\varphi\in\mathcal D(\R^d).
\]
Jeżeli tylko drugi czynnik ma zwarty nośnik, definiujemy
$S*T:=T*S$. Gdy oba czynniki mają zwarty nośnik,
zgodność obu możliwych kolejności zostanie dowiedziona
poniżej. W szczególności nie definiujemy tutaj splotu
dowolnych dwóch dystrybucji bez warunków na nośniki.
\end{definicja}

\begin{twierdzenie}[Istnienie, ciągłość i nośnik splotu]
\label{thm:convolution-properties21}
Niech $T\in\mathcal E'(\R^d)$ i
$S\in\mathcal D'(\R^d)$. Wówczas:
\begin{enumerate}
\item splot $T*S$ jest dobrze określoną dystrybucją;
\item $\supp(T*S)\subseteq\supp T+\supp S$;
\item dla ustalonego $T$ odwzorowanie
$S\mapsto T*S$ jest ciągłym operatorem liniowym na
$\mathcal D'(\R^d)$ zarówno dla topologii słabej-*, jak
i dla topologii silnej;
\item jeżeli także $S$ ma zwarty nośnik, to $T*S$ ma
zwarty nośnik.
\end{enumerate}
\end{twierdzenie}
\begin{proof}
Pierwszy punkt jest konsekwencją ciągłości $A_T$:
złożenie ciągłego funkcjonału $S$ z $A_T$ jest ciągłym
funkcjonałem na $\mathcal D$. Dla $\varphi$ o nośniku
rozłącznym z $\supp T+\supp S$, funkcja
$A_T\varphi$ zeruje się w otoczeniu $\supp S$;
stąd $\langle S,A_T\varphi\rangle=0$.
To daje inkluzję nośników; suma zbioru zwartego i zbioru
domkniętego w $\R^d$ jest domknięta.

Dla ustalonej funkcji testowej $\varphi$ zachodzi
$|\langle T*S,\varphi\rangle|=
|\langle S,A_T\varphi\rangle|$, co zapewnia
ciągłość słabą-*. Dla zbioru ograniczonego
$B\subseteq\mathcal D$ mamy
\[
 \sup_{\varphi\in B}|\langle T*S,\varphi\rangle|
 =\sup_{\psi\in A_T(B)}|\langle S,\psi\rangle|.
\]
Lemat~\ref{lem:convolution-transpose21} daje
ograniczoność $A_T(B)$, a więc ciągłość silną.
Jeśli oba nośniki są zwarte, ich suma jest zwarta,
co daje ostatni punkt.
\end{proof}

\subsection{Przemienność, łączność i różniczkowanie}

\begin{twierdzenie}[Przemienność i łączność splotu]
\label{thm:convolution-algebra21}
Niech $T,S,U\in\mathcal D'(\R^d)$.
Jeśli co najmniej jeden z czynników $T,S$ ma zwarty nośnik,
to $T*S=S*T$, a wynik nie zależy od sposobu użycia
Definicji~\ref{def:convolution21}. Jeśli co najmniej dwa
spośród $T,S,U$ mają zwarty nośnik, to
\[
 (T*S)*U=T*(S*U).
\]
W szczególności $\mathcal E'(\R^d)$ jest przemienną
algebrą z jedynką $\delta_0$ względem splotu,
a $\mathcal D'(\R^d)$ jest modułem nad tą algebrą.
\end{twierdzenie}
\begin{proof}
Jeżeli $T$ i $S$ mają zwarty nośnik, wybieramy funkcje
odcięcia $\chi_T,\chi_S$, równe $1$ na otoczeniach
odpowiednich nośników. Dla $\varphi\in\mathcal D$ funkcja
$\Phi(x,y)=\chi_T(x)\chi_S(y)\varphi(x+y)$
należy do $\mathcal D(\R^{2d})$. Oba możliwe
obliczenia splotu są iterowanymi działaniami na $\Phi$;
równość wynika z Lematu~\ref{lem:iterated-distributions21}.
Jeśli tylko jeden czynnik ma zwarty nośnik,
przemienność jest częścią definicji, a powyższy
argument zapewnia jej niesprzeczność.

Dla łączności można, bez utraty ogólności, założyć,
że $T$ i $S$ mają zwarty nośnik. Dla ustalonej
$\varphi\in\mathcal D$ obie strony są trzema
iterowanymi działaniami na $\varphi(x+y+z)$.
Wstawiamy funkcje odcięcia w zmiennych $x$ i $y$
równe $1$ na otoczeniach nośników $T$ i $S$,
a w zmiennej $z$ funkcję odcięcia równą $1$
na zbiorze $\supp\varphi-\supp T-\supp S$.
Uzyskujemy funkcję testową trzech zmiennych.
Kolejne zastosowania Lematu~\ref{lem:iterated-distributions21}
pozwalają dowolnie zmieniać kolejność działań.
To daje łączność; pozostałe ustawienia czynników
wynikają z przemienności. Wzory
$\delta_0*T=T$ oraz $T*(S+U)=T*S+T*U$
wynikają bezpośrednio z definicji.
\end{proof}

\begin{definicja}[Translacja dystrybucji]
\label{def:translation21}
Dla $a\in\R^d$ i $T\in\mathcal D'(\R^d)$ definiujemy
$\tau_aT$ wzorem
$\langle\tau_aT,\varphi\rangle=
\langle T,\varphi(\cdot+a)\rangle$.
Dla dystrybucji regularnej $T_f$ odpowiada to
przesunięciu funkcji $f(\cdot-a)$.
\end{definicja}

\begin{przyklad}[Dystrybucje Diraca i przesunięcia]
\label{ex:delta-convolution21}
Dla $a,b\in\R^d$ i $T\in\mathcal D'$ mamy
$\delta_a*T=\tau_aT$, $\delta_a*\delta_b=\delta_{a+b}$.
Dla wielowskaźników $\alpha,\beta\in\N_0^d$ otrzymujemy
$\partial^\alpha\delta_0*\partial^\beta\delta_0
=\partial^{\alpha+\beta}\delta_0$.
\end{przyklad}

\begin{twierdzenie}[Pochodne splotu]
\label{thm:convolution-derivative21}
Niech $T,S\in\mathcal D'(\R^d)$ i niech co najmniej
jeden z nich ma zwarty nośnik. Dla każdego
$\alpha\in\N_0^d$ zachodzi
\[
 \partial^\alpha(T*S)
 =(\partial^\alpha T)*S
 =T*(\partial^\alpha S).
\]
W szczególności dla wielomianu $P$ w $d$ zmiennych
i operatora stałowspółczynnikowego $P(\partial)$
operator splotu z ustaloną dystrybucją o zwartym
nośniku komutuje z $P(\partial)$.
\end{twierdzenie}
\begin{proof}
Załóżmy, że $T$ ma zwarty nośnik. Dla
$\varphi\in\mathcal D$ i $j\in\{1,\ldots,d\}$ mamy
\[
 \begin{aligned}
 \langle\partial_j(T*S),\varphi\rangle
 &=-\langle S_y,\langle T_x,\partial_j\varphi(x+y)\rangle\rangle\\
 &=\langle S_y,\langle(\partial_jT)_x,\varphi(x+y)\rangle\rangle
 =\langle(\partial_jT)*S,\varphi\rangle.
 \end{aligned}
\]
W drugim kroku różniczkowanie dotyczy zmiennej $x$.
Aby uzyskać $T*(\partial_jS)$, zauważamy, że
$\partial_{y_j}A_T\varphi=A_T(\partial_j\varphi)$,
więc definicja pochodnej $S$ daje tę samą wartość.
Indukcja dowodzi wzoru dla $\alpha$. Gdy zwarty
nośnik ma tylko $S$, stosujemy przemienność.
\end{proof}

\begin{przyklad}[Inkluzja nośników może być ostra]
\label{ex:support-strict21}
Niech $T=\delta_0-\delta_1\in\mathcal E'(\R)$,
a $S=T_1$ będzie dystrybucją odpowiadającą funkcji
stałej $1$. Wtedy $T*S=T_1-T_1=0$,
mimo że $\supp T+\supp S=\R$. Dlatego
inkluzji nośników nie można zastąpić równością.
\end{przyklad}

\subsection{Splot z funkcją testową i aproksymacja jedności}

\begin{twierdzenie}[Wygładzanie dystrybucji]
\label{thm:smoothing21}
Niech $T\in\mathcal D'(\R^d)$ i
$\rho\in\mathcal D(\R^d)$. Wówczas splot
$T*\rho$ jest dystrybucją regularną odpowiadającą
funkcji $u\in C^\infty(\R^d)$ danej przez
\[
 u(x)=\langle T_y,\rho(x-y)\rangle.
\]
Dla każdego wielowskaźnika $\alpha$ zachodzi
$\partial^\alpha u(x)
=\langle T_y,(\partial^\alpha\rho)(x-y)\rangle$.
Ponadto $\supp u\subseteq\supp T+\supp\rho$.
\end{twierdzenie}
\begin{proof}
Przy $x$ należącym do ustalonego zbioru zwartego
rodzina funkcji testowych $y\mapsto\rho(x-y)$
ma wspólny zwarty nośnik. Różniczkowanie po parametrze
$x$ zachodzi więc w topologii $\mathcal D$, skąd
gładkość $u$ i wzory na pochodne. Jeśli $x$
nie należy do $\supp T+\supp\rho$, funkcja
$y\mapsto\rho(x-y)$ znika na otoczeniu $\supp T$,
a zatem $u(x)=0$. Równość dystrybucji $T*\rho=T_u$
sprawdzamy, stosując Lemat~\ref{lem:iterated-distributions21}
do funkcji testowej $\rho(x-y)\varphi(x)$:
\[
 \begin{aligned}
 \int_{\R^d}u(x)\varphi(x)\,\mathrm d\lambda_d(x)
 &=\left\langle T_y,
      \int_{\R^d}\rho(x-y)\varphi(x)\,\mathrm d\lambda_d(x)
   \right\rangle\\
 &=\langle T,\check\rho*\varphi\rangle
 =\langle T*\rho,\varphi\rangle,
 \end{aligned}
\]
gdzie $\check\rho(z)=\rho(-z)$. W pierwszym
kroku zamiana całkowania i działania dystrybucji
jest dozwolona również bez globalnego zwartego nośnika
$T$, ponieważ $\varphi$ i $\rho$ mają zwarte nośniki.
\end{proof}

\begin{definicja}[Standardowa aproksymacja jedności]
\label{def:mollifier21}
Wybierzmy $\rho\in\mathcal D(\R^d)$ taką, że
$\rho\ge0$, $\supp\rho\subseteq\overline{B(0,1)}$
oraz $\int_{\R^d}\rho(x)\,\mathrm d\lambda_d(x)=1$.
Dla $\varepsilon>0$ piszemy
$\rho_\varepsilon(x)=\varepsilon^{-d}\rho(x/\varepsilon)$.
\end{definicja}

\begin{twierdzenie}[Regularyzacja jest silnie zbieżna]
\label{thm:strong-mollification21}
Dla każdego $T\in\mathcal D'(\R^d)$ i każdej
standardowej aproksymacji jedności z
Definicji~\ref{def:mollifier21} dystrybucje
$T*\rho_\varepsilon$ są reprezentowane przez funkcje
$C^\infty(\R^d)$ oraz
\[
 T*\rho_\varepsilon\longrightarrow T
 \quad\text{w topologii silnej }\beta(\mathcal D',\mathcal D)
 \quad(\varepsilon\downarrow0).
\]
Jeśli $T$ ma zwarty nośnik $K$, to
$\supp(T*\rho_\varepsilon)\subseteq K+
\overline{B(0,\varepsilon)}$, więc dla małych
$\varepsilon$ funkcje wygładzające mają nośniki
w jednym zbiorze zwartym.
\end{twierdzenie}
\begin{proof}
Z Twierdzenia~\ref{thm:smoothing21} mamy wygładzanie.
Dla $\varphi\in\mathcal D$ zachodzi
\[
 \langle T*\rho_\varepsilon-T,\varphi\rangle
 =\langle T,\check\rho_\varepsilon*\varphi-\varphi\rangle.
\]
Niech $B\subseteq\mathcal D$ będzie ograniczony.
Z twierdzenia o ograniczoności w $\mathcal D$
wynika istnienie $K\Subset\R^d$ zawierającego nośniki
wszystkich $\varphi\in B$, przy czym
$M_m(B)=\sup_{\varphi\in B}p_m(\varphi)<\infty$
dla wszystkich $m$. Dla $0<\varepsilon\le1$
różnice $\check\rho_\varepsilon*\varphi-\varphi$
mają nośniki w jednym zwartym zbiorze
$L=K+\overline{B(0,1)}$. Twierdzenie o wartości
średniej zastosowane do pochodnych rzędu $\le m$ daje
\[
 \sup_{\varphi\in B}p_m(\check\rho_\varepsilon*\varphi
 -\varphi)
 \le C_{\rho,d}\varepsilon\,M_{m+1}(B),
 \quad
 C_{\rho,d}=\int_{\R^d}\|z\|_1\rho(z)\,\mathrm d\lambda_d(z).
\]
Ciągłość $T|_{\mathcal D_L}$ oznacza, że dla pewnych
$m,C$ mamy $|\langle T,\psi\rangle|\le Cp_m(\psi)$
gdy $\psi\in\mathcal D_L$. Dlatego
\[
 \sup_{\varphi\in B}|\langle T*\rho_\varepsilon-T,
 \varphi\rangle|
 \le CC_{\rho,d}\varepsilon M_{m+1}(B)\longrightarrow0.
\]
Ponieważ $B$ jest dowolnym zbiorem ograniczonym,
jest to zbieżność silna. Oszacowanie nośnika wynika
z Twierdzenia~\ref{thm:smoothing21}.
\end{proof}

\begin{wniosek}[Gęstość funkcji testowych dla zwartych nośników]
\label{cor:dense-test-compact21}
Niech $T\in\mathcal E'(\R^d)$. Istnieje ciąg
$(f_n)\subseteq\mathcal D(\R^d)$ taki, że
$T_{f_n}\to T$ silnie w $\mathcal D'(\R^d)$.
Można wziąć $f_n=T*\rho_{1/n}$.
\end{wniosek}
\begin{proof}
Funkcje $f_n$ są gładkie i mają nośniki zwarte
z poprzedniego twierdzenia. Zbieżność silna także
wynika z tego twierdzenia.
\end{proof}

\begin{przyklad}[Regularyzacja delty Diraca]
\label{ex:delta-smooth21}
Dla $T=\delta_a$ zachodzi
$\delta_a*\rho_\varepsilon=\rho_\varepsilon(\cdot-a)$.
Funkcje te stają się coraz bardziej skupione
w otoczeniu $a$, a jako dystrybucje zbiegają
silnie do $\delta_a$. Nie ma tu zbieżności w
normie $L^1_{\mathrm{loc}}$ do funkcji:
$\delta_a$ nie jest dystrybucją regularną
reprezentowaną przez funkcję lokalnie całkowalną.
\end{przyklad}

\subsection{Ograniczenia operacji splotu i powiązania z dalszą teorią}

\begin{uwaga}[Kiedy splot nie jest zdefiniowany]
\label{rem:convolution-limits21}
Warunek zwartego nośnika przynajmniej jednego
czynnika jest wystarczający, ale niekonieczny.
Dla funkcji odpowiednio całkowalnych splot może istnieć
bez tego warunku. Natomiast dla dowolnych dwóch
dystrybucji nie istnieje jedna kanoniczna operacja
splotu o wszystkich powyższych własnościach.
Na przykład formalny splot dwóch funkcji stałych
na $\R^d$ wymagałby całkowania stałej $1$
po całej $\R^d$, co nie daje skończonej wartości.
Dalsze rozszerzenia opierają się między innymi
na warunkach na nośniki umożliwiających właściwe
rzutowanie ich sum lub na dodatkowych własnościach
analitycznych czynników.
\end{uwaga}

\begin{uwaga}[Dlaczego jest to zagadnienie analizy funkcjonalnej]
\label{rem:convolution-tvs21}
Kluczowe kroki nie polegały na rozwiązywaniu równań
różniczkowych. Użyliśmy ciągłości operatora
$A_T$ na przestrzeni funkcji testowych, transpozycji
na jej przestrzeń dualną, przenoszenia zbiorów ograniczonych
oraz zbieżności silnej w tej przestrzeni dualnej.
Są to konstrukcje z teorii przestrzeni liniowo-topologicznych.
Dalsze zastosowania splotu do rozwiązywania równań
różniczkowych pozostawiamy odrębnemu skryptowi.
\end{uwaga}


% Wersja źródłowa: audyt II, jednostka 22
\setcounter{section}{21}
\noindent\textbf{Fakty z wcześniejszych rozdziałów.}
Przyjmujemy, że przestrzeń Schwartza jest zupełną przestrzenią Frécheta,
$\mathcal D(\R^d)$ jest w niej gęsta, a transformacja Fouriera
jest jej ciągłym automorfizmem z odwrotnością i wzorem
$\mathcal F^2\varphi=\varphi(-\cdot)$.
Dystrybucja o zwartym nośniku $K$ działa również na funkcjach
$C^\infty$ poprzez $\langle T,f\rangle=\langle T,\chi f\rangle$,
gdzie $\chi=1$ w otoczeniu $K$, i spełnia oszacowanie
skończonego rzędu na ustalonym zwartym otoczeniu $K$.
Te rezultaty stanowią punkt wyjścia niniejszego rozdziału.

\section{Dystrybucje temperowane i transformacja Fouriera}
\label{sec:temperowane22}

W całym rozdziale $d\ge 1$ i rozważamy funkcje oraz dystrybucje
zespolone na $\R^d$. Przestrzeń Schwartza
$\mathcal S=\mathcal S(\R^d;\C)$ otrzymała w
wcześniejszym rozdziale o przestrzeni Schwartza swoją topologię Frécheta.
Dla $N,m\in\N_0$ przyjmujemy seminormy
$Q_{N,m}(\varphi)=\max_{|\alpha|\le m}\sup_{x\in\R^d}
(1+|x|)^N|\partial^\alpha\varphi(x)|$.
Przez $\mathcal S'$ rozumiemy przestrzeń \emph{ciągłych,
zespolono-liniowych} funkcjonałów na $\mathcal S$.
W parowaniu $\langle T,\varphi\rangle$ nie sprzęgamy zespolenie
funkcji testowej. Przypominamy, że
$\mathcal D(\R^d)=C_c^\infty(\R^d)$ jest gęsta w $\mathcal S$.

\subsection{Dystrybucje temperowane jako przestrzeń dualna}

\begin{definicja}[Dystrybucja temperowana]
\label{def:tempered22}
Dystrybucją temperowaną na $\R^d$ nazywamy element
$T\in\mathcal S'$, czyli ciągły liniowy funkcjonał
$T:\mathcal S\to\C$. Rozpatrujemy na $\mathcal S'$ topologię
słabą-* $\sigma(\mathcal S',\mathcal S)$, wyznaczoną przez
seminormy $T\mapsto|\langle T,\varphi\rangle|$
($\varphi\in\mathcal S$), oraz topologię silną
$\beta(\mathcal S',\mathcal S)$, wyznaczoną przez
$P_B(T)=\sup_{\varphi\in B}|\langle T,\varphi\rangle|$,
gdzie $B$ przebiega wszystkie zbiory ograniczone w $\mathcal S$.
\end{definicja}

\begin{twierdzenie}[Kryterium wzrostu seminormowego]
\label{thm:tempered-estimate22}
Niech $T:\mathcal S(\R^d;\C)\to\C$ będzie liniowy.
Następujące warunki są równoważne:
\begin{enumerate}
\item $T\in\mathcal S'(\R^d)$;
\item istnieją $C>0$ oraz $N,m\in\N_0$ takie, że
$|\langle T,\varphi\rangle|\le C Q_{N,m}(\varphi)$
dla wszystkich $\varphi\in\mathcal S$.
\end{enumerate}
W szczególności ograniczenie $T|_{\mathcal D}$ jest dystrybucją,
a odwzorowanie $\mathcal S'\to\mathcal D'$ dane przez
ograniczenie jest liniowe i różnowartościowe.
\end{twierdzenie}
\begin{proof}
Rodzina $Q_{j,j}$, $j\in\N_0$, jest rosnąca
i wyznacza topologię $\mathcal S$. Jeśli $T$ jest ciągły,
to dla pewnych $j$ i $\varepsilon>0$ warunek
$Q_{j,j}(\varphi)<\varepsilon$ implikuje
$|T(\varphi)|<1$. Liniowość i skalowanie dają
$|T(\varphi)|\le (2/\varepsilon)Q_{j,j}(\varphi)$;
gdy seminorma znika, wniosek wynika ze skalowania przez
dowolną liczbę. Odwrotna implikacja jest natychmiastowa.
Włożenie $\mathcal D\hookrightarrow\mathcal S$ jest ciągłe:
na każdym etapie $\mathcal D_K$ seminormy Schwartza
są oszacowane przez seminormy pochodnych na $K$.
Stąd ograniczenie $T|_{\mathcal D}$ jest ciągłe.
Jest też różnowartościowe jako odwzorowanie na przestrzeni
funkcjonałów, ponieważ $\mathcal D$ jest gęsta w $\mathcal S$.
\end{proof}

\begin{przyklad}[Dystrybucje regularne i miary]
\label{ex:tempered-regular22}
Jeżeli $f\in L^1_{\mathrm{loc}}(\R^d)$ i
$|f(x)|\le C(1+|x|)^M$ prawie wszędzie,
to $T_f(\varphi)=\int_{\R^d}f(x)\varphi(x)\,\mathrm d\lambda_d(x)$
określa element $\mathcal S'$. Istotnie, dla całkowitego
$N>M+d$ moduł tej całki jest ograniczony przez
$C'Q_{N,0}(\varphi)$.
Bardziej ogólnie, jeżeli $\mu$ jest zespoloną lokalnie skończoną
miarą borelowską i dla pewnego $N\ge0$ zachodzi
$\int_{\R^d}(1+|x|)^{-N}\,\mathrm d|\mu|(x)<\infty$,
to $\langle T_\mu,\varphi\rangle=\int\varphi(x)\,\mathrm d\mu(x)$
definiuje dystrybucję temperowaną.
Przykładami są funkcja stała $1$, wielomiany i miary Diraca.
\end{przyklad}

\begin{przyklad}[Dystrybucja nietemperowana]
\label{ex:non-tempered22}
Dystrybucja regularna wyznaczona przez
$f(x)=e^{|x|^2}$ należy do $\mathcal D'(\R^d)$,
ale nie jest ograniczeniem żadnego elementu $\mathcal S'$.
Aby to wykazać, weźmy $\chi\in\mathcal D(\R^d)$,
$0\le\chi\le1$, równą $1$ na $\overline B(0,1)$,
i funkcje $\varphi_R(x)=\chi(x/R)e^{-|x|^2/2}$.
Dla $R\to\infty$ zachodzi
$\varphi_R\to e^{-|x|^2/2}$ w $\mathcal S$
(na mocy oszacowania pochodnych przez szybki zanik Gaussa).
Jednak
$\langle T_f,\varphi_R\rangle
=\int_{\R^d}\chi(x/R)e^{|x|^2/2}\,\mathrm d\lambda_d(x)
\to+\infty$.
Gdyby istniało ciągłe przedłużenie na $\mathcal S$,
wartości te musiałyby zbiegać do skończonej liczby.
\end{przyklad}

\begin{stwierdzenie}[Podprzestrzenie klasycznych dystrybucji]
\label{prop:schwartz-embeddings22}
Istnieją naturalne, różnowartościowe i ciągłe odwzorowania
$\mathcal D\hookrightarrow\mathcal S\hookrightarrow
L^2(\R^d)\hookrightarrow\mathcal S'\hookrightarrow\mathcal D'$,
gdzie funkcję $f\in L^2$ identyfikujemy z funkcjonałem
$\langle T_f,\varphi\rangle=\int_{\R^d}f(x)\varphi(x)
\,\mathrm d\lambda_d(x)$.
Ponadto każda dystrybucja o zwartym nośniku jest temperowana:
$\mathcal E'(\R^d)\subseteq\mathcal S'(\R^d)$.
\end{stwierdzenie}
\begin{proof}
Pierwsze włożenie i jego ciągłość wynikają z definicji
topologii $\mathcal D$. Drugie jest ciągłe dzięki
oszacowaniom wielomianowym definiującym $\mathcal S$.
Dla $f\in L^2$ nierówność Cauchy'ego--Schwarza daje
$|\langle T_f,\varphi\rangle|\le\|f\|_2\|\varphi\|_2$;
także $\|\varphi\|_2\le C_dQ_{d+1,0}(\varphi)$.
Różnowartościowość wynika z gęstości $\mathcal D$
w $L^2$. Ostatnie włożenie wykazano w
Twierdzeniu~\ref{thm:tempered-estimate22}.
Jeżeli $T\in\mathcal E'$ ma zwarty nośnik,
to z przedłużenia jego działania na $C^\infty$ i oszacowania
skończonego rzędu (lemat o przedłużaniu dystrybucji o zwartym nośniku)
otrzymujemy $|T(\varphi)|\le C Q_{0,m}(\varphi)$.
\end{proof}

\subsection{Operatory na przestrzeni dystrybucji temperowanych}

\begin{twierdzenie}[Różniczkowanie i mnożenie przez wielomian]
\label{thm:tempered-operators22}
Dla każdego $\alpha,\beta\in\N_0^d$ operatory
$T\mapsto\partial^\alpha T$ i $T\mapsto x^\beta T$
są liniowe i ciągłe na $\mathcal S'$ zarówno dla topologii
$\sigma(\mathcal S',\mathcal S)$, jak i
$\beta(\mathcal S',\mathcal S)$, gdzie
\[
\langle\partial^\alpha T,\varphi\rangle
=(-1)^{|\alpha|}\langle T,\partial^\alpha\varphi\rangle,
\qquad\langle x^\beta T,\varphi\rangle
=\langle T,x^\beta\varphi\rangle.
\]
To samo dotyczy mnożenia przez każdą funkcję
$a\in C^\infty(\R^d)$, której każda pochodna
$\partial^\alpha a$ ma co najwyżej wzrost wielomianowy
(stopień wzrostu może zależeć od $\alpha$).
\end{twierdzenie}
\begin{proof}
Różniczkowanie i mnożenie przez $x^\beta$ są ciągłymi
operatorami liniowymi $\mathcal S\to\mathcal S$,
co wykazano w wcześniejszym rozdziale o przestrzeni Schwartza.
Dla ogólnego mnożnika $a$ wzór Leibniza prowadzi,
dla każdych $N,m$, do oszacowania
$Q_{N,m}(a\varphi)\le C_{N,m}Q_{N+M_m,m}(\varphi)$
z odpowiednim $M_m\in\N_0$. Zatem
$\varphi\mapsto a\varphi$ jest ciągły na $\mathcal S$.
Transpozycja ciągłego operatora liniowego na $\mathcal S$
jest ciągła w topologii słabej-* (przez ocenę na
ustalonych testach), a w silnej, ponieważ obrazy
zbiorów ograniczonych są ograniczone.
Znaki pochodnych wynikają z definicji transpozycji.
\end{proof}

\begin{lemat}[Względna zwartość zbiorów ograniczonych w $\mathcal S$]
\label{lem:montel-schwartz22}
Każdy zbiór ograniczony w $\mathcal S(\R^d)$ ma zwarte
domknięcie w jej topologii Frécheta.
\end{lemat}
\begin{proof}
Niech $B\subseteq\mathcal S$ będzie ograniczony i
$(\varphi_n)\subseteq B$ dowolnym ciągiem.
Dla każdej liczby całkowitej $R\ge1$ i każdego $m\ge0$
rodziny pochodnych $\partial^\alpha\varphi_n$ dla
$|\alpha|\le m$ są jednostajnie ograniczone i
jednostajnie równociągłe na $\overline B(0,R)$:
równoodległościowość wynika z ograniczenia pochodnych
rzędu $m+1$. Z twierdzenia Arzelà--Ascoliego
i argumentu diagonalnego uzyskujemy podciąg
$(\varphi_{n_k})$, którego wszystkie pochodne zbiegają
jednostajnie na wszystkich kulach domkniętych.
Granicą jest pewna funkcja $\varphi\in C^\infty$.
Dla ustalonych $N,m$ i $R\ge1$ poza $B(0,R)$ zachodzi
\[
(1+|x|)^N|\partial^\alpha(\varphi_{n_k}-\varphi_{n_l})(x)|
\le R^{-1}\bigl(Q_{N+1,m}(\varphi_{n_k})
+Q_{N+1,m}(\varphi_{n_l})\bigr)\quad(|\alpha|\le m).
\]
Prawa strona jest jednostajnie mała przy dużym $R$;
wewnątrz kuli $B(0,R)$ zbieżność wszystkich pochodnych
jest jednostajna. Dlatego $(\varphi_{n_k})$ jest ciągiem
Cauchy'ego we wszystkich seminormach $Q_{N,m}$.
Zupełność $\mathcal S$ daje zbieżność do elementu
$\mathcal S$. Ponieważ przestrzeń jest metryzowalna,
sekwencyjna względna zwartość ograniczonego zbioru
jest równoważna względnej zwartości.
\end{proof}

\begin{twierdzenie}[Zbieżność słaba-* i silna ciągów]
\label{thm:seq-topologies-schwartz22}
Niech $(T_n)$ będzie ciągiem dystrybucji temperowanych,
a $T\in\mathcal S'$. Wtedy
\[
T_n\longrightarrow T\text{ w }\sigma(\mathcal S',\mathcal S)
\quad\Longleftrightarrow\quad
T_n\longrightarrow T\text{ w }\beta(\mathcal S',\mathcal S).
\]
Mimo to topologie $\sigma(\mathcal S',\mathcal S)$ i
$\beta(\mathcal S',\mathcal S)$ nie są równe.
\end{twierdzenie}
\begin{proof}
Implikacja ze zbieżności silnej jest natychmiastowa.
Załóżmy zbieżność słabą-*. Rodzina
$\{T_n-T:n\ge1\}$ jest punktowo ograniczona.
Ponieważ $\mathcal S$ jest przestrzenią Frécheta,
a więc beczkową, twierdzenie Banacha--Steinhausa
implikuje jej równociągłość.
Dla dowolnego ograniczonego $B\subseteq\mathcal S$
Lemat~\ref{lem:montel-schwartz22} daje zwartość
$K=\overline B$. Równociągłość funkcjonałów i
punktowa zbieżność do zera zapewniają zbieżność
jednostajną na $K$: dla danego $\varepsilon>0$
wybieramy otoczenie $V$ z
$\sup_n|\langle T_n-T,v\rangle|<\varepsilon/2$
dla $v\in V$ i skończone pokrycie $K$ przesunięciami
$\varphi_j+V$. Następnie dla dostatecznie dużego $n$
wszystkie $|\langle T_n-T,\varphi_j\rangle|<\varepsilon/2$.
Wówczas $P_B(T_n-T)<\varepsilon$.

Aby wykazać różność topologii, wybierzmy niezerową funkcję
$\psi\in\mathcal D$ o nośniku w $B(0,1)$ i zdefiniujmy
$\psi_n(x)=e^{-n^2}\psi(x-3ne_1)$ dla $n\ge1$.
Zbiór $B=\{\psi_n:n\ge1\}$ jest ograniczony w
$\mathcal S$, a jego liniowa otoczka jest
nieskończenie wymiarowa (nośniki są rozłączne).
Otoczenie silne $\{T:P_B(T)<1/2\}$ nie zawiera
żadnego słabo-* otoczenia zera: dla dowolnie wielu
$\varphi_1,\ldots,\varphi_r\in\mathcal S$
można wybrać $\psi_n$ poza ich liniową otoczką.
Twierdzenie Hahna--Banacha pozwala wtedy skonstruować
ciągły funkcjonał $T\in\mathcal S'$ taki, że
$T(\varphi_j)=0$ dla wszystkich $j$, lecz
$T(\psi_n)=1$. Dowodzi to ścisłej nierówności topologii.
\end{proof}

\subsection{Rozszerzenie transformacji Fouriera}

\begin{definicja}[Transformacja Fouriera dystrybucji temperowanej]
\label{def:fourier-tempered22}
Dla $T\in\mathcal S'(\R^d)$ definiujemy
$\widehat T=\mathcal F T\in\mathcal S'$ przez
\[
\langle\mathcal F T,\varphi\rangle
=\langle T,\mathcal F\varphi\rangle
\qquad(\varphi\in\mathcal S).
\]
Występuje tu transformacja Fouriera na $\mathcal S$
z normalizacją $(2\pi)^{-d/2}$ i jądrem
$e^{-ix\cdot\xi}$ z wcześniejszego rozdziału o przestrzeni Schwartza.
\end{definicja}

\begin{twierdzenie}[Transformacja Fouriera jest automorfizmem]
\label{thm:fourier-auto-tempered22}
Transformacja $\mathcal F:\mathcal S'\to\mathcal S'$
jest ciągłym izomorfizmem liniowo-topologicznym
zarówno dla topologii $\sigma(\mathcal S',\mathcal S)$,
jak i $\beta(\mathcal S',\mathcal S)$.
Jej odwrotność dana jest transpozycją
$\mathcal F^{-1}$ na $\mathcal S$. Zachodzą wzory
\[
\mathcal F(\partial_jT)=i\xi_j\,\mathcal F T,
\qquad
\mathcal F(x_jT)=i\partial_{\xi_j}(\mathcal F T),
\qquad\mathcal F^2T=T(-\cdot),
\]
gdzie $\langle T(-\cdot),\varphi\rangle=
\langle T,\varphi(-\cdot)\rangle$.
\end{twierdzenie}
\begin{proof}
Twierdzenie o transformacji Fouriera na przestrzeni Schwartza mówi, że
$\mathcal F:\mathcal S\to\mathcal S$ jest ciągłym
izomorfizmem liniowo-topologicznym, a jego odwrotność
także jest ciągła. Transponowanie obu operatorów
zapewnia wzajemnie odwrotne odwzorowania na $\mathcal S'$.
Ciągłość słaba-* wynika z działania na pojedynczej
funkcji testowej. Dla silnej mamy
$P_B(\mathcal FT)=P_{\mathcal F(B)}(T)$;
$\mathcal F(B)$ jest ograniczony, gdy $B$ jest ograniczony,
co daje ciągłość. Odwrotność rozpatrujemy analogicznie.
Dla $\varphi\in\mathcal S$ i współrzędnej $j$ otrzymujemy
\[
\begin{aligned}
\langle\mathcal F(\partial_jT),\varphi\rangle
&=-\langle T,\partial_j\mathcal F\varphi\rangle
=-\langle T,\mathcal F(-ix_j\varphi)\rangle
=\langle i\xi_j\mathcal FT,\varphi\rangle,\\
\langle\mathcal F(x_jT),\varphi\rangle
&=\langle T,x_j\mathcal F\varphi\rangle
=\langle T,\mathcal F(-i\partial_j\varphi)\rangle
=\langle i\partial_{\xi_j}\mathcal FT,\varphi\rangle.
\end{aligned}
\]
W tych rachunkach zmienna w funkcji testowej i
wektor w przestrzeni dualnej są identyfikowane
tylko notacyjnie; kluczowe są wzory
$\partial\mathcal F\varphi=\mathcal F(-ix\varphi)$
i $x\mathcal F\varphi=\mathcal F(-i\partial\varphi)$.
Ostatnia tożsamość wynika z
$\mathcal F^2\varphi(x)=\varphi(-x)$ i transpozycji.
\end{proof}

\begin{przyklad}[Funkcja stała, delta Diraca i pochodne]
\label{ex:fourier-delta22}
Dla $a\in\R^d$ mamy w $\mathcal S'$ wzory
\[
\mathcal F\delta_a=(2\pi)^{-d/2}e^{-ia\cdot\xi},
\quad\mathcal F T_1=(2\pi)^{d/2}\delta_0,
\quad\mathcal F(\partial^\alpha\delta_0)
=(2\pi)^{-d/2}(i\xi)^\alpha.
\]
Pierwszy wynika z
$\langle\mathcal F\delta_a,\varphi\rangle
=(2\pi)^{-d/2}\int e^{-ia\cdot\xi}\varphi(\xi)
\,\mathrm d\lambda_d(\xi)$.
Drugi można uzyskać, stosując $\mathcal F^2T=T(-\cdot)$
do pierwszego wzoru dla $a=0$. Trzeci wynika z
Twierdzenia~\ref{thm:fourier-auto-tempered22}.
W szczególności delta Diraca nie jest funkcją,
ale jej transformata Fouriera jest funkcją gładką.
\end{przyklad}

\begin{twierdzenie}[Twierdzenie Plancherela]
\label{thm:plancherel22}
Transformacja Fouriera z normalizacją $(2\pi)^{-d/2}$
przedłuża się jednoznacznie do operatora unitarnego
$\mathcal F:L^2(\R^d;\C)\to L^2(\R^d;\C)$,
czyli dla $f,g\in L^2$ zachodzi
$\int_{\R^d}\widehat f(\xi)\overline{\widehat g(\xi)}
\,\mathrm d\lambda_d(\xi)
=\int_{\R^d}f(x)\overline{g(x)}\,\mathrm d\lambda_d(x)$.
Dla $f\in L^2$ transformata rozumiana w $L^2$
zgadza się z transformatą dystrybucyjną $T_f$.
\end{twierdzenie}
\begin{proof}
Dla $f,g\in\mathcal S$ wzór odwrotny oraz Fubini
(dopuszczalny dzięki $f,\widehat g\in L^1$) dają
\[
\int\widehat f(\xi)\overline{\widehat g(\xi)}
\,\mathrm d\lambda_d(\xi)
=\int f(x)\overline{g(x)}\,\mathrm d\lambda_d(x),
\]
gdzie w obu całkach obszarem całkowania jest $\R^d$.
Zatem $\mathcal F$ jest izometrią $\mathcal S$ w $L^2$.
Przestrzeń $\mathcal S$ jest gęsta w $L^2$,
więc operator przedłuża się jednoznacznie do izometrii
na $L^2$. Jego obraz jest domknięty, a zawiera
$\mathcal F(\mathcal S)=\mathcal S$, zbiór gęsty;
stąd jest to operator unitarny. Tożsamość polaryzacyjna
przedłuża zachowanie iloczynu skalarnego na całe $L^2$.
Jeśli $f_n\in\mathcal S$ i $f_n\to f$ w $L^2$,
to $T_{f_n}\to T_f$ w $\mathcal S'$, ponieważ
$|\langle T_{f_n-f},\varphi\rangle|
\le\|f_n-f\|_2\|\varphi\|_2$.
Ciągłość transformaty w $\mathcal S'$ i jej zgodność
z klasyczną na $\mathcal S$ zapewniają ostatnią tezę.
\end{proof}

\subsection{Splot z funkcją szybko malejącą i transformata dystrybucji o zwartym nośniku}

\begin{twierdzenie}[Splot z elementem przestrzeni Schwartza]
\label{thm:schwartz-convolution22}
Niech $T\in\mathcal S'(\R^d)$ i
$\rho\in\mathcal S(\R^d)$. Funkcja
$u(x)=\langle T_y,\rho(x-y)\rangle$ należy do
$C^\infty(\R^d)$, a każda jej pochodna ma
co najwyżej wzrost wielomianowy. Dystrybucja
regularna $T_u$ jest temperowana i oznaczamy ją $T*\rho$.
Zachodzi
\[
\mathcal F(T*\rho)=(2\pi)^{d/2}(\mathcal FT)\,(\mathcal F\rho),
\qquad \partial^\alpha(T*\rho)=T*(\partial^\alpha\rho).
\]
\end{twierdzenie}
\begin{proof}
Dla $x$ w zbiorze zwartym odwzorowanie
$x\mapsto\rho(x-\cdot)$ jest gładkie jako funkcja
z wartościami w $\mathcal S$; jego pochodne to
$x\mapsto(\partial^\alpha\rho)(x-\cdot)$.
Zatem można różniczkować pod działaniem $T$.
Z Twierdzenia~\ref{thm:tempered-estimate22} mamy
$|T(\varphi)|\le C Q_{N,m}(\varphi)$.
Wstawiając $\varphi(y)=\partial^\alpha\rho(x-y)$
i korzystając z
$1+|y|\le(1+|x|)(1+|x-y|)$,
otrzymujemy
$|\partial^\alpha u(x)|\le C_\alpha(1+|x|)^N$.
Stąd $u$ jest temperowana jako funkcja.
Dla $\psi\in\mathcal S$ przez zamianę działania $T$
i całkowania funkcji $\psi(x)\rho(x-y)$
(uzasadnioną przez zbieżność w seminormach Schwartza)
otrzymujemy
$\langle T_u,\psi\rangle=
\langle T,\check\rho*\psi\rangle$,
gdzie $\check\rho(x)=\rho(-x)$.
To także dowodzi ciągłości konstrukcji w $\mathcal S'$.
Z twierdzenia splotowego w $\mathcal S$ wynika
$\check\rho*\mathcal F\psi
=(2\pi)^{d/2}\mathcal F((\mathcal F\rho)\psi)$.
Zatem
\[
\begin{aligned}
\langle\mathcal F(T*\rho),\psi\rangle
&=\langle T,\check\rho*\mathcal F\psi\rangle\\
&=(2\pi)^{d/2}\langle\mathcal FT,
 (\mathcal F\rho)\psi\rangle,
\end{aligned}
\]
co daje pierwszy wzór. Drugi wynika z różniczkowania
pod działaniem dystrybucji.
\end{proof}

\begin{twierdzenie}[Transformata dystrybucji o zwartym nośniku]
\label{thm:compact-fourier22}
Jeżeli $T\in\mathcal E'(\R^d)$ ma zwarty nośnik $K$,
to jego transformata Fouriera jest dystrybucją regularną
reprezentowaną przez gładką funkcję
\[
F_T(\xi)=(2\pi)^{-d/2}\langle T_x,e^{-ix\cdot\xi}\rangle.
\]
Dla pewnego $m\in\N_0$ i każdego wielowskaźnika $\beta$
istnieje $C_\beta>0$ takie, że
$|\partial_\xi^\beta F_T(\xi)|
\le C_\beta(1+|\xi|)^m$
dla wszystkich $\xi\in\R^d$.
W działaniu na funkcji wykładniczej używamy
przedłużenia dystrybucji o zwartym nośniku
z lematu o przedłużaniu dystrybucji o zwartym nośniku.
\end{twierdzenie}
\begin{proof}
Wybierzmy $\chi\in\mathcal D$ równą $1$ na
otoczeniu $K$. Dla każdej $\beta$ możemy różniczkować
$x\mapsto\chi(x)(-ix)^\beta e^{-ix\cdot\xi}$
po parametrze $\xi$, gdyż funkcje mają wspólny zwarty nośnik.
Dlatego
$\partial_\xi^\beta F_T(\xi)
=(2\pi)^{-d/2}\langle T,(-ix)^\beta e^{-ix\cdot\xi}\rangle$.
Oszacowanie z lematu o przedłużaniu dystrybucji o zwartym nośniku
przy ustalonym rzędzie $m$ i zwartym otoczeniu nośnika
$T$, razem ze wzorem Leibniza, daje
$|\partial_\xi^\beta F_T(\xi)|
\le C_\beta(1+|\xi|)^m$.
Jeśli $\varphi\in\mathcal S$, to
$\mathcal F\varphi\in\mathcal S$; zamiana działania
$T$ i całkowania szybko malejącego parametru jest
dozwolona z tego samego oszacowania. Otrzymujemy
\[
\langle\mathcal FT,\varphi\rangle
=\int_{\R^d}F_T(\xi)\varphi(\xi)
\,\mathrm d\lambda_d(\xi),
\]
co identyfikuje dystrybucję.
\end{proof}

\begin{uwaga}[Zakres rozdziału]
\label{rem:scope22}
Transformacja Fouriera jest tu automorfizmem przestrzeni
dystrybucji temperowanych i podstawowym narzędziem teorii
operatorów. Jej zastosowania do rozwiązywania konkretnych
równań różniczkowych pozostawiamy osobnemu skryptowi.
Dalsze rozdziały rozwiną teorię przestrzeni Montela,
przestrzeni nuklearnych i produktów tensorowych topologicznych.
\end{uwaga}


% Wersja źródłowa: audyt II, jednostka 23
\setcounter{section}{22}
\noindent\textbf{Wymagane wyniki.} Korzystamy z topologii przestrzeni funkcji testowych i Schwartza, twierdzenia Banacha--Steinhausa, Arzelà--Ascoliego, rachunku Hilbertowskiego, transformaty Fouriera oraz znanych rezultatów o przestrzeniach testowych.
\medskip
\section{Przestrzenie Montela i nuklearne. Twierdzenie o jądrze Schwartza}
\label{sec:montel-nuklearne23}

W poprzednich rozdziałach zbudowaliśmy przestrzenie
$\mathcal E(\Omega)=C^\infty(\Omega)$,
$\mathcal D(\Omega)$ i $\mathcal S(\R^d)$ wraz z ich
przestrzeniami dualnymi. Obecnie wyjaśnimy, dlaczego w tych
przestrzeniach ograniczoność prowadzi do zwartości oraz w jakim
sensie operator liniowy może być reprezentowany przez jedno jądro
dystrybucyjne. Wszystkie przestrzenie liniowo-topologiczne w tym
rozdziale są nad $\K\in\{\R,\C\}$ i są Hausdorffa. W argumentach
wykorzystujących rozwinięcia Fouriera pracujemy nad $\C$;
wersje rzeczywiste wynikają przez kompleksyfikację.

\subsection{Własność Montela i jej znaczenie}

\begin{definicja}[Przestrzeń Montela]
\label{def:montel23}
Przestrzenią Montela nazywamy beczkową przestrzeń lokalnie wypukłą
$E$, w której każdy zbiór ograniczony ma zwarte domknięcie.
Jeżeli dodatkowo $E$ jest przestrzenią Frécheta, mówimy o
\emph{przestrzeni Frécheta--Montela}. W tej klasie beczkowość
wynika już z twierdzenia Baire'a, więc wystarczy wymagać
względnej zwartości zbiorów ograniczonych.
\end{definicja}

\begin{twierdzenie}[Klasyczne przestrzenie Montela]
\label{thm:montel-test23}
Dla każdego niepustego zbioru otwartego $\Omega\subseteq\R^d$,
$d\ge1$, przestrzenie $\mathcal E(\Omega)$,
$\mathcal S(\R^d)$ oraz $\mathcal D(\Omega)$, wyposażone
w swoje standardowe topologie lokalnie wypukłe, są przestrzeniami
Montela.
\end{twierdzenie}
\begin{proof}
Dla $\mathcal S$ względną zwartość ograniczonych zbiorów
udowodniono w Lemacie~\textup{22.7}; ponadto
$\mathcal S$ jest przestrzenią Frécheta.
Niech $(K_j)_{j\ge1}$ będzie wyczerpaniem $\Omega$ zbiorami
zwartymi, spełniającym $K_j\subseteq\operatorname{int}K_{j+1}$,
a każdy zwarty podzbiór $\Omega$ niech będzie zawarty w pewnym
$K_j$. Jeżeli $B\subseteq\mathcal E(\Omega)$ jest ograniczony,
to wszystkie pochodne funkcji z $B$ są jednostajnie ograniczone
na $K_j$. Ich równociągłość wynika z oszacowań pochodnych
rzędu o jeden wyższego, stosowanych na nieco większych
otoczeniach zwartych. Twierdzenie Arzelà--Ascoliego oraz
argument przekątniowy dają z dowolnego ciągu w $B$ podciąg,
którego wszystkie pochodne zbiegają jednostajnie na każdym
$K_j$. Spójność pochodnych z granicami uzasadnia, że
funkcja graniczna należy do $\mathcal E(\Omega)$.
Zupełność i metryzowalność tej przestrzeni kończą dowód.

Dla $\mathcal D(\Omega)$ korzystamy z udowodnionego w
Rozdziale~\textup{18} kryterium ograniczoności:
każdy ograniczony zbiór $B\subseteq\mathcal D(\Omega)$
jest zawarty w pewnym $\mathcal D_K(\Omega)$ i tam ograniczony.
Poprzedni argument Arzelà--Ascoliego, tym razem dla
funkcji o nośnikach w ustalonym $K$, daje jego względną
zwartość w $\mathcal D_K(\Omega)$, a więc również
w $\mathcal D(\Omega)$. Przestrzeń $\mathcal D(\Omega)$
jest beczkowa jako ścisła, przeliczalna granica induktywna
przestrzeni Frécheta (dowód tej własności podano wcześniej).
\end{proof}

\begin{wniosek}[Zbieżność na ograniczonych zbiorach testowych]
\label{cor:montel-dual-seq23}
Niech $E$ będzie przestrzenią Montela, a $(u_n)$ ciągiem w
$E'$ zbieżnym punktowo do $u\in E'$. Jeśli rodzina
$\{u_n:n\ge1\}$ jest równociągła, to $u_n\to u$
w topologii silnej $\beta(E',E)$. W szczególności tak jest,
gdy $E$ jest jedną z przestrzeni z
Twierdzenia~\ref{thm:montel-test23}.
\end{wniosek}
\begin{proof}
Dla ograniczonego $B\subseteq E$ zbiór $\overline B$
jest zwarty. Z równociągłości i skończonego pokrycia
$\overline B$ przesunięciami jednego otoczenia zera
wynika jednostajna na $B$ zbieżność funkcjonałów.
Wskazane przestrzenie są beczkowe, a zbieżność punktowa
ciągu implikuje jego punktową ograniczoność;
twierdzenie Banacha--Steinhausa daje równociągłość.
\end{proof}

\begin{przyklad}[Dlaczego nie wystarczy przestrzeń Banacha]
\label{ex:banach-not-montel23}
Jeżeli $X$ jest nieskończenie wymiarową przestrzenią Banacha,
to jej domknięta kula jednostkowa nie jest zwarta w topologii
normowej, na mocy lematu Riesza. Dlatego taka przestrzeń
nie jest przestrzenią Montela. Własność Montela
jest tu skutkiem szczególnie mocnej topologii przestrzeni
funkcji gładkich, a nie samej zupełności.
\end{przyklad}

\subsection{Operatory nuklearne w przestrzeniach Banacha}

\begin{definicja}[Operator nuklearny]
\label{def:nuclear-banach23}
Niech $X,Y$ będą przestrzeniami Banacha. Operator liniowy
$T:X\to Y$ nazywamy \emph{nuklearnym}, jeżeli istnieją
$\ell_j\in X'$ i $y_j\in Y$, $j\ge1$, takie, że
\[
 Tx=\sum_{j=1}^\infty \ell_j(x)y_j\quad(x\in X),
 \qquad \sum_{j=1}^\infty\|\ell_j\|\,\|y_j\|<\infty.
\]
Normą nuklearną $\|T\|_{\mathcal N}$ nazywamy infimum
ostatniej sumy po wszystkich takich przedstawieniach.
Przedstawienie dopuszcza także skończenie wiele składników.
\end{definicja}

\begin{twierdzenie}[Nuklearność implikuje zwartość]
\label{thm:nuclear-compact23}
Jeśli $X,Y$ są przestrzeniami Banacha i $T:X\to Y$
jest operatorem nuklearnym, to $T$ jest ograniczony i zwarty.
\end{twierdzenie}
\begin{proof}
Dla przedstawienia z Definicji~\ref{def:nuclear-banach23}
niech $T_Nx=\sum_{j=1}^N\ell_j(x)y_j$.
Każdy $T_N$ ma skończony rząd. Ponadto
$\|T-T_N\|\le\sum_{j>N}\|\ell_j\|\|y_j\|\to0$.
Przeto $T$ jest ograniczony i jest granicą operatorów
skończonego rzędu w normie operatorowej. Granica
operatorów zwartych w normie operatorowej jest zwarta:
obraz kuli jednostkowej można dla dowolnego $\varepsilon>0$
pokryć skończoną liczbą kul o promieniu $\varepsilon$
za pomocą dostatecznie dobrego przybliżenia $T_N$.
Zupełność $Y$ zapewnia zwartość domknięcia tego obrazu.
\end{proof}

\begin{przyklad}[Operator zwarty, który nie jest nuklearny]
\label{ex:compact-nonnuclear23}
Na przestrzeni Hilberta $\ell^2(\N)$ rozważmy
$De_n=n^{-1}e_n$. Operator $D$ jest zwarty, gdyż
$D_N=\sum_{n\le N}n^{-1}\langle\cdot,e_n\rangle e_n$
zbiegają do $D$ w normie operatorowej. Nie jest jednak
nuklearny. Istotnie, dla każdego przedstawienia nuklearnego
$D=\sum_j\ell_j(\cdot)y_j$ i każdego $N$ mamy
\[
\sum_{n=1}^N|\langle De_n,e_n\rangle|
\le\sum_j
\Big(\sum_{n=1}^N|\ell_j(e_n)|^2\Big)^{1/2}
\Big(\sum_{n=1}^N|\langle y_j,e_n\rangle|^2\Big)^{1/2}
\le\sum_j\|\ell_j\|\|y_j\|.
\]
Wykorzystaliśmy twierdzenie Riesza oraz nierówność Bessela.
Lewa strona jest równa $\sum_{n=1}^N1/n$ i rośnie
nieograniczenie. Zatem zwartość jest własnością ściśle
słabszą niż nuklearność.
\end{przyklad}

\subsection{Nuklearność przestrzeni lokalnie wypukłych}

\begin{definicja}[Przestrzeń nuklearna]
\label{def:nuclear-lcs23}
Niech $E$ będzie przestrzenią lokalnie wypukłą Hausdorffa,
$\mathcal P$ rodziną wszystkich ciągłych seminorm na $E$.
Dla $p\in\mathcal P$ definiujemy przestrzeń Banacha $E_p$
jako uzupełnienie przestrzeni unormowanej $E/\ker p$
w normie $\|x+\ker p\|=p(x)$.
Jeżeli $p\le q$, naturalne odwzorowanie $E/\ker q\to E/\ker p$
przedłuża się do ograniczonego operatora
$j_{q,p}:E_q\to E_p$.
Przestrzeń $E$ nazywamy \emph{nuklearną}, gdy dla każdego
$p\in\mathcal P$ istnieje $q\in\mathcal P$, $q\ge p$,
takie, że $j_{q,p}$ jest operatorem nuklearnym.
\end{definicja}

\begin{uwaga}[Ważne rozróżnienie]
\label{rem:nuclear-different23}
Nuklearność przestrzeni lokalnie wypukłej jest własnością
\emph{całej rodziny seminorm}, a nie twierdzeniem,
że identyczność na tej przestrzeni ma jedno nuklearne
przedstawienie. Jeśli przestrzeń Banacha $X$ jest nuklearna
w sensie Definicji~\ref{def:nuclear-lcs23}, to $X$
ma skończony wymiar: dla seminormy normowej wyższa
seminorma ciągła jest z nią równoważna, więc identyczność
$X\to X$ byłaby zwarta na mocy
Twierdzenia~\ref{thm:nuclear-compact23}.
\end{uwaga}

\begin{twierdzenie}[Nuklearność przestrzeni szybko malejących ciągów]
\label{thm:nuclear-s23}
Niech $I=\N_0^d$ i
\[
 s(I)=\bigl\{a=(a_\nu)_{\nu\in I}:
 p_k(a)=\sup_{\nu\in I}(1+|\nu|)^k|a_\nu|<\infty
 \text{ dla każdego }k\in\N_0\bigr\},
 \qquad |\nu|=\nu_1+\cdots+\nu_d.
\]
Przestrzeń $s(I)$ z topologią seminorm $(p_k)$ jest
nuklearną przestrzenią Frécheta--Montela.
\end{twierdzenie}
\begin{proof}
Zupełność i metryzowalność wynikają z kryterium Cauchy'ego
w seminormach: granice współrzędnych istnieją, a oszacowania
jednostajne przechodzą na granice. Niech $m>k+d$.
Uzupełnienie $s(I)$ względem $p_m$ jest izometrycznie
izomorficzne z ważoną przestrzenią $c_0(I)$.
Dla $a\in E_m$ odwzorowanie do $E_k$ ma przedstawienie
\[
 j_{m,k}a=\sum_{\nu\in I}a_\nu e_\nu,
 \qquad
 \sum_{\nu\in I}\|\mathrm{ev}_\nu\|_{E_m'}
 \|e_\nu\|_{E_k}
 =\sum_{\nu\in I}(1+|\nu|)^{k-m}<\infty.
\]
Ostatni szereg jest zbieżny, ponieważ liczba
$\nu\in\N_0^d$ z $|\nu|=r$ ma wzrost rzędu $r^{d-1}$.
To daje nuklearność. Jeśli $B\subseteq s(I)$ jest
ograniczony, z ograniczoności $p_{k+1}$ wynika
$\sup_{a\in B,\,|\nu|>R}(1+|\nu|)^k|a_\nu|\to0$.
Przybliżenia przez skończenie wiele współrzędnych są
więc jednostajne na $B$ w każdej seminormie.
Przez diagonalizację i zupełność otrzymujemy względną
zwartość $B$, czyli własność Montela.
\end{proof}

\begin{twierdzenie}[Nuklearność $C^\infty(\Omega)$]
\label{thm:nuclear-cinfty23}
Dla każdego niepustego zbioru otwartego $\Omega\subseteq\R^d$
przestrzeń $\mathcal E(\Omega)=C^\infty(\Omega)$
ze standardową topologią Frécheta jest nuklearna.
\end{twierdzenie}
\begin{proof}
Niech $p_{K,m}(f)=\max_{|\alpha|\le m}\sup_{x\in K}
|\partial^\alpha f(x)|$ będzie seminormą standardową,
gdzie $K\Subset\Omega$. Wybierzmy
$\chi\in\mathcal D(\Omega)$ równą $1$ w otoczeniu $K$
oraz zamknięty sześcian $Q$ z $\supp\chi\subset\operatorname{int}Q$.
Funkcję $\chi f$ przedłużamy okresowo na torus otrzymany
przez identyfikację przeciwległych ścian $Q$.
Niech $\{e_n:n\in\mathbb Z^d\}$ będzie jego układem
znormalizowanych eksponentów trygonometrycznych,
a $c_n(f)$ współczynnikami Fouriera funkcji $\chi f$.
Z całkowania przez części, stosując dostatecznie dużą
potęgę operatora $(1-\Delta)$, dostajemy dla każdego $r$
\[
 |c_n(f)|\le C_r(1+|n|)^{-2r}
 p_{L,2r}(f),\qquad L=\supp\chi,
\]
gdzie w stałej uwzględniono pochodne $\chi$ i rozmiar
sześcianu. Ponieważ $\chi=1$ na $K$, w przestrzeni $E_{p_{K,m}}$
zachodzi szereg
$f|_K=\sum_{n\in\mathbb Z^d}c_n(f)e_n|_K$.
Wektor $e_n|_K$ w $E_{p_{K,m}}$ rozumiemy jako obraz
funkcji $\chi e_n$, przedłużonej przez zero poza $Q$;
jest ona gładka na $\Omega$ i równa $e_n$ w otoczeniu $K$.
Norma tego wektora nie przekracza $C(1+|n|)^m$.
Wybierając $2r>m+d$, otrzymujemy
\[
 \sum_{n\in\mathbb Z^d}
 \|c_n\|_{(E_{p_{L,2r}})'}\,
 \|e_n|_K\|_{E_{p_{K,m}}}<\infty.
\]
Zatem naturalne odwzorowanie
$E_{p_{L,2r}}\to E_{p_{K,m}}$ jest nuklearne.
W przypadku rzeczywistym bierzemy część rzeczywistą
rozwinięcia: każdy składnik $c_n(f)e_n$ jest sumą
$\operatorname{Re}c_n(f)\operatorname{Re}e_n-
\operatorname{Im}c_n(f)\operatorname{Im}e_n$.
Z nierówności norm otrzymujemy rzeczywiste przedstawienie
nuklearne, którego suma norm jest co najwyżej dwukrotnie
większa od zespolonej.
Każda ciągła seminorma $p$ na $\mathcal E(\Omega)$
spełnia $p\le C p_{K,m}$ dla pewnego $C>0$, $K\Subset\Omega$
i $m\in\N_0$. Zwiększając seminormę źródłową
z poprzedniej konstrukcji i składając operator nuklearny
z ograniczonym przejściem $E_{p_{K,m}}\to E_p$,
uzyskujemy warunek z Definicji~\ref{def:nuclear-lcs23}.
\end{proof}

\begin{uwaga}[Hermite i przestrzeń Schwartza]
\label{rem:hermite-nuclear23}
Przestrzeń Schwartza $\mathcal S(\R^d)$ jest nuklearna:
współczynniki w bazie funkcji Hermite'a wyznaczają
izomorfizm liniowo-topologiczny
$\mathcal S(\R^d)\simeq s(\N_0^d)$.
Dokładne oszacowania i dowód tego izomorfizmu podano
w Uzupełnieniu~\ref{sec:hermite-supp23}.
Nuklearność $\mathcal D_K(\Omega)$ oraz
$\mathcal D(\Omega)$ zostanie poniżej wyprowadzona
z udowodnionych twierdzeń trwałości; nie jest to
formalna konsekwencja samej nuklearności $\mathcal S$.
\end{uwaga}

\subsection{Kryterium hilbertowskie i trwałość nuklearności}
\label{subsec:nuclear-permanence23}

W zastosowaniach do funkcji testowych potrzebujemy dwóch operacji:
przejścia do podprzestrzeni i do granicy induktywnej. Dowiedziemy
odpowiednich twierdzeń, wychodząc wyłącznie z
Definicji~\ref{def:nuclear-lcs23}. Dla ciągłej seminormy
\emph{hilbertowskiej} $h$ na przestrzeni lokalnie wypukłej $E$
istnieje dodatnia forma półtoraliniowa (w przypadku rzeczywistym:
symetryczna forma dwuliniowa), której kwadratem normy jest $h^2$.
Przez $H_h$ oznaczamy przestrzeń Hilberta otrzymaną przez
uzupełnienie $E/\ker h$. Nie wymagamy, by $E$ była metryzowalna.

\begin{lemat}[Złożenie operatorów Hilberta--Schmidta]
\label{lem:hs-compose-nuclear23}
Niech $H_2,H_1,H_0$ będą przestrzeniami Hilberta nad $\K$, a
$A:H_2\to H_1$ oraz $B:H_1\to H_0$ operatorami
Hilberta--Schmidta. Wówczas $BA:H_2\to H_0$ jest
operatorem nuklearnym w sensie
Definicji~\ref{def:nuclear-banach23}, przy czym
$\|BA\|_{\mathcal N}\le\|A\|_{\mathrm{HS}}\|B\|_{\mathrm{HS}}$.
\end{lemat}
\begin{proof}
Niech $(e_j)_{j\in J}$ będzie dowolną bazą ortonormalną $H_1$.
Dla każdego $x\in H_2$ zachodzi
$BAx=\sum_{j\in J}\langle Ax,e_j\rangle Be_j$,
gdyż rozwinięcie $Ax$ w bazie jest zbieżne w $H_1$,
a $B$ jest ograniczony. Ponadto
\[
 \sum_{j\in J}\|A^*e_j\|\,\|Be_j\|
 \le\Bigl(\sum_{j\in J}\|A^*e_j\|^2\Bigr)^{1/2}
     \Bigl(\sum_{j\in J}\|Be_j\|^2\Bigr)^{1/2}
 =\|A\|_{\mathrm{HS}}\|B\|_{\mathrm{HS}}.
\]
Przy sumowaniu po dowolnym zbiorze indeksów rozumiemy supremum
sum po jego skończonych podzbiorach. Skończoność powyższych sum
kwadratów zapewnia, że współczynniki niezerowe tworzą
co najwyżej przeliczalną rodzinę, nawet gdy $H_1$ nie jest
ośrodkowa. Funkcjonały
$x\mapsto\langle Ax,e_j\rangle$ mają normy $\|A^*e_j\|$,
a wektory $Be_j$ normy $\|Be_j\|$. Otrzymujemy zatem
przedstawienie nuklearne $BA$ z wymaganym oszacowaniem.
\end{proof}

\begin{twierdzenie}[Hilbertowskie kryterium nuklearności]
\label{thm:hilbert-nuclear23}
Niech $E$ będzie przestrzenią lokalnie wypukłą Hausdorffa
nad $\K$. Następujące warunki są równoważne:
\begin{enumerate}[label=\textup{(\roman*)},leftmargin=*]
\item $E$ jest nuklearna w sensie
Definicji~\ref{def:nuclear-lcs23};
\item każda ciągła seminorma $p$ na $E$ jest ograniczona
z góry przez pewną ciągłą seminormę hilbertowską $h$,
a dla każdej ciągłej seminormy hilbertowskiej $h$ istnieje
ciągła seminorma hilbertowska $r\ge h$, dla której
kanoniczny operator $H_r\to H_h$ jest operatorem
Hilberta--Schmidta.
\end{enumerate}
\end{twierdzenie}
\begin{proof}
\emph{(i)$\Rightarrow$(ii).} Ustalmy dowolną ciągłą seminormę $p$.
Z nuklearności wybieramy $q\ge p$ i przedstawienie operatora
$j_{q,p}:E_q\to E_p$ w postaci
\[
 j_{q,p}z=\sum_{j\ge1}a_j\ell_j(z)y_j,\qquad
 A=\sum_{j\ge1}|a_j|<\infty,\quad
 \|\ell_j\|\le1,\quad\|y_j\|\le1.
\]
Takie przedstawienie uzyskujemy przez włączenie norm funkcjonałów
i wektorów do współczynników. Jeśli $A=0$, to $p=0$ i wystarcza
seminorma zerowa. W przeciwnym razie kładziemy
$h(x)^2=A\sum_j|a_j|\,|\ell_j([x]_q)|^2$.
Jest to ciągła seminorma hilbertowska: powstaje z dodatniej
formy półtoraliniowej, a $h(x)\le Aq(x)$.
Z nierówności Cauchy'ego--Schwarza wynika
\[
 p(x)=\bigl\|j_{q,p}[x]_q\bigr\|
 \le\sum_j|a_j|\,|\ell_j([x]_q)|
 \le h(x).
\]

Weźmy teraz dowolną ciągłą seminormę hilbertowską $h$.
Jeżeli $h=0$, żądany operator jest zerowy. W przeciwnym razie
wybieramy $q\ge h$ i nuklearne przejście $j_{q,h}:E_q\to H_h$.
Niech $j_{q,h}z=\sum_ja_j\ell_j(z)y_j$ będzie jego
znormalizowanym przedstawieniem, z $A=\sum_j|a_j|>0$.
Definiujemy $r(x)^2=A\sum_j|a_j|\,|\ell_j([x]_q)|^2$.
Powyższy argument daje $r\ge h$ i ciągłość $r$.
Na $E/\ker r$ określamy
\[
 U([x]_r)=(\sqrt{A|a_j|}\,\ell_j([x]_q))_j\in\ell^2,
\qquad
 Ve_j=A^{-1/2}\sqrt{|a_j|}\,\operatorname{sgn}(a_j)y_j,
\]
gdzie $\operatorname{sgn}(a_j)=a_j/|a_j|$ dla $a_j\ne0$,
a $\operatorname{sgn}(0)=0$. Pierwsze odwzorowanie jest
izometrią i przedłuża się do $U:H_r\to\ell^2$.
Drugie jest operatorem Hilberta--Schmidta, ponieważ
$\sum_j\|Ve_j\|^2\le A^{-1}\sum_j|a_j|=1$.
Na gęstym obrazie $E$ mamy $VU=j_{r,h}$, a więc
$j_{r,h}:H_r\to H_h$ jest operatorem Hilberta--Schmidta.

\emph{(ii)$\Rightarrow$(i).} Dla dowolnej ciągłej seminormy
$p$ wybieramy $h\ge p$ hilbertowską, a następnie
hilbertowskie $r\ge h$ oraz $s\ge r$ tak, aby
$H_s\to H_r$ i $H_r\to H_h$ były operatorami
Hilberta--Schmidta. Ich złożenie $H_s\to H_h$ jest
nuklearne na mocy Lematu~\ref{lem:hs-compose-nuclear23}.
Nierówność $p\le h$ określa ograniczone odwzorowanie
$H_h\to E_p$. Złożenie operatora nuklearnego z ograniczonym
jest nuklearne, zatem $j_{s,p}:E_s=H_s\to E_p$ jest
nuklearny. To jest warunek z Definicji~\ref{def:nuclear-lcs23}.
\end{proof}

\begin{twierdzenie}[Trwałość względem podprzestrzeni i ilorazów]
\label{thm:nuclear-subquot23}
Niech $E$ będzie nuklearną przestrzenią lokalnie wypukłą
Hausdorffa. Każda podprzestrzeń liniowa $F\subseteq E$
z topologią indukowaną jest nuklearna. Jeśli $M\subseteq E$
jest domkniętą podprzestrzenią liniową, to iloraz $E/M$
z topologią ilorazową także jest nuklearny.
\end{twierdzenie}
\begin{proof}
Wykorzystamy dwie własności operatorów Hilberta--Schmidta.
Jeżeli $J:H_1\to H_0$ jest operatorem Hilberta--Schmidta,
to jego ograniczenie do domkniętej podprzestrzeni $R\subseteq H_1$
z wartościami w $\overline{J(R)}$ też ma tę własność:
wystarczy dopełnić bazę ortonormalną $R$ do bazy $H_1$.
Jeżeli natomiast $N\subseteq H_1$ i $L\subseteq H_0$
są domkniętymi podprzestrzeniami oraz $J(N)\subseteq L$,
to operator na ilorazach $H_1/N\to H_0/L$ jest
Hilberta--Schmidta. Istotnie, po identyfikacjach ilorazów
z dopełnieniami ortogonalnymi jest to
$P_{L^\perp}J|_{N^\perp}$.

Niech $F\subseteq E$ będzie dowolną podprzestrzenią.
Dla każdej ciągłej seminormy $p_F$ na $F$ istnieje
ciągła seminorma $p_E$ na $E$ taka, że
$p_F\le p_E|_F$. Aby to wykazać, wybieramy absolutnie
wypukłe otoczenie zera $U$ w $E$ z
$U\cap F\subseteq\{x\in F:p_F(x)<1\}$ i bierzemy
jego funkcjonał Minkowskiego (ewentualnie mnożąc go
przez stałą). Z kryterium hilbertowskiego
(Twierdzenie~\ref{thm:hilbert-nuclear23}) wybieramy
hilbertowskie $h\ge p_E$, $r\ge h$, $s\ge r$,
dla których $H_s\to H_r$ i $H_r\to H_h$ są
Hilberta--Schmidta. Dla $t\in\{h,r,s\}$ przestrzeń
$H_{t|_F}$ identyfikuje się z domknięciem obrazu $F$
w $H_t$. Wskazana własność ograniczeń zapewnia,
że $H_{s|_F}\to H_{r|_F}$ oraz
$H_{r|_F}\to H_{h|_F}$ są Hilberta--Schmidta.
Ich złożenie jest nuklearne. Ponieważ
$p_F\le h|_F$, istnieje ograniczone odwzorowanie
$H_{h|_F}\to F_{p_F}$; jego złożenie z poprzednim
przejściem dowodzi nuklearności $F$.

Niech teraz $M\subseteq E$ będzie domknięta i
$Q=E/M$ z odwzorowaniem ilorazowym $\pi:E\to Q$.
Dla ciągłej seminormy $p_Q$ na $Q$ zastosujmy kryterium
do $p_Q\circ\pi$ i wybierzmy hilbertowskie seminormy
$h\ge p_Q\circ\pi$, $r\ge h$, $s\ge r$,
których kolejne przejścia są Hilberta--Schmidta.
Dla $t\in\{h,r,s\}$ kładziemy
\[
 \bar t(\pi x)=\inf_{m\in M}t(x+m).
\]
Jest to ciągła seminorma hilbertowska na $Q$.
Jej ciągłość wynika z równości
$\{y:\bar t(y)<\varepsilon\}=
\pi(\{x:t(x)<\varepsilon\})$ i otwartości $\pi$.
Jeśli $N_t$ oznacza domknięcie obrazu $M$ w $H_t$,
to kanonicznie $H_{\bar t}\simeq H_t/N_t$, a
$\bar s\ge\bar r\ge\bar h\ge p_Q$.
Operator $H_s\to H_r$ odwzorowuje $N_s$ w $N_r$,
podobnie $H_r\to H_h$ odwzorowuje $N_r$ w $N_h$.
Z własności operatorów na ilorazach otrzymujemy
przejścia Hilberta--Schmidta
$H_{\bar s}\to H_{\bar r}\to H_{\bar h}$.
Ich złożenie jest nuklearne, zaś
$H_{\bar h}\to Q_{p_Q}$ jest ograniczone. Zatem
$j_{\bar s,p_Q}:Q_{\bar s}\to Q_{p_Q}$ jest nuklearny,
co kończy dowód.
\end{proof}

\begin{lemat}[Nuklearność przeliczalnej sumy prostej]
\label{lem:nuclear-direct-sum23}
Niech $E_1,E_2,\ldots$ będą nuklearnymi przestrzeniami
lokalnie wypukłymi Hausdorffa. Wówczas ich lokalnie wypukła
suma prosta $E=\bigoplus_{n\ge1}E_n$, z najsilniejszą
topologią lokalnie wypukłą, dla której wszystkie kanoniczne
włożenia $E_n\to E$ są ciągłe, jest nuklearna.
\end{lemat}
\begin{proof}
Niech $p$ będzie ciągłą seminormą na $E$,
$p_n=p|_{E_n}$. Dla każdego $n$ z nuklearności $E_n$
wybieramy ciągłą seminormę $q_n^0\ge p_n$, dla której
$j_{q_n^0,p_n}$ jest nuklearny. Zastępując $q_n^0$
przez $q_n=c_nq_n^0$ dla dostatecznie dużego $c_n\ge1$,
możemy zapewnić
$\|j_{q_n,p_n}\|_{\mathcal N}<2^{-n}$:
przeskalowanie normy w przestrzeni źródłowej przez
$c_n$ dzieli normę nuklearną przejścia przez $c_n$.
Definiujemy na algebraicznej sumie prostej
$q((x_n))=\sum_nq_n(x_n)$; jest to suma skończona.
Seminorma $q$ jest ciągła, ponieważ jej ograniczenie
do każdego $E_n$ jest ciągłe, a topologia sumy prostej
jest końcową topologią lokalnie wypukłą. Z nierówności
trójkąta wynika $p\le\sum_n p_n\le q$.

Uzupełnienie $E/\ker q$ identyfikuje się izometrycznie
z $\bigoplus_n^{\ell^1}(E_n)_{q_n}$: algebraiczne sumy
proste są gęste w tej sumie $\ell^1$.
Dla ustalonego $n$ kanoniczny operator
$(E_n)_{q_n}\to E_p$ jest złożeniem nuklearnego
$j_{q_n,p_n}$ i izometrycznego osadzenia
$(E_n)_{p_n}\to E_p$.
Wybieramy jego przedstawienie nuklearne, którego suma
norm składników jest mniejsza niż $2^{-n}$.
Każdy funkcjonał występujący w tym przedstawieniu
przedłużamy na sumę $\ell^1$ przez złożenie
z rzutem współrzędnym normy $1$.
Zebrane w jeden szereg przedstawienia są absolutnie
sumowalne w sensie norm nuklearnych, gdyż
$\sum_{n\ge1}2^{-n}<\infty$. Ich suma jest dokładnie
kanonicznym operatorem $E_q\to E_p$:
najpierw sprawdzamy tę równość dla wektorów o skończonym
nośniku, a następnie korzystamy z ciągłości i gęstości.
To dowodzi nuklearności $E$.
\end{proof}

\begin{twierdzenie}[Trwałość względem przeliczalnych granic induktywnych]
\label{thm:nuclear-indlim23}
Niech $(E_n,u_n)_{n\ge1}$ będzie ciągiem nuklearnych
przestrzeni lokalnie wypukłych Hausdorffa i ciągłych
odwzorowań liniowych $u_n:E_n\to E_{n+1}$.
Jeśli lokalnie wypukła granica induktywna
$E=\varinjlim_n E_n$ jest Hausdorffa, to $E$ jest nuklearna.
Nie zakładamy metryzowalności $E$ ani ścisłości granicy.
\end{twierdzenie}
\begin{proof}
Niech $S=\bigoplus_{n\ge1}E_n$ będzie lokalnie wypukłą
sumą prostą, a $\iota_n:E_n\to E$ odwzorowaniami
kanonicznymi. Definiujemy
$\Pi:S\to E$ przez
$\Pi((x_n))=\sum_n\iota_n(x_n)$; suma jest skończona.
Z końcowej własności topologii sumy prostej wynika
ciągłość $\Pi$, a z definicji granicy induktywnej
wynika jej surjektywność.

Topologia ilorazowa na $S/\ker\Pi$ ma dokładnie tę samą
własność uniwersalną co lokalnie wypukła topologia
granicy induktywnej: odwzorowanie liniowe $T:E\to G$
do dowolnej przestrzeni lokalnie wypukłej $G$ jest ciągłe
wtedy i tylko wtedy, gdy wszystkie $T\circ\iota_n$ są ciągłe.
Zatem $E\simeq S/\ker\Pi$ jako przestrzenie lokalnie
wypukłe. Ponieważ $E$ jest Hausdorffa,
$\ker\Pi=\Pi^{-1}(\{0\})$ jest domknięte w $S$.
Suma prosta $S$ jest nuklearna na mocy
Lematu~\ref{lem:nuclear-direct-sum23}, a iloraz
przez domkniętą podprzestrzeń jest nuklearny na mocy
Twierdzenia~\ref{thm:nuclear-subquot23}. Wynika stąd
nuklearność $E$. Ścisłość i metryzowalność granicy
nie zostały wykorzystane.
\end{proof}

\begin{twierdzenie}[Nuklearność $\mathcal D_K$ i $\mathcal D$]
\label{thm:nuclear-D23}
Niech $d\ge1$, $\Omega\subseteq\R^d$ będzie niepustym
zbiorem otwartym, a $K\Subset\Omega$ zbiorem zwartym.
Wówczas $\mathcal D_K(\Omega)$ ze standardową topologią
Frécheta i $\mathcal D(\Omega)$ ze standardową topologią
LF są nuklearnymi przestrzeniami lokalnie wypukłymi.
\end{twierdzenie}
\begin{proof}
Rozważmy najpierw $\mathcal D_K(\Omega)$ jako
podprzestrzeń $\mathcal E(\Omega)=C^\infty(\Omega)$.
Ponieważ $K$ jest domknięty w $\Omega$,
warunek $\supp f\subseteq K$ jest równoważny
$f(x)=0$ dla wszystkich $x\in\Omega\setminus K$.
A zatem
\[
 \mathcal D_K(\Omega)=
 \bigcap_{x\in\Omega\setminus K}\ker(\mathrm{ev}_x),
 \qquad \mathrm{ev}_x(f)=f(x).
\]
Każda ewaluacja jest ciągła w topologii $C^\infty$,
więc $\mathcal D_K(\Omega)$ jest podprzestrzenią
domkniętą. Ponadto wszystkie pochodne jej elementów
znikają poza $K$, stąd dla dowolnego $L\Subset\Omega$
\[
 p_{L,m}(f)\le p_{K,m}(f),\qquad
 f\in\mathcal D_K(\Omega).
\]
Ponieważ $K$ samo jest zwarte, rodzina $p_{K,m}$
jest zarazem rodziną ograniczeń seminorm z $C^\infty(\Omega)$.
Topologia indukowana jest więc dokładnie standardową
topologią Frécheta na $\mathcal D_K(\Omega)$.
Twierdzenia~\ref{thm:nuclear-cinfty23} i
\ref{thm:nuclear-subquot23} dowodzą jej nuklearności.

Wybierzmy wyczerpanie $\Omega$ zbiorami zwartymi
$K_1\subseteq\operatorname{int}K_2\subseteq\cdots\Subset\Omega$
takie, że każdy zwarty podzbiór $\Omega$ zawiera się
w pewnym $K_n$. Rozdział o funkcjach testowych
identyfikuje standardową topologię LF na
$\mathcal D(\Omega)$ z topologią lokalnie wypukłej
granicy induktywnej $\varinjlim_n\mathcal D_{K_n}(\Omega)$
i dowodzi, że ta granica jest Hausdorffa.
Jej nuklearność wynika teraz z
Twierdzenia~\ref{thm:nuclear-indlim23}.
\end{proof}

\begin{uwaga}[Dlaczego sam argument z $C^\infty$ nie wystarcza]
\label{rem:nuclear-D-summary23}
Nuklearność $\mathcal D_K$ jest konsekwencją trwałości
względem podprzestrzeni, ale nuklearność
$\mathcal D=\varinjlim\mathcal D_{K_n}$ wymaga osobnego
wyniku dotyczącego granic induktywnych. Powyższe dowody
ujawniają wszystkie użyte przejścia. Nie wynika z tego
równość topologii induktywnej tensorowej $i$ z topologią
rzutową $\pi$: dotyczą one innych własności uniwersalnych.
\end{uwaga}

\subsection{Twierdzenie o jądrze Schwartza}

\begin{definicja}[Forma dwuliniowa zadana jądrem]
\label{def:kernel-pair23}
Niech $\Omega\subseteq\R^{d_1}$ i
$\Lambda\subseteq\R^{d_2}$ będą niepustymi zbiorami otwartymi.
Dla $K\in\mathcal D'(\Omega\times\Lambda)$ definiujemy
formę na parach funkcji testowych przez
$b_K(f,g)=\langle K,f\otimes g\rangle$,
gdzie $(f\otimes g)(x,y)=f(x)g(y)$.
\end{definicja}

\begin{lemat}[Rozkład testu na szeregi tensorowe]
\label{lem:tensor-density23}
Niech $d_1,d_2\ge1$, a $\Omega\subseteq\R^{d_1}$ i
$\Lambda\subseteq\R^{d_2}$ będą niepustymi zbiorami otwartymi.
Dla każdego $\Phi\in\mathcal D(\Omega\times\Lambda)$
istnieje przedstawienie $\Phi$ jako sumy skończenie
wielu szeregów postaci
$\sum_{n\in\mathbb Z^{d_2}}f_n\otimes g_n$,
zbieżnych w $\mathcal D(\Omega\times\Lambda)$,
gdzie wszystkie $f_n$ mają nośniki w jednym zwartym
$K\Subset\Omega$, wszystkie $g_n$ w jednym zwartym
$L\Subset\Lambda$ i dla każdego $m,r\in\N_0$
\[
 \sum_n p_{K,m}(f_n)\,p_{L,r}(g_n)<\infty.
\]
W szczególności algebraiczna przestrzeń tensorowa
$\mathcal D(\Omega)\otimes\mathcal D(\Lambda)$
jest gęsta w $\mathcal D(\Omega\times\Lambda)$.
\end{lemat}
\begin{proof}
Nośnik $\Phi$ pokrywamy skończenie wieloma iloczynami
stosunkowo zwartych sześcianów $U_j\times V_j$
z $\overline{U_j}\subset\Omega$ i
$\overline{V_j}\subset\Lambda$, a następnie używamy
gładkiego rozkładu jedności. Wystarczy więc rozpatrzyć
funkcję o nośniku wewnątrz jednego takiego iloczynu.
W zmiennej $y$ przedłużamy ją okresowo na większym
sześcianie. Jej współczynniki Fouriera $f_n(x)$ są
funkcjami gładkimi o wspólnym zwartym nośniku w $x$.
Wielokrotne całkowanie przez części w $y$ daje
$p_{K,m}(f_n)\le C_{m,N}(1+|n|)^{-N}$ dla każdego $N$.
Jeżeli $\chi\in\mathcal D(\Lambda)$ jest równa $1$
w otoczeniu projekcji nośnika $\Phi$ na zmienną $y$,
możemy wziąć $g_n(y)=\chi(y)e^{in\cdot y}$
(z odpowiednim przeskalowaniem częstotliwości).
Wówczas $p_{L,r}(g_n)\le C_r(1+|n|)^r$.
Dobierając $N>r+d_2+1$ dostajemy absolutną zbieżność
we wszystkich seminormach $p_{K,m}\,p_{L,r}$, a więc
zbieżność szeregu tensorów na ustalonym etapie
$\mathcal D_{K\times L}(\Omega\times\Lambda)$:
pochodne $f_n(x)g_n(y)$ są szacowane przez iloczyny
odpowiednich seminorm. Wspólny nośnik jest zwarty.
Na obszarze, na którym $\chi=1$, suma jest pierwotną
funkcją; poza nośnikiem funkcji wyjściowej jej
periodyczne rozwinięcie jest zerowe. Wobec tego suma
równa się $\Phi$ na całej przestrzeni. Zbieżność
sum częściowych w jednej przestrzeni $\mathcal D_{K\times L}$
uzasadnia w szczególności wskazaną gęstość.
\end{proof}

\begin{lemat}[Ciągłość formy na ustalonych etapach]
\label{lem:separate-joint23}
Niech $E,F$ będą przestrzeniami Frécheta nad
$\K\in\{\R,\C\}$, a $b:E\times F\to\K$
formą dwuliniową ciągłą osobno w każdej zmiennej.
Wtedy $b$ jest łącznie ciągła.
\end{lemat}
\begin{proof}
Wybierzmy malejącą przeliczalną bazę
$(V_n)$ otoczeń zera w $F$, złożoną ze zbiorów
zrównoważonych. Niech
$A_n=\{x\in E: |b(x,y)|\le1\ \text{dla wszystkich }y\in V_n\}$.
Każdy $A_n$ jest domknięty, ponieważ $x\mapsto b(x,y)$
jest ciągły dla ustalonego $y$. Z ciągłości w drugiej
zmiennej wynika $E=\bigcup_n A_n$.
Z twierdzenia Baire'a pewien $A_n$ ma niepuste wnętrze;
ponieważ $A_n$ jest wypukły, zrównoważony i domknięty,
ma wówczas otoczenie zera we wnętrzu (sumując
symetryczne przesunięcia kuli mieszczącej się w $A_n$).
Stąd $|b(x,y)|\le1$ na iloczynie pewnego otoczenia
zera w $E$ oraz $V_n$; skalowanie daje łączną ciągłość.
\end{proof}

\begin{uwaga}[Ciągłość osobna i łączna]
\label{rem:kernel-separate23}
Dla formy $b:\mathcal D(\Omega)\times\mathcal D(\Lambda)\to\K$
ciągłość osobna nie musi implikować łącznej ciągłości na
całym iloczynie przestrzeni LF. Implikacja ta jest jednak
prawdziwa po ograniczeniu formy do
$\mathcal D_K(\Omega)\times\mathcal D_L(\Lambda)$,
gdyż obie przestrzenie są przestrzeniami Frécheta.
To rozróżnienie jest istotne w twierdzeniu o jądrze;
właściwą globalną topologią tensorową jest topologia
induktywna, omówiona w następnym rozdziale.
\end{uwaga}

\begin{twierdzenie}[Twierdzenie o jądrze Schwartza]
\label{thm:schwartz-kernel23}
Niech $\Omega\subseteq\R^{d_1}$ i
$\Lambda\subseteq\R^{d_2}$ będą niepustymi zbiorami
otwartymi. Dla każdego ciągłego operatora liniowego
\[
 A:\mathcal D(\Omega)\longrightarrow
 \mathcal D'(\Lambda)_\beta
\]
istnieje dokładnie jedna dystrybucja
$K_A\in\mathcal D'(\Omega\times\Lambda)$ taka, że
\[
 \langle Af,g\rangle=\langle K_A,f\otimes g\rangle
 \qquad(f\in\mathcal D(\Omega),\ g\in\mathcal D(\Lambda)).
\]
Odwrotnie, każda dystrybucja
$K\in\mathcal D'(\Omega\times\Lambda)$ określa w ten
sposób ciągły operator liniowy
$A_K:\mathcal D(\Omega)\to\mathcal D'(\Lambda)_\beta$.
Indeks $\beta$ oznacza silną topologię przestrzeni dualnej,
czyli jednostajną zbieżność na ograniczonych zbiorach
funkcji testowych.
\end{twierdzenie}
\begin{proof}
Niech $b(f,g)=\langle Af,g\rangle$. Ciągłość $A$
i ciągłość ewaluacji na pojedynczej funkcji testowej
zapewniają osobną ciągłość $b$. Na przestrzeniach
$\mathcal D_{K_x}(\Omega)$ i $\mathcal D_{K_y}(\Lambda)$
forma $b$ jest wspólnie ciągła na mocy
Lematu~\ref{lem:separate-joint23}; zatem istnieją
$C,m,r$ takie, że
$|b(f,g)|\le C p_{K_x,m}(f)p_{K_y,r}(g)$.
Zbudujemy jeden liniowy funkcjonał na całej przestrzeni
testowej, a następnie sprawdzimy jego niezależność
od wyborów pomocniczych. Wybierzmy lokalnie skończone
rozkłady jedności $(\eta_i)$ na $\Omega$ oraz
$(\theta_j)$ na $\Lambda$, których elementy mają zwarte
nośniki w sześcianach stosunkowo zwartych. Produkty
$\eta_i(x)\theta_j(y)$ tworzą rozkład jedności na
$\Omega\times\Lambda$. Dla danej
$\Phi\in\mathcal D(\Omega\times\Lambda)$ tylko skończenie
wiele funkcji $\Phi_{ij}=\eta_i\theta_j\Phi$ jest niezerowych.
Dla każdej z nich ustalamy raz na zawsze odcięcia i układ
Fouriera z Lematu~\ref{lem:tensor-density23}, otrzymując
$\Phi_{ij}=\sum_n f_{ijn}\otimes g_{ijn}$.
Definiujemy
\[
 \langle K_A,\Phi\rangle=\sum_{i,j}\sum_n b(f_{ijn},g_{ijn}),
\]
gdzie suma po $i,j$ jest skończona, a wszystkie szeregi
wewnętrzne są bezwzględnie zbieżne na mocy oszacowań
seminormowych i ciągłości $b$ na odpowiednich etapach.
Konstrukcja jest liniowa, ponieważ wybory odcięć i baz
Fouriera nie zależą od $\Phi$, a współczynniki Fouriera
zależą od niej liniowo.

Aby sprawdzić ciągłość, ustalmy $M\Subset\Omega\times\Lambda$.
Lokalna skończoność rozkładów jedności daje skończoną
liczbę par $(i,j)$, które mogą wystąpić dla
$\Phi\in\mathcal D_M$. Oszacowania współczynników
Fouriera z Lematu~\ref{lem:tensor-density23}, wraz
z regułą Leibniza dla ustalonych funkcji odcięcia,
zapewniają stałe $C_M>0$ i $r_M\in\N_0$ takie, że
\[
 |\langle K_A,\Phi\rangle|\le C_M p_{M,r_M}(\Phi)
 \qquad (\Phi\in\mathcal D_M(\Omega\times\Lambda)).
\]
Jest to ciągłość na każdym etapie, więc własność
uniwersalna granicy induktywnej daje
$K_A\in\mathcal D'(\Omega\times\Lambda)$.

Dla $\Phi=f\otimes g$ lokalne szeregi Fourierowskie
rekonstruują skończone sumy
$(\eta_i f)\otimes(\theta_j g)$. Przez ciągłość osobną
$b$ można przejść do granicy w ich sumach, otrzymując
$\langle K_A,f\otimes g\rangle=b(f,g)$.
Jeżeli zmienimy rozkłady jedności i odcięcia, otrzymamy
inny ciągły funkcjonał na całej przestrzeni $\mathcal D$,
ale o tych samych wartościach na tensorach elementarnych.
Ich liniowa powłoka jest gęsta na mocy
Lematu~\ref{lem:tensor-density23}; dwa funkcjonały muszą
więc być równe. To zarazem uzasadnia jednoznaczność,
bez zakładania gęstości tensorów o nośnikach w arbitralnym
zwartym zbiorze $M$.

W odwrotną stronę, dla $K\in\mathcal D'(\Omega\times\Lambda)$
forma $b_K(f,g)=\langle K,f\otimes g\rangle$ jest
łącznie ciągła na każdym
$\mathcal D_{K_x}(\Omega)\times\mathcal D_{K_y}(\Lambda)$,
ponieważ odwzorowanie $(f,g)\mapsto f\otimes g$
ma tam ciągłe oszacowania seminormowe.
Dla ustalonego $f$ forma $g\mapsto b_K(f,g)$
jest ciągłym funkcjonałem, więc wyznacza $A_Kf\in\mathcal D'(\Lambda)$.
Aby sprawdzić ciągłość w topologii silnej, niech
$B\subseteq\mathcal D(\Lambda)$ będzie ograniczony.
Z kryterium ograniczoności w $\mathcal D$ istnieje
$K_y\Subset\Lambda$, dla którego $B$ jest ograniczony
w $\mathcal D_{K_y}$. Na każdym etapie $\mathcal D_{K_x}$
mamy, dzięki wspólnej ciągłości formy,
\[
 \sup_{g\in B}|\langle A_Kf,g\rangle|
 \le C\,p_{K_x,m}(f)\sup_{g\in B}p_{K_y,r}(g).
\]
To dowodzi ciągłości $A_K$ na każdym etapie, a więc,
przez własność uniwersalną granicy induktywnej,
na całej $\mathcal D(\Omega)$. Jednoznaczność $K_A$
wynika ponownie z gęstości tensorów elementarnych.
\end{proof}

\begin{przyklad}[Jądro operatora o gładkim jądrze]
\label{ex:smooth-kernel23}
Jeżeli $k\in C^\infty(\Omega\times\Lambda)$, to dystrybucja
regularna $T_k$ jest jądrem operatora
$(Af)(y)=\int_\Omega k(x,y)f(x)\,\mathrm d\lambda_{d_1}(x)$.
Dla $g\in\mathcal D(\Lambda)$ mamy
\[
\langle Af,g\rangle
=\int_{\Lambda}\int_{\Omega}k(x,y)f(x)g(y)
\,\mathrm d\lambda_{d_1}(x)\,\mathrm d\lambda_{d_2}(y).
\]
Nie wymagamy, żeby $k$ miało zwarty nośnik, ponieważ
$f$ i $g$ są funkcjami testowymi.
\end{przyklad}

\begin{przyklad}[Jądro osobliwe na przekątnej]
\label{ex:diagonal-kernel23}
Dla $\Omega=\Lambda\subseteq\R^d$ operator naturalnego
włożenia $\mathcal D(\Omega)\hookrightarrow\mathcal D'(\Omega)$
ma jądro $\delta_\Delta$ zadane przez
$\langle\delta_\Delta,\Phi\rangle
=\int_\Omega\Phi(x,x)\,\mathrm d\lambda_d(x)$.
Jest ono dystrybucją skupioną na przekątnej
$\Delta=\{(x,x):x\in\Omega\}$. Nie istnieje
reprezentacja tego jądra przez funkcję z
$L^1_{\mathrm{loc}}(\Omega\times\Omega)$,
gdyż przekątna ma miarę Lebesgue'a zero, a
$\delta_\Delta\ne0$.
\end{przyklad}

\subsection{Uzupełnienie: funkcje Hermite'a i nuklearność przestrzeni Schwartza}
\label{sec:hermite-supp23}

\begin{lemat}[Rozwinięcie Hermite'a i seminormy]
\label{lem:hermite-schwartz23}
Dla $\nu\in\N_0^d$ niech $h_\nu$ będzie iloczynem
znormalizowanych jednowymiarowych funkcji Hermite'a.
Wówczas $(h_\nu)_{\nu\in\N_0^d}$ jest bazą ortonormalną
przestrzeni Hilberta $L^2(\R^d)$, a odwzorowanie
\[
 H:\mathcal S(\R^d)\to s(\N_0^d),\qquad
 H(f)=(\langle f,h_\nu\rangle_{L^2})_\nu
\]
jest izomorfizmem liniowo-topologicznym.
\end{lemat}
\begin{proof}
W jednym wymiarze przyjmijmy
$h_n(x)=(2^n n!\sqrt\pi)^{-1/2}
(-1)^n e^{x^2/2}(\mathrm d/\mathrm dx)^n e^{-x^2}$.
Wzory Rodriguesa i całkowanie przez części dają
ortogonalność i normy $\|h_n\|_2=1$.
Zupełność wynika z gęstości wielomianów pomnożonych
przez $e^{-x^2/2}$ w $L^2$: jeśli $f\in L^2$ jest
ortogonalna do wszystkich takich funkcji, to
$F(z)=\int_{\R} f(x)e^{-x^2/2}e^{zx}\,\mathrm d\lambda_1(x)$
jest funkcją całkowitą o wszystkich pochodnych w zerze
równych zero. Stąd $F=0$, a na osi urojonej
jednoznaczność transformacji Fouriera daje $f=0$.
Iloczyny tensorowe zapewniają zupełność w wymiarze $d$.

Niech $L=-\Delta+|x|^2$. Funkcje $h_\nu$ spełniają
$Lh_\nu=(2|\nu|+d)h_\nu$.
Wprowadzamy operatory
$a_j=\partial_j+x_j$ i $a_j^*=-\partial_j+x_j$.
Wzory rekurencyjne dają
$a_jh_\nu=\sqrt{2\nu_j}\,h_{\nu-e_j}$ oraz
$a_j^*h_\nu=\sqrt{2(\nu_j+1)}\,h_{\nu+e_j}$.
Ponieważ $x_j=(a_j+a_j^*)/2$ i
$\partial_j=(a_j-a_j^*)/2$, każdy iloczyn $r$
operatorów mnożenia przez współrzędną i różniczkowania
ma na skończonych kombinacjach $h_\nu$ normę $L^2$
kontrolowaną przez $C_r\|(1+L)^r f\|_2$.
Przez nierówność Sobolewa w wymiarze $d$ uzyskujemy
w konsekwencji, dla każdych $N,m$, pewne $r$ i $C$ takie, że
\[
 Q_{N,m}(f)\le C\|(1+L)^r f\|_2.
\]
W drugą stronę rozwinięcie $(1+L)^r$ jako skończonej
sumy iloczynów $x_j$ i $\partial_j$ oraz szybki zanik
$f$ dają $\|(1+L)^r f\|_2\le C'Q_{N',m'}(f)$
dla pewnych $N',m'$. Na skończonych kombinacjach
Hermite'a normy grafowe są równoważne normom
$\big(\sum_\nu(1+|\nu|)^{2r}|c_\nu|^2\big)^{1/2}$.
Z kolei te ważone normy $\ell^2$ wyznaczają tę samą
topologię co seminormy $p_k$ przestrzeni $s(\N_0^d)$:
jedna nierówność wynika z supremum, druga ze zbieżności
$\sum_\nu(1+|\nu|)^{-2r}$ dla $2r>d$.
W rezultacie odwzorowanie $H$ jest ciągłe z ciągłą
odwrotnością; dla szybkiego ciągu współczynników
jego sumy częściowe są Cauchy'ego we wszystkich normach
grafowych, a więc zbiegają w seminormach Schwartza.
\end{proof}

\begin{wniosek}[Nuklearność przestrzeni Schwartza]
\label{cor:nuclear-schwartz23}
Przestrzeń Schwartza $\mathcal S(\R^d)$ jest nuklearną
przestrzenią Frécheta--Montela.
\end{wniosek}
\begin{proof}
Przez Lemat~\ref{lem:hermite-schwartz23} jest izomorficzna
z $s(\N_0^d)$, a nuklearność i własność Montela są
niezmiennicze względem izomorfizmów liniowo-topologicznych.
Stosujemy Twierdzenie~\ref{thm:nuclear-s23}.
\end{proof}

\begin{uwaga}[Zakres twierdzenia o jądrze]
\label{rem:kernel-scope23}
Twierdzenie~\ref{thm:schwartz-kernel23} opisuje wszystkie
ciągłe operatory $\mathcal D(\Omega)\to\mathcal D'(\Lambda)_\beta$.
Nie oznacza ono, że jądro jest funkcją ani że operacja
składania dowolnych takich operatorów jest automatycznie
zdefiniowana. Twierdzenia trwałości nuklearności zostały udowodnione
powyżej; następny rozdział omawia topologiczne produkty tensorowe. Pełne rozwiązywanie równań
różniczkowych pozostaje w osobnym skrypcie.
\end{uwaga}


% Wersja źródłowa: audyt II, jednostka 24
\setcounter{section}{23}
\noindent\textbf{Wymagane wyniki.} Korzystamy z twierdzenia Hahna--Banacha, definicji nuklearności przez przejścia seminormowe, rozwinięć Fouriera i Hermitego oraz rozkładu funkcji testowych w szeregi tensorowe. Twierdzenia trwałości nuklearności dowiedziono w poprzednim rozdziale; tu korzystamy z nich w teorii produktów tensorowych.
\medskip
\section{Topologiczne produkty tensorowe przestrzeni lokalnie wypukłych}
\label{sec:tensor24}

W rachunku tensorowym przestrzeni skończenie wymiarowych pojęcie
$E\otimes F$ jest czysto algebraiczne. W analizie funkcjonalnej
potrzebujemy dodatkowo topologii: od niej zależy, które szeregi
tensorów są zbieżne, które formy dwuliniowe są ciągłe i czy
otrzymana przestrzeń ma interpretację jako przestrzeń funkcji.
W tym rozdziale $\K\in\{\R,\C\}$, a przestrzenie lokalnie wypukłe
są Hausdorffa, chyba że zaznaczono inaczej. Zapis $E'$ oznacza
przestrzeń dualną topologiczną $E$.

\subsection{Produkt tensorowy algebraiczny i własność uniwersalna}

\begin{definicja}[Produkt tensorowy algebraiczny]
\label{def:alg-tensor24}
Niech $E,F$ będą przestrzeniami liniowymi nad $\K$.
Ich \emph{produkt tensorowy algebraiczny} $E\otimes F$ jest
przestrzenią liniową wyposażoną w odwzorowanie dwuliniowe
$\otimes:E\times F\to E\otimes F$, $ (x,y)\mapsto x\otimes y$,
takie, że dla każdej przestrzeni liniowej $G$ i każdego
odwzorowania dwuliniowego $b:E\times F\to G$ istnieje dokładnie
jedno odwzorowanie liniowe $\widetilde b:E\otimes F\to G$
spełniające $\widetilde b(x\otimes y)=b(x,y)$ dla wszystkich
$x\in E$, $y\in F$.
\end{definicja}

\begin{stwierdzenie}[Istnienie i jednoznaczność]
\label{prop:alg-existence24}
Dla dowolnych przestrzeni liniowych $E,F$ produkt tensorowy
algebraiczny istnieje i jest jednoznaczny z dokładnością do
jednoznacznego izomorfizmu zachowującego tensory elementarne.
Każdy jego element jest skończoną sumą $\sum_{j=1}^N x_j\otimes y_j$.
\end{stwierdzenie}
\begin{proof}
Niech $V$ będzie wolną przestrzenią liniową o bazie
$[x,y]$, indeksowanej parami $(x,y)\in E\times F$.
Przez $N\subseteq V$ oznaczmy podprzestrzeń rozpiętą przez
relacje $[x+x',y]-[x,y]-[x',y]$,
$[x,y+y']-[x,y]-[x,y']$, $[ax,y]-a[x,y]$
oraz $[x,ay]-a[x,y]$ dla $a\in\K$.
Przyjmujemy $E\otimes F=V/N$ i $x\otimes y=[x,y]+N$.
Każde $b$ przedłuża się liniowo na $V$ i zeruje się na $N$,
więc przechodzi do jednoznacznego odwzorowania na ilorazie.
To daje własność uniwersalną. Jeśli $T$ i $T'$ są dwiema
realizacjami, własności uniwersalne dają odwzorowania liniowe
w obu kierunkach, których złożenia są identycznościami na
tensorach elementarnych, a więc na całych przestrzeniach.
\end{proof}

\begin{przyklad}[Tensory elementarne a funkcje]
\label{ex:separable-functions24}
Dla zbiorów $X,Y$ możemy odwzorować
$\mathcal F(X)\otimes\mathcal F(Y)$ w przestrzeń funkcji na
$X\times Y$ przez $f\otimes g\mapsto ((x,y)\mapsto f(x)g(y))$.
Nie każda funkcja dwóch zmiennych jest skończoną sumą funkcji
rozdzielonych; właśnie do dopuszczenia granic takich sum
potrzebujemy odpowiednich topologii i uzupełnień.
\end{przyklad}

\subsection{Topologia rzutowa}

\begin{definicja}[Seminorma rzutowa]
\label{def:pi-seminorm24}
Niech $E,F$ będą przestrzeniami lokalnie wypukłymi, a $p,q$
ciągłymi seminormami odpowiednio na $E$ i $F$. Dla
$u\in E\otimes F$ definiujemy
\[
 \pi_{p,q}(u)=\inf\Bigl\{\sum_{j=1}^N p(x_j)q(y_j):
 u=\sum_{j=1}^Nx_j\otimes y_j\Bigr\}.
\]
Topologię lokalnie wypukłą generowaną przez wszystkie
$\pi_{p,q}$ nazywamy \emph{topologią rzutową};
odpowiednią przestrzeń oznaczamy $E\otimes_\pi F$,
a jej uzupełnienie $E\widehat\otimes_\pi F$.
\end{definicja}

\begin{lemat}[Własności seminorm rzutowych]
\label{lem:pi-properties24}
Każda funkcja $\pi_{p,q}$ jest seminormą. Dla $x\in E$,
$y\in F$ zachodzi $\pi_{p,q}(x\otimes y)=p(x)q(y)$.
Odwzorowanie kanoniczne $E\times F\to E\otimes_\pi F$
jest ciągłe i dwuliniowe. Jeśli $E,F$ są Hausdorffa,
to również $E\otimes_\pi F$ jest Hausdorffa.
\end{lemat}
\begin{proof}
Jednorodność i nierówność trójkąta wynikają ze skalowania
oraz łączenia przedstawień tensorów. Dla tensora elementarnego
jednoskładnikowe przedstawienie daje nierówność $\le$.
Gdy $p(x)q(y)>0$, twierdzenie Hahna--Banacha daje
$\ell\in E'$ i $m\in F'$ z $|\ell(z)|\le p(z)$,
$|m(w)|\le q(w)$ oraz $\ell(x)=p(x)$, $m(y)=q(y)$
(po odpowiednim wyborze fazy w przypadku zespolonym).
Dla każdego przedstawienia $x\otimes y=\sum_jx_j\otimes y_j$
mamy $p(x)q(y)=|\sum_j\ell(x_j)m(y_j)|
\le\sum_jp(x_j)q(y_j)$. Gdy któryś czynnik zanika,
równość jest natychmiastowa. Ponieważ seminormy
$\pi_{p,q}(x\otimes y)$ są iloczynami $p(x)q(y)$,
ciągłość odwzorowania dwuliniowego wynika z ciągłości
seminorm $p,q$. Dla $u\ne0$ istnieje para funkcjonałów
liniowych $\ell\in E'$, $m\in F'$ z
$(\ell\otimes m)(u)\ne0$: ograniczamy się do skończenie
wymiarowych przestrzeni rozpiętych przez współczynniki $u$
i przedłużamy funkcjonały współrzędnych twierdzeniem
Hahna--Banacha. Wybieramy $p=|\ell|$, $q=|m|$;
wtedy $\pi_{p,q}(u)\ge|(\ell\otimes m)(u)|>0$.
\end{proof}

\begin{twierdzenie}[Własność uniwersalna topologii rzutowej]
\label{thm:pi-universal24}
Niech $E,F,G$ będą przestrzeniami lokalnie wypukłymi
Hausdorffa. Przyporządkowanie $b\mapsto\widetilde b$ jest bijekcją
między ciągłymi odwzorowaniami dwuliniowymi
$b:E\times F\to G$ i ciągłymi odwzorowaniami liniowymi
$\widetilde b:E\otimes_\pi F\to G$. Jeśli $G$ jest
zupełna, każde takie $\widetilde b$ przedłuża się jednoznacznie
do $E\widehat\otimes_\pi F$.
\end{twierdzenie}
\begin{proof}
Niech $s$ będzie ciągłą seminormą na $G$ i niech $b$
będzie ciągłe. Istnieją zrównoważone otoczenia zera
$U\subset E$, $V\subset F$ takie, że $s(b(x,y))\le1$
na $U\times V$. Po zastąpieniu ich mniejszymi otwartymi
otoczeniami absolutnie wypukłymi i użyciu funkcjonałów
Minkowskiego otrzymujemy ciągłe seminormy $p,q$,
dla których $s(b(x,y))\le p(x)q(y)$ dla wszystkich
$x,y$. Stąd dla $u=\sum_jx_j\otimes y_j$ mamy
$s(\widetilde b(u))\le\sum_jp(x_j)q(y_j)$,
a po wzięciu infimum $s(\widetilde b(u))\le\pi_{p,q}(u)$.
To jest ciągłość $\widetilde b$. Odwrotnie, złożenie
ciągłego $\widetilde b$ z ciągłym odwzorowaniem
$(x,y)\mapsto x\otimes y$ jest ciągłe.
Jednoznaczność wynika z rozpinania produktu tensorowego
przez tensory elementarne. Przedłużenie do uzupełnienia
wynika z ciągłości jednostajnej odwzorowania liniowego
oraz zupełności $G$.
\end{proof}

\begin{wniosek}[Forma dwuliniowa i przestrzeń dualna]
\label{cor:pi-dual24}
Dla lokalnie wypukłych przestrzeni Hausdorffa $E,F$
przestrzeń liniowa $(E\otimes_\pi F)'$ jest w naturalnej
bijekcji z przestrzenią wszystkich łącznie ciągłych form
dwuliniowych $b:E\times F\to\K$.
\end{wniosek}

\subsection{Topologia injektywna}

\begin{definicja}[Seminorma injektywna]
\label{def:eps-seminorm24}
Dla ciągłych seminorm $p$ na $E$ i $q$ na $F$ definiujemy
\[
 \varepsilon_{p,q}(u)=
 \sup\Bigl\{\big|\sum_{j=1}^N\ell(x_j)m(y_j)\big|:
 \substack{|\ell(x)|\le p(x)\ (x\in E),\\
             |m(y)|\le q(y)\ (y\in F)}\Bigr\},
 \quad u=\sum_{j=1}^N x_j\otimes y_j.
\]
Topologię wyznaczoną przez te seminormy nazywamy
\emph{topologią injektywną}; piszemy
$E\otimes_\varepsilon F$ i
$E\widehat\otimes_\varepsilon F$ dla jej uzupełnienia.
Występujące w definicji zbiory funkcjonałów są zbiorami
polarnymi kul jednostkowych seminorm $p$ i $q$.
\end{definicja}

\begin{twierdzenie}[Porównanie topologii]
\label{thm:eps-pi24}
Dla dowolnych lokalnie wypukłych przestrzeni Hausdorffa
$E,F$ mamy $\varepsilon_{p,q}\le\pi_{p,q}$ dla każdej
pary ciągłych seminorm $p,q$. Obie są seminormami i
$\varepsilon_{p,q}(x\otimes y)=p(x)q(y)$.
Topologia injektywna jest nie silniejsza od rzutowej;
kanoniczne odwzorowanie
$E\otimes_\pi F\to E\otimes_\varepsilon F$ jest ciągłe.
\end{twierdzenie}
\begin{proof}
Dla każdego przedstawienia $u=\sum_jx_j\otimes y_j$
i każdej dopuszczalnej pary $(\ell,m)$ zachodzi
$|\sum_j\ell(x_j)m(y_j)|\le\sum_jp(x_j)q(y_j)$.
Bierzemy supremum, a następnie infimum. Własności
seminormy są bezpośrednią konsekwencją definicji,
a równość dla tensora elementarnego wynika z osobnego
zastosowania twierdzenia Hahna--Banacha w obu czynnikach.
\end{proof}

\begin{twierdzenie}[Dziedziczność normy injektywnej]
\label{thm:eps-injective24}
Niech $X_0\subseteq X$ i $Y_0\subseteq Y$ będą izometrycznie
włożonymi podprzestrzeniami przestrzeni unormowanych.
Dla każdego $u\in X_0\otimes Y_0$ jego norma injektywna
liczona w $X_0\otimes_\varepsilon Y_0$ jest równa normie
liczonej po utożsamieniu z elementem $X\otimes_\varepsilon Y$.
\end{twierdzenie}
\begin{proof}
Ograniczenie funkcjonałów z kul jednostkowych $X'$ i $Y'$
daje jedną z nierówności. Odwrotnie, twierdzenie
Hahna--Banacha przedłuża dowolne $\ell\in X_0'$
i $m\in Y_0'$ z zachowaniem norm. Otrzymujemy więc
identyczny zbiór dopuszczalnych wartości na tensorze $u$.
\end{proof}

\begin{przyklad}[Przypadek skończenie wymiarowy]
\label{ex:finite-dim24}
Dla $E=\K^n$ z dowolną normą oraz przestrzeni Banacha $F$
obydwa uzupełnione produkty tensorowe są liniowo-topologicznie
izomorficzne z $F^n$. Rzeczywiście, przy wyborze bazy
$(e_i)$ każdy tensor ma postać $\sum_{i=1}^ne_i\otimes y_i$,
a każda z rozważanych topologii jest równoważna topologii
produktu skończenie wielu kopii $F$: projekcje współrzędne
są ciągłe, a skończenie wiele tensorów elementarnych
odtwarza cały element.
\end{przyklad}

\begin{twierdzenie}[Rozbieżność norm w nieskończonym wymiarze]
\label{thm:hilbert-norms24}
Niech $H$ będzie nieskończenie wymiarową przestrzenią Hilberta
nad $\K$ i $(e_n)$ ortonormalnym ciągiem w $H$. Dla
$u_N=\sum_{n=1}^N e_n\otimes e_n\in H\otimes H$ zachodzi
\[ \varepsilon(u_N)=1,\qquad \pi(u_N)=N. \]
W szczególności na $H\otimes H$ topologie rzutowa i
injektywna są różne.
\end{twierdzenie}
\begin{proof}
Dla $\ell,m$ z kul jednostkowych $H'$ nierówność
Cauchy'ego--Schwarza i Bessela dają
$|\sum_{n=1}^N\ell(e_n)m(e_n)|\le1$. Wybór funkcjonałów
współrzędnych dla $e_1$ pokazuje równość. Z definicji
$\pi(u_N)\le N$. Aby uzyskać nierówność przeciwną,
rozważmy formę dwuliniową
$b(x,y)=\sum_{n=1}^N\langle x,e_n\rangle\langle y,e_n\rangle$.
Przy konwencji liniowości iloczynu skalarnego w pierwszym
argumencie jest ona dwuliniowa. Na mocy Bessela i
nierówności Cauchy'ego--Schwarza ma normę nie większą niż $1$,
a $\widetilde b(u_N)=N$. Dla każdej reprezentacji $u_N$ suma
$\sum_j\|x_j\|\|y_j\|$ jest więc nie mniejsza niż $N$.
\end{proof}

\subsection{Nuklearność i równość topologii tensorowych}

\begin{twierdzenie}[Nuklearność usuwa różnicę między $\pi$ i $\varepsilon$]
\label{thm:nuclear-pi-eps24}
Niech $E$ będzie nuklearną przestrzenią lokalnie wypukłą
Hausdorffa w sensie definicji nuklearności z \hyperref[sec:montel-nuklearne23]{rozdziału~23}
(Definicja~\textup{23.8}; kryterium hilbertowskie: Twierdzenie~\textup{23.14}), a $F$
dowolną przestrzenią lokalnie wypukłą Hausdorffa.
Wtedy topologie rzutowa i injektywna na
$E\otimes F$ są równe. Ich uzupełnienia są więc
kanonicznie izomorficzne:
$E\widehat\otimes_\pi F\simeq E\widehat\otimes_\varepsilon F$.
\end{twierdzenie}
\begin{proof}
Ustalmy seminormę rzutową $\pi_{p,r}$, gdzie $p$ jest
ciągłą seminormą na $E$, a $r$ na $F$. Z nuklearności
istnieje ciągła seminorma $q\ge p$, taka że kanoniczny
operator $j_{q,p}:E_q\to E_p$ między przestrzeniami
Banacha jest nuklearny. Wybierzmy jego przedstawienie
$j_{q,p}(z)=\sum_{k\ge1}\ell_k(z)a_k$ z
$\sum_k\|\ell_k\|\,\|a_k\|=C<\infty$.
Jeżeli $u=\sum_{j=1}^N x_j\otimes y_j$, to w
$E_p\widehat\otimes_\pi F_r$ jego obraz jest granicą
szeregu
\[
 \sum_{k\ge1}a_k\otimes
 \Bigl(\sum_{j=1}^N\ell_k([x_j]_q)[y_j]_r\Bigr).
\]
Definicja seminormy injektywnej i nierówność
$|\ell_k(z)|\le\|\ell_k\|\,\|z\|_{E_q}$ dają
\[
 \left\|\sum_j\ell_k([x_j]_q)[y_j]_r\right\|_{F_r}
 \le\|\ell_k\|\,\varepsilon_{q,r}(u).
\]
Stąd $\pi_{p,r}(u)\le C\varepsilon_{q,r}(u)$.
(To samo oszacowanie można najpierw przeprowadzić na
unormowanych ilorazach wyznaczonych przez seminormy,
a następnie przejść do ich uzupełnień.)
Otrzymaliśmy, że każda seminorma rzutowa jest
kontrolowana przez seminormę injektywną (z potencjalnie
większą seminormą w pierwszym czynniku). Odwrotna
ciągłość wynika z Twierdzenia~\ref{thm:eps-pi24}.
\end{proof}

\begin{uwaga}[Granica stosowalności]
\label{rem:nuclear-characterization24}
Twierdzenie~\ref{thm:nuclear-pi-eps24} wyjaśnia,
dlaczego dla przestrzeni nuklearnych można niekiedy
pisać $E\widehat\otimes F$ bez wskazywania, czy
użyto topologii $\pi$, czy $\varepsilon$.
Nie oznacza to równości wszystkich możliwych topologii
tensorowych: topologia induktywna $i$, która pojawi się
poniżej, może być silniejsza nawet dla przestrzeni nuklearnych.
Pełne twierdzenie odwrotne Grothendiecka (charakterystyka
nuklearności przez równość $\pi$ i $\varepsilon$ dla wszystkich
czynników) wymaga dodatkowej teorii i nie jest używane w dowodach.
\end{uwaga}

\begin{wniosek}[Przestrzenie Banacha nuklearne tylko w skończonym wymiarze]
\label{cor:banach-finite-nuclear24}
Przestrzeń Banacha $X$ jest nuklearna jako przestrzeń
lokalnie wypukła wtedy i tylko wtedy, gdy ma skończony wymiar.
\end{wniosek}
\begin{proof}
Konieczność wykazano w poprzednim rozdziale, korzystając ze
zwartości operatorów nuklearnych i lematu Riesza.
Jeśli $X$ ma skończony wymiar, identyczność na $X$ jest
operatorem skończonego rzędu, a zatem nuklearnym.
Każda ciągła seminorma jest kontrolowana przez wielokrotność
normy, więc definicja nuklearności jest spełniona.
\end{proof}


\begin{stwierdzenie}[Nuklearność produktów]
\label{prop:nuclear-products24}
Niech $(E_\alpha)_{\alpha\in A}$ będzie dowolną rodziną
nuklearnych przestrzeni lokalnie wypukłych Hausdorffa.
Wówczas przestrzeń $E=\prod_{\alpha\in A}E_\alpha$
z topologią produktu jest nuklearna.
\end{stwierdzenie}
\begin{proof}
Niech $p$ będzie ciągłą seminormą na $E$.
Z definicji topologii produktu istnieje skończony zbiór
$J\subseteq A$ taki, że $p(x)=0$, gdy $x_\alpha=0$
dla każdego $\alpha\in J$. Istotnie, w otoczeniu zera
$\{x:p(x)<1\}$ mieści się otoczenie cylindryczne
wyznaczone przez skończenie wiele współrzędnych.
Wektor zerujący się na nich można pomnożyć przez dowolny
skalar bez opuszczania cylindra, więc jego seminorma
musi być zerowa. W konsekwencji $p$ zależy tylko od
współrzędnych z $J$ i jest kontrolowana przez sumę
$\sum_{\alpha\in J}p_\alpha(x_\alpha)$ dla pewnych
ciągłych seminorm $p_\alpha$ na $E_\alpha$.
Dla każdego $\alpha\in J$ wybieramy $q_\alpha\ge p_\alpha$
takie, że $j_{q_\alpha,p_\alpha}$ jest nuklearny,
i przyjmujemy $q(x)=\sum_{\alpha\in J}q_\alpha(x_\alpha)$.
Operator z uzupełnienia $E_q$ do uzupełnienia
$E_{\sum p_\alpha}$ jest skończoną sumą operatorów
nuklearnych działających na poszczególnych współrzędnych,
a więc jest nuklearny. Operator do $E_p$ jest jego
złożeniem z ograniczonym odwzorowaniem wynikającym z
$p\le\sum p_\alpha$; także jest nuklearny.
Ponadto $q\ge p$, co kończy dowód.
\end{proof}

\begin{uwaga}[Trwałość względem podprzestrzeni i granic induktywnych]
\label{rem:nuclear-permanence24}
Trwałość nuklearności względem podprzestrzeni i ilorazów
została udowodniona w
Twierdzeniu~\textup{23.15}, a względem
przeliczalnych granic induktywnych --- w
Twierdzeniu~\textup{23.17}.
Dlatego nuklearność $\mathcal D_K(\Omega)$ i
$\mathcal D(\Omega)$ jest już wynikiem
Twierdzenia~\textup{23.18}, a nie osobnym
nieudowodnionym założeniem. Wyniki te nie uprawniają
do utożsamienia topologii $i$ i $\pi$ na
$\mathcal D(\Omega)\otimes\mathcal D(\Lambda)$:
porównują one odmienne własności uniwersalne.
\end{uwaga}

\subsection{Przykłady produktów tensorowych przestrzeni funkcyjnych}

\begin{twierdzenie}[Produkt tensorowy $\ell^1$]
\label{thm:l1-tensor24}
Niech $F$ będzie przestrzenią Banacha. Wówczas istnieje
kanoniczny izometryczny izomorfizm
$\ell^1(\N)\widehat\otimes_\pi F\simeq\ell^1(\N;F)$,
gdzie prawa strona składa się z ciągów $(y_n)$ spełniających
$\sum_n\|y_n\|_F<\infty$, z normą tej sumy.
\end{twierdzenie}
\begin{proof}
Najpierw definiujemy liniowe odwzorowanie
$J:\ell^1\otimes F\to\ell^1(\N;F)$ wzorem
\[
 J\Bigl(\sum_{j=1}^M a^{(j)}\otimes y_j\Bigr)
 =\Bigl(\sum_{j=1}^M a_n^{(j)}y_j\Bigr)_{n\ge1}.
\]
Jest ono dobrze określone dzięki własności uniwersalnej
produktu algebraicznego. Dla dowolnego przedstawienia
$u=\sum_j a^{(j)}\otimes y_j$ zachodzi
$\|Ju\|_{\ell^1(F)}\le\sum_j\|a^{(j)}\|_1\|y_j\|$,
a więc $\|Ju\|_{\ell^1(F)}\le\pi(u)$.

Jeżeli $u=\sum_{n=1}^N e_n\otimes y_n$, to nierówność
przeciwna wynika z twierdzenia Hahna--Banacha.
Wybieramy $m_n\in F'$ z $\|m_n\|\le1$ i
$m_n(y_n)=\|y_n\|$ (dla $y_n=0$ bierzemy $m_n=0$).
Forma $b(a,y)=\sum_{n=1}^Na_nm_n(y)$ ma normę
nie większą niż $1$ na $\ell^1\times F$ oraz
$\widetilde b(u)=\sum_n\|y_n\|$. Zatem
$\pi(u)=\sum_n\|y_n\|=\|Ju\|_{\ell^1(F)}$
dla tensorów o skończonym nośniku współrzędnych.

Pozostaje uzasadnić, że wyliczenie normy obejmuje
\emph{dowolny} tensor algebraiczny. Jeśli
$u=\sum_{j=1}^M a^{(j)}\otimes y_j$, niech $P_N$
obcina ciąg po $N$ współrzędnych i połóżmy
$u_N=\sum_j(P_Na^{(j)})\otimes y_j$. Wtedy
\[
 \pi(u-u_N)\le\sum_{j=1}^M
 \|a^{(j)}-P_Na^{(j)}\|_1\,\|y_j\|\longrightarrow0.
\]
Ponadto $Ju_N\to Ju$ w $\ell^1(\N;F)$. Równość
norm dla $u_N$, wraz z ciągłością norm obu przestrzeni,
daje $\pi(u)=\|Ju\|_{\ell^1(F)}$.
Odwzorowanie $J$ jest więc izometrią o obrazie
zawierającym wszystkie ciągi o skończonym nośniku.
Przedłuża się do izometrii z zupełnego
$\ell^1\widehat\otimes_\pi F$ na $\ell^1(\N;F)$:
jej obraz jest domknięty i zarazem gęsty.
\end{proof}

\begin{twierdzenie}[Jądra gładkie i szybko malejące]
\label{thm:smooth-schwartz-tensor24}
Dla niepustych zbiorów otwartych $\Omega\subseteq\R^{d_1}$,
$\Lambda\subseteq\R^{d_2}$ oraz $d_1,d_2\ge1$ mamy
naturalne izomorfizmy liniowo-topologiczne
\[
 \mathcal E(\Omega)\widehat\otimes_\pi\mathcal E(\Lambda)
 \simeq \mathcal E(\Omega\times\Lambda),\qquad
 \mathcal S(\R^{d_1})\widehat\otimes_\pi\mathcal S(\R^{d_2})
 \simeq \mathcal S(\R^{d_1+d_2}).
\]
W obu wzorach można zastąpić $\pi$ przez $\varepsilon$.
\end{twierdzenie}
\begin{proof}
W obu przypadkach odwzorowanie kanoniczne na tensorach
elementarnych przypisuje $f\otimes g$ funkcję
$(x,y)\mapsto f(x)g(y)$.

Dla przestrzeni Schwartza stosujemy izomorfizmy z
przestrzeniami szybko malejących ciągów, udowodnione
w Lemacie~\textup{23.27}. W tym modelu
odwzorowanie kanoniczne wysyła $(a_i)\otimes(b_j)$
na macierz $(a_i b_j)_{i,j}$. Dla indeksów wielowymiarowych
$I=\N_0^{d_1}$, $J=\N_0^{d_2}$ oznaczmy
$w(i,j)=1+|i|+|j|$. Seminormy
$P_m(c)=\sup_{i,j}w(i,j)^m|c_{i,j}|$ wyznaczają
topologię przestrzeni $s(I\times J)$. Na tensorach
skończonych $P_m(c)\le C_m\pi_{p_m,p_m}(u)$.
W drugą stronę, rozwinięcie według współrzędnych daje
$u=\sum_{i,j}c_{ij}e_i\otimes e_j$ zbieżnie w
odpowiednim uzupełnieniu rzutowym, gdyż dla
$M>2m+d_1+d_2$ zachodzi
\[
 \sum_{i,j}|c_{ij}|(1+|i|)^m(1+|j|)^m
 \le C_{m,M}P_M(c).
\]
To jednocześnie wykazuje ciągłość odwrotności i gęstość
tensorów skończonego rzędu.

Dla $\mathcal E$ najpierw zauważmy, że odwzorowanie
$u\mapsto\Phi_u$, $\Phi_{f\otimes g}(x,y)=f(x)g(y)$,
jest ciągłe z topologii rzutowej do topologii
$\mathcal E(\Omega\times\Lambda)$. Istotnie, dla dowolnych
$K\Subset\Omega$, $L\Subset\Lambda$ i $m\in\N_0$ mamy
$p_{K\times L,m}(\Phi_u)\le C_m\pi_{p_{K,m},p_{L,m}}(u)$.

Aby skonstruować odwrotność, wybieramy lokalnie skończone
rozkłady jedności $(\eta_i)_{i\ge1}$ na $\Omega$
i $(\theta_j)_{j\ge1}$ na $\Lambda$, podporządkowane
stosunkowo zwartym sześcianom współrzędnych.
Każda funkcja $\Phi\in\mathcal E(\Omega\times\Lambda)$
jest lokalnie skończoną sumą funkcji
$\Phi_{ij}(x,y)=\eta_i(x)\theta_j(y)\Phi(x,y)$,
które mają nośniki zawarte w ustalonych zwartych
iloczynach $K_i\times L_j$. Każdą $\Phi_{ij}$
przedłużamy przez zero na nieco większy iloczyn
sześcianów, a następnie okresowo. Jej podwójny szereg
Fouriera, po pomnożeniu przez funkcje odcięcia w obu
zmiennych, daje przedstawienie
$\Phi_{ij}=\sum_{k,l}c_{kl}^{ij}f_{ik}\otimes g_{jl}$.
Dla wszystkich $r,s\in\N_0$ istnieją $M$ i $C$ takie, że
\[
 \sum_{k,l}|c_{kl}^{ij}|p_{K,r}(f_{ik})p_{L,s}(g_{jl})
 \le C\,p_{K'_{ij}\times L'_{ij},M}(\Phi),
\]
gdzie $K'_{ij}\Subset\Omega$ i $L'_{ij}\Subset\Lambda$
są nieco większe od nośników odcięć. Wynika to z
wielokrotnego całkowania przez części w zmiennych
periodycznych i sumowalności
$\sum_{k,l}(1+|k|+|l|)^{-q}$ dla $q>d_1+d_2$.

W ustalonej seminormie rzutowej
$\pi_{p_{K,r},p_{L,s}}$ niezerowy wkład mają wyłącznie
te pary $(i,j)$, dla których $K\cap\supp\eta_i\ne\varnothing$
i $L\cap\supp\theta_j\ne\varnothing$; takich par
jest skończenie wiele. Dlatego podwójna suma po $i,j$
ma sumę w uzupełnieniu rzutowym, a jej wartość
zależy ciągle od $\Phi$: dla ustalonych $K,L,r,s$
wystarcza skończenie wiele powyższych oszacowań.
Otrzymaliśmy ciągłe odwzorowanie
$\mathcal E(\Omega\times\Lambda)\to
\mathcal E(\Omega)\widehat\otimes_\pi\mathcal E(\Lambda)$,
które jest odwrotne do odwzorowania kanonicznego na
tensorach elementarnych. Ponieważ rozwinięcia lokalne
zbiegały w odpowiednich seminormach, odwzorowania te są
odwrotne również na uzupełnieniach.
Ponieważ obydwie przestrzenie w pierwszym czynniku są
nuklearne, Twierdzenie~\ref{thm:nuclear-pi-eps24}
daje również izomorfizmy dla $\varepsilon$.
\end{proof}

\subsection{Topologia induktywna: istotne zastrzeżenie dla \texorpdfstring{$\mathcal D(\Omega)$}{D(Omega)}}

\begin{definicja}[Induktywna topologia tensorowa]
\label{def:i-topology24}
Niech $E,F$ będą przestrzeniami lokalnie wypukłymi.
\emph{Induktywną topologią tensorową} $i$ na $E\otimes F$
nazywamy najsilniejszą topologię lokalnie wypukłą,
dla której kanoniczne odwzorowanie
$E\times F\to E\otimes F$ jest osobno ciągłe.
Piszemy $E\otimes_i F$, a przez
$E\widehat\otimes_i F$ oznaczamy jego uzupełnienie.
\end{definicja}

\begin{stwierdzenie}[Rola topologii induktywnej]
\label{prop:i-vs-pi24}
Dla dowolnych przestrzeni lokalnie wypukłych $E,F$
topologia induktywna $i$ jest nie słabsza od rzutowej
$\pi$. Jej ciągłe funkcjonały liniowe odpowiadają
dokładnie osobno ciągłym formom dwuliniowym
$E\times F\to\K$. W szczególności do identyfikacji
z twierdzeniem o jądrze Schwartza należy używać
odpowiedniej topologii induktywnej, a nie automatycznie
rzutowej.
\end{stwierdzenie}
\begin{proof}
Każda topologia, dla której odwzorowanie kanoniczne
jest łącznie ciągłe, zapewnia też osobną ciągłość,
więc topologia $\pi$ jest nie silniejsza od $i$.
Jeżeli $b$ jest osobno ciągła, jej algebraiczna
linearyzacja $\widetilde b$ jest ciągła dla topologii $i$:
rzeczywiście, do rodziny topologii dopuszczalnych w definicji
$i$ można dołączyć topologię wyznaczoną przez seminormę
$u\mapsto|\widetilde b(u)|$, gdyż jej ograniczenie do
$x\otimes y$ jest osobno ciągłe. Odwrotna implikacja
wynika z osobnej ciągłości kanonicznego odwzorowania.
\end{proof}

\begin{twierdzenie}[Produkt tensorowy funkcji testowych]
\label{thm:D-i-product24}
Niech $\Omega\subseteq\R^{d_1}$ i $\Lambda\subseteq\R^{d_2}$
będą niepustymi zbiorami otwartymi. Wówczas odwzorowanie
$f\otimes g\mapsto ((x,y)\mapsto f(x)g(y))$ przedłuża się
do kanonicznego izomorfizmu liniowo-topologicznego
\[
 \mathcal D(\Omega)\widehat\otimes_i\mathcal D(\Lambda)
 \simeq \mathcal D(\Omega\times\Lambda).
\]
\end{twierdzenie}
\begin{proof}
Najpierw dowodzimy następującego faktu. Jeżeli $G$ jest
zupełną przestrzenią lokalnie wypukłą Hausdorffa, to każde
osobno ciągłe odwzorowanie dwuliniowe o wartościach w $G$,
$b:\mathcal D(\Omega)\times\mathcal D(\Lambda)\to G$,
ma jednoznaczne ciągłe liniowe przedłużenie
$B:\mathcal D(\Omega\times\Lambda)\to G$, spełniające
$B(f\otimes g)=b(f,g)$.

Dla ustalonych zwartych $K\Subset\Omega$,
$L\Subset\Lambda$ ograniczenie $b$ do
$\mathcal D_K(\Omega)\times\mathcal D_L(\Lambda)$ jest
łącznie ciągłe. Aby to zobaczyć, dla dowolnej ciągłej
seminormy $s$ na $G$ stosujemy argument Baire'a
do domkniętych zbiorów
$A_n=\{f:s(b(f,g))\le1\text{ dla wszystkich }g\in V_n\}$,
gdzie $(V_n)$ jest malejącą bazą zrównoważonych otoczeń
zera w $\mathcal D_L$. Osobna ciągłość daje
$\mathcal D_K=\bigcup_n A_n$; pewien $A_n$ ma niepuste
wnętrze. Ponieważ jest absolutnie wypukły, zawiera
otoczenie zera. Skalując oba argumenty, otrzymujemy oszacowanie
$s(b(f,g))\le C p_{K,m}(f)p_{L,r}(g)$ dla pewnych
$m,r,C$, zależnych od $s,K,L$.

Dla każdego $\Phi\in\mathcal D(\Omega\times\Lambda)$
wcześniejszy lemat o rozkładzie tensorowym funkcji
testowych (Lemat~\textup{23.21}) zapewnia
przedstawienie $\Phi=\sum_n f_n\otimes g_n$, w którym
nośniki wszystkich $f_n$ i $g_n$ leżą w ustalonych
zwartych zbiorach, a szeregi
$\sum_n p_{K,m}(f_n)p_{L,r}(g_n)$ są zbieżne.
Zatem $\sum_n b(f_n,g_n)$ jest zbieżny w każdej
ciągłej seminormie na $G$, a dzięki zupełności $G$
ma sumę. Wspomniany lemat daje oszacowania współczynników
Fouriera na każdym ustalonym etapie
$\mathcal D_M(\Omega\times\Lambda)$, gdzie $M$ jest zwarty:
po zastosowaniu skończonego rozkładu jedności i wyborze
większych ustalonych nośników $K,L$ uzyskujemy dla
każdej seminormy $s$ na $G$ oszacowanie wartości
$\sum_n b(f_n,g_n)$ przez $C_s p_{M,r_s}(\Phi)$
dla pewnych $r_s\in\N_0$ i $C_s>0$.
Oszacowanie to daje ciągłość skonstruowanego odwzorowania
na każdym etapie $\mathcal D_M$.

Niezależność od przedstawienia trzeba uzasadnić bez
zakładania z góry, że obie sumy definiują funkcjonały
na całej przestrzeni. Dla każdego $\lambda\in G'$ forma
skalarna $(f,g)\mapsto\lambda(b(f,g))$ jest osobno ciągła.
Aby zastosować twierdzenie o jądrze Schwartza w jego
wersji operatorowej, zauważmy, że forma ta wyznacza
ciągłe odwzorowanie $\mathcal D(\Omega)\to
\mathcal D'(\Lambda)_\beta$: dla ograniczonego
$B\subset\mathcal D(\Lambda)$ istnieje etap
$\mathcal D_L(\Lambda)$ zawierający $B$ jako zbiór
ograniczony, a łączna ciągłość formy na
$\mathcal D_K(\Omega)\times\mathcal D_L(\Lambda)$
daje dla każdego $K\Subset\Omega$ oszacowanie
$\sup_{g\in B}|\lambda(b(f,g))|\le C_B p_{K,m}(f)$
na $\mathcal D_K(\Omega)$. Ciągłość na wszystkich
etapach jest równoważna ciągłości na $\mathcal D(\Omega)$.
Twierdzenie o jądrze Schwartza z \hyperref[sec:montel-nuklearne23]{rozdziału~23} daje
jednoznaczną dystrybucję $K_\lambda$, dla której
$\lambda(b(f,g))=\langle K_\lambda,f\otimes g\rangle$.
Dla dowolnego rozkładu tensorowego $\Phi$ ciągłość
$\lambda$ i zbieżność szeregu w $G$ pokazują, że
$\lambda(\sum_n b(f_n,g_n))=\langle K_\lambda,\Phi\rangle$.
Prawa strona nie zależy od rozkładu. Ponieważ ciągła
dualna $G'$ rozdziela punkty przestrzeni lokalnie
wypukłej Hausdorffa $G$, suma nie zależy od wyboru
rozwinięcia. Ten sam argument zapewnia liniowość
oraz jednoznaczność $B$.
Z własności uniwersalnej granicy induktywnej otrzymujemy
ciągłość $B$ na całej $\mathcal D(\Omega\times\Lambda)$.

Z drugiej strony odwzorowanie $(f,g)\mapsto f\otimes g$
z $\mathcal D(\Omega)\times\mathcal D(\Lambda)$ do
$\mathcal D(\Omega\times\Lambda)$ jest osobno ciągłe.
Wykazany fakt stwierdza zatem, że zupełna przestrzeń
$\mathcal D(\Omega\times\Lambda)$ ma tę samą własność
uniwersalną co uzupełnienie produktu
$\mathcal D(\Omega)\otimes_i\mathcal D(\Lambda)$.
Własność uniwersalna wyznacza jednoznaczny izomorfizm
liniowo-topologiczny, którego wartość na tensorach
elementarnych jest wskazana w tezie.
\end{proof}

\begin{przyklad}[Forma osobno ciągła, która nie jest łącznie ciągła]
\label{ex:separate-not-joint24}
Na $\mathcal D(\R)\times\mathcal D(\R)$ definiujemy
\[
 b(f,g)=\sum_{m=1}^\infty (-1)^m f^{(m)}(0)g(m).
\]
Suma jest skończona dla ustalonej funkcji $g$, ponieważ
$\supp g$ jest zwarty. Dla ustalonej $g$ funkcjonał
$f\mapsto b(f,g)$ jest więc ciągłą kombinacją skończenie
wielu pochodnych w zerze. Dla ustalonej $f$ otrzymujemy
lokalnie skończoną sumę dystrybucji punktowych
$\sum_m(-1)^mf^{(m)}(0)\delta_m$, a zatem funkcjonał
ciągły na $\mathcal D(\R)$. Forma $b$ jest osobno ciągła.

Przypuśćmy, że jest także łącznie ciągła. Wtedy istnieją
ciągłe seminormy $P,Q$ na $\mathcal D(\R)$ takie, że
$|b(f,g)|\le P(f)Q(g)$ dla wszystkich $f,g$.
Wybierzmy przedział zwarty $K\Subset\R$ zawierający
otoczenie zera. Ciągłość $P$ na $\mathcal D_K(\R)$
daje pewne $r\in\N_0$ i $C>0$ takie, że
$P(f)\le Cp_{K,r}(f)$ dla $f\in\mathcal D_K(\R)$.
Dla $m>r$ wybieramy $g_m\in\mathcal D(\R)$ o nośniku
w $(m-1/3,m+1/3)$ i $g_m(m)=1$. Wtedy
$|f^{(m)}(0)|\le CQ(g_m)p_{K,r}(f)$ dla wszystkich
$f\in\mathcal D_K(\R)$. Niech $\chi\in\mathcal D_K(\R)$
będzie równa $1$ w otoczeniu zera i niech
$f_N(x)=\chi(x)e^{iNx}$ (w przypadku rzeczywistym
można rozpatrzyć części rzeczywiste i urojone).
Wówczas $|f_N^{(m)}(0)|=N^m$, natomiast
$p_{K,r}(f_N)=O(N^r)$. Otrzymujemy sprzeczność.

Forma $b$ odpowiada dystrybucji
$\sum_{m\ge1}\partial_x^m\delta_0\otimes\delta_m$
na $\R^2$; suma jest lokalnie skończona ze względu
na drugą zmienną. Jest to jawny przykład, dla którego
własność uniwersalna topologii $i$ jest silniejsza
niż własność uniwersalna topologii $\pi$.
\end{przyklad}

\begin{uwaga}[Nie wolno utożsamiać trzech topologii]
\label{rem:D-not-projective24}
Na $\mathcal D(\Omega)\otimes\mathcal D(\Lambda)$
zawsze zachodzi $\varepsilon\le\pi\le i$, lecz ogólnie
nierówności mogą być ostre. W szczególności
$\mathcal D(\Omega)\widehat\otimes_\pi\mathcal D(\Lambda)$
nie jest \emph{kanonicznie izomorficzna topologicznie}
z $\mathcal D(\Omega\times\Lambda)$ przez odwzorowanie
$f\otimes g\mapsto f(x)g(y)$. Nuklearność
$\mathcal D(\Omega)$ może zapewniać równość $\pi=\varepsilon$,
ale nie wymusza $\pi=i$. Jest to powód, dla którego
w twierdzeniu o jądrze Schwartza posługiwaliśmy się
formami osobno ciągłymi, a nie wyłącznie łącznie ciągłymi.
\end{uwaga}

\subsection{Powiązanie z teorią jąder i dalsza geometria tensorowa}

\begin{wniosek}[Interpretacja jądra dystrybucyjnego]
\label{cor:kernel-i24}
Przy założeniach Twierdzenia~\ref{thm:D-i-product24}
forma $b:\mathcal D(\Omega)\times\mathcal D(\Lambda)\to\K$
jest osobno ciągła wtedy i tylko wtedy, gdy istnieje
jednoznaczna dystrybucja $K\in\mathcal D'(\Omega\times\Lambda)$
spełniająca $b(f,g)=\langle K,f\otimes g\rangle$.
\end{wniosek}
\begin{proof}
Przez własność uniwersalną produktu induktywnego
$b$ wyznacza ciągły funkcjonał na
$\mathcal D(\Omega)\otimes_i\mathcal D(\Lambda)$.
Przedłużamy go na uzupełnienie i stosujemy
Twierdzenie~\ref{thm:D-i-product24}. Odwrotnie, dla
$K\in\mathcal D'(\Omega\times\Lambda)$ odwzorowanie
$(f,g)\mapsto\langle K,f\otimes g\rangle$ jest osobno
ciągłe, gdyż dla ustalonego jednego czynnika
odwzorowanie drugiego czynnika w funkcję testową
na produkcie jest liniowe i ciągłe.
\end{proof}

\begin{przyklad}[Operator o jądrze na przekątnej]
\label{ex:diag-tensor24}
Dla $\Omega=\Lambda\subset\R^d$ dystrybucja
$\delta_\Delta$ na przekątnej $\Delta=\{(x,x):x\in\Omega\}$
zadaje formę $b(f,g)=\int_\Omega f(x)g(x)\,\mathrm d\lambda_d(x)$.
Jej interpretacja tensorowa nie wymaga utożsamienia
$\mathcal D(\Omega)\widehat\otimes_\pi\mathcal D(\Omega)$
z $\mathcal D(\Omega\times\Omega)$.
\end{przyklad}

\begin{uwaga}[Przyszłe rozszerzenia]
\label{rem:tensor-next24}
Topologiczne produkty tensorowe pojawią się ponownie
przy operatorach zwartych i nuklearnych, własności
aproksymacji oraz algebrach operatorowych.
W szczególności normy $\pi$ i $\varepsilon$ są dwiema
różnymi normami tensorowymi już na iloczynach dwóch
nieskończenie wymiarowych przestrzeni Hilberta.
Nie rozwijamy tu zastosowań do rozwiązywania równań
różniczkowych, które należą do odrębnego skryptu.
\end{uwaga}


% Wersja źródłowa: audyt II, jednostka 25
\setcounter{section}{24}
\noindent\textbf{Wyniki wykorzystywane.} Przyjmujemy znane
struktury przestrzeni $\mathcal D$, $\mathcal E$, $\mathcal S$,
ich zupełność i własność Montela, twierdzenie Banacha--Steinhausa,
a także podstawowe identyfikacje produktów tensorowych.
\par\medskip
\section{Przestrzenie funkcji testowych i dystrybucji: synteza topologiczna}
\label{sec:synteza25}

Po omówieniu przestrzeni lokalnie wypukłych, dystrybucji i produktów
 tensorowych możemy porównać najważniejsze przestrzenie funkcyjne bez
 utożsamiania ich topologii. W tym rozdziale $d\ge1$, $\K=\C$ lub
 $\K=\R$, a $\Omega\subseteq\R^d$ jest niepustym zbiorem otwartym.
 Funkcje testowe i dystrybucje mogą mieć wartości rzeczywiste lub
 zespolone; w twierdzeniach dotyczących transformaty Fouriera
 przyjmujemy $\K=\C$. Wszędzie $\supp f$ oznacza nośnik funkcji
 lub dystrybucji.

\subsection{Trzy przestrzenie funkcji: definicje i zanurzenia}

\begin{definicja}[Przestrzenie $\mathcal E$, $\mathcal D$ i $\mathcal S$]
\label{def:trzy-przestrzenie25}
Dla niepustego zbioru otwartego $\Omega\subseteq\R^d$ oznaczamy
$\mathcal E(\Omega)=C^\infty(\Omega)$, wyposażoną w topologię
jednostajnej zbieżności wszystkich pochodnych na zbiorach zwartych,
wyznaczoną przez seminormy
$p_{K,m}(f)=\max_{|\alpha|\le m}\sup_{x\in K}|\partial^\alpha f(x)|$,
gdzie $K\Subset\Omega$, $m\in\N_0$, oraz $\mathcal D(\Omega)=C_c^\infty(\Omega)$, wyposażoną w zwykłą
ścisłą topologię granicy induktywnej przestrzeni
$\mathcal D_K(\Omega)$, gdzie $K\Subset\Omega$.
Dla $\Omega=\R^d$ przez $\mathcal S(\R^d)$ oznaczamy przestrzeń
Schwartza z seminormami
$p_{N,m}(f)=\max_{|\alpha|\le m}\sup_{x\in\R^d}
(1+|x|)^N|\partial^\alpha f(x)|$, gdzie $N,m\in\N_0$.
Przestrzenie dualne topologiczne oznaczamy odpowiednio
$\mathcal E'(\Omega)$, $\mathcal D'(\Omega)$ i
$\mathcal S'(\R^d)$. Jeżeli podajemy topologię na przestrzeni
 dualnej, indeks $\sigma$ oznacza topologię słabą-*, zaś
 $\beta$ topologię jednostajnej zbieżności na zbiorach ograniczonych
 przestrzeni wyjściowej.
\end{definicja}

\begin{twierdzenie}[Ciągłe i gęste zanurzenia]
\label{thm:zanurzenia25}
Niech $d\ge1$. Odwzorowania identycznościowe
\[
 \mathcal D(\R^d)\longrightarrow\mathcal S(\R^d)
 \longrightarrow\mathcal E(\R^d)
\]
są ciągłymi, różnowartościowymi odwzorowaniami liniowymi o gęstych
obrazach. Dla dowolnego niepustego zbioru otwartego $\Omega\subseteq\R^d$
takie samo stwierdzenie zachodzi dla zanurzenia
$\mathcal D(\Omega)\to\mathcal E(\Omega)$.
Żadne z trzech wskazanych zanurzeń nie musi być izomorfizmem
liniowo-topologicznym na swój obraz z topologią podprzestrzeni.
\end{twierdzenie}
\begin{proof}
Dla ustalonego zwartego $K\Subset\R^d$ oraz $N,m\in\N_0$
mamy $p_{N,m}(f)\le (1+\sup_{x\in K}|x|)^N p_{K,m}(f)$
dla $f\in\mathcal D_K(\R^d)$. Stąd ciągłość
$\mathcal D_K\to\mathcal S$; własność uniwersalna granicy
induktywnej daje ciągłość $\mathcal D\to\mathcal S$.
Ciągłość $\mathcal S\to\mathcal E$ wynika z oszacowania
$p_{K,m}(f)\le p_{0,m}(f)$. Analogicznie sprawdzamy
ciągłość $\mathcal D(\Omega)\to\mathcal E(\Omega)$.

Aby wykazać gęstość, wybierzmy $\chi\in\mathcal D(\R^d)$
równe $1$ na kuli jednostkowej i zdefiniujmy
$\chi_R(x)=\chi(x/R)$ dla $R\ge1$. Dla
$f\in\mathcal S(\R^d)$ pochodne $(1-\chi_R)f$ są sumami
iloczynów pochodnych $f$ i $\chi_R$; na nośniku
$1-\chi_R$ zachodzi $|x|\ge R$ (ponieważ $\chi=1$ na kuli jednostkowej). Dla każdego $N,m$ ich seminormy są
kontrolowane przez $C_{N,m}R^{-1}p_{N+m+2,m}(f)$,
gdy $R$ jest duże. Zatem $\chi_R f\to f$ w $\mathcal S$.
Dla $f\in\mathcal E(\R^d)$ mamy $\chi_Rf=f$ na każdym
ustalonym zbiorze zwartym przy dostatecznie dużym $R$,
a więc $\chi_Rf\to f$ w $\mathcal E$. W ogólnej
$\Omega$ używamy funkcji odcięcia $\chi_j\in\mathcal D(\Omega)$,
równej $1$ w otoczeniu $K_j$, gdzie $(K_j)$ jest
wyczerpaniem $\Omega$ zbiorami zwartymi. Wtedy
$\chi_j f\to f$ w $\mathcal E(\Omega)$.

Na $\R^d$ niech $\psi\in\mathcal D(\R^d)$ będzie niezerowa
i $f_n(x)=\psi(x-ne_1)$. Wówczas $f_n\to0$ w
$\mathcal E(\R^d)$, ale nie w $\mathcal S(\R^d)$.
Z kolei $g_n(x)=e^{-n^2}\psi(x-ne_1)$ zbiega do zera w
$\mathcal S(\R^d)$, ale nie w $\mathcal D(\R^d)$:
nośniki $g_n$ nie mieszczą się w żadnym wspólnym zbiorze
zwartym. Także topologia $\mathcal D(\R^d)$ nie jest
topologią podprzestrzeni odziedziczoną z $\mathcal E(\R^d)$.
\end{proof}

\begin{uwaga}[Zupełność i lokalna wypukłość]
\label{rem:tabela-struktur25}
$\mathcal E(\Omega)$ i $\mathcal S(\R^d)$ są nuklearnymi
przestrzeniami Frécheta, a $\mathcal D(\Omega)$ jest nuklearną,
zupełną, zwykle niemetryzowalną przestrzenią LF. Wszystkie trzy są
przestrzeniami Montela. Właściwości te zostały opracowane we wcześniejszych
rozdziałach. W szczególności trwałość nuklearności
względem podprzestrzeni i ilorazów oraz odpowiednich
granic induktywnych została udowodniona w \hyperref[sec:montel-nuklearne23]{rozdziale~23}
(Twierdzenia~\textup{23.15} i~\textup{23.17});
nuklearność $\mathcal D_K$ i $\mathcal D$ wynika z
Twierdzenia~\textup{23.18}. Nie należy mylić nuklearności
z własnością Montela: te dwie własności były w skrypcie
uzasadniane oddzielnie.
\end{uwaga}

\subsection{Przechodzenie do przestrzeni dualnych topologicznych}

\begin{lemat}[Transpozycja gęstego zanurzenia]
\label{lem:transpozycja25}
Niech $X,Y$ będą przestrzeniami lokalnie wypukłymi
Hausdorffa, a $j:X\to Y$ ciągłym odwzorowaniem liniowym
o gęstym obrazie. Wówczas transpozycja
$j':Y'\to X'$, $j'(T)=T\circ j$, jest różnowartościowa
i ciągła zarówno dla topologii słabych-*, jak i dla
silnych topologii przestrzeni dualnych.
\end{lemat}
\begin{proof}
Jeżeli $T\circ j=0$, to $T=0$ przez gęstość $j(X)$
i ciągłość $T$. Dla $x\in X$ mamy
$|\langle j'T,x\rangle|=|\langle T,jx\rangle|$,
co daje ciągłość w topologiach słabych-*.
Jeśli $B\subseteq X$ jest ograniczony, to $j(B)$ jest
ograniczony w $Y$, zatem dla seminorm silnych zachodzi
$\sup_{x\in B}|\langle j'T,x\rangle|
\le \sup_{y\in j(B)}|\langle T,y\rangle|$.
To dowodzi ciągłości w topologiach silnych.
\end{proof}

\begin{wniosek}[Łańcuch przestrzeni dualnych]
\label{cor:dual-lancuch25}
Dla każdego $d\ge1$ zachodzą kanoniczne, ciągłe i
różnowartościowe odwzorowania liniowe
\[
 \mathcal E'(\R^d)\longrightarrow\mathcal S'(\R^d)
 \longrightarrow\mathcal D'(\R^d),
\]
zarówno w topologiach słabych-*, jak i silnych.
Dla każdego niepustego zbioru otwartego $\Omega\subseteq\R^d$
kanoniczne odwzorowanie
$\mathcal E'(\Omega)\to\mathcal D'(\Omega)$
ma te same własności. Samo stwierdzenie ciągłości i
różnowartościowości \emph{nie oznacza}, że są to zanurzenia
topologiczne na obrazy z topologią podprzestrzeni.
\end{wniosek}
\begin{proof}
Stosujemy Lemat~\ref{lem:transpozycja25} do zanurzeń
z Twierdzenia~\ref{thm:zanurzenia25}.
\end{proof}

\begin{twierdzenie}[Dystrybucje o zwartym nośniku]
\label{thm:compact-support-E25}
Niech $\Omega\subseteq\R^d$ będzie niepustym zbiorem otwartym.
Obraz kanonicznego odwzorowania
$\mathcal E'(\Omega)\to\mathcal D'(\Omega)$
składa się dokładnie z dystrybucji o zwartym nośniku.
Dokładniej, dla każdego $F\in\mathcal E'(\Omega)$ istnieją
$K\Subset\Omega$, $m\in\N_0$ i $C>0$ takie, że
\[
 |F(f)|\le C\max_{|\alpha|\le m}\sup_{x\in K}
 |\partial^\alpha f(x)|\qquad(f\in\mathcal E(\Omega)).
\]
Odwrotnie, jeżeli $T\in\mathcal D'(\Omega)$ ma zwarty nośnik,
to istnieje dokładnie jeden funkcjonał ciągły
$F\in\mathcal E'(\Omega)$ równy $T$ na $\mathcal D(\Omega)$.
\end{twierdzenie}
\begin{proof}
Każdy funkcjonał ciągły na przestrzeni Frécheta
$\mathcal E(\Omega)$ jest ograniczony przez skończenie wiele
seminorm $p_{K_j,m_j}$. Łącząc zbiory zwarte i zwiększając
rząd pochodnych, otrzymujemy oszacowanie z tezy. Gdy
$f\in\mathcal D(\Omega)$ ma nośnik rozłączny z $K$, jego
wartość wynosi $0$, więc $\supp F\subseteq K$.

Niech teraz $K=\supp T\Subset\Omega$ i wybierzmy
$\chi\in\mathcal D(\Omega)$ równą $1$ na otoczeniu $K$.
Kładziemy $F(f)=\langle T,\chi f\rangle$ dla
$f\in\mathcal E(\Omega)$. Jest to funkcjonał ciągły,
gdyż mnożenie przez $\chi$ jest ciągłym odwzorowaniem
$\mathcal E(\Omega)\to\mathcal D_{\supp\chi}(\Omega)$.
Wartość nie zależy od wyboru $\chi$: różnica dwóch takich
odcięć po pomnożeniu przez $f$ jest funkcją testową
zerową w otoczeniu $K$, a $T$ zeruje się na takich
funkcjach. Jeżeli $f\in\mathcal D(\Omega)$, to
$T(f)=T(\chi f)$, ponieważ $(1-\chi)f$ znika w
otoczeniu $K$. Jednoznaczność przedłużenia wynika
z gęstości $\mathcal D(\Omega)$ w $\mathcal E(\Omega)$.
\end{proof}

\subsection{Dystrybucje temperowane jako podprzestrzeń \texorpdfstring{$\mathcal D'$}{D'}}

\begin{twierdzenie}[Kryterium temperowania]
\label{thm:temp-criterion25}
Niech $T\in\mathcal D'(\R^d)$. Następujące warunki są
równoważne:
\begin{enumerate}[label=\textup{(\roman*)}]
\item $T$ należy do obrazu kanonicznego odwzorowania
$\mathcal S'(\R^d)\to\mathcal D'(\R^d)$;
\item istnieją $C>0$, $N,m\in\N_0$ takie, że
$|\langle T,\varphi\rangle|\le C p_{N,m}(\varphi)$
dla każdej $\varphi\in\mathcal D(\R^d)$.
\end{enumerate}
W szczególności każda dystrybucja o zwartym nośniku jest
temperowana.
\end{twierdzenie}
\begin{proof}
Gdy $T$ jest ograniczeniem ciągłego funkcjonału na
$\mathcal S$, jego ciągłość względem przeliczalnej rodziny
seminorm $p_{N,m}$ daje oszacowanie (ii), po zwiększeniu
$N$ i $m$. Odwrotnie, nierówność (ii) pozwala
przedłużyć $T$ jednoznacznie przez ciągłość z gęstej
podprzestrzeni $\mathcal D$ do $\mathcal S$; granica
wartości nie zależy od aproksymującego ciągu.
Jeżeli $T$ ma zwarty nośnik, to z
Twierdzenia~\ref{thm:compact-support-E25} otrzymujemy
$|T(\varphi)|\le C p_{K,m}(\varphi)\le C p_{0,m}(\varphi)$
na funkcjach testowych, a zatem (ii) zachodzi.
\end{proof}

\begin{przyklad}[Dystrybucja nietemperowana]
\label{ex:non-tempered25}
Na $\R$ funkcja $f(x)=e^x$ jest lokalnie całkowalna,
więc wyznacza dystrybucję regularną $T_f\in\mathcal D'(\R)$.
Nie jest ona temperowana. Wybierzmy nieujemną, niezerową
$\psi\in\mathcal D((-1,1))$ i niech
$\varphi_n(x)=e^{-n/2}\psi(x-n)$. Każda seminorma
$\mathcal S$ funkcji $\varphi_n$ dąży do zera, ale
\[
 \langle T_f,\varphi_n\rangle=
 e^{n/2}\int_\R e^y\psi(y)\,\mathrm d\lambda_1(y)
 \longrightarrow+\infty.
\]
Ciągłość przedłużenia na $\mathcal S$ byłaby z tym
sprzeczna.
\end{przyklad}

\subsection{Topologie, zbieżność i charakterystyczne przykłady}

\begin{twierdzenie}[Zbieżność ciągów w przestrzeniach dualnych]
\label{thm:weak-strong-sequences25}
Niech $X$ będzie beczkową przestrzenią Montela nad $\K$.
Dla ciągu $(T_n)$ w $X'$ i $T\in X'$ następujące
warunki są równoważne: $T_n\to T$ w topologii
$\sigma(X',X)$ oraz $T_n\to T$ w topologii $\beta(X',X)$.
Stwierdzenie stosuje się w szczególności do
$X=\mathcal D(\Omega)$, $X=\mathcal E(\Omega)$ i
$X=\mathcal S(\R^d)$.
\end{twierdzenie}
\begin{proof}
Zbieżność silna implikuje słabą-*, gdyż zbiory
jednoelementowe są ograniczone. W przeciwną stronę,
zbieżność skalarna $\langle T_n,x\rangle\to\langle T,x\rangle$
dla każdego $x$ daje punktową ograniczoność rodziny
$\{T_n-T:n\ge1\}$. Z beczkowości i twierdzenia
Banacha--Steinhausa wynika jej równociągłość. Niech
$B\subseteq X$ będzie ograniczony. Domknięcie $B$ jest
zwarte na mocy własności Montela. Dla $\varepsilon>0$
wybieramy skończenie wiele punktów $x_1,\ldots,x_k\in B$
i otoczenie zera $U$ tak, aby $B\subseteq\bigcup_{j=1}^k(x_j+U)$
oraz $|\langle T_n-T,u\rangle|<\varepsilon/2$
dla wszystkich $n$ i $u\in U$. Dla dostatecznie dużych $n$
mamy $|\langle T_n-T,x_j\rangle|<\varepsilon/2$
dla $j=1,\ldots,k$. W konsekwencji
$\sup_{x\in B}|\langle T_n-T,x\rangle|<\varepsilon$,
co kończy dowód.
\end{proof}

\begin{przyklad}[Zbieżność i ucieczka nośnika]
\label{ex:escaping-deltas25}
Na $\R$ niech $\delta_n$ oznacza deltę Diraca w punkcie
$n$. Wtedy $\delta_n\to0$ silnie w $\mathcal D'(\R)$:
każdy ograniczony zbiór funkcji testowych ma nośniki
zawarte w jednym zwartym przedziale. Również
$\delta_n\to0$ silnie w $\mathcal S'(\R)$, bo dla
ograniczonego zbioru $B\subseteq\mathcal S$ mamy
$\sup_{\varphi\in B}|\varphi(n)|\le C_B(1+n)^{-2}$.
Jednak $\delta_n\not\to0$ nawet słabo-* w
$\mathcal E'(\R)$, ponieważ $\langle\delta_n,1\rangle=1$.

Z kolei $T_n=e^n\delta_n$ zbiega silnie do zera w
$\mathcal D'(\R)$, lecz nie jest zbieżny do zera nawet
słabo-* w $\mathcal S'(\R)$: funkcja
$\varphi(x)=\exp(-\sqrt{1+x^2}/2)$ należy do
$\mathcal S(\R)$ i $e^n\varphi(n)\to\infty$.
Zatem ciągłe i różnowartościowe odwzorowania z
Wniosku~\ref{cor:dual-lancuch25} nie są ogólnie
zanurzeniami topologicznymi na obraz z topologią
podprzestrzeni, także przy silnych topologiach dualnych.
\end{przyklad}

\subsection{Operatory naturalne i granice ich działania}

\begin{stwierdzenie}[Różniczkowanie i mnożenie]
\label{prop:ops25}
Niech $\Omega\subseteq\R^d$ będzie niepustym zbiorem otwartym.
Dla każdego wielowskaźnika $\alpha\in\N_0^d$
operator $\partial^\alpha$ działa liniowo i ciągle na
$\mathcal E(\Omega)$ i $\mathcal D(\Omega)$, a przez
transpozycję także na $\mathcal E'(\Omega)$ i
$\mathcal D'(\Omega)$ w topologiach słabych-* i silnych.
Dla $a\in C^\infty(\Omega)$ mnożenie $f\mapsto af$
i odpowiadające mu mnożenie dystrybucji działają ciągle
na tych przestrzeniach.
Na $\mathcal S(\R^d)$ i $\mathcal S'(\R^d)$ te same
własności mnożenia zachodzą, jeżeli dla każdego
wielowskaźnika $\alpha$ istnieją $C_\alpha>0$ i
$N_\alpha\in\N_0$ takie, że
$|\partial^\alpha a(x)|\le C_\alpha(1+|x|)^{N_\alpha}$
dla wszystkich $x\in\R^d$.
\end{stwierdzenie}
\begin{proof}
Ciągłość operatorów na przestrzeniach testowych wynika z
nierówności Leibniza i z definicji seminorm. W przestrzeni
$\mathcal D$ mnożenie nie powiększa nośnika, więc
wystarcza sprawdzenie ciągłości na każdym etapie
$\mathcal D_K$. Dla $\mathcal S$ ustalmy $N,m$.
Po zastosowaniu wzoru Leibniza każdą seminormę
$p_{N,m}(af)$ można oszacować przez
$C_{N,m,a}p_{N+M,m}(f)$, gdzie $M$ jest maksimum
wykładników wzrostu pochodnych $a$ rzędu co najwyżej $m$.
Operator transponowany działa ciągle w topologii słabej-*,
a w silnej dlatego, że ciągłe odwzorowania liniowe
przenoszą zbiory ograniczone na ograniczone. Dla
pochodnej dystrybucyjnej pojawia się zwykły czynnik
$(-1)^{|\alpha|}$ w formule transpozycji.
\end{proof}

\begin{przyklad}[Nie każdy mnożnik gładki zachowuje przestrzeń Schwartza]
\label{ex:not-multiplier25}
Funkcja $a(x)=e^{|x|^2}$ jest gładka, ale mnożenie przez nią
nie zachowuje przestrzeni Schwartza: dla
$f(x)=e^{-|x|^2}\in\mathcal S(\R^d)$ iloczyn $af=1$
nie należy do $\mathcal S(\R^d)$. W przestrzeniach
$\mathcal D$ i $\mathcal E$ mnożenie przez $a$ nie
powoduje takiego problemu.
\end{przyklad}

\begin{uwaga}[Transformacja Fouriera i splot]
\label{rem:fourier-conv25}
Transformacja Fouriera jest automorfizmem
$\mathcal S(\R^d)$ i $\mathcal S'(\R^d)$.
Nie jest automorfizmem $\mathcal D'(\R^d)$, ponieważ
nie każda dystrybucja należy do $\mathcal S'$, a jej
transformata Fouriera nie musi być zdefiniowana przez
zwykłą transpozycję na przestrzeni testowej $\mathcal D$.
Splot z dystrybucją o zwartym nośniku jest natomiast
określony zarówno na $\mathcal D'$ jak i na $\mathcal S'$;
splot dwóch dowolnych dystrybucji nie musi być określony.
Teoria mnożenia dystrybucji przez funkcje gładkie i splotu
została udowodniona we wcześniejszych rozdziałach.
\end{uwaga}

\subsection{Lokalność i problemy przedłużania dystrybucji}

\begin{twierdzenie}[Ograniczenie i brak automatycznego przedłużenia]
\label{thm:restriction25}
Jeżeli $U\subseteq\Omega\subseteq\R^d$ są niepustymi
zbiorami otwartymi, to odwzorowanie ograniczenia
$\mathcal D'(\Omega)\to\mathcal D'(U)$ jest liniowe
i ciągłe w topologiach słabych-* oraz silnych.
Nie każda dystrybucja na $U$ przedłuża się do
dystrybucji na $\Omega$.
\end{twierdzenie}
\begin{proof}
Funkcję testową $\varphi\in\mathcal D(U)$ przedłużamy
przez zero do $\Omega$ i oznaczamy przedłużenie przez
$\widetilde\varphi$. Operator
$\mathcal D(U)\to\mathcal D(\Omega)$,
$\varphi\mapsto\widetilde\varphi$, jest ciągły: na każdym
etapie z ustalonym zwartym nośnikiem wynika to z równości
pochodnych i seminorm, a na całej przestrzeni z własności
uniwersalnej granicy induktywnej. Stąd transpozycja
daje ciągłe ograniczenie w obydwu topologiach dualnych.

Aby wykazać drugie stwierdzenie, weźmy
$U=(0,1)$, $\Omega=(-1,1)$ i dystrybucję regularną
$T_f$ wyznaczoną przez $f(x)=\exp(1/x^2)$ na $U$.
Gdyby miała przedłużenie $\widetilde T\in\mathcal D'(\Omega)$,
to dla pewnego zwartego $K=[-1/2,1/2]\Subset\Omega$,
$m\in\N_0$ i $C>0$ zachodziłoby
$|\langle\widetilde T,\varphi\rangle|\le C p_{K,m}(\varphi)$
dla wszystkich $\varphi\in\mathcal D_K(\Omega)$.
Niech $\psi\in\mathcal D((-1,1))$ będzie nieujemna i niezerowa,
$r_n=1/(10n^2)$ oraz
$\varphi_n(x)=\psi((x-1/n)/r_n)$ dla dużych $n$.
Wówczas $\supp\varphi_n\subset U\cap K$,
$p_{K,m}(\varphi_n)=O(n^{2m})$, natomiast
$\langle T_f,\varphi_n\rangle\ge c n^{-2}e^{c'n^2}$
dla pewnych $c,c'>0$. To sprzeczność.
\end{proof}

\subsection{Schemat zależności i produkty tensorowe}

\begin{stwierdzenie}[Zestawienie kanonicznych odwzorowań]
\label{prop:summary25}
Dla każdego $d\ge1$ istnieją ciągłe i gęste zanurzenia
liniowe $\mathcal D(\R^d)\to\mathcal S(\R^d)\to
\mathcal E(\R^d)$ oraz dualne odwzorowania ciągłe
i różnowartościowe
$\mathcal E'(\R^d)\to\mathcal S'(\R^d)\to
\mathcal D'(\R^d)$. Dla otwartych niepustych
$\Omega\subseteq\R^{d_1}$, $\Lambda\subseteq\R^{d_2}$
kanoniczne izomorfizmy
\[
\mathcal E(\Omega)\widehat\otimes_\pi\mathcal E(\Lambda)
\simeq\mathcal E(\Omega\times\Lambda),\quad
\mathcal D(\Omega)\widehat\otimes_i\mathcal D(\Lambda)
\simeq\mathcal D(\Omega\times\Lambda)
\]
łącznie z analogicznym izomorfizmem dla przestrzeni Schwartza,
porządkują związek między przestrzeniami testowymi a
ciągłymi formami dwuliniowymi. W szczególności nie wolno
bez dodatkowych założeń zastępować topologii $i$ przez $\pi$
w drugim wzorze.
\end{stwierdzenie}
\begin{proof}
Stwierdzenia o zanurzeniach i transpozycjach wynikają z
Twierdzenia~\ref{thm:zanurzenia25} i
Wniosku~\ref{cor:dual-lancuch25}.
Identyfikacje tensorowe udowodniono we wcześniejszym
rozdziale o produktach tensorowych; ich zestawienie
nie wymaga nowego przejścia granicznego.
\end{proof}

\begin{uwaga}[Co pozostaje do rozwinięcia]
\label{rem:future25}
Zastosowania do rozwiązywania równań różniczkowych
pozostają w odrębnym skrypcie. W tym opracowaniu
będziemy rozwijać dalej geometrię przestrzeni Banacha,
refleksywność i bazy Schaudera, własności aproksymacji,
teorię operatorów i analizę spektralną.
Moduły 15A (całkowanie w przestrzeniach lokalnie wypukłych)
i 15B (całki Bochnera, Pettisa, Dunforda i Gelfanda)
zostały już umieszczone bezpośrednio po \hyperref[sec:wyzsze-pochodne15]{rozdziale 15}.
Rozdział bieżący korzysta z wyników tych modułów jedynie
tam, gdzie jest to jawnie zaznaczone.
\end{uwaga}


% Wersja źródłowa: audyt II, jednostka 26
\setcounter{section}{25}
\noindent\textbf{Wyniki wykorzystywane.} Zakładamy udowodnione
w poprzednich rozdziałach twierdzenia Hahna--Banacha,
Banacha--Alaoglu i Banacha--Steinhausa oraz podstawy topologii słabej.
Twierdzenia Goldstine'a, Eberleina--\v Smuliana i Milmana--Pettisa
udowodnimy tutaj.
\par\medskip
\section{Geometria przestrzeni Banacha: refleksywność, zwartość słaba i wypukłość jednostajna}
\label{sec:geometria26}

W tym rozdziale $X$ oznacza przestrzeń Banacha nad $\K\in\{\R,\C\}$,
$X'$ jej przestrzeń dualną topologiczną, a $X''=(X')'$ bidual.
Przez $B_X=\{x\in X:\|x\|\le1\}$ rozumiemy domkniętą kulę
jednostkową, natomiast $\sigma(X,X')$ oznacza topologię słabą.
Wszystkie domknięcia oznaczone indeksem $w$ dotyczą tej topologii.
Celem jest ustalenie, kiedy ograniczoność daje zwartość słabą,
co wolno wówczas wnioskować o ciągach oraz jak geometria normy
wpływa na refleksywność.

\subsection{Kanoniczne zanurzenie i twierdzenie Goldstine'a}

\begin{definicja}[Refleksywność]
\label{def:refleksywnosc26}
Przestrzeń Banacha $X$ nazywamy \emph{refleksywną}, jeżeli
kanoniczne odwzorowanie liniowe
$J_X:X\to X''$, określone wzorem
$(J_Xx)(f)=f(x)$ dla $x\in X$ i $f\in X'$, jest surjektywne.
\end{definicja}

\begin{stwierdzenie}[Kanoniczne zanurzenie]
\label{prop:canonical26}
Dla każdej przestrzeni unormowanej $X$ odwzorowanie
$J_X:X\to X''$ jest liniową izometrią. Jeżeli $X$ jest zupełna,
to $J_X(X)$ jest domkniętą podprzestrzenią $X''$.
\end{stwierdzenie}
\begin{proof}
Liniowość wynika z definicji. Dla $x\in X$ mamy
$\|J_Xx\|=\sup_{\|f\|\le1}|f(x)|\le\|x\|$.
Dla $x\ne0$ twierdzenie Hahna--Banacha zapewnia funkcjonał
$f\in X'$ o normie $1$ i wartości $f(x)=\|x\|$.
Daje to odwrotną nierówność. Obraz przestrzeni zupełnej
przez izometrię jest zupełny, a zupełna podprzestrzeń
przestrzeni metrycznej jest domknięta.
\end{proof}

\begin{twierdzenie}[Goldstine]
\label{thm:goldstine26}
Niech $X$ będzie przestrzenią unormowaną nad $\K$.
Wówczas $J_X(B_X)$ jest gęsty w $B_{X''}$ względem topologii
$\sigma(X'',X')$.
\end{twierdzenie}
\begin{proof}
Przypuśćmy, że $\Phi\in B_{X''}$ nie należy do
słabo-* domkniętego zbioru $J_X(B_X)$.
Istnieją wówczas $f_1,\ldots,f_m\in X'$ takie, że wektor
$(\Phi(f_1),\ldots,\Phi(f_m))$ nie należy do domknięcia
w $\K^m$ wypukłego zbioru
$\{(f_1(x),\ldots,f_m(x)):x\in B_X\}$.
Twierdzenie o separacji w skończonym wymiarze daje skalary
$c_1,\ldots,c_m\in\K$ i liczbę rzeczywistą $a$ spełniające
\[
\sup_{x\in B_X}\operatorname{Re}\Big(\sum_{j=1}^m c_jf_j(x)\Big)
<a<\operatorname{Re}\Big(\sum_{j=1}^m c_j\Phi(f_j)\Big).
\]
Dla $f=\sum_jc_jf_j$ lewa strona ma wartość $\|f\|$,
a prawa nie przekracza $\|\Phi\|\|f\|\le\|f\|$.
Otrzymana sprzeczność kończy dowód.
\end{proof}

\begin{twierdzenie}[Refleksywność a zwartość słaba kuli]
\label{thm:reflexive-ball26}
Niech $X$ będzie przestrzenią Banacha. Następujące warunki
są równoważne:
\begin{enumerate}[label=\textup{(\roman*)}]
\item $X$ jest refleksywna;
\item $B_X$ jest zwarta w topologii $\sigma(X,X')$;
\item każdy niepusty, ograniczony, normowo domknięty i wypukły
podzbiór $X$ jest słabo zwarty.
\end{enumerate}
\end{twierdzenie}
\begin{proof}
Jeżeli $J_X$ jest surjektywne, to jego ograniczenie do $B_X$
identyfikuje słabą topologię kuli z topologią słabą-* na $B_{X''}$.
Twierdzenie Banacha--Alaoglu daje zwartość, więc (i) pociąga (ii).
Jeżeli zachodzi (ii), to dla każdego ograniczonego normowo,
zamkniętego i wypukłego $C\subseteq X$ istnieje $R>0$ z
$C\subseteq RB_X$. Z twierdzenia Hahna--Banacha wynika, że
normowo domknięty zbiór wypukły jest słabo domknięty.
Zatem $C$ jest domkniętym podzbiorem słabo zwartej kuli $RB_X$,
co daje (iii). Przypadek $C=\varnothing$ nie wpływa na argument.
Warunek (iii) zastosowany do $B_X$ daje (ii).
Wreszcie jeśli $B_X$ jest słabo zwarta, to $J_X(B_X)$ jest
słabo-* zwarta, a zatem słabo-* domknięta w $X''$.
Twierdzenie~\ref{thm:goldstine26} mówi, że jest ona gęsta
w $B_{X''}$; stąd $J_X(B_X)=B_{X''}$.
Jednorodność implikuje surjektywność $J_X$, czyli (i).
\end{proof}

\begin{wniosek}[Podprzestrzenie i przestrzeń dualna]
\label{cor:reflexivity-stability26}
Jeżeli $X$ jest refleksywną przestrzenią Banacha, to każda
jej domknięta podprzestrzeń liniowa jest refleksywna.
Ponadto $X$ jest refleksywna wtedy i tylko wtedy, gdy $X'$
jest refleksywna.
\end{wniosek}
\begin{proof}
Niech $Y\subseteq X$ będzie domkniętą podprzestrzenią.
Jako zbiór domknięty i wypukły $Y$ jest słabo domknięta.
Zatem $B_Y=B_X\cap Y$ jest słabo zwarta w $X$.
Z twierdzenia Hahna--Banacha każda $g\in Y'$ przedłuża się
do $X'$, toteż topologia słaba $Y$ jest dokładnie topologią
indukowaną z $X$. Twierdzenie~\ref{thm:reflexive-ball26}
daje refleksywność $Y$.

Jeżeli $X$ jest refleksywna, to $J_X:X\to X''$ jest
izomorfizmem izometrycznym. Dla $G\in X'''$ określamy
$f(x)=G(J_Xx)$. Wtedy $f\in X'$ i dla $\Phi=J_Xx\in X''$
mamy $(J_{X'}f)(\Phi)=\Phi(f)=f(x)=G(\Phi)$.
Surjektywność $J_X$ daje $J_{X'}f=G$, więc $X'$ jest refleksywna.
Odwrotnie, gdy $X'$ jest refleksywna i $J_X(X)\ne X''$,
twierdzenie Hahna--Banacha daje niezerowy $G\in X'''$
zerujący się na domkniętej podprzestrzeni $J_X(X)$.
Z refleksywności $X'$ wynika $G=J_{X'}f$ dla pewnego $f\in X'$.
Wtedy $0=G(J_Xx)=f(x)$ dla wszystkich $x\in X$, więc
$f=0$ i $G=0$, sprzeczność.
\end{proof}

\subsection{Twierdzenie Eberleina--\v Smuliana}

W dowolnej przestrzeni topologicznej zwartości nie można
charakteryzować wyłącznie ciągami. Dla słabej topologii
przestrzeni Banacha zachodzi jednak ważny wyjątek.

\begin{lemat}[Skończona rodzina normująca]
\label{lem:finite-norming26}
Niech $X$ będzie niezerową przestrzenią Banacha,
a $M\subseteq X''$ skończenie wymiarową podprzestrzenią.
Istnieje skończony niepusty zbiór $F\subseteq B_{X'}$ taki, że
$\|\Psi\|\le 2\max_{f\in F}|\Psi(f)|$ dla wszystkich $\Psi\in M$.
\end{lemat}
\begin{proof}
Dla $M=\{0\}$ wystarczy wybrać dowolny $f\in B_{X'}$.
W pozostałym przypadku dla każdego $\Psi\in M$ o normie $1$
wybieramy $f_\Psi\in B_{X'}$ z
$|\Psi(f_\Psi)|>3/4$. Kula jednostkowa w skończenie
wymiarowej przestrzeni $M$ ma zwartą sferę jednostkową.
Pokrywamy ją kulami $\{\Theta\in M:\|\Theta-\Psi\|<1/4\}$
i wybieramy skończone podpokrycie o środkach
$\Psi_1,\ldots,\Psi_r$. Dla $\Theta$ na sferze i odpowiedniego
$j$ mamy $|\Theta(f_{\Psi_j})|>3/4-1/4=1/2$.
Jednorodność daje tezę dla dowolnego $\Theta\in M$.
\end{proof}

\begin{twierdzenie}[Eberlein--\v Smulian]
\label{thm:eberlein-smulian26}
Niech $X$ będzie przestrzenią Banacha, a $A\subseteq X$
dowolnym podzbiorem. Następujące warunki są równoważne:
\begin{enumerate}[label=\textup{(\roman*)}]
\item domknięcie $\overline A^{\,w}$ jest słabo zwarte;
\item z każdego ciągu punktów $A$ można wybrać podciąg
zbieżny słabo do pewnego punktu $X$;
\item każdy ciąg w $A$ ma słaby punkt skupienia w $X$
(w każdym jego słabym otoczeniu leżą wyrazy ciągu o dowolnie dużych indeksach).
\end{enumerate}
\end{twierdzenie}
\begin{proof}
Gdy $A=\varnothing$, wszystkie warunki zachodzą.
Zakładamy więc $A\ne\varnothing$; przypadek $X=\{0\}$
jest natychmiastowy.

\emph{Krok pierwszy: (i) $\Rightarrow$ (ii).}
Niech $(x_n)\subseteq A$ i
$Y=\overline{\lin\{x_n:n\ge1\}}^{\,\|\cdot\|}$.
Jest to ośrodkowa, domknięta podprzestrzeń $X$,
a więc jest słabo domknięta. Zbiór
$K=\overline A^{\,w}\cap Y$ jest słabo zwarty i zawiera $(x_n)$.
Wybieramy ciąg gęsty $(y_j)$ w $Y$ i funkcjonały
$f_j\in B_{X'}$ takie, że $f_j(y_j)=\|y_j\|$.
Taka rodzina rozdziela punkty $Y$: dla $y\ne0$ wystarczy
wybrać $y_j$ w odległości mniejszej niż $\|y\|/3$ od $y$.
Wówczas $|f_j(y)|\ge\|y_j\|-\|y-y_j\|>0$.
Na $K$ funkcja
\[
d(u,v)=\sum_{j=1}^\infty 2^{-j}\min\{1,|f_j(u-v)|\}
\]
jest metryką. Topologia tej metryki jest słabsza od topologii
słabej, ale jest Hausdorffa. Ponieważ $K$ jest słabo zwarty,
identyczność z $K$ słabo zwartą do $K$ metryczną jest
homeomorfizmem: ciągła bijekcja z przestrzeni zwartej
do Hausdorffa ma ciągłą odwrotność. W konsekwencji $K$
jest zwartą przestrzenią metryczną, więc $(x_n)$ ma w $K$
podciąg zbieżny słabo.

\emph{Krok drugi: (ii) $\Rightarrow$ (iii).}
Granica podciągu słabo zbieżnego jest słabym punktem skupienia
całego ciągu, gdyż indeksy podciągu rosną do nieskończoności.

\emph{Krok trzeci: (iii) $\Rightarrow$ (i).}
Najpierw $A$ jest ograniczony w normie. W przeciwnym razie
rodzina $\{J_Xx:x\in A\}\subseteq X''$, rozpatrywana jako
rodzina funkcjonałów na przestrzeni Banacha $X'$, byłaby
nieograniczona jednostajnie w normie. Z twierdzenia
Banacha--Steinhausa istniałby $f\in X'$ z
$\sup_{x\in A}|f(x)|=\infty$. Wybierając $x_n\in A$ z
$|f(x_n)|\ge n$, dostalibyśmy ciąg bez słabego punktu
skupienia: obrazy przez $f$ nie mają punktu skupienia
w $\K$.

Niech $\Phi\in\overline{J_X(A)}^{\,w^\ast}\subseteq X''$.
Skonstruujemy ciąg $(x_n)\subseteq A$ oraz skończone
zbiory $F_n\subseteq B_{X'}$. Wybieramy $x_1\in A$ i zbiór
$F_1$ z Lematu~\ref{lem:finite-norming26}, normujący z dokładnością
do czynnika $2$ przestrzeń
$M_1=\lin\{\Phi,J_Xx_1\}\subseteq X''$.
Gdy $x_1,\ldots,x_n$ i $F_1,\ldots,F_n$ są ustalone,
wybieramy $x_{n+1}\in A$ tak, by
$|\Phi(f)-f(x_{n+1})|<1/(n+1)$ dla każdego
$f\in F_1\cup\cdots\cup F_n$; jest to możliwe, ponieważ
$\Phi$ należy do słabo-* domkniętego zbioru $J_X(A)$.
Następnie wybieramy $F_{n+1}$ normujący z dokładnością
do czynnika $2$ podprzestrzeń
$M_{n+1}=\lin\{\Phi,J_Xx_1,\ldots,J_Xx_{n+1}\}$.

Na mocy (iii) istnieje słaby punkt skupienia $x\in X$
ciągu $(x_n)$. Dla każdego $f\in\bigcup_nF_n$ mamy
$f(x_n)\to\Phi(f)$, zatem $f(x)=\Phi(f)$.
Ponadto $x$ należy do słabego domknięcia zbioru
$\{x_n\}$, a stąd do normowo domkniętej podprzestrzeni
$Y=\overline{\lin\{x_n:n\ge1\}}^{\,\|\cdot\|}$,
która jest słabo domknięta na mocy twierdzenia Hahna--Banacha.
W szczególności $J_Xx$ należy do domknięcia normowego
przestrzeni $\lin\{J_Xx_n:n\ge1\}$, zawartej w
$\bigcup_nM_n$. Ponieważ $\Phi\in M_1$, również
$\Psi=\Phi-J_Xx$ należy do normowego domknięcia
$\bigcup_nM_n$. Niech $G\in M_n$ będzie jego
przybliżeniem w normie. Z konstrukcji dla wszystkich
$f\in F_n$ zachodzi $\Psi(f)=0$, a więc
$G(f)=(G-\Psi)(f)$.
Ponieważ $G\in M_n$ oraz $F_n\subseteq B_{X'}$,
z własności rodziny normującej mamy
$\|G\|\le2\max_{f\in F_n}|G(f)|\le2\|G-\Psi\|$.
Stąd $\|\Psi\|\le3\|\Psi-G\|$, a przybliżając
$\Psi$ coraz dokładniej elementami $G\in\bigcup_nM_n$
otrzymujemy $\Psi=0$. Zatem $\Phi=J_Xx$.

Pokazaliśmy, że słabo-* domknięcie $J_X(A)$ zawiera się
w $J_X(X)$. Z pierwszej części kroku trzeciego zbiór $A$
jest normowo ograniczony, a więc jego obraz ma słabo-* zwarte
domknięcie w $X''$ na mocy twierdzenia Banacha--Alaoglu.
Ponieważ $J_X:(X,\sigma(X,X'))\to
(J_X(X),\sigma(X'',X')|_{J_X(X)})$ jest homeomorfizmem,
odpowiadający temu domknięciu zbiór w $X$ jest słabo zwarty
i zawiera $A$. Domknięcie słabe $A$ jest jego domkniętym
podzbiorem, zatem również jest słabo zwarte.
\end{proof}

\begin{wniosek}[Ograniczone ciągi w przestrzeni refleksywnej]
\label{cor:bounded-subsequence26}
Jeżeli $X$ jest refleksywną przestrzenią Banacha,
to każdy ograniczony ciąg w $X$ ma podciąg słabo zbieżny.
\end{wniosek}
\begin{proof}
Wszystkie wyrazy ciągu należą do pewnej słabo zwartej kuli
$RB_X$ na mocy Twierdzenia~\ref{thm:reflexive-ball26}.
Stosujemy Twierdzenie~\ref{thm:eberlein-smulian26}.
\end{proof}

\subsection{Własność Schura i geometria przestrzeni \texorpdfstring{$\ell^1$}{ell1}}

\begin{definicja}[Własność Schura]
\label{def:schur26}
Przestrzeń Banacha $X$ ma \emph{własność Schura}, jeżeli dla
każdego ciągu $(x_n)\subseteq X$ i $x\in X$ ze zbieżności
$x_n\to x$ słabo wynika $\|x_n-x\|\to0$.
\end{definicja}

\begin{twierdzenie}[Własność Schura w $\ell^1$]
\label{thm:schur-l1-26}
Przestrzeń Banacha $\ell^1(\N,\K)$ ma własność Schura.
\end{twierdzenie}
\begin{proof}
Wystarczy rozważyć ciąg $(x_n)$ zbieżny słabo do zera.
Wówczas dla każdej współrzędnej $j$ mamy $x_n(j)\to0$.
Przypuśćmy, że normy nie dążą do zera. Po przejściu do
podciągu możemy założyć, że $\|x_n\|_1\ge\varepsilon>0$
dla wszystkich $n$. Rekurencyjnie wybieramy podciąg
$y_k=x_{n_k}$ oraz liczby całkowite
$0=m_0<m_1<m_2<\cdots$ takie, że
\[
\sum_{j\le m_{k-1}}|y_k(j)|<\varepsilon/8,
\qquad
\sum_{j>m_k}|y_k(j)|<\varepsilon/8.
\]
Pierwszą nierówność można osiągnąć przez zbieżność
współrzędnych, a drugą przez sumowalność wybranego $y_k$.
Na rozłącznym bloku $I_k=\{m_{k-1}+1,\ldots,m_k\}$
mamy zatem $\sum_{j\in I_k}|y_k(j)|\ge3\varepsilon/4$.
Definiujemy $a=(a_j)\in\ell^\infty$ tak, że dla $j\in I_k$
wybieramy $a_j=\overline{y_k(j)}/|y_k(j)|$, gdy $y_k(j)\ne0$,
i $a_j=0$ w przeciwnym wypadku. Wówczas
$f(x)=\sum_{j\ge1}a_jx(j)$ jest ciągłym funkcjonałem
liniowym na $\ell^1$. Dla każdego $k$,
\[
\operatorname{Re}f(y_k)\ge
\sum_{j\in I_k}|y_k(j)|-\sum_{j\notin I_k}|y_k(j)|
\ge\varepsilon/2.
\]
Jest to sprzeczne ze słabą zbieżnością $y_k\to0$.
\end{proof}

\begin{wniosek}[Nierefleksywność przestrzeni $\ell^1$]
\label{cor:l1-nonreflexive26}
Nieskończenie wymiarowa przestrzeń Banacha $\ell^1(\N,\K)$
nie jest refleksywna.
\end{wniosek}
\begin{proof}
Wektory jednostkowe $e_n$ spełniają
$\|e_n-e_m\|_1=2$ dla $n\ne m$, więc żaden ich podciąg
nie jest zbieżny w normie. Gdyby $\ell^1$ była refleksywna,
Wniosek~\ref{cor:bounded-subsequence26} zapewniałby podciąg
słabo zbieżny, co przeczyłoby własności Schura.
\end{proof}

\begin{przyklad}[Słabo zbieżny ciąg bez zbieżności normowej]
\label{ex:en-hilbert26}
Niech $H$ będzie nieskończenie wymiarową przestrzenią Hilberta,
a $(e_n)$ przeliczalną rodziną ortonormalną.
Dla każdego $x\in H$ nierówność Bessela daje
$\langle e_n,x\rangle\to0$, a więc $e_n\to0$ słabo
w $H$. Jednak $\|e_n\|=1$ dla wszystkich $n$.
Dlatego nieskończenie wymiarowa przestrzeń Hilberta
nie ma własności Schura.
\end{przyklad}

\subsection{Jednostajna wypukłość i twierdzenie Milmana--Pettisa}

\begin{definicja}[Ścisła i jednostajna wypukłość]
\label{def:uniform-convex26}
Niech $X$ będzie przestrzenią unormowaną.
Mówimy, że norma jest \emph{ściśle wypukła}, jeżeli
dla różnych $x,y\in X$ o normach $1$ zachodzi
$\|(x+y)/2\|<1$.
Przestrzeń $X$ jest \emph{jednostajnie wypukła}, jeżeli
dla każdego $\varepsilon\in(0,2]$ istnieje
$\delta(\varepsilon)>0$ takie, że dla wszystkich
$x,y\in X$ o normach $1$ i $\|x-y\|\ge\varepsilon$
zachodzi $\|(x+y)/2\|\le1-\delta(\varepsilon)$.
\end{definicja}

\begin{przyklad}[Moduł wypukłości przestrzeni Hilberta]
\label{ex:modulus-hilbert26}
Niech $H$ będzie rzeczywistą lub zespoloną przestrzenią
Hilberta. Dla $x,y\in H$ o normach $1$ tożsamość
równoległoboku daje
$\|(x+y)/2\|^2=1-\|x-y\|^2/4$.
Zatem $H$ jest jednostajnie wypukła, a można przyjąć
$\delta_H(\varepsilon)=1-\sqrt{1-\varepsilon^2/4}$.
Natomiast przestrzeń $\ell^1$ nie jest ściśle wypukła:
$\|e_1\|_1=\|e_2\|_1=\|(e_1+e_2)/2\|_1=1$.
\end{przyklad}

\begin{twierdzenie}[Milman--Pettis]
\label{thm:milman-pettis26}
Każda jednostajnie wypukła przestrzeń Banacha jest refleksywna.
\end{twierdzenie}
\begin{proof}
Niech $\Phi\in X''$ oraz $\|\Phi\|=1$.
Z Twierdzenia~\ref{thm:goldstine26} istnieje sieć
$(x_\alpha)\subseteq B_X$ taka, że
$J_Xx_\alpha\to\Phi$ słabo-*.
Ustalmy $\varepsilon\in(0,2]$ i wybierzmy
$\delta>0$ mniejsze od $\min\{\delta_X(\varepsilon/2),\varepsilon/4\}$.
Z definicji normy $\Phi$ istnieje $f\in X'$ o normie $1$
takie, że $\operatorname{Re}\Phi(f)>1-\delta$,
po ewentualnym pomnożeniu $f$ przez skalar modułu $1$.
Dla dostatecznie dużych $\alpha$ zachodzi
$\operatorname{Re}f(x_\alpha)>1-\delta$.
Jeżeli $\alpha,\beta$ są dostatecznie duże, to
\[
\left\|\frac{x_\alpha+x_\beta}{2}\right\|
\ge\operatorname{Re}f\left(\frac{x_\alpha+x_\beta}{2}\right)
>1-\delta.
\]
Ponieważ $\operatorname{Re}f(x_\alpha)>1-\delta$, dla tych
indeksów $\|x_\alpha\|>1-\delta$ i możemy położyć
$u_\alpha=x_\alpha/\|x_\alpha\|$. Wektory $u_\alpha$
mają normę $1$, a $\operatorname{Re}f(u_\alpha)>1-\delta$.
Wobec tego $\|(u_\alpha+u_\beta)/2\|>1-\delta$.
Dzięki wybranemu $\delta$ jednostajna wypukłość daje
$\|u_\alpha-u_\beta\|<\varepsilon/2$.
Ponadto $\|u_\alpha-x_\alpha\|=1-\|x_\alpha\|<\delta$,
zatem $\|x_\alpha-x_\beta\|<\varepsilon/2+2\delta<\varepsilon$.
Stąd $(x_\alpha)$ jest siecią Cauchy'ego w normie.
Zupełność $X$ daje $x_\alpha\to x$ w normie, więc
$\Phi=J_Xx$, gdyż granica słaba-* jest jednoznaczna.
Zatem $J_X$ jest surjektywne.
\end{proof}

\begin{twierdzenie}[Własność Radona--Riesza]
\label{thm:radon-riesz26}
Niech $X$ będzie jednostajnie wypukłą przestrzenią Banacha,
a $(x_n)\subseteq X$ i $x\in X$. Jeżeli $x_n\to x$
słabo oraz $\|x_n\|\to\|x\|$, to $\|x_n-x\|\to0$.
\end{twierdzenie}
\begin{proof}
Gdy $x=0$, teza wynika wprost ze zbieżności norm.
W przeciwnym razie dla dużych $n$ kładziemy
$u_n=x_n/\|x_n\|$ i $u=x/\|x\|$.
Wtedy $\|u_n\|=\|u\|=1$ i $u_n\to u$ słabo.
Gdyby $\|u_n-u\|\not\to0$, istniałyby
$\varepsilon>0$ i podciąg z $\|u_{n_k}-u\|\ge\varepsilon$.
Z jednostajnej wypukłości wynika
$\|(u_{n_k}+u)/2\|\le1-\delta_X(\varepsilon)$.
Tymczasem $(u_{n_k}+u)/2\to u$ słabo, a norma jest
słabo półciągła z dołu (jest supremum modułów
ciągłych funkcjonałów liniowych).
Dostalibyśmy $1\le1-\delta_X(\varepsilon)$,
sprzeczność. Stąd $\|u_n-u\|\to0$, a powrót do
oryginalnej normalizacji kończy dowód.
\end{proof}

\subsection{Granice i kontrprzykłady}

\begin{twierdzenie}[Słabe domknięcie sfery jednostkowej]
\label{thm:sphere-weak-dense26}
Niech $X$ będzie nieskończenie wymiarową przestrzenią
unormowaną. Wówczas słabe domknięcie sfery jednostkowej
$S_X=\{x\in X:\|x\|=1\}$ jest równe $B_X$.
\end{twierdzenie}
\begin{proof}
Kula $B_X$ jest wypukła i normowo domknięta,
a więc słabo domknięta. Niech $x\in B_X$ ma normę mniejszą
od $1$, zaś $U$ będzie jego podstawowym słabym otoczeniem,
wyznaczonym przez funkcjonały $f_1,\ldots,f_m\in X'$.
Ponieważ $X$ jest nieskończenie wymiarowa, istnieje
niezerowy $v\in\bigcap_{j=1}^m\ker f_j$.
Funkcja rzeczywista $t\mapsto\|x+tv\|$ jest ciągła,
ma wartość mniejszą od $1$ dla $t=0$ i dąży do
nieskończoności wraz z $t\to\infty$.
Dla pewnego $t>0$ mamy więc $\|x+tv\|=1$,
a jednocześnie $f_j(x+tv)=f_j(x)$ dla każdego $j$.
Zatem $U\cap S_X\ne\varnothing$.
\end{proof}

\begin{uwaga}[Co wynika z twierdzenia Eberleina--\v Smuliana]
\label{rem:weak-compact-what26}
Dla podzbiorów przestrzeni Banacha zwartość słaba
jest równoważna zwartości ciągowej w topologii słabej,
ale nie oznacza
zwartości w normie. W przestrzeniach innych niż przestrzenie
Banacha takiej równoważności nie wolno zakładać bez dowodu.
Również jednostajna wypukłość implikuje refleksywność,
lecz implikacja odwrotna nie jest ogólnie prawdziwa dla
ustalonej normy.
\end{uwaga}


% Wersja źródłowa: audyt II, jednostka 27
\setcounter{section}{26}
\noindent\textbf{Wyniki wykorzystywane.}
Przyjmujemy znane twierdzenia Hahna--Banacha,
Banacha--Steinhausa, o odwzorowaniu otwartym,
twierdzenie Eberleina--\v Smuliana oraz podstawy
ortogonalności i bazy ortonormalnej przestrzeni Hilberta.
\par\medskip
\section{Bazy Schaudera, układy biortogonalne i zbieżność bezwarunkowa}
\label{sec:bazy27}

W tym rozdziale $X$ oznacza przestrzeń Banacha nad $\K\in\{\R,\C\}$,
$X'$ jej przestrzeń dualną topologiczną, a $\mathcal L(X,Y)$
przestrzeń operatorów liniowych ciągłych. Bazy ortonormalne przestrzeni
Hilberta omówiono we wcześniejszym rozdziale o przestrzeniach Hilberta; tutaj badamy bazy
\emph{uporządkowane}, właściwe ogólnym przestrzeniom Banacha.
Nie należy mylić baz Schaudera z bazami Hamela: w przypadku tych
ostatnich dopuszcza się wyłącznie \emph{skończone} kombinacje liniowe.

\subsection{Definicja, współczynniki i ciągłość rzutów częściowych}

\begin{definicja}[Baza Schaudera i współczynniki]
\label{def:schauder27}
Niech $X$ będzie przestrzenią Banacha, a $(e_n)_{n\ge1}\subset X$
ciągiem wektorów niezerowych. Nazywamy go \emph{bazą Schaudera}
przestrzeni $X$, jeśli dla każdego $x\in X$ istnieje dokładnie jeden
ciąg skalarów $(a_n)_{n\ge1}$, dla którego
$x=\lim_{N\to\infty}\sum_{n=1}^N a_ne_n$ w normie $X$.
Skalar $a_n$ oznaczamy $e_n^*(x)$ i nazywamy $n$-tym
\emph{współczynnikiem} wektora $x$. Odwzorowanie
$P_Nx=\sum_{n=1}^Ne_n^*(x)e_n$ nazywamy
\emph{rzutem częściowym}; kładziemy też $P_0=0$.
\end{definicja}

Jednoznaczność implikuje liniowość $e_n^*$ i $P_N$; nie zakładaliśmy
natomiast jeszcze ich ciągłości. Jest to istotne: ciągłość
współczynników stanowi treść następnego twierdzenia.

\begin{lemat}[Przestrzeń ciągów współczynników]
\label{lem:coefficient-space27}
Niech $X$ będzie przestrzenią Banacha, a $(e_n)_{n\ge1}\subset X$
ciągiem wektorów niezerowych. Zbiór
\[
E_e=\bigl\{a=(a_n)\in\K^{\N}:\sum_{n=1}^{\infty}a_ne_n
\text{ jest zbieżny w }X\bigr\}
\]
z normą $\|a\|_{E_e}=\sup_{N\ge1}\|\sum_{n=1}^Na_ne_n\|_X$
jest przestrzenią Banacha. Ponadto odwzorowanie
$S:E_e\to X$, $S(a)=\sum_{n\ge1}a_ne_n$, jest liniowe
ciągłe i spełnia $\|S\|\le1$.
\end{lemat}
\begin{proof}
Norma jest skończona dzięki zbieżności sum częściowych.
Jeżeli wszystkie sumy częściowe znikają, to z różnic kolejnych
sum i z $e_n\ne0$ wynika $a_n=0$ dla wszystkich $n$.
Niech $(a^{(k)})$ będzie ciągiem Cauchy'ego w $E_e$.
Dla każdego $n$ mamy
$|a_n^{(k)}-a_n^{(\ell)}|\,\|e_n\|
\le2\|a^{(k)}-a^{(\ell)}\|_{E_e}$,
więc istnieje granica $a_n=\lim_k a_n^{(k)}$.
Przechodząc do granicy przy ustalonym $N$, otrzymujemy
$\|\sum_{n=1}^N(a_n^{(k)}-a_n)e_n\|\le\eta$
jednostajnie w $N$, gdy $k$ jest dostatecznie duże,
bo analogiczne oszacowanie z $a^{(\ell)}$ jest
prawdziwe dla każdego $\ell$.
Wybieramy takie $k$. Dla $M>N$ mamy
\[
\Big\|\sum_{n=N+1}^Ma_ne_n\Big\|
\le2\eta+\Big\|\sum_{n=N+1}^Ma_n^{(k)}e_n\Big\|.
\]
W ostatnim wyrażeniu, dla ustalonego $k$, sumy ogonowe
są dowolnie małe przy dużym $N$. Ponieważ $\eta>0$
można przyjąć dowolnie małe, szereg dla $a$ jest zbieżny.
Ponadto $\|a^{(k)}-a\|_{E_e}\to0$.
Dla każdego $a\in E_e$ mamy
$\|S(a)\|\le\|a\|_{E_e}$, co dowodzi tezy o $S$.
\end{proof}

\begin{twierdzenie}[Ciągłość współczynników i stała bazy]
\label{thm:basis-continuity27}
Jeżeli $(e_n)_{n\ge1}$ jest bazą Schaudera przestrzeni
Banacha $X$, to $e_n^*\in X'$ dla każdego $n$, a każdy
$P_N\in\mathcal L(X,X)$. Ponadto
\[
K_e:=\sup_{N\ge1}\|P_N\|<\infty,\qquad K_e\ge1.
\]
Liczbę $K_e$ nazywamy \emph{stałą bazy}.
Dla każdego $n\ge1$ zachodzi $e_n^*(e_m)=\delta_{nm}$
oraz $\|e_n^*\|\,\|e_n\|\le2K_e$.
\end{twierdzenie}
\begin{proof}
Dla bazy Schaudera odwzorowanie $S:E_e\to X$
z Lematu~\ref{lem:coefficient-space27} jest liniową ciągłą
bijekcją. Twierdzenie o odwzorowaniu otwartym zapewnia
ciągłość $S^{-1}$. Dla każdego $N$ zachodzi
$\|P_Nx\|=\|\sum_{n=1}^Ne_n^*(x)e_n\|
\le\|S^{-1}x\|_{E_e}$,
stąd $K_e\le\|S^{-1}\|<\infty$.
Dla $n\ge1$ tożsamość
$(P_n-P_{n-1})x=e_n^*(x)e_n$
daje $\|e_n^*\|\,\|e_n\|\le\|P_n\|+\|P_{n-1}\|
\le2K_e$. Zatem wszystkie współczynniki są ciągłe.
Biortogonalność wynika z jednoznaczności reprezentacji $e_m$.
Każdy $P_N$ jest niezerowym rzutem (pozostawia $e_1$),
a więc $\|P_N\|\ge1$.
\end{proof}

\begin{wniosek}[Normalizacja i ośrodkowość]
\label{cor:separable27}
Każda przestrzeń Banacha z bazą Schaudera jest ośrodkowa.
Jeżeli $(e_n)$ jest bazą Schaudera, to
$(e_n/\|e_n\|)$ również jest bazą Schaudera o tej samej stałej
bazy; jej współczynniki są równe $\|e_n\|e_n^*(x)$.
\end{wniosek}
\begin{proof}
Skończone kombinacje $e_n$ o współczynnikach w $\mathbb Q$
(dla $\K=\C$ w $\mathbb Q+i\mathbb Q$) tworzą zbiór przeliczalny
i gęsty w $X$, bo sumy częściowe bazy przybliżają każdy $x$.
Zmiana normalizacji nie zmienia ani sum częściowych,
ani ich granic, zatem nie zmienia operatorów $P_N$.
\end{proof}

\begin{definicja}[Układ biortogonalny i ciąg bazowy]
\label{def:biorthogonal27}
Niech $X$ będzie przestrzenią Banacha.
Układ $(e_n,f_n)_{n\ge1}\subset X\times X'$ nazywamy
\emph{biortogonalnym}, jeżeli $f_n(e_m)=\delta_{nm}$
dla wszystkich $m,n$. Ciąg $(e_n)$ nazywamy
\emph{ciągiem bazowym}, jeżeli jest bazą Schaudera swojej
normowo domkniętej otoczki liniowej $[e_n]
=\overline{\operatorname{span}\{e_n:n\ge1\}}^{\,\|\cdot\|}$.
\end{definicja}

\begin{twierdzenie}[Kryterium bazy przez rzuty]
\label{thm:biorthogonal-criterion27}
Niech $X$ będzie przestrzenią Banacha, zaś
$(e_n,f_n)_{n\ge1}\subset X\times X'$ układem biortogonalnym,
dla którego $\operatorname{span}\{e_n:n\ge1\}$ jest gęsta w $X$.
Oznaczmy $P_Nx=\sum_{n=1}^Nf_n(x)e_n$.
Wówczas $(e_n)$ jest bazą Schaudera $X$, której
funkcjonałami współrzędnych są $f_n$, wtedy i tylko wtedy,
gdy $\sup_N\|P_N\|<\infty$.
\end{twierdzenie}
\begin{proof}
Warunek jest konieczny na mocy
Twierdzenia~\ref{thm:basis-continuity27}.
W przeciwną stronę niech $K=\sup_N\|P_N\|<\infty$.
Jeśli $y$ jest skończoną kombinacją $e_n$, to $P_Ny=y$
dla dostatecznie dużych $N$. Dla dowolnego $x\in X$
wybieramy takie $y$ bliskie $x$; wówczas
$\|x-P_Nx\|\le(1+K)\|x-y\|$ dla dużych $N$.
Zatem $P_Nx\to x$ w normie.
Jeżeli również $x=\sum_na_ne_n$, to przez ciągłość $f_j$
i biortogonalność dostajemy $f_j(x)=a_j$.
To dowodzi jednoznaczności rozwinięcia.
\end{proof}

\begin{wniosek}[Kryterium ciągu bazowego]
\label{cor:basic-sequence27}
Niech $(e_n)_{n\ge1}$ będzie ciągiem wektorów niezerowych
w przestrzeni Banacha $X$. Jest to ciąg bazowy wtedy i tylko wtedy,
gdy jest liniowo niezależny i istnieje $K<\infty$ takie, że
$\|\sum_{n=1}^Na_ne_n\|
\le K\|\sum_{n=1}^Ma_ne_n\|$
dla wszystkich $M\ge N\ge1$ i wszystkich $a_1,\ldots,a_M\in\K$.
\end{wniosek}
\begin{proof}
Dla ciągu bazowego nierówność wynika z ograniczoności
rzutów częściowych na $[e_n]$.
Odwrotnie, na algebraicznej otoczce liniowej $e_n$
jednoznacznie definiujemy $P_N$ przez ucięcie współczynników.
Założona nierówność daje $\|P_N\|\le K$.
Każdy $P_N$ przedłuża się do $[e_n]$ przez ciągłość;
wówczas $P_Nx\to x$ przez gęstość skończonych kombinacji.
Współczynniki są ciągłe, ponieważ
$(P_n-P_{n-1})x$ należy do jednowymiarowej przestrzeni
$\K e_n$, a ciągłość funkcjonałów współrzędnych
wynika z podzielenia przez stały niezerowy wektor $e_n$.
Stosujemy Twierdzenie~\ref{thm:biorthogonal-criterion27}
do przestrzeni $[e_n]$.
\end{proof}

\subsection{Bazy monotoniczne i klasyczne przykłady}

\begin{definicja}[Baza monotoniczna]
\label{def:monotonic27}
Bazę Schaudera $(e_n)$ nazywamy \emph{monotoniczną}, gdy
$\|P_N\|=1$ dla wszystkich $N\ge1$.
\end{definicja}

\begin{przyklad}[Przestrzenie ciągowe]
\label{ex:sequence-bases27}
W przestrzeniach Banacha $c_0(\N,\K)$ i $\ell^p(\N,\K)$
dla $1\le p<\infty$ wektory standardowe
$e_n=(\delta_{nj})_{j\ge1}$ tworzą bazę Schaudera.
Współczynnik $e_n^*(x)$ to $n$-ta współrzędna,
a $P_N(x)$ zeruje wszystkie współrzędne po $N$.
Dla $\ell^p$ zachodzi
$\|x-P_Nx\|_p^p=\sum_{j>N}|x_j|^p\to0$,
a dla $c_0$ mamy $\|x-P_Nx\|_\infty=\sup_{j>N}|x_j|\to0$.
Stąd $\|P_N\|=1$ i bazy są monotoniczne.
W $\ell^\infty$ ciąg standardowy nie jest bazą Schaudera:
jego domknięta otoczka liniowa równa jest $c_0\ne\ell^\infty$.
\end{przyklad}

\begin{przyklad}[Baza sumująca w $c_0$]
\label{ex:summing-basis27}
Niech $s_n=\sum_{j=1}^ne_j\in c_0$.
Dla $x=(x_j)\in c_0$ przyjmijmy $a_n=x_n-x_{n+1}$.
Wtedy $\sum_{n=1}^Na_ns_n$ ma na pozycjach $j\le N$
wartość $x_j-x_{N+1}$, a po pozycjach $j>N$ wartość $0$.
Zatem
$\|x-\sum_{n=1}^Na_ns_n\|_\infty
\le2\sup_{j>N}|x_j|\to0$.
Jednoznaczność współczynników wynika z tożsamości
$a_n=x_n-x_{n+1}$; układ $(s_n)$ jest więc bazą Schaudera.
Jej rzuty częściowe mają normę $2$ dla każdego $N\ge1$,
ponieważ $\|P_Nx\|_\infty
\le\|x\|_\infty+|x_{N+1}|\le2\|x\|_\infty$,
a równość otrzymujemy przez wybór $x_1=1$, $x_{N+1}=-1$
i pozostałych współrzędnych równych zeru (dla $N\ge1$).
Ta baza nie jest monotoniczna.
\end{przyklad}

\subsection{Bazy bezwarunkowe i zmiana kolejności sumowania}

\begin{definicja}[Baza bezwarunkowa]
\label{def:unconditional27}
Bazę Schaudera $(e_n)$ przestrzeni Banacha $X$ nazywamy
\emph{bezwarunkową}, jeżeli dla każdego $x\in X$ i każdej
permutacji $\pi:\N\to\N$ szereg
$\sum_{n=1}^\infty e_{\pi(n)}^*(x)e_{\pi(n)}$
jest zbieżny w normie $X$ do $x$.
Dla skończonego $F\subseteq\N$ piszemy
$P_Fx=\sum_{n\in F}e_n^*(x)e_n$.
\end{definicja}

\begin{lemat}[Permutacje a sumy po zbiorach skończonych]
\label{lem:net-unconditional27}
Niech $(u_n)$ będzie ciągiem w przestrzeni Banacha $X$
i załóżmy, że szereg $\sum_nu_n$ jest zbieżny w zadanej
kolejności do $u\in X$. Następujące warunki są równoważne:
\begin{enumerate}[label=\textup{(\roman*)}]
\item wszystkie permutacje szeregu są zbieżne do $u$;
\item sieć $(\sum_{n\in F}u_n)_{F\Subset\N}$, skierowana przez inkluzję,
 zbiega do $u$;
\item dla każdego $\varepsilon>0$ istnieje skończony $F_0\subseteq\N$
 taki, że $\|\sum_{n\in F}u_n\|<\varepsilon$
 dla każdego skończonego $F\subseteq\N\setminus F_0$.
\end{enumerate}
\end{lemat}
\begin{proof}
Jeśli (iii) zachodzi, to dla skończonych
$F,G\supseteq F_0$ różnica ich sum jest różnicą sum po
$F\setminus G$ i $G\setminus F$. Oba zbiory są rozłączne
z $F_0$, zatem nierówność trójkąta daje normę różnicy
mniejszą niż $2\varepsilon$. Zatem sieć jest
siecią Cauchy'ego. W przestrzeni Banacha ma granicę,
która musi być równa $u$, ponieważ zawiera ona jako podsieć
sumy po $\{1,\ldots,N\}$.
Implikacja (ii) $\Rightarrow$ (i) wynika stąd, że
zbiory $\{\pi(1),\ldots,\pi(N)\}$ wyczerpują $\N$.
Załóżmy, że (iii) nie zachodzi. Istnieje $\varepsilon>0$
takie, że po usunięciu dowolnego zbioru skończonego
można znaleźć zbiór skończony $F$ z normą sumy
przynajmniej $\varepsilon$. Rekurencyjnie konstruujemy
rozłączne zbiory skończone $F_k$, których elementy są coraz
większe, i których sumy mają normę przynajmniej $\varepsilon$.
Każde $F_k$ wybieramy poza poprzednimi zbiorami i poza
$\{1,\ldots,k\}$. Permutację $\N$ budujemy blokami,
wyliczając kolejno elementy $F_k$ i wstawiając
między blokami najmniejszy dotąd niewykorzystany indeks.
Wyrazy tak przestawionego szeregu mają sumy częściowe,
które na końcach każdego bloku $F_k$ różnią się o co najmniej
$\varepsilon$. Szereg po tej permutacji nie jest więc
zbieżny, co przeczy (i).
\end{proof}

\begin{twierdzenie}[Kryterium bezwarunkowości przez rzuty]
\label{thm:unconditional-projections27}
Niech $(e_n)$ będzie bazą Schaudera przestrzeni Banacha $X$.
Wówczas następujące warunki są równoważne:
\begin{enumerate}[label=\textup{(\roman*)}]
\item $(e_n)$ jest bazą bezwarunkową;
\item istnieje $U<\infty$ takie, że $\|P_F\|\le U$
 dla każdego skończonego $F\subseteq\N$;
\item dla każdego $x\in X$ rodzina
$\{P_Fx:F\Subset\N\}$ jest ograniczona w normie.
\end{enumerate}
Najmniejsze $U$ z (ii) nazywamy \emph{stałą bezwarunkowości
przez opuszczanie współczynników}.
\end{twierdzenie}
\begin{proof}
Z (i) i Lematu~\ref{lem:net-unconditional27} wynika, że
$P_Fx\to x$ w sieci zbiorów skończonych, zatem rodzina
$\{P_Fx\}$ jest ograniczona. To daje (iii).
Twierdzenie Banacha--Steinhausa zastosowane do operatorów
$P_F\in\mathcal L(X,X)$ daje (iii) $\Rightarrow$ (ii).
Jeżeli zachodzi (ii), to dla $x\in X$ i dla przybliżenia
$y\in\operatorname{span}\{e_n\}$ mamy
$\|P_F(x-y)\|\le U\|x-y\|$.
Gdy $F$ zawiera wszystkie indeksy występujące w $y$,
to $P_Fy=y$, skąd $\|P_Fx-x\|\le(U+1)\|x-y\|$.
Gęstość skończonych kombinacji implikuje zbieżność
sieci $P_Fx\to x$, a Lemat~\ref{lem:net-unconditional27}
daje (i).
\end{proof}

\begin{wniosek}[Bazy ortonormalne i bezwarunkowość]
\label{cor:hilbert-unconditional27}
Niech $H$ będzie ośrodkową, nieskończenie wymiarową
przestrzenią Hilberta, a $(e_n)_{n\ge1}$ dowolną numeracją
jej bazy ortonormalnej. Wtedy $(e_n)$ jest monotoniczną
bazą Schaudera i bezwarunkową o stałej $U=1$.
\end{wniosek}
\begin{proof}
Twierdzenie Parsevala daje zbieżność rozwinięcia
$\sum_n\langle x,e_n\rangle e_n$ do $x$.
Dla każdego skończonego $F$ twierdzenie Pitagorasa
daje $\|P_Fx\|^2=\sum_{n\in F}|\langle x,e_n\rangle|^2
\le\|x\|^2$. Zarówno $P_F$, jak i $P_N$ są
niezerowymi rzutami, więc mają normę dokładnie $1$.
Twierdzenie~\ref{thm:unconditional-projections27} kończy dowód.
\end{proof}

\begin{przyklad}[Baza warunkowa w przestrzeni $c_0$]
\label{ex:conditional-summing27}
Baza sumująca $(s_n)$ z Przykładu~\ref{ex:summing-basis27}
nie jest bezwarunkowa. Weźmy $x\in c_0$ o współrzędnych
$x_n=(-1)^n/\sqrt n$. Wtedy
$a_n=s_n^*(x)=x_n-x_{n+1}=(-1)^n(1/\sqrt n+1/\sqrt{n+1})$.
Szereg skalarów $\sum_n a_n$ jest zbieżny (jego suma to
$x_1$), ale nie jest bezwzględnie zbieżny.
Gdyby szereg wektorowy $\sum_n a_ns_n$ był bezwarunkowo
zbieżny w $c_0$, to po zastosowaniu ciągłego
funkcjonału $\delta_1(y)=y_1$ uzyskalibyśmy bezwarunkową
zbieżność $\sum_n a_n$ w $\K$, co jest równoważne
zbieżności bezwzględnej. Otrzymujemy sprzeczność.
\end{przyklad}

\begin{przyklad}[Bezwarunkowość nie oznacza absolutnej zbieżności]
\label{ex:not-absolute27}
W przestrzeni Hilberta $\ell^2$ szereg
$\sum_{n\ge1}n^{-1}e_n$ jest zbieżny bezwarunkowo
(każdy wybór skończonej rodziny wyrazów ma kwadrat
normy równy sumie $n^{-2}$), ale
$\sum_{n\ge1}\|n^{-1}e_n\|=\sum_{n\ge1}n^{-1}=\infty$.
\end{przyklad}

\subsection{Bazy ściągające i zupełne w sensie ograniczonym}

\begin{definicja}[Dwa warunki na bazę]
\label{def:shrinking-bcomplete27}
Niech $(e_n)$ będzie bazą Schaudera przestrzeni Banacha $X$,
a $(e_n^*)\subset X'$ jej funkcjonałami współrzędnych.
Bazę nazywamy \emph{ściągającą} (ang. \emph{shrinking}), gdy
$\overline{\operatorname{span}\{e_n^*:n\ge1\}}^{\,\|\cdot\|}=X'$.
Nazywamy ją \emph{zupełną w sensie ograniczonym}
(ang. \emph{boundedly complete}), jeżeli dla każdego ciągu
skalarów $(a_n)$ warunek
$\sup_N\|\sum_{n=1}^Na_ne_n\|<\infty$
implikuje zbieżność szeregu $\sum_na_ne_n$ w normie $X$.
\end{definicja}

\begin{twierdzenie}[Kryterium refleksywności przez bazę]
\label{thm:james-basis27}
Niech $X$ będzie przestrzenią Banacha z bazą Schaudera $(e_n)$.
Wtedy $X$ jest refleksywna wtedy i tylko wtedy, gdy
$(e_n)$ jest jednocześnie ściągająca i zupełna w sensie
ograniczonym.
\end{twierdzenie}
\begin{proof}
Załóżmy najpierw refleksywność $X$.
Aby dowieść zupełności w sensie ograniczonym,
niech $s_N=\sum_{n=1}^Na_ne_n$ będzie ograniczony.
Kula zawierająca wszystkie $s_N$ jest słabo zwarta.
Istnieje zatem podciąg słabo zbieżny do pewnego $x\in X$
na mocy twierdzenia Eberleina--\v Smuliana.
Dla każdego $j$ ciąg $e_j^*(s_N)$ jest od pewnego miejsca
stały i równy $a_j$, więc $e_j^*(x)=a_j$.
Stąd $P_Nx=s_N$ dla wszystkich $N$, a definicja bazy
daje $s_N\to x$ w normie.

Niech teraz $f\in X'$. Oznaczmy
$Y_N=\overline{\operatorname{span}\{e_j:j>N\}}^{\,\|\cdot\|}$.
Pokażemy, że $\|f|_{Y_N}\|\to0$, gdzie norma oznacza
normę ograniczenia $f$ do domkniętej podprzestrzeni $Y_N$.
Gdyby tak nie było, istniałyby $\varepsilon>0$, liczby
$N_k\nearrow\infty$ i wektory $y_k\in Y_{N_k}$
o normie $1$ takie, że $|f(y_k)|\ge\varepsilon$.
Wektory te można wybierać rekurencyjnie, gdyż $Y_{N+1}\subseteq Y_N$,
a więc normy ograniczeń tworzą ciąg nierosnący.
Refleksywność i twierdzenie Eberleina--\v Smuliana
dają podciąg $(y_{k_\ell})$ słabo zbieżny do $y\in X$.
Dla każdego ustalonego $j$ mamy
$e_j^*(y_{k_\ell})=0$ dla dużych $\ell$, więc $e_j^*(y)=0$.
Jednoznaczność rozwinięcia wymusza $y=0$,
co jest sprzeczne z $|f(y_{k_\ell})|\ge\varepsilon$.
Ponieważ $(I-P_N)X\subseteq Y_N$ i $\|I-P_N\|\le1+K_e$,
dostajemy $\|f-P_N'f\|\le(1+K_e)\|f|_{Y_N}\|\to0$.
Funkcjonał $P_N'f=f\circ P_N$ jest skończoną kombinacją
$e_1^*,\ldots,e_N^*$; baza jest zatem ściągająca.

Odwrotnie, niech baza ma obie własności i niech
$\Phi\in X''$. Kładziemy $a_n=\Phi(e_n^*)$ oraz
$s_N=\sum_{n=1}^Na_ne_n$. Dla $f\in X'$ mamy
$f(s_N)=\Phi(P_N'f)$, a więc
$|f(s_N)|\le K_e\|\Phi\|\,\|f\|$.
Z definicji normy wynika $\|s_N\|\le K_e\|\Phi\|$.
Zupełność w sensie ograniczonym daje
$x=\lim_Ns_N\in X$. Dla każdego $n$ otrzymujemy
$e_n^*(x)=\Phi(e_n^*)$.
Zatem $J_Xx$ i $\Phi$ pokrywają się na liniowej otoczce
$\{e_n^*\}$, która jest gęsta w $X'$ dzięki ściąganiu bazy.
Ciągłość obu funkcjonałów na $X'$ daje $J_Xx=\Phi$.
Kanoniczne odwzorowanie $J_X:X\to X''$ jest surjektywne,
czyli $X$ jest refleksywna.
\end{proof}

\begin{przyklad}[Przykłady rozdzielające oba warunki]
\label{ex:james-examples27}
Standardowa baza przestrzeni $c_0$ jest ściągająca,
gdyż $c_0'=\ell^1$, a jej funkcjonały współrzędnych
odpowiadają wektorom standardowym w $\ell^1$.
Nie jest zupełna w sensie ograniczonym: sumy
$\sum_{n=1}^Ne_n$ mają normę $1$ w $c_0$, ale
nie tworzą ciągu Cauchy'ego.
Standardowa baza przestrzeni $\ell^1$ jest zupełna w sensie
ograniczonym, ponieważ ograniczoność sum częściowych oznacza
$\sup_N\sum_{n=1}^N|a_n|<\infty$.
Nie jest ściągająca, gdyż $(\ell^1)'=\ell^\infty$,
a domknięcie liniowej otoczki funkcjonałów współrzędnych
to właściwa podprzestrzeń $c_0$ przestrzeni $\ell^\infty$.
\end{przyklad}

\subsection{Czego istnienie bazy nie gwarantuje?}

\begin{uwaga}[Baza ortonormalna a baza Hamela]
\label{rem:hamel-vs-schauder27}
Jeżeli przestrzeń Hilberta jest ośrodkowa i nieskończenie
wymiarowa, to przeliczalna baza ortonormalna jest bazą
Schaudera, lecz nigdy bazą Hamela całej przestrzeni:
wektory z bazy Hamela dają skończone kombinacje liniowe.
Co więcej, dla nieskończenie wymiarowej przestrzeni Hilberta
każda baza Hamela jest nieprzeliczalna. Dowody tych faktów,
a także przeliczalność niezerowych współczynników
względem dowolnej, nawet nieprzeliczalnej rodziny
ortonormalnej, znajdują się we wcześniejszym rozdziale o przestrzeniach Hilberta.
\end{uwaga}

\begin{uwaga}[Ośrodkowość nie wystarcza]
\label{rem:enflo27}
Istnieją ośrodkowe przestrzenie Banacha bez bazy Schaudera.
Wynik ten ma charakter głęboki i nie wynika z samej
nierówności Banacha--Steinhausa. Konstrukcje tego typu
wiążą się z niepowodzeniem własności aproksymacji,
którą omówimy w następnym rozdziale. Istnienie
\emph{ciągu bazowego} w każdej nieskończenie wymiarowej
przestrzeni Banacha nie oznacza, że jest on bazą całej
przestrzeni.
\end{uwaga}


% Wersja źródłowa: audyt II, jednostka 28
\setcounter{section}{27}
\noindent\textbf{Wyniki wykorzystywane.}
Przyjmujemy twierdzenia Hahna--Banacha i Banacha--Steinhausa,
teorię baz Schaudera i norm operatorowych, podstawy przestrzeni
Hilberta oraz kryterium ośrodkowości. Wszystkie nowe konstrukcje
i ich zależności są formułowane i dowodzone w tym rozdziale.
\par\medskip
\section{Własności aproksymacji przestrzeni Banacha}
\label{sec:aproksymacja28}

W całym rozdziale $X,Y,Z$ są przestrzeniami Banacha nad
$\K\in\{\R,\C\}$. Przez $\mathcal L(X,Y)$ oznaczamy
przestrzeń operatorów liniowych ciągłych z normą operatorową, zaś
$\mathcal F(X,Y)$ -- zbiór tych $T\in\mathcal L(X,Y)$, których
obraz ma skończony wymiar. Wyrażenie \emph{operator skończonego rzędu}
nie oznacza operatora będącego jedynie odwzorowaniem algebraicznym:
zawsze wymagana jest jego ciągłość. W rozdziale o bazach Schaudera
wykazano, że rzuty częściowe takiej bazy są jednostajnie ograniczonymi
operatorami skończonego rzędu. Tutaj wyjaśnimy, jaka własność
przestrzeni za tym stoi i dlaczego nie każda przestrzeń Banacha ją ma.

\subsection{Aproksymowanie jednostajnie na zbiorach zwartych}

\begin{definicja}[AP, BAP i MAP]
\label{def:ap-bap-map28}
Niech $X$ będzie przestrzenią Banacha, a $\lambda\ge1$.
Mówimy, że $X$ ma \emph{własność aproksymacji} (AP), jeżeli dla
każdego zwartego $K\subseteq X$ i każdego $\varepsilon>0$
istnieje $T\in\mathcal F(X,X)$ spełniający
$\sup_{x\in K}\|Tx-x\|<\varepsilon$.
Mówimy, że $X$ ma \emph{$\lambda$-ograniczoną własność aproksymacji}
($\lambda$-BAP), jeżeli można ponadto zawsze wybrać taki $T$
ze $\|T\|\le\lambda$. Własność BAP oznacza istnienie
$\lambda<\infty$ spełniającego ten warunek, natomiast
\emph{metryczna własność aproksymacji} (MAP) jest równoznaczna
z $1$-BAP.
\end{definicja}

Przestrzenie skończenie wymiarowe mają MAP: wystarczy wziąć
$T=I_X$. W przestrzeniach nieskończenie wymiarowych wymaganie
zbieżności w normie operatorowej byłoby zbyt silne.

\begin{stwierdzenie}[Brak aproksymacji identyczności w normie operatorowej]
\label{prop:no-norm-identity28}
Jeżeli $X$ jest nieskończenie wymiarową przestrzenią Banacha,
to dla każdego $T\in\mathcal F(X,X)$ zachodzi
$\|I_X-T\|\ge1$.
\end{stwierdzenie}
\begin{proof}
Ponieważ $T(X)$ ma wymiar skończony, $T$ nie może być
różnowartościowy na nieskończenie wymiarowej przestrzeni $X$.
Wybieramy $x\in\ker T$ o $\|x\|=1$. Wtedy
$\|(I_X-T)x\|=1$, co daje tezę.
\end{proof}

\begin{lemat}[Kryterium BAP na zbiorach skończonych]
\label{lem:finite-bap28}
Niech $X$ będzie przestrzenią Banacha, $\lambda\ge1$.
Następujące warunki są równoważne:
\begin{enumerate}[label=\textup{(\roman*)}]
\item $X$ ma $\lambda$-BAP;
\item dla wszystkich $m\ge1$, $x_1,\ldots,x_m\in X$ i
$\varepsilon>0$ istnieje $T\in\mathcal F(X,X)$ takie,
że $\|T\|\le\lambda$ oraz
$\max_{1\le j\le m}\|Tx_j-x_j\|<\varepsilon$.
\end{enumerate}
\end{lemat}
\begin{proof}
Pierwszy warunek implikuje drugi, ponieważ zbiór skończony
jest zwarty. W przeciwną stronę ustalmy zwarty $K\subset X$
i $\varepsilon>0$. Dla $K=\varnothing$ teza jest oczywista.
W przeciwnym razie istnieją $x_1,\ldots,x_m\in K$
takie, że dla każdego $x\in K$ pewien $j$ spełnia
$\|x-x_j\|<\varepsilon/[3(1+\lambda)]$.
Wybierzmy $T$ jak w (ii) z błędem mniejszym niż
$\varepsilon/3$. Wówczas
\[
\|Tx-x\|\le\|T(x-x_j)\|+\|Tx_j-x_j\|+\|x_j-x\|
<\varepsilon.
\]
W dowodzie korzystamy jednocześnie ze skończoności siatki
$\varepsilon$ i ze \emph{wspólnej granicy norm} aproksymantów.
\end{proof}

\begin{wniosek}[Hierarchia własności]
\label{cor:hierarchy28}
Dla każdej przestrzeni Banacha zachodzi
$\mathrm{MAP}\Rightarrow\mathrm{BAP}\Rightarrow\mathrm{AP}$.
Jeśli $X$ ma $\lambda$-BAP, to ma też $\mu$-BAP
dla każdego $\mu\ge\lambda$.
\end{wniosek}
\begin{proof}
Dwie pierwsze implikacje i monotoniczność stałej wynikają
bezpośrednio z Definicji~\ref{def:ap-bap-map28}.
\end{proof}

\begin{uwaga}[Dlaczego AP nie definiuje się na zbiorach skończonych?]
\label{rem:finite-exact28}
Dla \emph{dowolnej} przestrzeni Banacha $X$ i skończonego
zbioru $\{x_1,\ldots,x_m\}\subset X$ można znaleźć
$T\in\mathcal F(X,X)$ takie, że $Tx_j=x_j$ dla każdego $j$.
Istotnie, niech $E=\operatorname{span}\{x_1,\ldots,x_m\}$
i wybierzmy bazę $u_1,\ldots,u_r$ przestrzeni $E$.
Funkcjonały współrzędnych na $E$ są ciągłe, a twierdzenie
Hahna--Banacha przedłuża je do $f_1,\ldots,f_r\in X'$.
Operator $Tx=\sum_{j=1}^rf_j(x)u_j$ jest ciągły,
skończonego rzędu i jest identycznością na $E$.
Jego norma może jednak zależeć w sposób niekontrolowany
od wyboru $E$. To właśnie odróżnia powyższy fakt od BAP.
\end{uwaga}

\subsection{Topologia zbieżności na zbiorach zwartych i granice operatorów}

\begin{definicja}[Zbieżność zwarta i silna operatorów]
\label{def:topologies-ap28}
Niech $X,Y$ będą przestrzeniami Banacha. Topologia
\emph{zbieżności jednostajnej na zbiorach zwartych}
na $\mathcal L(X,Y)$ jest wyznaczona przez seminormy
$p_K(T)=\sup_{x\in K}\|Tx\|$, gdzie $K$ przebiega
zbiory zwarte w $X$. \emph{Silna topologia operatorowa}
(SOT) jest wyznaczona przez seminormy $p_x(T)=\|Tx\|$
indeksowane przez $x\in X$. W obu przypadkach
zbieżność rozumiemy także dla sieci, a nie tylko ciągów.
\end{definicja}

\begin{twierdzenie}[Operatorowa charakterystyka AP i BAP]
\label{thm:operator-character-ap28}
Niech $X$ będzie przestrzenią Banacha i $\lambda\ge1$.
\begin{enumerate}[label=\textup{(\roman*)}]
\item $X$ ma AP wtedy i tylko wtedy, gdy
$I_X$ należy do domknięcia $\mathcal F(X,X)$
w topologii zbieżności jednostajnej na zbiorach zwartych.
\item $X$ ma $\lambda$-BAP wtedy i tylko wtedy, gdy
$I_X$ należy do domknięcia zbioru
$\{T\in\mathcal F(X,X):\|T\|\le\lambda\}$ w SOT.
Dla tego ograniczonego zbioru aproksymantów można
równoważnie użyć topologii zbieżności zwartej.
\end{enumerate}
\end{twierdzenie}
\begin{proof}
Pierwszy punkt jest dokładnym przepisaniem definicji AP
na język otoczeń w topologii $p_K$; skończone przecięcia
otoczeń odpowiadają sumie skończonej liczby zbiorów zwartych.
Pierwsza część (ii) wynika z Lematu~\ref{lem:finite-bap28},
ponieważ otoczenia w SOT wyznaczają warunki na skończenie
wielu wektorach. Na zbiorze operatorów o normach
nieprzekraczających $\lambda$ zbieżność punktowa do
$I_X$ implikuje jednostajną aproksymację na każdym
ustalonym zbiorze zwartym: stosujemy argument ze
skończoną siatką z dowodu Lematu~\ref{lem:finite-bap28}.
Implikacja odwrotna jest oczywista.
\end{proof}

\begin{stwierdzenie}[Silna granica sieci skończonego rzędu istnieje zawsze]
\label{prop:universal-sot28}
Dla \emph{każdej} przestrzeni Banacha $X$ istnieje sieć
$(T_E)_{E\in\mathfrak F(X)}$ operatorów skończonego rzędu,
która zbiega w SOT do $I_X$. Tutaj $\mathfrak F(X)$
oznacza zbiór wszystkich skończenie wymiarowych podprzestrzeni
$E\subseteq X$, skierowany inkluzją.
\end{stwierdzenie}
\begin{proof}
Dla każdego $E\in\mathfrak F(X)$ wybierzmy
ciągły rzut $T_E:X\to E$, konstruowany jak
w Uwadze~\ref{rem:finite-exact28}. Jeżeli $x\in X$,
to dla wszystkich $E\supseteq\operatorname{span}\{x\}$
zachodzi $T_Ex=x$. Oznacza to $T_Ex\to x$
w normie dla każdego $x\in X$, czyli zbieżność w SOT.
\end{proof}

\begin{uwaga}[Niebezpieczne utożsamienie AP ze zbieżnością punktową]
\label{rem:net-vs-ap28}
Twierdzenie~\ref{prop:universal-sot28} nie daje AP:
konstruowana sieć może nie mieć \emph{wspólnego ograniczenia}
na normy operatorowe, a nawet nie musi być punktowo
ograniczona jako cała rodzina. Zbieżność sieci nie implikuje
ograniczoności całego jej zbioru wartości. Dlatego
zastosowanie do takiej sieci twierdzenia Banacha--Steinhausa
bez dodatkowych przesłanek byłoby błędem.
Istnieją przestrzenie Banacha bez AP, mimo że również
w nich zachodzi Twierdzenie~\ref{prop:universal-sot28}.
\end{uwaga}

\begin{twierdzenie}[Ciągi aproksymantów i ośrodkowość]
\label{thm:sequence-bap28}
Niech $X$ będzie ośrodkową przestrzenią Banacha.
Wówczas $X$ ma BAP wtedy i tylko wtedy, gdy istnieje
\emph{ciąg} $(T_n)_{n\ge1}\subset\mathcal F(X,X)$
taki, że $T_nx\to x$ dla każdego $x\in X$.
Każda przestrzeń Banacha, dla której istnieje taki ciąg,
jest ośrodkowa, bez żadnego założenia ośrodkowości z góry.
\end{twierdzenie}
\begin{proof}
Załóżmy $\lambda$-BAP i ośrodkowość.
Niech $(x_j)_{j\ge1}$ będzie gęstym ciągiem w $X$.
Na mocy Lematu~\ref{lem:finite-bap28} dla każdego $n$
istnieje $T_n\in\mathcal F(X,X)$ o normie nie większej
niż $\lambda$ i takim, że $\|T_nx_j-x_j\|<1/n$
dla $1\le j\le n$. Zatem $T_nx_j\to x_j$ dla każdego $j$.
Dla dowolnego $x\in X$ i $x_j$ bliskiego $x$ mamy
$\|T_nx-x\|\le(\lambda+1)\|x-x_j\|+\|T_nx_j-x_j\|$,
co daje $T_nx\to x$.
Odwrotnie, jeżeli $T_nx\to x$ dla każdego $x$, to rodzina
$\{T_n:n\ge1\}$ jest punktowo ograniczona. Twierdzenie
Banacha--Steinhausa daje $\sup_n\|T_n\|<\infty$.
Przybliżenie na skończonych zbiorach, a następnie
Lemat~\ref{lem:finite-bap28}, zapewniają BAP.
Ponadto $\bigcup_nT_n(X)$ leży w przeliczalnej sumie
skończenie wymiarowych podprzestrzeni. Jej liniowa otoczka
ma przeliczalny zbiór gęsty, a zawiera gęsty zbiór
$\{\lim_nT_nx:x\in X\}=X$ w swoim domknięciu.
Wobec tego $X$ jest ośrodkowa.
\end{proof}

\subsection{Bazy Schaudera a metryczna własność aproksymacji}

\begin{twierdzenie}[Baza Schaudera implikuje BAP]
\label{thm:basis-implies-bap28}
Niech $X$ będzie przestrzenią Banacha z bazą Schaudera
$(e_n)_{n\ge1}$ oraz rzutami częściowymi $P_N$.
Jeżeli $K_e=\sup_N\|P_N\|$, to $X$ ma $K_e$-BAP.
W szczególności baza monotoniczna daje MAP.
\end{twierdzenie}
\begin{proof}
Każdy $P_N$ ma skończony rząd i spełnia
$\|P_N\|\le K_e$, a definicja bazy daje
$P_Nx\to x$ w normie dla każdego $x\in X$.
Lemat~\ref{lem:finite-bap28} wraz z oszacowaniem
ze skończoną siatką dowodzi $K_e$-BAP.
Jeżeli baza jest monotoniczna, $K_e=1$.
\end{proof}

\begin{przyklad}[Przestrzenie ciągowe]
\label{ex:sequence-map28}
Przestrzenie Banacha $c_0$ i $\ell^p$ dla
$1\le p<\infty$ mają MAP: rzuty obcinające ciąg
po $N$ wyrazach mają normę $1$ i zbieżność punktową
do identyczności udowodniono już w rozdziale o bazach Schaudera.
Nie powtarzamy tamtego dowodu. Inaczej wygląda sytuacja
w przestrzeni $\ell^\infty$, mimo że i ona ma MAP.
\end{przyklad}

\begin{twierdzenie}[MAP w dowolnej przestrzeni Hilberta]
\label{thm:hilbert-map28}
Każda rzeczywista lub zespolona przestrzeń Hilberta $H$,
bez założenia ośrodkowości, ma MAP.
\end{twierdzenie}
\begin{proof}
Niech $K\subseteq H$ będzie zwarty i $\varepsilon>0$.
Dla $K=\varnothing$ nie ma nic do dowodzenia.
W przeciwnym razie wybieramy skończoną $\varepsilon/2$-siatkę
$\{x_1,\ldots,x_m\}\subset K$ i oznaczamy
$E=\operatorname{span}\{x_1,\ldots,x_m\}$.
Niech $P_E$ będzie rzutem ortogonalnym na $E$.
Jeśli $E\ne\{0\}$, to $\|P_E\|=1$, a w przeciwnym
razie $\|P_E\|=0$. Ponieważ $P_Ex_j=x_j$,
dla $x\in K$ i odpowiedniego $j$ zachodzi
$\|x-P_Ex\|\le\|x-x_j\|+\|P_E(x_j-x)\|
<\varepsilon$. Operator $P_E$ ma skończony rząd.
\end{proof}

\begin{twierdzenie}[MAP przestrzeni funkcji ciągłych na odcinku]
\label{thm:c01-map28}
Przestrzeń Banacha $C([0,1],\K)$ z normą supremum
ma MAP. Istnieje ciąg operatorów skończonego rzędu
$L_n$ o normach równych $1$, zbieżny punktowo do
identyczności.
\end{twierdzenie}
\begin{proof}
Dla $n\ge1$ przyjmijmy $t_k=k/n$ ($0\le k\le n$).
W szczególności operator $L_n$ zależy liniowo od wartości
$f(t_0),\ldots,f(t_n)$, co jawnie zapewnia skończoność jego rzędu.
Niech $L_nf$ będzie funkcją ciągłą, liniową na każdym
$[t_k,t_{k+1}]$ i spełniającą $L_nf(t_k)=f(t_k)$.
Operator $L_n$ jest liniowy; jego obraz leży w
$(n+1)$-wymiarowej przestrzeni funkcji odcinkami liniowych.
Ponieważ $L_nf(t)$ jest kombinacją wypukłą
$f(t_k)$ i $f(t_{k+1})$, mamy
$\|L_nf\|_\infty\le\|f\|_\infty$. Dla funkcji stałej
$1$ zachodzi $L_n1=1$, więc $\|L_n\|=1$.
Z jednostajnej ciągłości $f$ na $[0,1]$ wynika
$\|L_nf-f\|_\infty\to0$. Stosujemy
Twierdzenie~\ref{thm:sequence-bap28} lub bezpośrednio
Lemat~\ref{lem:finite-bap28} z $\lambda=1$.
\end{proof}

\begin{twierdzenie}[MAP bez bazy Schaudera]
\label{thm:linfty-map28}
Przestrzeń Banacha $\ell^\infty(\N,\K)$ ma MAP,
lecz nie ma żadnej bazy Schaudera.
\end{twierdzenie}
\begin{proof}
Ustalmy $x^{(1)},\ldots,x^{(m)}\in\ell^\infty$ oraz
$\varepsilon>0$. Zbiór
$A=\{(x^{(1)}_n,\ldots,x^{(m)}_n):n\in\N\}\subset\K^m$
jest ograniczony, a więc ma skończoną
$\varepsilon/2$-siatkę. Przypisujemy każdemu indeksowi $n$ pierwszy z wybranych
punktów siatki leżący w odległości mniejszej niż $\varepsilon/2$
od odpowiadającego mu wektora. Otrzymujemy podział $\N$
na skończenie wiele niepustych zbiorów $A_1,\ldots,A_r$ tak,
że dla $n,k\in A_j$ zachodzi
$|x_n^{(s)}-x_k^{(s)}|<\varepsilon$
dla każdego $1\le s\le m$ (na mocy nierówności trójkąta).
Wybierzmy po jednym $n_j\in A_j$ i połóżmy
$(Ty)_n=y_{n_j}$, jeżeli $n\in A_j$.
Jest to liniowy operator skończonego rzędu,
$\|T\|=1$ i
$\|Tx^{(s)}-x^{(s)}\|_\infty<\varepsilon$ dla każdego $s$.
Lemat~\ref{lem:finite-bap28} dowodzi MAP.

Przestrzeń $\ell^\infty$ nie jest ośrodkowa:
nieprzeliczalna rodzina ciągów $\chi_A=(\chi_A(n))_{n\ge1}$,
gdzie $A$ przebiega wszystkie podzbiory $\N$,
ma parami odległość $1$ w normie supremum.
Tymczasem każda przestrzeń Banacha z bazą Schaudera
jest ośrodkowa. Zatem $\ell^\infty$ nie ma takiej bazy.
\end{proof}

\begin{uwaga}[Brak uniwersalnego ciągu w przestrzeni nieośrodkowej]
\label{rem:linfty-sequence28}
Twierdzenia~\ref{thm:linfty-map28} oraz
\ref{thm:sequence-bap28} pokazują jednocześnie, że
w $\ell^\infty$ nie istnieje \emph{ciąg} operatorów
skończonego rzędu zbieżny silnie do $I$.
Istnieje natomiast jednostajnie ograniczona \emph{sieć}
takich operatorów zbieżna na wszystkich zbiorach zwartych.
Jest to powód, dla którego definicje AP i BAP nie powinny
opierać się wyłącznie na ciągach.
\end{uwaga}

\subsection{Aproksymowanie operatorów zwartych}

\begin{definicja}[Operatory zwarte i aproksymowalne]
\label{def:compact-approximable28}
Niech $X,Y$ będą przestrzeniami Banacha.
Operator $T\in\mathcal L(X,Y)$ jest \emph{zwarty},
jeżeli $T(B_X)$ ma zwarte domknięcie w $Y$, gdzie
$B_X=\{x\in X:\|x\|\le1\}$.
Oznaczamy zbiór takich operatorów przez
$\mathcal K(X,Y)$. Operator jest \emph{aproksymowalny},
jeżeli należy do domknięcia $\mathcal F(X,Y)$
w normie operatorowej; ich zbiór oznaczamy
$\mathcal A(X,Y)$.
\end{definicja}

\begin{twierdzenie}[Operatory aproksymowalne są zwarte]
\label{thm:approximable-compact28}
Dla wszystkich przestrzeni Banacha $X,Y$ zachodzi
$\mathcal F(X,Y)\subseteq\mathcal A(X,Y)\subseteq
\mathcal K(X,Y)\subseteq\mathcal L(X,Y)$.
Ponadto $\mathcal K(X,Y)$ jest domkniętą podprzestrzenią
liniową $\mathcal L(X,Y)$.
\end{twierdzenie}
\begin{proof}
Każdy operator skończonego rzędu jest zwarty, ponieważ
jego ograniczony obraz kuli jednostkowej ma zwarte
domknięcie w skończenie wymiarowej przestrzeni.
Jeżeli $T_n\to T$ w normie operatorowej i każdy $T_n$
jest zwarty, to dla $\varepsilon>0$ wybieramy $n$ z
$\|T-T_n\|<\varepsilon/3$. Zbiór $T_n(B_X)$ jest
całkowicie ograniczony, więc pokrywają go skończenie
wiele kul promienia $\varepsilon/3$.
Otrzymujemy skończone pokrycie $T(B_X)$ kulami
promienia $\varepsilon$, a zatem całkowitą ograniczoność.
Ponieważ $Y$ jest zupełna, $\overline{T(B_X)}$ jest zwarty.
To pokazuje domkniętość zbioru operatorów zwartych.
Liniowość wynika ze zwartości sum i skalarnych
wielokrotności zbiorów względnie zwartych.
\end{proof}

\begin{twierdzenie}[AP w przestrzeni wartości i aproksymacja operatorów]
\label{thm:ap-compact-operators28}
Niech $X,Y$ będą przestrzeniami Banacha, przy czym
$Y$ ma AP. Wówczas każdy operator zwarty
$T\in\mathcal K(X,Y)$ jest aproksymowalny, tzn.
$\mathcal K(X,Y)=\mathcal A(X,Y)$.
Jeżeli $Y$ ma $\lambda$-BAP, to dla każdego
$T\in\mathcal K(X,Y)$ i każdego $\varepsilon>0$
istnieje $F\in\mathcal F(X,Y)$ z
$\|F-T\|<\varepsilon$ oraz $\|F\|\le\lambda\|T\|$.
\end{twierdzenie}
\begin{proof}
Niech $K=\overline{T(B_X)}$; z założenia jest zwarty.
Z AP przestrzeni $Y$ otrzymujemy dla każdego
$\varepsilon>0$ operator skończonego rzędu
$S:Y\to Y$ taki, że
$\sup_{y\in K}\|Sy-y\|<\varepsilon$.
Operator $F=S\circ T$ ma skończony rząd i
$\|F-T\|=\sup_{x\in B_X}\|S(Tx)-Tx\|<\varepsilon$.
Przy $\lambda$-BAP można ponadto wybrać $\|S\|\le\lambda$,
stąd $\|F\|\le\lambda\|T\|$.
\end{proof}

\begin{uwaga}[Właściwy kierunek w Twierdzeniu~\ref{thm:ap-compact-operators28}]
\label{rem:range-not-domain28}
Założenie AP dotyczy tutaj \emph{przestrzeni wartości $Y$},
a nie automatycznie przestrzeni wyjściowej $X$:
aproksymujemy zwarty zbiór $T(B_X)\subseteq Y$.
Utożsamienie tych dwóch ról bez dodatkowego twierdzenia
o dualności byłoby nieuzasadnione.
\end{uwaga}

\subsection{Trwałość własności i ograniczenia twierdzeń}

\begin{twierdzenie}[Trwałość dla podprzestrzeni dopełnionych]
\label{thm:complemented28}
Niech $X$ będzie przestrzenią Banacha i $Y\subseteq X$
domkniętą podprzestrzenią. Załóżmy, że istnieje
ciągły rzut $Q:X\to Y$, tj. $Q|_Y=I_Y$.
Jeżeli $X$ ma AP, to $Y$ ma AP.
Jeżeli $X$ ma $\lambda$-BAP, to $Y$ ma
$\mu$-BAP dla $\mu=\max\{1,\lambda\|Q\|\}$.
W szczególności, gdy $X$ ma MAP i $\|Q\|=1$,
to $Y$ ma MAP.
\end{twierdzenie}
\begin{proof}
Ustalmy zwarty $K\subseteq Y$. Jest on również
zwarty w $X$. Wybierzmy operator skończonego rzędu
$T:X\to X$ przybliżający identyczność jednostajnie
na $K$ i połóżmy $S=Q\circ T|_Y:Y\to Y$.
Operator $S$ ma skończony rząd, a dla $y\in K$
$\|Sy-y\|=\|Q(Ty-y)\|\le\|Q\|\,\|Ty-y\|$.
Ponieważ dokładność aproksymacji $T$ możemy wybrać
dowolnie dużą, otrzymujemy AP na $Y$.
Przy $\lambda$-BAP zachodzi dodatkowo
$\|S\|\le\|Q\|\lambda$, co dowodzi drugiej części.
\end{proof}

\begin{twierdzenie}[Niezmienniczość AP i BAP przy izomorfizmach]
\label{thm:isomorphic28}
Niech $X,Y$ będą przestrzeniami Banacha i niech
$U:X\to Y$ będzie liniowym izomorfizmem topologicznym.
Wtedy $X$ ma AP wtedy i tylko wtedy, gdy $Y$ ma AP.
Jeżeli $X$ ma $\lambda$-BAP, to $Y$ ma
$\mu$-BAP z $\mu=\max\{1,\lambda\|U\|\|U^{-1}\|\}$.
W szczególności istnienie BAP jest niezmiennicze
względem równoważnych norm, natomiast liczba $\lambda$
może się zmienić.
\end{twierdzenie}
\begin{proof}
Dla zwartego $K\subseteq Y$ jego przeciwobraz
$U^{-1}(K)$ jest zwarty w $X$. Niech
$T\in\mathcal F(X,X)$ przybliża identyczność na
$U^{-1}(K)$. Operator $S=UTU^{-1}\in\mathcal F(Y,Y)$
spełnia $\|Sy-y\|\le\|U\|\,\|TU^{-1}y-U^{-1}y\|$.
W przypadku BAP mamy dodatkowo
$\|S\|\le\|U\|\|U^{-1}\|\lambda$.
Argument jest symetryczny względem $X,Y$.
\end{proof}

\begin{uwaga}[Granice implikacji: kontrprzykłady głębokie]
\label{rem:enflo-figiel28}
Nie każda ośrodkowa, a nawet nie każda \emph{refleksywna}
przestrzeń Banacha ma AP: istnieje taki kontrprzykład
Enflo. Istnieją też przestrzenie Banacha mające AP,
ale niemające BAP (konstrukcja Figiela--Johnsona),
oraz mające BAP, lecz nie MAP. Klasyczne źródła
dwóch pierwszych kontrprzykładów to prace Enflo
(\emph{Acta Mathematica} 130, 1973, s.~309--317) oraz
Figiela i Johnsona (\emph{Proceedings of the AMS} 41,
1973, s.~197--200). Wyniki te \emph{przytaczamy bez dowodu}, ponieważ szczegółowe
konstrukcje wykraczają poza zakres rozdziału.
W przeciwieństwie do tych rezultatów kontrprzykład
``MAP nie implikuje istnienia bazy Schaudera'' z
Twierdzenia~\ref{thm:linfty-map28} jest całkowicie
elementarny i został udowodniony powyżej.
\end{uwaga}

\begin{uwaga}[Co naprawdę daje istnienie bazy?]
\label{rem:recap28}
Zestawiając Twierdzenia~\ref{thm:basis-implies-bap28}
i~\ref{thm:linfty-map28}, otrzymujemy ściśle
jednokierunkową zależność między własnościami bazowymi
i aproksymacyjnymi. Z drugiej strony nie wolno
zastępować \emph{jednostajnej} aproksymacji na
zbiorach zwartych dowolną aproksymacją punktową przez
sieć operatorów skończonego rzędu: ta ostatnia zachodzi
w każdej przestrzeni Banacha na mocy
Twierdzenia~\ref{prop:universal-sot28}.
\end{uwaga}


% Wersja źródłowa: audyt II, jednostka 29
\setcounter{section}{28}
\noindent\textbf{Wyniki wykorzystywane.} Przyjmujemy znane podstawy teorii przestrzeni Banacha, twierdzenie Hahna--Banacha, twierdzenie Banacha--Steinhausa, zwartość słabą kuli w refleksywnej przestrzeni Banacha, twierdzenie Eberleina--\v Smuliana, twierdzenie o odwzorowaniu otwartym oraz charakteryzację operatorów zwartych z poprzedniego rozdziału.
\par\medskip
\section{Operatory ograniczone i zwarte: kryteria, transpozycja i ideały}
\label{sec:operatory29}

W rozdziale przyjmujemy, że $X,Y,Z$ są przestrzeniami Banacha nad
$\K\in\{\R,\C\}$. Zapis $\mathcal L(X,Y)$ oznacza przestrzeń
ograniczonych operatorów liniowych, $\mathcal F(X,Y)$ -- operatory
ciągłe skończonego rzędu, a $\mathcal K(X,Y)$ -- operatory zwarte.
Te pojęcia, a także domkniętość $\mathcal K(X,Y)$ w normie operatorowej,
ustanowiono w rozdziale o własnościach aproksymacji. Nie powtarzamy
tamtych dowodów. Dla operatora $T:X\to Y$ przez $T':Y'\to X'$
oznaczamy jego operator sprzężony (transpozycję), dany wzorem
$(T'y')(x)=y'(Tx)$. Normy operatora i jego transpozycji są równe:
$\|T'\|=\|T\|$; wynika to z charakteryzacji normy przez
funkcjonały liniowe ciągłe.

\subsection{Zwartość i kryteria ciągowe}

\begin{twierdzenie}[Trzy kryteria zwartości]
\label{thm:compact-criteria29}
Dla $T\in\mathcal L(X,Y)$ następujące warunki są równoważne:
\begin{enumerate}[label=\textup{(\roman*)}]
\item $T$ jest zwarty, tj. $\overline{T(B_X)}$ jest zwarty w $Y$;
\item dla każdego ograniczonego ciągu $(x_n)$ w $X$ ciąg $(Tx_n)$
      ma podciąg zbieżny w normie $Y$;
\item dla każdego $\varepsilon>0$ zbiór $T(B_X)$ ma skończoną
      $\varepsilon$-siatkę w $Y$.
\end{enumerate}
\end{twierdzenie}
\begin{proof}
W przestrzeni metrycznej zwartość jest równoważna zwartości ciągowej.
Zbiór $\overline{T(B_X)}$ jest więc zwarty dokładnie wtedy,
gdy każdy ciąg z tego domknięcia ma podciąg zbieżny do punktu
domknięcia. W dowodzie wystarczy sprawdzić ciągi z $T(B_X)$:
każdy punkt domknięcia można przybliżyć punktami tego zbioru
z błędem mniejszym niż $1/n$.
Dowolny ograniczony ciąg $(x_n)$ mieści się w $rB_X$ dla pewnego
$r>0$, więc (i) implikuje (ii). Dla implikacji odwrotnej stosujemy
(ii) do ciągów w $B_X$; własność podciągu zbieżnego daje całkowitą
ograniczoność $T(B_X)$ (w przeciwnym razie można indukcyjnie
wybrać punkty odległe parami co najmniej $\varepsilon$).
Przestrzeń $Y$ jest zupełna, zatem domknięcie zbioru całkowicie
ograniczonego jest zwarte. Tym samym otrzymujemy (iii) i (i).
\end{proof}

\begin{lemat}[Lemat Riesza o prawie prostopadłym wektorze]
\label{lem:riesz29}
Niech $M$ będzie właściwą domkniętą podprzestrzenią przestrzeni
unormowanej $X$, a $0<\theta<1$. Istnieje $x\in X$ takie, że
$\|x\|=1$ oraz $\operatorname{dist}(x,M)>\theta$.
\end{lemat}
\begin{proof}
Wybieramy $z\in X\setminus M$. Ponieważ $M$ jest domknięta,
$d:=\operatorname{dist}(z,M)>0$. Wybieramy $y\in M$ z
$\|z-y\|<d/\theta$ i kładziemy $x=(z-y)/\|z-y\|$.
Dla każdego $m\in M$ zachodzi
$\|x-m\|=\|z-y-\|z-y\|m\|/\|z-y\|\ge d/\|z-y\|>\theta$.
\end{proof}

\begin{twierdzenie}[Zwartość identyczności]
\label{thm:identity-compact29}
Dla przestrzeni Banacha $X$ operator $I_X$ jest zwarty wtedy
i tylko wtedy, gdy $X$ ma skończony wymiar.
\end{twierdzenie}
\begin{proof}
W skończonym wymiarze ograniczone zbiory mają zwarte domknięcie.
Gdy $X$ ma nieskończony wymiar, wybieramy indukcyjnie wektory
$x_n\in B_X$ o normie $1$ w taki sposób, że
$\operatorname{dist}(x_{n+1},\operatorname{span}\{x_1,\ldots,x_n\})>1/2$.
Lemat~\ref{lem:riesz29} pozwala wykonać każdy krok, gdyż
podprzestrzeń skończenie wymiarowa jest domknięta i właściwa.
Wówczas $\|x_n-x_m\|>1/2$ dla $n\ne m$, więc kula jednostkowa
nie jest względnie zwarta.
\end{proof}

\begin{twierdzenie}[Zwartość a składanie operatorów]
\label{thm:ideal-compact29}
Jeżeli $S\in\mathcal L(Y,Z)$ i $T\in\mathcal L(X,Y)$, to:
\begin{enumerate}[label=\textup{(\roman*)}]
\item $T\in\mathcal K(X,Y)$ implikuje
      $ST\in\mathcal K(X,Z)$;
\item $S\in\mathcal K(Y,Z)$ implikuje
      $ST\in\mathcal K(X,Z)$.
\end{enumerate}
Ponadto $\mathcal K(X,Y)$ jest podprzestrzenią liniową zamkniętą
w normie operatorowej.
\end{twierdzenie}
\begin{proof}
W (i) zbiór $T(B_X)$ jest względnie zwarty, a ciągłość $S$
zapewnia względną zwartość jego obrazu. W (ii) zbiór $T(B_X)$
jest ograniczony, więc zawiera się w $rB_Y$ dla pewnego $r>0$;
$S(rB_Y)$ jest względnie zwarty. Ostatnie zdanie udowodniono
w rozdziale o własnościach aproksymacji i przywołujemy je tutaj
bez powtarzania dowodu.
\end{proof}

\begin{przyklad}[Zwartość nie wymaga skończonego rzędu]
\label{ex:volterra29}
W przestrzeni Banacha $C([0,1],\K)$ definiujemy operator Volterry
$(Vf)(t)=\int_{0}^{t}f(s)\,\mathrm d\lambda_1(s)$.
Dla $\|f\|_\infty\le1$ mamy $\|Vf\|_\infty\le1$ i
$|(Vf)(t)-(Vf)(u)|\le|t-u|$. Zatem rodzina $V(B_X)$
jest jednostajnie ograniczona i jednakowo ciągła;
z twierdzenia Arzeli--Ascolego ma zwarte domknięcie.
Operator $V$ jest zatem zwarty. Jednocześnie obrazy wielomianów
$f_n(t)=t^n$, równe $t^{n+1}/(n+1)$, są liniowo niezależne,
więc rząd $V$ jest nieskończony. Ponadto $V$ jest różnowartościowy,
ale jego obraz nie jest domknięty: gdyby był domknięty, odwrotność
$V^{-1}:V(X)\to X$ byłaby ciągła z twierdzenia o odwzorowaniu
otwartym, a stąd $I_X=V^{-1}V$ byłaby zwarta, wbrew
Twierdzeniu~\ref{thm:identity-compact29}.
\end{przyklad}

\subsection{Twierdzenie Schaudera o operatorze sprzężonym}

\begin{lemat}[Jednostajna aproksymacja ograniczonych funkcjonałów]
\label{lem:restrictions29}
Niech $K$ będzie zwartym podzbiorem przestrzeni Banacha $Y$.
Rodzina $\{y'|_K:\|y'\|\le1\}$ jest względnie zwarta
w $C(K,\K)$ z normą supremum.
\end{lemat}
\begin{proof}
Dla $k\in K$ mamy $|y'(k)|\le\|k\|$, więc rodzina jest
jednostajnie ograniczona. Ponadto
$|y'(k)-y'(l)|\le\|k-l\|$ dla dowolnych $k,l\in K$:
rodzina jest jednakowo ciągła. Twierdzenie Arzeli--Ascolego
(dla funkcji zespolonych stosujemy je do części rzeczywistych
i urojonych) daje względną zwartość w normie supremum.
\end{proof}

\begin{twierdzenie}[Schauder: zwartość transpozycji]
\label{thm:schauder-adjoint29}
Niech $X,Y$ będą przestrzeniami Banacha nad tym samym ciałem
$\K$, a $T\in\mathcal L(X,Y)$. Wówczas
\[
T\in\mathcal K(X,Y)\quad\Longleftrightarrow\quad
T'\in\mathcal K(Y',X').
\]
\end{twierdzenie}
\begin{proof}
Załóżmy najpierw, że $T$ jest zwarty, i niech
$K=\overline{T(B_X)}\subseteq Y$. Dla $y'_1,y'_2\in B_{Y'}$
z definicji transpozycji dostajemy
\[
\|T'y'_1-T'y'_2\|_{X'}
=\sup_{x\in B_X}|(y'_1-y'_2)(Tx)|
\le\sup_{k\in K}|(y'_1-y'_2)(k)|.
\]
Lemat~\ref{lem:restrictions29} pokazuje, że rodzina
restrykcji $y'|_K$ ma zwarte domknięcie w $C(K,\K)$.
Dla każdego $\varepsilon>0$ wybieramy skończoną
$\varepsilon$-siatkę tej rodziny, z centrami będącymi
restrykcjami pewnych $y'_1,\ldots,y'_m\in B_{Y'}$
(po ewentualnym podwojeniu promienia siatki).
Wówczas $T'y'_1,\ldots,T'y'_m$ tworzą skończoną
siatkę $T'(B_{Y'})$, więc $T'$ jest zwarty.

Jeżeli natomiast $T'$ jest zwarty, to właśnie udowodniona
implikacja, zastosowana do $T':Y'\to X'$, daje zwartość
$T'':X''\to Y''$. Oznaczmy przez $J_X:X\to X''$ i
$J_Y:Y\to Y''$ kanoniczne izometryczne zanurzenia.
Dla każdego $x\in X$ zachodzi tożsamość
$T''J_Xx=J_YTx$. Ponieważ $J_X(B_X)\subseteq B_{X''}$,
zbiór $J_YT(B_X)$ jest względnie zwarty w $Y''$.
Obraz $J_Y(Y)$ jest domknięty (ponieważ $Y$ jest zupełna),
a $J_Y:Y\to J_Y(Y)$ jest izometrią. Stąd $T(B_X)$
jest względnie zwarty w $Y$, co dowodzi tezy.
\end{proof}

\begin{wniosek}[Sprzężenie zachowuje zwartość]
\label{cor:schauder-iterated29}
Dla każdego $T\in\mathcal L(X,Y)$ zwartość $T$ jest
równoważna zwartości $T'$ i $T''$; nie wymagamy
refleksywności żadnej z przestrzeni.
\end{wniosek}
\begin{proof}
Pierwsza równoważność jest treścią
Twierdzenia~\ref{thm:schauder-adjoint29}; druga wynika
z zastosowania tego twierdzenia do $T'$ oraz z tożsamości
$T''J_X=J_YT$.
\end{proof}

\subsection{Zbieżność słaba, operatory całkowicie ciągłe i słabo zwarte}

\begin{twierdzenie}[Operator zwarty na ciągu słabo zbieżnym]
\label{thm:weak-norm29}
Niech $T\in\mathcal K(X,Y)$. Jeżeli ciąg $(x_n)$ zbiega
słabo do $x\in X$, to $\|Tx_n-Tx\|\to0$.
To samo zachodzi dla \emph{ograniczonej} sieci słabo zbieżnej.
\end{twierdzenie}
\begin{proof}
Dla ciągu słabo zbieżnego twierdzenie Banacha--Steinhausa
zastosowane do obrazów $J_Xx_n\in X''$ zapewnia
$\sup_n\|x_n\|<\infty$. Jeżeli $\|Tx_n-Tx\|$ nie zbiega
do zera, można wybrać podciąg, którego odległość od $Tx$
jest ograniczona od dołu. Ze zwartości $T$ mamy podpodciąg
zbieżny w normie do pewnego $y\in Y$. Ponieważ $T$
jest słabo-słabo ciągły, granica słaba tego podpodciągu
musi być równa $Tx$, więc $y=Tx$ -- sprzeczność.
Dla ograniczonej sieci jej obraz jest względnie zwarty
w normie. Gdyby obrazy nie zbiegały do $Tx$, istniałaby
kula $B(Tx,\varepsilon)$, poza którą obrazy leżą
na współkońcowym zbiorze indeksów. Możemy wówczas
wybrać podsieć o wszystkich wyrazach poza tą kulą.
Ze zwartości można wybrać dalszą podsieć normowo zbieżną,
a jej granica słaba musiałaby być równa $Tx$.
\end{proof}

\begin{definicja}[Operator całkowicie ciągły i słabo zwarty]
\label{def:completely-continuous29}
Operator $T\in\mathcal L(X,Y)$ nazywamy
\emph{całkowicie ciągłym}, jeżeli odwzorowuje każdy ciąg
słabo zbieżny w $X$ na ciąg zbieżny w normie $Y$.
Operator $T$ jest \emph{słabo zwarty}, jeżeli
$T(B_X)$ ma słabo zwarte domknięcie w $Y$.
\end{definicja}

\begin{przyklad}[Implikacje są ostre]
\label{ex:weak-vs-compact29}
Każdy operator zwarty jest całkowicie ciągły i słabo zwarty.
Jednak operator identyczności na $\ell^1$ jest całkowicie
ciągły na mocy własności Schura, choć nie jest zwarty,
ponieważ $\ell^1$ ma nieskończony wymiar.
Z drugiej strony naturalne zanurzenie
$J:\ell^1\to\ell^2$, $J(x)=x$, jest ograniczone
i słabo zwarte, gdyż przestrzeń $\ell^2$ jest refleksywna,
ale nie jest zwarte: $Je_n=e_n$ i $(e_n)$ nie ma
podciągu zbieżnego w normie $\ell^2$.
\end{przyklad}

\begin{twierdzenie}[Kryterium w refleksywnej przestrzeni Banacha]
\label{thm:reflexive-complete-cont29}
Jeżeli $X$ jest refleksywną przestrzenią Banacha, to
$T\in\mathcal L(X,Y)$ jest zwarty wtedy i tylko wtedy,
gdy jest całkowicie ciągły.
\end{twierdzenie}
\begin{proof}
Jedna implikacja wynika z
Twierdzenia~\ref{thm:weak-norm29}. Odwrotnie,
niech $(x_n)$ będzie dowolnym ciągiem z $B_X$.
Refleksywność daje słabą zwartość $B_X$, a twierdzenie
Eberleina--\v Smuliana zapewnia podciąg $(x_{n_k})$
słabo zbieżny do pewnego $x\in B_X$.
Całkowita ciągłość $T$ daje normową zbieżność
$Tx_{n_k}\to Tx$. Stosujemy
Twierdzenie~\ref{thm:compact-criteria29}.
\end{proof}

\subsection{Aproksymacja operatorów zwartych i przykłady funkcyjne}

\begin{twierdzenie}[Kryterium ogonów w przestrzeniach ciągowych]
\label{thm:tails29}
Niech $Y=c_0$ lub $Y=\ell^p$ dla $1\le p<\infty$,
a $P_N:Y\to Y$ oznacza obcięcie ciągu po $N$ współrzędnych.
Dla każdego $T\in\mathcal L(X,Y)$ następujące warunki są
równoważne:
\begin{enumerate}[label=\textup{(\roman*)}]
\item $T$ jest zwarty;
\item $\lim_{N\to\infty}\|(I_Y-P_N)T\|=0$;
\item $\lim_{N\to\infty}\sup_{x\in B_X}
      \|(I_Y-P_N)Tx\|=0$.
\end{enumerate}
\end{twierdzenie}
\begin{proof}
Równoważność (ii) i (iii) jest definicją normy operatorowej.
Niech $T$ będzie zwarty i $K=\overline{T(B_X)}$.
Ponieważ $P_Ny\to y$ dla każdego $y\in Y$, a
$\|I_Y-P_N\|\le1$, zbieżność jest jednostajna na $K$:
wybieramy skończoną siatkę w $K$ i stosujemy oszacowanie
$\|(I-P_N)y\|\le\|(I-P_N)(y-y_j)\|+
\|(I-P_N)y_j\|$. Stąd (ii).
W drugą stronę $P_NT$ jest operatorem skończonego rzędu,
a więc jest zwarty. Z założenia $P_NT\to T$ w normie
operatorowej, a domkniętość klasy operatorów zwartych
daje (i).
\end{proof}

\begin{przyklad}[Operatory diagonalne]
\label{ex:diagonal29}
Niech $Y=c_0$ albo $Y=\ell^p$, $1\le p<\infty$,
oraz $a=(a_n)\in\ell^\infty$. Operator
$D_a:Y\to Y$, $(D_ax)_n=a_nx_n$, jest ograniczony
i $\|D_a\|=\|a\|_\infty$. Twierdzenie~\ref{thm:tails29}
pokazuje, że $D_a$ jest zwarty dokładnie wtedy, gdy
$a_n\to0$, ponieważ
$\|(I-P_N)D_a\|=\sup_{n>N}|a_n|$.
Na przykład $D_a$ dla $a_n=1/n$ jest zwartym
operatorem różnowartościowym o niedomkniętym obrazie.
\end{przyklad}

\begin{twierdzenie}[Operatory całkowe o jądrze kwadratowalnym]
\label{thm:hilbert-schmidt29}
Niech $k\in L^2([0,1]^2,\lambda_2)$ i niech
$H=L^2([0,1],\lambda_1)$ będzie przestrzenią Hilberta.
Wówczas wzór
\[
(T_kf)(t)=\int_{[0,1]}k(t,s)f(s)\,\mathrm d\lambda_1(s)
\]
definiuje dla prawie każdego $t$ operator liniowy ciągły
$T_k:H\to H$. Jest on zwarty i spełnia
$\|T_k\|\le\|k\|_{L^2(\lambda_2)}$.
\end{twierdzenie}
\begin{proof}
Z twierdzenia Fubiniego i nierówności Cauchy'ego--Schwarza
mamy dla prawie każdego $t$:
$|(T_kf)(t)|^2\le
(\int_{[0,1]}|k(t,s)|^2\,\mathrm d\lambda_1(s))\|f\|_2^2$.
Całkując po $t$ względem $\lambda_1$ otrzymujemy
$\|T_kf\|_2\le\|k\|_2\|f\|_2$.
Aby uzasadnić potrzebną aproksymację jądra, zauważmy, że
wskaźniki prostokątów $A\times B$ mają postać
$\chi_A(t)\chi_B(s)$, a ich skończone kombinacje
liniowe są gęste w $L^2([0,1]^2,\lambda_2)$.
Wynika to z gęstości funkcji prostych oraz aproksymacji
zbiorów mierzalnych przez skończone sumy prostokątów
w pseudometryce $A,B\mapsto\lambda_2(A\triangle B)$
(identyfikującej zbiory różniące się o zbiór miary zero).
Wybieramy zatem jądra $k_n$ będące skończonymi sumami
$g_j(t)h_j(s)$ i zbiegające do $k$ w $L^2$.
Każdy $T_{k_n}$ ma skończony rząd, bo jego obraz zawiera się
w skończenie wymiarowej otoczce liniowej odpowiednich
funkcji $g$. Z oszacowania normy wynika
$\|T_k-T_{k_n}\|\le\|k-k_n\|_2\to0$.
Zatem $T_k$ jest zwarty jako granica normowa operatorów
skończonego rzędu.
\end{proof}

\begin{uwaga}[Nie wszystkie operatory zwarte są aproksymowalne]
\label{rem:nonapproximable29}
W rozdziale o AP wykazano, że jeżeli przestrzeń wartości
$Y$ ma AP, to $\mathcal K(X,Y)=\mathcal A(X,Y)$.
Bez założenia AP w przestrzeni wartości takie twierdzenie
nie zachodzi automatycznie. Dlatego aproksymacja przez
operator skończonego rzędu nie jest równoważną
\emph{definicją} zwartości w dowolnych przestrzeniach Banacha.
\end{uwaga}

\subsection{Ideały operatorowe i przygotowanie teorii Fredholma}

\begin{definicja}[Ideał operatorowy]
\label{def:operator-ideal29}
Przez \emph{ideał operatorowy} w klasie przestrzeni Banacha
rozumiemy przyporządkowanie każdej parze $(X,Y)$
podprzestrzeni liniowej $\mathcal I(X,Y)\subseteq\mathcal L(X,Y)$,
zawierającej wszystkie operatory skończonego rzędu,
takie, że dla wszystkich przestrzeni Banacha $W,Z$
i operatorów $R\in\mathcal L(W,X)$, $T\in\mathcal I(X,Y)$,
$S\in\mathcal L(Y,Z)$ zachodzi
$STR\in\mathcal I(W,Z)$.
\end{definicja}

\begin{twierdzenie}[Trzy ideały i domkniętość]
\label{thm:ideals29}
Rodziny $\mathcal F$, $\mathcal A$ i $\mathcal K$ są
ideałami operatorowymi. Przy tym $\mathcal A(X,Y)$
i $\mathcal K(X,Y)$ są domknięte w normie operatorowej,
a $\mathcal F(X,Y)$ nie jest domknięta, jeżeli
$X=Y=\ell^2$. Dla wszystkich $X,Y$ zachodzi
$\mathcal F(X,Y)\subseteq\mathcal A(X,Y)\subseteq
\mathcal K(X,Y)$.
\end{twierdzenie}
\begin{proof}
Składanie operatora skończonego rzędu z operatorami
ograniczonymi z lewej i prawej strony zachowuje
skończoność wymiaru obrazu; dlatego $\mathcal F$ jest ideałem.
Jeżeli $T_n\to T$ w normie i każdy $T_n$ ma skończony
rząd, to $ST_nR\to STR$ w normie na mocy
$\|ST_nR-STR\|\le\|S\|\|T_n-T\|\|R\|$.
To daje własność ideału $\mathcal A$ i jego domkniętość
z definicji. Twierdzenie~\ref{thm:ideal-compact29}
daje własność ideału $\mathcal K$; jego domkniętość
udowodniono wcześniej. Operator diagonalny
$D_{(1/n)}:\ell^2\to\ell^2$ ma nieskończony rząd,
a jego obcięcia skończonego rzędu zbiegają do niego
w normie. To pokazuje, że $\mathcal F(\ell^2,\ell^2)$
nie jest domknięta. Zawierania między trzema ideałami
wynikają z definicji i zwartości operatorów
skończonego rzędu.
\end{proof}

\begin{definicja}[Norma istotna operatora]
\label{def:essential-norm29}
Dla $T\in\mathcal L(X,Y)$ definiujemy
$\|T\|_{\mathrm{ess}}=
\inf\{\|T-K\|:K\in\mathcal K(X,Y)\}$.
Jest to odległość operatora $T$ od domkniętej
podprzestrzeni operatorów zwartych; może wynosić zero
dla operatora niezerowego. Nie jest więc normą
na całej przestrzeni $\mathcal L(X,Y)$, ale wyznacza normę
na przestrzeni ilorazowej
$\mathcal L(X,Y)/\mathcal K(X,Y)$. Gdy $X=Y$,
przestrzeń ta jest dodatkowo algebrą ilorazową.
\end{definicja}

\begin{twierdzenie}[Norma istotna identyczności]
\label{thm:essential-identity29}
Jeżeli $X$ jest nieskończenie wymiarową przestrzenią
Banacha, to $\|I_X\|_{\mathrm{ess}}=1$.
\end{twierdzenie}
\begin{proof}
Wybór $K=0$ pokazuje, że odległość jest nie większa
niż $1$. Załóżmy, że istnieje $K\in\mathcal K(X,X)$
z $\|I_X-K\|<1$. Wtedy szereg geometryczny operatorów
$\sum_{n=0}^{\infty}(I_X-K)^n$ jest zbieżny w normie
i daje odwrotność $K$. W szczególności
$I_X=K^{-1}K$ byłby zwarty na mocy
Twierdzenia~\ref{thm:ideal-compact29}, wbrew
Twierdzeniu~\ref{thm:identity-compact29}.
\end{proof}

\begin{lemat}[Pierwszy krok teorii Fredholma]
\label{lem:compact-eigenvalue29}
Niech $T\in\mathcal K(X,X)$ i $\lambda\in\K\setminus\{0\}$.
Jeśli operator $\lambda I_X-T$ jest różnowartościowy,
to istnieje $c>0$ takie, że
$\|(\lambda I_X-T)x\|\ge c\|x\|$ dla wszystkich $x\in X$.
W szczególności obraz $\operatorname{ran}(\lambda I_X-T)$
jest domknięty.
\end{lemat}
\begin{proof}
Gdyby nie było takiego $c$, istniałby ciąg $(x_n)$
o normach równych $1$, spełniający
$(\lambda I_X-T)x_n\to0$. Ze zwartości $T$ możemy
wybrać podciąg, dla którego $Tx_{n_k}\to y$ w normie.
Wówczas $\lambda x_{n_k}=Tx_{n_k}+
(\lambda I_X-T)x_{n_k}\to y$, więc $x_{n_k}\to x=y/\lambda$
i $\|x\|=1$. Przez ciągłość operatora otrzymujemy
$(\lambda I_X-T)x=0$, co przeczy jego różnowartościowości.
Jeżeli obrazy $(\lambda I_X-T)x_n$ tworzą ciąg Cauchy'ego,
to uzyskane oszacowanie implikuje, że $(x_n)$ jest
ciągiem Cauchy'ego. Zupełność $X$ i ciągłość operatora
pokazują, że obraz jest domknięty.
\end{proof}

\begin{uwaga}[Granica zakresu rozdziału]
\label{rem:fredholm-next29}
Ostatni lemat nie stwierdza jeszcze surjektywności
$\lambda I_X-T$ ani nie opisuje krotności wartości
własnych. W następnym rozdziale dowiedziemy twierdzenia
Riesza--Schaudera, alternatywy Fredholma i stabilności
indeksu przy zaburzeniach zwartych, nie powtarzając
kryteriów zwartości i twierdzenia o transpozycji.
\end{uwaga}


% Wersja źródłowa: audyt II, jednostka 30
\setcounter{section}{29}
\noindent\textbf{Wyniki wykorzystywane.} Przyjmujemy znane twierdzenie Hahna--Banacha,
lemat Riesza, twierdzenie o odwzorowaniu otwartym, własności operatorów
zwartych i twierdzenie Schaudera o zwartości transpozycji.
\par\medskip

\section{Teoria Fredholma i twierdzenie Riesza--Schaudera}
\label{sec:fredholm30}

W tym rozdziale $X,Y$ są przestrzeniami Banacha nad
$\K\in\{\R,\C\}$, a $\mathcal L(X,Y)$ i $\mathcal K(X,Y)$
oznaczają odpowiednio operatory liniowe ciągłe i zwarte.
Kryteria zwartości i twierdzenie Schaudera o zwartości transpozycji
zostały już udowodnione. Najpierw przedstawimy alternatywę
Fredholma w przestrzeni Hilberta, w wersji korzystającej z
przybliżeń przez rzuty ortogonalne, a następnie argument
działający w dowolnych przestrzeniach Banacha.

\subsection{Operatory Fredholma i ich indeks}

\begin{definicja}[Operator Fredholma]
\label{def:fredholm30}
Niech $X,Y$ będą przestrzeniami Banacha i $T\in\mathcal L(X,Y)$.
Operator $T$ nazywamy \emph{operatorem Fredholma}, jeżeli
$\operatorname{ran}T$ jest domknięty, $\dim\ker T<\infty$ i
$\operatorname{codim}\operatorname{ran}T:=\dim(Y/\operatorname{ran}T)<\infty$.
Jego \emph{indeks Fredholma} wynosi
\[\operatorname{ind}T=\dim\ker T-\operatorname{codim}\operatorname{ran}T.\]
Symbol $X'$ oznacza przestrzeń dualną topologiczną $X$.
\end{definicja}

\begin{lemat}[Uzupełniające podprzestrzenie]
\label{lem:fredholm-complements30}
Każda skończenie wymiarowa podprzestrzeń $N$ przestrzeni Banacha
$X$ jest domknięta i ma domknięte dopełnienie topologiczne
$X=N\oplus X_1$, z ciągłymi rzutami. Jeżeli $R\subseteq Y$
jest domkniętą podprzestrzenią Banacha o skończonym kowymiarze,
to $Y=R\oplus C$, gdzie $C$ ma skończony wymiar i oba rzuty
są ciągłe.
\end{lemat}
\begin{proof}
Dla bazy $u_1,\ldots,u_n$ przestrzeni $N$ funkcjonały
współrzędnych są ciągłe na $N$ i przedłużają się na $X$
przez twierdzenie Hahna--Banacha do $f_1,\ldots,f_n\in X'$.
Wówczas $P(x)=\sum_{j=1}^n f_j(x)u_j$ jest ciągłym rzutem
na $N$, a $X_1=\ker P$. W drugim przypadku wybieramy bazę
$[v_1],\ldots,[v_m]$ skończenie wymiarowego ilorazu $Y/R$
i definiujemy $C=\operatorname{span}\{v_1,\ldots,v_m\}$.
Odwzorowanie przypisujące $y\in Y$ jego składową w $C$
przez ciągłe funkcjonały współrzędnych $Y/R$ jest ciągłym
rzutem o jądrze $R$. Obie dekompozycje są topologiczne.
\end{proof}

\begin{twierdzenie}[Obraz domknięty i warunek zgodności]
\label{thm:fredholm-compatibility30}
Jeżeli $T\in\mathcal L(X,Y)$ ma domknięty obraz, to
\[\operatorname{ran}T=\{y\in Y:\ f(y)=0\ \text{dla każdego }f\in\ker T'\}.\]
Jeżeli ponadto $T$ jest operatorem Fredholma, to
$\dim\ker T'=\operatorname{codim}\operatorname{ran}T$.
Zatem $Tx=y$ ma rozwiązanie wtedy i tylko wtedy, gdy
$f(y)=0$ dla wszystkich $f\in\ker T'$, a zbiór
rozwiązań ma postać $x_0+\ker T$.
\end{twierdzenie}
\begin{proof}
Z definicji $T'$ mamy $\ker T'=(\operatorname{ran}T)^\perp$,
gdzie symbol $R^\perp=\{f\in Y':f|_R=0\}$ oznacza
anihilator. Gdy $R=\operatorname{ran}T$ jest domknięty,
twierdzenie Hahna--Banacha oddziela każdy $y\notin R$
od $R$ ciągłym funkcjonałem zerującym się na $R$.
Stąd $R=(R^\perp)_\perp$. Ponadto $R^\perp$ jest
izomorficzna z $(Y/R)'$, które przy skończonym wymiarze
$Y/R$ ma ten sam wymiar. Reszta wynika z liniowości.
\end{proof}

\subsection{Alternatywa Fredholma w przestrzeni Hilberta}

\begin{twierdzenie}[Alternatywa Fredholma dla operatora zwartego]
\label{thm:fredholm-hilbert30}
Niech $H$ będzie zespoloną przestrzenią Hilberta, a
$K\in\mathcal K(H,H)$. Wówczas $T=I_H-K$ ma domknięty
obraz oraz skończenie wymiarowe jądra $\ker T$ i $\ker T^*$,
przy czym $\dim\ker T=\dim\ker T^*$.
Dla $f\in H$ równanie $Tu=f$ ma rozwiązanie dokładnie wtedy,
gdy $f\perp\ker T^*$; wszystkie rozwiązania tworzą
klasę $u_0+\ker T$.
\end{twierdzenie}
\begin{proof}
\emph{Krok 1: skończony wymiar jądra.}
Gdyby $N=\ker T$ był nieskończenie wymiarowy, zawierałby
nieskończony układ ortonormalny $(e_n)$. Równość $Ke_n=e_n$
przeczyłaby zwartości $K$.

\emph{Krok 2: domkniętość obrazu.}
Niech $M=N^\perp$. Istnieje $c>0$ takie, że
$\|Tx\|\ge c\|x\|$ dla $x\in M$: w przeciwnym przypadku
istniałby ciąg $x_n\in M$ o normach $1$, dla którego
$Tx_n\to0$. Ze zwartości $K$ wybralibyśmy podciąg
$(x_{n_j})$ z $Kx_{n_j}\to z$, a wtedy
$x_{n_j}=Tx_{n_j}+Kx_{n_j}\to z\in M\cap N$, przy $\|z\|=1$.
To sprzeczność. Jeżeli $Tx_n\to y$, zastępujemy $x_n$
jego rzutem ortogonalnym na $M$; oszacowanie dolne
czyni ten ciąg ciągiem Cauchy'ego. Jego granica $x\in M$
spełnia $Tx=y$.

\emph{Krok 3: indeks zerowy.}
Zwartość $K(B_H)$ daje skończenie wymiarową podprzestrzeń
$W\subseteq H$, dla której
$\|K-P_WK\|<1$: wystarczy wziąć podprzestrzeń
rozpiętą na punktach odpowiednio dokładnej skończonej siatki.
Piszemy $F=P_WK$, $E=K-F$, $\|E\|<1$ i
$G=(I-E)^{-1}F$. Szereg Neumanna dowodzi, że $I-E$
jest odwracalny, a
\[I-K=(I-E)(I-G).\]
Operator $G$ ma skończony rząd. Niech $W'=\operatorname{ran}G$
i rozłóżmy $H=W'\oplus(W')^\perp$.
W tych współrzędnych
\[(I-G)(w,z)=((I_{W'}-G|_{W'})w-Gz,\ z).\]
Składowa druga identyczności pozwala eliminować $-Gz$
i redukuje obliczenie wymiaru jądra i kowymiaru obrazu
do skończenie wymiarowego operatora $I_{W'}-G|_{W'}$.
Na mocy twierdzenia o wymiarze jądra i obrazu obie liczby są
równe. Mnożenie z lewej przez $I-E$ nie zmienia
indeksu, więc $\operatorname{ind}T=0$.

\emph{Krok 4: rozwiązywalność.}
Ponieważ $\operatorname{ran}T$ jest domknięty,
$(\operatorname{ran}T)^\perp=\ker T^*$ daje
$\operatorname{ran}T=(\ker T^*)^\perp$. Z kroku 3 wynika
$\dim\ker T=\dim\ker T^*$, zaś opis rozwiązań
wynika z liniowości.
\end{proof}

\begin{uwaga}[Zakres powyższego dowodu]
\label{rem:fredholm-hilbert-scope30}
W dowodzie wykorzystano rzut ortogonalny na skończenie
wymiarową podprzestrzeń. Dowolna przestrzeń Banacha
nie musi mieć własności aproksymacji, dlatego nie wolno
przenosić tego argumentu bez zmian. Alternatywę Fredholma
w przestrzeniach Banacha udowodnimy bez takiej aproksymacji.
\end{uwaga}

\subsection{Twierdzenie Atkinsona i indeks w przestrzeniach Banacha}

\begin{twierdzenie}[Twierdzenie Atkinsona]
\label{thm:atkinson30}
Dla $T\in\mathcal L(X,Y)$ następujące warunki są równoważne:
\begin{enumerate}[label=\textup{(\roman*)}]
\item $T$ jest operatorem Fredholma;
\item istnieje $S\in\mathcal L(Y,X)$, dla którego
$ST-I_X\in\mathcal K(X,X)$ i $TS-I_Y\in\mathcal K(Y,Y)$.
\end{enumerate}
Gdy (i) zachodzi, obie różnice można wybrać skończonego rzędu.
\end{twierdzenie}
\begin{proof}
Załóżmy (i) i wybierzmy rozkłady $X=N\oplus X_1$,
$Y=R\oplus C$, gdzie $N=\ker T$, $R=\operatorname{ran}T$.
Ograniczenie $A=T|_{X_1}:X_1\to R$ jest bijekcją między
przestrzeniami Banacha, więc $A^{-1}$ jest ograniczony.
Dla $y=r+c$ definiujemy $S(y)=A^{-1}r$.
Wtedy $ST=I_X-P_N$ i $TS=I_Y-P_C$, gdzie oba rzuty
mają skończony rząd.

Odwrotnie niech $ST=I_X-K_1$ i $TS=I_Y-K_2$,
gdzie $K_1,K_2$ są zwarte. Na $N=\ker T$ operator
$K_1$ działa jak identyczność, więc $N$ jest skończenie
wymiarowe (inaczej lemat Riesza przeczyłby zwartości).
Wybierzmy dopełnienie $X=N\oplus X_1$.
Jeśli $T|_{X_1}$ nie byłby ograniczony z dołu, istniałby
ciąg jednostkowy $x_n\in X_1$ z $Tx_n\to0$.
Równość $x_n=STx_n+K_1x_n$ oraz zwartość $K_1$
dawałyby normowo zbieżny podciąg do niezerowego wektora
w $N\cap X_1$, co jest niemożliwe. Zatem obraz $T$
jest domknięty.
Z równości $S'T'=I_{Y'}-K_2'$ (otrzymanej przez
transpozycję $TS=I_Y-K_2$) wynika, że
$\ker T'\subseteq\ker(I_{Y'}-K_2')$.
Operator $K_2'$ jest zwarty na mocy twierdzenia Schaudera,
więc drugie jądro ma skończony wymiar.
Twierdzenie~\ref{thm:fredholm-compatibility30} daje
$\operatorname{codim}\operatorname{ran}T=\dim\ker T'<\infty$.
\end{proof}

\begin{twierdzenie}[Lokalna stałość indeksu]
\label{thm:index-local30}
Niech $T\in\mathcal L(X,Y)$ będzie operatorem Fredholma.
Istnieje $\delta>0$ takie, że każdy $U\in\mathcal L(X,Y)$
z $\|U-T\|<\delta$ jest Fredholma i
$\operatorname{ind}U=\operatorname{ind}T$.
\end{twierdzenie}
\begin{proof}
Wybierzmy rozkłady $X=N\oplus X_1$, $Y=R\oplus C$
z dowodu twierdzenia Atkinsona. W ich współrzędnych
$T(n,x)=(Ax,0)$, gdzie $A:X_1\to R$ jest izomorfizmem.
Dla $U$ bliskiego $T$ mamy
\[U(n,x)=(Bn+Dx,\ En+Fx),\]
gdzie $B:N\to R$, $D:X_1\to R$, $E:N\to C$,
$F:X_1\to C$. Operator $D$ jest odwracalny, gdy
$\|A^{-1}(D-A)\|<1$, co zachodzi dla $U$ dostatecznie
bliskiego $T$.
Dokonujemy dwóch \emph{odwracalnych} zmian współrzędnych.
Najpierw w dziedzinie podstawiamy
$(n,x)=(n,x'-D^{-1}Bn)$. Otrzymujemy
\[U(n,x)=(Dx',\ (E-FD^{-1}B)n+Fx').\]
Następnie w przeciwdziedzinie stosujemy
$(r,c)\mapsto(r,c-FD^{-1}r)$. Wynik ma postać
\[(n,x')\mapsto(Dx',\ (E-FD^{-1}B)n).\]
Obraz jest domknięty, ponieważ $D$ jest izomorfizmem,
a $E-FD^{-1}B:N\to C$ działa między przestrzeniami
skończenie wymiarowymi. Indeks drugiego bloku to
$\dim N-\dim C$ niezależnie od jego rzędu, gdyż
\[\dim\ker M-\dim(C/\operatorname{ran}M)=\dim N-\dim C\]
dla każdego $M:N\to C$ w skończonym wymiarze.
Stąd $\operatorname{ind}U=\dim N-\dim C=\operatorname{ind}T$.
\end{proof}

\begin{twierdzenie}[Zaburzenia zwarte zachowują indeks]
\label{thm:index-compact-perturb30}
Niech $T\in\mathcal L(X,Y)$ będzie operatorem Fredholma,
a $K\in\mathcal K(X,Y)$. Wówczas $T+K$ jest Fredholma i
$\operatorname{ind}(T+K)=\operatorname{ind}T$.
\end{twierdzenie}
\begin{proof}
Z twierdzenia Atkinsona istnieje $S:Y\to X$ takie, że
$ST-I_X$ i $TS-I_Y$ są zwarte. Dla każdego $t\in[0,1]$
obie różnice $S(T+tK)-I_X$ i $(T+tK)S-I_Y$ są zwarte.
Zatem $T+tK$ jest Fredholma dla wszystkich $t$.
Jego indeks zależy lokalnie stale od $t$ na mocy
Twierdzenia~\ref{thm:index-local30}; ponieważ $[0,1]$
jest spójny, jest stały.
\end{proof}

\begin{wniosek}[Alternatywa Fredholma w przestrzeni Banacha]
\label{cor:fredholm-banach30}
Niech $X$ będzie przestrzenią Banacha i
$K\in\mathcal K(X,X)$. Wówczas $I_X-K$ jest operatorem
Fredholma indeksu zero. Równanie $(I_X-K)x=y$ jest
rozwiązywalne dokładnie wtedy, gdy $f(y)=0$ dla każdego
$f\in\ker(I_{X'}-K')$. Jeśli $\ker(I_X-K)=\{0\}$,
to $I_X-K$ jest izomorfizmem liniowo-topologicznym.
\end{wniosek}
\begin{proof}
Operator $I_X$ jest Fredholma indeksu zero; zaburzenie
$-K$ zachowuje indeks. Warunek zgodności wynika z
Twierdzenia~\ref{thm:fredholm-compatibility30}. Gdy jądro
jest zerowe, indeks zero i domkniętość obrazu dają
surjektywność. Ciągłość odwrotności wynika z twierdzenia
o odwzorowaniu otwartym.
\end{proof}

\begin{przyklad}[Indeks może być niezerowy]
\label{ex:shift-index30}
W $X=\ell^2(\N)$ przesunięcie w prawo
$R(x_1,x_2,\ldots)=(0,x_1,x_2,\ldots)$ ma zerowe jądro
i obraz kowymiaru jeden, a więc indeks $-1$.
Przesunięcie w lewo $L(x_1,x_2,\ldots)=(x_2,x_3,\ldots)$
ma indeks $1$. Nie każdy operator Fredholma ma więc
indeks zero.
\end{przyklad}

\subsection{Twierdzenie Riesza--Schaudera i widmo operatora zwartego}

\begin{definicja}[Widmo i rezolwenta]
\label{def:spectrum30}
Niech $X$ będzie zespoloną przestrzenią Banacha i
$K\in\mathcal L(X,X)$. Definiujemy
$\rho(K)=\{\lambda\in\C:\lambda I_X-K\text{ jest odwracalny}\}$
oraz $\sigma(K)=\C\setminus\rho(K)$.
Liczbę $\lambda$ nazywamy wartością własną, gdy
$\ker(\lambda I_X-K)\ne\{0\}$.
\end{definicja}

\begin{twierdzenie}[Riesz--Schauder]
\label{thm:riesz-schauder30}
Niech $X$ będzie nieskończenie wymiarową zespoloną
przestrzenią Banacha, a $K\in\mathcal K(X,X)$.
Wtedy:
\begin{enumerate}[label=\textup{(\roman*)}]
\item $0\in\sigma(K)$, a $\sigma(K)$ jest niepustym
zwartym podzbiorem $\C$;
\item każdy $\lambda\in\sigma(K)\setminus\{0\}$ jest
wartością własną o skończenie wymiarowym jądrze;
\item dla każdego $\varepsilon>0$ zbiór
$\{\lambda\in\sigma(K):|\lambda|\ge\varepsilon\}$
jest skończony;
\item dla $\lambda\ne0$ uogólniona przestrzeń własna
$\bigcup_{m\ge1}\ker(\lambda I_X-K)^m$ jest
skończenie wymiarowa i równa
$\ker(\lambda I_X-K)^{m_0}$ dla pewnego $m_0\ge1$.
\end{enumerate}
\end{twierdzenie}
\begin{proof}
Dla $|\lambda|>\|K\|$ szereg Neumanna dowodzi
odwracalności $\lambda I_X-K$. Zbiór rezolwenty jest otwarty:
jeśli $\lambda_0 I_X-K$ ma odwrotność, to dla małego
$|\lambda-\lambda_0|$ stosujemy szereg Neumanna do
$I+(\lambda-\lambda_0)(\lambda_0 I_X-K)^{-1}$.
Widmo jest więc domknięte i ograniczone. Ponieważ
$K$ nie może być odwracalny (inaczej
$I_X=K^{-1}K$ byłby zwarty), $0\in\sigma(K)$.

Dla $\lambda\ne0$ operator $\lambda I_X-K$ jest Fredholma
indeksu zero jako niezerowa wielokrotność $I_X-\lambda^{-1}K$.
Jeżeli nie jest odwracalny, nie może być różnowartościowy:
zerowe jądro i indeks zero wymuszałyby surjektywność.
Jądro jest skończenie wymiarowe. To dowodzi (ii).

Gdyby (iii) zawiodło, wybralibyśmy parami różne
niezerowe wartości własne $\lambda_n$ o
$|\lambda_n|\ge\varepsilon$ i odpowiadające im wektory
własne $v_n\ne0$. Wektory te są liniowo niezależne.
Niech $M_n=\operatorname{span}\{v_1,\ldots,v_n\}$.
Podprzestrzenie $M_n$ są niezmiennicze względem $K$.
Lemat Riesza dostarcza $x_n\in M_n$ o normie $1$ i
$\operatorname{dist}(x_n,M_{n-1})>1/2$ dla $n\ge2$.
Na ilorazie $M_n/M_{n-1}$ operator $K$ działa przez
mnożenie przez $\lambda_n$, więc
$Kx_n-\lambda_nx_n\in M_{n-1}$. Jeśli $m<n$,
to $Kx_m\in M_{n-1}$, a zatem
\[\|Kx_n-Kx_m\|\ge
|\lambda_n|\operatorname{dist}(x_n,M_{n-1})>\varepsilon/2.\]
To przeczy zwartości $K$.

Dla (iv) połóżmy $N_m=\ker(\lambda I_X-K)^m$.
Wielomian $\lambda^{-m}(\lambda I_X-K)^m$ ma postać
$I_X-Q_m$ z $Q_m$ zwartym (każdy niezerowy składnik
w $Q_m$ zawiera dodatnią potęgę $K$). Na mocy
alternatywy Fredholma każde $N_m$ ma skończony wymiar.
Gdyby rosnący ciąg $N_m$ nie stabilizował się, można
byłoby wybrać rosnące indeksy $n_j$ z
$N_{n_j-1}\subsetneq N_{n_j}$ oraz wektory
$x_j\in N_{n_j}$ o normie $1$ takie, że
$\operatorname{dist}(x_j,N_{n_j-1})>1/2$.
Dla $i<j$ zachodzą $Kx_i\in N_{n_j-1}$ i
$Kx_j-\lambda x_j\in N_{n_j-1}$; stąd
$\|Kx_j-Kx_i\|>|\lambda|/2$, co przeczy zwartości.
Zatem ciąg jąder stabilizuje się.
\end{proof}

\begin{przyklad}[Widmo operatora diagonalnego]
\label{ex:diagonal-spectrum30}
W $\ell^2(\N)$ operator
$K(x_1,x_2,\ldots)=(x_1,x_2/2,x_3/3,\ldots)$
jest zwarty jako granica operatorów skończonego rzędu.
Ma widmo $\sigma(K)=\{0\}\cup\{1/n:n\ge1\}$;
każda niezerowa wartość własna jest prosta.
\end{przyklad}

\begin{przyklad}[Operator Volterry: widmo bez niezerowych wartości własnych]
\label{ex:volterra-spectrum30}
W $X=C([0,1],\C)$ definiujemy
$(Vf)(t)=\int_{[0,t]}f(s)\,\mathrm d\lambda_1(s)$.
Operator $V$ jest zwarty. Iterując wzór na całkę,
dostajemy $\|V^n\|\le 1/n!$, ponieważ jądro $V^n$
jest równe $(t-s)^{n-1}/(n-1)!$ dla $0\le s\le t$.
Dla $\lambda\ne0$ szereg
$\lambda^{-1}\sum_{n\ge0}(V/\lambda)^n$
jest zbieżny w normie operatorowej i określa
$(\lambda I_X-V)^{-1}$. Zatem $\sigma(V)=\{0\}$,
mimo że $V\ne0$.
\end{przyklad}

\subsection{Zastosowanie w teorii równań różniczkowych cząstkowych}

\begin{uwaga}[Od operatora eliptycznego do operatora zwartego]
\label{rem:elliptic-fredholm30}
W skrypcie \emph{Równania różniczkowe cząstkowe, tom I}
rozpatruje się faktoryzację $L=L_\mu(I_V-K)$,
w której $L_\mu:V\to V^*$ jest izomorfizmem,
a $K:V\to V$ operatorem zwartym otrzymanym dzięki
zwartemu zanurzeniu Sobolewa. Wówczas $\operatorname{ind}L=0$,
a warunek rozwiązywalności wynika z anihilatora obrazu.
Wyprowadzenie odpowiednich nierówności koercywności
i zagadnienia brzegowe pozostają w skrypcie PDE.
\end{uwaga}


% Wersja źródłowa: audyt II, jednostka 31
\setcounter{section}{30}
\noindent\textbf{Wyniki wykorzystywane.} Zakładamy znane twierdzenia Hahna--Banacha, Banacha--Alaoglu oraz podstawowe twierdzenia skalarnej analizy zespolonej (Cauchy, Liouville). Twierdzenie Stone'a--Weierstrassa jest udowodnione w tym rozdziale.\par\medskip
\section{Algebry Banacha, widmo i holomorficzny rachunek funkcyjny}
\label{sec:banach-algebras31}

W poprzednim rozdziale wprowadziliśmy widmo ograniczonego operatora
na zespolonej przestrzeni Banacha, przede wszystkim w związku z operatorami
zwartymi. Teraz rozpatrzymy widmo elementu \emph{dowolnej} zespolonej algebry
Banacha. Otrzymamy w ten sposób wspólne podstawy teorii operatorów, algebr
funkcyjnych oraz algebr $C^*$. Nie powtarzamy twierdzenia Riesza--Schaudera:
jego twierdzenia będą przykładami ogólnego rachunku spektralnego.

\subsection{Algebra Banacha i elementy odwracalne}

\begin{definicja}[Zespolona algebra Banacha]
\label{def:banach-algebra31}
\emph{Zespoloną algebrą Banacha} nazywamy zespoloną przestrzeń Banacha
$A$ z łącznym, dwuliniowym mnożeniem $A\times A\to A$ spełniającym
$\|ab\|\le\|a\|\|b\|$ dla wszystkich $a,b\in A$.
W tym rozdziale, o ile nie zaznaczono inaczej, zakładamy, że istnieje
element jednostkowy $1_A$ i $\|1_A\|=1$. Algebra nie musi być przemienna.
Element $a\in A$ nazywamy odwracalnym, gdy istnieje $b\in A$ takie,
że $ab=ba=1_A$; oznaczamy wtedy $b=a^{-1}$.
\end{definicja}

\begin{przyklad}[Trzy podstawowe algebry]
\label{ex:banach-algebras31}
Jeśli $X$ jest niezerową zespoloną przestrzenią Banacha, to $\mathcal L(X,X)$
z normą operatorową i składaniem jest algebrą Banacha.
Jeśli $K$ jest niepustą zwartą przestrzenią Hausdorffa, to $C(K,\C)$
z normą supremum i mnożeniem punktowym jest przemienną algebrą Banacha.
Przestrzeń $\ell^1(\N_0)$ z iloczynem splotowym
$(a*b)_n=\sum_{j=0}^n a_jb_{n-j}$ jest przemienną algebrą Banacha,
której jednostką jest $(1,0,0,\ldots)$.
W ostatnim przypadku oszacowanie $\|a*b\|_1\le\|a\|_1\|b\|_1$
wynika ze zbieżności bezwzględnej sum podwójnych.
\end{przyklad}

\begin{lemat}[Szereg Neumanna i otwartość grupy odwracalnych]
\label{lem:neumann31}
Niech $A$ będzie zespoloną algebrą Banacha z jedynką.
Jeżeli $\|u\|<1$, to $1_A-u$ jest odwracalny i
$(1_A-u)^{-1}=\sum_{n=0}^\infty u^n$ (zbieżność w normie).
Zbiór $A^\times$ wszystkich elementów odwracalnych jest otwarty,
a odwzorowanie $a\mapsto a^{-1}$ jest na nim ciągłe.
\end{lemat}
\begin{proof}
Sumy częściowe spełniają
$(1_A-u)\sum_{n=0}^N u^n=1_A-u^{N+1}$,
a analogiczna tożsamość zachodzi po przemnożeniu z prawej.
Ponieważ $\|u^{N+1}\|\le\|u\|^{N+1}\to0$, uzyskujemy odwrotność.
Jeżeli $a$ jest odwracalny oraz $\|a^{-1}(b-a)\|<1$, to
$b=a(1_A+a^{-1}(b-a))$ jest odwracalny. Ponadto
$b^{-1}-a^{-1}=b^{-1}(a-b)a^{-1}$; lokalne oszacowanie
normy $b^{-1}$ przez szereg Neumanna daje ciągłość.
\end{proof}

\subsection{Widmo i rezolwenta}

\begin{definicja}[Widmo elementu algebry]
\label{def:spectrum31}
Dla $a\in A$ definiujemy \emph{zbiór rezolwenty}
$\rho_A(a)=\{z\in\C:z1_A-a\in A^\times\}$ oraz
\emph{widmo} $\sigma_A(a)=\C\setminus\rho_A(a)$.
Na $\rho_A(a)$ definiujemy \emph{rezolwentę}
$R_a(z)=(z1_A-a)^{-1}$.
Podpis $A$ pomijamy, gdy algebra jest ustalona. W przypadku
$A=\mathcal L(X,X)$ jest to definicja z poprzedniego rozdziału.
\end{definicja}

\begin{twierdzenie}[Zwartość i niepustość widma]
\label{thm:spectrum-nonempty31}
Niech $A$ będzie zespoloną algebrą Banacha z jedynką i $a\in A$.
Wtedy $\sigma_A(a)$ jest niepustym zwartym podzbiorem
$\{z\in\C:|z|\le\|a\|\}$. Rezolwenta $R_a$ jest holomorficzna
na $\rho_A(a)$ jako funkcja o wartościach w przestrzeni Banacha $A$.
Dla $z,w\in\rho_A(a)$ zachodzi tożsamość rezolwenty
\[
 R_a(z)-R_a(w)=(w-z)R_a(z)R_a(w).
\]
\end{twierdzenie}
\begin{proof}
Dla $|z|>\|a\|$ lemat~\ref{lem:neumann31} daje
$R_a(z)=\sum_{n\ge0}a^n z^{-n-1}$. Zatem widmo jest ograniczone,
a jego dopełnienie otwarte. Ten sam szereg, zastosowany do
$R_a(z_0)$, daje lokalne rozwinięcie rezolwenty w potęgi $z-z_0$,
a więc holomorficzność. Tożsamość rezolwenty wynika z przemnożenia
równości $(z1_A-a)^{-1}-(w1_A-a)^{-1}$ przez odpowiednie czynniki.

Gdyby widmo było puste, $R_a$ byłaby całkowitą funkcją holomorficzną
$\C\to A$. Dla $|z|>2\|a\|$ mielibyśmy
$\|R_a(z)\|\le 2/|z|$. Jest więc ograniczona na całej płaszczyźnie.
Dla dowolnego $\varphi\in A'$ funkcja skalarna
$\varphi\circ R_a$ jest całkowita i ograniczona, a z twierdzenia
Liouville'a stała. Granica przy $|z|\to\infty$ daje
$\varphi(R_a(z))=0$. Funkcjonały ciągłe rozdzielają punkty $A$,
więc $R_a(z)=0$, co przeczy $(z1_A-a)R_a(z)=1_A$.
\end{proof}

\begin{przyklad}[Widmo operatora mnożenia]
\label{ex:multiplication-spectrum31}
Dla $f\in C(K,\C)$ zachodzi $\sigma_{C(K)}(f)=f(K)$.
Istotnie, $z1-f$ ma odwrotność w $C(K)$ dokładnie wtedy,
gdy nigdzie nie znika; odwrotnością jest wówczas $1/(z-f)$.
W szczególności $\sigma(f)$ może być dowolnym niepustym zwartym
podzbiorem $\C$ (weźmy $K$ równy temu zbiorowi i $f(z)=z$).
\end{przyklad}

\begin{uwaga}[Widmo zależy od algebry]
Jeżeli $B\subseteq A$ jest domkniętą podalgebrą z tą samą jednostką,
to $\sigma_A(b)\subseteq\sigma_B(b)$ dla $b\in B$:
odwrotność w $B$ jest również odwrotnością w $A$.
Równość nie musi zachodzić. Niech
$\mathbb D=\{z\in\C:|z|<1\}$ i
$\mathbb T=\{z\in\C:|z|=1\}$.
Rozważmy algebrę dyskową
$\mathcal A(\overline{\mathbb D})$, tj. funkcje ciągłe
na $\overline{\mathbb D}$ i holomorficzne w $\mathbb D$.
Ograniczanie do brzegu daje izometryczną podalgebrę
$B\subseteq C(\mathbb T)$: izometrię zapewnia zasada maksimum,
a domkniętość wynika z jednostajnej zbieżności funkcji
holomorficznych na podzbiorach zwartych dysku.
Dla funkcji współrzędnej $u(z)=z$ na $\mathbb T$ zachodzi
$\sigma_{C(\mathbb T)}(u)=\mathbb T$.
Jednak $\sigma_B(u)=\overline{\mathbb D}$: dla
$|\lambda|>1$ funkcja $1/(u-\lambda)$ jest w $B$,
a dla $|\lambda|<1$ nie mogłaby być odwrotnością
$u-\lambda$ w $B$, gdyż jej holomorficzne przedłużenie
musiałoby spełniać $(z-\lambda)g(z)=1$ w dysku.
Dla $|\lambda|=1$ odwrotność nie istnieje już w
$C(\mathbb T)$.
\end{uwaga}

\begin{twierdzenie}[Wzór na promień spektralny]
\label{thm:spectral-radius31}
Niech $A$ będzie zespoloną algebrą Banacha z jedynką, $a\in A$.
Wówczas granica $\lim_{n\to\infty}\|a^n\|^{1/n}$ istnieje
i jest równa promieniowi spektralnemu
\[
 r(a):=\max\{|z|:z\in\sigma_A(a)\}
 =\lim_{n\to\infty}\|a^n\|^{1/n}
 =\inf_{n\ge1}\|a^n\|^{1/n}.
\]
\end{twierdzenie}
\begin{proof}
Nierówność $\|a^{n+m}\|\le\|a^n\|\|a^m\|$ daje istnienie granicy
$L$ i równość z infimum. Rzeczywiście, jeśli $a^m=0$ dla
pewnego $m$, teza o granicy jest natychmiastowa.
W przeciwnym razie, ustalając $m\ge1$ i zapisując $n=qm+r$,
$0\le r<m$, otrzymujemy
$\|a^n\|\le\|a^m\|^q\max_{0\le j<m}\|a^j\|$.
Po wyciągnięciu pierwiastka stopnia $n$ i przejściu do granicy
z góry dostajemy
$\limsup_n\|a^n\|^{1/n}\le\|a^m\|^{1/m}$.
Ponieważ $m$ jest dowolne, a każdy wyraz ciągu pierwiastków
jest nie mniejszy niż ich infimum, otrzymujemy żądaną granicę.
Jeżeli $|z|>L$, z definicji granicy szereg
$\sum_{n\ge0}a^n z^{-n-1}$ jest normowo zbieżny,
a pomnożenie przez $z1_A-a$ dowodzi odwracalności.
Zatem $r(a)\le L$.

Dla $R>\|a\|$ z szeregu Neumanna i skalarnego wzoru Cauchy'ego
(stosowanego po złożeniu z dowolnym $\varphi\in A'$) otrzymujemy
\[
 a^n=\frac{1}{2\pi i}\int_{|z|=R}z^nR_a(z)\,\mathrm dz.
\]
Całka jest granicą sum Riemanna w zupełnej przestrzeni Banacha $A$;
krzywa ma dodatnią orientację. Holomorficzność rezolwenty pozwala
przesunąć okrąg do dowolnego promienia $R>r(a)$, gdyż pierścień
między okręgami nie przecina widma. W konsekwencji
$\|a^n\|\le R^{n+1}\sup_{|z|=R}\|R_a(z)\|$.
Po wyciągnięciu pierwiastka stopnia $n$ uzyskujemy $L\le R$.
Przechodząc z $R$ do $r(a)$ dostajemy tezę.
\end{proof}

\begin{przyklad}[Element quasi-nilpotentny]
Operator Volterry $V$ na $C([0,1],\C)$ z
poprzedniego rozdziału ma $\sigma(V)=\{0\}$,
mimo że $V\ne0$. Zatem $r(V)=0$ i
$\lim_n\|V^n\|^{1/n}=0$. Nie oznacza to $V^m=0$
dla jakiegokolwiek $m\ge1$.
\end{przyklad}

\subsection{Charaktery, ideały maksymalne i transformacja Gelfanda}

\begin{definicja}[Charakter i przestrzeń charakterów]
\label{def:character31}
Niech $A$ będzie \emph{przemienną} zespoloną algebrą Banacha
z jedynką. Charakterem nazywamy niezerowy homomorfizm
zespolonych algebr $\chi:A\to\C$ (a priori bez założenia ciągłości).
Zbiór wszystkich charakterów oznaczamy $\Delta(A)$.
\emph{Transformacja Gelfanda} przyporządkowuje elementowi $a\in A$
funkcję $\widehat a:\Delta(A)\to\C$, $\widehat a(\chi)=\chi(a)$.
\end{definicja}

\begin{lemat}[Twierdzenie Gelfanda--Mazura]
\label{lem:gelfand-mazur31}
Każda zespolona algebra Banacha z jedynką, w której każdy niezerowy
element jest odwracalny, jest izomorficzna jako unitalna algebra z $\C$.
\end{lemat}
\begin{proof}
Dla $a\in A$ widmo $\sigma(a)$ jest niepuste na mocy
twierdzenia~\ref{thm:spectrum-nonempty31}. Wybierzmy
$z\in\sigma(a)$. Element $a-z1_A$ nie jest odwracalny,
więc przy założeniu o $A$ musi być równy zeru.
Zatem każdy $a$ jest skalarną wielokrotnością $1_A$.
\end{proof}

\begin{twierdzenie}[Charaktery i widmo elementu]
\label{thm:gelfand-spectrum31}
Niech $A$ będzie przemienną zespoloną algebrą Banacha z jedynką.
Wówczas każdy charakter $\chi\in\Delta(A)$ jest ciągły,
$\|\chi\|=1$, a $\Delta(A)$ jest niepustą zwartą przestrzenią
w topologii słabej-* $\sigma(A',A)$. Dla każdego $a\in A$ zachodzi
\[
 \sigma_A(a)=\{\chi(a):\chi\in\Delta(A)\},\qquad
 \|\widehat a\|_\infty=r(a)\le\|a\|.
\]
W szczególności $a\mapsto\widehat a$ jest ciągłym homomorfizmem
unitalnym $A\to C(\Delta(A),\C)$.
\end{twierdzenie}
\begin{proof}
Dla charakteru $\chi$ mamy $\chi(1_A)=1$.
Jeżeli $\chi(a)\notin\sigma(a)$, element $a-\chi(a)1_A$
byłby odwracalny, chociaż jego obraz przez $\chi$ jest zerem.
To niemożliwe, więc $|\chi(a)|\le r(a)\le\|a\|$.
Zatem $\chi$ jest ciągły i ma normę $1$.

Jeżeli $z\in\sigma(a)$, element $a-z1_A$ jest nieodwracalny.
Ideał przez niego generowany jest właściwy i, na mocy lematu Zorna,
zawiera się w pewnym ideale maksymalnym $M$. Ideał $M$ jest domknięty:
gdyby $\overline M=A$, zawierałby element $m$ z $\|m-1_A\|<1$,
a więc odwracalny, co przeczy właściwości ideału.
Iloraz $A/M$ jest zespoloną algebrą Banacha, której każdy niezerowy
element jest odwracalny. Z lematu~\ref{lem:gelfand-mazur31}
$A/M\simeq\C$, a odwzorowanie ilorazowe złożone z tym
izomorfizmem daje charakter $\chi$ spełniający $\chi(a)=z$.
To dowodzi formuły na widmo i niepustości $\Delta(A)$.

Topologia słaba-* czyni $\chi\mapsto\chi(a)$ ciągłym.
Zbiór charakterów jest w kuli jednostkowej $A'$ domknięty,
ponieważ równania $\chi(1_A)=1$ i
$\chi(ab)=\chi(a)\chi(b)$ są zachowane przy zbieżności słabej-*.
Z twierdzenia Banacha--Alaoglu kula jest zwarta.
Norma $\widehat a$ jest maksimum $|\chi(a)|$, czyli $r(a)$.
\end{proof}

\begin{uwaga}[Algebry niepółproste]
Transformacja Gelfanda nie musi być różnowartościowa.
Na przykład w dwuwymiarowej algebrze
$A=\C[\epsilon]/(\epsilon^2)$, wyposażonej w normę
$\|a+b\epsilon\|=|a|+|b|$, istnieje tylko jeden charakter,
$\chi(a+b\epsilon)=a$. Niezerowy element $\epsilon$ przechodzi na
funkcję zerową. Jądro transformacji Gelfanda jest radykałem
Jacobsona, a w przemiennej algebrze Banacha składa się dokładnie
z elementów o promieniu spektralnym zero.
\end{uwaga}

\subsection{Holomorficzny rachunek funkcyjny}

\begin{definicja}[Rachunek funkcyjny Dunforda]
\label{def:holomorphic-calculus31}
Niech $A$ będzie zespoloną algebrą Banacha z jedynką, $a\in A$,
a $U\subseteq\C$ otwartym otoczeniem $\sigma_A(a)$.
Dla $f\in\mathcal O(U)$ wybieramy skończoną, zorientowaną
rodzinę zamkniętych krzywych kawałkami $C^1$, oznaczoną $\Gamma$,
leżących w $U\setminus\sigma(a)$ i stanowiących brzeg
skończonej sumy obszarów o domknięciach zawartych w $U$.
Wymagamy, aby indeks $\operatorname{Ind}_\Gamma(w)$ był równy $1$
dla każdego $w\in\sigma(a)$.
Takie kontury istnieją dzięki zwartości widma.
Definiujemy
\[
 f(a):=\frac{1}{2\pi i}\int_\Gamma f(z)(z1_A-a)^{-1}\,\mathrm dz.
\]
Całka oznacza granicę sum Riemanna po parametryzacjach konturów
w normie przestrzeni Banacha $A$; zapis obejmuje sumę całek
po wszystkich składowych konturu. Kontury wybieramy tak, aby część
obszaru między nimi a widmem mieściła się w $U$.
\end{definicja}

\begin{twierdzenie}[Własności holomorficznego rachunku funkcyjnego]
\label{thm:holomorphic-calculus31}
Niech $A,a,U$ będą jak w definicji~\ref{def:holomorphic-calculus31}.
Odwzorowanie $f\mapsto f(a)$ z $\mathcal O(U)$ do $A$ jest unitalnym
homomorfizmem algebr. Dla wielomianu $p$ otrzymujemy zwykły element
$p(a)$, a dla $z\notin\sigma(a)$ rachunek funkcyjny funkcji
$w\mapsto(z-w)^{-1}$ daje $R_a(z)$. Wartość $f(a)$ zależy wyłącznie
od kiełka $f$ w otoczeniu widma i jest zgodna z ograniczaniem $U$.
Ponadto
\[
 \sigma_A(f(a))=f(\sigma_A(a)).
\]
\end{twierdzenie}
\begin{proof}
Niezależność od wyboru konturu wynika ze skalarnego twierdzenia
Cauchy'ego po złożeniu z każdym $\varphi\in A'$; twierdzenie
Hahna--Banacha przenosi równość funkcjonałów na równość w $A$.
Liniowość jest bezpośrednia. Dla stałej i funkcji $f(z)=z$
korzystamy z rozwinięcia rezolwenty na dużym okręgu i wzoru
Cauchy'ego na współczynniki, otrzymując $1_A$ i $a$.

Dla mnożenia bierzemy dwa zespoły konturów, z których zewnętrzny
otacza wewnętrzny, i używamy tożsamości
$R_a(z)R_a(w)=(R_a(w)-R_a(z))/(z-w)$.
Po podstawieniu do podwójnej całki dla $f(a)g(a)$
i całkowaniu najpierw po konturze zewnętrznym w wyrazie z $R_a(w)$
otrzymujemy przez wzór Cauchy'ego $f(w)g(w)R_a(w)$;
wyraz z $R_a(z)$ znika, ponieważ odpowiadająca mu funkcja
$w\mapsto g(w)/(z-w)$ jest holomorficzna wewnątrz konturu
wewnętrznego. Gdy $U$ jest niespójny, zwartość widma pozwala wybrać
skończenie wiele rozłącznych obszarów o domknięciach w $U$;
na każdym z nich rachunek jest identyczny.
Daje to $(fg)(a)=f(a)g(a)$.

Dla funkcji $g(w)=(z-w)^{-1}$, $z\notin\sigma(a)$,
rozważanej na dostatecznie małym otoczeniu widma, zachodzi $(z1_A-a)g(a)=1_A$ na mocy własności homomorfizmu;
stąd $g(a)=R_a(z)$.
Jeśli $\lambda\notin f(\sigma(a))$, funkcja
$1/(f-\lambda)$ jest holomorficzna na pewnym otoczeniu widma.
Z własności homomorfizmu wynika odwracalność
$f(a)-\lambda1_A$, czyli $\sigma(f(a))\subseteq f(\sigma(a))$.
Jeśli $\lambda=f(\mu)$ dla $\mu\in\sigma(a)$,
funkcja $h(w)=(f(w)-f(\mu))/(w-\mu)$ ma usuwalną
osobliwość w $\mu$ i jest holomorficzna na otoczeniu widma.
Gdyby $f(a)-f(\mu)1_A=(a-\mu1_A)h(a)$ było odwracalne,
to czynniki, które komutują, miałyby odwrotności
(w unitalnej algebrze odwracalność iloczynu komutujących czynników
pociąga odwracalność obu). Otrzymalibyśmy sprzeczność
z $\mu\in\sigma(a)$; stąd przeciwne zawieranie.
\end{proof}

\begin{przyklad}[Eksponenta i funkcje wymierne]
Jeśli $a\in A$, szereg
$e^a=\sum_{n\ge0}a^n/n!$ zbiega w normie i odpowiada
holomorficznemu rachunkowi funkcyjnemu funkcji $e^z$.
Twierdzenie~\ref{thm:holomorphic-calculus31} daje
$\sigma(e^a)=e^{\sigma(a)}$.
Jeśli $q$ jest wielomianem bez zer na $\sigma(a)$,
a $p$ wielomianem, to $(p/q)(a)=p(a)q(a)^{-1}$.
\end{przyklad}

\subsection{Algebry \texorpdfstring{$C^*$}{C*} i ich przemienny model}

\begin{definicja}[Algebra $C^*$]
\label{def:cstar31}
\emph{Algebrą $C^*$} nazywamy zespoloną algebrę Banacha $A$
z odwzorowaniem $a\mapsto a^*$ spełniającym
$(a+b)^*=a^*+b^*$, $(\lambda a)^*=\bar\lambda a^*$,
$(ab)^*=b^*a^*$, $(a^*)^*=a$ oraz
$\|a^*a\|=\|a\|^2$. W tym rozdziale algebry $C^*$
są unitalne. Element $a$ jest \emph{samosprzężony}, gdy $a=a^*$,
a \emph{normalny}, gdy $a^*a=aa^*$.
\end{definicja}

\begin{przyklad}[Dwie algebry $C^*$]
Dla zespolonej przestrzeni Hilberta $H$ algebra $\mathcal L(H,H)$
z operacją operatora sprzężonego jest algebrą $C^*$;
tożsamość $\|T^*T\|=\|T\|^2$ wynika z definicji
normy operatorowej i iloczynu skalarnego.
Dla przestrzeni zwartej Hausdorffa $K$ algebra $C(K,\C)$
z operacją $f^*(x)=\overline{f(x)}$ jest przemienną algebrą $C^*$.
\end{przyklad}

\begin{lemat}[Norma i widmo elementu normalnego]
\label{lem:cstar-normal31}
Niech $A$ będzie unitalną algebrą $C^*$, a $a\in A$ elementem
normalnym. Wtedy $r(a)=\|a\|$.
Jeżeli $a=a^*$, to $\sigma(a)\subseteq\R$.
\end{lemat}
\begin{proof}
Z własności $C^*$ wynika $\|b^*\|=\|b\|$:
$\|b\|^2=\|b^*b\|\le\|b^*\|\|b\|$,
a analogicznie po zamianie $b$ na $b^*$.
Dla samosprzężonego $h$ mamy
$\|h^2\|=\|h^*h\|=\|h\|^2$.
Jeżeli $a$ jest normalny, to
$\|(a^2)^*a^2\|=\|(a^*a)^2\|=\|a^*a\|^2$;
przez własność $C^*$ otrzymujemy $\|a^2\|=\|a\|^2$.
Indukcja daje $\|a^{2^n}\|=\|a\|^{2^n}$.
Wzór na promień spektralny implikuje $r(a)=\|a\|$.

Jeśli $h=h^*$, to $u_t=e^{ith}$ spełnia
$u_t^*u_t=u_tu_t^*=1_A$, więc $\|u_t\|=1$.
Na mocy twierdzenia o odwzorowaniu widma
$e^{it\lambda}\in\sigma(u_t)$ dla $\lambda\in\sigma(h)$.
Zatem $|e^{it\lambda}|\le1$ dla wszystkich $t\in\R$,
a po uwzględnieniu również $-t$ otrzymujemy
$\operatorname{Im}\lambda=0$.
\end{proof}

\begin{lemat}[Zespolone twierdzenie Stone'a--Weierstrassa]
\label{lem:stone-weierstrass31}
Niech $K$ będzie niepustą zwartą przestrzenią Hausdorffa, a
$B\subseteq C(K,\C)$ podalgebrą zawierającą funkcje stałe,
rozdzielającą punkty $K$ i zamkniętą na sprzężenie zespolone.
Wtedy $B$ jest gęsta w $C(K,\C)$ w normie supremum.
\end{lemat}
\begin{proof}
Najpierw ustalmy, że każdą rzeczywistą funkcję ciągłą na odcinku
można jednostajnie aproksymować wielomianami rzeczywistymi.
Po zmianie zmiennej wystarczy odcinek $[0,1]$. Dla $h\in C([0,1],\R)$
rozważamy wielomiany Bernsteina
$B_nh(t)=\sum_{k=0}^nh(k/n)\binom nk t^k(1-t)^{n-k}$.
Ponieważ wagi są nieujemne, sumują się do $1$ i mają wariancję
$t(1-t)/n\le1/(4n)$, jednostajna ciągłość $h$ i nierówność
Czebyszewa dają, przy oznaczeniu
$\omega_h(\delta)=\sup\{|h(s)-h(t)|:s,t\in[0,1],\ |s-t|\le\delta\}$,
oszacowanie
$\|B_nh-h\|_\infty\le\omega_h(\delta)
+\|h\|_\infty/(2n\delta^2)$ dla każdego $\delta>0$.
Wybierając najpierw małe $\delta$, a następnie duże $n$,
otrzymujemy zbieżność jednostajną.

Niech $B_\R=\{f\in B:f(K)\subseteq\R\}$.
Dzięki zamknięciu na sprzężenie, część rzeczywista i urojona
dowolnego elementu $B$ należą do $B_\R$.
Oznaczmy przez $D$ domknięcie $B_\R$ w $C(K,\R)$.
Z aproksymacji wielomianowej wynika, że $|u|\in D$ dla $u\in D$.
Stąd $\max(u,v)=(u+v+|u-v|)/2$ i
$\min(u,v)=(u+v-|u-v|)/2$ należą do $D$.

Ustalmy $h\in C(K,\R)$ i $\varepsilon>0$.
Dla każdych $x,y\in K$ rozdzielanie punktów i obecność stałych
dostarczają $u_{xy}\in B_\R$ z
$u_{xy}(x)=h(x)$ i $u_{xy}(y)=h(y)$.
Dla ustalonego $x$ z ciągłości wynika, że dla każdego $y$
istnieje otoczenie $V_y$ punktu $y$, na którym
$u_{xy}<h+\varepsilon$. Po wybraniu skończonego podpokrycia
odpowiednia funkcja $g_x=\min_j u_{xy_j}\in D$ spełnia
$g_x<h+\varepsilon$ na całym $K$ i $g_x(x)=h(x)$.
W pewnym otoczeniu $U_x$ punktu $x$ zachodzi także
$g_x>h-\varepsilon$. Wybierając skończenie wiele takich
otoczeń pokrywających $K$ i biorąc $g=\max_j g_{x_j}\in D$,
otrzymujemy $|g-h|<\varepsilon$ na $K$.
Zatem $D=C(K,\R)$. Każdą funkcję zespoloną można rozłożyć
na część rzeczywistą i urojoną, więc $B$ jest gęsta w $C(K,\C)$.
\end{proof}

\begin{twierdzenie}[Przemienne twierdzenie Gelfanda--Naimarka]
\label{thm:commutative-cstar31}
Niech $A$ będzie przemienną unitalną algebrą $C^*$.
Wówczas transformacja Gelfanda jest izometrycznym izomorfizmem
algebr z involucją
$A\simeq C(\Delta(A),\C)$,
przy czym $\widehat{a^*}=\overline{\widehat a}$.
\end{twierdzenie}
\begin{proof}
Dla elementu samosprzężonego $h$ mamy
$\chi(h)\in\sigma(h)\subseteq\R$ dla każdego charakteru $\chi$.
Każdy $a$ rozkłada się jako $a=h+ik$ z
$h=(a+a^*)/2$, $k=(a-a^*)/(2i)$ samosprzężonymi.
Zatem $\chi(a^*)=\overline{\chi(a)}$.
Ponieważ w przemiennej algebrze każdy element jest normalny,
lemat~\ref{lem:cstar-normal31} i twierdzenie~\ref{thm:gelfand-spectrum31}
dają $\|\widehat a\|_\infty=r(a)=\|a\|$.
Obraz transformacji jest więc domkniętą podalgebrą $C(\Delta(A))$,
zawierającą stałe, rozdzielającą punkty i zamkniętą na sprzężenie.
Z lematu~\ref{lem:stone-weierstrass31} wynika,
że jest gęsty w $C(\Delta(A))$, a więc równy całej tej przestrzeni.
Zatem odwzorowanie Gelfanda jest na $C(\Delta(A),\C)$.
\end{proof}

\begin{uwaga}[Granica zakresu]
Twierdzenie~\ref{thm:commutative-cstar31} dotyczy algebr $C^*$
\emph{przemiennych}. Pełne twierdzenie Gelfanda--Naimarka o
wiernej reprezentacji dowolnej algebry $C^*$ na przestrzeni Hilberta
wymaga m.in. konstrukcji GNS i będzie osobnym rozszerzeniem.
Nie utożsamiamy też rachunku holomorficznego z pełnym
ciągłym rachunkiem funkcyjnym dla operatorów normalnych;
ten drugi zostanie opracowany w \hyperref[sec:spectral-hilbert32]{rozdziale~32}.
\end{uwaga}


\subsection{Dwa przykłady pokazujące znaczenie struktury algebry}

\begin{przyklad}[Charaktery algebry splotowej $\ell^1(\N_0)$]
\label{ex:l1-characters31}
Niech $A=\ell^1(\N_0)$ z mnożeniem splotowym opisanym
w przykładzie~\ref{ex:banach-algebras31}. Oznaczmy przez $e_n$ ciąg,
którego $n$-ty wyraz jest równy $1$, a wszystkie pozostałe są zerowe.
Wówczas $e_n=e_1^{*n}$ i $e_0=1_A$.
Jeżeli $\chi\in\Delta(A)$, to $z:=\chi(e_1)$ spełnia
$|z|\le\|e_1\|_1=1$. Przez multiplikatywność i ciągłość charakteru
otrzymujemy dla dowolnego $a=(a_n)\in A$
\[
 \chi(a)=\sum_{n=0}^\infty a_nz^n.
\]
Odwrotnie, dla każdego $|z|\le1$ ten szereg określa charakter:
jest absolutnie zbieżny i zachowuje splot dzięki możliwości zamiany
kolejności sum bezwzględnie zbieżnych. Zatem przestrzeń charakterów
$\Delta(A)$ jest homeomorficzna z domkniętym dyskiem jednostkowym
(w topologii słabej-*), a transformacja Gelfanda przyporządkowuje
ciągowi $a$ funkcję $\widehat a(z)=\sum_{n\ge0}a_nz^n$.
W szczególności $\sigma_A(e_1)=\overline{\mathbb D}$,
choć $e_1$ nie jest elementem odwracalnym.
\end{przyklad}

\begin{przyklad}[Przesunięcie jednostronne — widmo niebędące zbiorem wartości własnych]
\label{ex:unilateral-shift31}
Na zespolonej przestrzeni Hilberta $H=\ell^2(\N)$ niech
$S(x_1,x_2,\ldots)=(0,x_1,x_2,\ldots)$.
Jest to izometria o normie $1$, więc $\sigma(S)\subseteq\overline{\mathbb D}$.
Dla każdego $|\lambda|<1$ operator sprzężony $S^*$ ma niezerowy
wektor własny $v=(1,\overline\lambda,\overline\lambda^2,\ldots)$
z wartością własną $\overline\lambda$.
Zatem obraz $\lambda I-S$ nie jest gęsty (jest zawarty w
$\{v\}^\perp$), a więc $\lambda\in\sigma(S)$.
Z domkniętości widma dostajemy $\sigma(S)=\overline{\mathbb D}$.
Jednak $S$ nie ma żadnych wartości własnych: równanie $Sx=\lambda x$
wymusza kolejno $x_1=0,x_2=0,\ldots$.
Pokazuje to, dlaczego w ogólnej teorii operatorów nie należy
utożsamiać widma ze zbiorem wartości własnych, co było możliwe
poza zerem w twierdzeniu Riesza--Schaudera dla operatorów zwartych.
\end{przyklad}


% Wersja źródłowa: audyt II, jednostka 32
\setcounter{section}{31}
\noindent\textbf{Wyniki wykorzystywane.} Zakładamy przemienne twierdzenie Gelfanda--Naimarka, twierdzenie Stone'a--Weierstrassa, twierdzenie Riesza--Markowa o reprezentacji dodatnich funkcjonałów na $C(K)$, własności całki Lebesgue'a i twierdzenie Riesza--Schaudera. Wszystkie inne kroki konstrukcji miary spektralnej są dowodzone w rozdziale.\par\medskip
\section{Teoria spektralna operatorów ograniczonych na przestrzeniach Hilberta}
\label{sec:spectral-hilbert32}

Poprzedni rozdział dostarczył holomorficznego rachunku funkcyjnego
w dowolnej zespolonej algebrze Banacha i twierdzenia Gelfanda--Naimarka
dla przemiennych algebr $C^*$. Dla operatorów normalnych na przestrzeniach
Hilberta otrzymamy znacznie więcej: rachunek wszystkich ograniczonych
funkcji borelowskich oraz reprezentację przez \emph{miarę spektralną}.
Nie zakładamy ośrodkowości przestrzeni Hilberta. Iloczyn skalarny jest
liniowy w pierwszym argumencie, tak jak w \hyperref[sec:hilbert13]{rozdziale~13}.

\subsection{Operatory normalne i ciągły rachunek funkcyjny}

\begin{definicja}[Operator normalny, samosprzężony i unitarny]
\label{def:normal32}
Niech $H$ będzie niezerową zespoloną przestrzenią Hilberta,
a $T\in\mathcal L(H,H)$. Operator $T$ jest \emph{normalny}, gdy
$TT^*=T^*T$, \emph{samosprzężony}, gdy $T=T^*$, oraz
\emph{unitarny}, gdy $T^*T=TT^*=I_H$. Operator samosprzężony
$T$ jest \emph{dodatni}, co zapisujemy $T\geq0$, gdy
$\langle Tx,x\rangle\geq0$ dla każdego $x\in H$.
\end{definicja}

\begin{lemat}[Zgodność widma z podalgebrą $C^*$]
\label{lem:spectral-permanence32}
Niech $H$ będzie niezerową zespoloną przestrzenią Hilberta,
$T\in\mathcal L(H,H)$ operatorem normalnym i
$A=C^*(I_H,T)$ domkniętą w normie podalgebrą generowaną
przez $I_H,T,T^*$. Wówczas $A$ jest przemienną unitalną
algebrą $C^*$ i $\sigma_A(T)=\sigma_{\mathcal L(H,H)}(T)$.
\end{lemat}
\begin{proof}
Normalność daje przemienność $T$ i $T^*$, a więc całej
podalgebry $A$. Zawsze $\sigma_{\mathcal L(H,H)}(T)\subseteq\sigma_A(T)$.
Niech odwrotnie $\lambda\notin\sigma_{\mathcal L(H,H)}(T)$ i
$B=T-\lambda I_H$. Wówczas istnieje $m>0$ takie, że
$\|Bx\|^2\geq m\|x\|^2$ dla wszystkich $x\in H$.
Operator $S=B^*B$ należy do $A$, jest dodatni i spełnia
$mI_H\leq S\leq\|S\|I_H$. Wykorzystując to, że dla operatora
samosprzężonego $R$ norma wynosi
$\sup_{\|x\|=1}|\langle Rx,x\rangle|$, otrzymujemy
$\|I_H-S/\|S\|\|\leq1-m/\|S\|<1$.
Szereg Neumanna dowodzi, że $S^{-1}\in A$. Ponieważ $B$ jest
normalny, $B$ i $B^*$ komutują z $S^{-1}$; zatem
$S^{-1}B^*$ jest obustronną odwrotnością $B$ w $A$.
W konsekwencji $\lambda\notin\sigma_A(T)$.
\end{proof}

\begin{twierdzenie}[Ciągły rachunek funkcyjny operatora normalnego]
\label{thm:continuous-calculus32}
Niech $H$ i $T$ będą jak w lemacie~\ref{lem:spectral-permanence32},
a $K=\sigma(T)\subset\C$. Istnieje dokładnie jeden izometryczny
unitalny homomorfizm algebr z involucją
$\Phi_T:C(K,\C)\to\mathcal L(H,H)$ taki, że
$\Phi_T(\mathrm{id}_K)=T$. Dla $f\in C(K,\C)$ piszemy $f(T)=\Phi_T(f)$.
Zachodzą
$\|f(T)\|=\max_{z\in K}|f(z)|$,
$\overline f(T)=f(T)^*$ i $\sigma(f(T))=f(K)$.
Jeżeli $p(z,\overline z)$ jest wielomianem w dwóch zmiennych,
to $p(T, T^*)=p|_K(T)$.
\end{twierdzenie}
\begin{proof}
Przemienne twierdzenie Gelfanda--Naimarka z
poprzedniego rozdziału daje izometryczny
izomorfizm $A\simeq C(\Delta(A),\C)$, zaś lemat
\ref{lem:spectral-permanence32} identyfikuje $K$ z widmem $T$
w $A$. Odwzorowanie
$\Delta(A)\ni\chi\mapsto\chi(T)\in K$ jest ciągłą surjekcją.
Jest różnowartościowe: charakter jest określony przez wartości
na $T,T^*$ i jednostce, a $\chi(T^*)=\overline{\chi(T)}$.
Jest więc homeomorfizmem zwartych przestrzeni Hausdorffa.
Przeniesienie transformacji Gelfanda daje $\Phi_T$.
Jednoznaczność wynika z gęstości wielomianów w $z,\overline z$
w $C(K)$, czyli z twierdzenia Stone'a--Weierstrassa.
Izometrię i równość widm otrzymujemy z twierdzenia
Gelfanda--Naimarka i zachowania widma w unitalnej
podalgebrze $C^*$ (ostatnią własność stosujemy do normalnego
operatora $f(T)$, jak w poprzednim lemacie).
\end{proof}

\begin{uwaga}[Związek z rachunkiem holomorficznym]
Jeżeli $f$ jest holomorficzna na otoczeniu $K$, to $f(T)$
z twierdzenia~\ref{thm:continuous-calculus32} zgadza się z całką
Dunforda z poprzedniego rozdziału.
Istotnie, dla $w\notin K$ rachunek ciągły daje
$(w-z)^{-1}(T)=(wI_H-T)^{-1}$. Jeżeli $\Gamma$ jest
konturem dopuszczalnym w definicji całki Dunforda,
ciągłość homomorfizmu $\Phi_T$ pozwala przenieść go
pod znak całki (jako granicę sum Riemanna). Wtedy
\[
 \frac1{2\pi i}\int_\Gamma f(w)(wI_H-T)^{-1}\,\mathrm dw
 =\Phi_T\!\left(\frac1{2\pi i}\int_\Gamma
 \frac{f(w)}{w-z}\,\mathrm dw\right)=\Phi_T(f|_K),
\]
gdzie ostatnia równość jest skalarnym wzorem Cauchy'ego
stosowanym punktowo dla $z\in K$.
Dla operatora nienormalnego sam rachunek
ciągły na $C(\sigma(T))$ z własnością $\overline f(T)=f(T)^*$
na ogół nie istnieje.
\end{uwaga}

\subsection{Miary spektralne i twierdzenie spektralne}

\begin{definicja}[Miara spektralna]
\label{def:pvm32}
Niech $H$ będzie zespoloną przestrzenią Hilberta i $K\subset\C$
niepustym zwartym zbiorem. \emph{Miarą spektralną} na
$\mathcal B(K)$ nazywamy odwzorowanie
$E:\mathcal B(K)\to\mathcal L(H,H)$ spełniające:
$E(B)$ jest rzutem ortogonalnym dla każdego zbioru borelowskiego $B$,
$E(\varnothing)=0$, $E(K)=I_H$ oraz
$E(B\cap C)=E(B)E(C)$ dla $B,C\in\mathcal B(K)$.
Ponadto dla parami rozłącznych $(B_n)_{n\geq1}$ wymaga się
\emph{ciągowej addytywności w silnej topologii operatorowej}, tj.
$E(\bigcup_nB_n)x=\lim_{N\to\infty}\sum_{n=1}^NE(B_n)x$
dla każdego $x\in H$.
Dla $x\in H$ związana z $E$ dodatnia skończona miara skalarna
to $\mu_x(B):=\langle E(B)x,x\rangle=\|E(B)x\|^2$.
\end{definicja}

\begin{lemat}[Model cykliczny]
\label{lem:cyclic-model32}
Niech $T\in\mathcal L(H,H)$ będzie normalny, $K=\sigma(T)$,
$a\in H\setminus\{0\}$ i
$H_a=\overline{\{f(T)a:f\in C(K,\C)\}}$.
Wówczas $H_a$ jest podprzestrzenią redukującą $T$
(niezmienniczą względem $T$ i $T^*$).
Istnieje jedyna skończona dodatnia regularna miara borelowska
$\nu_a$ na $K$ i izometria liniowa na
$U_a:L^2(K,\nu_a)\to H_a$ takie, że
$U_a1=a$ i $U_aM_z=(T|_{H_a})U_a$, gdzie
$(M_zf)(w)=wf(w)$.
\end{lemat}
\begin{proof}
Z domknięcia algebry $C(K)$ na sprzężenie otrzymujemy
$T H_a\subseteq H_a$ i $T^*H_a\subseteq H_a$.
Definiujemy $L_a(f)=\langle f(T)a,a\rangle$. Jest to ciągły
dodatni funkcjonał na $C(K)$, ponieważ dla $f\geq0$
możemy zapisać $f=|\sqrt f|^2$, a więc
$f(T)=(\sqrt f(T))^*\sqrt f(T)\geq0$.
Z twierdzenia Riesza--Markowa o reprezentacji dodatnich
funkcjonałów na $C(K)$ otrzymujemy jednoznaczną dodatnią,
skończoną regularną miarę borelowską $\nu_a$ spełniającą
$L_a(f)=\int_K f(z)\,\mathrm d\nu_a(z)$.
Mamy $\nu_a(K)=\|a\|^2$ oraz
\[
 \langle f(T)a,g(T)a\rangle
 =\langle (\overline g f)(T)a,a\rangle
 =\int_K f(z)\overline{g(z)}\,\mathrm d\nu_a(z).
\]
Zatem odwzorowanie $f\mapsto f(T)a$ jest izometrią z
$C(K)$ wyposażonej w normę $L^2(\nu_a)$ na podprzestrzeń $H$.
Funkcje ciągłe są gęste w $L^2(K,\nu_a)$: regularność
miary i lemat Urysohna pozwalają najpierw przybliżać
funkcje charakterystyczne zbiorów borelowskich, a następnie
funkcje proste. Izometria przedłuża się do $U_a$ na całym
$L^2(K,\nu_a)$; jej obraz jest domknięty i równy $H_a$.
Ponieważ $U_a(zf)=T U_a(f)$ dla funkcji ciągłych,
przez gęstość otrzymujemy równość operatorową.
\end{proof}

\begin{twierdzenie}[Twierdzenie spektralne dla operatora normalnego]
\label{thm:spectral-normal32}
Niech $H$ będzie niezerową zespoloną przestrzenią Hilberta,
$T\in\mathcal L(H,H)$ operatorem normalnym, a $K=\sigma(T)$.
Istnieje dokładnie jedna miara spektralna
$E:\mathcal B(K)\to\mathcal L(H,H)$ taka, że
$T=\int_Kz\,\mathrm dE(z)$.
Dla każdej ograniczonej borelowskiej funkcji $f:K\to\C$
istnieje operator $f(T)=\int_K f(z)\,\mathrm dE(z)$ spełniający
$\langle f(T)x,y\rangle=\int_K f(z)\,\mathrm d\mu_{x,y}(z)$,
gdzie $\mu_{x,y}(B):=\langle E(B)x,y\rangle$
jest skończoną miarą zespoloną, otrzymaną przez
polaryzację miar dodatnich $\mu_x$.
Zachodzą: $(fg)(T)=f(T)g(T)$,
$f(T)^*=\overline f(T)$ oraz
$\|f(T)x\|^2=\int_K|f(z)|^2\,\mathrm d\mu_x(z)$.
Rachunek ten przedłuża ciągły rachunek funkcyjny z
twierdzenia~\ref{thm:continuous-calculus32}.
\end{twierdzenie}
\begin{proof}
\textbf{Krok 1: rozkład na podprzestrzenie cykliczne.}
Rozważamy rodziny parami ortogonalnych, niezerowych podprzestrzeni
postaci $H_a$ z lematu~\ref{lem:cyclic-model32}. Są one redukujące.
Lemat Zorna daje rodzinę maksymalną $(H_{a_j})_{j\in J}$.
Gdyby ortogonalne dopełnienie ich sumy było niezerowe,
byłoby redukujące, a każdy niezerowy wektor z niego generowałby
kolejną podprzestrzeń cykliczną, wbrew maksymalności.
Zatem $H=\bigoplus_{j\in J}H_{a_j}$, także dla nieprzeliczalnego $J$.

\textbf{Krok 2: konstrukcja miary.}
Z lematu~\ref{lem:cyclic-model32} każdy $T|_{H_{a_j}}$
jest unitarnie równoważny operatorowi mnożenia $M_z$ na
$L^2(K,\nu_{a_j})$. Oznaczmy tę równoważność przez $U_j$.
Dla borelowskiego $B\subseteq K$ definiujemy
$E(B)|_{H_{a_j}}=U_jM_{\chi_B}U_j^{-1}$.
Są to rzuty ortogonalne, komutują i spełniają
$E(B\cap C)=E(B)E(C)$, $E(K)=I_H$. Jeżeli $(B_n)$ jest
rodziną parami rozłączną, to dla każdego $x_j\in H_{a_j}$
\[
 \bigl\|\bigl(E(\textstyle\bigcup_n B_n)
 -\sum_{n=1}^{N}E(B_n)\bigr)x_j\bigr\|^2
 =\int_{\bigcup_{n>N}B_n}|U_j^{-1}x_j|^2\,\mathrm d\nu_{a_j}
 \longrightarrow0.
\]
Każdy wektor $x\in\bigoplus_{j\in J}H_{a_j}$ ma co najwyżej
przeliczalnie wiele niezerowych składowych; sumując powyższe
oszacowania i stosując zbieżność zmajoryzowaną dla szeregu
nieujemnych składników, otrzymujemy silną ciągową addytywność.

\textbf{Krok 3: rachunek funkcyjny.}
Dla ograniczonej funkcji borelowskiej $f$ definiujemy
$f(T)|_{H_{a_j}}=U_jM_fU_j^{-1}$.
Normy tych operatorów nie przekraczają $\|f\|_\infty$,
więc określają jeden operator na całej sumie ortogonalnej.
Tożsamości algebraiczne i wzór na normę wektora wynikają
bezpośrednio z własności operatorów mnożenia.
Dla $f(z)=z$ dostajemy $T$. Dla funkcji ciągłych
otrzymujemy wcześniejszy rachunek, gdyż oba homomorfizmy
zgadzają się na $z$ i $\overline z$, a te generują
$C(K)$ jako zamkniętą algebrę.

\textbf{Krok 4: jednoznaczność.}
Jeżeli $F$ jest inną miarą spektralną o podanej własności,
to rachunek oparty na $F$ przyporządkowuje wielomianom
w $z,\overline z$ operatory $p(T,T^*)$.
Z twierdzenia Stone'a--Weierstrassa to samo zachodzi
dla wszystkich $f\in C(K)$. Dla dowolnych $x,y\in H$
miary zespolone $\langle E(\cdot)x,y\rangle$ i
$\langle F(\cdot)x,y\rangle$ mają więc identyczne całki
z wszystkich funkcji ciągłych. Każda z nich jest regularna:
wynika to przez polaryzację ze skończoności i regularności
miar dodatnich $\mu_x$ na zwartym zbiorze metrycznym $K$.
Jednoznaczność w twierdzeniu Riesza--Markowa daje
równość miar, a więc $E(B)=F(B)$ dla każdego zbioru
borelowskiego $B$.
\end{proof}

\begin{uwaga}[Zakres wykorzystanego twierdzenia miarowego]
Dowód twierdzenia spektralnego sprowadza jego część miarową
do \emph{twierdzenia Riesza--Markowa}: każdy dodatni ciągły
funkcjonał na $C(K)$ dla zwartej przestrzeni Hausdorffa $K$
ma jednoznaczną reprezentację przez skończoną regularną miarę
borelowską. Jest to rezultat skalarnej teorii miary, a nie
wniosek z samego twierdzenia Hahna--Banacha. W pozostałych
krokach dowodu konstrukcja miary projekcyjnej i jej jednoznaczność
zostały przeprowadzone jawnie. Nie używaliśmy założenia
ośrodkowości przestrzeni Hilberta.
\end{uwaga}

\subsection{Konsekwencje miary spektralnej}

\begin{twierdzenie}[Nośnik, widmo i wartości własne]
\label{thm:spectral-support32}
Niech $T$ i $E$ będą jak w twierdzeniu~\ref{thm:spectral-normal32}.
Wtedy $\sigma(T)$ jest najmniejszym domkniętym zbiorem
$F\subseteq K$ takim, że $E(F)=I_H$.
Dla każdego $\lambda\in\C$:
$\lambda$ jest wartością własną $T$ wtedy i tylko wtedy,
gdy $E(\{\lambda\}\cap K)\ne0$, a wówczas
$\ker(T-\lambda I_H)=E(\{\lambda\}\cap K)H$.
Dla $\lambda\in\sigma(T)$ i $\varepsilon>0$ istnieje
$\|x\|=1$ takie, że $\|(T-\lambda I_H)x\|<\varepsilon$.
\end{twierdzenie}
\begin{proof}
Jeśli $U\subseteq\C$ jest otwarty i $U\cap K\ne\varnothing$,
wybieramy niezerową funkcję $f\in C(K)$ o nośniku
zawartym w $U\cap K$ (z odpowiednią interpretacją
nośnika względnego w $K$).
Izometria rachunku ciągłego daje $f(T)\ne0$.
Gdyby $E(U\cap K)=0$, tożsamość
$f(T)=f(T)E(U\cap K)$ dawałaby sprzeczność.
Każdy względnie otwarty niepusty podzbiór $K$
ma zatem niezerową projekcję spektralną. Stąd
$K$ jest nośnikiem $E$.

Jeśli $Tx=\lambda x$, to
$0=\|(T-\lambda I_H)x\|^2
=\int_K|z-\lambda|^2\,\mathrm d\mu_x(z)$,
więc $x=E(\{\lambda\}\cap K)x$.
Odwrotnie, na obrazie tego rzutu operator $T$ działa
jako mnożenie przez $\lambda$.
Wybierając niezerowy wektor z
$E(K\cap B(\lambda,\varepsilon))H$ i normalizując go,
otrzymujemy wymagany wektor przybliżony dzięki wzorowi
na normę z twierdzenia~\ref{thm:spectral-normal32}.
\end{proof}

\begin{twierdzenie}[Borelowski rachunek funkcyjny i zbieżność]
\label{thm:borel-convergence32}
Niech $T\in\mathcal L(H,H)$ będzie normalny z miarą spektralną $E$
na $K=\sigma(T)$. Dla ograniczonej funkcji borelowskiej
$f:K\to\C$ operator $f(T)$ zależy jedynie od wartości $f$
poza zbiorem borelowskim $N$ z $E(N)=0$, a
\[
 \|f(T)\|=\inf\{M\ge0:E(\{z\in K:|f(z)|>M\})=0\}.
\]
Jeżeli $f_n$ są ograniczonymi funkcjami borelowskimi,
$\sup_n\|f_n\|_\infty<\infty$ oraz $f_n(z)\to f(z)$
poza zbiorem $E$-zerowym, to $f_n(T)x\to f(T)x$
w normie $H$ dla każdego $x\in H$.
W szczególności $0\le f\le g$ pociąga $f(T)\le g(T)$
dla rzeczywistych ograniczonych funkcji borelowskich $f,g$.
\end{twierdzenie}
\begin{proof}
Dla $x\in H$ wzór na normę daje
$\|(f_n(T)-f(T))x\|^2
=\int_K|f_n(z)-f(z)|^2\,\mathrm d\mu_x(z)$.
Zbieżność zmajoryzowana dowodzi tezy o zbieżności.
Jeżeli $E(\{|f|>M\})=0$, to całki normowe dają
$\|f(T)\|\le M$. Odwrotnie, gdy $E(\{|f|>M\})\ne0$,
możemy wybrać wektor jednostkowy z obrazu tego rzutu;
wtedy $\|f(T)x\|>M$. Przechodząc do infimum,
otrzymujemy wzór na normę.
Pozytywność wynika z
$\langle f(T)x,x\rangle=\int_Kf\,\mathrm d\mu_x$.
\end{proof}

\begin{wniosek}[Operatory samosprzężone, dodatnie i unitarne]
\label{cor:spectral-types32}
Niech $T\in\mathcal L(H,H)$ będzie normalny.
Wówczas $T=T^*$ wtedy i tylko wtedy, gdy
$\sigma(T)\subseteq\R$; jest dodatni wtedy i tylko wtedy, gdy
$\sigma(T)\subseteq[0,\infty)$, a unitarny wtedy i tylko wtedy,
gdy $\sigma(T)\subseteq\{z\in\C:|z|=1\}$.
Dla samosprzężonego $T$ projekcje
$E(( -\infty,t]\cap\sigma(T))$, $t\in\R$, tworzą
monotoniczną rodzinę rzutów opisującą rozkład spektralny.
\end{wniosek}
\begin{proof}
W modelu mnożenia $T^*$ odpowiada mnożeniu przez
$\overline z$, a $T^*T$ przez $|z|^2$. Dlatego odpowiednie
równości operatorowe i nierówność kwadratowa są równoważne
z warunkami na nośnik miary $E$, który jest równy $\sigma(T)$.
Dla operatora samosprzężonego monotoniczność wynika z
$E(B)E(C)=E(B)$ dla $B\subseteq C$.
\end{proof}

\begin{przyklad}[Widmo ciągłe operatora mnożenia]
\label{ex:mult-spectrum32}
W zespolonej przestrzeni Hilberta $H=L^2([0,1],\lambda_1)$
niech $(Tf)(t)=tf(t)$. Operator jest samosprzężony,
$\sigma(T)=[0,1]$, a miara spektralna ma postać
$(E(B)f)(t)=\chi_{B\cap[0,1]}(t)f(t)$.
Każda pojedyncza wartość $\lambda\in[0,1]$ ma miarę Lebesgue'a
równą $0$, więc $E(\{\lambda\})=0$ i $T$ nie ma wartości
własnych. Dowodzi to, że nawet dla operatora samosprzężonego
pełny rozkład na wektory własne nie musi istnieć.
\end{przyklad}

\begin{przyklad}[Nieośrodkowa przestrzeń Hilberta]
\label{ex:nonsep-spectrum32}
Niech $I$ będzie dowolnym niepustym zbiorem i $H=\ell^2(I)$.
Dla ograniczonej funkcji $a:I\to\C$ definiujemy
$Te_i=a(i)e_i$ na kanonicznej bazie ortonormalnej.
Wtedy $T$ jest normalny, $\sigma(T)=\overline{a(I)}$,
a $E(B)$ jest rzutem na domkniętą powłokę liniową
$\{e_i:a(i)\in B\}$. Każdy wektor ma co najwyżej
przeliczalnie wiele niezerowych współczynników ortonormalnych,
choć sam zbiór indeksów $I$ może być nieprzeliczalny.
\end{przyklad}

\subsection{Rodzina spektralna i rozkład polarny}

\begin{twierdzenie}[Rodzina spektralna operatora samosprzężonego]
\label{thm:spectral-family32}
Niech $T\in\mathcal L(H,H)$ będzie samosprzężony, $E$ jego
miarą spektralną na $K=\sigma(T)\subseteq\R$, a
$P_t=E(K\cap(-\infty,t])$ dla $t\in\R$.
Wtedy $P_sP_t=P_{\min(s,t)}$, $P_s\le P_t$ dla $s\le t$,
$\lim_{t\to-\infty}P_tx=0$ i
$\lim_{t\to+\infty}P_tx=x$ dla wszystkich $x\in H$.
Dla każdego $t\in\R$ i każdego ciągu $t_n\downarrow t$
zachodzi $P_{t_n}x\to P_tx$ w normie; dla $t_n\uparrow t$
granicą jest $E(K\cap(-\infty,t))x$.
Ponadto
\[
 \langle Tx,y\rangle=\int_{K}t\,\mathrm d\mu_{x,y}(t)
 \quad(x,y\in H),
\]
a różnica $P_t-E(K\cap(-\infty,t))=E(K\cap\{t\})$
jest rzutem na $\ker(T-tI_H)$; jej rząd (być może
nieskończony) jest krotnością geometryczną wartości własnej $t$.
\end{twierdzenie}
\begin{proof}
Monotoniczność i równość iloczynów wynikają z własności
$E(B\cap C)=E(B)E(C)$. Ponieważ $K$ jest zwarty,
dla dostatecznie małego $t$ mamy $P_t=0$, a dla dostatecznie
dużego $t$ mamy $P_t=I_H$.
Jeżeli $t_n\downarrow t$, to zbiory $K\cap(-\infty,t_n]$
maleją do $K\cap(-\infty,t]$. Dla każdego $x\in H$
\[
 \|(P_{t_n}-P_t)x\|^2
 =\mu_x(K\cap(t,t_n])\longrightarrow0
\]
na mocy ciągłości miary skończonej z góry.
Analogicznie dla $t_n<t$ i $t_n\uparrow t$ zbiory
$K\cap(-\infty,t_n]$ rosną do $K\cap(-\infty,t)$.
Wzór na parowanie jest szczególnym przypadkiem
rachunku spektralnego, a ostatnia równość wynika
z addytywności miary i z twierdzenia~\ref{thm:spectral-support32}.
\end{proof}

\begin{twierdzenie}[Rozkład polarny operatora ograniczonego]
\label{thm:polar-decomp32}
Niech $H$ będzie zespoloną przestrzenią Hilberta i
$T\in\mathcal L(H,H)$ dowolnym operatorem (niekoniecznie normalnym).
Istnieje dokładnie jedna \emph{częściowa izometria}
$U\in\mathcal L(H,H)$, której przestrzenią początkową jest
$(\ker T)^\perp$ i przestrzenią końcową
$\overline{\operatorname{Ran}T}$, taka że
$T=U|T|$, gdzie $|T|=(T^*T)^{1/2}$ jest jedynym dodatnim
pierwiastkiem z $T^*T$. Co więcej, $\ker U=\ker T$
i $U^*U$ jest rzutem ortogonalnym na $(\ker T)^\perp$.
Dla operatora normalnego $T$ czynniki $U$ i $|T|$ komutują.
\end{twierdzenie}
\begin{proof}
Ponieważ $T^*T\ge0$, ciągły rachunek funkcyjny
dostarcza dodatniego operatora $|T|$ spełniającego
$|T|^2=T^*T$. Wówczas
$\||T|x\|^2=\langle T^*Tx,x\rangle=\|Tx\|^2$, więc
$\ker|T|=\ker T$ oraz
$\overline{\operatorname{Ran}|T|}=(\ker T)^\perp$.
Na $\operatorname{Ran}|T|$ definiujemy $U_0(|T|x)=Tx$.
Jest to dobrze określona izometria, którą można przedłużyć
przez ciągłość na $(\ker T)^\perp$.
Na $\ker T$ kładziemy $U=0$. Otrzymujemy częściową
izometrię o wskazanych przestrzeniach, normie co najwyżej $1$
i z własnością $U|T|=T$.
Każda inna częściowa izometria spełniająca warunki tezy
musi zgadzać się z $U$ na gęstym zbiorze
$\operatorname{Ran}|T|$ w przestrzeni początkowej oraz znikać
na $\ker T$, co dowodzi jednoznaczności.
Jeżeli $T$ jest normalny, w modelu spektralnym
$|T|$ odpowiada mnożeniu przez $|z|$, a $U$ przez
funkcję $u(z)=z/|z|$ dla $z\ne0$ i $u(0)=0$.
Ponieważ funkcje $u(z)$ i $|z|$ komutują punktowo,
operatory $U$ i $|T|$ również komutują.
\end{proof}

\begin{przyklad}[Rozkład polarny macierzy]
Dla $T=\left(\begin{smallmatrix}0&2\\0&0\end{smallmatrix}\right)$
na $\C^2$ mamy
$|T|=\operatorname{diag}(0,2)$ oraz
$U=\left(\begin{smallmatrix}0&1\\0&0\end{smallmatrix}\right)$.
Zatem $T=U|T|$, a $U$ jest częściową izometrią;
nie jest operatorem unitarnym, gdyż $U^*U\ne I$.
\end{przyklad}

\subsection{Operator normalny zwarty jako przypadek dyskretny}

\begin{twierdzenie}[Twierdzenie spektralne dla zwartego operatora normalnego]
\label{thm:compact-normal32}
Niech $H$ będzie niezerową zespoloną przestrzenią Hilberta,
a $T\in\mathcal L(H,H)$ operatorem zwartym i normalnym.
Niezerowe wartości widma tworzą zbiór co najwyżej przeliczalny,
każda z nich jest wartością własną o skończonej krotności,
a jedynym możliwym punktem skupienia tego zbioru jest $0$.
Przestrzeń $H$ jest sumą ortogonalną przestrzeni własnych $T$,
włącznie z $\ker T$. Dla różnych niezerowych wartości
własnych $\lambda$ oznaczmy przez $P_\lambda$ rzut na
$\ker(T-\lambda I_H)$. Wtedy
$T=\sum_{\lambda\in\sigma(T)\setminus\{0\}}
\lambda P_\lambda$, przy czym szereg jest zbieżny w normie
operatorowej (sumy rozumiemy po skończonych podzbiorach indeksów).
\end{twierdzenie}
\begin{proof}
Z twierdzenia Riesza--Schaudera we wcześniejszym rozdziale o teorii Fredholma
niezerowe punkty widma operatora zwartego są izolowanymi
wartościami własnymi skończonej krotności i mają jedyny
możliwy punkt skupienia w zerze.
Z twierdzenia~\ref{thm:spectral-support32} każda niepusta
względnie otwarta część widma ma niezerową projekcję spektralną.
Poza zerem widmo jest zatem sumą pojedynczych punktów,
a odpowiadające im projekcje są ortogonalne.
Dopełnienie ortogonalne ich obrazów jest obrazem
$E(K\cap\{0\})$ i równa się $\ker T$.
Dla skończonego zbioru wartości własnych $F$ mamy
$\|T-\sum_{\lambda\in F}\lambda P_\lambda\|
=\sup\{|\lambda|:\lambda\in\sigma(T)\setminus(F\cup\{0\})\}$.
Ta wielkość dąży do zera, gdy $F$ obejmuje wszystkie
wartości własne o module większym od ustalonego $\varepsilon>0$.
W rezultacie otrzymujemy zbieżność w normie operatorowej.
\end{proof}

\begin{uwaga}[Granica twierdzenia]
Dla operatorów nienormalnych nie można w ogólności znaleźć
miary spektralnej projekcyjnej reprezentującej operator.
Na przykład niezerowa macierz
$\left(\begin{smallmatrix}0&1\\0&0\end{smallmatrix}\right)$
ma widmo $\{0\}$, lecz nie jest zerowa; gdyby miała reprezentację
$T=\int z\,\mathrm dE(z)$ z miarą spektralną na tym widmie,
byłaby równa $0$. Rozszerzenie rachunku spektralnego na
operatory nieograniczone opracujemy w \hyperref[sec:unbounded33]{rozdziale~33}.
\end{uwaga}


% Wersja źródłowa: audyt II, jednostka 33
\begingroup
\setlength{\parskip}{3pt}
\setcounter{section}{32}
\noindent\textbf{Wyniki wykorzystywane.} Z poprzednich rozdziałów korzystamy z twierdzenia o wykresie domkniętym, teorii operatorów sprzężonych na przestrzeniach Hilberta, twierdzenia Plancherela i twierdzenia spektralnego dla operatorów ograniczonych normalnych wraz z rachunkiem borelowskim. Dowody twierdzenia spektralnego dla operatorów nieograniczonych, twierdzenia von Neumanna o rozszerzeniach i twierdzenia Stone'a podajemy w tym rozdziale.\par\medskip
\section{Operatory nieograniczone na przestrzeniach Hilberta}
\label{sec:unbounded33}

W \hyperref[sec:spectral-hilbert32]{rozdziale~32} skonstruowaliśmy miarę spektralną
operatora \emph{ograniczonego} i normalnego. Operator różniczkowania na
$L^2$, operator mnożenia przez zmienną oraz generatory jednoparametrowych
grup unitarnych są na ogół nieograniczone. Zasadnicza różnica polega na tym,
że \emph{dziedzina operatora jest częścią jego definicji}. W całym rozdziale
$H$ oznacza niezerową zespoloną przestrzeń Hilberta, a iloczyn skalarny jest liniowy
w pierwszym argumencie. Zapis $\mathcal L(H)$ oznacza operatory ograniczone
na całym $H$.

\subsection{Dziedziny, wykresy, operatory sprzężone}

\begin{definicja}[Operator liniowy i jego wykres]
\label{def:unbounded-operator33}
Operatorem liniowym w $H$ nazywamy liniowe odwzorowanie
$A:D(A)\to H$, gdzie $D(A)\subseteq H$ jest podprzestrzenią liniową.
Jego \emph{wykresem} jest $G(A)=\{(x,Ax):x\in D(A)\}\subseteq H\times H$.
Operator jest \emph{domknięty}, gdy $G(A)$ jest domknięty, oraz
\emph{domykalny}, gdy domknięcie $G(A)$ jest wykresem pewnego operatora;
operator ten nazywamy \emph{domknięciem} $\overline A$.
Dla operatorów $A,B$ piszemy $A\subseteq B$, jeżeli
$D(A)\subseteq D(B)$ i $Bx=Ax$ dla $x\in D(A)$.
\end{definicja}

\begin{stwierdzenie}[Kryteria domkniętości i norma wykresowa]
\label{prop:graph33}
Niech $A:D(A)\subseteq H\to H$ będzie liniowy.
Wówczas $A$ jest domknięty wtedy i tylko wtedy, gdy dla każdego
ciągu $x_n\in D(A)$ z $x_n\to x$ i $Ax_n\to y$ w $H$ zachodzi
$x\in D(A)$ i $Ax=y$. Ponadto $A$ jest domknięty wtedy i tylko wtedy,
gdy $D(A)$ z normą $\|x\|_A=(\|x\|^2+\|Ax\|^2)^{1/2}$ jest
przestrzenią Hilberta z iloczynem skalarnym
$\langle x,y\rangle_A=\langle x,y\rangle+\langle Ax,Ay\rangle$.
Jeżeli $D(A)=H$, to operator domknięty jest ograniczony.
\end{stwierdzenie}
\begin{proof}
Pierwsza równoważność jest ciągowym kryterium domkniętości
podprzestrzeni przestrzeni metrycznej $H\times H$.
Odwzorowanie $D(A)\ni x\mapsto(x,Ax)\in G(A)$ jest izometrią
przy wskazanych iloczynach skalarnych, a $G(A)$ jest zupełny dokładnie
wtedy, gdy jest domknięty. Ostatnia teza wynika z twierdzenia
o wykresie domkniętym dla przestrzeni Banacha
(\hyperref[sec:frechet-open-graph]{rozdział~11}).
\end{proof}

\begin{definicja}[Zbiór rezolwenty operatora domkniętego]
Niech $A:D(A)\subseteq H\to H$ będzie operatorem domkniętym.
Do jego \emph{zbioru rezolwenty} $\rho(A)$ należą te $z\in\C$,
dla których $A-zI:D(A)\to H$ jest bijekcją, a jej
odwrotność $(A-zI)^{-1}:H\to H$ jest ograniczona.
\emph{Widmo} to $\sigma(A)=\C\setminus\rho(A)$.
Dla operatora domkniętego ograniczoność odwrotności wynika już
z jej istnienia i twierdzenia o wykresie domkniętym,
ale podajemy ją jawnie dla odróżnienia od operatorów niezamkniętych.
\end{definicja}

\begin{definicja}[Operator sprzężony]
\label{def:adjoint33}
Niech $A:D(A)\subseteq H\to H$ będzie liniowy i gęsto określony,
czyli $\overline{D(A)}=H$. Definiujemy
\[
 D(A^*)=\{y\in H:\text{istnieje }z\in H\text{ takie, że }
 \langle Ax,y\rangle=\langle x,z\rangle\ (x\in D(A))\},
 \qquad A^*y=z.
\]
Wektor $z$ jest jednoznaczny na mocy gęstości $D(A)$.
Operator $A$ jest \emph{symetryczny}, gdy $A\subseteq A^*$,
a \emph{samosprzężony}, gdy $A=A^*$, włącznie z równością dziedzin.
\end{definicja}

\begin{twierdzenie}[Sprzężenie i domykalność]
\label{thm:closable-adjoint33}
Niech $A:D(A)\subseteq H\to H$ będzie gęsto określony.
Wówczas $A^*$ jest operatorem domkniętym (jego dziedzina
nie musi być gęsta). Operator $A$ jest domykalny wtedy i tylko wtedy,
gdy $D(A^*)$ jest gęsta w $H$. W tym przypadku
$\overline A=A^{**}$. Każdy gęsto określony operator symetryczny
jest domykalny, a jego domknięcie również jest symetryczne.
\end{twierdzenie}
\begin{proof}
Jeżeli $y_n\to y$ i $A^*y_n\to z$, to dla każdego $x\in D(A)$
$\langle Ax,y\rangle=\lim_n\langle Ax,y_n\rangle
=\lim_n\langle x,A^*y_n\rangle=\langle x,z\rangle$;
stąd $y\in D(A^*)$ i $A^*y=z$.

Sprawdźmy kryterium domykalności. Dla wektora $(u,v)\in H\times H$
warunek $(u,v)\perp G(A)$ jest równoważny
$\langle u,x\rangle+\langle v,Ax\rangle=0$ dla wszystkich
$x\in D(A)$. Po sprzężeniu tej równości otrzymujemy
$v\in D(A^*)$ oraz $u=-A^*v$. Zatem
$G(A)^\perp=\{(-A^*v,v):v\in D(A^*)\}$.
W szczególności $(0,w)\in\overline{G(A)}$ wtedy i tylko wtedy,
gdy $\langle w,v\rangle=0$ dla wszystkich $v\in D(A^*)$.
Liniowa podprzestrzeń $\overline{G(A)}$ jest wykresem operatora
dokładnie wtedy, gdy nie zawiera niezerowych wektorów $(0,w)$.
Otrzymujemy równoważność z gęstością $D(A^*)$.
Gdy ta gęstość zachodzi, obliczenie prostopadłej do wykresu
$A^*$ daje $G(A^{**})=G(A)^{\perp\perp}=\overline{G(A)}$.
Jeżeli $A$ jest symetryczny, to $D(A)\subseteq D(A^*)$, a więc
$D(A^*)$ jest gęsta; przejście do granic w tożsamości
$\langle Ax,y\rangle=\langle x,Ay\rangle$ zachowuje symetrię.
\end{proof}

\begin{uwaga}[Operator ograniczony na niewłaściwej dziedzinie]
Nawet operator zerowy na gęstej właściwej podprzestrzeni $D\subsetneq H$
nie jest samosprzężony: jego operator sprzężony jest zerowy na całym
$H$, natomiast jego domknięcie jest operatorem zerowym na $H$.
Nie wystarczy zatem sprawdzać samego wzoru działania operatora.
\end{uwaga}

\subsection{Symetria, samosprzężoność i rezolwenta}

\begin{lemat}[Oszacowanie poza osią rzeczywistą]
\label{lem:symmetric-resolvent33}
Niech $A$ będzie gęsto określonym operatorem symetrycznym w $H$.
Dla $z=a+ib\in\C$ z $b\ne0$ i $x\in D(A)$ zachodzi
$\|(A-zI)x\|\ge |b|\|x\|$. Jeśli $A$ jest ponadto domknięty,
$\operatorname{Ran}(A-zI)$ jest podprzestrzenią domkniętą.
\end{lemat}
\begin{proof}
Ponieważ $\langle Ax,x\rangle$ jest rzeczywiste,
$|\operatorname{Im}\langle(A-zI)x,x\rangle|=|b|\|x\|^2$.
Nierówność Cauchy'ego--Schwarza daje oszacowanie.
Gdy $A$ jest domknięty, oszacowanie implikuje, że
$(A-zI)x_n\to y$ pociąga zbieżność $x_n\to x$;
domkniętość operatora daje $x\in D(A)$ i $(A-zI)x=y$.
\end{proof}

\begin{twierdzenie}[Kryterium samosprzężoności]
\label{thm:selfadjoint-range33}
Niech $A$ będzie gęsto określonym, symetrycznym operatorem w $H$.
Wówczas następujące warunki są równoważne:
\begin{enumerate}[label=\textup{(\roman*)}]
\item $A$ jest samosprzężony;
\item $\operatorname{Ran}(A+iI)=H$ i $\operatorname{Ran}(A-iI)=H$.
\end{enumerate}
Jeśli warunki te zachodzą, dla każdego $z\in\C\setminus\R$
operator $A-zI:D(A)\to H$ jest bijekcją, a jego odwrotność
$R(z,A)=(A-zI)^{-1}\in\mathcal L(H)$ spełnia
$\|R(z,A)\|\le1/|\operatorname{Im}z|$.
\end{twierdzenie}
\begin{proof}
Jeżeli $A=A^*$, to $A$ jest domknięty. Z lematu
\ref{lem:symmetric-resolvent33} oba obrazy $A\pm iI$ są domknięte.
Ich ortogonalne dopełnienia wynoszą odpowiednio
$\ker(A^*\mp iI)$, a te są zerowe: dla $y\in D(A)$
równość $(A\mp iI)y=0$ i oszacowanie z lematu wymuszają $y=0$.
Obrazy są więc również gęste, a zatem równe $H$.

Odwrotnie, załóżmy surjektywność obu operatorów.
Wówczas $A$ jest domknięty: jeśli $x_n\to x$ i $Ax_n\to y$,
to $(A+iI)x_n\to y+ix$; wybierając $u\in D(A)$
z $(A+iI)u=y+ix$ i stosując oszacowanie z lematu,
otrzymujemy $x_n\to u$, więc $x=u\in D(A)$ i $Ax=y$.
Dla $y\in D(A^*)$ wybieramy $u\in D(A)$ z
$(A+iI)u=(A^*+iI)y$. Wówczas
$(A^*+iI)(y-u)=0$, a
$\ker(A^*+iI)=\operatorname{Ran}(A-iI)^\perp=\{0\}$.
Stąd $y=u\in D(A)$ i $A=A^*$.

Dla dowolnego $z\notin\R$ ponawiamy argument z domkniętością
obrazu i jego gęstością, używając $A=A^*$ oraz oszacowania
z lematu. Oszacowanie normy odwrotności jest natychmiastowe.
\end{proof}

\begin{przyklad}[Mnożenie przez funkcję rzeczywistą]
\label{ex:mult-unbounded33}
Niech $(X,\Sigma,\mu)$ będzie przestrzenią z miarą, a
$m:X\to\R$ funkcją mierzalną, skończoną prawie wszędzie.
W przestrzeni Hilberta $H=L^2(X,\mu)$ operator
$(M_mf)(x)=m(x)f(x)$ z dziedziną
$D(M_m)=\{f\in L^2(X,\mu):mf\in L^2(X,\mu)\}$
jest gęsto określony i samosprzężony. Istotnie, funkcje
$f\chi_{\{|m|\le n\}}$ należą do $D(M_m)$ i zbiegają do $f$
w $L^2$. Dla każdego $g\in H$ funkcje
$f_\pm(x)=g(x)/(m(x)\pm i)$ należą do $D(M_m)$,
ponieważ $|m/(m\pm i)|\le1$, i spełniają
$(M_m\pm iI)f_\pm=g$. Operator jest symetryczny, więc
poprzednie twierdzenie daje samosprzężoność.
Dla $X=\R$, $m(x)=x$ operator jest nieograniczony.
\end{przyklad}

\subsection{Miary spektralne i rachunek funkcji nieograniczonych}

\begin{lemat}[Całkowanie nieograniczonych funkcji względem miary spektralnej]
\label{lem:unbounded-integral33}
Niech $K\subseteq\C$ będzie niepustym zbiorem borelowskim
(niekoniecznie zwartym), $H$ zespoloną przestrzenią Hilberta,
a $E:\mathcal B(K)\to\mathcal L(H)$ miarą spektralną:
$E(K)=I_H$, każde $E(B)$ jest rzutem ortogonalnym,
$E(B\cap C)=E(B)E(C)$ i $E$ jest przeliczalnie addytywna
w silnej topologii operatorowej. Niech $f:K\to\C$ będzie
funkcją borelowską, niekoniecznie ograniczoną.
Wówczas istnieje gęsto określony operator domknięty $f(E)$
o dziedzinie
\[
 D(f(E))=\left\{x\in H:\int_K|f(z)|^2\,\mathrm d\mu_x(z)<\infty\right\},
 \qquad \|f(E)x\|^2=\int_K|f(z)|^2\,\mathrm d\mu_x(z),
\]
gdzie $\mu_x(B)=\langle E(B)x,x\rangle$.
Zachodzi $f(E)^*=\overline f(E)$; w szczególności dla
$f$ rzeczywistej $E$-prawie wszędzie operator $f(E)$ jest
samosprzężony. Dla ograniczonej borelowskiej funkcji $g$
operator $g(E)$ zachowuje $D(f(E))$ oraz
$f(E)g(E)x=g(E)f(E)x=(fg)(E)x$ dla $x\in D(f(E))$.
Dla ograniczonej $f$ konstrukcja jest zwykłym borelowskim
rachunkiem funkcyjnym.
\end{lemat}
\begin{proof}
\textbf{Krok 1: rachunek funkcji ograniczonych.}
Dla funkcji prostej $s=\sum_{j=1}^m c_j\chi_{B_j}$, gdzie
$B_j$ są parami rozłączne, definiujemy
$s(E)=\sum_{j=1}^m c_jE(B_j)$.
Własności rzutów dają
$\|s(E)x\|^2=\int_K|s|^2\,\mathrm d\mu_x$,
$(st)(E)=s(E)t(E)$ i $s(E)^*=\overline s(E)$
dla funkcji prostych. Każdą ograniczoną funkcję borelowską
można jednostajnie aproksymować funkcjami prostymi.
Z nierówności $\|s(E)\|\le\|s\|_\infty$ otrzymujemy
przez granicę w normie operatorowej rachunek wszystkich
ograniczonych funkcji borelowskich, z tymi samymi
własnościami algebraicznymi i normowymi.

\textbf{Krok 2: dziedzina i domkniętość.}
Funkcja $f$ jest skończona punktowo; rachunek
nie zmieni się po modyfikacji $f$ na zbiorze $E$-zerowym.
Kładziemy $B_n=\{z\in K:|f(z)|\le n\}$,
$P_n=E(B_n)$ i $f_n=f\chi_{B_n}$.
Wtedy $P_n\to I_H$ silnie, a $f_n$ są ograniczone.
Jeżeli $x$ należy do przestrzeni $D$ określonej całką w tezie,
to
\[
 \|(f_n(E)-f_m(E))x\|^2
 =\int_K|f_n-f_m|^2\,\mathrm d\mu_x\longrightarrow0.
\]
Definiujemy $f(E)x=\lim_n f_n(E)x$.
Odwrotnie, gdy $f_n(E)x$ ma granicę, to jego normy są
ograniczone, a monotoniczna zbieżność całek
$\int|f_n|^2\,\mathrm d\mu_x$ dowodzi $x\in D$.
Ponieważ $P_nH\subseteq D$ i $P_n\to I_H$ silnie,
dziedzina $D$ jest gęsta. Tożsamość na normę wynika
z monotonicznej zbieżności całek.

Aby wykazać domkniętość, przypuśćmy, że $x_j\in D$,
$x_j\to x$ oraz $f(E)x_j\to y$. Dla każdego $n$
z rachunku ograniczonego wynika
$P_nf(E)x_j=f_n(E)x_j$.
Przechodząc do granicy przy ustalonym $n$ otrzymujemy
$P_ny=f_n(E)x$, a więc $\|f_n(E)x\|\le\|y\|$.
Monotoniczna zbieżność całek pokazuje, że $x\in D$.
Następnie $f_n(E)x=P_ny\to y$, czyli $f(E)x=y$.

\textbf{Krok 3: operator sprzężony i mnożenie.}
Dla $x\in D(f(E))$ i $y\in D(\overline f(E))$,
przechodząc do granicy w tożsamości dla $f_n(E)$,
dostajemy $\langle f(E)x,y\rangle
=\langle x,\overline f(E)y\rangle$.
Zatem $\overline f(E)\subseteq f(E)^*$.
Jeśli zaś $y\in D(f(E)^*)$ i $z=f(E)^*y$, to
testując tożsamość sprzężenia na $P_nx$ dla dowolnego
$x\in H$ otrzymujemy
$\overline{f_n}(E)y=P_nz$.
Stąd $\|\overline{f_n}(E)y\|\le\|z\|$.
Monotoniczna zbieżność całek implikuje
$y\in D(\overline f(E))$, a granica daje
$\overline f(E)y=z$. Zatem $f(E)^*=\overline f(E)$.
Jeżeli $g$ jest ograniczona, to
$\int|fg|^2\,\mathrm d\mu_x\le
\|g\|_\infty^2\int|f|^2\,\mathrm d\mu_x$.
Przechodząc do granicy w tożsamościach dla
$f_n(E)$ oraz $g(E)$ uzyskujemy ostatnią tezę.
\end{proof}

\begin{twierdzenie}[Twierdzenie spektralne dla operatora samosprzężonego]
\label{thm:unbounded-spectral33}
Niech $A:D(A)\subseteq H\to H$ będzie gęsto określonym
operatorem samosprzężonym. Istnieje dokładnie jedna miara spektralna
$E:\mathcal B(\R)\to\mathcal L(H)$ taka, że
\[
 D(A)=\left\{x\in H:\int_\R t^2\,\mathrm d\mu_x(t)<\infty\right\},
 \qquad Ax=\int_\R t\,\mathrm dE(t)\,x\quad(x\in D(A)).
\]
Dla każdego $z\in\C\setminus\R$ zachodzi
$(A-zI)^{-1}=\int_\R(t-z)^{-1}\,\mathrm dE(t)$.
Dla dowolnej borelowskiej funkcji $f:\R\to\C$
definiujemy $f(A)$ według
lematu~\ref{lem:unbounded-integral33}; operator ten jest domknięty,
a jego dziedzina jest wyznaczona przez całkę $\int|f|^2\,\mathrm d\mu_x$.
\end{twierdzenie}
\begin{proof}
\textbf{Krok 1: transformacja Cayleya.}
Z twierdzenia~\ref{thm:selfadjoint-range33} operatory
$A\pm iI$ odwzorowują bijektywnie $D(A)$ na $H$.
Definiujemy ograniczony operator
$U=(A-iI)(A+iI)^{-1}=I_H-2i(A+iI)^{-1}$.
Dla $y=(A+iI)x$ symetria daje
$\|Uy\|=\|(A-iI)x\|=\|(A+iI)x\|=\|y\|$.
Surjektywność $A-iI$ pokazuje, że $U$ jest unitarny.
Ponadto $I_H-U=2i(A+iI)^{-1}$ jest różnowartościowy,
a zatem $\ker(I_H-U)=\{0\}$.

\textbf{Krok 2: przeniesienie miary spektralnej.}
Niech $P$ będzie jedyną miarą spektralną operatora
unitarnego $U$ na $\sigma(U)\subseteq\mathbb T$.
Ponieważ $\ker(I_H-U)=P(\{1\})H=\{0\}$,
mamy $P(\{1\})=0$. Definiujemy rzeczywistą funkcję borelowską
$h(w)=i(1+w)/(1-w)$ dla $w\ne1$ na okręgu jednostkowym.
Jej wartości są rzeczywiste: dla $w=e^{i\theta}\ne1$
liczba $h(w)=-\cot(\theta/2)\in\R$.
Wartość $h(1)$ ustalamy dowolnie, ponieważ $P(\{1\})=0$.
Kładziemy $E(B)=P(h^{-1}(B))$ dla $B\in\mathcal B(\R)$.
Otrzymujemy miarę spektralną na $\R$.

\textbf{Krok 3: identyfikacja operatora.}
Operator $B=h(P)$ z lematu~\ref{lem:unbounded-integral33}
jest samosprzężony. Dla $w\ne1$ zachodzi
$(h(w)+i)^{-1}=(1-w)/(2i)$.
Funkcja po prawej jest ograniczona i mnożenie przez nią
odwzorowuje $H$ na $D(B)$: istotnie, iloczyn z $h+i$
jest równy $1$ poza zbiorem $P$-zerowym.
Rachunek spektralny daje więc obustronną odwrotność
$(B+iI)^{-1}=(I_H-U)/(2i)=(A+iI)^{-1}$.
Odwrotności wyznaczają jednoznacznie dziedziny i działania
operatorów, więc $A=B$. Wzór na dziedzinę i całkę
spektralną wynika bezpośrednio z lematu.

\textbf{Krok 4: jednoznaczność.}
Jeżeli $F$ jest inną miarą spektralną reprezentującą $A$,
to ograniczony rachunek borelowski daje
$U=(A-iI)(A+iI)^{-1}=\int_\R (t-i)/(t+i)\,\mathrm dF(t)$.
Jednoznaczność miary spektralnej operatora normalnego
$U$ z \hyperref[sec:spectral-hilbert32]{rozdziału~32} wymusza, że obraz $F$ przez funkcję
$t\mapsto(t-i)/(t+i)$ jest równy $P$.
Funkcja ta jest różnowartościowa z $\R$ na
$\mathbb T\setminus\{1\}$, o odwrotności $h$, więc $F=E$.
Równość na rezolwencie jest wzorem rachunku borelowskiego
z funkcją $(t-z)^{-1}$, która jest ograniczona dla $z\notin\R$.
\end{proof}

\begin{wniosek}[Widmo i przestrzenie spektralne]
\label{cor:unbounded-spectrum33}
Niech $A$ będzie samosprzężony z miarą spektralną $E$.
Wówczas $\sigma(A)\subseteq\R$ i jest równe nośnikowi $E$,
tj. najmniejszemu domkniętemu zbiorowi $F\subseteq\R$
z $E(F)=I_H$. Dla $\lambda\in\R$ zachodzi
$\ker(A-\lambda I)=E(\{\lambda\})H$.
Jeżeli $\lambda\in\sigma(A)$, to dla każdego $\varepsilon>0$
istnieje $x\in D(A)$ o normie $1$ taki, że
$\|(A-\lambda I)x\|<\varepsilon$.
\end{wniosek}
\begin{proof}
Dla $\lambda\notin\operatorname{supp}E$ istnieje $\delta>0$
takie, że $E((\lambda-\delta,\lambda+\delta))=0$.
Funkcja $(t-\lambda)^{-1}$ jest wtedy ograniczona
$E$-prawie wszędzie; jej rachunek spektralny daje ograniczoną
odwrotność $A-\lambda I$. Odwrotnie, jeśli
$\lambda\in\operatorname{supp}E$, to
$E((\lambda-\varepsilon,\lambda+\varepsilon))\ne0$.
Wybieramy jednostkowy wektor $x$ z obrazu tego rzutu;
ma on ograniczone spektralnie wsparcie, a więc należy do $D(A)$,
i spełnia żądane oszacowanie. Operator nie może mieć wtedy
ograniczonej odwrotności. Tożsamość dotycząca jądra wynika
z $\|(A-\lambda I)x\|^2=\int|t-\lambda|^2\,\mathrm d\mu_x(t)$.
\end{proof}

\begin{przyklad}[Różne dziedziny tego samego wzoru]
Niech $H=L^2(\R,\lambda_1)$.
Operator $Af=xf$ z dziedziną $D(A)=\{f:xf\in L^2\}$ jest
samosprzężony i ma widmo $\R$. Dla operatora
$Bf=-if'$ z dziedziną $H^1(\R)=\{f\in L^2(\R):f'\in L^2(\R)\}$
(pochodna rozumiana w sensie dystrybucyjnym) transformacja Fouriera
z \hyperref[sec:gladkie-schwartz17]{rozdziału~17} daje
$\widehat{Bf}(\xi)=\xi\widehat f(\xi)$.
Ponieważ transformacja jest unitarna na $L^2$, $B$ jest
samosprzężony i ma widmo $\R$. Tymczasem ograniczenie
operatora $-if'$ do przestrzeni $C_c^\infty(\R)$
jest tylko symetryczne i domykalne, a nie domknięte.
\end{przyklad}

\subsection{Rozszerzenia samosprzężone i indeksy defektu}

\begin{definicja}[Przestrzenie defektu]
\label{def:deficiency33}
Dla gęsto określonego domkniętego operatora symetrycznego
$S$ definiujemy przestrzenie defektu
$N_+=\ker(S^*-iI)$ i $N_-=\ker(S^*+iI)$
oraz indeksy defektu $n_\pm=\dim N_\pm$ rozumiane jako liczby
kardynalne, tj. moce baz ortonormalnych (w szczególności
dopuszczamy wymiary nieprzeliczalne).
\end{definicja}

\begin{twierdzenie}[Rozkład von Neumanna i rozszerzenia]
\label{thm:von-neumann33}
Niech $S$ będzie gęsto określonym domkniętym operatorem
symetrycznym w zespolonej przestrzeni Hilberta $H$.
W normie wykresowej $D(S^*)$ zachodzi ortogonalny rozkład
\[
 D(S^*)=D(S)\oplus N_+\oplus N_-.
\]
Operator $S$ jest samosprzężony wtedy i tylko wtedy,
gdy $N_+=N_-=\{0\}$. Istnieje rozszerzenie samosprzężone
$S\subseteq\widetilde S=\widetilde S^*$ wtedy i tylko wtedy,
gdy $\dim N_+=\dim N_-$. Wszystkie takie rozszerzenia są
parametryzowane operatorami unitarnymi $V:N_+\to N_-$:
\[
 D(S_V)=D(S)\dotplus\{u+Vu:u\in N_+\},\qquad
 S_V(x+u+Vu)=Sx+iu-iVu.
\]
\end{twierdzenie}
\begin{proof}
Operator $S^*$ jest domknięty, a $D(S)$ jest domkniętą
podprzestrzenią przestrzeni Hilberta $D(S^*)$ z normą wykresową.
Dla $x\in D(S)$ i $u\in N_\pm$ otrzymujemy, korzystając
z $S\subseteq S^*$,
$\langle x,u\rangle_{S^*}=\langle x,u\rangle+
\langle Sx,\pm i u\rangle=0$.
Również $N_+\perp N_-$ w iloczynie wykresowym.

Niech $z\in D(S^*)$ będzie prostopadły w normie wykresowej
do $D(S)$. To znaczy
$\langle x,z\rangle+\langle Sx,S^*z\rangle=0$
dla każdego $x\in D(S)$. Z definicji operatora sprzężonego
wynika $S^*z\in D(S^*)$ oraz $S^{*2}z=-z$.
Stąd $z_\pm=(z\mp iS^*z)/2$ należą odpowiednio do $N_\pm$
i $z=z_++z_-$. Otrzymujemy cały rozkład.

Dla wektorów $u,v\in N_+$ i $p,q\in N_-$
forma brzegowa
$\mathfrak b(w,z)=\langle S^*w,z\rangle-
\langle w,S^*z\rangle$ spełnia
$\mathfrak b(u+p,v+q)=2i(\langle u,v\rangle-\langle p,q\rangle)$.
Każde domknięte rozszerzenie $S\subseteq T\subseteq S^*$
odpowiada domkniętej podprzestrzeni
$L\subseteq N_+\oplus N_-$ przez
$D(T)=D(S)\oplus L$. Operator jest symetryczny dokładnie
wtedy, gdy $\mathfrak b$ znika na $L\times L$.
Ponadto $D(T^*)$ odpowiada anihilatorowi
$L^{\perp_{\mathfrak b}}$ (wynika to wprost z definicji
sprzężenia), więc samosprzężoność oznacza
$L=L^{\perp_{\mathfrak b}}$.

Niech najpierw $L$ odpowiada operatorowi samosprzężonemu.
Jeżeli $u+p\in L$, $u\in N_+$, $p\in N_-$, to
$\mathfrak b(u+p,u+p)=0$ daje $\|u\|=\|p\|$.
Rzut $L$ na $N_+$ jest więc różnowartościowy, a $L$
ma postać wykresu pewnej izometrii
$V:D(V)\subseteq N_+\to N_-$. Podprzestrzeń $D(V)$
jest domknięta: jeśli $u_n\to u$ w $N_+$, to
$Vu_n$ jest ciągiem Cauchy'ego i zbiega do pewnego $p\in N_-$.
Ponieważ $L$ jest domknięta w $N_+\oplus N_-$,
$u+p\in L$, a więc $u\in D(V)$ i $Vu=p$. Jeśli $D(V)\ne N_+$,
istniałby $0\ne w\in D(V)^\perp\subset N_+$.
Wówczas $\mathfrak b(w,l)=0$ dla każdego $l\in L$, ale
$w\notin L$, wbrew $L=L^{\perp_{\mathfrak b}}$.
A więc $D(V)=N_+$. Analogicznie obraz $V$ jest domknięty,
a gdyby nie był całym $N_-$, wektor
$0\ne q\in (\operatorname{Ran}V)^\perp$ należałby do
$L^{\perp_{\mathfrak b}}\setminus L$. Zatem $V$ jest unitarny.
Odwrotnie, dla dowolnego unitarnego $V:N_+\to N_-$
jego wykres $L=\{u+Vu:u\in N_+\}$ jest domknięty
i spełnia $L=L^{\perp_{\mathfrak b}}$, co daje
samosprzężone rozszerzenie o wskazanej dziedzinie.
Gdy $N_+=N_- =\{0\}$, rozkład dziedziny daje
$D(S^*)=D(S)$, a więc $S=S^*$.
\end{proof}

\begin{przyklad}[Pochodna na przedziale i warunki brzegowe]
\label{ex:momentum-interval33}
Na $H=L^2(0,1)$ rozważmy
$Sf=-if'$ z dziedziną
$D(S)=\{f\in H^1(0,1):f(0)=f(1)=0\}$.
Przestrzeń $H^1(0,1)$ oznacza funkcje z $L^2(0,1)$
o słabej pochodnej w $L^2$, identyfikowane z ich
absolutnie ciągłymi reprezentantami.
Całkowanie przez części daje symetrię, a domkniętość
wynika z zupełności normy $H^1$.
Operator sprzężony ma ten sam wzór $S^*f=-if'$, lecz
$D(S^*)=H^1(0,1)$. Istotnie, testowanie warunku
sprzężenia na funkcjach z $C_c^\infty(0,1)$ wymusza istnienie
słabej pochodnej w $L^2$, a zatem przynależność do $H^1$.
Odwrotnie, całkowanie przez części dla dwóch funkcji z $H^1$
pokazuje, że wyrazy brzegowe znikają, ponieważ
wektor z $D(S)$ ma zerowe wartości na obu końcach.
Równania $S^*f=if$ i $S^*f=-if$ mają odpowiednio
rozwiązania $f(x)=ce^{-x}$ i $f(x)=ce^x$.
Zatem $n_+=n_-=1$. Wzór na formę brzegową brzmi
$\mathfrak b(f,g)=-i(f(1)\overline{g(1)}-f(0)\overline{g(0)})$.
Rozszerzenia samosprzężone są dokładnie operatorami
$-if'$ o dziedzinach
$D(S_\theta)=\{f\in H^1(0,1):f(1)=e^{i\theta}f(0)\}$,
$\theta\in[0,2\pi)$. Są to różne realizacje \emph{tego samego wzoru}.
\end{przyklad}

\begin{przyklad}[Brak rozszerzenia samosprzężonego]
Na $L^2(0,\infty)$ operator $-if'$ z dziedziną
$H_0^1(0,\infty)=\{f\in H^1(0,\infty):f(0)=0\}$
jest domknięty i symetryczny. Przyjmujemy, że $H^1(0,\infty)$
składa się z funkcji $L^2$ o słabej pochodnej w $L^2$;
każda z nich ma lokalnie absolutnie ciągłego reprezentanta. Jego operator sprzężony ma
dziedzinę $H^1(0,\infty)$. Z równań defektu otrzymujemy
$N_+=\lin\{e^{-x}\}$ i $N_- =\{0\}$, ponieważ $e^x\notin L^2$.
Indeksy defektu różnią się, więc operator nie ma rozszerzenia
samosprzężonego w tej samej przestrzeni Hilberta.
\end{przyklad}

\subsection{Jednoparametrowe grupy unitarne i twierdzenie Stone'a}

\begin{definicja}[Silnie ciągła grupa unitarna]
Niech $(U_t)_{t\in\R}$ będzie rodziną operatorów unitarnych
na $H$ z $U_0=I_H$ i $U_{t+s}=U_tU_s$ dla $s,t\in\R$.
Jest to \emph{silnie ciągła jednoparametrowa grupa unitarna}, gdy
$\lim_{t\to0}\|U_tx-x\|=0$ dla każdego $x\in H$.
Jej generator definiujemy na dziedzinie
$D(G)=\{x:\lim_{t\to0}(U_tx-x)/t\text{ istnieje w }H\}$
przez $Gx=\lim_{t\to0}(U_tx-x)/t$.
\end{definicja}

\begin{twierdzenie}[Twierdzenie Stone'a]
\label{thm:stone33}
Dla każdej silnie ciągłej jednoparametrowej grupy unitarnej
$(U_t)_{t\in\R}$ na zespolonej przestrzeni Hilberta $H$
istnieje dokładnie jeden operator samosprzężony $A$ taki, że
$U_t=e^{itA}$ dla wszystkich $t\in\R$, gdzie
$e^{itA}=\int_\R e^{it\lambda}\,\mathrm dE_A(\lambda)$.
Ponadto $D(A)=D(G)$ i $G=iA$.
Odwrotnie, dla każdego samosprzężonego $A$ rachunek spektralny
wyznacza silnie ciągłą grupę unitarną $(e^{itA})_{t\in\R}$,
a jej generator jest równy $iA$.
\end{twierdzenie}
\begin{proof}
\textbf{Krok 1: konstrukcja grupy z operatora.}
Dla samosprzężonego $A$ funkcja $\lambda\mapsto e^{it\lambda}$
jest ograniczona, zatem definiuje operator unitarny
$U_t=e^{itA}$ poprzez rachunek borelowski.
Równości $U_{t+s}=U_tU_s$ i $U_0=I_H$ są punktowymi
tożsamościami funkcji w rachunku spektralnym.
Dla każdego $x\in H$ mamy
$\|(U_t-I_H)x\|^2=\int_\R|e^{it\lambda}-1|^2\,\mathrm d\mu_x(\lambda)
\to0$ na mocy zbieżności zmajoryzowanej.
Dla $x\in D(A)$ nierówność
$|(e^{it\lambda}-1)/t|\le|\lambda|$ i zbieżność punktowa
ilorazu różnicowego do $i\lambda$ dają
$\lim_{t\to0}(U_tx-x)/t=iAx$ w normie.
Odwrotnie, gdy dla $x$ iloraz ma granicę, lemat Fatou
dla miary $\mu_x$ i
$\lim_{t\to0}|(e^{it\lambda}-1)/t|^2=\lambda^2$
dają $\int_\R\lambda^2\,\mathrm d\mu_x<\infty$.
Zatem $D(G)=D(A)$.

\textbf{Krok 2: gęstość i skośna symetria generatora.}
Niech teraz $(U_t)$ będzie dowolną silnie ciągłą grupą unitarną.
Dla $\phi\in C_c^\infty(\R)$ i $x\in H$ definiujemy
$x_\phi=\int_\R\phi(t)U_tx\,\mathrm d\lambda_1(t)$.
Jest to całka w normie przestrzeni Hilberta, zdefiniowana
przez granice sum Riemanna na zwartym nośniku $\phi$;
nie potrzeba tu ogólnej teorii całki Pettisa.
Przesunięcie zmiennej pod całką daje
$(U_sx_\phi-x_\phi)/s\to
-\int_\R\phi'(t)U_tx\,\mathrm d\lambda_1(t)$ w $H$.
W szczególności $x_\phi\in D(G)$. Wybierając
nieujemną aproksymację jedności $\phi_n$ wokół zera,
otrzymujemy $x_{\phi_n}\to x$ z silnej ciągłości grupy.
Dlatego $D(G)$ jest gęsta.
Dla $x,y\in D(G)$ różniczkowanie
$\langle U_tx,U_ty\rangle=\langle x,y\rangle$ w zerze daje
$\langle Gx,y\rangle+\langle x,Gy\rangle=0$.
Stąd operator $A=-iG$ jest symetryczny.

\textbf{Krok 3: surjektywność $A\pm iI$.}
Dla $x\in H$ kładziemy
$R_+x=\int_0^\infty e^{-t}U_tx\,\mathrm d\lambda_1(t)$,
$R_-x=\int_0^\infty e^{-t}U_{-t}x\,\mathrm d\lambda_1(t)$.
Są to granice w normie całek na odcinkach $[0,N]$;
$\|R_\pm\|\le1$, ponieważ $\|U_tx\|=\|x\|$.
Dla $s\in\R$ zmiana zmiennej daje
$U_sR_+x=e^s\int_s^\infty e^{-r}U_rx\,\mathrm d\lambda_1(r)$.
Iloraz różnicowy w zerze zbiega więc w normie do
$R_+x-x$ (z silnej ciągłości $U_rx$ w $r=0$).
Stąd $R_+x\in D(G)$ i $(I_H-G)R_+x=x$.
Podobnie
$U_sR_-x=e^{-s}\int_{-s}^\infty e^{-r}U_{-r}x\,\mathrm d\lambda_1(r)$,
zatem $GR_-x=x-R_-x$ i $(I_H+G)R_-x=x$.
W konsekwencji
$\operatorname{Ran}(I_H\pm G)=H$, co jest równoważne
$\operatorname{Ran}(A\pm iI)=H$.
Twierdzenie~\ref{thm:selfadjoint-range33} daje samosprzężoność $A$.

\textbf{Krok 4: identyfikacja i jednoznaczność grupy.}
Z własności grupowej wynika
$U_tD(G)\subseteq D(G)$ i $GU_tx=U_tGx$:
istotnie, iloraz różnicowy w punkcie $t$ jest obrazem
ilorazu różnicowego w zerze przez ograniczony operator $U_t$.
Analogiczna własność zachodzi dla $e^{itA}$
na mocy rachunku spektralnego.
Dla $x\in D(A)$ zarówno $t\mapsto U_tx$, jak i
$t\mapsto e^{itA}x$ są więc różniczkowalne i spełniają
równanie $v'(t)=iAv(t)$ oraz warunek $v(0)=x$.
Ich różnica $w$ ma stałą normę, ponieważ
$\frac{\mathrm d}{\mathrm dt}\|w(t)\|^2
=2\operatorname{Re}\langle iAw(t),w(t)\rangle=0$.
Ponieważ $w(0)=0$, mamy $w=0$.
Gęstość $D(A)$ i unitarność obu grup rozszerzają równość
na całe $H$. Jednoznaczność $A$ wynika z definicji generatora.
\end{proof}
\endgroup

\end{document}
